\documentclass[
	% Opções da classe memoir
	12pt,				        % tamanho da fonte
	openany,			        % capítulos começam em pág ímpar (insere página vazia caso preciso)
	oneside,			        % para impressão em recto e verso. Oposto a oneside
	a4paper,			        % tamanho do papel. 
	% Opções da classe abntex2
	%chapter=TITLE,		        % títulos de capítulos convertidos em letras maiúsculas
	%section=TITLE,		        % títulos de seções convertidos em letras maiúsculas
	%subsection=TITLE,	        % títulos de subseções convertidos em letras maiúsculas
	%subsubsection=TITLE,       % títulos de subsubseções convertidos em letras maiúsculas
	% Opções do pacote babel
	english,			        % idioma adicional para hifenização
	brazil				        % o último idioma é o principal do documento
	]{abntex2}

% ---
% Pacotes básicos 
% ---
\usepackage{lmodern}			% Usa a fonte Latin Modern			
\usepackage[T1]{fontenc}		% Selecao de codigos de fonte.
\usepackage[utf8]{inputenc}		% Codificacao do documento (conversão automática dos acentos)
\usepackage{indentfirst}		% Indenta o primeiro parágrafo de cada seção.
\usepackage{color}				% Controle das cores
\usepackage{graphicx}			% Inclusão de gráficos
\usepackage{microtype} 			% para melhorias de justificação
% ---
		

% ---
% Pacotes de citações
% ---
% \usepackage[alf]{abntex2cite}						% Citações padrão ABNT
\usepackage[
  alf,
  abnt-nbr10520=1988,                     % Citações em sentence case
  abnt-emphasize=bf
  ]{abntex2cite}


% ---
% Pacotes adicionais
% ---
\usepackage{pdfpages}	% Inclusão de páginas em PDF


% ---
% Informações para CAPA e FOLHA DE ROSTO
% ---
\titulo{APRENDIZAGEM DE MÁQUINA PARA PREDIÇÃO DE DESFECHOS EM SAÚDE MENTAL: \textmd{UMA INTRODUÇÃO PARA PSICÓLOGOS}}
\autor{GUILHERME AUGUSTO DIAS}
\local{Porto Alegre}
\data{2024}
\orientador{Prof. Dr. Wagner de Lara Machado}
% \coorientador{}
% \instituicao{
%   Pontifícia Universidade Católica do Rio Grande do Sul
%   \par
%   Escola de Ciências da Saúde e da Vida
%   \par
%   Curso de Psicologia}
\tipotrabalho{TCC}
\preambulo{Trabalho de Conclusão de Curso apresentado como requisito para a obtenção do grau de Bacharel em Psicologia pela Escola de Ciências da Saúde e da Vida da Pontifícia Universidade Católica do Rio Grande do Sul.}
% ---

% ---
% Configurações de aparência do PDF final
% ---

% Alterando o aspecto da cor azul
\definecolor{blue}{RGB}{41,5,195}

% Informações do PDF
\makeatletter
\hypersetup{
		pdftitle={\@title}, 
		pdfauthor={\@author},
    	pdfsubject={\imprimirpreambulo},
	    pdfcreator={LaTeX with abnTeX2},
		pdfkeywords={abnt}{latex}{abntex}{abntex2}{trabalho acadêmico}, 
		colorlinks=true,		% true: colored links; false: boxed links
    	linkcolor=black,		% color of internal links
    	citecolor=black,		% color of links to bibliography
    	filecolor=black,		% color of file links
		urlcolor=black,
		bookmarksdepth=4
}
\makeatother
% --- 


% ---
% Posiciona figuras e tabelas no topo da página quando adicionadas sozinhas
% em um página em branco. Ver https://github.com/abntex/abntex2/issues/170
\makeatletter
\setlength{\@fptop}{5pt} % Set distance from top of page to first float
\makeatother
% ---

% ---
% Possibilita criação de Quadros e Lista de quadros.
% Ver https://github.com/abntex/abntex2/issues/176
%
\newcommand{\quadroname}{Quadro}
\newcommand{\listofquadrosname}{Lista de quadros}

\newfloat[chapter]{quadro}{loq}{\quadroname}
\newlistof{listofquadros}{loq}{\listofquadrosname}
\newlistentry{quadro}{loq}{0}

% Configurações para atender às regras da ABNT
\setfloatadjustment{quadro}{\centering}
\counterwithout{quadro}{chapter}
\renewcommand{\cftquadroname}{\quadroname\space} 
\renewcommand*{\cftquadroaftersnum}{\hfill--\hfill}

\setfloatlocations{quadro}{hbtp} % Ver https://github.com/abntex/abntex2/issues/176
% ---


% --- 
% Espaçamentos entre linhas e parágrafos 
% --- 

% O tamanho do parágrafo é dado por:
\setlength{\parindent}{1.3cm}

% Controle do espaçamento entre um parágrafo e outro:
\setlength{\parskip}{0.2cm}  % tente também \onelineskip

% ---
% Compila o índice
% ---
\makeindex
% ---


% ----
% Início do documento
% ----
\begin{document}

% Seleciona o idioma do documento (conforme pacotes do babel)
\selectlanguage{brazil}

% Retira espaço extra obsoleto entre as frases.
\frenchspacing 

% ----------------------------------------------------------
% ELEMENTOS PRÉ-TEXTUAIS
% ----------------------------------------------------------
\pretextual

% ---
% CAPA
% ---
\includepdf{./05-extras/capa.pdf}
% ---


% ---
% FOLHA DE ROSTO
% (o * indica que haverá a ficha bibliográfica)
% ---
\imprimirfolhaderosto
% ---


% ---
% Ficha bibliografica
% ---

% Isto é um exemplo de Ficha Catalográfica, ou ``Dados internacionais de
% catalogação-na-publicação''. Você pode utilizar este modelo como referência. 
% Porém, provavelmente a biblioteca da sua universidade lhe fornecerá um PDF
% com a ficha catalográfica definitiva após a defesa do trabalho. Quando estiver
% com o documento, salve-o como PDF no diretório do seu projeto e substitua todo
% o conteúdo de implementação deste arquivo pelo comando abaixo:
%
% \usepackage{pdfpages}		% necessário para comando \includepdf
% \begin{fichacatalografica}
%     \includepdf{fig_ficha_catalografica.pdf}
% \end{fichacatalografica}

% \begin{fichacatalografica}
% 	\sffamily
% 	\vspace*{\fill}                             % Posição vertical
% 	\begin{center}                              % Minipage Centralizado
% 	\fbox{\begin{minipage}[c][8cm]{13.5cm}		% Largura
% 	\small
% 	\imprimirautor
% 	%Sobrenome, Nome do autor
	
% 	\hspace{0.5cm} \imprimirtitulo  / \imprimirautor. --
% 	\imprimirlocal, \imprimirdata-
	
% 	\hspace{0.5cm} \thelastpage p. : il. (algumas color.) ; 30 cm.\\
	
% 	\hspace{0.5cm} \imprimirorientadorRotulo~\imprimirorientador\\
	
% 	\hspace{0.5cm}
% 	\parbox[t]{\textwidth}{\imprimirtipotrabalho~--~\imprimirinstituicao,
% 	\imprimirdata.}\\
	
% 	\hspace{0.5cm}
% 		1. Palavra-chave1.
% 		2. Palavra-chave2.
% 		2. Palavra-chave3.
% 		I. Orientador.
% 		II. Universidade xxx.
% 		III. Faculdade de xxx.
% 		IV. Título 			
% 	\end{minipage}}
% 	\end{center}
% \end{fichacatalografica}
% ---


% ---
% FOLHA DE APROVAÇÃO
% ---

% Isto é um exemplo de Folha de aprovação, elemento obrigatório da NBR
% 14724/2011 (seção 4.2.1.3). Você pode utilizar este modelo até a aprovação
% do trabalho. Após isso, substitua todo o conteúdo deste arquivo por uma
% imagem da página assinada pela banca com o comando abaixo:
%
% \begin{folhadeaprovacao}
% \includepdf{folhadeaprovacao_final.pdf}
% \end{folhadeaprovacao}
%
\setlength{\ABNTEXsignskip}{2cm}
\setlength{\ABNTEXsignwidth}{10cm}

\begin{folhadeaprovacao}
  \begin{center}
    {\ABNTEXchapterfont\large\imprimirautor}

    \vspace*{1cm}

    \begin{center}
      \ABNTEXchapterfont\bfseries\Large\imprimirtitulo
    \end{center}

    \vspace*{1cm}
    
    \hspace{.45\textwidth}
    \begin{minipage}{.5\textwidth}
        \imprimirpreambulo
    \end{minipage}%
   \end{center}

   \vspace*{1cm}
        
   \begin{center}
   Aprovado em 17 de julho de 2024.

   \vspace*{1cm}

   Banca examinadora:
   \end{center}

   \vspace*{-1cm}

   \assinatura{Prof. Ma. Caroline Santa Maria Rodrigues -- PUCRS}
   \assinatura{Prof. Dra. Clarissa Pinto Pizarro de Freitas -- PUCRS}
   \assinatura{\imprimirorientador -- PUCRS}
\end{folhadeaprovacao}
% ---


% ---
% Agradecimentos
% ---
\begin{agradecimentos}
Valeu!
\end{agradecimentos}
% ---


% ---
% RESUMOS
% ---
% Resumo em português
\setlength{\absparsep}{18pt} % ajusta o espaçamento dos parágrafos do resumo
\begin{resumo}
Desenvolvimentos na área de aprendizagem de máquina representam uma série de oportunidades para diversas áreas de pesquisa e de atuação profissional.
O uso de modelos de aprendizagem de máquina na predição de desfecho para tratamentos em saúde mental tem sido investigado e apresenta resultados promissores.
A aplicação desse tipo de modelo possibilitaria, por exemplo, maior assertividade na personalização de tratamentos psicoterápicos a nível individual.
Apesar dos benefícios potenciais, o tema ainda é pouco discutido nos espaços de pesquisa, formação e atuação em psicologia clínica. Este trabalho tem
por objetivo introduzir o tema de aprendizagem de máquina a psicólogos. Busca-se apresentar, de maneira acessível, os conceitos e terminologia básicos
em aprendizagem de máquina, proporcionando um panorama geral da área. Para complementar a exposição conceitual, é construído um exemplo a partir de dados
de um estudo sobre eficácia de intervenções em psicologia positiva para depressão via internet. Por fim, discutem-se algumas implicações do uso desse tipo
de tecnologia para a pesquisa e para a prática em psicoterapia.

 \textbf{Palavras-chave}: aprendizagem de máquina, psicoterapia, psicologia positiva, depressão.
\end{resumo}

% Resumo em inglês
\begin{resumo}[Abstract]
 \begin{otherlanguage*}{english}
   Advances in machine learning have opened up a multitude of opportunities for various research and professional fields. The use of machine learning models to
   predict treatment outcomes in mental health has been investigated showing promising results. The application of such models would allow greater assertiveness
   in tailoring psychotherapeutic treatments at the individual level. Despite its potential benefits, the topic remains unexplored in the fields of research, training
   and practice in clinical psychology. This work aims to introduce the topic of machine learning to psychologists. It presents the basic machine learning concepts and
   terminology in an accessible manner, providing a comprehensive overview of the field. To complement the conceptual discussion, an example is constructed using data
   from a study on the effectiveness of internet-based positive psychology interventions for depression. Finally, implications of using this type of technology in
   psychotherapy research and practice are discussed.
   \vspace{\onelineskip}
 
   \noindent 
   \textbf{Keywords}: machine learning, psychotherapy, positive psychology, depression.
 \end{otherlanguage*}
\end{resumo}


% ---
% LISTA DE ILUSTRAÇÕES
% ---
\pdfbookmark[0]{\listfigurename}{lof}
\listoffigures*
\cleardoublepage
% ---


% ---
% LISTA DE QUADROS
% ---
% \pdfbookmark[0]{\listofquadrosname}{loq}
% \listofquadros*
% \cleardoublepage
% ---


% ---
% LISTA DE TABELAS
% ---
% \pdfbookmark[0]{\listtablename}{lot}
% \listoftables*
% \cleardoublepage
% ---


% ---
% LISTA DE ABREVIATURAS E SIGLAS
% ---
% \begin{siglas}
%   \item[ABNT] Associação Brasileira de Normas Técnicas
%   \item[abnTeX] ABsurdas Normas para TeX
% \end{siglas}
% ---


% ---
% LISTA DE SÍMBOLOS
% ---
% \begin{simbolos}
%   \item[$ \Gamma $] Letra grega Gama
%   \item[$ \Lambda $] Lambda
%   \item[$ \zeta $] Letra grega minúscula zeta
%   \item[$ \in $] Pertence
% \end{simbolos}
% ---


% ---
% Sumário
% ---
\pdfbookmark[0]{\contentsname}{toc}
\tableofcontents*
\cleardoublepage
% ---


% ----------------------------------------------------------
% ELEMENTOS TEXTUAIS
% ----------------------------------------------------------
\textual

\documentclass[
	% Opções da classe memoir
	12pt,				        % tamanho da fonte
	openany,			        % capítulos começam em pág ímpar (insere página vazia caso preciso)
	oneside,			        % para impressão em recto e verso. Oposto a oneside
	a4paper,			        % tamanho do papel. 
	% Opções da classe abntex2
	%chapter=TITLE,		        % títulos de capítulos convertidos em letras maiúsculas
	%section=TITLE,		        % títulos de seções convertidos em letras maiúsculas
	%subsection=TITLE,	        % títulos de subseções convertidos em letras maiúsculas
	%subsubsection=TITLE,       % títulos de subsubseções convertidos em letras maiúsculas
	% Opções do pacote babel
	english,			        % idioma adicional para hifenização
	brazil				        % o último idioma é o principal do documento
	]{abntex2}

% ---
% Pacotes básicos 
% ---
\usepackage{lmodern}			% Usa a fonte Latin Modern			
\usepackage[T1]{fontenc}		% Selecao de codigos de fonte.
\usepackage[utf8]{inputenc}		% Codificacao do documento (conversão automática dos acentos)
\usepackage{indentfirst}		% Indenta o primeiro parágrafo de cada seção.
\usepackage{color}				% Controle das cores
\usepackage{graphicx}			% Inclusão de gráficos
\usepackage{microtype} 			% para melhorias de justificação
% ---
		

% ---
% Pacotes de citações
% ---
% \usepackage[alf]{abntex2cite}						% Citações padrão ABNT
\usepackage[
  alf,
  abnt-nbr10520=1988,                     % Citações em sentence case
  abnt-emphasize=bf
  ]{abntex2cite}


% ---
% Pacotes adicionais
% ---
\usepackage{pdfpages}	% Inclusão de páginas em PDF


% ---
% Informações para CAPA e FOLHA DE ROSTO
% ---
\titulo{APRENDIZAGEM DE MÁQUINA PARA PREDIÇÃO DE DESFECHOS EM SAÚDE MENTAL: \textmd{UMA INTRODUÇÃO PARA PSICÓLOGOS}}
\autor{GUILHERME AUGUSTO DIAS}
\local{Porto Alegre}
\data{2024}
\orientador{Prof. Dr. Wagner de Lara Machado}
% \coorientador{}
% \instituicao{
%   Pontifícia Universidade Católica do Rio Grande do Sul
%   \par
%   Escola de Ciências da Saúde e da Vida
%   \par
%   Curso de Psicologia}
\tipotrabalho{TCC}
\preambulo{Trabalho de Conclusão de Curso apresentado como requisito para a obtenção do grau de Bacharel em Psicologia pela Escola de Ciências da Saúde e da Vida da Pontifícia Universidade Católica do Rio Grande do Sul.}
% ---

% ---
% Configurações de aparência do PDF final
% ---

% Alterando o aspecto da cor azul
\definecolor{blue}{RGB}{41,5,195}

% Informações do PDF
\makeatletter
\hypersetup{
		pdftitle={Aprendizagem de máquina para predição de desfechos em saúde mental: uma introdução para psicólogos}, 
		pdfauthor={Guilherme Augusto Dias},
    	pdfsubject={\imprimirpreambulo},
	    pdfcreator={LaTeX with abnTeX2},
		pdfkeywords={aprendizagem de máquina}{psicoterapia}{psicologia positiva}{depressão}, 
		colorlinks=true,		% true: colored links; false: boxed links
    	linkcolor=black,		% color of internal links
    	citecolor=black,		% color of links to bibliography
    	filecolor=black,		% color of file links
		urlcolor=black,
		bookmarksdepth=4
}
\makeatother
% --- 


% ---
% Posiciona figuras e tabelas no topo da página quando adicionadas sozinhas
% em um página em branco. Ver https://github.com/abntex/abntex2/issues/170
\makeatletter
\setlength{\@fptop}{5pt} % Set distance from top of page to first float
\makeatother
% ---

% ---
% Possibilita criação de Quadros e Lista de quadros.
% Ver https://github.com/abntex/abntex2/issues/176
%
\newcommand{\quadroname}{Quadro}
\newcommand{\listofquadrosname}{Lista de quadros}

\newfloat[chapter]{quadro}{loq}{\quadroname}
\newlistof{listofquadros}{loq}{\listofquadrosname}
\newlistentry{quadro}{loq}{0}

% Configurações para atender às regras da ABNT
\setfloatadjustment{quadro}{\centering}
\counterwithout{quadro}{chapter}
\renewcommand{\cftquadroname}{\quadroname\space} 
\renewcommand*{\cftquadroaftersnum}{\hfill--\hfill}

\setfloatlocations{quadro}{hbtp} % Ver https://github.com/abntex/abntex2/issues/176
% ---


% --- 
% Espaçamentos entre linhas e parágrafos 
% --- 

% O tamanho do parágrafo é dado por:
\setlength{\parindent}{1.3cm}

% Controle do espaçamento entre um parágrafo e outro:
\setlength{\parskip}{0.2cm}  % tente também \onelineskip

% ---
% Compila o índice
% ---
\makeindex
% ---


% ----
% Início do documento
% ----
\begin{document}

% Seleciona o idioma do documento (conforme pacotes do babel)
\selectlanguage{brazil}

% Retira espaço extra obsoleto entre as frases.
\frenchspacing 

% ----------------------------------------------------------
% ELEMENTOS PRÉ-TEXTUAIS
% ----------------------------------------------------------
\pretextual

% ---
% CAPA
% ---
\includepdf{./05-extras/capa.pdf}
% ---


% ---
% FOLHA DE ROSTO
% (o * indica que haverá a ficha bibliográfica)
% ---
\imprimirfolhaderosto
% ---


% ---
% Ficha bibliografica
% ---

% Isto é um exemplo de Ficha Catalográfica, ou ``Dados internacionais de
% catalogação-na-publicação''. Você pode utilizar este modelo como referência. 
% Porém, provavelmente a biblioteca da sua universidade lhe fornecerá um PDF
% com a ficha catalográfica definitiva após a defesa do trabalho. Quando estiver
% com o documento, salve-o como PDF no diretório do seu projeto e substitua todo
% o conteúdo de implementação deste arquivo pelo comando abaixo:
%
% \usepackage{pdfpages}		% necessário para comando \includepdf
% \begin{fichacatalografica}
%     \includepdf{fig_ficha_catalografica.pdf}
% \end{fichacatalografica}

% \begin{fichacatalografica}
% 	\sffamily
% 	\vspace*{\fill}                             % Posição vertical
% 	\begin{center}                              % Minipage Centralizado
% 	\fbox{\begin{minipage}[c][8cm]{13.5cm}		% Largura
% 	\small
% 	\imprimirautor
% 	%Sobrenome, Nome do autor
	
% 	\hspace{0.5cm} \imprimirtitulo  / \imprimirautor. --
% 	\imprimirlocal, \imprimirdata-
	
% 	\hspace{0.5cm} \thelastpage p. : il. (algumas color.) ; 30 cm.\\
	
% 	\hspace{0.5cm} \imprimirorientadorRotulo~\imprimirorientador\\
	
% 	\hspace{0.5cm}
% 	\parbox[t]{\textwidth}{\imprimirtipotrabalho~--~\imprimirinstituicao,
% 	\imprimirdata.}\\
	
% 	\hspace{0.5cm}
% 		1. Palavra-chave1.
% 		2. Palavra-chave2.
% 		2. Palavra-chave3.
% 		I. Orientador.
% 		II. Universidade xxx.
% 		III. Faculdade de xxx.
% 		IV. Título 			
% 	\end{minipage}}
% 	\end{center}
% \end{fichacatalografica}
% ---


% ---
% FOLHA DE APROVAÇÃO
% ---

% Isto é um exemplo de Folha de aprovação, elemento obrigatório da NBR
% 14724/2011 (seção 4.2.1.3). Você pode utilizar este modelo até a aprovação
% do trabalho. Após isso, substitua todo o conteúdo deste arquivo por uma
% imagem da página assinada pela banca com o comando abaixo:
%
% \begin{folhadeaprovacao}
% \includepdf{folhadeaprovacao_final.pdf}
% \end{folhadeaprovacao}
%
\setlength{\ABNTEXsignskip}{2cm}
\setlength{\ABNTEXsignwidth}{10cm}

\begin{folhadeaprovacao}
  \begin{center}
    {\ABNTEXchapterfont\large\imprimirautor}

    \vspace*{1cm}

    \begin{center}
      \ABNTEXchapterfont\bfseries\Large\imprimirtitulo
    \end{center}

    \vspace*{1cm}
    
    \hspace{.45\textwidth}
    \begin{minipage}{.5\textwidth}
        \imprimirpreambulo
    \end{minipage}%
   \end{center}

   \vspace*{1cm}
        
   \begin{center}
   Aprovado em 17 de julho de 2024.

   \vspace*{1cm}

   Banca examinadora:
   \end{center}

   \vspace*{-1cm}

   \assinatura{Prof. Ma. Caroline Santa Maria Rodrigues -- PUCRS}
   \assinatura{Prof. Dra. Clarissa Pinto Pizarro de Freitas -- PUCRS}
   \assinatura{\imprimirorientador -- PUCRS}
\end{folhadeaprovacao}
% ---


% ---
% Agradecimentos
% ---
\begin{agradecimentos}
Agradeço aos meus pais, Ademilson José Dias e Irinês Porfírio Dias, pelo apoio incondicional que me deram ao longo da vida,
incentivando mesmo as ideias questionáveis como aquela de fazer uma segunda graduação.

A minha irmã, Maria Eduarda Dias, cuja trajetória acadêmica é fonte de inspiração.

A minha namorada, Nathalia Fallavena Ceratti, que acompanhou a esta jornada de perto, oferecendo apoio e compreensão que
fizeram toda a diferença.

Aos meus amigos, Guilherme da Silva Petró, Norton Neves e Uilinton Luciano Marques da Silva, que foram compreensíveis nos
finais de semana em que não pude jogar RPG.

Ao professor Wagner de Lara Machado pela orientação e apoio no desenvolvimento deste trabalho.
\end{agradecimentos}
% ---


% ---
% RESUMOS
% ---
% Resumo em português
\setlength{\absparsep}{18pt} % ajusta o espaçamento dos parágrafos do resumo
\begin{resumo}
Desenvolvimentos na área de aprendizagem de máquina representam uma série de oportunidades para diversas áreas de pesquisa e de atuação profissional.
O uso de modelos de aprendizagem de máquina na predição de desfecho para tratamentos em saúde mental tem sido investigado e apresenta resultados promissores.
A aplicação desse tipo de modelo possibilitaria, por exemplo, maior assertividade na personalização de tratamentos psicoterápicos a nível individual.
Apesar dos benefícios potenciais, o tema ainda é pouco discutido nos espaços de pesquisa, formação e atuação em psicologia clínica. Este trabalho tem
por objetivo introduzir o tema de aprendizagem de máquina a psicólogos. Busca-se apresentar, de maneira acessível, os conceitos e terminologia básicos
em aprendizagem de máquina, proporcionando um panorama geral da área. Para complementar a exposição conceitual, é construído um exemplo a partir de dados
de um estudo sobre eficácia de intervenções em psicologia positiva para depressão via internet. Por fim, discutem-se algumas implicações do uso desse tipo
de tecnologia para a pesquisa e para a prática em psicoterapia.

 \textbf{Palavras-chave}: aprendizagem de máquina, psicoterapia, psicologia positiva, depressão.
\end{resumo}

% Resumo em inglês
\begin{resumo}[Abstract]
 \begin{otherlanguage*}{english}
   Advances in machine learning have opened up a multitude of opportunities for various research and professional fields. The use of machine learning models to
   predict treatment outcomes in mental health has been investigated showing promising results. The application of such models would allow greater assertiveness
   in tailoring psychotherapeutic treatments at the individual level. Despite its potential benefits, the topic remains unexplored in the fields of research, training
   and practice in clinical psychology. This work aims to introduce the topic of machine learning to psychologists. It presents the basic machine learning concepts and
   terminology in an accessible manner, providing a comprehensive overview of the field. To complement the conceptual discussion, an example is constructed using data
   from a study on the effectiveness of internet-based positive psychology interventions for depression. Finally, implications of using this type of technology in
   psychotherapy research and practice are discussed.
   \vspace{\onelineskip}
 
   \noindent 
   \textbf{Keywords}: machine learning, psychotherapy, positive psychology, depression.
 \end{otherlanguage*}
\end{resumo}


% ---
% LISTA DE ILUSTRAÇÕES
% ---
\pdfbookmark[0]{\listfigurename}{lof}
\listoffigures*
\cleardoublepage
% ---


% ---
% LISTA DE QUADROS
% ---
% \pdfbookmark[0]{\listofquadrosname}{loq}
% \listofquadros*
% \cleardoublepage
% ---


% ---
% LISTA DE TABELAS
% ---
% \pdfbookmark[0]{\listtablename}{lot}
% \listoftables*
% \cleardoublepage
% ---


% ---
% LISTA DE ABREVIATURAS E SIGLAS
% ---
% \begin{siglas}
%   \item[ABNT] Associação Brasileira de Normas Técnicas
%   \item[abnTeX] ABsurdas Normas para TeX
% \end{siglas}
% ---


% ---
% LISTA DE SÍMBOLOS
% ---
% \begin{simbolos}
%   \item[$ \Gamma $] Letra grega Gama
%   \item[$ \Lambda $] Lambda
%   \item[$ \zeta $] Letra grega minúscula zeta
%   \item[$ \in $] Pertence
% \end{simbolos}
% ---


% ---
% Sumário
% ---
\pdfbookmark[0]{\contentsname}{toc}
\tableofcontents*
\cleardoublepage
% ---


% ----------------------------------------------------------
% ELEMENTOS TEXTUAIS
% ----------------------------------------------------------
\textual

\documentclass[
	% Opções da classe memoir
	12pt,				        % tamanho da fonte
	openany,			        % capítulos começam em pág ímpar (insere página vazia caso preciso)
	oneside,			        % para impressão em recto e verso. Oposto a oneside
	a4paper,			        % tamanho do papel. 
	% Opções da classe abntex2
	%chapter=TITLE,		        % títulos de capítulos convertidos em letras maiúsculas
	%section=TITLE,		        % títulos de seções convertidos em letras maiúsculas
	%subsection=TITLE,	        % títulos de subseções convertidos em letras maiúsculas
	%subsubsection=TITLE,       % títulos de subsubseções convertidos em letras maiúsculas
	% Opções do pacote babel
	english,			        % idioma adicional para hifenização
	brazil				        % o último idioma é o principal do documento
	]{abntex2}

% ---
% Pacotes básicos 
% ---
\usepackage{lmodern}			% Usa a fonte Latin Modern			
\usepackage[T1]{fontenc}		% Selecao de codigos de fonte.
\usepackage[utf8]{inputenc}		% Codificacao do documento (conversão automática dos acentos)
\usepackage{indentfirst}		% Indenta o primeiro parágrafo de cada seção.
\usepackage{color}				% Controle das cores
\usepackage{graphicx}			% Inclusão de gráficos
\usepackage{microtype} 			% para melhorias de justificação
% ---
		

% ---
% Pacotes de citações
% ---
% \usepackage[alf]{abntex2cite}						% Citações padrão ABNT
\usepackage[
  alf,
  abnt-nbr10520=1988,                     % Citações em sentence case
  abnt-emphasize=bf
  ]{abntex2cite}


% ---
% Pacotes adicionais
% ---
\usepackage{pdfpages}	% Inclusão de páginas em PDF


% ---
% Informações para CAPA e FOLHA DE ROSTO
% ---
\titulo{APRENDIZAGEM DE MÁQUINA PARA PREDIÇÃO DE DESFECHOS EM SAÚDE MENTAL: \textmd{UMA INTRODUÇÃO PARA PSICÓLOGOS}}
\autor{GUILHERME AUGUSTO DIAS}
\local{Porto Alegre}
\data{2024}
\orientador{Prof. Dr. Wagner de Lara Machado}
% \coorientador{}
% \instituicao{
%   Pontifícia Universidade Católica do Rio Grande do Sul
%   \par
%   Escola de Ciências da Saúde e da Vida
%   \par
%   Curso de Psicologia}
\tipotrabalho{TCC}
\preambulo{Trabalho de Conclusão de Curso apresentado como requisito para a obtenção do grau de Bacharel em Psicologia pela Escola de Ciências da Saúde e da Vida da Pontifícia Universidade Católica do Rio Grande do Sul.}
% ---

% ---
% Configurações de aparência do PDF final
% ---

% Alterando o aspecto da cor azul
\definecolor{blue}{RGB}{41,5,195}

% Informações do PDF
\makeatletter
\hypersetup{
		pdftitle={Aprendizagem de máquina para predição de desfechos em saúde mental: uma introdução para psicólogos}, 
		pdfauthor={Guilherme Augusto Dias},
    	pdfsubject={\imprimirpreambulo},
	    pdfcreator={LaTeX with abnTeX2},
		pdfkeywords={aprendizagem de máquina}{psicoterapia}{psicologia positiva}{depressão}, 
		colorlinks=true,		% true: colored links; false: boxed links
    	linkcolor=black,		% color of internal links
    	citecolor=black,		% color of links to bibliography
    	filecolor=black,		% color of file links
		urlcolor=black,
		bookmarksdepth=4
}
\makeatother
% --- 


% ---
% Posiciona figuras e tabelas no topo da página quando adicionadas sozinhas
% em um página em branco. Ver https://github.com/abntex/abntex2/issues/170
\makeatletter
\setlength{\@fptop}{5pt} % Set distance from top of page to first float
\makeatother
% ---

% ---
% Possibilita criação de Quadros e Lista de quadros.
% Ver https://github.com/abntex/abntex2/issues/176
%
\newcommand{\quadroname}{Quadro}
\newcommand{\listofquadrosname}{Lista de quadros}

\newfloat[chapter]{quadro}{loq}{\quadroname}
\newlistof{listofquadros}{loq}{\listofquadrosname}
\newlistentry{quadro}{loq}{0}

% Configurações para atender às regras da ABNT
\setfloatadjustment{quadro}{\centering}
\counterwithout{quadro}{chapter}
\renewcommand{\cftquadroname}{\quadroname\space} 
\renewcommand*{\cftquadroaftersnum}{\hfill--\hfill}

\setfloatlocations{quadro}{hbtp} % Ver https://github.com/abntex/abntex2/issues/176
% ---


% --- 
% Espaçamentos entre linhas e parágrafos 
% --- 

% O tamanho do parágrafo é dado por:
\setlength{\parindent}{1.3cm}

% Controle do espaçamento entre um parágrafo e outro:
\setlength{\parskip}{0.2cm}  % tente também \onelineskip

% ---
% Compila o índice
% ---
\makeindex
% ---


% ----
% Início do documento
% ----
\begin{document}

% Seleciona o idioma do documento (conforme pacotes do babel)
\selectlanguage{brazil}

% Retira espaço extra obsoleto entre as frases.
\frenchspacing 

% ----------------------------------------------------------
% ELEMENTOS PRÉ-TEXTUAIS
% ----------------------------------------------------------
\pretextual

% ---
% CAPA
% ---
\includepdf{./05-extras/capa.pdf}
% ---


% ---
% FOLHA DE ROSTO
% (o * indica que haverá a ficha bibliográfica)
% ---
\imprimirfolhaderosto
% ---


% ---
% Ficha bibliografica
% ---

% Isto é um exemplo de Ficha Catalográfica, ou ``Dados internacionais de
% catalogação-na-publicação''. Você pode utilizar este modelo como referência. 
% Porém, provavelmente a biblioteca da sua universidade lhe fornecerá um PDF
% com a ficha catalográfica definitiva após a defesa do trabalho. Quando estiver
% com o documento, salve-o como PDF no diretório do seu projeto e substitua todo
% o conteúdo de implementação deste arquivo pelo comando abaixo:
%
% \usepackage{pdfpages}		% necessário para comando \includepdf
% \begin{fichacatalografica}
%     \includepdf{fig_ficha_catalografica.pdf}
% \end{fichacatalografica}

% \begin{fichacatalografica}
% 	\sffamily
% 	\vspace*{\fill}                             % Posição vertical
% 	\begin{center}                              % Minipage Centralizado
% 	\fbox{\begin{minipage}[c][8cm]{13.5cm}		% Largura
% 	\small
% 	\imprimirautor
% 	%Sobrenome, Nome do autor
	
% 	\hspace{0.5cm} \imprimirtitulo  / \imprimirautor. --
% 	\imprimirlocal, \imprimirdata-
	
% 	\hspace{0.5cm} \thelastpage p. : il. (algumas color.) ; 30 cm.\\
	
% 	\hspace{0.5cm} \imprimirorientadorRotulo~\imprimirorientador\\
	
% 	\hspace{0.5cm}
% 	\parbox[t]{\textwidth}{\imprimirtipotrabalho~--~\imprimirinstituicao,
% 	\imprimirdata.}\\
	
% 	\hspace{0.5cm}
% 		1. Palavra-chave1.
% 		2. Palavra-chave2.
% 		2. Palavra-chave3.
% 		I. Orientador.
% 		II. Universidade xxx.
% 		III. Faculdade de xxx.
% 		IV. Título 			
% 	\end{minipage}}
% 	\end{center}
% \end{fichacatalografica}
% ---


% ---
% FOLHA DE APROVAÇÃO
% ---

% Isto é um exemplo de Folha de aprovação, elemento obrigatório da NBR
% 14724/2011 (seção 4.2.1.3). Você pode utilizar este modelo até a aprovação
% do trabalho. Após isso, substitua todo o conteúdo deste arquivo por uma
% imagem da página assinada pela banca com o comando abaixo:
%
% \begin{folhadeaprovacao}
% \includepdf{folhadeaprovacao_final.pdf}
% \end{folhadeaprovacao}
%
\setlength{\ABNTEXsignskip}{2cm}
\setlength{\ABNTEXsignwidth}{10cm}

\begin{folhadeaprovacao}
  \begin{center}
    {\ABNTEXchapterfont\large\imprimirautor}

    \vspace*{1cm}

    \begin{center}
      \ABNTEXchapterfont\bfseries\Large\imprimirtitulo
    \end{center}

    \vspace*{1cm}
    
    \hspace{.45\textwidth}
    \begin{minipage}{.5\textwidth}
        \imprimirpreambulo
    \end{minipage}%
   \end{center}

   \vspace*{1cm}
        
   \begin{center}
   Aprovado em 17 de julho de 2024.

   \vspace*{1cm}

   Banca examinadora:
   \end{center}

   \vspace*{-1cm}

   \assinatura{Prof. Ma. Caroline Santa Maria Rodrigues -- PUCRS}
   \assinatura{Prof. Dra. Clarissa Pinto Pizarro de Freitas -- PUCRS}
   \assinatura{\imprimirorientador -- PUCRS}
\end{folhadeaprovacao}
% ---


% ---
% Agradecimentos
% ---
\begin{agradecimentos}
Agradeço aos meus pais, Ademilson José Dias e Irinês Porfírio Dias, pelo apoio incondicional que me deram ao longo da vida,
incentivando mesmo as ideias questionáveis como aquela de fazer uma segunda graduação.

A minha irmã, Maria Eduarda Dias, cuja trajetória acadêmica é fonte de inspiração.

A minha namorada, Nathalia Fallavena Ceratti, que acompanhou a esta jornada de perto, oferecendo apoio e compreensão que
fizeram toda a diferença.

Aos meus amigos, Guilherme da Silva Petró, Norton Neves e Uilinton Luciano Marques da Silva, que foram compreensíveis nos
finais de semana em que não pude jogar RPG.

Ao professor Wagner de Lara Machado pela orientação e apoio no desenvolvimento deste trabalho.
\end{agradecimentos}
% ---


% ---
% RESUMOS
% ---
% Resumo em português
\setlength{\absparsep}{18pt} % ajusta o espaçamento dos parágrafos do resumo
\begin{resumo}
Desenvolvimentos na área de aprendizagem de máquina representam uma série de oportunidades para diversas áreas de pesquisa e de atuação profissional.
O uso de modelos de aprendizagem de máquina na predição de desfecho para tratamentos em saúde mental tem sido investigado e apresenta resultados promissores.
A aplicação desse tipo de modelo possibilitaria, por exemplo, maior assertividade na personalização de tratamentos psicoterápicos a nível individual.
Apesar dos benefícios potenciais, o tema ainda é pouco discutido nos espaços de pesquisa, formação e atuação em psicologia clínica. Este trabalho tem
por objetivo introduzir o tema de aprendizagem de máquina a psicólogos. Busca-se apresentar, de maneira acessível, os conceitos e terminologia básicos
em aprendizagem de máquina, proporcionando um panorama geral da área. Para complementar a exposição conceitual, é construído um exemplo a partir de dados
de um estudo sobre eficácia de intervenções em psicologia positiva para depressão via internet. Por fim, discutem-se algumas implicações do uso desse tipo
de tecnologia para a pesquisa e para a prática em psicoterapia.

 \textbf{Palavras-chave}: aprendizagem de máquina, psicoterapia, psicologia positiva, depressão.
\end{resumo}

% Resumo em inglês
\begin{resumo}[Abstract]
 \begin{otherlanguage*}{english}
   Advances in machine learning have opened up a multitude of opportunities for various research and professional fields. The use of machine learning models to
   predict treatment outcomes in mental health has been investigated showing promising results. The application of such models would allow greater assertiveness
   in tailoring psychotherapeutic treatments at the individual level. Despite its potential benefits, the topic remains unexplored in the fields of research, training
   and practice in clinical psychology. This work aims to introduce the topic of machine learning to psychologists. It presents the basic machine learning concepts and
   terminology in an accessible manner, providing a comprehensive overview of the field. To complement the conceptual discussion, an example is constructed using data
   from a study on the effectiveness of internet-based positive psychology interventions for depression. Finally, implications of using this type of technology in
   psychotherapy research and practice are discussed.
   \vspace{\onelineskip}
 
   \noindent 
   \textbf{Keywords}: machine learning, psychotherapy, positive psychology, depression.
 \end{otherlanguage*}
\end{resumo}


% ---
% LISTA DE ILUSTRAÇÕES
% ---
\pdfbookmark[0]{\listfigurename}{lof}
\listoffigures*
\cleardoublepage
% ---


% ---
% LISTA DE QUADROS
% ---
% \pdfbookmark[0]{\listofquadrosname}{loq}
% \listofquadros*
% \cleardoublepage
% ---


% ---
% LISTA DE TABELAS
% ---
% \pdfbookmark[0]{\listtablename}{lot}
% \listoftables*
% \cleardoublepage
% ---


% ---
% LISTA DE ABREVIATURAS E SIGLAS
% ---
% \begin{siglas}
%   \item[ABNT] Associação Brasileira de Normas Técnicas
%   \item[abnTeX] ABsurdas Normas para TeX
% \end{siglas}
% ---


% ---
% LISTA DE SÍMBOLOS
% ---
% \begin{simbolos}
%   \item[$ \Gamma $] Letra grega Gama
%   \item[$ \Lambda $] Lambda
%   \item[$ \zeta $] Letra grega minúscula zeta
%   \item[$ \in $] Pertence
% \end{simbolos}
% ---


% ---
% Sumário
% ---
\pdfbookmark[0]{\contentsname}{toc}
\tableofcontents*
\cleardoublepage
% ---


% ----------------------------------------------------------
% ELEMENTOS TEXTUAIS
% ----------------------------------------------------------
\textual

\documentclass[
	% Opções da classe memoir
	12pt,				        % tamanho da fonte
	openany,			        % capítulos começam em pág ímpar (insere página vazia caso preciso)
	oneside,			        % para impressão em recto e verso. Oposto a oneside
	a4paper,			        % tamanho do papel. 
	% Opções da classe abntex2
	%chapter=TITLE,		        % títulos de capítulos convertidos em letras maiúsculas
	%section=TITLE,		        % títulos de seções convertidos em letras maiúsculas
	%subsection=TITLE,	        % títulos de subseções convertidos em letras maiúsculas
	%subsubsection=TITLE,       % títulos de subsubseções convertidos em letras maiúsculas
	% Opções do pacote babel
	english,			        % idioma adicional para hifenização
	brazil				        % o último idioma é o principal do documento
	]{abntex2}

% ---
% Pacotes básicos 
% ---
\usepackage{lmodern}			% Usa a fonte Latin Modern			
\usepackage[T1]{fontenc}		% Selecao de codigos de fonte.
\usepackage[utf8]{inputenc}		% Codificacao do documento (conversão automática dos acentos)
\usepackage{indentfirst}		% Indenta o primeiro parágrafo de cada seção.
\usepackage{color}				% Controle das cores
\usepackage{graphicx}			% Inclusão de gráficos
\usepackage{microtype} 			% para melhorias de justificação
% ---
		

% ---
% Pacotes de citações
% ---
% \usepackage[alf]{abntex2cite}						% Citações padrão ABNT
\usepackage[
  alf,
  abnt-nbr10520=1988,                     % Citações em sentence case
  abnt-emphasize=bf
  ]{abntex2cite}


% ---
% Pacotes adicionais
% ---
\usepackage{pdfpages}	% Inclusão de páginas em PDF


% ---
% Informações para CAPA e FOLHA DE ROSTO
% ---
\titulo{APRENDIZAGEM DE MÁQUINA PARA PREDIÇÃO DE DESFECHOS EM SAÚDE MENTAL: \textmd{UMA INTRODUÇÃO PARA PSICÓLOGOS}}
\autor{GUILHERME AUGUSTO DIAS}
\local{Porto Alegre}
\data{2024}
\orientador{Prof. Dr. Wagner de Lara Machado}
% \coorientador{}
% \instituicao{
%   Pontifícia Universidade Católica do Rio Grande do Sul
%   \par
%   Escola de Ciências da Saúde e da Vida
%   \par
%   Curso de Psicologia}
\tipotrabalho{TCC}
\preambulo{Trabalho de Conclusão de Curso apresentado como requisito para a obtenção do grau de Bacharel em Psicologia pela Escola de Ciências da Saúde e da Vida da Pontifícia Universidade Católica do Rio Grande do Sul.}
% ---

% ---
% Configurações de aparência do PDF final
% ---

% Alterando o aspecto da cor azul
\definecolor{blue}{RGB}{41,5,195}

% Informações do PDF
\makeatletter
\hypersetup{
		pdftitle={Aprendizagem de máquina para predição de desfechos em saúde mental: uma introdução para psicólogos}, 
		pdfauthor={Guilherme Augusto Dias},
    	pdfsubject={\imprimirpreambulo},
	    pdfcreator={LaTeX with abnTeX2},
		pdfkeywords={aprendizagem de máquina}{psicoterapia}{psicologia positiva}{depressão}, 
		colorlinks=true,		% true: colored links; false: boxed links
    	linkcolor=black,		% color of internal links
    	citecolor=black,		% color of links to bibliography
    	filecolor=black,		% color of file links
		urlcolor=black,
		bookmarksdepth=4
}
\makeatother
% --- 


% ---
% Posiciona figuras e tabelas no topo da página quando adicionadas sozinhas
% em um página em branco. Ver https://github.com/abntex/abntex2/issues/170
\makeatletter
\setlength{\@fptop}{5pt} % Set distance from top of page to first float
\makeatother
% ---

% ---
% Possibilita criação de Quadros e Lista de quadros.
% Ver https://github.com/abntex/abntex2/issues/176
%
\newcommand{\quadroname}{Quadro}
\newcommand{\listofquadrosname}{Lista de quadros}

\newfloat[chapter]{quadro}{loq}{\quadroname}
\newlistof{listofquadros}{loq}{\listofquadrosname}
\newlistentry{quadro}{loq}{0}

% Configurações para atender às regras da ABNT
\setfloatadjustment{quadro}{\centering}
\counterwithout{quadro}{chapter}
\renewcommand{\cftquadroname}{\quadroname\space} 
\renewcommand*{\cftquadroaftersnum}{\hfill--\hfill}

\setfloatlocations{quadro}{hbtp} % Ver https://github.com/abntex/abntex2/issues/176
% ---


% --- 
% Espaçamentos entre linhas e parágrafos 
% --- 

% O tamanho do parágrafo é dado por:
\setlength{\parindent}{1.3cm}

% Controle do espaçamento entre um parágrafo e outro:
\setlength{\parskip}{0.2cm}  % tente também \onelineskip

% ---
% Compila o índice
% ---
\makeindex
% ---


% ----
% Início do documento
% ----
\begin{document}

% Seleciona o idioma do documento (conforme pacotes do babel)
\selectlanguage{brazil}

% Retira espaço extra obsoleto entre as frases.
\frenchspacing 

% ----------------------------------------------------------
% ELEMENTOS PRÉ-TEXTUAIS
% ----------------------------------------------------------
\pretextual

% ---
% CAPA
% ---
\includepdf{./05-extras/capa.pdf}
% ---


% ---
% FOLHA DE ROSTO
% (o * indica que haverá a ficha bibliográfica)
% ---
\imprimirfolhaderosto
% ---


% ---
% Ficha bibliografica
% ---

% Isto é um exemplo de Ficha Catalográfica, ou ``Dados internacionais de
% catalogação-na-publicação''. Você pode utilizar este modelo como referência. 
% Porém, provavelmente a biblioteca da sua universidade lhe fornecerá um PDF
% com a ficha catalográfica definitiva após a defesa do trabalho. Quando estiver
% com o documento, salve-o como PDF no diretório do seu projeto e substitua todo
% o conteúdo de implementação deste arquivo pelo comando abaixo:
%
% \usepackage{pdfpages}		% necessário para comando \includepdf
% \begin{fichacatalografica}
%     \includepdf{fig_ficha_catalografica.pdf}
% \end{fichacatalografica}

% \begin{fichacatalografica}
% 	\sffamily
% 	\vspace*{\fill}                             % Posição vertical
% 	\begin{center}                              % Minipage Centralizado
% 	\fbox{\begin{minipage}[c][8cm]{13.5cm}		% Largura
% 	\small
% 	\imprimirautor
% 	%Sobrenome, Nome do autor
	
% 	\hspace{0.5cm} \imprimirtitulo  / \imprimirautor. --
% 	\imprimirlocal, \imprimirdata-
	
% 	\hspace{0.5cm} \thelastpage p. : il. (algumas color.) ; 30 cm.\\
	
% 	\hspace{0.5cm} \imprimirorientadorRotulo~\imprimirorientador\\
	
% 	\hspace{0.5cm}
% 	\parbox[t]{\textwidth}{\imprimirtipotrabalho~--~\imprimirinstituicao,
% 	\imprimirdata.}\\
	
% 	\hspace{0.5cm}
% 		1. Palavra-chave1.
% 		2. Palavra-chave2.
% 		2. Palavra-chave3.
% 		I. Orientador.
% 		II. Universidade xxx.
% 		III. Faculdade de xxx.
% 		IV. Título 			
% 	\end{minipage}}
% 	\end{center}
% \end{fichacatalografica}
% ---


% ---
% FOLHA DE APROVAÇÃO
% ---

% Isto é um exemplo de Folha de aprovação, elemento obrigatório da NBR
% 14724/2011 (seção 4.2.1.3). Você pode utilizar este modelo até a aprovação
% do trabalho. Após isso, substitua todo o conteúdo deste arquivo por uma
% imagem da página assinada pela banca com o comando abaixo:
%
% \begin{folhadeaprovacao}
% \includepdf{folhadeaprovacao_final.pdf}
% \end{folhadeaprovacao}
%
\setlength{\ABNTEXsignskip}{2cm}
\setlength{\ABNTEXsignwidth}{10cm}

\begin{folhadeaprovacao}
  \begin{center}
    {\ABNTEXchapterfont\large\imprimirautor}

    \vspace*{1cm}

    \begin{center}
      \ABNTEXchapterfont\bfseries\Large\imprimirtitulo
    \end{center}

    \vspace*{1cm}
    
    \hspace{.45\textwidth}
    \begin{minipage}{.5\textwidth}
        \imprimirpreambulo
    \end{minipage}%
   \end{center}

   \vspace*{1cm}
        
   \begin{center}
   Aprovado em 17 de julho de 2024.

   \vspace*{1cm}

   Banca examinadora:
   \end{center}

   \vspace*{-1cm}

   \assinatura{Prof. Ma. Caroline Santa Maria Rodrigues -- PUCRS}
   \assinatura{Prof. Dra. Clarissa Pinto Pizarro de Freitas -- PUCRS}
   \assinatura{\imprimirorientador -- PUCRS}
\end{folhadeaprovacao}
% ---


% ---
% Agradecimentos
% ---
\begin{agradecimentos}
Agradeço aos meus pais, Ademilson José Dias e Irinês Porfírio Dias, pelo apoio incondicional que me deram ao longo da vida,
incentivando mesmo as ideias questionáveis como aquela de fazer uma segunda graduação.

A minha irmã, Maria Eduarda Dias, cuja trajetória acadêmica é fonte de inspiração.

A minha namorada, Nathalia Fallavena Ceratti, que acompanhou a esta jornada de perto, oferecendo apoio e compreensão que
fizeram toda a diferença.

Aos meus amigos, Guilherme da Silva Petró, Norton Neves e Uilinton Luciano Marques da Silva, que foram compreensíveis nos
finais de semana em que não pude jogar RPG.

Ao professor Wagner de Lara Machado pela orientação e apoio no desenvolvimento deste trabalho.
\end{agradecimentos}
% ---


% ---
% RESUMOS
% ---
% Resumo em português
\setlength{\absparsep}{18pt} % ajusta o espaçamento dos parágrafos do resumo
\begin{resumo}
Desenvolvimentos na área de aprendizagem de máquina representam uma série de oportunidades para diversas áreas de pesquisa e de atuação profissional.
O uso de modelos de aprendizagem de máquina na predição de desfecho para tratamentos em saúde mental tem sido investigado e apresenta resultados promissores.
A aplicação desse tipo de modelo possibilitaria, por exemplo, maior assertividade na personalização de tratamentos psicoterápicos a nível individual.
Apesar dos benefícios potenciais, o tema ainda é pouco discutido nos espaços de pesquisa, formação e atuação em psicologia clínica. Este trabalho tem
por objetivo introduzir o tema de aprendizagem de máquina a psicólogos. Busca-se apresentar, de maneira acessível, os conceitos e terminologia básicos
em aprendizagem de máquina, proporcionando um panorama geral da área. Para complementar a exposição conceitual, é construído um exemplo a partir de dados
de um estudo sobre eficácia de intervenções em psicologia positiva para depressão via internet. Por fim, discutem-se algumas implicações do uso desse tipo
de tecnologia para a pesquisa e para a prática em psicoterapia.

 \textbf{Palavras-chave}: aprendizagem de máquina, psicoterapia, psicologia positiva, depressão.
\end{resumo}

% Resumo em inglês
\begin{resumo}[Abstract]
 \begin{otherlanguage*}{english}
   Advances in machine learning have opened up a multitude of opportunities for various research and professional fields. The use of machine learning models to
   predict treatment outcomes in mental health has been investigated showing promising results. The application of such models would allow greater assertiveness
   in tailoring psychotherapeutic treatments at the individual level. Despite its potential benefits, the topic remains unexplored in the fields of research, training
   and practice in clinical psychology. This work aims to introduce the topic of machine learning to psychologists. It presents the basic machine learning concepts and
   terminology in an accessible manner, providing a comprehensive overview of the field. To complement the conceptual discussion, an example is constructed using data
   from a study on the effectiveness of internet-based positive psychology interventions for depression. Finally, implications of using this type of technology in
   psychotherapy research and practice are discussed.
   \vspace{\onelineskip}
 
   \noindent 
   \textbf{Keywords}: machine learning, psychotherapy, positive psychology, depression.
 \end{otherlanguage*}
\end{resumo}


% ---
% LISTA DE ILUSTRAÇÕES
% ---
\pdfbookmark[0]{\listfigurename}{lof}
\listoffigures*
\cleardoublepage
% ---


% ---
% LISTA DE QUADROS
% ---
% \pdfbookmark[0]{\listofquadrosname}{loq}
% \listofquadros*
% \cleardoublepage
% ---


% ---
% LISTA DE TABELAS
% ---
% \pdfbookmark[0]{\listtablename}{lot}
% \listoftables*
% \cleardoublepage
% ---


% ---
% LISTA DE ABREVIATURAS E SIGLAS
% ---
% \begin{siglas}
%   \item[ABNT] Associação Brasileira de Normas Técnicas
%   \item[abnTeX] ABsurdas Normas para TeX
% \end{siglas}
% ---


% ---
% LISTA DE SÍMBOLOS
% ---
% \begin{simbolos}
%   \item[$ \Gamma $] Letra grega Gama
%   \item[$ \Lambda $] Lambda
%   \item[$ \zeta $] Letra grega minúscula zeta
%   \item[$ \in $] Pertence
% \end{simbolos}
% ---


% ---
% Sumário
% ---
\pdfbookmark[0]{\contentsname}{toc}
\tableofcontents*
\cleardoublepage
% ---


% ----------------------------------------------------------
% ELEMENTOS TEXTUAIS
% ----------------------------------------------------------
\textual

\include{./01-introducao/main.tex}

\include{./02-machine-learning/main.tex}

\include{./03-exemplo/main.tex}

\include{./04-conclusao/main.tex}

% ----------------------------------------------------------
% ELEMENTOS PÓS-TEXTUAIS
% ----------------------------------------------------------
\postextual
% ----------------------------------------------------------


% ----------------------------------------------------------
% REFERÊNCIAS
% ----------------------------------------------------------
\bibliography{references.bib}


% ----------------------------------------------------------
% GLOSSÁRIO
% ----------------------------------------------------------
%
% Consulte o manual da classe abntex2 para orientações sobre o glossário.
%
%\glossary

% ----------------------------------------------------------
% APÊNDICES
% ----------------------------------------------------------

% ---
% Inicia os apêndices
% ---
% \begin{apendicesenv}

% % Imprime uma página indicando o início dos apêndices
% \partapendices

% % ----------------------------------------------------------
% \chapter{Quisque libero justo}
% % ----------------------------------------------------------

% % ----------------------------------------------------------
% \chapter{Nullam elementum urna}
% % ----------------------------------------------------------

% \end{apendicesenv}
% ---


% ----------------------------------------------------------
% ANEXOS
% ----------------------------------------------------------

% ---
% Inicia os anexos
% ---
% \begin{anexosenv}

% % Imprime uma página indicando o início dos anexos
% \partanexos

% % ---
% \chapter{Morbi ultrices rutrum lorem}
% % ---

% % ---
% \chapter{Cras non urna sed feugiat cum}
% % ---


% % ---
% \chapter{Fusce facilisis lacinia dui}
% % ---

% \end{anexosenv}

%---------------------------------------------------------------------
% INDICE REMISSIVO
%---------------------------------------------------------------------
% \phantompart
% \printindex
%---------------------------------------------------------------------
\includepdf{./05-extras/contracapa.pdf}
\end{document}


\documentclass[
	% Opções da classe memoir
	12pt,				        % tamanho da fonte
	openany,			        % capítulos começam em pág ímpar (insere página vazia caso preciso)
	oneside,			        % para impressão em recto e verso. Oposto a oneside
	a4paper,			        % tamanho do papel. 
	% Opções da classe abntex2
	%chapter=TITLE,		        % títulos de capítulos convertidos em letras maiúsculas
	%section=TITLE,		        % títulos de seções convertidos em letras maiúsculas
	%subsection=TITLE,	        % títulos de subseções convertidos em letras maiúsculas
	%subsubsection=TITLE,       % títulos de subsubseções convertidos em letras maiúsculas
	% Opções do pacote babel
	english,			        % idioma adicional para hifenização
	brazil				        % o último idioma é o principal do documento
	]{abntex2}

% ---
% Pacotes básicos 
% ---
\usepackage{lmodern}			% Usa a fonte Latin Modern			
\usepackage[T1]{fontenc}		% Selecao de codigos de fonte.
\usepackage[utf8]{inputenc}		% Codificacao do documento (conversão automática dos acentos)
\usepackage{indentfirst}		% Indenta o primeiro parágrafo de cada seção.
\usepackage{color}				% Controle das cores
\usepackage{graphicx}			% Inclusão de gráficos
\usepackage{microtype} 			% para melhorias de justificação
% ---
		

% ---
% Pacotes de citações
% ---
% \usepackage[alf]{abntex2cite}						% Citações padrão ABNT
\usepackage[
  alf,
  abnt-nbr10520=1988,                     % Citações em sentence case
  abnt-emphasize=bf
  ]{abntex2cite}


% ---
% Pacotes adicionais
% ---
\usepackage{pdfpages}	% Inclusão de páginas em PDF


% ---
% Informações para CAPA e FOLHA DE ROSTO
% ---
\titulo{APRENDIZAGEM DE MÁQUINA PARA PREDIÇÃO DE DESFECHOS EM SAÚDE MENTAL: \textmd{UMA INTRODUÇÃO PARA PSICÓLOGOS}}
\autor{GUILHERME AUGUSTO DIAS}
\local{Porto Alegre}
\data{2024}
\orientador{Prof. Dr. Wagner de Lara Machado}
% \coorientador{}
% \instituicao{
%   Pontifícia Universidade Católica do Rio Grande do Sul
%   \par
%   Escola de Ciências da Saúde e da Vida
%   \par
%   Curso de Psicologia}
\tipotrabalho{TCC}
\preambulo{Trabalho de Conclusão de Curso apresentado como requisito para a obtenção do grau de Bacharel em Psicologia pela Escola de Ciências da Saúde e da Vida da Pontifícia Universidade Católica do Rio Grande do Sul.}
% ---

% ---
% Configurações de aparência do PDF final
% ---

% Alterando o aspecto da cor azul
\definecolor{blue}{RGB}{41,5,195}

% Informações do PDF
\makeatletter
\hypersetup{
		pdftitle={Aprendizagem de máquina para predição de desfechos em saúde mental: uma introdução para psicólogos}, 
		pdfauthor={Guilherme Augusto Dias},
    	pdfsubject={\imprimirpreambulo},
	    pdfcreator={LaTeX with abnTeX2},
		pdfkeywords={aprendizagem de máquina}{psicoterapia}{psicologia positiva}{depressão}, 
		colorlinks=true,		% true: colored links; false: boxed links
    	linkcolor=black,		% color of internal links
    	citecolor=black,		% color of links to bibliography
    	filecolor=black,		% color of file links
		urlcolor=black,
		bookmarksdepth=4
}
\makeatother
% --- 


% ---
% Posiciona figuras e tabelas no topo da página quando adicionadas sozinhas
% em um página em branco. Ver https://github.com/abntex/abntex2/issues/170
\makeatletter
\setlength{\@fptop}{5pt} % Set distance from top of page to first float
\makeatother
% ---

% ---
% Possibilita criação de Quadros e Lista de quadros.
% Ver https://github.com/abntex/abntex2/issues/176
%
\newcommand{\quadroname}{Quadro}
\newcommand{\listofquadrosname}{Lista de quadros}

\newfloat[chapter]{quadro}{loq}{\quadroname}
\newlistof{listofquadros}{loq}{\listofquadrosname}
\newlistentry{quadro}{loq}{0}

% Configurações para atender às regras da ABNT
\setfloatadjustment{quadro}{\centering}
\counterwithout{quadro}{chapter}
\renewcommand{\cftquadroname}{\quadroname\space} 
\renewcommand*{\cftquadroaftersnum}{\hfill--\hfill}

\setfloatlocations{quadro}{hbtp} % Ver https://github.com/abntex/abntex2/issues/176
% ---


% --- 
% Espaçamentos entre linhas e parágrafos 
% --- 

% O tamanho do parágrafo é dado por:
\setlength{\parindent}{1.3cm}

% Controle do espaçamento entre um parágrafo e outro:
\setlength{\parskip}{0.2cm}  % tente também \onelineskip

% ---
% Compila o índice
% ---
\makeindex
% ---


% ----
% Início do documento
% ----
\begin{document}

% Seleciona o idioma do documento (conforme pacotes do babel)
\selectlanguage{brazil}

% Retira espaço extra obsoleto entre as frases.
\frenchspacing 

% ----------------------------------------------------------
% ELEMENTOS PRÉ-TEXTUAIS
% ----------------------------------------------------------
\pretextual

% ---
% CAPA
% ---
\includepdf{./05-extras/capa.pdf}
% ---


% ---
% FOLHA DE ROSTO
% (o * indica que haverá a ficha bibliográfica)
% ---
\imprimirfolhaderosto
% ---


% ---
% Ficha bibliografica
% ---

% Isto é um exemplo de Ficha Catalográfica, ou ``Dados internacionais de
% catalogação-na-publicação''. Você pode utilizar este modelo como referência. 
% Porém, provavelmente a biblioteca da sua universidade lhe fornecerá um PDF
% com a ficha catalográfica definitiva após a defesa do trabalho. Quando estiver
% com o documento, salve-o como PDF no diretório do seu projeto e substitua todo
% o conteúdo de implementação deste arquivo pelo comando abaixo:
%
% \usepackage{pdfpages}		% necessário para comando \includepdf
% \begin{fichacatalografica}
%     \includepdf{fig_ficha_catalografica.pdf}
% \end{fichacatalografica}

% \begin{fichacatalografica}
% 	\sffamily
% 	\vspace*{\fill}                             % Posição vertical
% 	\begin{center}                              % Minipage Centralizado
% 	\fbox{\begin{minipage}[c][8cm]{13.5cm}		% Largura
% 	\small
% 	\imprimirautor
% 	%Sobrenome, Nome do autor
	
% 	\hspace{0.5cm} \imprimirtitulo  / \imprimirautor. --
% 	\imprimirlocal, \imprimirdata-
	
% 	\hspace{0.5cm} \thelastpage p. : il. (algumas color.) ; 30 cm.\\
	
% 	\hspace{0.5cm} \imprimirorientadorRotulo~\imprimirorientador\\
	
% 	\hspace{0.5cm}
% 	\parbox[t]{\textwidth}{\imprimirtipotrabalho~--~\imprimirinstituicao,
% 	\imprimirdata.}\\
	
% 	\hspace{0.5cm}
% 		1. Palavra-chave1.
% 		2. Palavra-chave2.
% 		2. Palavra-chave3.
% 		I. Orientador.
% 		II. Universidade xxx.
% 		III. Faculdade de xxx.
% 		IV. Título 			
% 	\end{minipage}}
% 	\end{center}
% \end{fichacatalografica}
% ---


% ---
% FOLHA DE APROVAÇÃO
% ---

% Isto é um exemplo de Folha de aprovação, elemento obrigatório da NBR
% 14724/2011 (seção 4.2.1.3). Você pode utilizar este modelo até a aprovação
% do trabalho. Após isso, substitua todo o conteúdo deste arquivo por uma
% imagem da página assinada pela banca com o comando abaixo:
%
% \begin{folhadeaprovacao}
% \includepdf{folhadeaprovacao_final.pdf}
% \end{folhadeaprovacao}
%
\setlength{\ABNTEXsignskip}{2cm}
\setlength{\ABNTEXsignwidth}{10cm}

\begin{folhadeaprovacao}
  \begin{center}
    {\ABNTEXchapterfont\large\imprimirautor}

    \vspace*{1cm}

    \begin{center}
      \ABNTEXchapterfont\bfseries\Large\imprimirtitulo
    \end{center}

    \vspace*{1cm}
    
    \hspace{.45\textwidth}
    \begin{minipage}{.5\textwidth}
        \imprimirpreambulo
    \end{minipage}%
   \end{center}

   \vspace*{1cm}
        
   \begin{center}
   Aprovado em 17 de julho de 2024.

   \vspace*{1cm}

   Banca examinadora:
   \end{center}

   \vspace*{-1cm}

   \assinatura{Prof. Ma. Caroline Santa Maria Rodrigues -- PUCRS}
   \assinatura{Prof. Dra. Clarissa Pinto Pizarro de Freitas -- PUCRS}
   \assinatura{\imprimirorientador -- PUCRS}
\end{folhadeaprovacao}
% ---


% ---
% Agradecimentos
% ---
\begin{agradecimentos}
Agradeço aos meus pais, Ademilson José Dias e Irinês Porfírio Dias, pelo apoio incondicional que me deram ao longo da vida,
incentivando mesmo as ideias questionáveis como aquela de fazer uma segunda graduação.

A minha irmã, Maria Eduarda Dias, cuja trajetória acadêmica é fonte de inspiração.

A minha namorada, Nathalia Fallavena Ceratti, que acompanhou a esta jornada de perto, oferecendo apoio e compreensão que
fizeram toda a diferença.

Aos meus amigos, Guilherme da Silva Petró, Norton Neves e Uilinton Luciano Marques da Silva, que foram compreensíveis nos
finais de semana em que não pude jogar RPG.

Ao professor Wagner de Lara Machado pela orientação e apoio no desenvolvimento deste trabalho.
\end{agradecimentos}
% ---


% ---
% RESUMOS
% ---
% Resumo em português
\setlength{\absparsep}{18pt} % ajusta o espaçamento dos parágrafos do resumo
\begin{resumo}
Desenvolvimentos na área de aprendizagem de máquina representam uma série de oportunidades para diversas áreas de pesquisa e de atuação profissional.
O uso de modelos de aprendizagem de máquina na predição de desfecho para tratamentos em saúde mental tem sido investigado e apresenta resultados promissores.
A aplicação desse tipo de modelo possibilitaria, por exemplo, maior assertividade na personalização de tratamentos psicoterápicos a nível individual.
Apesar dos benefícios potenciais, o tema ainda é pouco discutido nos espaços de pesquisa, formação e atuação em psicologia clínica. Este trabalho tem
por objetivo introduzir o tema de aprendizagem de máquina a psicólogos. Busca-se apresentar, de maneira acessível, os conceitos e terminologia básicos
em aprendizagem de máquina, proporcionando um panorama geral da área. Para complementar a exposição conceitual, é construído um exemplo a partir de dados
de um estudo sobre eficácia de intervenções em psicologia positiva para depressão via internet. Por fim, discutem-se algumas implicações do uso desse tipo
de tecnologia para a pesquisa e para a prática em psicoterapia.

 \textbf{Palavras-chave}: aprendizagem de máquina, psicoterapia, psicologia positiva, depressão.
\end{resumo}

% Resumo em inglês
\begin{resumo}[Abstract]
 \begin{otherlanguage*}{english}
   Advances in machine learning have opened up a multitude of opportunities for various research and professional fields. The use of machine learning models to
   predict treatment outcomes in mental health has been investigated showing promising results. The application of such models would allow greater assertiveness
   in tailoring psychotherapeutic treatments at the individual level. Despite its potential benefits, the topic remains unexplored in the fields of research, training
   and practice in clinical psychology. This work aims to introduce the topic of machine learning to psychologists. It presents the basic machine learning concepts and
   terminology in an accessible manner, providing a comprehensive overview of the field. To complement the conceptual discussion, an example is constructed using data
   from a study on the effectiveness of internet-based positive psychology interventions for depression. Finally, implications of using this type of technology in
   psychotherapy research and practice are discussed.
   \vspace{\onelineskip}
 
   \noindent 
   \textbf{Keywords}: machine learning, psychotherapy, positive psychology, depression.
 \end{otherlanguage*}
\end{resumo}


% ---
% LISTA DE ILUSTRAÇÕES
% ---
\pdfbookmark[0]{\listfigurename}{lof}
\listoffigures*
\cleardoublepage
% ---


% ---
% LISTA DE QUADROS
% ---
% \pdfbookmark[0]{\listofquadrosname}{loq}
% \listofquadros*
% \cleardoublepage
% ---


% ---
% LISTA DE TABELAS
% ---
% \pdfbookmark[0]{\listtablename}{lot}
% \listoftables*
% \cleardoublepage
% ---


% ---
% LISTA DE ABREVIATURAS E SIGLAS
% ---
% \begin{siglas}
%   \item[ABNT] Associação Brasileira de Normas Técnicas
%   \item[abnTeX] ABsurdas Normas para TeX
% \end{siglas}
% ---


% ---
% LISTA DE SÍMBOLOS
% ---
% \begin{simbolos}
%   \item[$ \Gamma $] Letra grega Gama
%   \item[$ \Lambda $] Lambda
%   \item[$ \zeta $] Letra grega minúscula zeta
%   \item[$ \in $] Pertence
% \end{simbolos}
% ---


% ---
% Sumário
% ---
\pdfbookmark[0]{\contentsname}{toc}
\tableofcontents*
\cleardoublepage
% ---


% ----------------------------------------------------------
% ELEMENTOS TEXTUAIS
% ----------------------------------------------------------
\textual

\include{./01-introducao/main.tex}

\include{./02-machine-learning/main.tex}

\include{./03-exemplo/main.tex}

\include{./04-conclusao/main.tex}

% ----------------------------------------------------------
% ELEMENTOS PÓS-TEXTUAIS
% ----------------------------------------------------------
\postextual
% ----------------------------------------------------------


% ----------------------------------------------------------
% REFERÊNCIAS
% ----------------------------------------------------------
\bibliography{references.bib}


% ----------------------------------------------------------
% GLOSSÁRIO
% ----------------------------------------------------------
%
% Consulte o manual da classe abntex2 para orientações sobre o glossário.
%
%\glossary

% ----------------------------------------------------------
% APÊNDICES
% ----------------------------------------------------------

% ---
% Inicia os apêndices
% ---
% \begin{apendicesenv}

% % Imprime uma página indicando o início dos apêndices
% \partapendices

% % ----------------------------------------------------------
% \chapter{Quisque libero justo}
% % ----------------------------------------------------------

% % ----------------------------------------------------------
% \chapter{Nullam elementum urna}
% % ----------------------------------------------------------

% \end{apendicesenv}
% ---


% ----------------------------------------------------------
% ANEXOS
% ----------------------------------------------------------

% ---
% Inicia os anexos
% ---
% \begin{anexosenv}

% % Imprime uma página indicando o início dos anexos
% \partanexos

% % ---
% \chapter{Morbi ultrices rutrum lorem}
% % ---

% % ---
% \chapter{Cras non urna sed feugiat cum}
% % ---


% % ---
% \chapter{Fusce facilisis lacinia dui}
% % ---

% \end{anexosenv}

%---------------------------------------------------------------------
% INDICE REMISSIVO
%---------------------------------------------------------------------
% \phantompart
% \printindex
%---------------------------------------------------------------------
\includepdf{./05-extras/contracapa.pdf}
\end{document}


\documentclass[
	% Opções da classe memoir
	12pt,				        % tamanho da fonte
	openany,			        % capítulos começam em pág ímpar (insere página vazia caso preciso)
	oneside,			        % para impressão em recto e verso. Oposto a oneside
	a4paper,			        % tamanho do papel. 
	% Opções da classe abntex2
	%chapter=TITLE,		        % títulos de capítulos convertidos em letras maiúsculas
	%section=TITLE,		        % títulos de seções convertidos em letras maiúsculas
	%subsection=TITLE,	        % títulos de subseções convertidos em letras maiúsculas
	%subsubsection=TITLE,       % títulos de subsubseções convertidos em letras maiúsculas
	% Opções do pacote babel
	english,			        % idioma adicional para hifenização
	brazil				        % o último idioma é o principal do documento
	]{abntex2}

% ---
% Pacotes básicos 
% ---
\usepackage{lmodern}			% Usa a fonte Latin Modern			
\usepackage[T1]{fontenc}		% Selecao de codigos de fonte.
\usepackage[utf8]{inputenc}		% Codificacao do documento (conversão automática dos acentos)
\usepackage{indentfirst}		% Indenta o primeiro parágrafo de cada seção.
\usepackage{color}				% Controle das cores
\usepackage{graphicx}			% Inclusão de gráficos
\usepackage{microtype} 			% para melhorias de justificação
% ---
		

% ---
% Pacotes de citações
% ---
% \usepackage[alf]{abntex2cite}						% Citações padrão ABNT
\usepackage[
  alf,
  abnt-nbr10520=1988,                     % Citações em sentence case
  abnt-emphasize=bf
  ]{abntex2cite}


% ---
% Pacotes adicionais
% ---
\usepackage{pdfpages}	% Inclusão de páginas em PDF


% ---
% Informações para CAPA e FOLHA DE ROSTO
% ---
\titulo{APRENDIZAGEM DE MÁQUINA PARA PREDIÇÃO DE DESFECHOS EM SAÚDE MENTAL: \textmd{UMA INTRODUÇÃO PARA PSICÓLOGOS}}
\autor{GUILHERME AUGUSTO DIAS}
\local{Porto Alegre}
\data{2024}
\orientador{Prof. Dr. Wagner de Lara Machado}
% \coorientador{}
% \instituicao{
%   Pontifícia Universidade Católica do Rio Grande do Sul
%   \par
%   Escola de Ciências da Saúde e da Vida
%   \par
%   Curso de Psicologia}
\tipotrabalho{TCC}
\preambulo{Trabalho de Conclusão de Curso apresentado como requisito para a obtenção do grau de Bacharel em Psicologia pela Escola de Ciências da Saúde e da Vida da Pontifícia Universidade Católica do Rio Grande do Sul.}
% ---

% ---
% Configurações de aparência do PDF final
% ---

% Alterando o aspecto da cor azul
\definecolor{blue}{RGB}{41,5,195}

% Informações do PDF
\makeatletter
\hypersetup{
		pdftitle={Aprendizagem de máquina para predição de desfechos em saúde mental: uma introdução para psicólogos}, 
		pdfauthor={Guilherme Augusto Dias},
    	pdfsubject={\imprimirpreambulo},
	    pdfcreator={LaTeX with abnTeX2},
		pdfkeywords={aprendizagem de máquina}{psicoterapia}{psicologia positiva}{depressão}, 
		colorlinks=true,		% true: colored links; false: boxed links
    	linkcolor=black,		% color of internal links
    	citecolor=black,		% color of links to bibliography
    	filecolor=black,		% color of file links
		urlcolor=black,
		bookmarksdepth=4
}
\makeatother
% --- 


% ---
% Posiciona figuras e tabelas no topo da página quando adicionadas sozinhas
% em um página em branco. Ver https://github.com/abntex/abntex2/issues/170
\makeatletter
\setlength{\@fptop}{5pt} % Set distance from top of page to first float
\makeatother
% ---

% ---
% Possibilita criação de Quadros e Lista de quadros.
% Ver https://github.com/abntex/abntex2/issues/176
%
\newcommand{\quadroname}{Quadro}
\newcommand{\listofquadrosname}{Lista de quadros}

\newfloat[chapter]{quadro}{loq}{\quadroname}
\newlistof{listofquadros}{loq}{\listofquadrosname}
\newlistentry{quadro}{loq}{0}

% Configurações para atender às regras da ABNT
\setfloatadjustment{quadro}{\centering}
\counterwithout{quadro}{chapter}
\renewcommand{\cftquadroname}{\quadroname\space} 
\renewcommand*{\cftquadroaftersnum}{\hfill--\hfill}

\setfloatlocations{quadro}{hbtp} % Ver https://github.com/abntex/abntex2/issues/176
% ---


% --- 
% Espaçamentos entre linhas e parágrafos 
% --- 

% O tamanho do parágrafo é dado por:
\setlength{\parindent}{1.3cm}

% Controle do espaçamento entre um parágrafo e outro:
\setlength{\parskip}{0.2cm}  % tente também \onelineskip

% ---
% Compila o índice
% ---
\makeindex
% ---


% ----
% Início do documento
% ----
\begin{document}

% Seleciona o idioma do documento (conforme pacotes do babel)
\selectlanguage{brazil}

% Retira espaço extra obsoleto entre as frases.
\frenchspacing 

% ----------------------------------------------------------
% ELEMENTOS PRÉ-TEXTUAIS
% ----------------------------------------------------------
\pretextual

% ---
% CAPA
% ---
\includepdf{./05-extras/capa.pdf}
% ---


% ---
% FOLHA DE ROSTO
% (o * indica que haverá a ficha bibliográfica)
% ---
\imprimirfolhaderosto
% ---


% ---
% Ficha bibliografica
% ---

% Isto é um exemplo de Ficha Catalográfica, ou ``Dados internacionais de
% catalogação-na-publicação''. Você pode utilizar este modelo como referência. 
% Porém, provavelmente a biblioteca da sua universidade lhe fornecerá um PDF
% com a ficha catalográfica definitiva após a defesa do trabalho. Quando estiver
% com o documento, salve-o como PDF no diretório do seu projeto e substitua todo
% o conteúdo de implementação deste arquivo pelo comando abaixo:
%
% \usepackage{pdfpages}		% necessário para comando \includepdf
% \begin{fichacatalografica}
%     \includepdf{fig_ficha_catalografica.pdf}
% \end{fichacatalografica}

% \begin{fichacatalografica}
% 	\sffamily
% 	\vspace*{\fill}                             % Posição vertical
% 	\begin{center}                              % Minipage Centralizado
% 	\fbox{\begin{minipage}[c][8cm]{13.5cm}		% Largura
% 	\small
% 	\imprimirautor
% 	%Sobrenome, Nome do autor
	
% 	\hspace{0.5cm} \imprimirtitulo  / \imprimirautor. --
% 	\imprimirlocal, \imprimirdata-
	
% 	\hspace{0.5cm} \thelastpage p. : il. (algumas color.) ; 30 cm.\\
	
% 	\hspace{0.5cm} \imprimirorientadorRotulo~\imprimirorientador\\
	
% 	\hspace{0.5cm}
% 	\parbox[t]{\textwidth}{\imprimirtipotrabalho~--~\imprimirinstituicao,
% 	\imprimirdata.}\\
	
% 	\hspace{0.5cm}
% 		1. Palavra-chave1.
% 		2. Palavra-chave2.
% 		2. Palavra-chave3.
% 		I. Orientador.
% 		II. Universidade xxx.
% 		III. Faculdade de xxx.
% 		IV. Título 			
% 	\end{minipage}}
% 	\end{center}
% \end{fichacatalografica}
% ---


% ---
% FOLHA DE APROVAÇÃO
% ---

% Isto é um exemplo de Folha de aprovação, elemento obrigatório da NBR
% 14724/2011 (seção 4.2.1.3). Você pode utilizar este modelo até a aprovação
% do trabalho. Após isso, substitua todo o conteúdo deste arquivo por uma
% imagem da página assinada pela banca com o comando abaixo:
%
% \begin{folhadeaprovacao}
% \includepdf{folhadeaprovacao_final.pdf}
% \end{folhadeaprovacao}
%
\setlength{\ABNTEXsignskip}{2cm}
\setlength{\ABNTEXsignwidth}{10cm}

\begin{folhadeaprovacao}
  \begin{center}
    {\ABNTEXchapterfont\large\imprimirautor}

    \vspace*{1cm}

    \begin{center}
      \ABNTEXchapterfont\bfseries\Large\imprimirtitulo
    \end{center}

    \vspace*{1cm}
    
    \hspace{.45\textwidth}
    \begin{minipage}{.5\textwidth}
        \imprimirpreambulo
    \end{minipage}%
   \end{center}

   \vspace*{1cm}
        
   \begin{center}
   Aprovado em 17 de julho de 2024.

   \vspace*{1cm}

   Banca examinadora:
   \end{center}

   \vspace*{-1cm}

   \assinatura{Prof. Ma. Caroline Santa Maria Rodrigues -- PUCRS}
   \assinatura{Prof. Dra. Clarissa Pinto Pizarro de Freitas -- PUCRS}
   \assinatura{\imprimirorientador -- PUCRS}
\end{folhadeaprovacao}
% ---


% ---
% Agradecimentos
% ---
\begin{agradecimentos}
Agradeço aos meus pais, Ademilson José Dias e Irinês Porfírio Dias, pelo apoio incondicional que me deram ao longo da vida,
incentivando mesmo as ideias questionáveis como aquela de fazer uma segunda graduação.

A minha irmã, Maria Eduarda Dias, cuja trajetória acadêmica é fonte de inspiração.

A minha namorada, Nathalia Fallavena Ceratti, que acompanhou a esta jornada de perto, oferecendo apoio e compreensão que
fizeram toda a diferença.

Aos meus amigos, Guilherme da Silva Petró, Norton Neves e Uilinton Luciano Marques da Silva, que foram compreensíveis nos
finais de semana em que não pude jogar RPG.

Ao professor Wagner de Lara Machado pela orientação e apoio no desenvolvimento deste trabalho.
\end{agradecimentos}
% ---


% ---
% RESUMOS
% ---
% Resumo em português
\setlength{\absparsep}{18pt} % ajusta o espaçamento dos parágrafos do resumo
\begin{resumo}
Desenvolvimentos na área de aprendizagem de máquina representam uma série de oportunidades para diversas áreas de pesquisa e de atuação profissional.
O uso de modelos de aprendizagem de máquina na predição de desfecho para tratamentos em saúde mental tem sido investigado e apresenta resultados promissores.
A aplicação desse tipo de modelo possibilitaria, por exemplo, maior assertividade na personalização de tratamentos psicoterápicos a nível individual.
Apesar dos benefícios potenciais, o tema ainda é pouco discutido nos espaços de pesquisa, formação e atuação em psicologia clínica. Este trabalho tem
por objetivo introduzir o tema de aprendizagem de máquina a psicólogos. Busca-se apresentar, de maneira acessível, os conceitos e terminologia básicos
em aprendizagem de máquina, proporcionando um panorama geral da área. Para complementar a exposição conceitual, é construído um exemplo a partir de dados
de um estudo sobre eficácia de intervenções em psicologia positiva para depressão via internet. Por fim, discutem-se algumas implicações do uso desse tipo
de tecnologia para a pesquisa e para a prática em psicoterapia.

 \textbf{Palavras-chave}: aprendizagem de máquina, psicoterapia, psicologia positiva, depressão.
\end{resumo}

% Resumo em inglês
\begin{resumo}[Abstract]
 \begin{otherlanguage*}{english}
   Advances in machine learning have opened up a multitude of opportunities for various research and professional fields. The use of machine learning models to
   predict treatment outcomes in mental health has been investigated showing promising results. The application of such models would allow greater assertiveness
   in tailoring psychotherapeutic treatments at the individual level. Despite its potential benefits, the topic remains unexplored in the fields of research, training
   and practice in clinical psychology. This work aims to introduce the topic of machine learning to psychologists. It presents the basic machine learning concepts and
   terminology in an accessible manner, providing a comprehensive overview of the field. To complement the conceptual discussion, an example is constructed using data
   from a study on the effectiveness of internet-based positive psychology interventions for depression. Finally, implications of using this type of technology in
   psychotherapy research and practice are discussed.
   \vspace{\onelineskip}
 
   \noindent 
   \textbf{Keywords}: machine learning, psychotherapy, positive psychology, depression.
 \end{otherlanguage*}
\end{resumo}


% ---
% LISTA DE ILUSTRAÇÕES
% ---
\pdfbookmark[0]{\listfigurename}{lof}
\listoffigures*
\cleardoublepage
% ---


% ---
% LISTA DE QUADROS
% ---
% \pdfbookmark[0]{\listofquadrosname}{loq}
% \listofquadros*
% \cleardoublepage
% ---


% ---
% LISTA DE TABELAS
% ---
% \pdfbookmark[0]{\listtablename}{lot}
% \listoftables*
% \cleardoublepage
% ---


% ---
% LISTA DE ABREVIATURAS E SIGLAS
% ---
% \begin{siglas}
%   \item[ABNT] Associação Brasileira de Normas Técnicas
%   \item[abnTeX] ABsurdas Normas para TeX
% \end{siglas}
% ---


% ---
% LISTA DE SÍMBOLOS
% ---
% \begin{simbolos}
%   \item[$ \Gamma $] Letra grega Gama
%   \item[$ \Lambda $] Lambda
%   \item[$ \zeta $] Letra grega minúscula zeta
%   \item[$ \in $] Pertence
% \end{simbolos}
% ---


% ---
% Sumário
% ---
\pdfbookmark[0]{\contentsname}{toc}
\tableofcontents*
\cleardoublepage
% ---


% ----------------------------------------------------------
% ELEMENTOS TEXTUAIS
% ----------------------------------------------------------
\textual

\include{./01-introducao/main.tex}

\include{./02-machine-learning/main.tex}

\include{./03-exemplo/main.tex}

\include{./04-conclusao/main.tex}

% ----------------------------------------------------------
% ELEMENTOS PÓS-TEXTUAIS
% ----------------------------------------------------------
\postextual
% ----------------------------------------------------------


% ----------------------------------------------------------
% REFERÊNCIAS
% ----------------------------------------------------------
\bibliography{references.bib}


% ----------------------------------------------------------
% GLOSSÁRIO
% ----------------------------------------------------------
%
% Consulte o manual da classe abntex2 para orientações sobre o glossário.
%
%\glossary

% ----------------------------------------------------------
% APÊNDICES
% ----------------------------------------------------------

% ---
% Inicia os apêndices
% ---
% \begin{apendicesenv}

% % Imprime uma página indicando o início dos apêndices
% \partapendices

% % ----------------------------------------------------------
% \chapter{Quisque libero justo}
% % ----------------------------------------------------------

% % ----------------------------------------------------------
% \chapter{Nullam elementum urna}
% % ----------------------------------------------------------

% \end{apendicesenv}
% ---


% ----------------------------------------------------------
% ANEXOS
% ----------------------------------------------------------

% ---
% Inicia os anexos
% ---
% \begin{anexosenv}

% % Imprime uma página indicando o início dos anexos
% \partanexos

% % ---
% \chapter{Morbi ultrices rutrum lorem}
% % ---

% % ---
% \chapter{Cras non urna sed feugiat cum}
% % ---


% % ---
% \chapter{Fusce facilisis lacinia dui}
% % ---

% \end{anexosenv}

%---------------------------------------------------------------------
% INDICE REMISSIVO
%---------------------------------------------------------------------
% \phantompart
% \printindex
%---------------------------------------------------------------------
\includepdf{./05-extras/contracapa.pdf}
\end{document}


\documentclass[
	% Opções da classe memoir
	12pt,				        % tamanho da fonte
	openany,			        % capítulos começam em pág ímpar (insere página vazia caso preciso)
	oneside,			        % para impressão em recto e verso. Oposto a oneside
	a4paper,			        % tamanho do papel. 
	% Opções da classe abntex2
	%chapter=TITLE,		        % títulos de capítulos convertidos em letras maiúsculas
	%section=TITLE,		        % títulos de seções convertidos em letras maiúsculas
	%subsection=TITLE,	        % títulos de subseções convertidos em letras maiúsculas
	%subsubsection=TITLE,       % títulos de subsubseções convertidos em letras maiúsculas
	% Opções do pacote babel
	english,			        % idioma adicional para hifenização
	brazil				        % o último idioma é o principal do documento
	]{abntex2}

% ---
% Pacotes básicos 
% ---
\usepackage{lmodern}			% Usa a fonte Latin Modern			
\usepackage[T1]{fontenc}		% Selecao de codigos de fonte.
\usepackage[utf8]{inputenc}		% Codificacao do documento (conversão automática dos acentos)
\usepackage{indentfirst}		% Indenta o primeiro parágrafo de cada seção.
\usepackage{color}				% Controle das cores
\usepackage{graphicx}			% Inclusão de gráficos
\usepackage{microtype} 			% para melhorias de justificação
% ---
		

% ---
% Pacotes de citações
% ---
% \usepackage[alf]{abntex2cite}						% Citações padrão ABNT
\usepackage[
  alf,
  abnt-nbr10520=1988,                     % Citações em sentence case
  abnt-emphasize=bf
  ]{abntex2cite}


% ---
% Pacotes adicionais
% ---
\usepackage{pdfpages}	% Inclusão de páginas em PDF


% ---
% Informações para CAPA e FOLHA DE ROSTO
% ---
\titulo{APRENDIZAGEM DE MÁQUINA PARA PREDIÇÃO DE DESFECHOS EM SAÚDE MENTAL: \textmd{UMA INTRODUÇÃO PARA PSICÓLOGOS}}
\autor{GUILHERME AUGUSTO DIAS}
\local{Porto Alegre}
\data{2024}
\orientador{Prof. Dr. Wagner de Lara Machado}
% \coorientador{}
% \instituicao{
%   Pontifícia Universidade Católica do Rio Grande do Sul
%   \par
%   Escola de Ciências da Saúde e da Vida
%   \par
%   Curso de Psicologia}
\tipotrabalho{TCC}
\preambulo{Trabalho de Conclusão de Curso apresentado como requisito para a obtenção do grau de Bacharel em Psicologia pela Escola de Ciências da Saúde e da Vida da Pontifícia Universidade Católica do Rio Grande do Sul.}
% ---

% ---
% Configurações de aparência do PDF final
% ---

% Alterando o aspecto da cor azul
\definecolor{blue}{RGB}{41,5,195}

% Informações do PDF
\makeatletter
\hypersetup{
		pdftitle={Aprendizagem de máquina para predição de desfechos em saúde mental: uma introdução para psicólogos}, 
		pdfauthor={Guilherme Augusto Dias},
    	pdfsubject={\imprimirpreambulo},
	    pdfcreator={LaTeX with abnTeX2},
		pdfkeywords={aprendizagem de máquina}{psicoterapia}{psicologia positiva}{depressão}, 
		colorlinks=true,		% true: colored links; false: boxed links
    	linkcolor=black,		% color of internal links
    	citecolor=black,		% color of links to bibliography
    	filecolor=black,		% color of file links
		urlcolor=black,
		bookmarksdepth=4
}
\makeatother
% --- 


% ---
% Posiciona figuras e tabelas no topo da página quando adicionadas sozinhas
% em um página em branco. Ver https://github.com/abntex/abntex2/issues/170
\makeatletter
\setlength{\@fptop}{5pt} % Set distance from top of page to first float
\makeatother
% ---

% ---
% Possibilita criação de Quadros e Lista de quadros.
% Ver https://github.com/abntex/abntex2/issues/176
%
\newcommand{\quadroname}{Quadro}
\newcommand{\listofquadrosname}{Lista de quadros}

\newfloat[chapter]{quadro}{loq}{\quadroname}
\newlistof{listofquadros}{loq}{\listofquadrosname}
\newlistentry{quadro}{loq}{0}

% Configurações para atender às regras da ABNT
\setfloatadjustment{quadro}{\centering}
\counterwithout{quadro}{chapter}
\renewcommand{\cftquadroname}{\quadroname\space} 
\renewcommand*{\cftquadroaftersnum}{\hfill--\hfill}

\setfloatlocations{quadro}{hbtp} % Ver https://github.com/abntex/abntex2/issues/176
% ---


% --- 
% Espaçamentos entre linhas e parágrafos 
% --- 

% O tamanho do parágrafo é dado por:
\setlength{\parindent}{1.3cm}

% Controle do espaçamento entre um parágrafo e outro:
\setlength{\parskip}{0.2cm}  % tente também \onelineskip

% ---
% Compila o índice
% ---
\makeindex
% ---


% ----
% Início do documento
% ----
\begin{document}

% Seleciona o idioma do documento (conforme pacotes do babel)
\selectlanguage{brazil}

% Retira espaço extra obsoleto entre as frases.
\frenchspacing 

% ----------------------------------------------------------
% ELEMENTOS PRÉ-TEXTUAIS
% ----------------------------------------------------------
\pretextual

% ---
% CAPA
% ---
\includepdf{./05-extras/capa.pdf}
% ---


% ---
% FOLHA DE ROSTO
% (o * indica que haverá a ficha bibliográfica)
% ---
\imprimirfolhaderosto
% ---


% ---
% Ficha bibliografica
% ---

% Isto é um exemplo de Ficha Catalográfica, ou ``Dados internacionais de
% catalogação-na-publicação''. Você pode utilizar este modelo como referência. 
% Porém, provavelmente a biblioteca da sua universidade lhe fornecerá um PDF
% com a ficha catalográfica definitiva após a defesa do trabalho. Quando estiver
% com o documento, salve-o como PDF no diretório do seu projeto e substitua todo
% o conteúdo de implementação deste arquivo pelo comando abaixo:
%
% \usepackage{pdfpages}		% necessário para comando \includepdf
% \begin{fichacatalografica}
%     \includepdf{fig_ficha_catalografica.pdf}
% \end{fichacatalografica}

% \begin{fichacatalografica}
% 	\sffamily
% 	\vspace*{\fill}                             % Posição vertical
% 	\begin{center}                              % Minipage Centralizado
% 	\fbox{\begin{minipage}[c][8cm]{13.5cm}		% Largura
% 	\small
% 	\imprimirautor
% 	%Sobrenome, Nome do autor
	
% 	\hspace{0.5cm} \imprimirtitulo  / \imprimirautor. --
% 	\imprimirlocal, \imprimirdata-
	
% 	\hspace{0.5cm} \thelastpage p. : il. (algumas color.) ; 30 cm.\\
	
% 	\hspace{0.5cm} \imprimirorientadorRotulo~\imprimirorientador\\
	
% 	\hspace{0.5cm}
% 	\parbox[t]{\textwidth}{\imprimirtipotrabalho~--~\imprimirinstituicao,
% 	\imprimirdata.}\\
	
% 	\hspace{0.5cm}
% 		1. Palavra-chave1.
% 		2. Palavra-chave2.
% 		2. Palavra-chave3.
% 		I. Orientador.
% 		II. Universidade xxx.
% 		III. Faculdade de xxx.
% 		IV. Título 			
% 	\end{minipage}}
% 	\end{center}
% \end{fichacatalografica}
% ---


% ---
% FOLHA DE APROVAÇÃO
% ---

% Isto é um exemplo de Folha de aprovação, elemento obrigatório da NBR
% 14724/2011 (seção 4.2.1.3). Você pode utilizar este modelo até a aprovação
% do trabalho. Após isso, substitua todo o conteúdo deste arquivo por uma
% imagem da página assinada pela banca com o comando abaixo:
%
% \begin{folhadeaprovacao}
% \includepdf{folhadeaprovacao_final.pdf}
% \end{folhadeaprovacao}
%
\setlength{\ABNTEXsignskip}{2cm}
\setlength{\ABNTEXsignwidth}{10cm}

\begin{folhadeaprovacao}
  \begin{center}
    {\ABNTEXchapterfont\large\imprimirautor}

    \vspace*{1cm}

    \begin{center}
      \ABNTEXchapterfont\bfseries\Large\imprimirtitulo
    \end{center}

    \vspace*{1cm}
    
    \hspace{.45\textwidth}
    \begin{minipage}{.5\textwidth}
        \imprimirpreambulo
    \end{minipage}%
   \end{center}

   \vspace*{1cm}
        
   \begin{center}
   Aprovado em 17 de julho de 2024.

   \vspace*{1cm}

   Banca examinadora:
   \end{center}

   \vspace*{-1cm}

   \assinatura{Prof. Ma. Caroline Santa Maria Rodrigues -- PUCRS}
   \assinatura{Prof. Dra. Clarissa Pinto Pizarro de Freitas -- PUCRS}
   \assinatura{\imprimirorientador -- PUCRS}
\end{folhadeaprovacao}
% ---


% ---
% Agradecimentos
% ---
\begin{agradecimentos}
Agradeço aos meus pais, Ademilson José Dias e Irinês Porfírio Dias, pelo apoio incondicional que me deram ao longo da vida,
incentivando mesmo as ideias questionáveis como aquela de fazer uma segunda graduação.

A minha irmã, Maria Eduarda Dias, cuja trajetória acadêmica é fonte de inspiração.

A minha namorada, Nathalia Fallavena Ceratti, que acompanhou a esta jornada de perto, oferecendo apoio e compreensão que
fizeram toda a diferença.

Aos meus amigos, Guilherme da Silva Petró, Norton Neves e Uilinton Luciano Marques da Silva, que foram compreensíveis nos
finais de semana em que não pude jogar RPG.

Ao professor Wagner de Lara Machado pela orientação e apoio no desenvolvimento deste trabalho.
\end{agradecimentos}
% ---


% ---
% RESUMOS
% ---
% Resumo em português
\setlength{\absparsep}{18pt} % ajusta o espaçamento dos parágrafos do resumo
\begin{resumo}
Desenvolvimentos na área de aprendizagem de máquina representam uma série de oportunidades para diversas áreas de pesquisa e de atuação profissional.
O uso de modelos de aprendizagem de máquina na predição de desfecho para tratamentos em saúde mental tem sido investigado e apresenta resultados promissores.
A aplicação desse tipo de modelo possibilitaria, por exemplo, maior assertividade na personalização de tratamentos psicoterápicos a nível individual.
Apesar dos benefícios potenciais, o tema ainda é pouco discutido nos espaços de pesquisa, formação e atuação em psicologia clínica. Este trabalho tem
por objetivo introduzir o tema de aprendizagem de máquina a psicólogos. Busca-se apresentar, de maneira acessível, os conceitos e terminologia básicos
em aprendizagem de máquina, proporcionando um panorama geral da área. Para complementar a exposição conceitual, é construído um exemplo a partir de dados
de um estudo sobre eficácia de intervenções em psicologia positiva para depressão via internet. Por fim, discutem-se algumas implicações do uso desse tipo
de tecnologia para a pesquisa e para a prática em psicoterapia.

 \textbf{Palavras-chave}: aprendizagem de máquina, psicoterapia, psicologia positiva, depressão.
\end{resumo}

% Resumo em inglês
\begin{resumo}[Abstract]
 \begin{otherlanguage*}{english}
   Advances in machine learning have opened up a multitude of opportunities for various research and professional fields. The use of machine learning models to
   predict treatment outcomes in mental health has been investigated showing promising results. The application of such models would allow greater assertiveness
   in tailoring psychotherapeutic treatments at the individual level. Despite its potential benefits, the topic remains unexplored in the fields of research, training
   and practice in clinical psychology. This work aims to introduce the topic of machine learning to psychologists. It presents the basic machine learning concepts and
   terminology in an accessible manner, providing a comprehensive overview of the field. To complement the conceptual discussion, an example is constructed using data
   from a study on the effectiveness of internet-based positive psychology interventions for depression. Finally, implications of using this type of technology in
   psychotherapy research and practice are discussed.
   \vspace{\onelineskip}
 
   \noindent 
   \textbf{Keywords}: machine learning, psychotherapy, positive psychology, depression.
 \end{otherlanguage*}
\end{resumo}


% ---
% LISTA DE ILUSTRAÇÕES
% ---
\pdfbookmark[0]{\listfigurename}{lof}
\listoffigures*
\cleardoublepage
% ---


% ---
% LISTA DE QUADROS
% ---
% \pdfbookmark[0]{\listofquadrosname}{loq}
% \listofquadros*
% \cleardoublepage
% ---


% ---
% LISTA DE TABELAS
% ---
% \pdfbookmark[0]{\listtablename}{lot}
% \listoftables*
% \cleardoublepage
% ---


% ---
% LISTA DE ABREVIATURAS E SIGLAS
% ---
% \begin{siglas}
%   \item[ABNT] Associação Brasileira de Normas Técnicas
%   \item[abnTeX] ABsurdas Normas para TeX
% \end{siglas}
% ---


% ---
% LISTA DE SÍMBOLOS
% ---
% \begin{simbolos}
%   \item[$ \Gamma $] Letra grega Gama
%   \item[$ \Lambda $] Lambda
%   \item[$ \zeta $] Letra grega minúscula zeta
%   \item[$ \in $] Pertence
% \end{simbolos}
% ---


% ---
% Sumário
% ---
\pdfbookmark[0]{\contentsname}{toc}
\tableofcontents*
\cleardoublepage
% ---


% ----------------------------------------------------------
% ELEMENTOS TEXTUAIS
% ----------------------------------------------------------
\textual

\include{./01-introducao/main.tex}

\include{./02-machine-learning/main.tex}

\include{./03-exemplo/main.tex}

\include{./04-conclusao/main.tex}

% ----------------------------------------------------------
% ELEMENTOS PÓS-TEXTUAIS
% ----------------------------------------------------------
\postextual
% ----------------------------------------------------------


% ----------------------------------------------------------
% REFERÊNCIAS
% ----------------------------------------------------------
\bibliography{references.bib}


% ----------------------------------------------------------
% GLOSSÁRIO
% ----------------------------------------------------------
%
% Consulte o manual da classe abntex2 para orientações sobre o glossário.
%
%\glossary

% ----------------------------------------------------------
% APÊNDICES
% ----------------------------------------------------------

% ---
% Inicia os apêndices
% ---
% \begin{apendicesenv}

% % Imprime uma página indicando o início dos apêndices
% \partapendices

% % ----------------------------------------------------------
% \chapter{Quisque libero justo}
% % ----------------------------------------------------------

% % ----------------------------------------------------------
% \chapter{Nullam elementum urna}
% % ----------------------------------------------------------

% \end{apendicesenv}
% ---


% ----------------------------------------------------------
% ANEXOS
% ----------------------------------------------------------

% ---
% Inicia os anexos
% ---
% \begin{anexosenv}

% % Imprime uma página indicando o início dos anexos
% \partanexos

% % ---
% \chapter{Morbi ultrices rutrum lorem}
% % ---

% % ---
% \chapter{Cras non urna sed feugiat cum}
% % ---


% % ---
% \chapter{Fusce facilisis lacinia dui}
% % ---

% \end{anexosenv}

%---------------------------------------------------------------------
% INDICE REMISSIVO
%---------------------------------------------------------------------
% \phantompart
% \printindex
%---------------------------------------------------------------------
\includepdf{./05-extras/contracapa.pdf}
\end{document}


% ----------------------------------------------------------
% ELEMENTOS PÓS-TEXTUAIS
% ----------------------------------------------------------
\postextual
% ----------------------------------------------------------


% ----------------------------------------------------------
% REFERÊNCIAS
% ----------------------------------------------------------
\bibliography{references.bib}


% ----------------------------------------------------------
% GLOSSÁRIO
% ----------------------------------------------------------
%
% Consulte o manual da classe abntex2 para orientações sobre o glossário.
%
%\glossary

% ----------------------------------------------------------
% APÊNDICES
% ----------------------------------------------------------

% ---
% Inicia os apêndices
% ---
% \begin{apendicesenv}

% % Imprime uma página indicando o início dos apêndices
% \partapendices

% % ----------------------------------------------------------
% \chapter{Quisque libero justo}
% % ----------------------------------------------------------

% % ----------------------------------------------------------
% \chapter{Nullam elementum urna}
% % ----------------------------------------------------------

% \end{apendicesenv}
% ---


% ----------------------------------------------------------
% ANEXOS
% ----------------------------------------------------------

% ---
% Inicia os anexos
% ---
% \begin{anexosenv}

% % Imprime uma página indicando o início dos anexos
% \partanexos

% % ---
% \chapter{Morbi ultrices rutrum lorem}
% % ---

% % ---
% \chapter{Cras non urna sed feugiat cum}
% % ---


% % ---
% \chapter{Fusce facilisis lacinia dui}
% % ---

% \end{anexosenv}

%---------------------------------------------------------------------
% INDICE REMISSIVO
%---------------------------------------------------------------------
% \phantompart
% \printindex
%---------------------------------------------------------------------
\includepdf{./05-extras/contracapa.pdf}
\end{document}


\documentclass[
	% Opções da classe memoir
	12pt,				        % tamanho da fonte
	openany,			        % capítulos começam em pág ímpar (insere página vazia caso preciso)
	oneside,			        % para impressão em recto e verso. Oposto a oneside
	a4paper,			        % tamanho do papel. 
	% Opções da classe abntex2
	%chapter=TITLE,		        % títulos de capítulos convertidos em letras maiúsculas
	%section=TITLE,		        % títulos de seções convertidos em letras maiúsculas
	%subsection=TITLE,	        % títulos de subseções convertidos em letras maiúsculas
	%subsubsection=TITLE,       % títulos de subsubseções convertidos em letras maiúsculas
	% Opções do pacote babel
	english,			        % idioma adicional para hifenização
	brazil				        % o último idioma é o principal do documento
	]{abntex2}

% ---
% Pacotes básicos 
% ---
\usepackage{lmodern}			% Usa a fonte Latin Modern			
\usepackage[T1]{fontenc}		% Selecao de codigos de fonte.
\usepackage[utf8]{inputenc}		% Codificacao do documento (conversão automática dos acentos)
\usepackage{indentfirst}		% Indenta o primeiro parágrafo de cada seção.
\usepackage{color}				% Controle das cores
\usepackage{graphicx}			% Inclusão de gráficos
\usepackage{microtype} 			% para melhorias de justificação
% ---
		

% ---
% Pacotes de citações
% ---
% \usepackage[alf]{abntex2cite}						% Citações padrão ABNT
\usepackage[
  alf,
  abnt-nbr10520=1988,                     % Citações em sentence case
  abnt-emphasize=bf
  ]{abntex2cite}


% ---
% Pacotes adicionais
% ---
\usepackage{pdfpages}	% Inclusão de páginas em PDF


% ---
% Informações para CAPA e FOLHA DE ROSTO
% ---
\titulo{APRENDIZAGEM DE MÁQUINA PARA PREDIÇÃO DE DESFECHOS EM SAÚDE MENTAL: \textmd{UMA INTRODUÇÃO PARA PSICÓLOGOS}}
\autor{GUILHERME AUGUSTO DIAS}
\local{Porto Alegre}
\data{2024}
\orientador{Prof. Dr. Wagner de Lara Machado}
% \coorientador{}
% \instituicao{
%   Pontifícia Universidade Católica do Rio Grande do Sul
%   \par
%   Escola de Ciências da Saúde e da Vida
%   \par
%   Curso de Psicologia}
\tipotrabalho{TCC}
\preambulo{Trabalho de Conclusão de Curso apresentado como requisito para a obtenção do grau de Bacharel em Psicologia pela Escola de Ciências da Saúde e da Vida da Pontifícia Universidade Católica do Rio Grande do Sul.}
% ---

% ---
% Configurações de aparência do PDF final
% ---

% Alterando o aspecto da cor azul
\definecolor{blue}{RGB}{41,5,195}

% Informações do PDF
\makeatletter
\hypersetup{
		pdftitle={Aprendizagem de máquina para predição de desfechos em saúde mental: uma introdução para psicólogos}, 
		pdfauthor={Guilherme Augusto Dias},
    	pdfsubject={\imprimirpreambulo},
	    pdfcreator={LaTeX with abnTeX2},
		pdfkeywords={aprendizagem de máquina}{psicoterapia}{psicologia positiva}{depressão}, 
		colorlinks=true,		% true: colored links; false: boxed links
    	linkcolor=black,		% color of internal links
    	citecolor=black,		% color of links to bibliography
    	filecolor=black,		% color of file links
		urlcolor=black,
		bookmarksdepth=4
}
\makeatother
% --- 


% ---
% Posiciona figuras e tabelas no topo da página quando adicionadas sozinhas
% em um página em branco. Ver https://github.com/abntex/abntex2/issues/170
\makeatletter
\setlength{\@fptop}{5pt} % Set distance from top of page to first float
\makeatother
% ---

% ---
% Possibilita criação de Quadros e Lista de quadros.
% Ver https://github.com/abntex/abntex2/issues/176
%
\newcommand{\quadroname}{Quadro}
\newcommand{\listofquadrosname}{Lista de quadros}

\newfloat[chapter]{quadro}{loq}{\quadroname}
\newlistof{listofquadros}{loq}{\listofquadrosname}
\newlistentry{quadro}{loq}{0}

% Configurações para atender às regras da ABNT
\setfloatadjustment{quadro}{\centering}
\counterwithout{quadro}{chapter}
\renewcommand{\cftquadroname}{\quadroname\space} 
\renewcommand*{\cftquadroaftersnum}{\hfill--\hfill}

\setfloatlocations{quadro}{hbtp} % Ver https://github.com/abntex/abntex2/issues/176
% ---


% --- 
% Espaçamentos entre linhas e parágrafos 
% --- 

% O tamanho do parágrafo é dado por:
\setlength{\parindent}{1.3cm}

% Controle do espaçamento entre um parágrafo e outro:
\setlength{\parskip}{0.2cm}  % tente também \onelineskip

% ---
% Compila o índice
% ---
\makeindex
% ---


% ----
% Início do documento
% ----
\begin{document}

% Seleciona o idioma do documento (conforme pacotes do babel)
\selectlanguage{brazil}

% Retira espaço extra obsoleto entre as frases.
\frenchspacing 

% ----------------------------------------------------------
% ELEMENTOS PRÉ-TEXTUAIS
% ----------------------------------------------------------
\pretextual

% ---
% CAPA
% ---
\includepdf{./05-extras/capa.pdf}
% ---


% ---
% FOLHA DE ROSTO
% (o * indica que haverá a ficha bibliográfica)
% ---
\imprimirfolhaderosto
% ---


% ---
% Ficha bibliografica
% ---

% Isto é um exemplo de Ficha Catalográfica, ou ``Dados internacionais de
% catalogação-na-publicação''. Você pode utilizar este modelo como referência. 
% Porém, provavelmente a biblioteca da sua universidade lhe fornecerá um PDF
% com a ficha catalográfica definitiva após a defesa do trabalho. Quando estiver
% com o documento, salve-o como PDF no diretório do seu projeto e substitua todo
% o conteúdo de implementação deste arquivo pelo comando abaixo:
%
% \usepackage{pdfpages}		% necessário para comando \includepdf
% \begin{fichacatalografica}
%     \includepdf{fig_ficha_catalografica.pdf}
% \end{fichacatalografica}

% \begin{fichacatalografica}
% 	\sffamily
% 	\vspace*{\fill}                             % Posição vertical
% 	\begin{center}                              % Minipage Centralizado
% 	\fbox{\begin{minipage}[c][8cm]{13.5cm}		% Largura
% 	\small
% 	\imprimirautor
% 	%Sobrenome, Nome do autor
	
% 	\hspace{0.5cm} \imprimirtitulo  / \imprimirautor. --
% 	\imprimirlocal, \imprimirdata-
	
% 	\hspace{0.5cm} \thelastpage p. : il. (algumas color.) ; 30 cm.\\
	
% 	\hspace{0.5cm} \imprimirorientadorRotulo~\imprimirorientador\\
	
% 	\hspace{0.5cm}
% 	\parbox[t]{\textwidth}{\imprimirtipotrabalho~--~\imprimirinstituicao,
% 	\imprimirdata.}\\
	
% 	\hspace{0.5cm}
% 		1. Palavra-chave1.
% 		2. Palavra-chave2.
% 		2. Palavra-chave3.
% 		I. Orientador.
% 		II. Universidade xxx.
% 		III. Faculdade de xxx.
% 		IV. Título 			
% 	\end{minipage}}
% 	\end{center}
% \end{fichacatalografica}
% ---


% ---
% FOLHA DE APROVAÇÃO
% ---

% Isto é um exemplo de Folha de aprovação, elemento obrigatório da NBR
% 14724/2011 (seção 4.2.1.3). Você pode utilizar este modelo até a aprovação
% do trabalho. Após isso, substitua todo o conteúdo deste arquivo por uma
% imagem da página assinada pela banca com o comando abaixo:
%
% \begin{folhadeaprovacao}
% \includepdf{folhadeaprovacao_final.pdf}
% \end{folhadeaprovacao}
%
\setlength{\ABNTEXsignskip}{2cm}
\setlength{\ABNTEXsignwidth}{10cm}

\begin{folhadeaprovacao}
  \begin{center}
    {\ABNTEXchapterfont\large\imprimirautor}

    \vspace*{1cm}

    \begin{center}
      \ABNTEXchapterfont\bfseries\Large\imprimirtitulo
    \end{center}

    \vspace*{1cm}
    
    \hspace{.45\textwidth}
    \begin{minipage}{.5\textwidth}
        \imprimirpreambulo
    \end{minipage}%
   \end{center}

   \vspace*{1cm}
        
   \begin{center}
   Aprovado em 17 de julho de 2024.

   \vspace*{1cm}

   Banca examinadora:
   \end{center}

   \vspace*{-1cm}

   \assinatura{Prof. Ma. Caroline Santa Maria Rodrigues -- PUCRS}
   \assinatura{Prof. Dra. Clarissa Pinto Pizarro de Freitas -- PUCRS}
   \assinatura{\imprimirorientador -- PUCRS}
\end{folhadeaprovacao}
% ---


% ---
% Agradecimentos
% ---
\begin{agradecimentos}
Agradeço aos meus pais, Ademilson José Dias e Irinês Porfírio Dias, pelo apoio incondicional que me deram ao longo da vida,
incentivando mesmo as ideias questionáveis como aquela de fazer uma segunda graduação.

A minha irmã, Maria Eduarda Dias, cuja trajetória acadêmica é fonte de inspiração.

A minha namorada, Nathalia Fallavena Ceratti, que acompanhou a esta jornada de perto, oferecendo apoio e compreensão que
fizeram toda a diferença.

Aos meus amigos, Guilherme da Silva Petró, Norton Neves e Uilinton Luciano Marques da Silva, que foram compreensíveis nos
finais de semana em que não pude jogar RPG.

Ao professor Wagner de Lara Machado pela orientação e apoio no desenvolvimento deste trabalho.
\end{agradecimentos}
% ---


% ---
% RESUMOS
% ---
% Resumo em português
\setlength{\absparsep}{18pt} % ajusta o espaçamento dos parágrafos do resumo
\begin{resumo}
Desenvolvimentos na área de aprendizagem de máquina representam uma série de oportunidades para diversas áreas de pesquisa e de atuação profissional.
O uso de modelos de aprendizagem de máquina na predição de desfecho para tratamentos em saúde mental tem sido investigado e apresenta resultados promissores.
A aplicação desse tipo de modelo possibilitaria, por exemplo, maior assertividade na personalização de tratamentos psicoterápicos a nível individual.
Apesar dos benefícios potenciais, o tema ainda é pouco discutido nos espaços de pesquisa, formação e atuação em psicologia clínica. Este trabalho tem
por objetivo introduzir o tema de aprendizagem de máquina a psicólogos. Busca-se apresentar, de maneira acessível, os conceitos e terminologia básicos
em aprendizagem de máquina, proporcionando um panorama geral da área. Para complementar a exposição conceitual, é construído um exemplo a partir de dados
de um estudo sobre eficácia de intervenções em psicologia positiva para depressão via internet. Por fim, discutem-se algumas implicações do uso desse tipo
de tecnologia para a pesquisa e para a prática em psicoterapia.

 \textbf{Palavras-chave}: aprendizagem de máquina, psicoterapia, psicologia positiva, depressão.
\end{resumo}

% Resumo em inglês
\begin{resumo}[Abstract]
 \begin{otherlanguage*}{english}
   Advances in machine learning have opened up a multitude of opportunities for various research and professional fields. The use of machine learning models to
   predict treatment outcomes in mental health has been investigated showing promising results. The application of such models would allow greater assertiveness
   in tailoring psychotherapeutic treatments at the individual level. Despite its potential benefits, the topic remains unexplored in the fields of research, training
   and practice in clinical psychology. This work aims to introduce the topic of machine learning to psychologists. It presents the basic machine learning concepts and
   terminology in an accessible manner, providing a comprehensive overview of the field. To complement the conceptual discussion, an example is constructed using data
   from a study on the effectiveness of internet-based positive psychology interventions for depression. Finally, implications of using this type of technology in
   psychotherapy research and practice are discussed.
   \vspace{\onelineskip}
 
   \noindent 
   \textbf{Keywords}: machine learning, psychotherapy, positive psychology, depression.
 \end{otherlanguage*}
\end{resumo}


% ---
% LISTA DE ILUSTRAÇÕES
% ---
\pdfbookmark[0]{\listfigurename}{lof}
\listoffigures*
\cleardoublepage
% ---


% ---
% LISTA DE QUADROS
% ---
% \pdfbookmark[0]{\listofquadrosname}{loq}
% \listofquadros*
% \cleardoublepage
% ---


% ---
% LISTA DE TABELAS
% ---
% \pdfbookmark[0]{\listtablename}{lot}
% \listoftables*
% \cleardoublepage
% ---


% ---
% LISTA DE ABREVIATURAS E SIGLAS
% ---
% \begin{siglas}
%   \item[ABNT] Associação Brasileira de Normas Técnicas
%   \item[abnTeX] ABsurdas Normas para TeX
% \end{siglas}
% ---


% ---
% LISTA DE SÍMBOLOS
% ---
% \begin{simbolos}
%   \item[$ \Gamma $] Letra grega Gama
%   \item[$ \Lambda $] Lambda
%   \item[$ \zeta $] Letra grega minúscula zeta
%   \item[$ \in $] Pertence
% \end{simbolos}
% ---


% ---
% Sumário
% ---
\pdfbookmark[0]{\contentsname}{toc}
\tableofcontents*
\cleardoublepage
% ---


% ----------------------------------------------------------
% ELEMENTOS TEXTUAIS
% ----------------------------------------------------------
\textual

\documentclass[
	% Opções da classe memoir
	12pt,				        % tamanho da fonte
	openany,			        % capítulos começam em pág ímpar (insere página vazia caso preciso)
	oneside,			        % para impressão em recto e verso. Oposto a oneside
	a4paper,			        % tamanho do papel. 
	% Opções da classe abntex2
	%chapter=TITLE,		        % títulos de capítulos convertidos em letras maiúsculas
	%section=TITLE,		        % títulos de seções convertidos em letras maiúsculas
	%subsection=TITLE,	        % títulos de subseções convertidos em letras maiúsculas
	%subsubsection=TITLE,       % títulos de subsubseções convertidos em letras maiúsculas
	% Opções do pacote babel
	english,			        % idioma adicional para hifenização
	brazil				        % o último idioma é o principal do documento
	]{abntex2}

% ---
% Pacotes básicos 
% ---
\usepackage{lmodern}			% Usa a fonte Latin Modern			
\usepackage[T1]{fontenc}		% Selecao de codigos de fonte.
\usepackage[utf8]{inputenc}		% Codificacao do documento (conversão automática dos acentos)
\usepackage{indentfirst}		% Indenta o primeiro parágrafo de cada seção.
\usepackage{color}				% Controle das cores
\usepackage{graphicx}			% Inclusão de gráficos
\usepackage{microtype} 			% para melhorias de justificação
% ---
		

% ---
% Pacotes de citações
% ---
% \usepackage[alf]{abntex2cite}						% Citações padrão ABNT
\usepackage[
  alf,
  abnt-nbr10520=1988,                     % Citações em sentence case
  abnt-emphasize=bf
  ]{abntex2cite}


% ---
% Pacotes adicionais
% ---
\usepackage{pdfpages}	% Inclusão de páginas em PDF


% ---
% Informações para CAPA e FOLHA DE ROSTO
% ---
\titulo{APRENDIZAGEM DE MÁQUINA PARA PREDIÇÃO DE DESFECHOS EM SAÚDE MENTAL: \textmd{UMA INTRODUÇÃO PARA PSICÓLOGOS}}
\autor{GUILHERME AUGUSTO DIAS}
\local{Porto Alegre}
\data{2024}
\orientador{Prof. Dr. Wagner de Lara Machado}
% \coorientador{}
% \instituicao{
%   Pontifícia Universidade Católica do Rio Grande do Sul
%   \par
%   Escola de Ciências da Saúde e da Vida
%   \par
%   Curso de Psicologia}
\tipotrabalho{TCC}
\preambulo{Trabalho de Conclusão de Curso apresentado como requisito para a obtenção do grau de Bacharel em Psicologia pela Escola de Ciências da Saúde e da Vida da Pontifícia Universidade Católica do Rio Grande do Sul.}
% ---

% ---
% Configurações de aparência do PDF final
% ---

% Alterando o aspecto da cor azul
\definecolor{blue}{RGB}{41,5,195}

% Informações do PDF
\makeatletter
\hypersetup{
		pdftitle={Aprendizagem de máquina para predição de desfechos em saúde mental: uma introdução para psicólogos}, 
		pdfauthor={Guilherme Augusto Dias},
    	pdfsubject={\imprimirpreambulo},
	    pdfcreator={LaTeX with abnTeX2},
		pdfkeywords={aprendizagem de máquina}{psicoterapia}{psicologia positiva}{depressão}, 
		colorlinks=true,		% true: colored links; false: boxed links
    	linkcolor=black,		% color of internal links
    	citecolor=black,		% color of links to bibliography
    	filecolor=black,		% color of file links
		urlcolor=black,
		bookmarksdepth=4
}
\makeatother
% --- 


% ---
% Posiciona figuras e tabelas no topo da página quando adicionadas sozinhas
% em um página em branco. Ver https://github.com/abntex/abntex2/issues/170
\makeatletter
\setlength{\@fptop}{5pt} % Set distance from top of page to first float
\makeatother
% ---

% ---
% Possibilita criação de Quadros e Lista de quadros.
% Ver https://github.com/abntex/abntex2/issues/176
%
\newcommand{\quadroname}{Quadro}
\newcommand{\listofquadrosname}{Lista de quadros}

\newfloat[chapter]{quadro}{loq}{\quadroname}
\newlistof{listofquadros}{loq}{\listofquadrosname}
\newlistentry{quadro}{loq}{0}

% Configurações para atender às regras da ABNT
\setfloatadjustment{quadro}{\centering}
\counterwithout{quadro}{chapter}
\renewcommand{\cftquadroname}{\quadroname\space} 
\renewcommand*{\cftquadroaftersnum}{\hfill--\hfill}

\setfloatlocations{quadro}{hbtp} % Ver https://github.com/abntex/abntex2/issues/176
% ---


% --- 
% Espaçamentos entre linhas e parágrafos 
% --- 

% O tamanho do parágrafo é dado por:
\setlength{\parindent}{1.3cm}

% Controle do espaçamento entre um parágrafo e outro:
\setlength{\parskip}{0.2cm}  % tente também \onelineskip

% ---
% Compila o índice
% ---
\makeindex
% ---


% ----
% Início do documento
% ----
\begin{document}

% Seleciona o idioma do documento (conforme pacotes do babel)
\selectlanguage{brazil}

% Retira espaço extra obsoleto entre as frases.
\frenchspacing 

% ----------------------------------------------------------
% ELEMENTOS PRÉ-TEXTUAIS
% ----------------------------------------------------------
\pretextual

% ---
% CAPA
% ---
\includepdf{./05-extras/capa.pdf}
% ---


% ---
% FOLHA DE ROSTO
% (o * indica que haverá a ficha bibliográfica)
% ---
\imprimirfolhaderosto
% ---


% ---
% Ficha bibliografica
% ---

% Isto é um exemplo de Ficha Catalográfica, ou ``Dados internacionais de
% catalogação-na-publicação''. Você pode utilizar este modelo como referência. 
% Porém, provavelmente a biblioteca da sua universidade lhe fornecerá um PDF
% com a ficha catalográfica definitiva após a defesa do trabalho. Quando estiver
% com o documento, salve-o como PDF no diretório do seu projeto e substitua todo
% o conteúdo de implementação deste arquivo pelo comando abaixo:
%
% \usepackage{pdfpages}		% necessário para comando \includepdf
% \begin{fichacatalografica}
%     \includepdf{fig_ficha_catalografica.pdf}
% \end{fichacatalografica}

% \begin{fichacatalografica}
% 	\sffamily
% 	\vspace*{\fill}                             % Posição vertical
% 	\begin{center}                              % Minipage Centralizado
% 	\fbox{\begin{minipage}[c][8cm]{13.5cm}		% Largura
% 	\small
% 	\imprimirautor
% 	%Sobrenome, Nome do autor
	
% 	\hspace{0.5cm} \imprimirtitulo  / \imprimirautor. --
% 	\imprimirlocal, \imprimirdata-
	
% 	\hspace{0.5cm} \thelastpage p. : il. (algumas color.) ; 30 cm.\\
	
% 	\hspace{0.5cm} \imprimirorientadorRotulo~\imprimirorientador\\
	
% 	\hspace{0.5cm}
% 	\parbox[t]{\textwidth}{\imprimirtipotrabalho~--~\imprimirinstituicao,
% 	\imprimirdata.}\\
	
% 	\hspace{0.5cm}
% 		1. Palavra-chave1.
% 		2. Palavra-chave2.
% 		2. Palavra-chave3.
% 		I. Orientador.
% 		II. Universidade xxx.
% 		III. Faculdade de xxx.
% 		IV. Título 			
% 	\end{minipage}}
% 	\end{center}
% \end{fichacatalografica}
% ---


% ---
% FOLHA DE APROVAÇÃO
% ---

% Isto é um exemplo de Folha de aprovação, elemento obrigatório da NBR
% 14724/2011 (seção 4.2.1.3). Você pode utilizar este modelo até a aprovação
% do trabalho. Após isso, substitua todo o conteúdo deste arquivo por uma
% imagem da página assinada pela banca com o comando abaixo:
%
% \begin{folhadeaprovacao}
% \includepdf{folhadeaprovacao_final.pdf}
% \end{folhadeaprovacao}
%
\setlength{\ABNTEXsignskip}{2cm}
\setlength{\ABNTEXsignwidth}{10cm}

\begin{folhadeaprovacao}
  \begin{center}
    {\ABNTEXchapterfont\large\imprimirautor}

    \vspace*{1cm}

    \begin{center}
      \ABNTEXchapterfont\bfseries\Large\imprimirtitulo
    \end{center}

    \vspace*{1cm}
    
    \hspace{.45\textwidth}
    \begin{minipage}{.5\textwidth}
        \imprimirpreambulo
    \end{minipage}%
   \end{center}

   \vspace*{1cm}
        
   \begin{center}
   Aprovado em 17 de julho de 2024.

   \vspace*{1cm}

   Banca examinadora:
   \end{center}

   \vspace*{-1cm}

   \assinatura{Prof. Ma. Caroline Santa Maria Rodrigues -- PUCRS}
   \assinatura{Prof. Dra. Clarissa Pinto Pizarro de Freitas -- PUCRS}
   \assinatura{\imprimirorientador -- PUCRS}
\end{folhadeaprovacao}
% ---


% ---
% Agradecimentos
% ---
\begin{agradecimentos}
Agradeço aos meus pais, Ademilson José Dias e Irinês Porfírio Dias, pelo apoio incondicional que me deram ao longo da vida,
incentivando mesmo as ideias questionáveis como aquela de fazer uma segunda graduação.

A minha irmã, Maria Eduarda Dias, cuja trajetória acadêmica é fonte de inspiração.

A minha namorada, Nathalia Fallavena Ceratti, que acompanhou a esta jornada de perto, oferecendo apoio e compreensão que
fizeram toda a diferença.

Aos meus amigos, Guilherme da Silva Petró, Norton Neves e Uilinton Luciano Marques da Silva, que foram compreensíveis nos
finais de semana em que não pude jogar RPG.

Ao professor Wagner de Lara Machado pela orientação e apoio no desenvolvimento deste trabalho.
\end{agradecimentos}
% ---


% ---
% RESUMOS
% ---
% Resumo em português
\setlength{\absparsep}{18pt} % ajusta o espaçamento dos parágrafos do resumo
\begin{resumo}
Desenvolvimentos na área de aprendizagem de máquina representam uma série de oportunidades para diversas áreas de pesquisa e de atuação profissional.
O uso de modelos de aprendizagem de máquina na predição de desfecho para tratamentos em saúde mental tem sido investigado e apresenta resultados promissores.
A aplicação desse tipo de modelo possibilitaria, por exemplo, maior assertividade na personalização de tratamentos psicoterápicos a nível individual.
Apesar dos benefícios potenciais, o tema ainda é pouco discutido nos espaços de pesquisa, formação e atuação em psicologia clínica. Este trabalho tem
por objetivo introduzir o tema de aprendizagem de máquina a psicólogos. Busca-se apresentar, de maneira acessível, os conceitos e terminologia básicos
em aprendizagem de máquina, proporcionando um panorama geral da área. Para complementar a exposição conceitual, é construído um exemplo a partir de dados
de um estudo sobre eficácia de intervenções em psicologia positiva para depressão via internet. Por fim, discutem-se algumas implicações do uso desse tipo
de tecnologia para a pesquisa e para a prática em psicoterapia.

 \textbf{Palavras-chave}: aprendizagem de máquina, psicoterapia, psicologia positiva, depressão.
\end{resumo}

% Resumo em inglês
\begin{resumo}[Abstract]
 \begin{otherlanguage*}{english}
   Advances in machine learning have opened up a multitude of opportunities for various research and professional fields. The use of machine learning models to
   predict treatment outcomes in mental health has been investigated showing promising results. The application of such models would allow greater assertiveness
   in tailoring psychotherapeutic treatments at the individual level. Despite its potential benefits, the topic remains unexplored in the fields of research, training
   and practice in clinical psychology. This work aims to introduce the topic of machine learning to psychologists. It presents the basic machine learning concepts and
   terminology in an accessible manner, providing a comprehensive overview of the field. To complement the conceptual discussion, an example is constructed using data
   from a study on the effectiveness of internet-based positive psychology interventions for depression. Finally, implications of using this type of technology in
   psychotherapy research and practice are discussed.
   \vspace{\onelineskip}
 
   \noindent 
   \textbf{Keywords}: machine learning, psychotherapy, positive psychology, depression.
 \end{otherlanguage*}
\end{resumo}


% ---
% LISTA DE ILUSTRAÇÕES
% ---
\pdfbookmark[0]{\listfigurename}{lof}
\listoffigures*
\cleardoublepage
% ---


% ---
% LISTA DE QUADROS
% ---
% \pdfbookmark[0]{\listofquadrosname}{loq}
% \listofquadros*
% \cleardoublepage
% ---


% ---
% LISTA DE TABELAS
% ---
% \pdfbookmark[0]{\listtablename}{lot}
% \listoftables*
% \cleardoublepage
% ---


% ---
% LISTA DE ABREVIATURAS E SIGLAS
% ---
% \begin{siglas}
%   \item[ABNT] Associação Brasileira de Normas Técnicas
%   \item[abnTeX] ABsurdas Normas para TeX
% \end{siglas}
% ---


% ---
% LISTA DE SÍMBOLOS
% ---
% \begin{simbolos}
%   \item[$ \Gamma $] Letra grega Gama
%   \item[$ \Lambda $] Lambda
%   \item[$ \zeta $] Letra grega minúscula zeta
%   \item[$ \in $] Pertence
% \end{simbolos}
% ---


% ---
% Sumário
% ---
\pdfbookmark[0]{\contentsname}{toc}
\tableofcontents*
\cleardoublepage
% ---


% ----------------------------------------------------------
% ELEMENTOS TEXTUAIS
% ----------------------------------------------------------
\textual

\include{./01-introducao/main.tex}

\include{./02-machine-learning/main.tex}

\include{./03-exemplo/main.tex}

\include{./04-conclusao/main.tex}

% ----------------------------------------------------------
% ELEMENTOS PÓS-TEXTUAIS
% ----------------------------------------------------------
\postextual
% ----------------------------------------------------------


% ----------------------------------------------------------
% REFERÊNCIAS
% ----------------------------------------------------------
\bibliography{references.bib}


% ----------------------------------------------------------
% GLOSSÁRIO
% ----------------------------------------------------------
%
% Consulte o manual da classe abntex2 para orientações sobre o glossário.
%
%\glossary

% ----------------------------------------------------------
% APÊNDICES
% ----------------------------------------------------------

% ---
% Inicia os apêndices
% ---
% \begin{apendicesenv}

% % Imprime uma página indicando o início dos apêndices
% \partapendices

% % ----------------------------------------------------------
% \chapter{Quisque libero justo}
% % ----------------------------------------------------------

% % ----------------------------------------------------------
% \chapter{Nullam elementum urna}
% % ----------------------------------------------------------

% \end{apendicesenv}
% ---


% ----------------------------------------------------------
% ANEXOS
% ----------------------------------------------------------

% ---
% Inicia os anexos
% ---
% \begin{anexosenv}

% % Imprime uma página indicando o início dos anexos
% \partanexos

% % ---
% \chapter{Morbi ultrices rutrum lorem}
% % ---

% % ---
% \chapter{Cras non urna sed feugiat cum}
% % ---


% % ---
% \chapter{Fusce facilisis lacinia dui}
% % ---

% \end{anexosenv}

%---------------------------------------------------------------------
% INDICE REMISSIVO
%---------------------------------------------------------------------
% \phantompart
% \printindex
%---------------------------------------------------------------------
\includepdf{./05-extras/contracapa.pdf}
\end{document}


\documentclass[
	% Opções da classe memoir
	12pt,				        % tamanho da fonte
	openany,			        % capítulos começam em pág ímpar (insere página vazia caso preciso)
	oneside,			        % para impressão em recto e verso. Oposto a oneside
	a4paper,			        % tamanho do papel. 
	% Opções da classe abntex2
	%chapter=TITLE,		        % títulos de capítulos convertidos em letras maiúsculas
	%section=TITLE,		        % títulos de seções convertidos em letras maiúsculas
	%subsection=TITLE,	        % títulos de subseções convertidos em letras maiúsculas
	%subsubsection=TITLE,       % títulos de subsubseções convertidos em letras maiúsculas
	% Opções do pacote babel
	english,			        % idioma adicional para hifenização
	brazil				        % o último idioma é o principal do documento
	]{abntex2}

% ---
% Pacotes básicos 
% ---
\usepackage{lmodern}			% Usa a fonte Latin Modern			
\usepackage[T1]{fontenc}		% Selecao de codigos de fonte.
\usepackage[utf8]{inputenc}		% Codificacao do documento (conversão automática dos acentos)
\usepackage{indentfirst}		% Indenta o primeiro parágrafo de cada seção.
\usepackage{color}				% Controle das cores
\usepackage{graphicx}			% Inclusão de gráficos
\usepackage{microtype} 			% para melhorias de justificação
% ---
		

% ---
% Pacotes de citações
% ---
% \usepackage[alf]{abntex2cite}						% Citações padrão ABNT
\usepackage[
  alf,
  abnt-nbr10520=1988,                     % Citações em sentence case
  abnt-emphasize=bf
  ]{abntex2cite}


% ---
% Pacotes adicionais
% ---
\usepackage{pdfpages}	% Inclusão de páginas em PDF


% ---
% Informações para CAPA e FOLHA DE ROSTO
% ---
\titulo{APRENDIZAGEM DE MÁQUINA PARA PREDIÇÃO DE DESFECHOS EM SAÚDE MENTAL: \textmd{UMA INTRODUÇÃO PARA PSICÓLOGOS}}
\autor{GUILHERME AUGUSTO DIAS}
\local{Porto Alegre}
\data{2024}
\orientador{Prof. Dr. Wagner de Lara Machado}
% \coorientador{}
% \instituicao{
%   Pontifícia Universidade Católica do Rio Grande do Sul
%   \par
%   Escola de Ciências da Saúde e da Vida
%   \par
%   Curso de Psicologia}
\tipotrabalho{TCC}
\preambulo{Trabalho de Conclusão de Curso apresentado como requisito para a obtenção do grau de Bacharel em Psicologia pela Escola de Ciências da Saúde e da Vida da Pontifícia Universidade Católica do Rio Grande do Sul.}
% ---

% ---
% Configurações de aparência do PDF final
% ---

% Alterando o aspecto da cor azul
\definecolor{blue}{RGB}{41,5,195}

% Informações do PDF
\makeatletter
\hypersetup{
		pdftitle={Aprendizagem de máquina para predição de desfechos em saúde mental: uma introdução para psicólogos}, 
		pdfauthor={Guilherme Augusto Dias},
    	pdfsubject={\imprimirpreambulo},
	    pdfcreator={LaTeX with abnTeX2},
		pdfkeywords={aprendizagem de máquina}{psicoterapia}{psicologia positiva}{depressão}, 
		colorlinks=true,		% true: colored links; false: boxed links
    	linkcolor=black,		% color of internal links
    	citecolor=black,		% color of links to bibliography
    	filecolor=black,		% color of file links
		urlcolor=black,
		bookmarksdepth=4
}
\makeatother
% --- 


% ---
% Posiciona figuras e tabelas no topo da página quando adicionadas sozinhas
% em um página em branco. Ver https://github.com/abntex/abntex2/issues/170
\makeatletter
\setlength{\@fptop}{5pt} % Set distance from top of page to first float
\makeatother
% ---

% ---
% Possibilita criação de Quadros e Lista de quadros.
% Ver https://github.com/abntex/abntex2/issues/176
%
\newcommand{\quadroname}{Quadro}
\newcommand{\listofquadrosname}{Lista de quadros}

\newfloat[chapter]{quadro}{loq}{\quadroname}
\newlistof{listofquadros}{loq}{\listofquadrosname}
\newlistentry{quadro}{loq}{0}

% Configurações para atender às regras da ABNT
\setfloatadjustment{quadro}{\centering}
\counterwithout{quadro}{chapter}
\renewcommand{\cftquadroname}{\quadroname\space} 
\renewcommand*{\cftquadroaftersnum}{\hfill--\hfill}

\setfloatlocations{quadro}{hbtp} % Ver https://github.com/abntex/abntex2/issues/176
% ---


% --- 
% Espaçamentos entre linhas e parágrafos 
% --- 

% O tamanho do parágrafo é dado por:
\setlength{\parindent}{1.3cm}

% Controle do espaçamento entre um parágrafo e outro:
\setlength{\parskip}{0.2cm}  % tente também \onelineskip

% ---
% Compila o índice
% ---
\makeindex
% ---


% ----
% Início do documento
% ----
\begin{document}

% Seleciona o idioma do documento (conforme pacotes do babel)
\selectlanguage{brazil}

% Retira espaço extra obsoleto entre as frases.
\frenchspacing 

% ----------------------------------------------------------
% ELEMENTOS PRÉ-TEXTUAIS
% ----------------------------------------------------------
\pretextual

% ---
% CAPA
% ---
\includepdf{./05-extras/capa.pdf}
% ---


% ---
% FOLHA DE ROSTO
% (o * indica que haverá a ficha bibliográfica)
% ---
\imprimirfolhaderosto
% ---


% ---
% Ficha bibliografica
% ---

% Isto é um exemplo de Ficha Catalográfica, ou ``Dados internacionais de
% catalogação-na-publicação''. Você pode utilizar este modelo como referência. 
% Porém, provavelmente a biblioteca da sua universidade lhe fornecerá um PDF
% com a ficha catalográfica definitiva após a defesa do trabalho. Quando estiver
% com o documento, salve-o como PDF no diretório do seu projeto e substitua todo
% o conteúdo de implementação deste arquivo pelo comando abaixo:
%
% \usepackage{pdfpages}		% necessário para comando \includepdf
% \begin{fichacatalografica}
%     \includepdf{fig_ficha_catalografica.pdf}
% \end{fichacatalografica}

% \begin{fichacatalografica}
% 	\sffamily
% 	\vspace*{\fill}                             % Posição vertical
% 	\begin{center}                              % Minipage Centralizado
% 	\fbox{\begin{minipage}[c][8cm]{13.5cm}		% Largura
% 	\small
% 	\imprimirautor
% 	%Sobrenome, Nome do autor
	
% 	\hspace{0.5cm} \imprimirtitulo  / \imprimirautor. --
% 	\imprimirlocal, \imprimirdata-
	
% 	\hspace{0.5cm} \thelastpage p. : il. (algumas color.) ; 30 cm.\\
	
% 	\hspace{0.5cm} \imprimirorientadorRotulo~\imprimirorientador\\
	
% 	\hspace{0.5cm}
% 	\parbox[t]{\textwidth}{\imprimirtipotrabalho~--~\imprimirinstituicao,
% 	\imprimirdata.}\\
	
% 	\hspace{0.5cm}
% 		1. Palavra-chave1.
% 		2. Palavra-chave2.
% 		2. Palavra-chave3.
% 		I. Orientador.
% 		II. Universidade xxx.
% 		III. Faculdade de xxx.
% 		IV. Título 			
% 	\end{minipage}}
% 	\end{center}
% \end{fichacatalografica}
% ---


% ---
% FOLHA DE APROVAÇÃO
% ---

% Isto é um exemplo de Folha de aprovação, elemento obrigatório da NBR
% 14724/2011 (seção 4.2.1.3). Você pode utilizar este modelo até a aprovação
% do trabalho. Após isso, substitua todo o conteúdo deste arquivo por uma
% imagem da página assinada pela banca com o comando abaixo:
%
% \begin{folhadeaprovacao}
% \includepdf{folhadeaprovacao_final.pdf}
% \end{folhadeaprovacao}
%
\setlength{\ABNTEXsignskip}{2cm}
\setlength{\ABNTEXsignwidth}{10cm}

\begin{folhadeaprovacao}
  \begin{center}
    {\ABNTEXchapterfont\large\imprimirautor}

    \vspace*{1cm}

    \begin{center}
      \ABNTEXchapterfont\bfseries\Large\imprimirtitulo
    \end{center}

    \vspace*{1cm}
    
    \hspace{.45\textwidth}
    \begin{minipage}{.5\textwidth}
        \imprimirpreambulo
    \end{minipage}%
   \end{center}

   \vspace*{1cm}
        
   \begin{center}
   Aprovado em 17 de julho de 2024.

   \vspace*{1cm}

   Banca examinadora:
   \end{center}

   \vspace*{-1cm}

   \assinatura{Prof. Ma. Caroline Santa Maria Rodrigues -- PUCRS}
   \assinatura{Prof. Dra. Clarissa Pinto Pizarro de Freitas -- PUCRS}
   \assinatura{\imprimirorientador -- PUCRS}
\end{folhadeaprovacao}
% ---


% ---
% Agradecimentos
% ---
\begin{agradecimentos}
Agradeço aos meus pais, Ademilson José Dias e Irinês Porfírio Dias, pelo apoio incondicional que me deram ao longo da vida,
incentivando mesmo as ideias questionáveis como aquela de fazer uma segunda graduação.

A minha irmã, Maria Eduarda Dias, cuja trajetória acadêmica é fonte de inspiração.

A minha namorada, Nathalia Fallavena Ceratti, que acompanhou a esta jornada de perto, oferecendo apoio e compreensão que
fizeram toda a diferença.

Aos meus amigos, Guilherme da Silva Petró, Norton Neves e Uilinton Luciano Marques da Silva, que foram compreensíveis nos
finais de semana em que não pude jogar RPG.

Ao professor Wagner de Lara Machado pela orientação e apoio no desenvolvimento deste trabalho.
\end{agradecimentos}
% ---


% ---
% RESUMOS
% ---
% Resumo em português
\setlength{\absparsep}{18pt} % ajusta o espaçamento dos parágrafos do resumo
\begin{resumo}
Desenvolvimentos na área de aprendizagem de máquina representam uma série de oportunidades para diversas áreas de pesquisa e de atuação profissional.
O uso de modelos de aprendizagem de máquina na predição de desfecho para tratamentos em saúde mental tem sido investigado e apresenta resultados promissores.
A aplicação desse tipo de modelo possibilitaria, por exemplo, maior assertividade na personalização de tratamentos psicoterápicos a nível individual.
Apesar dos benefícios potenciais, o tema ainda é pouco discutido nos espaços de pesquisa, formação e atuação em psicologia clínica. Este trabalho tem
por objetivo introduzir o tema de aprendizagem de máquina a psicólogos. Busca-se apresentar, de maneira acessível, os conceitos e terminologia básicos
em aprendizagem de máquina, proporcionando um panorama geral da área. Para complementar a exposição conceitual, é construído um exemplo a partir de dados
de um estudo sobre eficácia de intervenções em psicologia positiva para depressão via internet. Por fim, discutem-se algumas implicações do uso desse tipo
de tecnologia para a pesquisa e para a prática em psicoterapia.

 \textbf{Palavras-chave}: aprendizagem de máquina, psicoterapia, psicologia positiva, depressão.
\end{resumo}

% Resumo em inglês
\begin{resumo}[Abstract]
 \begin{otherlanguage*}{english}
   Advances in machine learning have opened up a multitude of opportunities for various research and professional fields. The use of machine learning models to
   predict treatment outcomes in mental health has been investigated showing promising results. The application of such models would allow greater assertiveness
   in tailoring psychotherapeutic treatments at the individual level. Despite its potential benefits, the topic remains unexplored in the fields of research, training
   and practice in clinical psychology. This work aims to introduce the topic of machine learning to psychologists. It presents the basic machine learning concepts and
   terminology in an accessible manner, providing a comprehensive overview of the field. To complement the conceptual discussion, an example is constructed using data
   from a study on the effectiveness of internet-based positive psychology interventions for depression. Finally, implications of using this type of technology in
   psychotherapy research and practice are discussed.
   \vspace{\onelineskip}
 
   \noindent 
   \textbf{Keywords}: machine learning, psychotherapy, positive psychology, depression.
 \end{otherlanguage*}
\end{resumo}


% ---
% LISTA DE ILUSTRAÇÕES
% ---
\pdfbookmark[0]{\listfigurename}{lof}
\listoffigures*
\cleardoublepage
% ---


% ---
% LISTA DE QUADROS
% ---
% \pdfbookmark[0]{\listofquadrosname}{loq}
% \listofquadros*
% \cleardoublepage
% ---


% ---
% LISTA DE TABELAS
% ---
% \pdfbookmark[0]{\listtablename}{lot}
% \listoftables*
% \cleardoublepage
% ---


% ---
% LISTA DE ABREVIATURAS E SIGLAS
% ---
% \begin{siglas}
%   \item[ABNT] Associação Brasileira de Normas Técnicas
%   \item[abnTeX] ABsurdas Normas para TeX
% \end{siglas}
% ---


% ---
% LISTA DE SÍMBOLOS
% ---
% \begin{simbolos}
%   \item[$ \Gamma $] Letra grega Gama
%   \item[$ \Lambda $] Lambda
%   \item[$ \zeta $] Letra grega minúscula zeta
%   \item[$ \in $] Pertence
% \end{simbolos}
% ---


% ---
% Sumário
% ---
\pdfbookmark[0]{\contentsname}{toc}
\tableofcontents*
\cleardoublepage
% ---


% ----------------------------------------------------------
% ELEMENTOS TEXTUAIS
% ----------------------------------------------------------
\textual

\include{./01-introducao/main.tex}

\include{./02-machine-learning/main.tex}

\include{./03-exemplo/main.tex}

\include{./04-conclusao/main.tex}

% ----------------------------------------------------------
% ELEMENTOS PÓS-TEXTUAIS
% ----------------------------------------------------------
\postextual
% ----------------------------------------------------------


% ----------------------------------------------------------
% REFERÊNCIAS
% ----------------------------------------------------------
\bibliography{references.bib}


% ----------------------------------------------------------
% GLOSSÁRIO
% ----------------------------------------------------------
%
% Consulte o manual da classe abntex2 para orientações sobre o glossário.
%
%\glossary

% ----------------------------------------------------------
% APÊNDICES
% ----------------------------------------------------------

% ---
% Inicia os apêndices
% ---
% \begin{apendicesenv}

% % Imprime uma página indicando o início dos apêndices
% \partapendices

% % ----------------------------------------------------------
% \chapter{Quisque libero justo}
% % ----------------------------------------------------------

% % ----------------------------------------------------------
% \chapter{Nullam elementum urna}
% % ----------------------------------------------------------

% \end{apendicesenv}
% ---


% ----------------------------------------------------------
% ANEXOS
% ----------------------------------------------------------

% ---
% Inicia os anexos
% ---
% \begin{anexosenv}

% % Imprime uma página indicando o início dos anexos
% \partanexos

% % ---
% \chapter{Morbi ultrices rutrum lorem}
% % ---

% % ---
% \chapter{Cras non urna sed feugiat cum}
% % ---


% % ---
% \chapter{Fusce facilisis lacinia dui}
% % ---

% \end{anexosenv}

%---------------------------------------------------------------------
% INDICE REMISSIVO
%---------------------------------------------------------------------
% \phantompart
% \printindex
%---------------------------------------------------------------------
\includepdf{./05-extras/contracapa.pdf}
\end{document}


\documentclass[
	% Opções da classe memoir
	12pt,				        % tamanho da fonte
	openany,			        % capítulos começam em pág ímpar (insere página vazia caso preciso)
	oneside,			        % para impressão em recto e verso. Oposto a oneside
	a4paper,			        % tamanho do papel. 
	% Opções da classe abntex2
	%chapter=TITLE,		        % títulos de capítulos convertidos em letras maiúsculas
	%section=TITLE,		        % títulos de seções convertidos em letras maiúsculas
	%subsection=TITLE,	        % títulos de subseções convertidos em letras maiúsculas
	%subsubsection=TITLE,       % títulos de subsubseções convertidos em letras maiúsculas
	% Opções do pacote babel
	english,			        % idioma adicional para hifenização
	brazil				        % o último idioma é o principal do documento
	]{abntex2}

% ---
% Pacotes básicos 
% ---
\usepackage{lmodern}			% Usa a fonte Latin Modern			
\usepackage[T1]{fontenc}		% Selecao de codigos de fonte.
\usepackage[utf8]{inputenc}		% Codificacao do documento (conversão automática dos acentos)
\usepackage{indentfirst}		% Indenta o primeiro parágrafo de cada seção.
\usepackage{color}				% Controle das cores
\usepackage{graphicx}			% Inclusão de gráficos
\usepackage{microtype} 			% para melhorias de justificação
% ---
		

% ---
% Pacotes de citações
% ---
% \usepackage[alf]{abntex2cite}						% Citações padrão ABNT
\usepackage[
  alf,
  abnt-nbr10520=1988,                     % Citações em sentence case
  abnt-emphasize=bf
  ]{abntex2cite}


% ---
% Pacotes adicionais
% ---
\usepackage{pdfpages}	% Inclusão de páginas em PDF


% ---
% Informações para CAPA e FOLHA DE ROSTO
% ---
\titulo{APRENDIZAGEM DE MÁQUINA PARA PREDIÇÃO DE DESFECHOS EM SAÚDE MENTAL: \textmd{UMA INTRODUÇÃO PARA PSICÓLOGOS}}
\autor{GUILHERME AUGUSTO DIAS}
\local{Porto Alegre}
\data{2024}
\orientador{Prof. Dr. Wagner de Lara Machado}
% \coorientador{}
% \instituicao{
%   Pontifícia Universidade Católica do Rio Grande do Sul
%   \par
%   Escola de Ciências da Saúde e da Vida
%   \par
%   Curso de Psicologia}
\tipotrabalho{TCC}
\preambulo{Trabalho de Conclusão de Curso apresentado como requisito para a obtenção do grau de Bacharel em Psicologia pela Escola de Ciências da Saúde e da Vida da Pontifícia Universidade Católica do Rio Grande do Sul.}
% ---

% ---
% Configurações de aparência do PDF final
% ---

% Alterando o aspecto da cor azul
\definecolor{blue}{RGB}{41,5,195}

% Informações do PDF
\makeatletter
\hypersetup{
		pdftitle={Aprendizagem de máquina para predição de desfechos em saúde mental: uma introdução para psicólogos}, 
		pdfauthor={Guilherme Augusto Dias},
    	pdfsubject={\imprimirpreambulo},
	    pdfcreator={LaTeX with abnTeX2},
		pdfkeywords={aprendizagem de máquina}{psicoterapia}{psicologia positiva}{depressão}, 
		colorlinks=true,		% true: colored links; false: boxed links
    	linkcolor=black,		% color of internal links
    	citecolor=black,		% color of links to bibliography
    	filecolor=black,		% color of file links
		urlcolor=black,
		bookmarksdepth=4
}
\makeatother
% --- 


% ---
% Posiciona figuras e tabelas no topo da página quando adicionadas sozinhas
% em um página em branco. Ver https://github.com/abntex/abntex2/issues/170
\makeatletter
\setlength{\@fptop}{5pt} % Set distance from top of page to first float
\makeatother
% ---

% ---
% Possibilita criação de Quadros e Lista de quadros.
% Ver https://github.com/abntex/abntex2/issues/176
%
\newcommand{\quadroname}{Quadro}
\newcommand{\listofquadrosname}{Lista de quadros}

\newfloat[chapter]{quadro}{loq}{\quadroname}
\newlistof{listofquadros}{loq}{\listofquadrosname}
\newlistentry{quadro}{loq}{0}

% Configurações para atender às regras da ABNT
\setfloatadjustment{quadro}{\centering}
\counterwithout{quadro}{chapter}
\renewcommand{\cftquadroname}{\quadroname\space} 
\renewcommand*{\cftquadroaftersnum}{\hfill--\hfill}

\setfloatlocations{quadro}{hbtp} % Ver https://github.com/abntex/abntex2/issues/176
% ---


% --- 
% Espaçamentos entre linhas e parágrafos 
% --- 

% O tamanho do parágrafo é dado por:
\setlength{\parindent}{1.3cm}

% Controle do espaçamento entre um parágrafo e outro:
\setlength{\parskip}{0.2cm}  % tente também \onelineskip

% ---
% Compila o índice
% ---
\makeindex
% ---


% ----
% Início do documento
% ----
\begin{document}

% Seleciona o idioma do documento (conforme pacotes do babel)
\selectlanguage{brazil}

% Retira espaço extra obsoleto entre as frases.
\frenchspacing 

% ----------------------------------------------------------
% ELEMENTOS PRÉ-TEXTUAIS
% ----------------------------------------------------------
\pretextual

% ---
% CAPA
% ---
\includepdf{./05-extras/capa.pdf}
% ---


% ---
% FOLHA DE ROSTO
% (o * indica que haverá a ficha bibliográfica)
% ---
\imprimirfolhaderosto
% ---


% ---
% Ficha bibliografica
% ---

% Isto é um exemplo de Ficha Catalográfica, ou ``Dados internacionais de
% catalogação-na-publicação''. Você pode utilizar este modelo como referência. 
% Porém, provavelmente a biblioteca da sua universidade lhe fornecerá um PDF
% com a ficha catalográfica definitiva após a defesa do trabalho. Quando estiver
% com o documento, salve-o como PDF no diretório do seu projeto e substitua todo
% o conteúdo de implementação deste arquivo pelo comando abaixo:
%
% \usepackage{pdfpages}		% necessário para comando \includepdf
% \begin{fichacatalografica}
%     \includepdf{fig_ficha_catalografica.pdf}
% \end{fichacatalografica}

% \begin{fichacatalografica}
% 	\sffamily
% 	\vspace*{\fill}                             % Posição vertical
% 	\begin{center}                              % Minipage Centralizado
% 	\fbox{\begin{minipage}[c][8cm]{13.5cm}		% Largura
% 	\small
% 	\imprimirautor
% 	%Sobrenome, Nome do autor
	
% 	\hspace{0.5cm} \imprimirtitulo  / \imprimirautor. --
% 	\imprimirlocal, \imprimirdata-
	
% 	\hspace{0.5cm} \thelastpage p. : il. (algumas color.) ; 30 cm.\\
	
% 	\hspace{0.5cm} \imprimirorientadorRotulo~\imprimirorientador\\
	
% 	\hspace{0.5cm}
% 	\parbox[t]{\textwidth}{\imprimirtipotrabalho~--~\imprimirinstituicao,
% 	\imprimirdata.}\\
	
% 	\hspace{0.5cm}
% 		1. Palavra-chave1.
% 		2. Palavra-chave2.
% 		2. Palavra-chave3.
% 		I. Orientador.
% 		II. Universidade xxx.
% 		III. Faculdade de xxx.
% 		IV. Título 			
% 	\end{minipage}}
% 	\end{center}
% \end{fichacatalografica}
% ---


% ---
% FOLHA DE APROVAÇÃO
% ---

% Isto é um exemplo de Folha de aprovação, elemento obrigatório da NBR
% 14724/2011 (seção 4.2.1.3). Você pode utilizar este modelo até a aprovação
% do trabalho. Após isso, substitua todo o conteúdo deste arquivo por uma
% imagem da página assinada pela banca com o comando abaixo:
%
% \begin{folhadeaprovacao}
% \includepdf{folhadeaprovacao_final.pdf}
% \end{folhadeaprovacao}
%
\setlength{\ABNTEXsignskip}{2cm}
\setlength{\ABNTEXsignwidth}{10cm}

\begin{folhadeaprovacao}
  \begin{center}
    {\ABNTEXchapterfont\large\imprimirautor}

    \vspace*{1cm}

    \begin{center}
      \ABNTEXchapterfont\bfseries\Large\imprimirtitulo
    \end{center}

    \vspace*{1cm}
    
    \hspace{.45\textwidth}
    \begin{minipage}{.5\textwidth}
        \imprimirpreambulo
    \end{minipage}%
   \end{center}

   \vspace*{1cm}
        
   \begin{center}
   Aprovado em 17 de julho de 2024.

   \vspace*{1cm}

   Banca examinadora:
   \end{center}

   \vspace*{-1cm}

   \assinatura{Prof. Ma. Caroline Santa Maria Rodrigues -- PUCRS}
   \assinatura{Prof. Dra. Clarissa Pinto Pizarro de Freitas -- PUCRS}
   \assinatura{\imprimirorientador -- PUCRS}
\end{folhadeaprovacao}
% ---


% ---
% Agradecimentos
% ---
\begin{agradecimentos}
Agradeço aos meus pais, Ademilson José Dias e Irinês Porfírio Dias, pelo apoio incondicional que me deram ao longo da vida,
incentivando mesmo as ideias questionáveis como aquela de fazer uma segunda graduação.

A minha irmã, Maria Eduarda Dias, cuja trajetória acadêmica é fonte de inspiração.

A minha namorada, Nathalia Fallavena Ceratti, que acompanhou a esta jornada de perto, oferecendo apoio e compreensão que
fizeram toda a diferença.

Aos meus amigos, Guilherme da Silva Petró, Norton Neves e Uilinton Luciano Marques da Silva, que foram compreensíveis nos
finais de semana em que não pude jogar RPG.

Ao professor Wagner de Lara Machado pela orientação e apoio no desenvolvimento deste trabalho.
\end{agradecimentos}
% ---


% ---
% RESUMOS
% ---
% Resumo em português
\setlength{\absparsep}{18pt} % ajusta o espaçamento dos parágrafos do resumo
\begin{resumo}
Desenvolvimentos na área de aprendizagem de máquina representam uma série de oportunidades para diversas áreas de pesquisa e de atuação profissional.
O uso de modelos de aprendizagem de máquina na predição de desfecho para tratamentos em saúde mental tem sido investigado e apresenta resultados promissores.
A aplicação desse tipo de modelo possibilitaria, por exemplo, maior assertividade na personalização de tratamentos psicoterápicos a nível individual.
Apesar dos benefícios potenciais, o tema ainda é pouco discutido nos espaços de pesquisa, formação e atuação em psicologia clínica. Este trabalho tem
por objetivo introduzir o tema de aprendizagem de máquina a psicólogos. Busca-se apresentar, de maneira acessível, os conceitos e terminologia básicos
em aprendizagem de máquina, proporcionando um panorama geral da área. Para complementar a exposição conceitual, é construído um exemplo a partir de dados
de um estudo sobre eficácia de intervenções em psicologia positiva para depressão via internet. Por fim, discutem-se algumas implicações do uso desse tipo
de tecnologia para a pesquisa e para a prática em psicoterapia.

 \textbf{Palavras-chave}: aprendizagem de máquina, psicoterapia, psicologia positiva, depressão.
\end{resumo}

% Resumo em inglês
\begin{resumo}[Abstract]
 \begin{otherlanguage*}{english}
   Advances in machine learning have opened up a multitude of opportunities for various research and professional fields. The use of machine learning models to
   predict treatment outcomes in mental health has been investigated showing promising results. The application of such models would allow greater assertiveness
   in tailoring psychotherapeutic treatments at the individual level. Despite its potential benefits, the topic remains unexplored in the fields of research, training
   and practice in clinical psychology. This work aims to introduce the topic of machine learning to psychologists. It presents the basic machine learning concepts and
   terminology in an accessible manner, providing a comprehensive overview of the field. To complement the conceptual discussion, an example is constructed using data
   from a study on the effectiveness of internet-based positive psychology interventions for depression. Finally, implications of using this type of technology in
   psychotherapy research and practice are discussed.
   \vspace{\onelineskip}
 
   \noindent 
   \textbf{Keywords}: machine learning, psychotherapy, positive psychology, depression.
 \end{otherlanguage*}
\end{resumo}


% ---
% LISTA DE ILUSTRAÇÕES
% ---
\pdfbookmark[0]{\listfigurename}{lof}
\listoffigures*
\cleardoublepage
% ---


% ---
% LISTA DE QUADROS
% ---
% \pdfbookmark[0]{\listofquadrosname}{loq}
% \listofquadros*
% \cleardoublepage
% ---


% ---
% LISTA DE TABELAS
% ---
% \pdfbookmark[0]{\listtablename}{lot}
% \listoftables*
% \cleardoublepage
% ---


% ---
% LISTA DE ABREVIATURAS E SIGLAS
% ---
% \begin{siglas}
%   \item[ABNT] Associação Brasileira de Normas Técnicas
%   \item[abnTeX] ABsurdas Normas para TeX
% \end{siglas}
% ---


% ---
% LISTA DE SÍMBOLOS
% ---
% \begin{simbolos}
%   \item[$ \Gamma $] Letra grega Gama
%   \item[$ \Lambda $] Lambda
%   \item[$ \zeta $] Letra grega minúscula zeta
%   \item[$ \in $] Pertence
% \end{simbolos}
% ---


% ---
% Sumário
% ---
\pdfbookmark[0]{\contentsname}{toc}
\tableofcontents*
\cleardoublepage
% ---


% ----------------------------------------------------------
% ELEMENTOS TEXTUAIS
% ----------------------------------------------------------
\textual

\include{./01-introducao/main.tex}

\include{./02-machine-learning/main.tex}

\include{./03-exemplo/main.tex}

\include{./04-conclusao/main.tex}

% ----------------------------------------------------------
% ELEMENTOS PÓS-TEXTUAIS
% ----------------------------------------------------------
\postextual
% ----------------------------------------------------------


% ----------------------------------------------------------
% REFERÊNCIAS
% ----------------------------------------------------------
\bibliography{references.bib}


% ----------------------------------------------------------
% GLOSSÁRIO
% ----------------------------------------------------------
%
% Consulte o manual da classe abntex2 para orientações sobre o glossário.
%
%\glossary

% ----------------------------------------------------------
% APÊNDICES
% ----------------------------------------------------------

% ---
% Inicia os apêndices
% ---
% \begin{apendicesenv}

% % Imprime uma página indicando o início dos apêndices
% \partapendices

% % ----------------------------------------------------------
% \chapter{Quisque libero justo}
% % ----------------------------------------------------------

% % ----------------------------------------------------------
% \chapter{Nullam elementum urna}
% % ----------------------------------------------------------

% \end{apendicesenv}
% ---


% ----------------------------------------------------------
% ANEXOS
% ----------------------------------------------------------

% ---
% Inicia os anexos
% ---
% \begin{anexosenv}

% % Imprime uma página indicando o início dos anexos
% \partanexos

% % ---
% \chapter{Morbi ultrices rutrum lorem}
% % ---

% % ---
% \chapter{Cras non urna sed feugiat cum}
% % ---


% % ---
% \chapter{Fusce facilisis lacinia dui}
% % ---

% \end{anexosenv}

%---------------------------------------------------------------------
% INDICE REMISSIVO
%---------------------------------------------------------------------
% \phantompart
% \printindex
%---------------------------------------------------------------------
\includepdf{./05-extras/contracapa.pdf}
\end{document}


\documentclass[
	% Opções da classe memoir
	12pt,				        % tamanho da fonte
	openany,			        % capítulos começam em pág ímpar (insere página vazia caso preciso)
	oneside,			        % para impressão em recto e verso. Oposto a oneside
	a4paper,			        % tamanho do papel. 
	% Opções da classe abntex2
	%chapter=TITLE,		        % títulos de capítulos convertidos em letras maiúsculas
	%section=TITLE,		        % títulos de seções convertidos em letras maiúsculas
	%subsection=TITLE,	        % títulos de subseções convertidos em letras maiúsculas
	%subsubsection=TITLE,       % títulos de subsubseções convertidos em letras maiúsculas
	% Opções do pacote babel
	english,			        % idioma adicional para hifenização
	brazil				        % o último idioma é o principal do documento
	]{abntex2}

% ---
% Pacotes básicos 
% ---
\usepackage{lmodern}			% Usa a fonte Latin Modern			
\usepackage[T1]{fontenc}		% Selecao de codigos de fonte.
\usepackage[utf8]{inputenc}		% Codificacao do documento (conversão automática dos acentos)
\usepackage{indentfirst}		% Indenta o primeiro parágrafo de cada seção.
\usepackage{color}				% Controle das cores
\usepackage{graphicx}			% Inclusão de gráficos
\usepackage{microtype} 			% para melhorias de justificação
% ---
		

% ---
% Pacotes de citações
% ---
% \usepackage[alf]{abntex2cite}						% Citações padrão ABNT
\usepackage[
  alf,
  abnt-nbr10520=1988,                     % Citações em sentence case
  abnt-emphasize=bf
  ]{abntex2cite}


% ---
% Pacotes adicionais
% ---
\usepackage{pdfpages}	% Inclusão de páginas em PDF


% ---
% Informações para CAPA e FOLHA DE ROSTO
% ---
\titulo{APRENDIZAGEM DE MÁQUINA PARA PREDIÇÃO DE DESFECHOS EM SAÚDE MENTAL: \textmd{UMA INTRODUÇÃO PARA PSICÓLOGOS}}
\autor{GUILHERME AUGUSTO DIAS}
\local{Porto Alegre}
\data{2024}
\orientador{Prof. Dr. Wagner de Lara Machado}
% \coorientador{}
% \instituicao{
%   Pontifícia Universidade Católica do Rio Grande do Sul
%   \par
%   Escola de Ciências da Saúde e da Vida
%   \par
%   Curso de Psicologia}
\tipotrabalho{TCC}
\preambulo{Trabalho de Conclusão de Curso apresentado como requisito para a obtenção do grau de Bacharel em Psicologia pela Escola de Ciências da Saúde e da Vida da Pontifícia Universidade Católica do Rio Grande do Sul.}
% ---

% ---
% Configurações de aparência do PDF final
% ---

% Alterando o aspecto da cor azul
\definecolor{blue}{RGB}{41,5,195}

% Informações do PDF
\makeatletter
\hypersetup{
		pdftitle={Aprendizagem de máquina para predição de desfechos em saúde mental: uma introdução para psicólogos}, 
		pdfauthor={Guilherme Augusto Dias},
    	pdfsubject={\imprimirpreambulo},
	    pdfcreator={LaTeX with abnTeX2},
		pdfkeywords={aprendizagem de máquina}{psicoterapia}{psicologia positiva}{depressão}, 
		colorlinks=true,		% true: colored links; false: boxed links
    	linkcolor=black,		% color of internal links
    	citecolor=black,		% color of links to bibliography
    	filecolor=black,		% color of file links
		urlcolor=black,
		bookmarksdepth=4
}
\makeatother
% --- 


% ---
% Posiciona figuras e tabelas no topo da página quando adicionadas sozinhas
% em um página em branco. Ver https://github.com/abntex/abntex2/issues/170
\makeatletter
\setlength{\@fptop}{5pt} % Set distance from top of page to first float
\makeatother
% ---

% ---
% Possibilita criação de Quadros e Lista de quadros.
% Ver https://github.com/abntex/abntex2/issues/176
%
\newcommand{\quadroname}{Quadro}
\newcommand{\listofquadrosname}{Lista de quadros}

\newfloat[chapter]{quadro}{loq}{\quadroname}
\newlistof{listofquadros}{loq}{\listofquadrosname}
\newlistentry{quadro}{loq}{0}

% Configurações para atender às regras da ABNT
\setfloatadjustment{quadro}{\centering}
\counterwithout{quadro}{chapter}
\renewcommand{\cftquadroname}{\quadroname\space} 
\renewcommand*{\cftquadroaftersnum}{\hfill--\hfill}

\setfloatlocations{quadro}{hbtp} % Ver https://github.com/abntex/abntex2/issues/176
% ---


% --- 
% Espaçamentos entre linhas e parágrafos 
% --- 

% O tamanho do parágrafo é dado por:
\setlength{\parindent}{1.3cm}

% Controle do espaçamento entre um parágrafo e outro:
\setlength{\parskip}{0.2cm}  % tente também \onelineskip

% ---
% Compila o índice
% ---
\makeindex
% ---


% ----
% Início do documento
% ----
\begin{document}

% Seleciona o idioma do documento (conforme pacotes do babel)
\selectlanguage{brazil}

% Retira espaço extra obsoleto entre as frases.
\frenchspacing 

% ----------------------------------------------------------
% ELEMENTOS PRÉ-TEXTUAIS
% ----------------------------------------------------------
\pretextual

% ---
% CAPA
% ---
\includepdf{./05-extras/capa.pdf}
% ---


% ---
% FOLHA DE ROSTO
% (o * indica que haverá a ficha bibliográfica)
% ---
\imprimirfolhaderosto
% ---


% ---
% Ficha bibliografica
% ---

% Isto é um exemplo de Ficha Catalográfica, ou ``Dados internacionais de
% catalogação-na-publicação''. Você pode utilizar este modelo como referência. 
% Porém, provavelmente a biblioteca da sua universidade lhe fornecerá um PDF
% com a ficha catalográfica definitiva após a defesa do trabalho. Quando estiver
% com o documento, salve-o como PDF no diretório do seu projeto e substitua todo
% o conteúdo de implementação deste arquivo pelo comando abaixo:
%
% \usepackage{pdfpages}		% necessário para comando \includepdf
% \begin{fichacatalografica}
%     \includepdf{fig_ficha_catalografica.pdf}
% \end{fichacatalografica}

% \begin{fichacatalografica}
% 	\sffamily
% 	\vspace*{\fill}                             % Posição vertical
% 	\begin{center}                              % Minipage Centralizado
% 	\fbox{\begin{minipage}[c][8cm]{13.5cm}		% Largura
% 	\small
% 	\imprimirautor
% 	%Sobrenome, Nome do autor
	
% 	\hspace{0.5cm} \imprimirtitulo  / \imprimirautor. --
% 	\imprimirlocal, \imprimirdata-
	
% 	\hspace{0.5cm} \thelastpage p. : il. (algumas color.) ; 30 cm.\\
	
% 	\hspace{0.5cm} \imprimirorientadorRotulo~\imprimirorientador\\
	
% 	\hspace{0.5cm}
% 	\parbox[t]{\textwidth}{\imprimirtipotrabalho~--~\imprimirinstituicao,
% 	\imprimirdata.}\\
	
% 	\hspace{0.5cm}
% 		1. Palavra-chave1.
% 		2. Palavra-chave2.
% 		2. Palavra-chave3.
% 		I. Orientador.
% 		II. Universidade xxx.
% 		III. Faculdade de xxx.
% 		IV. Título 			
% 	\end{minipage}}
% 	\end{center}
% \end{fichacatalografica}
% ---


% ---
% FOLHA DE APROVAÇÃO
% ---

% Isto é um exemplo de Folha de aprovação, elemento obrigatório da NBR
% 14724/2011 (seção 4.2.1.3). Você pode utilizar este modelo até a aprovação
% do trabalho. Após isso, substitua todo o conteúdo deste arquivo por uma
% imagem da página assinada pela banca com o comando abaixo:
%
% \begin{folhadeaprovacao}
% \includepdf{folhadeaprovacao_final.pdf}
% \end{folhadeaprovacao}
%
\setlength{\ABNTEXsignskip}{2cm}
\setlength{\ABNTEXsignwidth}{10cm}

\begin{folhadeaprovacao}
  \begin{center}
    {\ABNTEXchapterfont\large\imprimirautor}

    \vspace*{1cm}

    \begin{center}
      \ABNTEXchapterfont\bfseries\Large\imprimirtitulo
    \end{center}

    \vspace*{1cm}
    
    \hspace{.45\textwidth}
    \begin{minipage}{.5\textwidth}
        \imprimirpreambulo
    \end{minipage}%
   \end{center}

   \vspace*{1cm}
        
   \begin{center}
   Aprovado em 17 de julho de 2024.

   \vspace*{1cm}

   Banca examinadora:
   \end{center}

   \vspace*{-1cm}

   \assinatura{Prof. Ma. Caroline Santa Maria Rodrigues -- PUCRS}
   \assinatura{Prof. Dra. Clarissa Pinto Pizarro de Freitas -- PUCRS}
   \assinatura{\imprimirorientador -- PUCRS}
\end{folhadeaprovacao}
% ---


% ---
% Agradecimentos
% ---
\begin{agradecimentos}
Agradeço aos meus pais, Ademilson José Dias e Irinês Porfírio Dias, pelo apoio incondicional que me deram ao longo da vida,
incentivando mesmo as ideias questionáveis como aquela de fazer uma segunda graduação.

A minha irmã, Maria Eduarda Dias, cuja trajetória acadêmica é fonte de inspiração.

A minha namorada, Nathalia Fallavena Ceratti, que acompanhou a esta jornada de perto, oferecendo apoio e compreensão que
fizeram toda a diferença.

Aos meus amigos, Guilherme da Silva Petró, Norton Neves e Uilinton Luciano Marques da Silva, que foram compreensíveis nos
finais de semana em que não pude jogar RPG.

Ao professor Wagner de Lara Machado pela orientação e apoio no desenvolvimento deste trabalho.
\end{agradecimentos}
% ---


% ---
% RESUMOS
% ---
% Resumo em português
\setlength{\absparsep}{18pt} % ajusta o espaçamento dos parágrafos do resumo
\begin{resumo}
Desenvolvimentos na área de aprendizagem de máquina representam uma série de oportunidades para diversas áreas de pesquisa e de atuação profissional.
O uso de modelos de aprendizagem de máquina na predição de desfecho para tratamentos em saúde mental tem sido investigado e apresenta resultados promissores.
A aplicação desse tipo de modelo possibilitaria, por exemplo, maior assertividade na personalização de tratamentos psicoterápicos a nível individual.
Apesar dos benefícios potenciais, o tema ainda é pouco discutido nos espaços de pesquisa, formação e atuação em psicologia clínica. Este trabalho tem
por objetivo introduzir o tema de aprendizagem de máquina a psicólogos. Busca-se apresentar, de maneira acessível, os conceitos e terminologia básicos
em aprendizagem de máquina, proporcionando um panorama geral da área. Para complementar a exposição conceitual, é construído um exemplo a partir de dados
de um estudo sobre eficácia de intervenções em psicologia positiva para depressão via internet. Por fim, discutem-se algumas implicações do uso desse tipo
de tecnologia para a pesquisa e para a prática em psicoterapia.

 \textbf{Palavras-chave}: aprendizagem de máquina, psicoterapia, psicologia positiva, depressão.
\end{resumo}

% Resumo em inglês
\begin{resumo}[Abstract]
 \begin{otherlanguage*}{english}
   Advances in machine learning have opened up a multitude of opportunities for various research and professional fields. The use of machine learning models to
   predict treatment outcomes in mental health has been investigated showing promising results. The application of such models would allow greater assertiveness
   in tailoring psychotherapeutic treatments at the individual level. Despite its potential benefits, the topic remains unexplored in the fields of research, training
   and practice in clinical psychology. This work aims to introduce the topic of machine learning to psychologists. It presents the basic machine learning concepts and
   terminology in an accessible manner, providing a comprehensive overview of the field. To complement the conceptual discussion, an example is constructed using data
   from a study on the effectiveness of internet-based positive psychology interventions for depression. Finally, implications of using this type of technology in
   psychotherapy research and practice are discussed.
   \vspace{\onelineskip}
 
   \noindent 
   \textbf{Keywords}: machine learning, psychotherapy, positive psychology, depression.
 \end{otherlanguage*}
\end{resumo}


% ---
% LISTA DE ILUSTRAÇÕES
% ---
\pdfbookmark[0]{\listfigurename}{lof}
\listoffigures*
\cleardoublepage
% ---


% ---
% LISTA DE QUADROS
% ---
% \pdfbookmark[0]{\listofquadrosname}{loq}
% \listofquadros*
% \cleardoublepage
% ---


% ---
% LISTA DE TABELAS
% ---
% \pdfbookmark[0]{\listtablename}{lot}
% \listoftables*
% \cleardoublepage
% ---


% ---
% LISTA DE ABREVIATURAS E SIGLAS
% ---
% \begin{siglas}
%   \item[ABNT] Associação Brasileira de Normas Técnicas
%   \item[abnTeX] ABsurdas Normas para TeX
% \end{siglas}
% ---


% ---
% LISTA DE SÍMBOLOS
% ---
% \begin{simbolos}
%   \item[$ \Gamma $] Letra grega Gama
%   \item[$ \Lambda $] Lambda
%   \item[$ \zeta $] Letra grega minúscula zeta
%   \item[$ \in $] Pertence
% \end{simbolos}
% ---


% ---
% Sumário
% ---
\pdfbookmark[0]{\contentsname}{toc}
\tableofcontents*
\cleardoublepage
% ---


% ----------------------------------------------------------
% ELEMENTOS TEXTUAIS
% ----------------------------------------------------------
\textual

\include{./01-introducao/main.tex}

\include{./02-machine-learning/main.tex}

\include{./03-exemplo/main.tex}

\include{./04-conclusao/main.tex}

% ----------------------------------------------------------
% ELEMENTOS PÓS-TEXTUAIS
% ----------------------------------------------------------
\postextual
% ----------------------------------------------------------


% ----------------------------------------------------------
% REFERÊNCIAS
% ----------------------------------------------------------
\bibliography{references.bib}


% ----------------------------------------------------------
% GLOSSÁRIO
% ----------------------------------------------------------
%
% Consulte o manual da classe abntex2 para orientações sobre o glossário.
%
%\glossary

% ----------------------------------------------------------
% APÊNDICES
% ----------------------------------------------------------

% ---
% Inicia os apêndices
% ---
% \begin{apendicesenv}

% % Imprime uma página indicando o início dos apêndices
% \partapendices

% % ----------------------------------------------------------
% \chapter{Quisque libero justo}
% % ----------------------------------------------------------

% % ----------------------------------------------------------
% \chapter{Nullam elementum urna}
% % ----------------------------------------------------------

% \end{apendicesenv}
% ---


% ----------------------------------------------------------
% ANEXOS
% ----------------------------------------------------------

% ---
% Inicia os anexos
% ---
% \begin{anexosenv}

% % Imprime uma página indicando o início dos anexos
% \partanexos

% % ---
% \chapter{Morbi ultrices rutrum lorem}
% % ---

% % ---
% \chapter{Cras non urna sed feugiat cum}
% % ---


% % ---
% \chapter{Fusce facilisis lacinia dui}
% % ---

% \end{anexosenv}

%---------------------------------------------------------------------
% INDICE REMISSIVO
%---------------------------------------------------------------------
% \phantompart
% \printindex
%---------------------------------------------------------------------
\includepdf{./05-extras/contracapa.pdf}
\end{document}


% ----------------------------------------------------------
% ELEMENTOS PÓS-TEXTUAIS
% ----------------------------------------------------------
\postextual
% ----------------------------------------------------------


% ----------------------------------------------------------
% REFERÊNCIAS
% ----------------------------------------------------------
\bibliography{references.bib}


% ----------------------------------------------------------
% GLOSSÁRIO
% ----------------------------------------------------------
%
% Consulte o manual da classe abntex2 para orientações sobre o glossário.
%
%\glossary

% ----------------------------------------------------------
% APÊNDICES
% ----------------------------------------------------------

% ---
% Inicia os apêndices
% ---
% \begin{apendicesenv}

% % Imprime uma página indicando o início dos apêndices
% \partapendices

% % ----------------------------------------------------------
% \chapter{Quisque libero justo}
% % ----------------------------------------------------------

% % ----------------------------------------------------------
% \chapter{Nullam elementum urna}
% % ----------------------------------------------------------

% \end{apendicesenv}
% ---


% ----------------------------------------------------------
% ANEXOS
% ----------------------------------------------------------

% ---
% Inicia os anexos
% ---
% \begin{anexosenv}

% % Imprime uma página indicando o início dos anexos
% \partanexos

% % ---
% \chapter{Morbi ultrices rutrum lorem}
% % ---

% % ---
% \chapter{Cras non urna sed feugiat cum}
% % ---


% % ---
% \chapter{Fusce facilisis lacinia dui}
% % ---

% \end{anexosenv}

%---------------------------------------------------------------------
% INDICE REMISSIVO
%---------------------------------------------------------------------
% \phantompart
% \printindex
%---------------------------------------------------------------------
\includepdf{./05-extras/contracapa.pdf}
\end{document}


\documentclass[
	% Opções da classe memoir
	12pt,				        % tamanho da fonte
	openany,			        % capítulos começam em pág ímpar (insere página vazia caso preciso)
	oneside,			        % para impressão em recto e verso. Oposto a oneside
	a4paper,			        % tamanho do papel. 
	% Opções da classe abntex2
	%chapter=TITLE,		        % títulos de capítulos convertidos em letras maiúsculas
	%section=TITLE,		        % títulos de seções convertidos em letras maiúsculas
	%subsection=TITLE,	        % títulos de subseções convertidos em letras maiúsculas
	%subsubsection=TITLE,       % títulos de subsubseções convertidos em letras maiúsculas
	% Opções do pacote babel
	english,			        % idioma adicional para hifenização
	brazil				        % o último idioma é o principal do documento
	]{abntex2}

% ---
% Pacotes básicos 
% ---
\usepackage{lmodern}			% Usa a fonte Latin Modern			
\usepackage[T1]{fontenc}		% Selecao de codigos de fonte.
\usepackage[utf8]{inputenc}		% Codificacao do documento (conversão automática dos acentos)
\usepackage{indentfirst}		% Indenta o primeiro parágrafo de cada seção.
\usepackage{color}				% Controle das cores
\usepackage{graphicx}			% Inclusão de gráficos
\usepackage{microtype} 			% para melhorias de justificação
% ---
		

% ---
% Pacotes de citações
% ---
% \usepackage[alf]{abntex2cite}						% Citações padrão ABNT
\usepackage[
  alf,
  abnt-nbr10520=1988,                     % Citações em sentence case
  abnt-emphasize=bf
  ]{abntex2cite}


% ---
% Pacotes adicionais
% ---
\usepackage{pdfpages}	% Inclusão de páginas em PDF


% ---
% Informações para CAPA e FOLHA DE ROSTO
% ---
\titulo{APRENDIZAGEM DE MÁQUINA PARA PREDIÇÃO DE DESFECHOS EM SAÚDE MENTAL: \textmd{UMA INTRODUÇÃO PARA PSICÓLOGOS}}
\autor{GUILHERME AUGUSTO DIAS}
\local{Porto Alegre}
\data{2024}
\orientador{Prof. Dr. Wagner de Lara Machado}
% \coorientador{}
% \instituicao{
%   Pontifícia Universidade Católica do Rio Grande do Sul
%   \par
%   Escola de Ciências da Saúde e da Vida
%   \par
%   Curso de Psicologia}
\tipotrabalho{TCC}
\preambulo{Trabalho de Conclusão de Curso apresentado como requisito para a obtenção do grau de Bacharel em Psicologia pela Escola de Ciências da Saúde e da Vida da Pontifícia Universidade Católica do Rio Grande do Sul.}
% ---

% ---
% Configurações de aparência do PDF final
% ---

% Alterando o aspecto da cor azul
\definecolor{blue}{RGB}{41,5,195}

% Informações do PDF
\makeatletter
\hypersetup{
		pdftitle={Aprendizagem de máquina para predição de desfechos em saúde mental: uma introdução para psicólogos}, 
		pdfauthor={Guilherme Augusto Dias},
    	pdfsubject={\imprimirpreambulo},
	    pdfcreator={LaTeX with abnTeX2},
		pdfkeywords={aprendizagem de máquina}{psicoterapia}{psicologia positiva}{depressão}, 
		colorlinks=true,		% true: colored links; false: boxed links
    	linkcolor=black,		% color of internal links
    	citecolor=black,		% color of links to bibliography
    	filecolor=black,		% color of file links
		urlcolor=black,
		bookmarksdepth=4
}
\makeatother
% --- 


% ---
% Posiciona figuras e tabelas no topo da página quando adicionadas sozinhas
% em um página em branco. Ver https://github.com/abntex/abntex2/issues/170
\makeatletter
\setlength{\@fptop}{5pt} % Set distance from top of page to first float
\makeatother
% ---

% ---
% Possibilita criação de Quadros e Lista de quadros.
% Ver https://github.com/abntex/abntex2/issues/176
%
\newcommand{\quadroname}{Quadro}
\newcommand{\listofquadrosname}{Lista de quadros}

\newfloat[chapter]{quadro}{loq}{\quadroname}
\newlistof{listofquadros}{loq}{\listofquadrosname}
\newlistentry{quadro}{loq}{0}

% Configurações para atender às regras da ABNT
\setfloatadjustment{quadro}{\centering}
\counterwithout{quadro}{chapter}
\renewcommand{\cftquadroname}{\quadroname\space} 
\renewcommand*{\cftquadroaftersnum}{\hfill--\hfill}

\setfloatlocations{quadro}{hbtp} % Ver https://github.com/abntex/abntex2/issues/176
% ---


% --- 
% Espaçamentos entre linhas e parágrafos 
% --- 

% O tamanho do parágrafo é dado por:
\setlength{\parindent}{1.3cm}

% Controle do espaçamento entre um parágrafo e outro:
\setlength{\parskip}{0.2cm}  % tente também \onelineskip

% ---
% Compila o índice
% ---
\makeindex
% ---


% ----
% Início do documento
% ----
\begin{document}

% Seleciona o idioma do documento (conforme pacotes do babel)
\selectlanguage{brazil}

% Retira espaço extra obsoleto entre as frases.
\frenchspacing 

% ----------------------------------------------------------
% ELEMENTOS PRÉ-TEXTUAIS
% ----------------------------------------------------------
\pretextual

% ---
% CAPA
% ---
\includepdf{./05-extras/capa.pdf}
% ---


% ---
% FOLHA DE ROSTO
% (o * indica que haverá a ficha bibliográfica)
% ---
\imprimirfolhaderosto
% ---


% ---
% Ficha bibliografica
% ---

% Isto é um exemplo de Ficha Catalográfica, ou ``Dados internacionais de
% catalogação-na-publicação''. Você pode utilizar este modelo como referência. 
% Porém, provavelmente a biblioteca da sua universidade lhe fornecerá um PDF
% com a ficha catalográfica definitiva após a defesa do trabalho. Quando estiver
% com o documento, salve-o como PDF no diretório do seu projeto e substitua todo
% o conteúdo de implementação deste arquivo pelo comando abaixo:
%
% \usepackage{pdfpages}		% necessário para comando \includepdf
% \begin{fichacatalografica}
%     \includepdf{fig_ficha_catalografica.pdf}
% \end{fichacatalografica}

% \begin{fichacatalografica}
% 	\sffamily
% 	\vspace*{\fill}                             % Posição vertical
% 	\begin{center}                              % Minipage Centralizado
% 	\fbox{\begin{minipage}[c][8cm]{13.5cm}		% Largura
% 	\small
% 	\imprimirautor
% 	%Sobrenome, Nome do autor
	
% 	\hspace{0.5cm} \imprimirtitulo  / \imprimirautor. --
% 	\imprimirlocal, \imprimirdata-
	
% 	\hspace{0.5cm} \thelastpage p. : il. (algumas color.) ; 30 cm.\\
	
% 	\hspace{0.5cm} \imprimirorientadorRotulo~\imprimirorientador\\
	
% 	\hspace{0.5cm}
% 	\parbox[t]{\textwidth}{\imprimirtipotrabalho~--~\imprimirinstituicao,
% 	\imprimirdata.}\\
	
% 	\hspace{0.5cm}
% 		1. Palavra-chave1.
% 		2. Palavra-chave2.
% 		2. Palavra-chave3.
% 		I. Orientador.
% 		II. Universidade xxx.
% 		III. Faculdade de xxx.
% 		IV. Título 			
% 	\end{minipage}}
% 	\end{center}
% \end{fichacatalografica}
% ---


% ---
% FOLHA DE APROVAÇÃO
% ---

% Isto é um exemplo de Folha de aprovação, elemento obrigatório da NBR
% 14724/2011 (seção 4.2.1.3). Você pode utilizar este modelo até a aprovação
% do trabalho. Após isso, substitua todo o conteúdo deste arquivo por uma
% imagem da página assinada pela banca com o comando abaixo:
%
% \begin{folhadeaprovacao}
% \includepdf{folhadeaprovacao_final.pdf}
% \end{folhadeaprovacao}
%
\setlength{\ABNTEXsignskip}{2cm}
\setlength{\ABNTEXsignwidth}{10cm}

\begin{folhadeaprovacao}
  \begin{center}
    {\ABNTEXchapterfont\large\imprimirautor}

    \vspace*{1cm}

    \begin{center}
      \ABNTEXchapterfont\bfseries\Large\imprimirtitulo
    \end{center}

    \vspace*{1cm}
    
    \hspace{.45\textwidth}
    \begin{minipage}{.5\textwidth}
        \imprimirpreambulo
    \end{minipage}%
   \end{center}

   \vspace*{1cm}
        
   \begin{center}
   Aprovado em 17 de julho de 2024.

   \vspace*{1cm}

   Banca examinadora:
   \end{center}

   \vspace*{-1cm}

   \assinatura{Prof. Ma. Caroline Santa Maria Rodrigues -- PUCRS}
   \assinatura{Prof. Dra. Clarissa Pinto Pizarro de Freitas -- PUCRS}
   \assinatura{\imprimirorientador -- PUCRS}
\end{folhadeaprovacao}
% ---


% ---
% Agradecimentos
% ---
\begin{agradecimentos}
Agradeço aos meus pais, Ademilson José Dias e Irinês Porfírio Dias, pelo apoio incondicional que me deram ao longo da vida,
incentivando mesmo as ideias questionáveis como aquela de fazer uma segunda graduação.

A minha irmã, Maria Eduarda Dias, cuja trajetória acadêmica é fonte de inspiração.

A minha namorada, Nathalia Fallavena Ceratti, que acompanhou a esta jornada de perto, oferecendo apoio e compreensão que
fizeram toda a diferença.

Aos meus amigos, Guilherme da Silva Petró, Norton Neves e Uilinton Luciano Marques da Silva, que foram compreensíveis nos
finais de semana em que não pude jogar RPG.

Ao professor Wagner de Lara Machado pela orientação e apoio no desenvolvimento deste trabalho.
\end{agradecimentos}
% ---


% ---
% RESUMOS
% ---
% Resumo em português
\setlength{\absparsep}{18pt} % ajusta o espaçamento dos parágrafos do resumo
\begin{resumo}
Desenvolvimentos na área de aprendizagem de máquina representam uma série de oportunidades para diversas áreas de pesquisa e de atuação profissional.
O uso de modelos de aprendizagem de máquina na predição de desfecho para tratamentos em saúde mental tem sido investigado e apresenta resultados promissores.
A aplicação desse tipo de modelo possibilitaria, por exemplo, maior assertividade na personalização de tratamentos psicoterápicos a nível individual.
Apesar dos benefícios potenciais, o tema ainda é pouco discutido nos espaços de pesquisa, formação e atuação em psicologia clínica. Este trabalho tem
por objetivo introduzir o tema de aprendizagem de máquina a psicólogos. Busca-se apresentar, de maneira acessível, os conceitos e terminologia básicos
em aprendizagem de máquina, proporcionando um panorama geral da área. Para complementar a exposição conceitual, é construído um exemplo a partir de dados
de um estudo sobre eficácia de intervenções em psicologia positiva para depressão via internet. Por fim, discutem-se algumas implicações do uso desse tipo
de tecnologia para a pesquisa e para a prática em psicoterapia.

 \textbf{Palavras-chave}: aprendizagem de máquina, psicoterapia, psicologia positiva, depressão.
\end{resumo}

% Resumo em inglês
\begin{resumo}[Abstract]
 \begin{otherlanguage*}{english}
   Advances in machine learning have opened up a multitude of opportunities for various research and professional fields. The use of machine learning models to
   predict treatment outcomes in mental health has been investigated showing promising results. The application of such models would allow greater assertiveness
   in tailoring psychotherapeutic treatments at the individual level. Despite its potential benefits, the topic remains unexplored in the fields of research, training
   and practice in clinical psychology. This work aims to introduce the topic of machine learning to psychologists. It presents the basic machine learning concepts and
   terminology in an accessible manner, providing a comprehensive overview of the field. To complement the conceptual discussion, an example is constructed using data
   from a study on the effectiveness of internet-based positive psychology interventions for depression. Finally, implications of using this type of technology in
   psychotherapy research and practice are discussed.
   \vspace{\onelineskip}
 
   \noindent 
   \textbf{Keywords}: machine learning, psychotherapy, positive psychology, depression.
 \end{otherlanguage*}
\end{resumo}


% ---
% LISTA DE ILUSTRAÇÕES
% ---
\pdfbookmark[0]{\listfigurename}{lof}
\listoffigures*
\cleardoublepage
% ---


% ---
% LISTA DE QUADROS
% ---
% \pdfbookmark[0]{\listofquadrosname}{loq}
% \listofquadros*
% \cleardoublepage
% ---


% ---
% LISTA DE TABELAS
% ---
% \pdfbookmark[0]{\listtablename}{lot}
% \listoftables*
% \cleardoublepage
% ---


% ---
% LISTA DE ABREVIATURAS E SIGLAS
% ---
% \begin{siglas}
%   \item[ABNT] Associação Brasileira de Normas Técnicas
%   \item[abnTeX] ABsurdas Normas para TeX
% \end{siglas}
% ---


% ---
% LISTA DE SÍMBOLOS
% ---
% \begin{simbolos}
%   \item[$ \Gamma $] Letra grega Gama
%   \item[$ \Lambda $] Lambda
%   \item[$ \zeta $] Letra grega minúscula zeta
%   \item[$ \in $] Pertence
% \end{simbolos}
% ---


% ---
% Sumário
% ---
\pdfbookmark[0]{\contentsname}{toc}
\tableofcontents*
\cleardoublepage
% ---


% ----------------------------------------------------------
% ELEMENTOS TEXTUAIS
% ----------------------------------------------------------
\textual

\documentclass[
	% Opções da classe memoir
	12pt,				        % tamanho da fonte
	openany,			        % capítulos começam em pág ímpar (insere página vazia caso preciso)
	oneside,			        % para impressão em recto e verso. Oposto a oneside
	a4paper,			        % tamanho do papel. 
	% Opções da classe abntex2
	%chapter=TITLE,		        % títulos de capítulos convertidos em letras maiúsculas
	%section=TITLE,		        % títulos de seções convertidos em letras maiúsculas
	%subsection=TITLE,	        % títulos de subseções convertidos em letras maiúsculas
	%subsubsection=TITLE,       % títulos de subsubseções convertidos em letras maiúsculas
	% Opções do pacote babel
	english,			        % idioma adicional para hifenização
	brazil				        % o último idioma é o principal do documento
	]{abntex2}

% ---
% Pacotes básicos 
% ---
\usepackage{lmodern}			% Usa a fonte Latin Modern			
\usepackage[T1]{fontenc}		% Selecao de codigos de fonte.
\usepackage[utf8]{inputenc}		% Codificacao do documento (conversão automática dos acentos)
\usepackage{indentfirst}		% Indenta o primeiro parágrafo de cada seção.
\usepackage{color}				% Controle das cores
\usepackage{graphicx}			% Inclusão de gráficos
\usepackage{microtype} 			% para melhorias de justificação
% ---
		

% ---
% Pacotes de citações
% ---
% \usepackage[alf]{abntex2cite}						% Citações padrão ABNT
\usepackage[
  alf,
  abnt-nbr10520=1988,                     % Citações em sentence case
  abnt-emphasize=bf
  ]{abntex2cite}


% ---
% Pacotes adicionais
% ---
\usepackage{pdfpages}	% Inclusão de páginas em PDF


% ---
% Informações para CAPA e FOLHA DE ROSTO
% ---
\titulo{APRENDIZAGEM DE MÁQUINA PARA PREDIÇÃO DE DESFECHOS EM SAÚDE MENTAL: \textmd{UMA INTRODUÇÃO PARA PSICÓLOGOS}}
\autor{GUILHERME AUGUSTO DIAS}
\local{Porto Alegre}
\data{2024}
\orientador{Prof. Dr. Wagner de Lara Machado}
% \coorientador{}
% \instituicao{
%   Pontifícia Universidade Católica do Rio Grande do Sul
%   \par
%   Escola de Ciências da Saúde e da Vida
%   \par
%   Curso de Psicologia}
\tipotrabalho{TCC}
\preambulo{Trabalho de Conclusão de Curso apresentado como requisito para a obtenção do grau de Bacharel em Psicologia pela Escola de Ciências da Saúde e da Vida da Pontifícia Universidade Católica do Rio Grande do Sul.}
% ---

% ---
% Configurações de aparência do PDF final
% ---

% Alterando o aspecto da cor azul
\definecolor{blue}{RGB}{41,5,195}

% Informações do PDF
\makeatletter
\hypersetup{
		pdftitle={Aprendizagem de máquina para predição de desfechos em saúde mental: uma introdução para psicólogos}, 
		pdfauthor={Guilherme Augusto Dias},
    	pdfsubject={\imprimirpreambulo},
	    pdfcreator={LaTeX with abnTeX2},
		pdfkeywords={aprendizagem de máquina}{psicoterapia}{psicologia positiva}{depressão}, 
		colorlinks=true,		% true: colored links; false: boxed links
    	linkcolor=black,		% color of internal links
    	citecolor=black,		% color of links to bibliography
    	filecolor=black,		% color of file links
		urlcolor=black,
		bookmarksdepth=4
}
\makeatother
% --- 


% ---
% Posiciona figuras e tabelas no topo da página quando adicionadas sozinhas
% em um página em branco. Ver https://github.com/abntex/abntex2/issues/170
\makeatletter
\setlength{\@fptop}{5pt} % Set distance from top of page to first float
\makeatother
% ---

% ---
% Possibilita criação de Quadros e Lista de quadros.
% Ver https://github.com/abntex/abntex2/issues/176
%
\newcommand{\quadroname}{Quadro}
\newcommand{\listofquadrosname}{Lista de quadros}

\newfloat[chapter]{quadro}{loq}{\quadroname}
\newlistof{listofquadros}{loq}{\listofquadrosname}
\newlistentry{quadro}{loq}{0}

% Configurações para atender às regras da ABNT
\setfloatadjustment{quadro}{\centering}
\counterwithout{quadro}{chapter}
\renewcommand{\cftquadroname}{\quadroname\space} 
\renewcommand*{\cftquadroaftersnum}{\hfill--\hfill}

\setfloatlocations{quadro}{hbtp} % Ver https://github.com/abntex/abntex2/issues/176
% ---


% --- 
% Espaçamentos entre linhas e parágrafos 
% --- 

% O tamanho do parágrafo é dado por:
\setlength{\parindent}{1.3cm}

% Controle do espaçamento entre um parágrafo e outro:
\setlength{\parskip}{0.2cm}  % tente também \onelineskip

% ---
% Compila o índice
% ---
\makeindex
% ---


% ----
% Início do documento
% ----
\begin{document}

% Seleciona o idioma do documento (conforme pacotes do babel)
\selectlanguage{brazil}

% Retira espaço extra obsoleto entre as frases.
\frenchspacing 

% ----------------------------------------------------------
% ELEMENTOS PRÉ-TEXTUAIS
% ----------------------------------------------------------
\pretextual

% ---
% CAPA
% ---
\includepdf{./05-extras/capa.pdf}
% ---


% ---
% FOLHA DE ROSTO
% (o * indica que haverá a ficha bibliográfica)
% ---
\imprimirfolhaderosto
% ---


% ---
% Ficha bibliografica
% ---

% Isto é um exemplo de Ficha Catalográfica, ou ``Dados internacionais de
% catalogação-na-publicação''. Você pode utilizar este modelo como referência. 
% Porém, provavelmente a biblioteca da sua universidade lhe fornecerá um PDF
% com a ficha catalográfica definitiva após a defesa do trabalho. Quando estiver
% com o documento, salve-o como PDF no diretório do seu projeto e substitua todo
% o conteúdo de implementação deste arquivo pelo comando abaixo:
%
% \usepackage{pdfpages}		% necessário para comando \includepdf
% \begin{fichacatalografica}
%     \includepdf{fig_ficha_catalografica.pdf}
% \end{fichacatalografica}

% \begin{fichacatalografica}
% 	\sffamily
% 	\vspace*{\fill}                             % Posição vertical
% 	\begin{center}                              % Minipage Centralizado
% 	\fbox{\begin{minipage}[c][8cm]{13.5cm}		% Largura
% 	\small
% 	\imprimirautor
% 	%Sobrenome, Nome do autor
	
% 	\hspace{0.5cm} \imprimirtitulo  / \imprimirautor. --
% 	\imprimirlocal, \imprimirdata-
	
% 	\hspace{0.5cm} \thelastpage p. : il. (algumas color.) ; 30 cm.\\
	
% 	\hspace{0.5cm} \imprimirorientadorRotulo~\imprimirorientador\\
	
% 	\hspace{0.5cm}
% 	\parbox[t]{\textwidth}{\imprimirtipotrabalho~--~\imprimirinstituicao,
% 	\imprimirdata.}\\
	
% 	\hspace{0.5cm}
% 		1. Palavra-chave1.
% 		2. Palavra-chave2.
% 		2. Palavra-chave3.
% 		I. Orientador.
% 		II. Universidade xxx.
% 		III. Faculdade de xxx.
% 		IV. Título 			
% 	\end{minipage}}
% 	\end{center}
% \end{fichacatalografica}
% ---


% ---
% FOLHA DE APROVAÇÃO
% ---

% Isto é um exemplo de Folha de aprovação, elemento obrigatório da NBR
% 14724/2011 (seção 4.2.1.3). Você pode utilizar este modelo até a aprovação
% do trabalho. Após isso, substitua todo o conteúdo deste arquivo por uma
% imagem da página assinada pela banca com o comando abaixo:
%
% \begin{folhadeaprovacao}
% \includepdf{folhadeaprovacao_final.pdf}
% \end{folhadeaprovacao}
%
\setlength{\ABNTEXsignskip}{2cm}
\setlength{\ABNTEXsignwidth}{10cm}

\begin{folhadeaprovacao}
  \begin{center}
    {\ABNTEXchapterfont\large\imprimirautor}

    \vspace*{1cm}

    \begin{center}
      \ABNTEXchapterfont\bfseries\Large\imprimirtitulo
    \end{center}

    \vspace*{1cm}
    
    \hspace{.45\textwidth}
    \begin{minipage}{.5\textwidth}
        \imprimirpreambulo
    \end{minipage}%
   \end{center}

   \vspace*{1cm}
        
   \begin{center}
   Aprovado em 17 de julho de 2024.

   \vspace*{1cm}

   Banca examinadora:
   \end{center}

   \vspace*{-1cm}

   \assinatura{Prof. Ma. Caroline Santa Maria Rodrigues -- PUCRS}
   \assinatura{Prof. Dra. Clarissa Pinto Pizarro de Freitas -- PUCRS}
   \assinatura{\imprimirorientador -- PUCRS}
\end{folhadeaprovacao}
% ---


% ---
% Agradecimentos
% ---
\begin{agradecimentos}
Agradeço aos meus pais, Ademilson José Dias e Irinês Porfírio Dias, pelo apoio incondicional que me deram ao longo da vida,
incentivando mesmo as ideias questionáveis como aquela de fazer uma segunda graduação.

A minha irmã, Maria Eduarda Dias, cuja trajetória acadêmica é fonte de inspiração.

A minha namorada, Nathalia Fallavena Ceratti, que acompanhou a esta jornada de perto, oferecendo apoio e compreensão que
fizeram toda a diferença.

Aos meus amigos, Guilherme da Silva Petró, Norton Neves e Uilinton Luciano Marques da Silva, que foram compreensíveis nos
finais de semana em que não pude jogar RPG.

Ao professor Wagner de Lara Machado pela orientação e apoio no desenvolvimento deste trabalho.
\end{agradecimentos}
% ---


% ---
% RESUMOS
% ---
% Resumo em português
\setlength{\absparsep}{18pt} % ajusta o espaçamento dos parágrafos do resumo
\begin{resumo}
Desenvolvimentos na área de aprendizagem de máquina representam uma série de oportunidades para diversas áreas de pesquisa e de atuação profissional.
O uso de modelos de aprendizagem de máquina na predição de desfecho para tratamentos em saúde mental tem sido investigado e apresenta resultados promissores.
A aplicação desse tipo de modelo possibilitaria, por exemplo, maior assertividade na personalização de tratamentos psicoterápicos a nível individual.
Apesar dos benefícios potenciais, o tema ainda é pouco discutido nos espaços de pesquisa, formação e atuação em psicologia clínica. Este trabalho tem
por objetivo introduzir o tema de aprendizagem de máquina a psicólogos. Busca-se apresentar, de maneira acessível, os conceitos e terminologia básicos
em aprendizagem de máquina, proporcionando um panorama geral da área. Para complementar a exposição conceitual, é construído um exemplo a partir de dados
de um estudo sobre eficácia de intervenções em psicologia positiva para depressão via internet. Por fim, discutem-se algumas implicações do uso desse tipo
de tecnologia para a pesquisa e para a prática em psicoterapia.

 \textbf{Palavras-chave}: aprendizagem de máquina, psicoterapia, psicologia positiva, depressão.
\end{resumo}

% Resumo em inglês
\begin{resumo}[Abstract]
 \begin{otherlanguage*}{english}
   Advances in machine learning have opened up a multitude of opportunities for various research and professional fields. The use of machine learning models to
   predict treatment outcomes in mental health has been investigated showing promising results. The application of such models would allow greater assertiveness
   in tailoring psychotherapeutic treatments at the individual level. Despite its potential benefits, the topic remains unexplored in the fields of research, training
   and practice in clinical psychology. This work aims to introduce the topic of machine learning to psychologists. It presents the basic machine learning concepts and
   terminology in an accessible manner, providing a comprehensive overview of the field. To complement the conceptual discussion, an example is constructed using data
   from a study on the effectiveness of internet-based positive psychology interventions for depression. Finally, implications of using this type of technology in
   psychotherapy research and practice are discussed.
   \vspace{\onelineskip}
 
   \noindent 
   \textbf{Keywords}: machine learning, psychotherapy, positive psychology, depression.
 \end{otherlanguage*}
\end{resumo}


% ---
% LISTA DE ILUSTRAÇÕES
% ---
\pdfbookmark[0]{\listfigurename}{lof}
\listoffigures*
\cleardoublepage
% ---


% ---
% LISTA DE QUADROS
% ---
% \pdfbookmark[0]{\listofquadrosname}{loq}
% \listofquadros*
% \cleardoublepage
% ---


% ---
% LISTA DE TABELAS
% ---
% \pdfbookmark[0]{\listtablename}{lot}
% \listoftables*
% \cleardoublepage
% ---


% ---
% LISTA DE ABREVIATURAS E SIGLAS
% ---
% \begin{siglas}
%   \item[ABNT] Associação Brasileira de Normas Técnicas
%   \item[abnTeX] ABsurdas Normas para TeX
% \end{siglas}
% ---


% ---
% LISTA DE SÍMBOLOS
% ---
% \begin{simbolos}
%   \item[$ \Gamma $] Letra grega Gama
%   \item[$ \Lambda $] Lambda
%   \item[$ \zeta $] Letra grega minúscula zeta
%   \item[$ \in $] Pertence
% \end{simbolos}
% ---


% ---
% Sumário
% ---
\pdfbookmark[0]{\contentsname}{toc}
\tableofcontents*
\cleardoublepage
% ---


% ----------------------------------------------------------
% ELEMENTOS TEXTUAIS
% ----------------------------------------------------------
\textual

\include{./01-introducao/main.tex}

\include{./02-machine-learning/main.tex}

\include{./03-exemplo/main.tex}

\include{./04-conclusao/main.tex}

% ----------------------------------------------------------
% ELEMENTOS PÓS-TEXTUAIS
% ----------------------------------------------------------
\postextual
% ----------------------------------------------------------


% ----------------------------------------------------------
% REFERÊNCIAS
% ----------------------------------------------------------
\bibliography{references.bib}


% ----------------------------------------------------------
% GLOSSÁRIO
% ----------------------------------------------------------
%
% Consulte o manual da classe abntex2 para orientações sobre o glossário.
%
%\glossary

% ----------------------------------------------------------
% APÊNDICES
% ----------------------------------------------------------

% ---
% Inicia os apêndices
% ---
% \begin{apendicesenv}

% % Imprime uma página indicando o início dos apêndices
% \partapendices

% % ----------------------------------------------------------
% \chapter{Quisque libero justo}
% % ----------------------------------------------------------

% % ----------------------------------------------------------
% \chapter{Nullam elementum urna}
% % ----------------------------------------------------------

% \end{apendicesenv}
% ---


% ----------------------------------------------------------
% ANEXOS
% ----------------------------------------------------------

% ---
% Inicia os anexos
% ---
% \begin{anexosenv}

% % Imprime uma página indicando o início dos anexos
% \partanexos

% % ---
% \chapter{Morbi ultrices rutrum lorem}
% % ---

% % ---
% \chapter{Cras non urna sed feugiat cum}
% % ---


% % ---
% \chapter{Fusce facilisis lacinia dui}
% % ---

% \end{anexosenv}

%---------------------------------------------------------------------
% INDICE REMISSIVO
%---------------------------------------------------------------------
% \phantompart
% \printindex
%---------------------------------------------------------------------
\includepdf{./05-extras/contracapa.pdf}
\end{document}


\documentclass[
	% Opções da classe memoir
	12pt,				        % tamanho da fonte
	openany,			        % capítulos começam em pág ímpar (insere página vazia caso preciso)
	oneside,			        % para impressão em recto e verso. Oposto a oneside
	a4paper,			        % tamanho do papel. 
	% Opções da classe abntex2
	%chapter=TITLE,		        % títulos de capítulos convertidos em letras maiúsculas
	%section=TITLE,		        % títulos de seções convertidos em letras maiúsculas
	%subsection=TITLE,	        % títulos de subseções convertidos em letras maiúsculas
	%subsubsection=TITLE,       % títulos de subsubseções convertidos em letras maiúsculas
	% Opções do pacote babel
	english,			        % idioma adicional para hifenização
	brazil				        % o último idioma é o principal do documento
	]{abntex2}

% ---
% Pacotes básicos 
% ---
\usepackage{lmodern}			% Usa a fonte Latin Modern			
\usepackage[T1]{fontenc}		% Selecao de codigos de fonte.
\usepackage[utf8]{inputenc}		% Codificacao do documento (conversão automática dos acentos)
\usepackage{indentfirst}		% Indenta o primeiro parágrafo de cada seção.
\usepackage{color}				% Controle das cores
\usepackage{graphicx}			% Inclusão de gráficos
\usepackage{microtype} 			% para melhorias de justificação
% ---
		

% ---
% Pacotes de citações
% ---
% \usepackage[alf]{abntex2cite}						% Citações padrão ABNT
\usepackage[
  alf,
  abnt-nbr10520=1988,                     % Citações em sentence case
  abnt-emphasize=bf
  ]{abntex2cite}


% ---
% Pacotes adicionais
% ---
\usepackage{pdfpages}	% Inclusão de páginas em PDF


% ---
% Informações para CAPA e FOLHA DE ROSTO
% ---
\titulo{APRENDIZAGEM DE MÁQUINA PARA PREDIÇÃO DE DESFECHOS EM SAÚDE MENTAL: \textmd{UMA INTRODUÇÃO PARA PSICÓLOGOS}}
\autor{GUILHERME AUGUSTO DIAS}
\local{Porto Alegre}
\data{2024}
\orientador{Prof. Dr. Wagner de Lara Machado}
% \coorientador{}
% \instituicao{
%   Pontifícia Universidade Católica do Rio Grande do Sul
%   \par
%   Escola de Ciências da Saúde e da Vida
%   \par
%   Curso de Psicologia}
\tipotrabalho{TCC}
\preambulo{Trabalho de Conclusão de Curso apresentado como requisito para a obtenção do grau de Bacharel em Psicologia pela Escola de Ciências da Saúde e da Vida da Pontifícia Universidade Católica do Rio Grande do Sul.}
% ---

% ---
% Configurações de aparência do PDF final
% ---

% Alterando o aspecto da cor azul
\definecolor{blue}{RGB}{41,5,195}

% Informações do PDF
\makeatletter
\hypersetup{
		pdftitle={Aprendizagem de máquina para predição de desfechos em saúde mental: uma introdução para psicólogos}, 
		pdfauthor={Guilherme Augusto Dias},
    	pdfsubject={\imprimirpreambulo},
	    pdfcreator={LaTeX with abnTeX2},
		pdfkeywords={aprendizagem de máquina}{psicoterapia}{psicologia positiva}{depressão}, 
		colorlinks=true,		% true: colored links; false: boxed links
    	linkcolor=black,		% color of internal links
    	citecolor=black,		% color of links to bibliography
    	filecolor=black,		% color of file links
		urlcolor=black,
		bookmarksdepth=4
}
\makeatother
% --- 


% ---
% Posiciona figuras e tabelas no topo da página quando adicionadas sozinhas
% em um página em branco. Ver https://github.com/abntex/abntex2/issues/170
\makeatletter
\setlength{\@fptop}{5pt} % Set distance from top of page to first float
\makeatother
% ---

% ---
% Possibilita criação de Quadros e Lista de quadros.
% Ver https://github.com/abntex/abntex2/issues/176
%
\newcommand{\quadroname}{Quadro}
\newcommand{\listofquadrosname}{Lista de quadros}

\newfloat[chapter]{quadro}{loq}{\quadroname}
\newlistof{listofquadros}{loq}{\listofquadrosname}
\newlistentry{quadro}{loq}{0}

% Configurações para atender às regras da ABNT
\setfloatadjustment{quadro}{\centering}
\counterwithout{quadro}{chapter}
\renewcommand{\cftquadroname}{\quadroname\space} 
\renewcommand*{\cftquadroaftersnum}{\hfill--\hfill}

\setfloatlocations{quadro}{hbtp} % Ver https://github.com/abntex/abntex2/issues/176
% ---


% --- 
% Espaçamentos entre linhas e parágrafos 
% --- 

% O tamanho do parágrafo é dado por:
\setlength{\parindent}{1.3cm}

% Controle do espaçamento entre um parágrafo e outro:
\setlength{\parskip}{0.2cm}  % tente também \onelineskip

% ---
% Compila o índice
% ---
\makeindex
% ---


% ----
% Início do documento
% ----
\begin{document}

% Seleciona o idioma do documento (conforme pacotes do babel)
\selectlanguage{brazil}

% Retira espaço extra obsoleto entre as frases.
\frenchspacing 

% ----------------------------------------------------------
% ELEMENTOS PRÉ-TEXTUAIS
% ----------------------------------------------------------
\pretextual

% ---
% CAPA
% ---
\includepdf{./05-extras/capa.pdf}
% ---


% ---
% FOLHA DE ROSTO
% (o * indica que haverá a ficha bibliográfica)
% ---
\imprimirfolhaderosto
% ---


% ---
% Ficha bibliografica
% ---

% Isto é um exemplo de Ficha Catalográfica, ou ``Dados internacionais de
% catalogação-na-publicação''. Você pode utilizar este modelo como referência. 
% Porém, provavelmente a biblioteca da sua universidade lhe fornecerá um PDF
% com a ficha catalográfica definitiva após a defesa do trabalho. Quando estiver
% com o documento, salve-o como PDF no diretório do seu projeto e substitua todo
% o conteúdo de implementação deste arquivo pelo comando abaixo:
%
% \usepackage{pdfpages}		% necessário para comando \includepdf
% \begin{fichacatalografica}
%     \includepdf{fig_ficha_catalografica.pdf}
% \end{fichacatalografica}

% \begin{fichacatalografica}
% 	\sffamily
% 	\vspace*{\fill}                             % Posição vertical
% 	\begin{center}                              % Minipage Centralizado
% 	\fbox{\begin{minipage}[c][8cm]{13.5cm}		% Largura
% 	\small
% 	\imprimirautor
% 	%Sobrenome, Nome do autor
	
% 	\hspace{0.5cm} \imprimirtitulo  / \imprimirautor. --
% 	\imprimirlocal, \imprimirdata-
	
% 	\hspace{0.5cm} \thelastpage p. : il. (algumas color.) ; 30 cm.\\
	
% 	\hspace{0.5cm} \imprimirorientadorRotulo~\imprimirorientador\\
	
% 	\hspace{0.5cm}
% 	\parbox[t]{\textwidth}{\imprimirtipotrabalho~--~\imprimirinstituicao,
% 	\imprimirdata.}\\
	
% 	\hspace{0.5cm}
% 		1. Palavra-chave1.
% 		2. Palavra-chave2.
% 		2. Palavra-chave3.
% 		I. Orientador.
% 		II. Universidade xxx.
% 		III. Faculdade de xxx.
% 		IV. Título 			
% 	\end{minipage}}
% 	\end{center}
% \end{fichacatalografica}
% ---


% ---
% FOLHA DE APROVAÇÃO
% ---

% Isto é um exemplo de Folha de aprovação, elemento obrigatório da NBR
% 14724/2011 (seção 4.2.1.3). Você pode utilizar este modelo até a aprovação
% do trabalho. Após isso, substitua todo o conteúdo deste arquivo por uma
% imagem da página assinada pela banca com o comando abaixo:
%
% \begin{folhadeaprovacao}
% \includepdf{folhadeaprovacao_final.pdf}
% \end{folhadeaprovacao}
%
\setlength{\ABNTEXsignskip}{2cm}
\setlength{\ABNTEXsignwidth}{10cm}

\begin{folhadeaprovacao}
  \begin{center}
    {\ABNTEXchapterfont\large\imprimirautor}

    \vspace*{1cm}

    \begin{center}
      \ABNTEXchapterfont\bfseries\Large\imprimirtitulo
    \end{center}

    \vspace*{1cm}
    
    \hspace{.45\textwidth}
    \begin{minipage}{.5\textwidth}
        \imprimirpreambulo
    \end{minipage}%
   \end{center}

   \vspace*{1cm}
        
   \begin{center}
   Aprovado em 17 de julho de 2024.

   \vspace*{1cm}

   Banca examinadora:
   \end{center}

   \vspace*{-1cm}

   \assinatura{Prof. Ma. Caroline Santa Maria Rodrigues -- PUCRS}
   \assinatura{Prof. Dra. Clarissa Pinto Pizarro de Freitas -- PUCRS}
   \assinatura{\imprimirorientador -- PUCRS}
\end{folhadeaprovacao}
% ---


% ---
% Agradecimentos
% ---
\begin{agradecimentos}
Agradeço aos meus pais, Ademilson José Dias e Irinês Porfírio Dias, pelo apoio incondicional que me deram ao longo da vida,
incentivando mesmo as ideias questionáveis como aquela de fazer uma segunda graduação.

A minha irmã, Maria Eduarda Dias, cuja trajetória acadêmica é fonte de inspiração.

A minha namorada, Nathalia Fallavena Ceratti, que acompanhou a esta jornada de perto, oferecendo apoio e compreensão que
fizeram toda a diferença.

Aos meus amigos, Guilherme da Silva Petró, Norton Neves e Uilinton Luciano Marques da Silva, que foram compreensíveis nos
finais de semana em que não pude jogar RPG.

Ao professor Wagner de Lara Machado pela orientação e apoio no desenvolvimento deste trabalho.
\end{agradecimentos}
% ---


% ---
% RESUMOS
% ---
% Resumo em português
\setlength{\absparsep}{18pt} % ajusta o espaçamento dos parágrafos do resumo
\begin{resumo}
Desenvolvimentos na área de aprendizagem de máquina representam uma série de oportunidades para diversas áreas de pesquisa e de atuação profissional.
O uso de modelos de aprendizagem de máquina na predição de desfecho para tratamentos em saúde mental tem sido investigado e apresenta resultados promissores.
A aplicação desse tipo de modelo possibilitaria, por exemplo, maior assertividade na personalização de tratamentos psicoterápicos a nível individual.
Apesar dos benefícios potenciais, o tema ainda é pouco discutido nos espaços de pesquisa, formação e atuação em psicologia clínica. Este trabalho tem
por objetivo introduzir o tema de aprendizagem de máquina a psicólogos. Busca-se apresentar, de maneira acessível, os conceitos e terminologia básicos
em aprendizagem de máquina, proporcionando um panorama geral da área. Para complementar a exposição conceitual, é construído um exemplo a partir de dados
de um estudo sobre eficácia de intervenções em psicologia positiva para depressão via internet. Por fim, discutem-se algumas implicações do uso desse tipo
de tecnologia para a pesquisa e para a prática em psicoterapia.

 \textbf{Palavras-chave}: aprendizagem de máquina, psicoterapia, psicologia positiva, depressão.
\end{resumo}

% Resumo em inglês
\begin{resumo}[Abstract]
 \begin{otherlanguage*}{english}
   Advances in machine learning have opened up a multitude of opportunities for various research and professional fields. The use of machine learning models to
   predict treatment outcomes in mental health has been investigated showing promising results. The application of such models would allow greater assertiveness
   in tailoring psychotherapeutic treatments at the individual level. Despite its potential benefits, the topic remains unexplored in the fields of research, training
   and practice in clinical psychology. This work aims to introduce the topic of machine learning to psychologists. It presents the basic machine learning concepts and
   terminology in an accessible manner, providing a comprehensive overview of the field. To complement the conceptual discussion, an example is constructed using data
   from a study on the effectiveness of internet-based positive psychology interventions for depression. Finally, implications of using this type of technology in
   psychotherapy research and practice are discussed.
   \vspace{\onelineskip}
 
   \noindent 
   \textbf{Keywords}: machine learning, psychotherapy, positive psychology, depression.
 \end{otherlanguage*}
\end{resumo}


% ---
% LISTA DE ILUSTRAÇÕES
% ---
\pdfbookmark[0]{\listfigurename}{lof}
\listoffigures*
\cleardoublepage
% ---


% ---
% LISTA DE QUADROS
% ---
% \pdfbookmark[0]{\listofquadrosname}{loq}
% \listofquadros*
% \cleardoublepage
% ---


% ---
% LISTA DE TABELAS
% ---
% \pdfbookmark[0]{\listtablename}{lot}
% \listoftables*
% \cleardoublepage
% ---


% ---
% LISTA DE ABREVIATURAS E SIGLAS
% ---
% \begin{siglas}
%   \item[ABNT] Associação Brasileira de Normas Técnicas
%   \item[abnTeX] ABsurdas Normas para TeX
% \end{siglas}
% ---


% ---
% LISTA DE SÍMBOLOS
% ---
% \begin{simbolos}
%   \item[$ \Gamma $] Letra grega Gama
%   \item[$ \Lambda $] Lambda
%   \item[$ \zeta $] Letra grega minúscula zeta
%   \item[$ \in $] Pertence
% \end{simbolos}
% ---


% ---
% Sumário
% ---
\pdfbookmark[0]{\contentsname}{toc}
\tableofcontents*
\cleardoublepage
% ---


% ----------------------------------------------------------
% ELEMENTOS TEXTUAIS
% ----------------------------------------------------------
\textual

\include{./01-introducao/main.tex}

\include{./02-machine-learning/main.tex}

\include{./03-exemplo/main.tex}

\include{./04-conclusao/main.tex}

% ----------------------------------------------------------
% ELEMENTOS PÓS-TEXTUAIS
% ----------------------------------------------------------
\postextual
% ----------------------------------------------------------


% ----------------------------------------------------------
% REFERÊNCIAS
% ----------------------------------------------------------
\bibliography{references.bib}


% ----------------------------------------------------------
% GLOSSÁRIO
% ----------------------------------------------------------
%
% Consulte o manual da classe abntex2 para orientações sobre o glossário.
%
%\glossary

% ----------------------------------------------------------
% APÊNDICES
% ----------------------------------------------------------

% ---
% Inicia os apêndices
% ---
% \begin{apendicesenv}

% % Imprime uma página indicando o início dos apêndices
% \partapendices

% % ----------------------------------------------------------
% \chapter{Quisque libero justo}
% % ----------------------------------------------------------

% % ----------------------------------------------------------
% \chapter{Nullam elementum urna}
% % ----------------------------------------------------------

% \end{apendicesenv}
% ---


% ----------------------------------------------------------
% ANEXOS
% ----------------------------------------------------------

% ---
% Inicia os anexos
% ---
% \begin{anexosenv}

% % Imprime uma página indicando o início dos anexos
% \partanexos

% % ---
% \chapter{Morbi ultrices rutrum lorem}
% % ---

% % ---
% \chapter{Cras non urna sed feugiat cum}
% % ---


% % ---
% \chapter{Fusce facilisis lacinia dui}
% % ---

% \end{anexosenv}

%---------------------------------------------------------------------
% INDICE REMISSIVO
%---------------------------------------------------------------------
% \phantompart
% \printindex
%---------------------------------------------------------------------
\includepdf{./05-extras/contracapa.pdf}
\end{document}


\documentclass[
	% Opções da classe memoir
	12pt,				        % tamanho da fonte
	openany,			        % capítulos começam em pág ímpar (insere página vazia caso preciso)
	oneside,			        % para impressão em recto e verso. Oposto a oneside
	a4paper,			        % tamanho do papel. 
	% Opções da classe abntex2
	%chapter=TITLE,		        % títulos de capítulos convertidos em letras maiúsculas
	%section=TITLE,		        % títulos de seções convertidos em letras maiúsculas
	%subsection=TITLE,	        % títulos de subseções convertidos em letras maiúsculas
	%subsubsection=TITLE,       % títulos de subsubseções convertidos em letras maiúsculas
	% Opções do pacote babel
	english,			        % idioma adicional para hifenização
	brazil				        % o último idioma é o principal do documento
	]{abntex2}

% ---
% Pacotes básicos 
% ---
\usepackage{lmodern}			% Usa a fonte Latin Modern			
\usepackage[T1]{fontenc}		% Selecao de codigos de fonte.
\usepackage[utf8]{inputenc}		% Codificacao do documento (conversão automática dos acentos)
\usepackage{indentfirst}		% Indenta o primeiro parágrafo de cada seção.
\usepackage{color}				% Controle das cores
\usepackage{graphicx}			% Inclusão de gráficos
\usepackage{microtype} 			% para melhorias de justificação
% ---
		

% ---
% Pacotes de citações
% ---
% \usepackage[alf]{abntex2cite}						% Citações padrão ABNT
\usepackage[
  alf,
  abnt-nbr10520=1988,                     % Citações em sentence case
  abnt-emphasize=bf
  ]{abntex2cite}


% ---
% Pacotes adicionais
% ---
\usepackage{pdfpages}	% Inclusão de páginas em PDF


% ---
% Informações para CAPA e FOLHA DE ROSTO
% ---
\titulo{APRENDIZAGEM DE MÁQUINA PARA PREDIÇÃO DE DESFECHOS EM SAÚDE MENTAL: \textmd{UMA INTRODUÇÃO PARA PSICÓLOGOS}}
\autor{GUILHERME AUGUSTO DIAS}
\local{Porto Alegre}
\data{2024}
\orientador{Prof. Dr. Wagner de Lara Machado}
% \coorientador{}
% \instituicao{
%   Pontifícia Universidade Católica do Rio Grande do Sul
%   \par
%   Escola de Ciências da Saúde e da Vida
%   \par
%   Curso de Psicologia}
\tipotrabalho{TCC}
\preambulo{Trabalho de Conclusão de Curso apresentado como requisito para a obtenção do grau de Bacharel em Psicologia pela Escola de Ciências da Saúde e da Vida da Pontifícia Universidade Católica do Rio Grande do Sul.}
% ---

% ---
% Configurações de aparência do PDF final
% ---

% Alterando o aspecto da cor azul
\definecolor{blue}{RGB}{41,5,195}

% Informações do PDF
\makeatletter
\hypersetup{
		pdftitle={Aprendizagem de máquina para predição de desfechos em saúde mental: uma introdução para psicólogos}, 
		pdfauthor={Guilherme Augusto Dias},
    	pdfsubject={\imprimirpreambulo},
	    pdfcreator={LaTeX with abnTeX2},
		pdfkeywords={aprendizagem de máquina}{psicoterapia}{psicologia positiva}{depressão}, 
		colorlinks=true,		% true: colored links; false: boxed links
    	linkcolor=black,		% color of internal links
    	citecolor=black,		% color of links to bibliography
    	filecolor=black,		% color of file links
		urlcolor=black,
		bookmarksdepth=4
}
\makeatother
% --- 


% ---
% Posiciona figuras e tabelas no topo da página quando adicionadas sozinhas
% em um página em branco. Ver https://github.com/abntex/abntex2/issues/170
\makeatletter
\setlength{\@fptop}{5pt} % Set distance from top of page to first float
\makeatother
% ---

% ---
% Possibilita criação de Quadros e Lista de quadros.
% Ver https://github.com/abntex/abntex2/issues/176
%
\newcommand{\quadroname}{Quadro}
\newcommand{\listofquadrosname}{Lista de quadros}

\newfloat[chapter]{quadro}{loq}{\quadroname}
\newlistof{listofquadros}{loq}{\listofquadrosname}
\newlistentry{quadro}{loq}{0}

% Configurações para atender às regras da ABNT
\setfloatadjustment{quadro}{\centering}
\counterwithout{quadro}{chapter}
\renewcommand{\cftquadroname}{\quadroname\space} 
\renewcommand*{\cftquadroaftersnum}{\hfill--\hfill}

\setfloatlocations{quadro}{hbtp} % Ver https://github.com/abntex/abntex2/issues/176
% ---


% --- 
% Espaçamentos entre linhas e parágrafos 
% --- 

% O tamanho do parágrafo é dado por:
\setlength{\parindent}{1.3cm}

% Controle do espaçamento entre um parágrafo e outro:
\setlength{\parskip}{0.2cm}  % tente também \onelineskip

% ---
% Compila o índice
% ---
\makeindex
% ---


% ----
% Início do documento
% ----
\begin{document}

% Seleciona o idioma do documento (conforme pacotes do babel)
\selectlanguage{brazil}

% Retira espaço extra obsoleto entre as frases.
\frenchspacing 

% ----------------------------------------------------------
% ELEMENTOS PRÉ-TEXTUAIS
% ----------------------------------------------------------
\pretextual

% ---
% CAPA
% ---
\includepdf{./05-extras/capa.pdf}
% ---


% ---
% FOLHA DE ROSTO
% (o * indica que haverá a ficha bibliográfica)
% ---
\imprimirfolhaderosto
% ---


% ---
% Ficha bibliografica
% ---

% Isto é um exemplo de Ficha Catalográfica, ou ``Dados internacionais de
% catalogação-na-publicação''. Você pode utilizar este modelo como referência. 
% Porém, provavelmente a biblioteca da sua universidade lhe fornecerá um PDF
% com a ficha catalográfica definitiva após a defesa do trabalho. Quando estiver
% com o documento, salve-o como PDF no diretório do seu projeto e substitua todo
% o conteúdo de implementação deste arquivo pelo comando abaixo:
%
% \usepackage{pdfpages}		% necessário para comando \includepdf
% \begin{fichacatalografica}
%     \includepdf{fig_ficha_catalografica.pdf}
% \end{fichacatalografica}

% \begin{fichacatalografica}
% 	\sffamily
% 	\vspace*{\fill}                             % Posição vertical
% 	\begin{center}                              % Minipage Centralizado
% 	\fbox{\begin{minipage}[c][8cm]{13.5cm}		% Largura
% 	\small
% 	\imprimirautor
% 	%Sobrenome, Nome do autor
	
% 	\hspace{0.5cm} \imprimirtitulo  / \imprimirautor. --
% 	\imprimirlocal, \imprimirdata-
	
% 	\hspace{0.5cm} \thelastpage p. : il. (algumas color.) ; 30 cm.\\
	
% 	\hspace{0.5cm} \imprimirorientadorRotulo~\imprimirorientador\\
	
% 	\hspace{0.5cm}
% 	\parbox[t]{\textwidth}{\imprimirtipotrabalho~--~\imprimirinstituicao,
% 	\imprimirdata.}\\
	
% 	\hspace{0.5cm}
% 		1. Palavra-chave1.
% 		2. Palavra-chave2.
% 		2. Palavra-chave3.
% 		I. Orientador.
% 		II. Universidade xxx.
% 		III. Faculdade de xxx.
% 		IV. Título 			
% 	\end{minipage}}
% 	\end{center}
% \end{fichacatalografica}
% ---


% ---
% FOLHA DE APROVAÇÃO
% ---

% Isto é um exemplo de Folha de aprovação, elemento obrigatório da NBR
% 14724/2011 (seção 4.2.1.3). Você pode utilizar este modelo até a aprovação
% do trabalho. Após isso, substitua todo o conteúdo deste arquivo por uma
% imagem da página assinada pela banca com o comando abaixo:
%
% \begin{folhadeaprovacao}
% \includepdf{folhadeaprovacao_final.pdf}
% \end{folhadeaprovacao}
%
\setlength{\ABNTEXsignskip}{2cm}
\setlength{\ABNTEXsignwidth}{10cm}

\begin{folhadeaprovacao}
  \begin{center}
    {\ABNTEXchapterfont\large\imprimirautor}

    \vspace*{1cm}

    \begin{center}
      \ABNTEXchapterfont\bfseries\Large\imprimirtitulo
    \end{center}

    \vspace*{1cm}
    
    \hspace{.45\textwidth}
    \begin{minipage}{.5\textwidth}
        \imprimirpreambulo
    \end{minipage}%
   \end{center}

   \vspace*{1cm}
        
   \begin{center}
   Aprovado em 17 de julho de 2024.

   \vspace*{1cm}

   Banca examinadora:
   \end{center}

   \vspace*{-1cm}

   \assinatura{Prof. Ma. Caroline Santa Maria Rodrigues -- PUCRS}
   \assinatura{Prof. Dra. Clarissa Pinto Pizarro de Freitas -- PUCRS}
   \assinatura{\imprimirorientador -- PUCRS}
\end{folhadeaprovacao}
% ---


% ---
% Agradecimentos
% ---
\begin{agradecimentos}
Agradeço aos meus pais, Ademilson José Dias e Irinês Porfírio Dias, pelo apoio incondicional que me deram ao longo da vida,
incentivando mesmo as ideias questionáveis como aquela de fazer uma segunda graduação.

A minha irmã, Maria Eduarda Dias, cuja trajetória acadêmica é fonte de inspiração.

A minha namorada, Nathalia Fallavena Ceratti, que acompanhou a esta jornada de perto, oferecendo apoio e compreensão que
fizeram toda a diferença.

Aos meus amigos, Guilherme da Silva Petró, Norton Neves e Uilinton Luciano Marques da Silva, que foram compreensíveis nos
finais de semana em que não pude jogar RPG.

Ao professor Wagner de Lara Machado pela orientação e apoio no desenvolvimento deste trabalho.
\end{agradecimentos}
% ---


% ---
% RESUMOS
% ---
% Resumo em português
\setlength{\absparsep}{18pt} % ajusta o espaçamento dos parágrafos do resumo
\begin{resumo}
Desenvolvimentos na área de aprendizagem de máquina representam uma série de oportunidades para diversas áreas de pesquisa e de atuação profissional.
O uso de modelos de aprendizagem de máquina na predição de desfecho para tratamentos em saúde mental tem sido investigado e apresenta resultados promissores.
A aplicação desse tipo de modelo possibilitaria, por exemplo, maior assertividade na personalização de tratamentos psicoterápicos a nível individual.
Apesar dos benefícios potenciais, o tema ainda é pouco discutido nos espaços de pesquisa, formação e atuação em psicologia clínica. Este trabalho tem
por objetivo introduzir o tema de aprendizagem de máquina a psicólogos. Busca-se apresentar, de maneira acessível, os conceitos e terminologia básicos
em aprendizagem de máquina, proporcionando um panorama geral da área. Para complementar a exposição conceitual, é construído um exemplo a partir de dados
de um estudo sobre eficácia de intervenções em psicologia positiva para depressão via internet. Por fim, discutem-se algumas implicações do uso desse tipo
de tecnologia para a pesquisa e para a prática em psicoterapia.

 \textbf{Palavras-chave}: aprendizagem de máquina, psicoterapia, psicologia positiva, depressão.
\end{resumo}

% Resumo em inglês
\begin{resumo}[Abstract]
 \begin{otherlanguage*}{english}
   Advances in machine learning have opened up a multitude of opportunities for various research and professional fields. The use of machine learning models to
   predict treatment outcomes in mental health has been investigated showing promising results. The application of such models would allow greater assertiveness
   in tailoring psychotherapeutic treatments at the individual level. Despite its potential benefits, the topic remains unexplored in the fields of research, training
   and practice in clinical psychology. This work aims to introduce the topic of machine learning to psychologists. It presents the basic machine learning concepts and
   terminology in an accessible manner, providing a comprehensive overview of the field. To complement the conceptual discussion, an example is constructed using data
   from a study on the effectiveness of internet-based positive psychology interventions for depression. Finally, implications of using this type of technology in
   psychotherapy research and practice are discussed.
   \vspace{\onelineskip}
 
   \noindent 
   \textbf{Keywords}: machine learning, psychotherapy, positive psychology, depression.
 \end{otherlanguage*}
\end{resumo}


% ---
% LISTA DE ILUSTRAÇÕES
% ---
\pdfbookmark[0]{\listfigurename}{lof}
\listoffigures*
\cleardoublepage
% ---


% ---
% LISTA DE QUADROS
% ---
% \pdfbookmark[0]{\listofquadrosname}{loq}
% \listofquadros*
% \cleardoublepage
% ---


% ---
% LISTA DE TABELAS
% ---
% \pdfbookmark[0]{\listtablename}{lot}
% \listoftables*
% \cleardoublepage
% ---


% ---
% LISTA DE ABREVIATURAS E SIGLAS
% ---
% \begin{siglas}
%   \item[ABNT] Associação Brasileira de Normas Técnicas
%   \item[abnTeX] ABsurdas Normas para TeX
% \end{siglas}
% ---


% ---
% LISTA DE SÍMBOLOS
% ---
% \begin{simbolos}
%   \item[$ \Gamma $] Letra grega Gama
%   \item[$ \Lambda $] Lambda
%   \item[$ \zeta $] Letra grega minúscula zeta
%   \item[$ \in $] Pertence
% \end{simbolos}
% ---


% ---
% Sumário
% ---
\pdfbookmark[0]{\contentsname}{toc}
\tableofcontents*
\cleardoublepage
% ---


% ----------------------------------------------------------
% ELEMENTOS TEXTUAIS
% ----------------------------------------------------------
\textual

\include{./01-introducao/main.tex}

\include{./02-machine-learning/main.tex}

\include{./03-exemplo/main.tex}

\include{./04-conclusao/main.tex}

% ----------------------------------------------------------
% ELEMENTOS PÓS-TEXTUAIS
% ----------------------------------------------------------
\postextual
% ----------------------------------------------------------


% ----------------------------------------------------------
% REFERÊNCIAS
% ----------------------------------------------------------
\bibliography{references.bib}


% ----------------------------------------------------------
% GLOSSÁRIO
% ----------------------------------------------------------
%
% Consulte o manual da classe abntex2 para orientações sobre o glossário.
%
%\glossary

% ----------------------------------------------------------
% APÊNDICES
% ----------------------------------------------------------

% ---
% Inicia os apêndices
% ---
% \begin{apendicesenv}

% % Imprime uma página indicando o início dos apêndices
% \partapendices

% % ----------------------------------------------------------
% \chapter{Quisque libero justo}
% % ----------------------------------------------------------

% % ----------------------------------------------------------
% \chapter{Nullam elementum urna}
% % ----------------------------------------------------------

% \end{apendicesenv}
% ---


% ----------------------------------------------------------
% ANEXOS
% ----------------------------------------------------------

% ---
% Inicia os anexos
% ---
% \begin{anexosenv}

% % Imprime uma página indicando o início dos anexos
% \partanexos

% % ---
% \chapter{Morbi ultrices rutrum lorem}
% % ---

% % ---
% \chapter{Cras non urna sed feugiat cum}
% % ---


% % ---
% \chapter{Fusce facilisis lacinia dui}
% % ---

% \end{anexosenv}

%---------------------------------------------------------------------
% INDICE REMISSIVO
%---------------------------------------------------------------------
% \phantompart
% \printindex
%---------------------------------------------------------------------
\includepdf{./05-extras/contracapa.pdf}
\end{document}


\documentclass[
	% Opções da classe memoir
	12pt,				        % tamanho da fonte
	openany,			        % capítulos começam em pág ímpar (insere página vazia caso preciso)
	oneside,			        % para impressão em recto e verso. Oposto a oneside
	a4paper,			        % tamanho do papel. 
	% Opções da classe abntex2
	%chapter=TITLE,		        % títulos de capítulos convertidos em letras maiúsculas
	%section=TITLE,		        % títulos de seções convertidos em letras maiúsculas
	%subsection=TITLE,	        % títulos de subseções convertidos em letras maiúsculas
	%subsubsection=TITLE,       % títulos de subsubseções convertidos em letras maiúsculas
	% Opções do pacote babel
	english,			        % idioma adicional para hifenização
	brazil				        % o último idioma é o principal do documento
	]{abntex2}

% ---
% Pacotes básicos 
% ---
\usepackage{lmodern}			% Usa a fonte Latin Modern			
\usepackage[T1]{fontenc}		% Selecao de codigos de fonte.
\usepackage[utf8]{inputenc}		% Codificacao do documento (conversão automática dos acentos)
\usepackage{indentfirst}		% Indenta o primeiro parágrafo de cada seção.
\usepackage{color}				% Controle das cores
\usepackage{graphicx}			% Inclusão de gráficos
\usepackage{microtype} 			% para melhorias de justificação
% ---
		

% ---
% Pacotes de citações
% ---
% \usepackage[alf]{abntex2cite}						% Citações padrão ABNT
\usepackage[
  alf,
  abnt-nbr10520=1988,                     % Citações em sentence case
  abnt-emphasize=bf
  ]{abntex2cite}


% ---
% Pacotes adicionais
% ---
\usepackage{pdfpages}	% Inclusão de páginas em PDF


% ---
% Informações para CAPA e FOLHA DE ROSTO
% ---
\titulo{APRENDIZAGEM DE MÁQUINA PARA PREDIÇÃO DE DESFECHOS EM SAÚDE MENTAL: \textmd{UMA INTRODUÇÃO PARA PSICÓLOGOS}}
\autor{GUILHERME AUGUSTO DIAS}
\local{Porto Alegre}
\data{2024}
\orientador{Prof. Dr. Wagner de Lara Machado}
% \coorientador{}
% \instituicao{
%   Pontifícia Universidade Católica do Rio Grande do Sul
%   \par
%   Escola de Ciências da Saúde e da Vida
%   \par
%   Curso de Psicologia}
\tipotrabalho{TCC}
\preambulo{Trabalho de Conclusão de Curso apresentado como requisito para a obtenção do grau de Bacharel em Psicologia pela Escola de Ciências da Saúde e da Vida da Pontifícia Universidade Católica do Rio Grande do Sul.}
% ---

% ---
% Configurações de aparência do PDF final
% ---

% Alterando o aspecto da cor azul
\definecolor{blue}{RGB}{41,5,195}

% Informações do PDF
\makeatletter
\hypersetup{
		pdftitle={Aprendizagem de máquina para predição de desfechos em saúde mental: uma introdução para psicólogos}, 
		pdfauthor={Guilherme Augusto Dias},
    	pdfsubject={\imprimirpreambulo},
	    pdfcreator={LaTeX with abnTeX2},
		pdfkeywords={aprendizagem de máquina}{psicoterapia}{psicologia positiva}{depressão}, 
		colorlinks=true,		% true: colored links; false: boxed links
    	linkcolor=black,		% color of internal links
    	citecolor=black,		% color of links to bibliography
    	filecolor=black,		% color of file links
		urlcolor=black,
		bookmarksdepth=4
}
\makeatother
% --- 


% ---
% Posiciona figuras e tabelas no topo da página quando adicionadas sozinhas
% em um página em branco. Ver https://github.com/abntex/abntex2/issues/170
\makeatletter
\setlength{\@fptop}{5pt} % Set distance from top of page to first float
\makeatother
% ---

% ---
% Possibilita criação de Quadros e Lista de quadros.
% Ver https://github.com/abntex/abntex2/issues/176
%
\newcommand{\quadroname}{Quadro}
\newcommand{\listofquadrosname}{Lista de quadros}

\newfloat[chapter]{quadro}{loq}{\quadroname}
\newlistof{listofquadros}{loq}{\listofquadrosname}
\newlistentry{quadro}{loq}{0}

% Configurações para atender às regras da ABNT
\setfloatadjustment{quadro}{\centering}
\counterwithout{quadro}{chapter}
\renewcommand{\cftquadroname}{\quadroname\space} 
\renewcommand*{\cftquadroaftersnum}{\hfill--\hfill}

\setfloatlocations{quadro}{hbtp} % Ver https://github.com/abntex/abntex2/issues/176
% ---


% --- 
% Espaçamentos entre linhas e parágrafos 
% --- 

% O tamanho do parágrafo é dado por:
\setlength{\parindent}{1.3cm}

% Controle do espaçamento entre um parágrafo e outro:
\setlength{\parskip}{0.2cm}  % tente também \onelineskip

% ---
% Compila o índice
% ---
\makeindex
% ---


% ----
% Início do documento
% ----
\begin{document}

% Seleciona o idioma do documento (conforme pacotes do babel)
\selectlanguage{brazil}

% Retira espaço extra obsoleto entre as frases.
\frenchspacing 

% ----------------------------------------------------------
% ELEMENTOS PRÉ-TEXTUAIS
% ----------------------------------------------------------
\pretextual

% ---
% CAPA
% ---
\includepdf{./05-extras/capa.pdf}
% ---


% ---
% FOLHA DE ROSTO
% (o * indica que haverá a ficha bibliográfica)
% ---
\imprimirfolhaderosto
% ---


% ---
% Ficha bibliografica
% ---

% Isto é um exemplo de Ficha Catalográfica, ou ``Dados internacionais de
% catalogação-na-publicação''. Você pode utilizar este modelo como referência. 
% Porém, provavelmente a biblioteca da sua universidade lhe fornecerá um PDF
% com a ficha catalográfica definitiva após a defesa do trabalho. Quando estiver
% com o documento, salve-o como PDF no diretório do seu projeto e substitua todo
% o conteúdo de implementação deste arquivo pelo comando abaixo:
%
% \usepackage{pdfpages}		% necessário para comando \includepdf
% \begin{fichacatalografica}
%     \includepdf{fig_ficha_catalografica.pdf}
% \end{fichacatalografica}

% \begin{fichacatalografica}
% 	\sffamily
% 	\vspace*{\fill}                             % Posição vertical
% 	\begin{center}                              % Minipage Centralizado
% 	\fbox{\begin{minipage}[c][8cm]{13.5cm}		% Largura
% 	\small
% 	\imprimirautor
% 	%Sobrenome, Nome do autor
	
% 	\hspace{0.5cm} \imprimirtitulo  / \imprimirautor. --
% 	\imprimirlocal, \imprimirdata-
	
% 	\hspace{0.5cm} \thelastpage p. : il. (algumas color.) ; 30 cm.\\
	
% 	\hspace{0.5cm} \imprimirorientadorRotulo~\imprimirorientador\\
	
% 	\hspace{0.5cm}
% 	\parbox[t]{\textwidth}{\imprimirtipotrabalho~--~\imprimirinstituicao,
% 	\imprimirdata.}\\
	
% 	\hspace{0.5cm}
% 		1. Palavra-chave1.
% 		2. Palavra-chave2.
% 		2. Palavra-chave3.
% 		I. Orientador.
% 		II. Universidade xxx.
% 		III. Faculdade de xxx.
% 		IV. Título 			
% 	\end{minipage}}
% 	\end{center}
% \end{fichacatalografica}
% ---


% ---
% FOLHA DE APROVAÇÃO
% ---

% Isto é um exemplo de Folha de aprovação, elemento obrigatório da NBR
% 14724/2011 (seção 4.2.1.3). Você pode utilizar este modelo até a aprovação
% do trabalho. Após isso, substitua todo o conteúdo deste arquivo por uma
% imagem da página assinada pela banca com o comando abaixo:
%
% \begin{folhadeaprovacao}
% \includepdf{folhadeaprovacao_final.pdf}
% \end{folhadeaprovacao}
%
\setlength{\ABNTEXsignskip}{2cm}
\setlength{\ABNTEXsignwidth}{10cm}

\begin{folhadeaprovacao}
  \begin{center}
    {\ABNTEXchapterfont\large\imprimirautor}

    \vspace*{1cm}

    \begin{center}
      \ABNTEXchapterfont\bfseries\Large\imprimirtitulo
    \end{center}

    \vspace*{1cm}
    
    \hspace{.45\textwidth}
    \begin{minipage}{.5\textwidth}
        \imprimirpreambulo
    \end{minipage}%
   \end{center}

   \vspace*{1cm}
        
   \begin{center}
   Aprovado em 17 de julho de 2024.

   \vspace*{1cm}

   Banca examinadora:
   \end{center}

   \vspace*{-1cm}

   \assinatura{Prof. Ma. Caroline Santa Maria Rodrigues -- PUCRS}
   \assinatura{Prof. Dra. Clarissa Pinto Pizarro de Freitas -- PUCRS}
   \assinatura{\imprimirorientador -- PUCRS}
\end{folhadeaprovacao}
% ---


% ---
% Agradecimentos
% ---
\begin{agradecimentos}
Agradeço aos meus pais, Ademilson José Dias e Irinês Porfírio Dias, pelo apoio incondicional que me deram ao longo da vida,
incentivando mesmo as ideias questionáveis como aquela de fazer uma segunda graduação.

A minha irmã, Maria Eduarda Dias, cuja trajetória acadêmica é fonte de inspiração.

A minha namorada, Nathalia Fallavena Ceratti, que acompanhou a esta jornada de perto, oferecendo apoio e compreensão que
fizeram toda a diferença.

Aos meus amigos, Guilherme da Silva Petró, Norton Neves e Uilinton Luciano Marques da Silva, que foram compreensíveis nos
finais de semana em que não pude jogar RPG.

Ao professor Wagner de Lara Machado pela orientação e apoio no desenvolvimento deste trabalho.
\end{agradecimentos}
% ---


% ---
% RESUMOS
% ---
% Resumo em português
\setlength{\absparsep}{18pt} % ajusta o espaçamento dos parágrafos do resumo
\begin{resumo}
Desenvolvimentos na área de aprendizagem de máquina representam uma série de oportunidades para diversas áreas de pesquisa e de atuação profissional.
O uso de modelos de aprendizagem de máquina na predição de desfecho para tratamentos em saúde mental tem sido investigado e apresenta resultados promissores.
A aplicação desse tipo de modelo possibilitaria, por exemplo, maior assertividade na personalização de tratamentos psicoterápicos a nível individual.
Apesar dos benefícios potenciais, o tema ainda é pouco discutido nos espaços de pesquisa, formação e atuação em psicologia clínica. Este trabalho tem
por objetivo introduzir o tema de aprendizagem de máquina a psicólogos. Busca-se apresentar, de maneira acessível, os conceitos e terminologia básicos
em aprendizagem de máquina, proporcionando um panorama geral da área. Para complementar a exposição conceitual, é construído um exemplo a partir de dados
de um estudo sobre eficácia de intervenções em psicologia positiva para depressão via internet. Por fim, discutem-se algumas implicações do uso desse tipo
de tecnologia para a pesquisa e para a prática em psicoterapia.

 \textbf{Palavras-chave}: aprendizagem de máquina, psicoterapia, psicologia positiva, depressão.
\end{resumo}

% Resumo em inglês
\begin{resumo}[Abstract]
 \begin{otherlanguage*}{english}
   Advances in machine learning have opened up a multitude of opportunities for various research and professional fields. The use of machine learning models to
   predict treatment outcomes in mental health has been investigated showing promising results. The application of such models would allow greater assertiveness
   in tailoring psychotherapeutic treatments at the individual level. Despite its potential benefits, the topic remains unexplored in the fields of research, training
   and practice in clinical psychology. This work aims to introduce the topic of machine learning to psychologists. It presents the basic machine learning concepts and
   terminology in an accessible manner, providing a comprehensive overview of the field. To complement the conceptual discussion, an example is constructed using data
   from a study on the effectiveness of internet-based positive psychology interventions for depression. Finally, implications of using this type of technology in
   psychotherapy research and practice are discussed.
   \vspace{\onelineskip}
 
   \noindent 
   \textbf{Keywords}: machine learning, psychotherapy, positive psychology, depression.
 \end{otherlanguage*}
\end{resumo}


% ---
% LISTA DE ILUSTRAÇÕES
% ---
\pdfbookmark[0]{\listfigurename}{lof}
\listoffigures*
\cleardoublepage
% ---


% ---
% LISTA DE QUADROS
% ---
% \pdfbookmark[0]{\listofquadrosname}{loq}
% \listofquadros*
% \cleardoublepage
% ---


% ---
% LISTA DE TABELAS
% ---
% \pdfbookmark[0]{\listtablename}{lot}
% \listoftables*
% \cleardoublepage
% ---


% ---
% LISTA DE ABREVIATURAS E SIGLAS
% ---
% \begin{siglas}
%   \item[ABNT] Associação Brasileira de Normas Técnicas
%   \item[abnTeX] ABsurdas Normas para TeX
% \end{siglas}
% ---


% ---
% LISTA DE SÍMBOLOS
% ---
% \begin{simbolos}
%   \item[$ \Gamma $] Letra grega Gama
%   \item[$ \Lambda $] Lambda
%   \item[$ \zeta $] Letra grega minúscula zeta
%   \item[$ \in $] Pertence
% \end{simbolos}
% ---


% ---
% Sumário
% ---
\pdfbookmark[0]{\contentsname}{toc}
\tableofcontents*
\cleardoublepage
% ---


% ----------------------------------------------------------
% ELEMENTOS TEXTUAIS
% ----------------------------------------------------------
\textual

\include{./01-introducao/main.tex}

\include{./02-machine-learning/main.tex}

\include{./03-exemplo/main.tex}

\include{./04-conclusao/main.tex}

% ----------------------------------------------------------
% ELEMENTOS PÓS-TEXTUAIS
% ----------------------------------------------------------
\postextual
% ----------------------------------------------------------


% ----------------------------------------------------------
% REFERÊNCIAS
% ----------------------------------------------------------
\bibliography{references.bib}


% ----------------------------------------------------------
% GLOSSÁRIO
% ----------------------------------------------------------
%
% Consulte o manual da classe abntex2 para orientações sobre o glossário.
%
%\glossary

% ----------------------------------------------------------
% APÊNDICES
% ----------------------------------------------------------

% ---
% Inicia os apêndices
% ---
% \begin{apendicesenv}

% % Imprime uma página indicando o início dos apêndices
% \partapendices

% % ----------------------------------------------------------
% \chapter{Quisque libero justo}
% % ----------------------------------------------------------

% % ----------------------------------------------------------
% \chapter{Nullam elementum urna}
% % ----------------------------------------------------------

% \end{apendicesenv}
% ---


% ----------------------------------------------------------
% ANEXOS
% ----------------------------------------------------------

% ---
% Inicia os anexos
% ---
% \begin{anexosenv}

% % Imprime uma página indicando o início dos anexos
% \partanexos

% % ---
% \chapter{Morbi ultrices rutrum lorem}
% % ---

% % ---
% \chapter{Cras non urna sed feugiat cum}
% % ---


% % ---
% \chapter{Fusce facilisis lacinia dui}
% % ---

% \end{anexosenv}

%---------------------------------------------------------------------
% INDICE REMISSIVO
%---------------------------------------------------------------------
% \phantompart
% \printindex
%---------------------------------------------------------------------
\includepdf{./05-extras/contracapa.pdf}
\end{document}


% ----------------------------------------------------------
% ELEMENTOS PÓS-TEXTUAIS
% ----------------------------------------------------------
\postextual
% ----------------------------------------------------------


% ----------------------------------------------------------
% REFERÊNCIAS
% ----------------------------------------------------------
\bibliography{references.bib}


% ----------------------------------------------------------
% GLOSSÁRIO
% ----------------------------------------------------------
%
% Consulte o manual da classe abntex2 para orientações sobre o glossário.
%
%\glossary

% ----------------------------------------------------------
% APÊNDICES
% ----------------------------------------------------------

% ---
% Inicia os apêndices
% ---
% \begin{apendicesenv}

% % Imprime uma página indicando o início dos apêndices
% \partapendices

% % ----------------------------------------------------------
% \chapter{Quisque libero justo}
% % ----------------------------------------------------------

% % ----------------------------------------------------------
% \chapter{Nullam elementum urna}
% % ----------------------------------------------------------

% \end{apendicesenv}
% ---


% ----------------------------------------------------------
% ANEXOS
% ----------------------------------------------------------

% ---
% Inicia os anexos
% ---
% \begin{anexosenv}

% % Imprime uma página indicando o início dos anexos
% \partanexos

% % ---
% \chapter{Morbi ultrices rutrum lorem}
% % ---

% % ---
% \chapter{Cras non urna sed feugiat cum}
% % ---


% % ---
% \chapter{Fusce facilisis lacinia dui}
% % ---

% \end{anexosenv}

%---------------------------------------------------------------------
% INDICE REMISSIVO
%---------------------------------------------------------------------
% \phantompart
% \printindex
%---------------------------------------------------------------------
\includepdf{./05-extras/contracapa.pdf}
\end{document}


\documentclass[
	% Opções da classe memoir
	12pt,				        % tamanho da fonte
	openany,			        % capítulos começam em pág ímpar (insere página vazia caso preciso)
	oneside,			        % para impressão em recto e verso. Oposto a oneside
	a4paper,			        % tamanho do papel. 
	% Opções da classe abntex2
	%chapter=TITLE,		        % títulos de capítulos convertidos em letras maiúsculas
	%section=TITLE,		        % títulos de seções convertidos em letras maiúsculas
	%subsection=TITLE,	        % títulos de subseções convertidos em letras maiúsculas
	%subsubsection=TITLE,       % títulos de subsubseções convertidos em letras maiúsculas
	% Opções do pacote babel
	english,			        % idioma adicional para hifenização
	brazil				        % o último idioma é o principal do documento
	]{abntex2}

% ---
% Pacotes básicos 
% ---
\usepackage{lmodern}			% Usa a fonte Latin Modern			
\usepackage[T1]{fontenc}		% Selecao de codigos de fonte.
\usepackage[utf8]{inputenc}		% Codificacao do documento (conversão automática dos acentos)
\usepackage{indentfirst}		% Indenta o primeiro parágrafo de cada seção.
\usepackage{color}				% Controle das cores
\usepackage{graphicx}			% Inclusão de gráficos
\usepackage{microtype} 			% para melhorias de justificação
% ---
		

% ---
% Pacotes de citações
% ---
% \usepackage[alf]{abntex2cite}						% Citações padrão ABNT
\usepackage[
  alf,
  abnt-nbr10520=1988,                     % Citações em sentence case
  abnt-emphasize=bf
  ]{abntex2cite}


% ---
% Pacotes adicionais
% ---
\usepackage{pdfpages}	% Inclusão de páginas em PDF


% ---
% Informações para CAPA e FOLHA DE ROSTO
% ---
\titulo{APRENDIZAGEM DE MÁQUINA PARA PREDIÇÃO DE DESFECHOS EM SAÚDE MENTAL: \textmd{UMA INTRODUÇÃO PARA PSICÓLOGOS}}
\autor{GUILHERME AUGUSTO DIAS}
\local{Porto Alegre}
\data{2024}
\orientador{Prof. Dr. Wagner de Lara Machado}
% \coorientador{}
% \instituicao{
%   Pontifícia Universidade Católica do Rio Grande do Sul
%   \par
%   Escola de Ciências da Saúde e da Vida
%   \par
%   Curso de Psicologia}
\tipotrabalho{TCC}
\preambulo{Trabalho de Conclusão de Curso apresentado como requisito para a obtenção do grau de Bacharel em Psicologia pela Escola de Ciências da Saúde e da Vida da Pontifícia Universidade Católica do Rio Grande do Sul.}
% ---

% ---
% Configurações de aparência do PDF final
% ---

% Alterando o aspecto da cor azul
\definecolor{blue}{RGB}{41,5,195}

% Informações do PDF
\makeatletter
\hypersetup{
		pdftitle={Aprendizagem de máquina para predição de desfechos em saúde mental: uma introdução para psicólogos}, 
		pdfauthor={Guilherme Augusto Dias},
    	pdfsubject={\imprimirpreambulo},
	    pdfcreator={LaTeX with abnTeX2},
		pdfkeywords={aprendizagem de máquina}{psicoterapia}{psicologia positiva}{depressão}, 
		colorlinks=true,		% true: colored links; false: boxed links
    	linkcolor=black,		% color of internal links
    	citecolor=black,		% color of links to bibliography
    	filecolor=black,		% color of file links
		urlcolor=black,
		bookmarksdepth=4
}
\makeatother
% --- 


% ---
% Posiciona figuras e tabelas no topo da página quando adicionadas sozinhas
% em um página em branco. Ver https://github.com/abntex/abntex2/issues/170
\makeatletter
\setlength{\@fptop}{5pt} % Set distance from top of page to first float
\makeatother
% ---

% ---
% Possibilita criação de Quadros e Lista de quadros.
% Ver https://github.com/abntex/abntex2/issues/176
%
\newcommand{\quadroname}{Quadro}
\newcommand{\listofquadrosname}{Lista de quadros}

\newfloat[chapter]{quadro}{loq}{\quadroname}
\newlistof{listofquadros}{loq}{\listofquadrosname}
\newlistentry{quadro}{loq}{0}

% Configurações para atender às regras da ABNT
\setfloatadjustment{quadro}{\centering}
\counterwithout{quadro}{chapter}
\renewcommand{\cftquadroname}{\quadroname\space} 
\renewcommand*{\cftquadroaftersnum}{\hfill--\hfill}

\setfloatlocations{quadro}{hbtp} % Ver https://github.com/abntex/abntex2/issues/176
% ---


% --- 
% Espaçamentos entre linhas e parágrafos 
% --- 

% O tamanho do parágrafo é dado por:
\setlength{\parindent}{1.3cm}

% Controle do espaçamento entre um parágrafo e outro:
\setlength{\parskip}{0.2cm}  % tente também \onelineskip

% ---
% Compila o índice
% ---
\makeindex
% ---


% ----
% Início do documento
% ----
\begin{document}

% Seleciona o idioma do documento (conforme pacotes do babel)
\selectlanguage{brazil}

% Retira espaço extra obsoleto entre as frases.
\frenchspacing 

% ----------------------------------------------------------
% ELEMENTOS PRÉ-TEXTUAIS
% ----------------------------------------------------------
\pretextual

% ---
% CAPA
% ---
\includepdf{./05-extras/capa.pdf}
% ---


% ---
% FOLHA DE ROSTO
% (o * indica que haverá a ficha bibliográfica)
% ---
\imprimirfolhaderosto
% ---


% ---
% Ficha bibliografica
% ---

% Isto é um exemplo de Ficha Catalográfica, ou ``Dados internacionais de
% catalogação-na-publicação''. Você pode utilizar este modelo como referência. 
% Porém, provavelmente a biblioteca da sua universidade lhe fornecerá um PDF
% com a ficha catalográfica definitiva após a defesa do trabalho. Quando estiver
% com o documento, salve-o como PDF no diretório do seu projeto e substitua todo
% o conteúdo de implementação deste arquivo pelo comando abaixo:
%
% \usepackage{pdfpages}		% necessário para comando \includepdf
% \begin{fichacatalografica}
%     \includepdf{fig_ficha_catalografica.pdf}
% \end{fichacatalografica}

% \begin{fichacatalografica}
% 	\sffamily
% 	\vspace*{\fill}                             % Posição vertical
% 	\begin{center}                              % Minipage Centralizado
% 	\fbox{\begin{minipage}[c][8cm]{13.5cm}		% Largura
% 	\small
% 	\imprimirautor
% 	%Sobrenome, Nome do autor
	
% 	\hspace{0.5cm} \imprimirtitulo  / \imprimirautor. --
% 	\imprimirlocal, \imprimirdata-
	
% 	\hspace{0.5cm} \thelastpage p. : il. (algumas color.) ; 30 cm.\\
	
% 	\hspace{0.5cm} \imprimirorientadorRotulo~\imprimirorientador\\
	
% 	\hspace{0.5cm}
% 	\parbox[t]{\textwidth}{\imprimirtipotrabalho~--~\imprimirinstituicao,
% 	\imprimirdata.}\\
	
% 	\hspace{0.5cm}
% 		1. Palavra-chave1.
% 		2. Palavra-chave2.
% 		2. Palavra-chave3.
% 		I. Orientador.
% 		II. Universidade xxx.
% 		III. Faculdade de xxx.
% 		IV. Título 			
% 	\end{minipage}}
% 	\end{center}
% \end{fichacatalografica}
% ---


% ---
% FOLHA DE APROVAÇÃO
% ---

% Isto é um exemplo de Folha de aprovação, elemento obrigatório da NBR
% 14724/2011 (seção 4.2.1.3). Você pode utilizar este modelo até a aprovação
% do trabalho. Após isso, substitua todo o conteúdo deste arquivo por uma
% imagem da página assinada pela banca com o comando abaixo:
%
% \begin{folhadeaprovacao}
% \includepdf{folhadeaprovacao_final.pdf}
% \end{folhadeaprovacao}
%
\setlength{\ABNTEXsignskip}{2cm}
\setlength{\ABNTEXsignwidth}{10cm}

\begin{folhadeaprovacao}
  \begin{center}
    {\ABNTEXchapterfont\large\imprimirautor}

    \vspace*{1cm}

    \begin{center}
      \ABNTEXchapterfont\bfseries\Large\imprimirtitulo
    \end{center}

    \vspace*{1cm}
    
    \hspace{.45\textwidth}
    \begin{minipage}{.5\textwidth}
        \imprimirpreambulo
    \end{minipage}%
   \end{center}

   \vspace*{1cm}
        
   \begin{center}
   Aprovado em 17 de julho de 2024.

   \vspace*{1cm}

   Banca examinadora:
   \end{center}

   \vspace*{-1cm}

   \assinatura{Prof. Ma. Caroline Santa Maria Rodrigues -- PUCRS}
   \assinatura{Prof. Dra. Clarissa Pinto Pizarro de Freitas -- PUCRS}
   \assinatura{\imprimirorientador -- PUCRS}
\end{folhadeaprovacao}
% ---


% ---
% Agradecimentos
% ---
\begin{agradecimentos}
Agradeço aos meus pais, Ademilson José Dias e Irinês Porfírio Dias, pelo apoio incondicional que me deram ao longo da vida,
incentivando mesmo as ideias questionáveis como aquela de fazer uma segunda graduação.

A minha irmã, Maria Eduarda Dias, cuja trajetória acadêmica é fonte de inspiração.

A minha namorada, Nathalia Fallavena Ceratti, que acompanhou a esta jornada de perto, oferecendo apoio e compreensão que
fizeram toda a diferença.

Aos meus amigos, Guilherme da Silva Petró, Norton Neves e Uilinton Luciano Marques da Silva, que foram compreensíveis nos
finais de semana em que não pude jogar RPG.

Ao professor Wagner de Lara Machado pela orientação e apoio no desenvolvimento deste trabalho.
\end{agradecimentos}
% ---


% ---
% RESUMOS
% ---
% Resumo em português
\setlength{\absparsep}{18pt} % ajusta o espaçamento dos parágrafos do resumo
\begin{resumo}
Desenvolvimentos na área de aprendizagem de máquina representam uma série de oportunidades para diversas áreas de pesquisa e de atuação profissional.
O uso de modelos de aprendizagem de máquina na predição de desfecho para tratamentos em saúde mental tem sido investigado e apresenta resultados promissores.
A aplicação desse tipo de modelo possibilitaria, por exemplo, maior assertividade na personalização de tratamentos psicoterápicos a nível individual.
Apesar dos benefícios potenciais, o tema ainda é pouco discutido nos espaços de pesquisa, formação e atuação em psicologia clínica. Este trabalho tem
por objetivo introduzir o tema de aprendizagem de máquina a psicólogos. Busca-se apresentar, de maneira acessível, os conceitos e terminologia básicos
em aprendizagem de máquina, proporcionando um panorama geral da área. Para complementar a exposição conceitual, é construído um exemplo a partir de dados
de um estudo sobre eficácia de intervenções em psicologia positiva para depressão via internet. Por fim, discutem-se algumas implicações do uso desse tipo
de tecnologia para a pesquisa e para a prática em psicoterapia.

 \textbf{Palavras-chave}: aprendizagem de máquina, psicoterapia, psicologia positiva, depressão.
\end{resumo}

% Resumo em inglês
\begin{resumo}[Abstract]
 \begin{otherlanguage*}{english}
   Advances in machine learning have opened up a multitude of opportunities for various research and professional fields. The use of machine learning models to
   predict treatment outcomes in mental health has been investigated showing promising results. The application of such models would allow greater assertiveness
   in tailoring psychotherapeutic treatments at the individual level. Despite its potential benefits, the topic remains unexplored in the fields of research, training
   and practice in clinical psychology. This work aims to introduce the topic of machine learning to psychologists. It presents the basic machine learning concepts and
   terminology in an accessible manner, providing a comprehensive overview of the field. To complement the conceptual discussion, an example is constructed using data
   from a study on the effectiveness of internet-based positive psychology interventions for depression. Finally, implications of using this type of technology in
   psychotherapy research and practice are discussed.
   \vspace{\onelineskip}
 
   \noindent 
   \textbf{Keywords}: machine learning, psychotherapy, positive psychology, depression.
 \end{otherlanguage*}
\end{resumo}


% ---
% LISTA DE ILUSTRAÇÕES
% ---
\pdfbookmark[0]{\listfigurename}{lof}
\listoffigures*
\cleardoublepage
% ---


% ---
% LISTA DE QUADROS
% ---
% \pdfbookmark[0]{\listofquadrosname}{loq}
% \listofquadros*
% \cleardoublepage
% ---


% ---
% LISTA DE TABELAS
% ---
% \pdfbookmark[0]{\listtablename}{lot}
% \listoftables*
% \cleardoublepage
% ---


% ---
% LISTA DE ABREVIATURAS E SIGLAS
% ---
% \begin{siglas}
%   \item[ABNT] Associação Brasileira de Normas Técnicas
%   \item[abnTeX] ABsurdas Normas para TeX
% \end{siglas}
% ---


% ---
% LISTA DE SÍMBOLOS
% ---
% \begin{simbolos}
%   \item[$ \Gamma $] Letra grega Gama
%   \item[$ \Lambda $] Lambda
%   \item[$ \zeta $] Letra grega minúscula zeta
%   \item[$ \in $] Pertence
% \end{simbolos}
% ---


% ---
% Sumário
% ---
\pdfbookmark[0]{\contentsname}{toc}
\tableofcontents*
\cleardoublepage
% ---


% ----------------------------------------------------------
% ELEMENTOS TEXTUAIS
% ----------------------------------------------------------
\textual

\documentclass[
	% Opções da classe memoir
	12pt,				        % tamanho da fonte
	openany,			        % capítulos começam em pág ímpar (insere página vazia caso preciso)
	oneside,			        % para impressão em recto e verso. Oposto a oneside
	a4paper,			        % tamanho do papel. 
	% Opções da classe abntex2
	%chapter=TITLE,		        % títulos de capítulos convertidos em letras maiúsculas
	%section=TITLE,		        % títulos de seções convertidos em letras maiúsculas
	%subsection=TITLE,	        % títulos de subseções convertidos em letras maiúsculas
	%subsubsection=TITLE,       % títulos de subsubseções convertidos em letras maiúsculas
	% Opções do pacote babel
	english,			        % idioma adicional para hifenização
	brazil				        % o último idioma é o principal do documento
	]{abntex2}

% ---
% Pacotes básicos 
% ---
\usepackage{lmodern}			% Usa a fonte Latin Modern			
\usepackage[T1]{fontenc}		% Selecao de codigos de fonte.
\usepackage[utf8]{inputenc}		% Codificacao do documento (conversão automática dos acentos)
\usepackage{indentfirst}		% Indenta o primeiro parágrafo de cada seção.
\usepackage{color}				% Controle das cores
\usepackage{graphicx}			% Inclusão de gráficos
\usepackage{microtype} 			% para melhorias de justificação
% ---
		

% ---
% Pacotes de citações
% ---
% \usepackage[alf]{abntex2cite}						% Citações padrão ABNT
\usepackage[
  alf,
  abnt-nbr10520=1988,                     % Citações em sentence case
  abnt-emphasize=bf
  ]{abntex2cite}


% ---
% Pacotes adicionais
% ---
\usepackage{pdfpages}	% Inclusão de páginas em PDF


% ---
% Informações para CAPA e FOLHA DE ROSTO
% ---
\titulo{APRENDIZAGEM DE MÁQUINA PARA PREDIÇÃO DE DESFECHOS EM SAÚDE MENTAL: \textmd{UMA INTRODUÇÃO PARA PSICÓLOGOS}}
\autor{GUILHERME AUGUSTO DIAS}
\local{Porto Alegre}
\data{2024}
\orientador{Prof. Dr. Wagner de Lara Machado}
% \coorientador{}
% \instituicao{
%   Pontifícia Universidade Católica do Rio Grande do Sul
%   \par
%   Escola de Ciências da Saúde e da Vida
%   \par
%   Curso de Psicologia}
\tipotrabalho{TCC}
\preambulo{Trabalho de Conclusão de Curso apresentado como requisito para a obtenção do grau de Bacharel em Psicologia pela Escola de Ciências da Saúde e da Vida da Pontifícia Universidade Católica do Rio Grande do Sul.}
% ---

% ---
% Configurações de aparência do PDF final
% ---

% Alterando o aspecto da cor azul
\definecolor{blue}{RGB}{41,5,195}

% Informações do PDF
\makeatletter
\hypersetup{
		pdftitle={Aprendizagem de máquina para predição de desfechos em saúde mental: uma introdução para psicólogos}, 
		pdfauthor={Guilherme Augusto Dias},
    	pdfsubject={\imprimirpreambulo},
	    pdfcreator={LaTeX with abnTeX2},
		pdfkeywords={aprendizagem de máquina}{psicoterapia}{psicologia positiva}{depressão}, 
		colorlinks=true,		% true: colored links; false: boxed links
    	linkcolor=black,		% color of internal links
    	citecolor=black,		% color of links to bibliography
    	filecolor=black,		% color of file links
		urlcolor=black,
		bookmarksdepth=4
}
\makeatother
% --- 


% ---
% Posiciona figuras e tabelas no topo da página quando adicionadas sozinhas
% em um página em branco. Ver https://github.com/abntex/abntex2/issues/170
\makeatletter
\setlength{\@fptop}{5pt} % Set distance from top of page to first float
\makeatother
% ---

% ---
% Possibilita criação de Quadros e Lista de quadros.
% Ver https://github.com/abntex/abntex2/issues/176
%
\newcommand{\quadroname}{Quadro}
\newcommand{\listofquadrosname}{Lista de quadros}

\newfloat[chapter]{quadro}{loq}{\quadroname}
\newlistof{listofquadros}{loq}{\listofquadrosname}
\newlistentry{quadro}{loq}{0}

% Configurações para atender às regras da ABNT
\setfloatadjustment{quadro}{\centering}
\counterwithout{quadro}{chapter}
\renewcommand{\cftquadroname}{\quadroname\space} 
\renewcommand*{\cftquadroaftersnum}{\hfill--\hfill}

\setfloatlocations{quadro}{hbtp} % Ver https://github.com/abntex/abntex2/issues/176
% ---


% --- 
% Espaçamentos entre linhas e parágrafos 
% --- 

% O tamanho do parágrafo é dado por:
\setlength{\parindent}{1.3cm}

% Controle do espaçamento entre um parágrafo e outro:
\setlength{\parskip}{0.2cm}  % tente também \onelineskip

% ---
% Compila o índice
% ---
\makeindex
% ---


% ----
% Início do documento
% ----
\begin{document}

% Seleciona o idioma do documento (conforme pacotes do babel)
\selectlanguage{brazil}

% Retira espaço extra obsoleto entre as frases.
\frenchspacing 

% ----------------------------------------------------------
% ELEMENTOS PRÉ-TEXTUAIS
% ----------------------------------------------------------
\pretextual

% ---
% CAPA
% ---
\includepdf{./05-extras/capa.pdf}
% ---


% ---
% FOLHA DE ROSTO
% (o * indica que haverá a ficha bibliográfica)
% ---
\imprimirfolhaderosto
% ---


% ---
% Ficha bibliografica
% ---

% Isto é um exemplo de Ficha Catalográfica, ou ``Dados internacionais de
% catalogação-na-publicação''. Você pode utilizar este modelo como referência. 
% Porém, provavelmente a biblioteca da sua universidade lhe fornecerá um PDF
% com a ficha catalográfica definitiva após a defesa do trabalho. Quando estiver
% com o documento, salve-o como PDF no diretório do seu projeto e substitua todo
% o conteúdo de implementação deste arquivo pelo comando abaixo:
%
% \usepackage{pdfpages}		% necessário para comando \includepdf
% \begin{fichacatalografica}
%     \includepdf{fig_ficha_catalografica.pdf}
% \end{fichacatalografica}

% \begin{fichacatalografica}
% 	\sffamily
% 	\vspace*{\fill}                             % Posição vertical
% 	\begin{center}                              % Minipage Centralizado
% 	\fbox{\begin{minipage}[c][8cm]{13.5cm}		% Largura
% 	\small
% 	\imprimirautor
% 	%Sobrenome, Nome do autor
	
% 	\hspace{0.5cm} \imprimirtitulo  / \imprimirautor. --
% 	\imprimirlocal, \imprimirdata-
	
% 	\hspace{0.5cm} \thelastpage p. : il. (algumas color.) ; 30 cm.\\
	
% 	\hspace{0.5cm} \imprimirorientadorRotulo~\imprimirorientador\\
	
% 	\hspace{0.5cm}
% 	\parbox[t]{\textwidth}{\imprimirtipotrabalho~--~\imprimirinstituicao,
% 	\imprimirdata.}\\
	
% 	\hspace{0.5cm}
% 		1. Palavra-chave1.
% 		2. Palavra-chave2.
% 		2. Palavra-chave3.
% 		I. Orientador.
% 		II. Universidade xxx.
% 		III. Faculdade de xxx.
% 		IV. Título 			
% 	\end{minipage}}
% 	\end{center}
% \end{fichacatalografica}
% ---


% ---
% FOLHA DE APROVAÇÃO
% ---

% Isto é um exemplo de Folha de aprovação, elemento obrigatório da NBR
% 14724/2011 (seção 4.2.1.3). Você pode utilizar este modelo até a aprovação
% do trabalho. Após isso, substitua todo o conteúdo deste arquivo por uma
% imagem da página assinada pela banca com o comando abaixo:
%
% \begin{folhadeaprovacao}
% \includepdf{folhadeaprovacao_final.pdf}
% \end{folhadeaprovacao}
%
\setlength{\ABNTEXsignskip}{2cm}
\setlength{\ABNTEXsignwidth}{10cm}

\begin{folhadeaprovacao}
  \begin{center}
    {\ABNTEXchapterfont\large\imprimirautor}

    \vspace*{1cm}

    \begin{center}
      \ABNTEXchapterfont\bfseries\Large\imprimirtitulo
    \end{center}

    \vspace*{1cm}
    
    \hspace{.45\textwidth}
    \begin{minipage}{.5\textwidth}
        \imprimirpreambulo
    \end{minipage}%
   \end{center}

   \vspace*{1cm}
        
   \begin{center}
   Aprovado em 17 de julho de 2024.

   \vspace*{1cm}

   Banca examinadora:
   \end{center}

   \vspace*{-1cm}

   \assinatura{Prof. Ma. Caroline Santa Maria Rodrigues -- PUCRS}
   \assinatura{Prof. Dra. Clarissa Pinto Pizarro de Freitas -- PUCRS}
   \assinatura{\imprimirorientador -- PUCRS}
\end{folhadeaprovacao}
% ---


% ---
% Agradecimentos
% ---
\begin{agradecimentos}
Agradeço aos meus pais, Ademilson José Dias e Irinês Porfírio Dias, pelo apoio incondicional que me deram ao longo da vida,
incentivando mesmo as ideias questionáveis como aquela de fazer uma segunda graduação.

A minha irmã, Maria Eduarda Dias, cuja trajetória acadêmica é fonte de inspiração.

A minha namorada, Nathalia Fallavena Ceratti, que acompanhou a esta jornada de perto, oferecendo apoio e compreensão que
fizeram toda a diferença.

Aos meus amigos, Guilherme da Silva Petró, Norton Neves e Uilinton Luciano Marques da Silva, que foram compreensíveis nos
finais de semana em que não pude jogar RPG.

Ao professor Wagner de Lara Machado pela orientação e apoio no desenvolvimento deste trabalho.
\end{agradecimentos}
% ---


% ---
% RESUMOS
% ---
% Resumo em português
\setlength{\absparsep}{18pt} % ajusta o espaçamento dos parágrafos do resumo
\begin{resumo}
Desenvolvimentos na área de aprendizagem de máquina representam uma série de oportunidades para diversas áreas de pesquisa e de atuação profissional.
O uso de modelos de aprendizagem de máquina na predição de desfecho para tratamentos em saúde mental tem sido investigado e apresenta resultados promissores.
A aplicação desse tipo de modelo possibilitaria, por exemplo, maior assertividade na personalização de tratamentos psicoterápicos a nível individual.
Apesar dos benefícios potenciais, o tema ainda é pouco discutido nos espaços de pesquisa, formação e atuação em psicologia clínica. Este trabalho tem
por objetivo introduzir o tema de aprendizagem de máquina a psicólogos. Busca-se apresentar, de maneira acessível, os conceitos e terminologia básicos
em aprendizagem de máquina, proporcionando um panorama geral da área. Para complementar a exposição conceitual, é construído um exemplo a partir de dados
de um estudo sobre eficácia de intervenções em psicologia positiva para depressão via internet. Por fim, discutem-se algumas implicações do uso desse tipo
de tecnologia para a pesquisa e para a prática em psicoterapia.

 \textbf{Palavras-chave}: aprendizagem de máquina, psicoterapia, psicologia positiva, depressão.
\end{resumo}

% Resumo em inglês
\begin{resumo}[Abstract]
 \begin{otherlanguage*}{english}
   Advances in machine learning have opened up a multitude of opportunities for various research and professional fields. The use of machine learning models to
   predict treatment outcomes in mental health has been investigated showing promising results. The application of such models would allow greater assertiveness
   in tailoring psychotherapeutic treatments at the individual level. Despite its potential benefits, the topic remains unexplored in the fields of research, training
   and practice in clinical psychology. This work aims to introduce the topic of machine learning to psychologists. It presents the basic machine learning concepts and
   terminology in an accessible manner, providing a comprehensive overview of the field. To complement the conceptual discussion, an example is constructed using data
   from a study on the effectiveness of internet-based positive psychology interventions for depression. Finally, implications of using this type of technology in
   psychotherapy research and practice are discussed.
   \vspace{\onelineskip}
 
   \noindent 
   \textbf{Keywords}: machine learning, psychotherapy, positive psychology, depression.
 \end{otherlanguage*}
\end{resumo}


% ---
% LISTA DE ILUSTRAÇÕES
% ---
\pdfbookmark[0]{\listfigurename}{lof}
\listoffigures*
\cleardoublepage
% ---


% ---
% LISTA DE QUADROS
% ---
% \pdfbookmark[0]{\listofquadrosname}{loq}
% \listofquadros*
% \cleardoublepage
% ---


% ---
% LISTA DE TABELAS
% ---
% \pdfbookmark[0]{\listtablename}{lot}
% \listoftables*
% \cleardoublepage
% ---


% ---
% LISTA DE ABREVIATURAS E SIGLAS
% ---
% \begin{siglas}
%   \item[ABNT] Associação Brasileira de Normas Técnicas
%   \item[abnTeX] ABsurdas Normas para TeX
% \end{siglas}
% ---


% ---
% LISTA DE SÍMBOLOS
% ---
% \begin{simbolos}
%   \item[$ \Gamma $] Letra grega Gama
%   \item[$ \Lambda $] Lambda
%   \item[$ \zeta $] Letra grega minúscula zeta
%   \item[$ \in $] Pertence
% \end{simbolos}
% ---


% ---
% Sumário
% ---
\pdfbookmark[0]{\contentsname}{toc}
\tableofcontents*
\cleardoublepage
% ---


% ----------------------------------------------------------
% ELEMENTOS TEXTUAIS
% ----------------------------------------------------------
\textual

\include{./01-introducao/main.tex}

\include{./02-machine-learning/main.tex}

\include{./03-exemplo/main.tex}

\include{./04-conclusao/main.tex}

% ----------------------------------------------------------
% ELEMENTOS PÓS-TEXTUAIS
% ----------------------------------------------------------
\postextual
% ----------------------------------------------------------


% ----------------------------------------------------------
% REFERÊNCIAS
% ----------------------------------------------------------
\bibliography{references.bib}


% ----------------------------------------------------------
% GLOSSÁRIO
% ----------------------------------------------------------
%
% Consulte o manual da classe abntex2 para orientações sobre o glossário.
%
%\glossary

% ----------------------------------------------------------
% APÊNDICES
% ----------------------------------------------------------

% ---
% Inicia os apêndices
% ---
% \begin{apendicesenv}

% % Imprime uma página indicando o início dos apêndices
% \partapendices

% % ----------------------------------------------------------
% \chapter{Quisque libero justo}
% % ----------------------------------------------------------

% % ----------------------------------------------------------
% \chapter{Nullam elementum urna}
% % ----------------------------------------------------------

% \end{apendicesenv}
% ---


% ----------------------------------------------------------
% ANEXOS
% ----------------------------------------------------------

% ---
% Inicia os anexos
% ---
% \begin{anexosenv}

% % Imprime uma página indicando o início dos anexos
% \partanexos

% % ---
% \chapter{Morbi ultrices rutrum lorem}
% % ---

% % ---
% \chapter{Cras non urna sed feugiat cum}
% % ---


% % ---
% \chapter{Fusce facilisis lacinia dui}
% % ---

% \end{anexosenv}

%---------------------------------------------------------------------
% INDICE REMISSIVO
%---------------------------------------------------------------------
% \phantompart
% \printindex
%---------------------------------------------------------------------
\includepdf{./05-extras/contracapa.pdf}
\end{document}


\documentclass[
	% Opções da classe memoir
	12pt,				        % tamanho da fonte
	openany,			        % capítulos começam em pág ímpar (insere página vazia caso preciso)
	oneside,			        % para impressão em recto e verso. Oposto a oneside
	a4paper,			        % tamanho do papel. 
	% Opções da classe abntex2
	%chapter=TITLE,		        % títulos de capítulos convertidos em letras maiúsculas
	%section=TITLE,		        % títulos de seções convertidos em letras maiúsculas
	%subsection=TITLE,	        % títulos de subseções convertidos em letras maiúsculas
	%subsubsection=TITLE,       % títulos de subsubseções convertidos em letras maiúsculas
	% Opções do pacote babel
	english,			        % idioma adicional para hifenização
	brazil				        % o último idioma é o principal do documento
	]{abntex2}

% ---
% Pacotes básicos 
% ---
\usepackage{lmodern}			% Usa a fonte Latin Modern			
\usepackage[T1]{fontenc}		% Selecao de codigos de fonte.
\usepackage[utf8]{inputenc}		% Codificacao do documento (conversão automática dos acentos)
\usepackage{indentfirst}		% Indenta o primeiro parágrafo de cada seção.
\usepackage{color}				% Controle das cores
\usepackage{graphicx}			% Inclusão de gráficos
\usepackage{microtype} 			% para melhorias de justificação
% ---
		

% ---
% Pacotes de citações
% ---
% \usepackage[alf]{abntex2cite}						% Citações padrão ABNT
\usepackage[
  alf,
  abnt-nbr10520=1988,                     % Citações em sentence case
  abnt-emphasize=bf
  ]{abntex2cite}


% ---
% Pacotes adicionais
% ---
\usepackage{pdfpages}	% Inclusão de páginas em PDF


% ---
% Informações para CAPA e FOLHA DE ROSTO
% ---
\titulo{APRENDIZAGEM DE MÁQUINA PARA PREDIÇÃO DE DESFECHOS EM SAÚDE MENTAL: \textmd{UMA INTRODUÇÃO PARA PSICÓLOGOS}}
\autor{GUILHERME AUGUSTO DIAS}
\local{Porto Alegre}
\data{2024}
\orientador{Prof. Dr. Wagner de Lara Machado}
% \coorientador{}
% \instituicao{
%   Pontifícia Universidade Católica do Rio Grande do Sul
%   \par
%   Escola de Ciências da Saúde e da Vida
%   \par
%   Curso de Psicologia}
\tipotrabalho{TCC}
\preambulo{Trabalho de Conclusão de Curso apresentado como requisito para a obtenção do grau de Bacharel em Psicologia pela Escola de Ciências da Saúde e da Vida da Pontifícia Universidade Católica do Rio Grande do Sul.}
% ---

% ---
% Configurações de aparência do PDF final
% ---

% Alterando o aspecto da cor azul
\definecolor{blue}{RGB}{41,5,195}

% Informações do PDF
\makeatletter
\hypersetup{
		pdftitle={Aprendizagem de máquina para predição de desfechos em saúde mental: uma introdução para psicólogos}, 
		pdfauthor={Guilherme Augusto Dias},
    	pdfsubject={\imprimirpreambulo},
	    pdfcreator={LaTeX with abnTeX2},
		pdfkeywords={aprendizagem de máquina}{psicoterapia}{psicologia positiva}{depressão}, 
		colorlinks=true,		% true: colored links; false: boxed links
    	linkcolor=black,		% color of internal links
    	citecolor=black,		% color of links to bibliography
    	filecolor=black,		% color of file links
		urlcolor=black,
		bookmarksdepth=4
}
\makeatother
% --- 


% ---
% Posiciona figuras e tabelas no topo da página quando adicionadas sozinhas
% em um página em branco. Ver https://github.com/abntex/abntex2/issues/170
\makeatletter
\setlength{\@fptop}{5pt} % Set distance from top of page to first float
\makeatother
% ---

% ---
% Possibilita criação de Quadros e Lista de quadros.
% Ver https://github.com/abntex/abntex2/issues/176
%
\newcommand{\quadroname}{Quadro}
\newcommand{\listofquadrosname}{Lista de quadros}

\newfloat[chapter]{quadro}{loq}{\quadroname}
\newlistof{listofquadros}{loq}{\listofquadrosname}
\newlistentry{quadro}{loq}{0}

% Configurações para atender às regras da ABNT
\setfloatadjustment{quadro}{\centering}
\counterwithout{quadro}{chapter}
\renewcommand{\cftquadroname}{\quadroname\space} 
\renewcommand*{\cftquadroaftersnum}{\hfill--\hfill}

\setfloatlocations{quadro}{hbtp} % Ver https://github.com/abntex/abntex2/issues/176
% ---


% --- 
% Espaçamentos entre linhas e parágrafos 
% --- 

% O tamanho do parágrafo é dado por:
\setlength{\parindent}{1.3cm}

% Controle do espaçamento entre um parágrafo e outro:
\setlength{\parskip}{0.2cm}  % tente também \onelineskip

% ---
% Compila o índice
% ---
\makeindex
% ---


% ----
% Início do documento
% ----
\begin{document}

% Seleciona o idioma do documento (conforme pacotes do babel)
\selectlanguage{brazil}

% Retira espaço extra obsoleto entre as frases.
\frenchspacing 

% ----------------------------------------------------------
% ELEMENTOS PRÉ-TEXTUAIS
% ----------------------------------------------------------
\pretextual

% ---
% CAPA
% ---
\includepdf{./05-extras/capa.pdf}
% ---


% ---
% FOLHA DE ROSTO
% (o * indica que haverá a ficha bibliográfica)
% ---
\imprimirfolhaderosto
% ---


% ---
% Ficha bibliografica
% ---

% Isto é um exemplo de Ficha Catalográfica, ou ``Dados internacionais de
% catalogação-na-publicação''. Você pode utilizar este modelo como referência. 
% Porém, provavelmente a biblioteca da sua universidade lhe fornecerá um PDF
% com a ficha catalográfica definitiva após a defesa do trabalho. Quando estiver
% com o documento, salve-o como PDF no diretório do seu projeto e substitua todo
% o conteúdo de implementação deste arquivo pelo comando abaixo:
%
% \usepackage{pdfpages}		% necessário para comando \includepdf
% \begin{fichacatalografica}
%     \includepdf{fig_ficha_catalografica.pdf}
% \end{fichacatalografica}

% \begin{fichacatalografica}
% 	\sffamily
% 	\vspace*{\fill}                             % Posição vertical
% 	\begin{center}                              % Minipage Centralizado
% 	\fbox{\begin{minipage}[c][8cm]{13.5cm}		% Largura
% 	\small
% 	\imprimirautor
% 	%Sobrenome, Nome do autor
	
% 	\hspace{0.5cm} \imprimirtitulo  / \imprimirautor. --
% 	\imprimirlocal, \imprimirdata-
	
% 	\hspace{0.5cm} \thelastpage p. : il. (algumas color.) ; 30 cm.\\
	
% 	\hspace{0.5cm} \imprimirorientadorRotulo~\imprimirorientador\\
	
% 	\hspace{0.5cm}
% 	\parbox[t]{\textwidth}{\imprimirtipotrabalho~--~\imprimirinstituicao,
% 	\imprimirdata.}\\
	
% 	\hspace{0.5cm}
% 		1. Palavra-chave1.
% 		2. Palavra-chave2.
% 		2. Palavra-chave3.
% 		I. Orientador.
% 		II. Universidade xxx.
% 		III. Faculdade de xxx.
% 		IV. Título 			
% 	\end{minipage}}
% 	\end{center}
% \end{fichacatalografica}
% ---


% ---
% FOLHA DE APROVAÇÃO
% ---

% Isto é um exemplo de Folha de aprovação, elemento obrigatório da NBR
% 14724/2011 (seção 4.2.1.3). Você pode utilizar este modelo até a aprovação
% do trabalho. Após isso, substitua todo o conteúdo deste arquivo por uma
% imagem da página assinada pela banca com o comando abaixo:
%
% \begin{folhadeaprovacao}
% \includepdf{folhadeaprovacao_final.pdf}
% \end{folhadeaprovacao}
%
\setlength{\ABNTEXsignskip}{2cm}
\setlength{\ABNTEXsignwidth}{10cm}

\begin{folhadeaprovacao}
  \begin{center}
    {\ABNTEXchapterfont\large\imprimirautor}

    \vspace*{1cm}

    \begin{center}
      \ABNTEXchapterfont\bfseries\Large\imprimirtitulo
    \end{center}

    \vspace*{1cm}
    
    \hspace{.45\textwidth}
    \begin{minipage}{.5\textwidth}
        \imprimirpreambulo
    \end{minipage}%
   \end{center}

   \vspace*{1cm}
        
   \begin{center}
   Aprovado em 17 de julho de 2024.

   \vspace*{1cm}

   Banca examinadora:
   \end{center}

   \vspace*{-1cm}

   \assinatura{Prof. Ma. Caroline Santa Maria Rodrigues -- PUCRS}
   \assinatura{Prof. Dra. Clarissa Pinto Pizarro de Freitas -- PUCRS}
   \assinatura{\imprimirorientador -- PUCRS}
\end{folhadeaprovacao}
% ---


% ---
% Agradecimentos
% ---
\begin{agradecimentos}
Agradeço aos meus pais, Ademilson José Dias e Irinês Porfírio Dias, pelo apoio incondicional que me deram ao longo da vida,
incentivando mesmo as ideias questionáveis como aquela de fazer uma segunda graduação.

A minha irmã, Maria Eduarda Dias, cuja trajetória acadêmica é fonte de inspiração.

A minha namorada, Nathalia Fallavena Ceratti, que acompanhou a esta jornada de perto, oferecendo apoio e compreensão que
fizeram toda a diferença.

Aos meus amigos, Guilherme da Silva Petró, Norton Neves e Uilinton Luciano Marques da Silva, que foram compreensíveis nos
finais de semana em que não pude jogar RPG.

Ao professor Wagner de Lara Machado pela orientação e apoio no desenvolvimento deste trabalho.
\end{agradecimentos}
% ---


% ---
% RESUMOS
% ---
% Resumo em português
\setlength{\absparsep}{18pt} % ajusta o espaçamento dos parágrafos do resumo
\begin{resumo}
Desenvolvimentos na área de aprendizagem de máquina representam uma série de oportunidades para diversas áreas de pesquisa e de atuação profissional.
O uso de modelos de aprendizagem de máquina na predição de desfecho para tratamentos em saúde mental tem sido investigado e apresenta resultados promissores.
A aplicação desse tipo de modelo possibilitaria, por exemplo, maior assertividade na personalização de tratamentos psicoterápicos a nível individual.
Apesar dos benefícios potenciais, o tema ainda é pouco discutido nos espaços de pesquisa, formação e atuação em psicologia clínica. Este trabalho tem
por objetivo introduzir o tema de aprendizagem de máquina a psicólogos. Busca-se apresentar, de maneira acessível, os conceitos e terminologia básicos
em aprendizagem de máquina, proporcionando um panorama geral da área. Para complementar a exposição conceitual, é construído um exemplo a partir de dados
de um estudo sobre eficácia de intervenções em psicologia positiva para depressão via internet. Por fim, discutem-se algumas implicações do uso desse tipo
de tecnologia para a pesquisa e para a prática em psicoterapia.

 \textbf{Palavras-chave}: aprendizagem de máquina, psicoterapia, psicologia positiva, depressão.
\end{resumo}

% Resumo em inglês
\begin{resumo}[Abstract]
 \begin{otherlanguage*}{english}
   Advances in machine learning have opened up a multitude of opportunities for various research and professional fields. The use of machine learning models to
   predict treatment outcomes in mental health has been investigated showing promising results. The application of such models would allow greater assertiveness
   in tailoring psychotherapeutic treatments at the individual level. Despite its potential benefits, the topic remains unexplored in the fields of research, training
   and practice in clinical psychology. This work aims to introduce the topic of machine learning to psychologists. It presents the basic machine learning concepts and
   terminology in an accessible manner, providing a comprehensive overview of the field. To complement the conceptual discussion, an example is constructed using data
   from a study on the effectiveness of internet-based positive psychology interventions for depression. Finally, implications of using this type of technology in
   psychotherapy research and practice are discussed.
   \vspace{\onelineskip}
 
   \noindent 
   \textbf{Keywords}: machine learning, psychotherapy, positive psychology, depression.
 \end{otherlanguage*}
\end{resumo}


% ---
% LISTA DE ILUSTRAÇÕES
% ---
\pdfbookmark[0]{\listfigurename}{lof}
\listoffigures*
\cleardoublepage
% ---


% ---
% LISTA DE QUADROS
% ---
% \pdfbookmark[0]{\listofquadrosname}{loq}
% \listofquadros*
% \cleardoublepage
% ---


% ---
% LISTA DE TABELAS
% ---
% \pdfbookmark[0]{\listtablename}{lot}
% \listoftables*
% \cleardoublepage
% ---


% ---
% LISTA DE ABREVIATURAS E SIGLAS
% ---
% \begin{siglas}
%   \item[ABNT] Associação Brasileira de Normas Técnicas
%   \item[abnTeX] ABsurdas Normas para TeX
% \end{siglas}
% ---


% ---
% LISTA DE SÍMBOLOS
% ---
% \begin{simbolos}
%   \item[$ \Gamma $] Letra grega Gama
%   \item[$ \Lambda $] Lambda
%   \item[$ \zeta $] Letra grega minúscula zeta
%   \item[$ \in $] Pertence
% \end{simbolos}
% ---


% ---
% Sumário
% ---
\pdfbookmark[0]{\contentsname}{toc}
\tableofcontents*
\cleardoublepage
% ---


% ----------------------------------------------------------
% ELEMENTOS TEXTUAIS
% ----------------------------------------------------------
\textual

\include{./01-introducao/main.tex}

\include{./02-machine-learning/main.tex}

\include{./03-exemplo/main.tex}

\include{./04-conclusao/main.tex}

% ----------------------------------------------------------
% ELEMENTOS PÓS-TEXTUAIS
% ----------------------------------------------------------
\postextual
% ----------------------------------------------------------


% ----------------------------------------------------------
% REFERÊNCIAS
% ----------------------------------------------------------
\bibliography{references.bib}


% ----------------------------------------------------------
% GLOSSÁRIO
% ----------------------------------------------------------
%
% Consulte o manual da classe abntex2 para orientações sobre o glossário.
%
%\glossary

% ----------------------------------------------------------
% APÊNDICES
% ----------------------------------------------------------

% ---
% Inicia os apêndices
% ---
% \begin{apendicesenv}

% % Imprime uma página indicando o início dos apêndices
% \partapendices

% % ----------------------------------------------------------
% \chapter{Quisque libero justo}
% % ----------------------------------------------------------

% % ----------------------------------------------------------
% \chapter{Nullam elementum urna}
% % ----------------------------------------------------------

% \end{apendicesenv}
% ---


% ----------------------------------------------------------
% ANEXOS
% ----------------------------------------------------------

% ---
% Inicia os anexos
% ---
% \begin{anexosenv}

% % Imprime uma página indicando o início dos anexos
% \partanexos

% % ---
% \chapter{Morbi ultrices rutrum lorem}
% % ---

% % ---
% \chapter{Cras non urna sed feugiat cum}
% % ---


% % ---
% \chapter{Fusce facilisis lacinia dui}
% % ---

% \end{anexosenv}

%---------------------------------------------------------------------
% INDICE REMISSIVO
%---------------------------------------------------------------------
% \phantompart
% \printindex
%---------------------------------------------------------------------
\includepdf{./05-extras/contracapa.pdf}
\end{document}


\documentclass[
	% Opções da classe memoir
	12pt,				        % tamanho da fonte
	openany,			        % capítulos começam em pág ímpar (insere página vazia caso preciso)
	oneside,			        % para impressão em recto e verso. Oposto a oneside
	a4paper,			        % tamanho do papel. 
	% Opções da classe abntex2
	%chapter=TITLE,		        % títulos de capítulos convertidos em letras maiúsculas
	%section=TITLE,		        % títulos de seções convertidos em letras maiúsculas
	%subsection=TITLE,	        % títulos de subseções convertidos em letras maiúsculas
	%subsubsection=TITLE,       % títulos de subsubseções convertidos em letras maiúsculas
	% Opções do pacote babel
	english,			        % idioma adicional para hifenização
	brazil				        % o último idioma é o principal do documento
	]{abntex2}

% ---
% Pacotes básicos 
% ---
\usepackage{lmodern}			% Usa a fonte Latin Modern			
\usepackage[T1]{fontenc}		% Selecao de codigos de fonte.
\usepackage[utf8]{inputenc}		% Codificacao do documento (conversão automática dos acentos)
\usepackage{indentfirst}		% Indenta o primeiro parágrafo de cada seção.
\usepackage{color}				% Controle das cores
\usepackage{graphicx}			% Inclusão de gráficos
\usepackage{microtype} 			% para melhorias de justificação
% ---
		

% ---
% Pacotes de citações
% ---
% \usepackage[alf]{abntex2cite}						% Citações padrão ABNT
\usepackage[
  alf,
  abnt-nbr10520=1988,                     % Citações em sentence case
  abnt-emphasize=bf
  ]{abntex2cite}


% ---
% Pacotes adicionais
% ---
\usepackage{pdfpages}	% Inclusão de páginas em PDF


% ---
% Informações para CAPA e FOLHA DE ROSTO
% ---
\titulo{APRENDIZAGEM DE MÁQUINA PARA PREDIÇÃO DE DESFECHOS EM SAÚDE MENTAL: \textmd{UMA INTRODUÇÃO PARA PSICÓLOGOS}}
\autor{GUILHERME AUGUSTO DIAS}
\local{Porto Alegre}
\data{2024}
\orientador{Prof. Dr. Wagner de Lara Machado}
% \coorientador{}
% \instituicao{
%   Pontifícia Universidade Católica do Rio Grande do Sul
%   \par
%   Escola de Ciências da Saúde e da Vida
%   \par
%   Curso de Psicologia}
\tipotrabalho{TCC}
\preambulo{Trabalho de Conclusão de Curso apresentado como requisito para a obtenção do grau de Bacharel em Psicologia pela Escola de Ciências da Saúde e da Vida da Pontifícia Universidade Católica do Rio Grande do Sul.}
% ---

% ---
% Configurações de aparência do PDF final
% ---

% Alterando o aspecto da cor azul
\definecolor{blue}{RGB}{41,5,195}

% Informações do PDF
\makeatletter
\hypersetup{
		pdftitle={Aprendizagem de máquina para predição de desfechos em saúde mental: uma introdução para psicólogos}, 
		pdfauthor={Guilherme Augusto Dias},
    	pdfsubject={\imprimirpreambulo},
	    pdfcreator={LaTeX with abnTeX2},
		pdfkeywords={aprendizagem de máquina}{psicoterapia}{psicologia positiva}{depressão}, 
		colorlinks=true,		% true: colored links; false: boxed links
    	linkcolor=black,		% color of internal links
    	citecolor=black,		% color of links to bibliography
    	filecolor=black,		% color of file links
		urlcolor=black,
		bookmarksdepth=4
}
\makeatother
% --- 


% ---
% Posiciona figuras e tabelas no topo da página quando adicionadas sozinhas
% em um página em branco. Ver https://github.com/abntex/abntex2/issues/170
\makeatletter
\setlength{\@fptop}{5pt} % Set distance from top of page to first float
\makeatother
% ---

% ---
% Possibilita criação de Quadros e Lista de quadros.
% Ver https://github.com/abntex/abntex2/issues/176
%
\newcommand{\quadroname}{Quadro}
\newcommand{\listofquadrosname}{Lista de quadros}

\newfloat[chapter]{quadro}{loq}{\quadroname}
\newlistof{listofquadros}{loq}{\listofquadrosname}
\newlistentry{quadro}{loq}{0}

% Configurações para atender às regras da ABNT
\setfloatadjustment{quadro}{\centering}
\counterwithout{quadro}{chapter}
\renewcommand{\cftquadroname}{\quadroname\space} 
\renewcommand*{\cftquadroaftersnum}{\hfill--\hfill}

\setfloatlocations{quadro}{hbtp} % Ver https://github.com/abntex/abntex2/issues/176
% ---


% --- 
% Espaçamentos entre linhas e parágrafos 
% --- 

% O tamanho do parágrafo é dado por:
\setlength{\parindent}{1.3cm}

% Controle do espaçamento entre um parágrafo e outro:
\setlength{\parskip}{0.2cm}  % tente também \onelineskip

% ---
% Compila o índice
% ---
\makeindex
% ---


% ----
% Início do documento
% ----
\begin{document}

% Seleciona o idioma do documento (conforme pacotes do babel)
\selectlanguage{brazil}

% Retira espaço extra obsoleto entre as frases.
\frenchspacing 

% ----------------------------------------------------------
% ELEMENTOS PRÉ-TEXTUAIS
% ----------------------------------------------------------
\pretextual

% ---
% CAPA
% ---
\includepdf{./05-extras/capa.pdf}
% ---


% ---
% FOLHA DE ROSTO
% (o * indica que haverá a ficha bibliográfica)
% ---
\imprimirfolhaderosto
% ---


% ---
% Ficha bibliografica
% ---

% Isto é um exemplo de Ficha Catalográfica, ou ``Dados internacionais de
% catalogação-na-publicação''. Você pode utilizar este modelo como referência. 
% Porém, provavelmente a biblioteca da sua universidade lhe fornecerá um PDF
% com a ficha catalográfica definitiva após a defesa do trabalho. Quando estiver
% com o documento, salve-o como PDF no diretório do seu projeto e substitua todo
% o conteúdo de implementação deste arquivo pelo comando abaixo:
%
% \usepackage{pdfpages}		% necessário para comando \includepdf
% \begin{fichacatalografica}
%     \includepdf{fig_ficha_catalografica.pdf}
% \end{fichacatalografica}

% \begin{fichacatalografica}
% 	\sffamily
% 	\vspace*{\fill}                             % Posição vertical
% 	\begin{center}                              % Minipage Centralizado
% 	\fbox{\begin{minipage}[c][8cm]{13.5cm}		% Largura
% 	\small
% 	\imprimirautor
% 	%Sobrenome, Nome do autor
	
% 	\hspace{0.5cm} \imprimirtitulo  / \imprimirautor. --
% 	\imprimirlocal, \imprimirdata-
	
% 	\hspace{0.5cm} \thelastpage p. : il. (algumas color.) ; 30 cm.\\
	
% 	\hspace{0.5cm} \imprimirorientadorRotulo~\imprimirorientador\\
	
% 	\hspace{0.5cm}
% 	\parbox[t]{\textwidth}{\imprimirtipotrabalho~--~\imprimirinstituicao,
% 	\imprimirdata.}\\
	
% 	\hspace{0.5cm}
% 		1. Palavra-chave1.
% 		2. Palavra-chave2.
% 		2. Palavra-chave3.
% 		I. Orientador.
% 		II. Universidade xxx.
% 		III. Faculdade de xxx.
% 		IV. Título 			
% 	\end{minipage}}
% 	\end{center}
% \end{fichacatalografica}
% ---


% ---
% FOLHA DE APROVAÇÃO
% ---

% Isto é um exemplo de Folha de aprovação, elemento obrigatório da NBR
% 14724/2011 (seção 4.2.1.3). Você pode utilizar este modelo até a aprovação
% do trabalho. Após isso, substitua todo o conteúdo deste arquivo por uma
% imagem da página assinada pela banca com o comando abaixo:
%
% \begin{folhadeaprovacao}
% \includepdf{folhadeaprovacao_final.pdf}
% \end{folhadeaprovacao}
%
\setlength{\ABNTEXsignskip}{2cm}
\setlength{\ABNTEXsignwidth}{10cm}

\begin{folhadeaprovacao}
  \begin{center}
    {\ABNTEXchapterfont\large\imprimirautor}

    \vspace*{1cm}

    \begin{center}
      \ABNTEXchapterfont\bfseries\Large\imprimirtitulo
    \end{center}

    \vspace*{1cm}
    
    \hspace{.45\textwidth}
    \begin{minipage}{.5\textwidth}
        \imprimirpreambulo
    \end{minipage}%
   \end{center}

   \vspace*{1cm}
        
   \begin{center}
   Aprovado em 17 de julho de 2024.

   \vspace*{1cm}

   Banca examinadora:
   \end{center}

   \vspace*{-1cm}

   \assinatura{Prof. Ma. Caroline Santa Maria Rodrigues -- PUCRS}
   \assinatura{Prof. Dra. Clarissa Pinto Pizarro de Freitas -- PUCRS}
   \assinatura{\imprimirorientador -- PUCRS}
\end{folhadeaprovacao}
% ---


% ---
% Agradecimentos
% ---
\begin{agradecimentos}
Agradeço aos meus pais, Ademilson José Dias e Irinês Porfírio Dias, pelo apoio incondicional que me deram ao longo da vida,
incentivando mesmo as ideias questionáveis como aquela de fazer uma segunda graduação.

A minha irmã, Maria Eduarda Dias, cuja trajetória acadêmica é fonte de inspiração.

A minha namorada, Nathalia Fallavena Ceratti, que acompanhou a esta jornada de perto, oferecendo apoio e compreensão que
fizeram toda a diferença.

Aos meus amigos, Guilherme da Silva Petró, Norton Neves e Uilinton Luciano Marques da Silva, que foram compreensíveis nos
finais de semana em que não pude jogar RPG.

Ao professor Wagner de Lara Machado pela orientação e apoio no desenvolvimento deste trabalho.
\end{agradecimentos}
% ---


% ---
% RESUMOS
% ---
% Resumo em português
\setlength{\absparsep}{18pt} % ajusta o espaçamento dos parágrafos do resumo
\begin{resumo}
Desenvolvimentos na área de aprendizagem de máquina representam uma série de oportunidades para diversas áreas de pesquisa e de atuação profissional.
O uso de modelos de aprendizagem de máquina na predição de desfecho para tratamentos em saúde mental tem sido investigado e apresenta resultados promissores.
A aplicação desse tipo de modelo possibilitaria, por exemplo, maior assertividade na personalização de tratamentos psicoterápicos a nível individual.
Apesar dos benefícios potenciais, o tema ainda é pouco discutido nos espaços de pesquisa, formação e atuação em psicologia clínica. Este trabalho tem
por objetivo introduzir o tema de aprendizagem de máquina a psicólogos. Busca-se apresentar, de maneira acessível, os conceitos e terminologia básicos
em aprendizagem de máquina, proporcionando um panorama geral da área. Para complementar a exposição conceitual, é construído um exemplo a partir de dados
de um estudo sobre eficácia de intervenções em psicologia positiva para depressão via internet. Por fim, discutem-se algumas implicações do uso desse tipo
de tecnologia para a pesquisa e para a prática em psicoterapia.

 \textbf{Palavras-chave}: aprendizagem de máquina, psicoterapia, psicologia positiva, depressão.
\end{resumo}

% Resumo em inglês
\begin{resumo}[Abstract]
 \begin{otherlanguage*}{english}
   Advances in machine learning have opened up a multitude of opportunities for various research and professional fields. The use of machine learning models to
   predict treatment outcomes in mental health has been investigated showing promising results. The application of such models would allow greater assertiveness
   in tailoring psychotherapeutic treatments at the individual level. Despite its potential benefits, the topic remains unexplored in the fields of research, training
   and practice in clinical psychology. This work aims to introduce the topic of machine learning to psychologists. It presents the basic machine learning concepts and
   terminology in an accessible manner, providing a comprehensive overview of the field. To complement the conceptual discussion, an example is constructed using data
   from a study on the effectiveness of internet-based positive psychology interventions for depression. Finally, implications of using this type of technology in
   psychotherapy research and practice are discussed.
   \vspace{\onelineskip}
 
   \noindent 
   \textbf{Keywords}: machine learning, psychotherapy, positive psychology, depression.
 \end{otherlanguage*}
\end{resumo}


% ---
% LISTA DE ILUSTRAÇÕES
% ---
\pdfbookmark[0]{\listfigurename}{lof}
\listoffigures*
\cleardoublepage
% ---


% ---
% LISTA DE QUADROS
% ---
% \pdfbookmark[0]{\listofquadrosname}{loq}
% \listofquadros*
% \cleardoublepage
% ---


% ---
% LISTA DE TABELAS
% ---
% \pdfbookmark[0]{\listtablename}{lot}
% \listoftables*
% \cleardoublepage
% ---


% ---
% LISTA DE ABREVIATURAS E SIGLAS
% ---
% \begin{siglas}
%   \item[ABNT] Associação Brasileira de Normas Técnicas
%   \item[abnTeX] ABsurdas Normas para TeX
% \end{siglas}
% ---


% ---
% LISTA DE SÍMBOLOS
% ---
% \begin{simbolos}
%   \item[$ \Gamma $] Letra grega Gama
%   \item[$ \Lambda $] Lambda
%   \item[$ \zeta $] Letra grega minúscula zeta
%   \item[$ \in $] Pertence
% \end{simbolos}
% ---


% ---
% Sumário
% ---
\pdfbookmark[0]{\contentsname}{toc}
\tableofcontents*
\cleardoublepage
% ---


% ----------------------------------------------------------
% ELEMENTOS TEXTUAIS
% ----------------------------------------------------------
\textual

\include{./01-introducao/main.tex}

\include{./02-machine-learning/main.tex}

\include{./03-exemplo/main.tex}

\include{./04-conclusao/main.tex}

% ----------------------------------------------------------
% ELEMENTOS PÓS-TEXTUAIS
% ----------------------------------------------------------
\postextual
% ----------------------------------------------------------


% ----------------------------------------------------------
% REFERÊNCIAS
% ----------------------------------------------------------
\bibliography{references.bib}


% ----------------------------------------------------------
% GLOSSÁRIO
% ----------------------------------------------------------
%
% Consulte o manual da classe abntex2 para orientações sobre o glossário.
%
%\glossary

% ----------------------------------------------------------
% APÊNDICES
% ----------------------------------------------------------

% ---
% Inicia os apêndices
% ---
% \begin{apendicesenv}

% % Imprime uma página indicando o início dos apêndices
% \partapendices

% % ----------------------------------------------------------
% \chapter{Quisque libero justo}
% % ----------------------------------------------------------

% % ----------------------------------------------------------
% \chapter{Nullam elementum urna}
% % ----------------------------------------------------------

% \end{apendicesenv}
% ---


% ----------------------------------------------------------
% ANEXOS
% ----------------------------------------------------------

% ---
% Inicia os anexos
% ---
% \begin{anexosenv}

% % Imprime uma página indicando o início dos anexos
% \partanexos

% % ---
% \chapter{Morbi ultrices rutrum lorem}
% % ---

% % ---
% \chapter{Cras non urna sed feugiat cum}
% % ---


% % ---
% \chapter{Fusce facilisis lacinia dui}
% % ---

% \end{anexosenv}

%---------------------------------------------------------------------
% INDICE REMISSIVO
%---------------------------------------------------------------------
% \phantompart
% \printindex
%---------------------------------------------------------------------
\includepdf{./05-extras/contracapa.pdf}
\end{document}


\documentclass[
	% Opções da classe memoir
	12pt,				        % tamanho da fonte
	openany,			        % capítulos começam em pág ímpar (insere página vazia caso preciso)
	oneside,			        % para impressão em recto e verso. Oposto a oneside
	a4paper,			        % tamanho do papel. 
	% Opções da classe abntex2
	%chapter=TITLE,		        % títulos de capítulos convertidos em letras maiúsculas
	%section=TITLE,		        % títulos de seções convertidos em letras maiúsculas
	%subsection=TITLE,	        % títulos de subseções convertidos em letras maiúsculas
	%subsubsection=TITLE,       % títulos de subsubseções convertidos em letras maiúsculas
	% Opções do pacote babel
	english,			        % idioma adicional para hifenização
	brazil				        % o último idioma é o principal do documento
	]{abntex2}

% ---
% Pacotes básicos 
% ---
\usepackage{lmodern}			% Usa a fonte Latin Modern			
\usepackage[T1]{fontenc}		% Selecao de codigos de fonte.
\usepackage[utf8]{inputenc}		% Codificacao do documento (conversão automática dos acentos)
\usepackage{indentfirst}		% Indenta o primeiro parágrafo de cada seção.
\usepackage{color}				% Controle das cores
\usepackage{graphicx}			% Inclusão de gráficos
\usepackage{microtype} 			% para melhorias de justificação
% ---
		

% ---
% Pacotes de citações
% ---
% \usepackage[alf]{abntex2cite}						% Citações padrão ABNT
\usepackage[
  alf,
  abnt-nbr10520=1988,                     % Citações em sentence case
  abnt-emphasize=bf
  ]{abntex2cite}


% ---
% Pacotes adicionais
% ---
\usepackage{pdfpages}	% Inclusão de páginas em PDF


% ---
% Informações para CAPA e FOLHA DE ROSTO
% ---
\titulo{APRENDIZAGEM DE MÁQUINA PARA PREDIÇÃO DE DESFECHOS EM SAÚDE MENTAL: \textmd{UMA INTRODUÇÃO PARA PSICÓLOGOS}}
\autor{GUILHERME AUGUSTO DIAS}
\local{Porto Alegre}
\data{2024}
\orientador{Prof. Dr. Wagner de Lara Machado}
% \coorientador{}
% \instituicao{
%   Pontifícia Universidade Católica do Rio Grande do Sul
%   \par
%   Escola de Ciências da Saúde e da Vida
%   \par
%   Curso de Psicologia}
\tipotrabalho{TCC}
\preambulo{Trabalho de Conclusão de Curso apresentado como requisito para a obtenção do grau de Bacharel em Psicologia pela Escola de Ciências da Saúde e da Vida da Pontifícia Universidade Católica do Rio Grande do Sul.}
% ---

% ---
% Configurações de aparência do PDF final
% ---

% Alterando o aspecto da cor azul
\definecolor{blue}{RGB}{41,5,195}

% Informações do PDF
\makeatletter
\hypersetup{
		pdftitle={Aprendizagem de máquina para predição de desfechos em saúde mental: uma introdução para psicólogos}, 
		pdfauthor={Guilherme Augusto Dias},
    	pdfsubject={\imprimirpreambulo},
	    pdfcreator={LaTeX with abnTeX2},
		pdfkeywords={aprendizagem de máquina}{psicoterapia}{psicologia positiva}{depressão}, 
		colorlinks=true,		% true: colored links; false: boxed links
    	linkcolor=black,		% color of internal links
    	citecolor=black,		% color of links to bibliography
    	filecolor=black,		% color of file links
		urlcolor=black,
		bookmarksdepth=4
}
\makeatother
% --- 


% ---
% Posiciona figuras e tabelas no topo da página quando adicionadas sozinhas
% em um página em branco. Ver https://github.com/abntex/abntex2/issues/170
\makeatletter
\setlength{\@fptop}{5pt} % Set distance from top of page to first float
\makeatother
% ---

% ---
% Possibilita criação de Quadros e Lista de quadros.
% Ver https://github.com/abntex/abntex2/issues/176
%
\newcommand{\quadroname}{Quadro}
\newcommand{\listofquadrosname}{Lista de quadros}

\newfloat[chapter]{quadro}{loq}{\quadroname}
\newlistof{listofquadros}{loq}{\listofquadrosname}
\newlistentry{quadro}{loq}{0}

% Configurações para atender às regras da ABNT
\setfloatadjustment{quadro}{\centering}
\counterwithout{quadro}{chapter}
\renewcommand{\cftquadroname}{\quadroname\space} 
\renewcommand*{\cftquadroaftersnum}{\hfill--\hfill}

\setfloatlocations{quadro}{hbtp} % Ver https://github.com/abntex/abntex2/issues/176
% ---


% --- 
% Espaçamentos entre linhas e parágrafos 
% --- 

% O tamanho do parágrafo é dado por:
\setlength{\parindent}{1.3cm}

% Controle do espaçamento entre um parágrafo e outro:
\setlength{\parskip}{0.2cm}  % tente também \onelineskip

% ---
% Compila o índice
% ---
\makeindex
% ---


% ----
% Início do documento
% ----
\begin{document}

% Seleciona o idioma do documento (conforme pacotes do babel)
\selectlanguage{brazil}

% Retira espaço extra obsoleto entre as frases.
\frenchspacing 

% ----------------------------------------------------------
% ELEMENTOS PRÉ-TEXTUAIS
% ----------------------------------------------------------
\pretextual

% ---
% CAPA
% ---
\includepdf{./05-extras/capa.pdf}
% ---


% ---
% FOLHA DE ROSTO
% (o * indica que haverá a ficha bibliográfica)
% ---
\imprimirfolhaderosto
% ---


% ---
% Ficha bibliografica
% ---

% Isto é um exemplo de Ficha Catalográfica, ou ``Dados internacionais de
% catalogação-na-publicação''. Você pode utilizar este modelo como referência. 
% Porém, provavelmente a biblioteca da sua universidade lhe fornecerá um PDF
% com a ficha catalográfica definitiva após a defesa do trabalho. Quando estiver
% com o documento, salve-o como PDF no diretório do seu projeto e substitua todo
% o conteúdo de implementação deste arquivo pelo comando abaixo:
%
% \usepackage{pdfpages}		% necessário para comando \includepdf
% \begin{fichacatalografica}
%     \includepdf{fig_ficha_catalografica.pdf}
% \end{fichacatalografica}

% \begin{fichacatalografica}
% 	\sffamily
% 	\vspace*{\fill}                             % Posição vertical
% 	\begin{center}                              % Minipage Centralizado
% 	\fbox{\begin{minipage}[c][8cm]{13.5cm}		% Largura
% 	\small
% 	\imprimirautor
% 	%Sobrenome, Nome do autor
	
% 	\hspace{0.5cm} \imprimirtitulo  / \imprimirautor. --
% 	\imprimirlocal, \imprimirdata-
	
% 	\hspace{0.5cm} \thelastpage p. : il. (algumas color.) ; 30 cm.\\
	
% 	\hspace{0.5cm} \imprimirorientadorRotulo~\imprimirorientador\\
	
% 	\hspace{0.5cm}
% 	\parbox[t]{\textwidth}{\imprimirtipotrabalho~--~\imprimirinstituicao,
% 	\imprimirdata.}\\
	
% 	\hspace{0.5cm}
% 		1. Palavra-chave1.
% 		2. Palavra-chave2.
% 		2. Palavra-chave3.
% 		I. Orientador.
% 		II. Universidade xxx.
% 		III. Faculdade de xxx.
% 		IV. Título 			
% 	\end{minipage}}
% 	\end{center}
% \end{fichacatalografica}
% ---


% ---
% FOLHA DE APROVAÇÃO
% ---

% Isto é um exemplo de Folha de aprovação, elemento obrigatório da NBR
% 14724/2011 (seção 4.2.1.3). Você pode utilizar este modelo até a aprovação
% do trabalho. Após isso, substitua todo o conteúdo deste arquivo por uma
% imagem da página assinada pela banca com o comando abaixo:
%
% \begin{folhadeaprovacao}
% \includepdf{folhadeaprovacao_final.pdf}
% \end{folhadeaprovacao}
%
\setlength{\ABNTEXsignskip}{2cm}
\setlength{\ABNTEXsignwidth}{10cm}

\begin{folhadeaprovacao}
  \begin{center}
    {\ABNTEXchapterfont\large\imprimirautor}

    \vspace*{1cm}

    \begin{center}
      \ABNTEXchapterfont\bfseries\Large\imprimirtitulo
    \end{center}

    \vspace*{1cm}
    
    \hspace{.45\textwidth}
    \begin{minipage}{.5\textwidth}
        \imprimirpreambulo
    \end{minipage}%
   \end{center}

   \vspace*{1cm}
        
   \begin{center}
   Aprovado em 17 de julho de 2024.

   \vspace*{1cm}

   Banca examinadora:
   \end{center}

   \vspace*{-1cm}

   \assinatura{Prof. Ma. Caroline Santa Maria Rodrigues -- PUCRS}
   \assinatura{Prof. Dra. Clarissa Pinto Pizarro de Freitas -- PUCRS}
   \assinatura{\imprimirorientador -- PUCRS}
\end{folhadeaprovacao}
% ---


% ---
% Agradecimentos
% ---
\begin{agradecimentos}
Agradeço aos meus pais, Ademilson José Dias e Irinês Porfírio Dias, pelo apoio incondicional que me deram ao longo da vida,
incentivando mesmo as ideias questionáveis como aquela de fazer uma segunda graduação.

A minha irmã, Maria Eduarda Dias, cuja trajetória acadêmica é fonte de inspiração.

A minha namorada, Nathalia Fallavena Ceratti, que acompanhou a esta jornada de perto, oferecendo apoio e compreensão que
fizeram toda a diferença.

Aos meus amigos, Guilherme da Silva Petró, Norton Neves e Uilinton Luciano Marques da Silva, que foram compreensíveis nos
finais de semana em que não pude jogar RPG.

Ao professor Wagner de Lara Machado pela orientação e apoio no desenvolvimento deste trabalho.
\end{agradecimentos}
% ---


% ---
% RESUMOS
% ---
% Resumo em português
\setlength{\absparsep}{18pt} % ajusta o espaçamento dos parágrafos do resumo
\begin{resumo}
Desenvolvimentos na área de aprendizagem de máquina representam uma série de oportunidades para diversas áreas de pesquisa e de atuação profissional.
O uso de modelos de aprendizagem de máquina na predição de desfecho para tratamentos em saúde mental tem sido investigado e apresenta resultados promissores.
A aplicação desse tipo de modelo possibilitaria, por exemplo, maior assertividade na personalização de tratamentos psicoterápicos a nível individual.
Apesar dos benefícios potenciais, o tema ainda é pouco discutido nos espaços de pesquisa, formação e atuação em psicologia clínica. Este trabalho tem
por objetivo introduzir o tema de aprendizagem de máquina a psicólogos. Busca-se apresentar, de maneira acessível, os conceitos e terminologia básicos
em aprendizagem de máquina, proporcionando um panorama geral da área. Para complementar a exposição conceitual, é construído um exemplo a partir de dados
de um estudo sobre eficácia de intervenções em psicologia positiva para depressão via internet. Por fim, discutem-se algumas implicações do uso desse tipo
de tecnologia para a pesquisa e para a prática em psicoterapia.

 \textbf{Palavras-chave}: aprendizagem de máquina, psicoterapia, psicologia positiva, depressão.
\end{resumo}

% Resumo em inglês
\begin{resumo}[Abstract]
 \begin{otherlanguage*}{english}
   Advances in machine learning have opened up a multitude of opportunities for various research and professional fields. The use of machine learning models to
   predict treatment outcomes in mental health has been investigated showing promising results. The application of such models would allow greater assertiveness
   in tailoring psychotherapeutic treatments at the individual level. Despite its potential benefits, the topic remains unexplored in the fields of research, training
   and practice in clinical psychology. This work aims to introduce the topic of machine learning to psychologists. It presents the basic machine learning concepts and
   terminology in an accessible manner, providing a comprehensive overview of the field. To complement the conceptual discussion, an example is constructed using data
   from a study on the effectiveness of internet-based positive psychology interventions for depression. Finally, implications of using this type of technology in
   psychotherapy research and practice are discussed.
   \vspace{\onelineskip}
 
   \noindent 
   \textbf{Keywords}: machine learning, psychotherapy, positive psychology, depression.
 \end{otherlanguage*}
\end{resumo}


% ---
% LISTA DE ILUSTRAÇÕES
% ---
\pdfbookmark[0]{\listfigurename}{lof}
\listoffigures*
\cleardoublepage
% ---


% ---
% LISTA DE QUADROS
% ---
% \pdfbookmark[0]{\listofquadrosname}{loq}
% \listofquadros*
% \cleardoublepage
% ---


% ---
% LISTA DE TABELAS
% ---
% \pdfbookmark[0]{\listtablename}{lot}
% \listoftables*
% \cleardoublepage
% ---


% ---
% LISTA DE ABREVIATURAS E SIGLAS
% ---
% \begin{siglas}
%   \item[ABNT] Associação Brasileira de Normas Técnicas
%   \item[abnTeX] ABsurdas Normas para TeX
% \end{siglas}
% ---


% ---
% LISTA DE SÍMBOLOS
% ---
% \begin{simbolos}
%   \item[$ \Gamma $] Letra grega Gama
%   \item[$ \Lambda $] Lambda
%   \item[$ \zeta $] Letra grega minúscula zeta
%   \item[$ \in $] Pertence
% \end{simbolos}
% ---


% ---
% Sumário
% ---
\pdfbookmark[0]{\contentsname}{toc}
\tableofcontents*
\cleardoublepage
% ---


% ----------------------------------------------------------
% ELEMENTOS TEXTUAIS
% ----------------------------------------------------------
\textual

\include{./01-introducao/main.tex}

\include{./02-machine-learning/main.tex}

\include{./03-exemplo/main.tex}

\include{./04-conclusao/main.tex}

% ----------------------------------------------------------
% ELEMENTOS PÓS-TEXTUAIS
% ----------------------------------------------------------
\postextual
% ----------------------------------------------------------


% ----------------------------------------------------------
% REFERÊNCIAS
% ----------------------------------------------------------
\bibliography{references.bib}


% ----------------------------------------------------------
% GLOSSÁRIO
% ----------------------------------------------------------
%
% Consulte o manual da classe abntex2 para orientações sobre o glossário.
%
%\glossary

% ----------------------------------------------------------
% APÊNDICES
% ----------------------------------------------------------

% ---
% Inicia os apêndices
% ---
% \begin{apendicesenv}

% % Imprime uma página indicando o início dos apêndices
% \partapendices

% % ----------------------------------------------------------
% \chapter{Quisque libero justo}
% % ----------------------------------------------------------

% % ----------------------------------------------------------
% \chapter{Nullam elementum urna}
% % ----------------------------------------------------------

% \end{apendicesenv}
% ---


% ----------------------------------------------------------
% ANEXOS
% ----------------------------------------------------------

% ---
% Inicia os anexos
% ---
% \begin{anexosenv}

% % Imprime uma página indicando o início dos anexos
% \partanexos

% % ---
% \chapter{Morbi ultrices rutrum lorem}
% % ---

% % ---
% \chapter{Cras non urna sed feugiat cum}
% % ---


% % ---
% \chapter{Fusce facilisis lacinia dui}
% % ---

% \end{anexosenv}

%---------------------------------------------------------------------
% INDICE REMISSIVO
%---------------------------------------------------------------------
% \phantompart
% \printindex
%---------------------------------------------------------------------
\includepdf{./05-extras/contracapa.pdf}
\end{document}


% ----------------------------------------------------------
% ELEMENTOS PÓS-TEXTUAIS
% ----------------------------------------------------------
\postextual
% ----------------------------------------------------------


% ----------------------------------------------------------
% REFERÊNCIAS
% ----------------------------------------------------------
\bibliography{references.bib}


% ----------------------------------------------------------
% GLOSSÁRIO
% ----------------------------------------------------------
%
% Consulte o manual da classe abntex2 para orientações sobre o glossário.
%
%\glossary

% ----------------------------------------------------------
% APÊNDICES
% ----------------------------------------------------------

% ---
% Inicia os apêndices
% ---
% \begin{apendicesenv}

% % Imprime uma página indicando o início dos apêndices
% \partapendices

% % ----------------------------------------------------------
% \chapter{Quisque libero justo}
% % ----------------------------------------------------------

% % ----------------------------------------------------------
% \chapter{Nullam elementum urna}
% % ----------------------------------------------------------

% \end{apendicesenv}
% ---


% ----------------------------------------------------------
% ANEXOS
% ----------------------------------------------------------

% ---
% Inicia os anexos
% ---
% \begin{anexosenv}

% % Imprime uma página indicando o início dos anexos
% \partanexos

% % ---
% \chapter{Morbi ultrices rutrum lorem}
% % ---

% % ---
% \chapter{Cras non urna sed feugiat cum}
% % ---


% % ---
% \chapter{Fusce facilisis lacinia dui}
% % ---

% \end{anexosenv}

%---------------------------------------------------------------------
% INDICE REMISSIVO
%---------------------------------------------------------------------
% \phantompart
% \printindex
%---------------------------------------------------------------------
\includepdf{./05-extras/contracapa.pdf}
\end{document}


% ----------------------------------------------------------
% ELEMENTOS PÓS-TEXTUAIS
% ----------------------------------------------------------
\postextual
% ----------------------------------------------------------


% ----------------------------------------------------------
% REFERÊNCIAS
% ----------------------------------------------------------
\bibliography{references.bib}


% ----------------------------------------------------------
% GLOSSÁRIO
% ----------------------------------------------------------
%
% Consulte o manual da classe abntex2 para orientações sobre o glossário.
%
%\glossary

% ----------------------------------------------------------
% APÊNDICES
% ----------------------------------------------------------

% ---
% Inicia os apêndices
% ---
% \begin{apendicesenv}

% % Imprime uma página indicando o início dos apêndices
% \partapendices

% % ----------------------------------------------------------
% \chapter{Quisque libero justo}
% % ----------------------------------------------------------

% % ----------------------------------------------------------
% \chapter{Nullam elementum urna}
% % ----------------------------------------------------------

% \end{apendicesenv}
% ---


% ----------------------------------------------------------
% ANEXOS
% ----------------------------------------------------------

% ---
% Inicia os anexos
% ---
% \begin{anexosenv}

% % Imprime uma página indicando o início dos anexos
% \partanexos

% % ---
% \chapter{Morbi ultrices rutrum lorem}
% % ---

% % ---
% \chapter{Cras non urna sed feugiat cum}
% % ---


% % ---
% \chapter{Fusce facilisis lacinia dui}
% % ---

% \end{anexosenv}

%---------------------------------------------------------------------
% INDICE REMISSIVO
%---------------------------------------------------------------------
% \phantompart
% \printindex
%---------------------------------------------------------------------
\includepdf{./05-extras/contracapa.pdf}
\end{document}


\documentclass[
	% Opções da classe memoir
	12pt,				        % tamanho da fonte
	openany,			        % capítulos começam em pág ímpar (insere página vazia caso preciso)
	oneside,			        % para impressão em recto e verso. Oposto a oneside
	a4paper,			        % tamanho do papel. 
	% Opções da classe abntex2
	%chapter=TITLE,		        % títulos de capítulos convertidos em letras maiúsculas
	%section=TITLE,		        % títulos de seções convertidos em letras maiúsculas
	%subsection=TITLE,	        % títulos de subseções convertidos em letras maiúsculas
	%subsubsection=TITLE,       % títulos de subsubseções convertidos em letras maiúsculas
	% Opções do pacote babel
	english,			        % idioma adicional para hifenização
	brazil				        % o último idioma é o principal do documento
	]{abntex2}

% ---
% Pacotes básicos 
% ---
\usepackage{lmodern}			% Usa a fonte Latin Modern			
\usepackage[T1]{fontenc}		% Selecao de codigos de fonte.
\usepackage[utf8]{inputenc}		% Codificacao do documento (conversão automática dos acentos)
\usepackage{indentfirst}		% Indenta o primeiro parágrafo de cada seção.
\usepackage{color}				% Controle das cores
\usepackage{graphicx}			% Inclusão de gráficos
\usepackage{microtype} 			% para melhorias de justificação
% ---
		

% ---
% Pacotes de citações
% ---
% \usepackage[alf]{abntex2cite}						% Citações padrão ABNT
\usepackage[
  alf,
  abnt-nbr10520=1988,                     % Citações em sentence case
  abnt-emphasize=bf
  ]{abntex2cite}


% ---
% Pacotes adicionais
% ---
\usepackage{pdfpages}	% Inclusão de páginas em PDF


% ---
% Informações para CAPA e FOLHA DE ROSTO
% ---
\titulo{APRENDIZAGEM DE MÁQUINA PARA PREDIÇÃO DE DESFECHOS EM SAÚDE MENTAL: \textmd{UMA INTRODUÇÃO PARA PSICÓLOGOS}}
\autor{GUILHERME AUGUSTO DIAS}
\local{Porto Alegre}
\data{2024}
\orientador{Prof. Dr. Wagner de Lara Machado}
% \coorientador{}
% \instituicao{
%   Pontifícia Universidade Católica do Rio Grande do Sul
%   \par
%   Escola de Ciências da Saúde e da Vida
%   \par
%   Curso de Psicologia}
\tipotrabalho{TCC}
\preambulo{Trabalho de Conclusão de Curso apresentado como requisito para a obtenção do grau de Bacharel em Psicologia pela Escola de Ciências da Saúde e da Vida da Pontifícia Universidade Católica do Rio Grande do Sul.}
% ---

% ---
% Configurações de aparência do PDF final
% ---

% Alterando o aspecto da cor azul
\definecolor{blue}{RGB}{41,5,195}

% Informações do PDF
\makeatletter
\hypersetup{
		pdftitle={Aprendizagem de máquina para predição de desfechos em saúde mental: uma introdução para psicólogos}, 
		pdfauthor={Guilherme Augusto Dias},
    	pdfsubject={\imprimirpreambulo},
	    pdfcreator={LaTeX with abnTeX2},
		pdfkeywords={aprendizagem de máquina}{psicoterapia}{psicologia positiva}{depressão}, 
		colorlinks=true,		% true: colored links; false: boxed links
    	linkcolor=black,		% color of internal links
    	citecolor=black,		% color of links to bibliography
    	filecolor=black,		% color of file links
		urlcolor=black,
		bookmarksdepth=4
}
\makeatother
% --- 


% ---
% Posiciona figuras e tabelas no topo da página quando adicionadas sozinhas
% em um página em branco. Ver https://github.com/abntex/abntex2/issues/170
\makeatletter
\setlength{\@fptop}{5pt} % Set distance from top of page to first float
\makeatother
% ---

% ---
% Possibilita criação de Quadros e Lista de quadros.
% Ver https://github.com/abntex/abntex2/issues/176
%
\newcommand{\quadroname}{Quadro}
\newcommand{\listofquadrosname}{Lista de quadros}

\newfloat[chapter]{quadro}{loq}{\quadroname}
\newlistof{listofquadros}{loq}{\listofquadrosname}
\newlistentry{quadro}{loq}{0}

% Configurações para atender às regras da ABNT
\setfloatadjustment{quadro}{\centering}
\counterwithout{quadro}{chapter}
\renewcommand{\cftquadroname}{\quadroname\space} 
\renewcommand*{\cftquadroaftersnum}{\hfill--\hfill}

\setfloatlocations{quadro}{hbtp} % Ver https://github.com/abntex/abntex2/issues/176
% ---


% --- 
% Espaçamentos entre linhas e parágrafos 
% --- 

% O tamanho do parágrafo é dado por:
\setlength{\parindent}{1.3cm}

% Controle do espaçamento entre um parágrafo e outro:
\setlength{\parskip}{0.2cm}  % tente também \onelineskip

% ---
% Compila o índice
% ---
\makeindex
% ---


% ----
% Início do documento
% ----
\begin{document}

% Seleciona o idioma do documento (conforme pacotes do babel)
\selectlanguage{brazil}

% Retira espaço extra obsoleto entre as frases.
\frenchspacing 

% ----------------------------------------------------------
% ELEMENTOS PRÉ-TEXTUAIS
% ----------------------------------------------------------
\pretextual

% ---
% CAPA
% ---
\includepdf{./05-extras/capa.pdf}
% ---


% ---
% FOLHA DE ROSTO
% (o * indica que haverá a ficha bibliográfica)
% ---
\imprimirfolhaderosto
% ---


% ---
% Ficha bibliografica
% ---

% Isto é um exemplo de Ficha Catalográfica, ou ``Dados internacionais de
% catalogação-na-publicação''. Você pode utilizar este modelo como referência. 
% Porém, provavelmente a biblioteca da sua universidade lhe fornecerá um PDF
% com a ficha catalográfica definitiva após a defesa do trabalho. Quando estiver
% com o documento, salve-o como PDF no diretório do seu projeto e substitua todo
% o conteúdo de implementação deste arquivo pelo comando abaixo:
%
% \usepackage{pdfpages}		% necessário para comando \includepdf
% \begin{fichacatalografica}
%     \includepdf{fig_ficha_catalografica.pdf}
% \end{fichacatalografica}

% \begin{fichacatalografica}
% 	\sffamily
% 	\vspace*{\fill}                             % Posição vertical
% 	\begin{center}                              % Minipage Centralizado
% 	\fbox{\begin{minipage}[c][8cm]{13.5cm}		% Largura
% 	\small
% 	\imprimirautor
% 	%Sobrenome, Nome do autor
	
% 	\hspace{0.5cm} \imprimirtitulo  / \imprimirautor. --
% 	\imprimirlocal, \imprimirdata-
	
% 	\hspace{0.5cm} \thelastpage p. : il. (algumas color.) ; 30 cm.\\
	
% 	\hspace{0.5cm} \imprimirorientadorRotulo~\imprimirorientador\\
	
% 	\hspace{0.5cm}
% 	\parbox[t]{\textwidth}{\imprimirtipotrabalho~--~\imprimirinstituicao,
% 	\imprimirdata.}\\
	
% 	\hspace{0.5cm}
% 		1. Palavra-chave1.
% 		2. Palavra-chave2.
% 		2. Palavra-chave3.
% 		I. Orientador.
% 		II. Universidade xxx.
% 		III. Faculdade de xxx.
% 		IV. Título 			
% 	\end{minipage}}
% 	\end{center}
% \end{fichacatalografica}
% ---


% ---
% FOLHA DE APROVAÇÃO
% ---

% Isto é um exemplo de Folha de aprovação, elemento obrigatório da NBR
% 14724/2011 (seção 4.2.1.3). Você pode utilizar este modelo até a aprovação
% do trabalho. Após isso, substitua todo o conteúdo deste arquivo por uma
% imagem da página assinada pela banca com o comando abaixo:
%
% \begin{folhadeaprovacao}
% \includepdf{folhadeaprovacao_final.pdf}
% \end{folhadeaprovacao}
%
\setlength{\ABNTEXsignskip}{2cm}
\setlength{\ABNTEXsignwidth}{10cm}

\begin{folhadeaprovacao}
  \begin{center}
    {\ABNTEXchapterfont\large\imprimirautor}

    \vspace*{1cm}

    \begin{center}
      \ABNTEXchapterfont\bfseries\Large\imprimirtitulo
    \end{center}

    \vspace*{1cm}
    
    \hspace{.45\textwidth}
    \begin{minipage}{.5\textwidth}
        \imprimirpreambulo
    \end{minipage}%
   \end{center}

   \vspace*{1cm}
        
   \begin{center}
   Aprovado em 17 de julho de 2024.

   \vspace*{1cm}

   Banca examinadora:
   \end{center}

   \vspace*{-1cm}

   \assinatura{Prof. Ma. Caroline Santa Maria Rodrigues -- PUCRS}
   \assinatura{Prof. Dra. Clarissa Pinto Pizarro de Freitas -- PUCRS}
   \assinatura{\imprimirorientador -- PUCRS}
\end{folhadeaprovacao}
% ---


% ---
% Agradecimentos
% ---
\begin{agradecimentos}
Agradeço aos meus pais, Ademilson José Dias e Irinês Porfírio Dias, pelo apoio incondicional que me deram ao longo da vida,
incentivando mesmo as ideias questionáveis como aquela de fazer uma segunda graduação.

A minha irmã, Maria Eduarda Dias, cuja trajetória acadêmica é fonte de inspiração.

A minha namorada, Nathalia Fallavena Ceratti, que acompanhou a esta jornada de perto, oferecendo apoio e compreensão que
fizeram toda a diferença.

Aos meus amigos, Guilherme da Silva Petró, Norton Neves e Uilinton Luciano Marques da Silva, que foram compreensíveis nos
finais de semana em que não pude jogar RPG.

Ao professor Wagner de Lara Machado pela orientação e apoio no desenvolvimento deste trabalho.
\end{agradecimentos}
% ---


% ---
% RESUMOS
% ---
% Resumo em português
\setlength{\absparsep}{18pt} % ajusta o espaçamento dos parágrafos do resumo
\begin{resumo}
Desenvolvimentos na área de aprendizagem de máquina representam uma série de oportunidades para diversas áreas de pesquisa e de atuação profissional.
O uso de modelos de aprendizagem de máquina na predição de desfecho para tratamentos em saúde mental tem sido investigado e apresenta resultados promissores.
A aplicação desse tipo de modelo possibilitaria, por exemplo, maior assertividade na personalização de tratamentos psicoterápicos a nível individual.
Apesar dos benefícios potenciais, o tema ainda é pouco discutido nos espaços de pesquisa, formação e atuação em psicologia clínica. Este trabalho tem
por objetivo introduzir o tema de aprendizagem de máquina a psicólogos. Busca-se apresentar, de maneira acessível, os conceitos e terminologia básicos
em aprendizagem de máquina, proporcionando um panorama geral da área. Para complementar a exposição conceitual, é construído um exemplo a partir de dados
de um estudo sobre eficácia de intervenções em psicologia positiva para depressão via internet. Por fim, discutem-se algumas implicações do uso desse tipo
de tecnologia para a pesquisa e para a prática em psicoterapia.

 \textbf{Palavras-chave}: aprendizagem de máquina, psicoterapia, psicologia positiva, depressão.
\end{resumo}

% Resumo em inglês
\begin{resumo}[Abstract]
 \begin{otherlanguage*}{english}
   Advances in machine learning have opened up a multitude of opportunities for various research and professional fields. The use of machine learning models to
   predict treatment outcomes in mental health has been investigated showing promising results. The application of such models would allow greater assertiveness
   in tailoring psychotherapeutic treatments at the individual level. Despite its potential benefits, the topic remains unexplored in the fields of research, training
   and practice in clinical psychology. This work aims to introduce the topic of machine learning to psychologists. It presents the basic machine learning concepts and
   terminology in an accessible manner, providing a comprehensive overview of the field. To complement the conceptual discussion, an example is constructed using data
   from a study on the effectiveness of internet-based positive psychology interventions for depression. Finally, implications of using this type of technology in
   psychotherapy research and practice are discussed.
   \vspace{\onelineskip}
 
   \noindent 
   \textbf{Keywords}: machine learning, psychotherapy, positive psychology, depression.
 \end{otherlanguage*}
\end{resumo}


% ---
% LISTA DE ILUSTRAÇÕES
% ---
\pdfbookmark[0]{\listfigurename}{lof}
\listoffigures*
\cleardoublepage
% ---


% ---
% LISTA DE QUADROS
% ---
% \pdfbookmark[0]{\listofquadrosname}{loq}
% \listofquadros*
% \cleardoublepage
% ---


% ---
% LISTA DE TABELAS
% ---
% \pdfbookmark[0]{\listtablename}{lot}
% \listoftables*
% \cleardoublepage
% ---


% ---
% LISTA DE ABREVIATURAS E SIGLAS
% ---
% \begin{siglas}
%   \item[ABNT] Associação Brasileira de Normas Técnicas
%   \item[abnTeX] ABsurdas Normas para TeX
% \end{siglas}
% ---


% ---
% LISTA DE SÍMBOLOS
% ---
% \begin{simbolos}
%   \item[$ \Gamma $] Letra grega Gama
%   \item[$ \Lambda $] Lambda
%   \item[$ \zeta $] Letra grega minúscula zeta
%   \item[$ \in $] Pertence
% \end{simbolos}
% ---


% ---
% Sumário
% ---
\pdfbookmark[0]{\contentsname}{toc}
\tableofcontents*
\cleardoublepage
% ---


% ----------------------------------------------------------
% ELEMENTOS TEXTUAIS
% ----------------------------------------------------------
\textual

\documentclass[
	% Opções da classe memoir
	12pt,				        % tamanho da fonte
	openany,			        % capítulos começam em pág ímpar (insere página vazia caso preciso)
	oneside,			        % para impressão em recto e verso. Oposto a oneside
	a4paper,			        % tamanho do papel. 
	% Opções da classe abntex2
	%chapter=TITLE,		        % títulos de capítulos convertidos em letras maiúsculas
	%section=TITLE,		        % títulos de seções convertidos em letras maiúsculas
	%subsection=TITLE,	        % títulos de subseções convertidos em letras maiúsculas
	%subsubsection=TITLE,       % títulos de subsubseções convertidos em letras maiúsculas
	% Opções do pacote babel
	english,			        % idioma adicional para hifenização
	brazil				        % o último idioma é o principal do documento
	]{abntex2}

% ---
% Pacotes básicos 
% ---
\usepackage{lmodern}			% Usa a fonte Latin Modern			
\usepackage[T1]{fontenc}		% Selecao de codigos de fonte.
\usepackage[utf8]{inputenc}		% Codificacao do documento (conversão automática dos acentos)
\usepackage{indentfirst}		% Indenta o primeiro parágrafo de cada seção.
\usepackage{color}				% Controle das cores
\usepackage{graphicx}			% Inclusão de gráficos
\usepackage{microtype} 			% para melhorias de justificação
% ---
		

% ---
% Pacotes de citações
% ---
% \usepackage[alf]{abntex2cite}						% Citações padrão ABNT
\usepackage[
  alf,
  abnt-nbr10520=1988,                     % Citações em sentence case
  abnt-emphasize=bf
  ]{abntex2cite}


% ---
% Pacotes adicionais
% ---
\usepackage{pdfpages}	% Inclusão de páginas em PDF


% ---
% Informações para CAPA e FOLHA DE ROSTO
% ---
\titulo{APRENDIZAGEM DE MÁQUINA PARA PREDIÇÃO DE DESFECHOS EM SAÚDE MENTAL: \textmd{UMA INTRODUÇÃO PARA PSICÓLOGOS}}
\autor{GUILHERME AUGUSTO DIAS}
\local{Porto Alegre}
\data{2024}
\orientador{Prof. Dr. Wagner de Lara Machado}
% \coorientador{}
% \instituicao{
%   Pontifícia Universidade Católica do Rio Grande do Sul
%   \par
%   Escola de Ciências da Saúde e da Vida
%   \par
%   Curso de Psicologia}
\tipotrabalho{TCC}
\preambulo{Trabalho de Conclusão de Curso apresentado como requisito para a obtenção do grau de Bacharel em Psicologia pela Escola de Ciências da Saúde e da Vida da Pontifícia Universidade Católica do Rio Grande do Sul.}
% ---

% ---
% Configurações de aparência do PDF final
% ---

% Alterando o aspecto da cor azul
\definecolor{blue}{RGB}{41,5,195}

% Informações do PDF
\makeatletter
\hypersetup{
		pdftitle={Aprendizagem de máquina para predição de desfechos em saúde mental: uma introdução para psicólogos}, 
		pdfauthor={Guilherme Augusto Dias},
    	pdfsubject={\imprimirpreambulo},
	    pdfcreator={LaTeX with abnTeX2},
		pdfkeywords={aprendizagem de máquina}{psicoterapia}{psicologia positiva}{depressão}, 
		colorlinks=true,		% true: colored links; false: boxed links
    	linkcolor=black,		% color of internal links
    	citecolor=black,		% color of links to bibliography
    	filecolor=black,		% color of file links
		urlcolor=black,
		bookmarksdepth=4
}
\makeatother
% --- 


% ---
% Posiciona figuras e tabelas no topo da página quando adicionadas sozinhas
% em um página em branco. Ver https://github.com/abntex/abntex2/issues/170
\makeatletter
\setlength{\@fptop}{5pt} % Set distance from top of page to first float
\makeatother
% ---

% ---
% Possibilita criação de Quadros e Lista de quadros.
% Ver https://github.com/abntex/abntex2/issues/176
%
\newcommand{\quadroname}{Quadro}
\newcommand{\listofquadrosname}{Lista de quadros}

\newfloat[chapter]{quadro}{loq}{\quadroname}
\newlistof{listofquadros}{loq}{\listofquadrosname}
\newlistentry{quadro}{loq}{0}

% Configurações para atender às regras da ABNT
\setfloatadjustment{quadro}{\centering}
\counterwithout{quadro}{chapter}
\renewcommand{\cftquadroname}{\quadroname\space} 
\renewcommand*{\cftquadroaftersnum}{\hfill--\hfill}

\setfloatlocations{quadro}{hbtp} % Ver https://github.com/abntex/abntex2/issues/176
% ---


% --- 
% Espaçamentos entre linhas e parágrafos 
% --- 

% O tamanho do parágrafo é dado por:
\setlength{\parindent}{1.3cm}

% Controle do espaçamento entre um parágrafo e outro:
\setlength{\parskip}{0.2cm}  % tente também \onelineskip

% ---
% Compila o índice
% ---
\makeindex
% ---


% ----
% Início do documento
% ----
\begin{document}

% Seleciona o idioma do documento (conforme pacotes do babel)
\selectlanguage{brazil}

% Retira espaço extra obsoleto entre as frases.
\frenchspacing 

% ----------------------------------------------------------
% ELEMENTOS PRÉ-TEXTUAIS
% ----------------------------------------------------------
\pretextual

% ---
% CAPA
% ---
\includepdf{./05-extras/capa.pdf}
% ---


% ---
% FOLHA DE ROSTO
% (o * indica que haverá a ficha bibliográfica)
% ---
\imprimirfolhaderosto
% ---


% ---
% Ficha bibliografica
% ---

% Isto é um exemplo de Ficha Catalográfica, ou ``Dados internacionais de
% catalogação-na-publicação''. Você pode utilizar este modelo como referência. 
% Porém, provavelmente a biblioteca da sua universidade lhe fornecerá um PDF
% com a ficha catalográfica definitiva após a defesa do trabalho. Quando estiver
% com o documento, salve-o como PDF no diretório do seu projeto e substitua todo
% o conteúdo de implementação deste arquivo pelo comando abaixo:
%
% \usepackage{pdfpages}		% necessário para comando \includepdf
% \begin{fichacatalografica}
%     \includepdf{fig_ficha_catalografica.pdf}
% \end{fichacatalografica}

% \begin{fichacatalografica}
% 	\sffamily
% 	\vspace*{\fill}                             % Posição vertical
% 	\begin{center}                              % Minipage Centralizado
% 	\fbox{\begin{minipage}[c][8cm]{13.5cm}		% Largura
% 	\small
% 	\imprimirautor
% 	%Sobrenome, Nome do autor
	
% 	\hspace{0.5cm} \imprimirtitulo  / \imprimirautor. --
% 	\imprimirlocal, \imprimirdata-
	
% 	\hspace{0.5cm} \thelastpage p. : il. (algumas color.) ; 30 cm.\\
	
% 	\hspace{0.5cm} \imprimirorientadorRotulo~\imprimirorientador\\
	
% 	\hspace{0.5cm}
% 	\parbox[t]{\textwidth}{\imprimirtipotrabalho~--~\imprimirinstituicao,
% 	\imprimirdata.}\\
	
% 	\hspace{0.5cm}
% 		1. Palavra-chave1.
% 		2. Palavra-chave2.
% 		2. Palavra-chave3.
% 		I. Orientador.
% 		II. Universidade xxx.
% 		III. Faculdade de xxx.
% 		IV. Título 			
% 	\end{minipage}}
% 	\end{center}
% \end{fichacatalografica}
% ---


% ---
% FOLHA DE APROVAÇÃO
% ---

% Isto é um exemplo de Folha de aprovação, elemento obrigatório da NBR
% 14724/2011 (seção 4.2.1.3). Você pode utilizar este modelo até a aprovação
% do trabalho. Após isso, substitua todo o conteúdo deste arquivo por uma
% imagem da página assinada pela banca com o comando abaixo:
%
% \begin{folhadeaprovacao}
% \includepdf{folhadeaprovacao_final.pdf}
% \end{folhadeaprovacao}
%
\setlength{\ABNTEXsignskip}{2cm}
\setlength{\ABNTEXsignwidth}{10cm}

\begin{folhadeaprovacao}
  \begin{center}
    {\ABNTEXchapterfont\large\imprimirautor}

    \vspace*{1cm}

    \begin{center}
      \ABNTEXchapterfont\bfseries\Large\imprimirtitulo
    \end{center}

    \vspace*{1cm}
    
    \hspace{.45\textwidth}
    \begin{minipage}{.5\textwidth}
        \imprimirpreambulo
    \end{minipage}%
   \end{center}

   \vspace*{1cm}
        
   \begin{center}
   Aprovado em 17 de julho de 2024.

   \vspace*{1cm}

   Banca examinadora:
   \end{center}

   \vspace*{-1cm}

   \assinatura{Prof. Ma. Caroline Santa Maria Rodrigues -- PUCRS}
   \assinatura{Prof. Dra. Clarissa Pinto Pizarro de Freitas -- PUCRS}
   \assinatura{\imprimirorientador -- PUCRS}
\end{folhadeaprovacao}
% ---


% ---
% Agradecimentos
% ---
\begin{agradecimentos}
Agradeço aos meus pais, Ademilson José Dias e Irinês Porfírio Dias, pelo apoio incondicional que me deram ao longo da vida,
incentivando mesmo as ideias questionáveis como aquela de fazer uma segunda graduação.

A minha irmã, Maria Eduarda Dias, cuja trajetória acadêmica é fonte de inspiração.

A minha namorada, Nathalia Fallavena Ceratti, que acompanhou a esta jornada de perto, oferecendo apoio e compreensão que
fizeram toda a diferença.

Aos meus amigos, Guilherme da Silva Petró, Norton Neves e Uilinton Luciano Marques da Silva, que foram compreensíveis nos
finais de semana em que não pude jogar RPG.

Ao professor Wagner de Lara Machado pela orientação e apoio no desenvolvimento deste trabalho.
\end{agradecimentos}
% ---


% ---
% RESUMOS
% ---
% Resumo em português
\setlength{\absparsep}{18pt} % ajusta o espaçamento dos parágrafos do resumo
\begin{resumo}
Desenvolvimentos na área de aprendizagem de máquina representam uma série de oportunidades para diversas áreas de pesquisa e de atuação profissional.
O uso de modelos de aprendizagem de máquina na predição de desfecho para tratamentos em saúde mental tem sido investigado e apresenta resultados promissores.
A aplicação desse tipo de modelo possibilitaria, por exemplo, maior assertividade na personalização de tratamentos psicoterápicos a nível individual.
Apesar dos benefícios potenciais, o tema ainda é pouco discutido nos espaços de pesquisa, formação e atuação em psicologia clínica. Este trabalho tem
por objetivo introduzir o tema de aprendizagem de máquina a psicólogos. Busca-se apresentar, de maneira acessível, os conceitos e terminologia básicos
em aprendizagem de máquina, proporcionando um panorama geral da área. Para complementar a exposição conceitual, é construído um exemplo a partir de dados
de um estudo sobre eficácia de intervenções em psicologia positiva para depressão via internet. Por fim, discutem-se algumas implicações do uso desse tipo
de tecnologia para a pesquisa e para a prática em psicoterapia.

 \textbf{Palavras-chave}: aprendizagem de máquina, psicoterapia, psicologia positiva, depressão.
\end{resumo}

% Resumo em inglês
\begin{resumo}[Abstract]
 \begin{otherlanguage*}{english}
   Advances in machine learning have opened up a multitude of opportunities for various research and professional fields. The use of machine learning models to
   predict treatment outcomes in mental health has been investigated showing promising results. The application of such models would allow greater assertiveness
   in tailoring psychotherapeutic treatments at the individual level. Despite its potential benefits, the topic remains unexplored in the fields of research, training
   and practice in clinical psychology. This work aims to introduce the topic of machine learning to psychologists. It presents the basic machine learning concepts and
   terminology in an accessible manner, providing a comprehensive overview of the field. To complement the conceptual discussion, an example is constructed using data
   from a study on the effectiveness of internet-based positive psychology interventions for depression. Finally, implications of using this type of technology in
   psychotherapy research and practice are discussed.
   \vspace{\onelineskip}
 
   \noindent 
   \textbf{Keywords}: machine learning, psychotherapy, positive psychology, depression.
 \end{otherlanguage*}
\end{resumo}


% ---
% LISTA DE ILUSTRAÇÕES
% ---
\pdfbookmark[0]{\listfigurename}{lof}
\listoffigures*
\cleardoublepage
% ---


% ---
% LISTA DE QUADROS
% ---
% \pdfbookmark[0]{\listofquadrosname}{loq}
% \listofquadros*
% \cleardoublepage
% ---


% ---
% LISTA DE TABELAS
% ---
% \pdfbookmark[0]{\listtablename}{lot}
% \listoftables*
% \cleardoublepage
% ---


% ---
% LISTA DE ABREVIATURAS E SIGLAS
% ---
% \begin{siglas}
%   \item[ABNT] Associação Brasileira de Normas Técnicas
%   \item[abnTeX] ABsurdas Normas para TeX
% \end{siglas}
% ---


% ---
% LISTA DE SÍMBOLOS
% ---
% \begin{simbolos}
%   \item[$ \Gamma $] Letra grega Gama
%   \item[$ \Lambda $] Lambda
%   \item[$ \zeta $] Letra grega minúscula zeta
%   \item[$ \in $] Pertence
% \end{simbolos}
% ---


% ---
% Sumário
% ---
\pdfbookmark[0]{\contentsname}{toc}
\tableofcontents*
\cleardoublepage
% ---


% ----------------------------------------------------------
% ELEMENTOS TEXTUAIS
% ----------------------------------------------------------
\textual

\documentclass[
	% Opções da classe memoir
	12pt,				        % tamanho da fonte
	openany,			        % capítulos começam em pág ímpar (insere página vazia caso preciso)
	oneside,			        % para impressão em recto e verso. Oposto a oneside
	a4paper,			        % tamanho do papel. 
	% Opções da classe abntex2
	%chapter=TITLE,		        % títulos de capítulos convertidos em letras maiúsculas
	%section=TITLE,		        % títulos de seções convertidos em letras maiúsculas
	%subsection=TITLE,	        % títulos de subseções convertidos em letras maiúsculas
	%subsubsection=TITLE,       % títulos de subsubseções convertidos em letras maiúsculas
	% Opções do pacote babel
	english,			        % idioma adicional para hifenização
	brazil				        % o último idioma é o principal do documento
	]{abntex2}

% ---
% Pacotes básicos 
% ---
\usepackage{lmodern}			% Usa a fonte Latin Modern			
\usepackage[T1]{fontenc}		% Selecao de codigos de fonte.
\usepackage[utf8]{inputenc}		% Codificacao do documento (conversão automática dos acentos)
\usepackage{indentfirst}		% Indenta o primeiro parágrafo de cada seção.
\usepackage{color}				% Controle das cores
\usepackage{graphicx}			% Inclusão de gráficos
\usepackage{microtype} 			% para melhorias de justificação
% ---
		

% ---
% Pacotes de citações
% ---
% \usepackage[alf]{abntex2cite}						% Citações padrão ABNT
\usepackage[
  alf,
  abnt-nbr10520=1988,                     % Citações em sentence case
  abnt-emphasize=bf
  ]{abntex2cite}


% ---
% Pacotes adicionais
% ---
\usepackage{pdfpages}	% Inclusão de páginas em PDF


% ---
% Informações para CAPA e FOLHA DE ROSTO
% ---
\titulo{APRENDIZAGEM DE MÁQUINA PARA PREDIÇÃO DE DESFECHOS EM SAÚDE MENTAL: \textmd{UMA INTRODUÇÃO PARA PSICÓLOGOS}}
\autor{GUILHERME AUGUSTO DIAS}
\local{Porto Alegre}
\data{2024}
\orientador{Prof. Dr. Wagner de Lara Machado}
% \coorientador{}
% \instituicao{
%   Pontifícia Universidade Católica do Rio Grande do Sul
%   \par
%   Escola de Ciências da Saúde e da Vida
%   \par
%   Curso de Psicologia}
\tipotrabalho{TCC}
\preambulo{Trabalho de Conclusão de Curso apresentado como requisito para a obtenção do grau de Bacharel em Psicologia pela Escola de Ciências da Saúde e da Vida da Pontifícia Universidade Católica do Rio Grande do Sul.}
% ---

% ---
% Configurações de aparência do PDF final
% ---

% Alterando o aspecto da cor azul
\definecolor{blue}{RGB}{41,5,195}

% Informações do PDF
\makeatletter
\hypersetup{
		pdftitle={Aprendizagem de máquina para predição de desfechos em saúde mental: uma introdução para psicólogos}, 
		pdfauthor={Guilherme Augusto Dias},
    	pdfsubject={\imprimirpreambulo},
	    pdfcreator={LaTeX with abnTeX2},
		pdfkeywords={aprendizagem de máquina}{psicoterapia}{psicologia positiva}{depressão}, 
		colorlinks=true,		% true: colored links; false: boxed links
    	linkcolor=black,		% color of internal links
    	citecolor=black,		% color of links to bibliography
    	filecolor=black,		% color of file links
		urlcolor=black,
		bookmarksdepth=4
}
\makeatother
% --- 


% ---
% Posiciona figuras e tabelas no topo da página quando adicionadas sozinhas
% em um página em branco. Ver https://github.com/abntex/abntex2/issues/170
\makeatletter
\setlength{\@fptop}{5pt} % Set distance from top of page to first float
\makeatother
% ---

% ---
% Possibilita criação de Quadros e Lista de quadros.
% Ver https://github.com/abntex/abntex2/issues/176
%
\newcommand{\quadroname}{Quadro}
\newcommand{\listofquadrosname}{Lista de quadros}

\newfloat[chapter]{quadro}{loq}{\quadroname}
\newlistof{listofquadros}{loq}{\listofquadrosname}
\newlistentry{quadro}{loq}{0}

% Configurações para atender às regras da ABNT
\setfloatadjustment{quadro}{\centering}
\counterwithout{quadro}{chapter}
\renewcommand{\cftquadroname}{\quadroname\space} 
\renewcommand*{\cftquadroaftersnum}{\hfill--\hfill}

\setfloatlocations{quadro}{hbtp} % Ver https://github.com/abntex/abntex2/issues/176
% ---


% --- 
% Espaçamentos entre linhas e parágrafos 
% --- 

% O tamanho do parágrafo é dado por:
\setlength{\parindent}{1.3cm}

% Controle do espaçamento entre um parágrafo e outro:
\setlength{\parskip}{0.2cm}  % tente também \onelineskip

% ---
% Compila o índice
% ---
\makeindex
% ---


% ----
% Início do documento
% ----
\begin{document}

% Seleciona o idioma do documento (conforme pacotes do babel)
\selectlanguage{brazil}

% Retira espaço extra obsoleto entre as frases.
\frenchspacing 

% ----------------------------------------------------------
% ELEMENTOS PRÉ-TEXTUAIS
% ----------------------------------------------------------
\pretextual

% ---
% CAPA
% ---
\includepdf{./05-extras/capa.pdf}
% ---


% ---
% FOLHA DE ROSTO
% (o * indica que haverá a ficha bibliográfica)
% ---
\imprimirfolhaderosto
% ---


% ---
% Ficha bibliografica
% ---

% Isto é um exemplo de Ficha Catalográfica, ou ``Dados internacionais de
% catalogação-na-publicação''. Você pode utilizar este modelo como referência. 
% Porém, provavelmente a biblioteca da sua universidade lhe fornecerá um PDF
% com a ficha catalográfica definitiva após a defesa do trabalho. Quando estiver
% com o documento, salve-o como PDF no diretório do seu projeto e substitua todo
% o conteúdo de implementação deste arquivo pelo comando abaixo:
%
% \usepackage{pdfpages}		% necessário para comando \includepdf
% \begin{fichacatalografica}
%     \includepdf{fig_ficha_catalografica.pdf}
% \end{fichacatalografica}

% \begin{fichacatalografica}
% 	\sffamily
% 	\vspace*{\fill}                             % Posição vertical
% 	\begin{center}                              % Minipage Centralizado
% 	\fbox{\begin{minipage}[c][8cm]{13.5cm}		% Largura
% 	\small
% 	\imprimirautor
% 	%Sobrenome, Nome do autor
	
% 	\hspace{0.5cm} \imprimirtitulo  / \imprimirautor. --
% 	\imprimirlocal, \imprimirdata-
	
% 	\hspace{0.5cm} \thelastpage p. : il. (algumas color.) ; 30 cm.\\
	
% 	\hspace{0.5cm} \imprimirorientadorRotulo~\imprimirorientador\\
	
% 	\hspace{0.5cm}
% 	\parbox[t]{\textwidth}{\imprimirtipotrabalho~--~\imprimirinstituicao,
% 	\imprimirdata.}\\
	
% 	\hspace{0.5cm}
% 		1. Palavra-chave1.
% 		2. Palavra-chave2.
% 		2. Palavra-chave3.
% 		I. Orientador.
% 		II. Universidade xxx.
% 		III. Faculdade de xxx.
% 		IV. Título 			
% 	\end{minipage}}
% 	\end{center}
% \end{fichacatalografica}
% ---


% ---
% FOLHA DE APROVAÇÃO
% ---

% Isto é um exemplo de Folha de aprovação, elemento obrigatório da NBR
% 14724/2011 (seção 4.2.1.3). Você pode utilizar este modelo até a aprovação
% do trabalho. Após isso, substitua todo o conteúdo deste arquivo por uma
% imagem da página assinada pela banca com o comando abaixo:
%
% \begin{folhadeaprovacao}
% \includepdf{folhadeaprovacao_final.pdf}
% \end{folhadeaprovacao}
%
\setlength{\ABNTEXsignskip}{2cm}
\setlength{\ABNTEXsignwidth}{10cm}

\begin{folhadeaprovacao}
  \begin{center}
    {\ABNTEXchapterfont\large\imprimirautor}

    \vspace*{1cm}

    \begin{center}
      \ABNTEXchapterfont\bfseries\Large\imprimirtitulo
    \end{center}

    \vspace*{1cm}
    
    \hspace{.45\textwidth}
    \begin{minipage}{.5\textwidth}
        \imprimirpreambulo
    \end{minipage}%
   \end{center}

   \vspace*{1cm}
        
   \begin{center}
   Aprovado em 17 de julho de 2024.

   \vspace*{1cm}

   Banca examinadora:
   \end{center}

   \vspace*{-1cm}

   \assinatura{Prof. Ma. Caroline Santa Maria Rodrigues -- PUCRS}
   \assinatura{Prof. Dra. Clarissa Pinto Pizarro de Freitas -- PUCRS}
   \assinatura{\imprimirorientador -- PUCRS}
\end{folhadeaprovacao}
% ---


% ---
% Agradecimentos
% ---
\begin{agradecimentos}
Agradeço aos meus pais, Ademilson José Dias e Irinês Porfírio Dias, pelo apoio incondicional que me deram ao longo da vida,
incentivando mesmo as ideias questionáveis como aquela de fazer uma segunda graduação.

A minha irmã, Maria Eduarda Dias, cuja trajetória acadêmica é fonte de inspiração.

A minha namorada, Nathalia Fallavena Ceratti, que acompanhou a esta jornada de perto, oferecendo apoio e compreensão que
fizeram toda a diferença.

Aos meus amigos, Guilherme da Silva Petró, Norton Neves e Uilinton Luciano Marques da Silva, que foram compreensíveis nos
finais de semana em que não pude jogar RPG.

Ao professor Wagner de Lara Machado pela orientação e apoio no desenvolvimento deste trabalho.
\end{agradecimentos}
% ---


% ---
% RESUMOS
% ---
% Resumo em português
\setlength{\absparsep}{18pt} % ajusta o espaçamento dos parágrafos do resumo
\begin{resumo}
Desenvolvimentos na área de aprendizagem de máquina representam uma série de oportunidades para diversas áreas de pesquisa e de atuação profissional.
O uso de modelos de aprendizagem de máquina na predição de desfecho para tratamentos em saúde mental tem sido investigado e apresenta resultados promissores.
A aplicação desse tipo de modelo possibilitaria, por exemplo, maior assertividade na personalização de tratamentos psicoterápicos a nível individual.
Apesar dos benefícios potenciais, o tema ainda é pouco discutido nos espaços de pesquisa, formação e atuação em psicologia clínica. Este trabalho tem
por objetivo introduzir o tema de aprendizagem de máquina a psicólogos. Busca-se apresentar, de maneira acessível, os conceitos e terminologia básicos
em aprendizagem de máquina, proporcionando um panorama geral da área. Para complementar a exposição conceitual, é construído um exemplo a partir de dados
de um estudo sobre eficácia de intervenções em psicologia positiva para depressão via internet. Por fim, discutem-se algumas implicações do uso desse tipo
de tecnologia para a pesquisa e para a prática em psicoterapia.

 \textbf{Palavras-chave}: aprendizagem de máquina, psicoterapia, psicologia positiva, depressão.
\end{resumo}

% Resumo em inglês
\begin{resumo}[Abstract]
 \begin{otherlanguage*}{english}
   Advances in machine learning have opened up a multitude of opportunities for various research and professional fields. The use of machine learning models to
   predict treatment outcomes in mental health has been investigated showing promising results. The application of such models would allow greater assertiveness
   in tailoring psychotherapeutic treatments at the individual level. Despite its potential benefits, the topic remains unexplored in the fields of research, training
   and practice in clinical psychology. This work aims to introduce the topic of machine learning to psychologists. It presents the basic machine learning concepts and
   terminology in an accessible manner, providing a comprehensive overview of the field. To complement the conceptual discussion, an example is constructed using data
   from a study on the effectiveness of internet-based positive psychology interventions for depression. Finally, implications of using this type of technology in
   psychotherapy research and practice are discussed.
   \vspace{\onelineskip}
 
   \noindent 
   \textbf{Keywords}: machine learning, psychotherapy, positive psychology, depression.
 \end{otherlanguage*}
\end{resumo}


% ---
% LISTA DE ILUSTRAÇÕES
% ---
\pdfbookmark[0]{\listfigurename}{lof}
\listoffigures*
\cleardoublepage
% ---


% ---
% LISTA DE QUADROS
% ---
% \pdfbookmark[0]{\listofquadrosname}{loq}
% \listofquadros*
% \cleardoublepage
% ---


% ---
% LISTA DE TABELAS
% ---
% \pdfbookmark[0]{\listtablename}{lot}
% \listoftables*
% \cleardoublepage
% ---


% ---
% LISTA DE ABREVIATURAS E SIGLAS
% ---
% \begin{siglas}
%   \item[ABNT] Associação Brasileira de Normas Técnicas
%   \item[abnTeX] ABsurdas Normas para TeX
% \end{siglas}
% ---


% ---
% LISTA DE SÍMBOLOS
% ---
% \begin{simbolos}
%   \item[$ \Gamma $] Letra grega Gama
%   \item[$ \Lambda $] Lambda
%   \item[$ \zeta $] Letra grega minúscula zeta
%   \item[$ \in $] Pertence
% \end{simbolos}
% ---


% ---
% Sumário
% ---
\pdfbookmark[0]{\contentsname}{toc}
\tableofcontents*
\cleardoublepage
% ---


% ----------------------------------------------------------
% ELEMENTOS TEXTUAIS
% ----------------------------------------------------------
\textual

\include{./01-introducao/main.tex}

\include{./02-machine-learning/main.tex}

\include{./03-exemplo/main.tex}

\include{./04-conclusao/main.tex}

% ----------------------------------------------------------
% ELEMENTOS PÓS-TEXTUAIS
% ----------------------------------------------------------
\postextual
% ----------------------------------------------------------


% ----------------------------------------------------------
% REFERÊNCIAS
% ----------------------------------------------------------
\bibliography{references.bib}


% ----------------------------------------------------------
% GLOSSÁRIO
% ----------------------------------------------------------
%
% Consulte o manual da classe abntex2 para orientações sobre o glossário.
%
%\glossary

% ----------------------------------------------------------
% APÊNDICES
% ----------------------------------------------------------

% ---
% Inicia os apêndices
% ---
% \begin{apendicesenv}

% % Imprime uma página indicando o início dos apêndices
% \partapendices

% % ----------------------------------------------------------
% \chapter{Quisque libero justo}
% % ----------------------------------------------------------

% % ----------------------------------------------------------
% \chapter{Nullam elementum urna}
% % ----------------------------------------------------------

% \end{apendicesenv}
% ---


% ----------------------------------------------------------
% ANEXOS
% ----------------------------------------------------------

% ---
% Inicia os anexos
% ---
% \begin{anexosenv}

% % Imprime uma página indicando o início dos anexos
% \partanexos

% % ---
% \chapter{Morbi ultrices rutrum lorem}
% % ---

% % ---
% \chapter{Cras non urna sed feugiat cum}
% % ---


% % ---
% \chapter{Fusce facilisis lacinia dui}
% % ---

% \end{anexosenv}

%---------------------------------------------------------------------
% INDICE REMISSIVO
%---------------------------------------------------------------------
% \phantompart
% \printindex
%---------------------------------------------------------------------
\includepdf{./05-extras/contracapa.pdf}
\end{document}


\documentclass[
	% Opções da classe memoir
	12pt,				        % tamanho da fonte
	openany,			        % capítulos começam em pág ímpar (insere página vazia caso preciso)
	oneside,			        % para impressão em recto e verso. Oposto a oneside
	a4paper,			        % tamanho do papel. 
	% Opções da classe abntex2
	%chapter=TITLE,		        % títulos de capítulos convertidos em letras maiúsculas
	%section=TITLE,		        % títulos de seções convertidos em letras maiúsculas
	%subsection=TITLE,	        % títulos de subseções convertidos em letras maiúsculas
	%subsubsection=TITLE,       % títulos de subsubseções convertidos em letras maiúsculas
	% Opções do pacote babel
	english,			        % idioma adicional para hifenização
	brazil				        % o último idioma é o principal do documento
	]{abntex2}

% ---
% Pacotes básicos 
% ---
\usepackage{lmodern}			% Usa a fonte Latin Modern			
\usepackage[T1]{fontenc}		% Selecao de codigos de fonte.
\usepackage[utf8]{inputenc}		% Codificacao do documento (conversão automática dos acentos)
\usepackage{indentfirst}		% Indenta o primeiro parágrafo de cada seção.
\usepackage{color}				% Controle das cores
\usepackage{graphicx}			% Inclusão de gráficos
\usepackage{microtype} 			% para melhorias de justificação
% ---
		

% ---
% Pacotes de citações
% ---
% \usepackage[alf]{abntex2cite}						% Citações padrão ABNT
\usepackage[
  alf,
  abnt-nbr10520=1988,                     % Citações em sentence case
  abnt-emphasize=bf
  ]{abntex2cite}


% ---
% Pacotes adicionais
% ---
\usepackage{pdfpages}	% Inclusão de páginas em PDF


% ---
% Informações para CAPA e FOLHA DE ROSTO
% ---
\titulo{APRENDIZAGEM DE MÁQUINA PARA PREDIÇÃO DE DESFECHOS EM SAÚDE MENTAL: \textmd{UMA INTRODUÇÃO PARA PSICÓLOGOS}}
\autor{GUILHERME AUGUSTO DIAS}
\local{Porto Alegre}
\data{2024}
\orientador{Prof. Dr. Wagner de Lara Machado}
% \coorientador{}
% \instituicao{
%   Pontifícia Universidade Católica do Rio Grande do Sul
%   \par
%   Escola de Ciências da Saúde e da Vida
%   \par
%   Curso de Psicologia}
\tipotrabalho{TCC}
\preambulo{Trabalho de Conclusão de Curso apresentado como requisito para a obtenção do grau de Bacharel em Psicologia pela Escola de Ciências da Saúde e da Vida da Pontifícia Universidade Católica do Rio Grande do Sul.}
% ---

% ---
% Configurações de aparência do PDF final
% ---

% Alterando o aspecto da cor azul
\definecolor{blue}{RGB}{41,5,195}

% Informações do PDF
\makeatletter
\hypersetup{
		pdftitle={Aprendizagem de máquina para predição de desfechos em saúde mental: uma introdução para psicólogos}, 
		pdfauthor={Guilherme Augusto Dias},
    	pdfsubject={\imprimirpreambulo},
	    pdfcreator={LaTeX with abnTeX2},
		pdfkeywords={aprendizagem de máquina}{psicoterapia}{psicologia positiva}{depressão}, 
		colorlinks=true,		% true: colored links; false: boxed links
    	linkcolor=black,		% color of internal links
    	citecolor=black,		% color of links to bibliography
    	filecolor=black,		% color of file links
		urlcolor=black,
		bookmarksdepth=4
}
\makeatother
% --- 


% ---
% Posiciona figuras e tabelas no topo da página quando adicionadas sozinhas
% em um página em branco. Ver https://github.com/abntex/abntex2/issues/170
\makeatletter
\setlength{\@fptop}{5pt} % Set distance from top of page to first float
\makeatother
% ---

% ---
% Possibilita criação de Quadros e Lista de quadros.
% Ver https://github.com/abntex/abntex2/issues/176
%
\newcommand{\quadroname}{Quadro}
\newcommand{\listofquadrosname}{Lista de quadros}

\newfloat[chapter]{quadro}{loq}{\quadroname}
\newlistof{listofquadros}{loq}{\listofquadrosname}
\newlistentry{quadro}{loq}{0}

% Configurações para atender às regras da ABNT
\setfloatadjustment{quadro}{\centering}
\counterwithout{quadro}{chapter}
\renewcommand{\cftquadroname}{\quadroname\space} 
\renewcommand*{\cftquadroaftersnum}{\hfill--\hfill}

\setfloatlocations{quadro}{hbtp} % Ver https://github.com/abntex/abntex2/issues/176
% ---


% --- 
% Espaçamentos entre linhas e parágrafos 
% --- 

% O tamanho do parágrafo é dado por:
\setlength{\parindent}{1.3cm}

% Controle do espaçamento entre um parágrafo e outro:
\setlength{\parskip}{0.2cm}  % tente também \onelineskip

% ---
% Compila o índice
% ---
\makeindex
% ---


% ----
% Início do documento
% ----
\begin{document}

% Seleciona o idioma do documento (conforme pacotes do babel)
\selectlanguage{brazil}

% Retira espaço extra obsoleto entre as frases.
\frenchspacing 

% ----------------------------------------------------------
% ELEMENTOS PRÉ-TEXTUAIS
% ----------------------------------------------------------
\pretextual

% ---
% CAPA
% ---
\includepdf{./05-extras/capa.pdf}
% ---


% ---
% FOLHA DE ROSTO
% (o * indica que haverá a ficha bibliográfica)
% ---
\imprimirfolhaderosto
% ---


% ---
% Ficha bibliografica
% ---

% Isto é um exemplo de Ficha Catalográfica, ou ``Dados internacionais de
% catalogação-na-publicação''. Você pode utilizar este modelo como referência. 
% Porém, provavelmente a biblioteca da sua universidade lhe fornecerá um PDF
% com a ficha catalográfica definitiva após a defesa do trabalho. Quando estiver
% com o documento, salve-o como PDF no diretório do seu projeto e substitua todo
% o conteúdo de implementação deste arquivo pelo comando abaixo:
%
% \usepackage{pdfpages}		% necessário para comando \includepdf
% \begin{fichacatalografica}
%     \includepdf{fig_ficha_catalografica.pdf}
% \end{fichacatalografica}

% \begin{fichacatalografica}
% 	\sffamily
% 	\vspace*{\fill}                             % Posição vertical
% 	\begin{center}                              % Minipage Centralizado
% 	\fbox{\begin{minipage}[c][8cm]{13.5cm}		% Largura
% 	\small
% 	\imprimirautor
% 	%Sobrenome, Nome do autor
	
% 	\hspace{0.5cm} \imprimirtitulo  / \imprimirautor. --
% 	\imprimirlocal, \imprimirdata-
	
% 	\hspace{0.5cm} \thelastpage p. : il. (algumas color.) ; 30 cm.\\
	
% 	\hspace{0.5cm} \imprimirorientadorRotulo~\imprimirorientador\\
	
% 	\hspace{0.5cm}
% 	\parbox[t]{\textwidth}{\imprimirtipotrabalho~--~\imprimirinstituicao,
% 	\imprimirdata.}\\
	
% 	\hspace{0.5cm}
% 		1. Palavra-chave1.
% 		2. Palavra-chave2.
% 		2. Palavra-chave3.
% 		I. Orientador.
% 		II. Universidade xxx.
% 		III. Faculdade de xxx.
% 		IV. Título 			
% 	\end{minipage}}
% 	\end{center}
% \end{fichacatalografica}
% ---


% ---
% FOLHA DE APROVAÇÃO
% ---

% Isto é um exemplo de Folha de aprovação, elemento obrigatório da NBR
% 14724/2011 (seção 4.2.1.3). Você pode utilizar este modelo até a aprovação
% do trabalho. Após isso, substitua todo o conteúdo deste arquivo por uma
% imagem da página assinada pela banca com o comando abaixo:
%
% \begin{folhadeaprovacao}
% \includepdf{folhadeaprovacao_final.pdf}
% \end{folhadeaprovacao}
%
\setlength{\ABNTEXsignskip}{2cm}
\setlength{\ABNTEXsignwidth}{10cm}

\begin{folhadeaprovacao}
  \begin{center}
    {\ABNTEXchapterfont\large\imprimirautor}

    \vspace*{1cm}

    \begin{center}
      \ABNTEXchapterfont\bfseries\Large\imprimirtitulo
    \end{center}

    \vspace*{1cm}
    
    \hspace{.45\textwidth}
    \begin{minipage}{.5\textwidth}
        \imprimirpreambulo
    \end{minipage}%
   \end{center}

   \vspace*{1cm}
        
   \begin{center}
   Aprovado em 17 de julho de 2024.

   \vspace*{1cm}

   Banca examinadora:
   \end{center}

   \vspace*{-1cm}

   \assinatura{Prof. Ma. Caroline Santa Maria Rodrigues -- PUCRS}
   \assinatura{Prof. Dra. Clarissa Pinto Pizarro de Freitas -- PUCRS}
   \assinatura{\imprimirorientador -- PUCRS}
\end{folhadeaprovacao}
% ---


% ---
% Agradecimentos
% ---
\begin{agradecimentos}
Agradeço aos meus pais, Ademilson José Dias e Irinês Porfírio Dias, pelo apoio incondicional que me deram ao longo da vida,
incentivando mesmo as ideias questionáveis como aquela de fazer uma segunda graduação.

A minha irmã, Maria Eduarda Dias, cuja trajetória acadêmica é fonte de inspiração.

A minha namorada, Nathalia Fallavena Ceratti, que acompanhou a esta jornada de perto, oferecendo apoio e compreensão que
fizeram toda a diferença.

Aos meus amigos, Guilherme da Silva Petró, Norton Neves e Uilinton Luciano Marques da Silva, que foram compreensíveis nos
finais de semana em que não pude jogar RPG.

Ao professor Wagner de Lara Machado pela orientação e apoio no desenvolvimento deste trabalho.
\end{agradecimentos}
% ---


% ---
% RESUMOS
% ---
% Resumo em português
\setlength{\absparsep}{18pt} % ajusta o espaçamento dos parágrafos do resumo
\begin{resumo}
Desenvolvimentos na área de aprendizagem de máquina representam uma série de oportunidades para diversas áreas de pesquisa e de atuação profissional.
O uso de modelos de aprendizagem de máquina na predição de desfecho para tratamentos em saúde mental tem sido investigado e apresenta resultados promissores.
A aplicação desse tipo de modelo possibilitaria, por exemplo, maior assertividade na personalização de tratamentos psicoterápicos a nível individual.
Apesar dos benefícios potenciais, o tema ainda é pouco discutido nos espaços de pesquisa, formação e atuação em psicologia clínica. Este trabalho tem
por objetivo introduzir o tema de aprendizagem de máquina a psicólogos. Busca-se apresentar, de maneira acessível, os conceitos e terminologia básicos
em aprendizagem de máquina, proporcionando um panorama geral da área. Para complementar a exposição conceitual, é construído um exemplo a partir de dados
de um estudo sobre eficácia de intervenções em psicologia positiva para depressão via internet. Por fim, discutem-se algumas implicações do uso desse tipo
de tecnologia para a pesquisa e para a prática em psicoterapia.

 \textbf{Palavras-chave}: aprendizagem de máquina, psicoterapia, psicologia positiva, depressão.
\end{resumo}

% Resumo em inglês
\begin{resumo}[Abstract]
 \begin{otherlanguage*}{english}
   Advances in machine learning have opened up a multitude of opportunities for various research and professional fields. The use of machine learning models to
   predict treatment outcomes in mental health has been investigated showing promising results. The application of such models would allow greater assertiveness
   in tailoring psychotherapeutic treatments at the individual level. Despite its potential benefits, the topic remains unexplored in the fields of research, training
   and practice in clinical psychology. This work aims to introduce the topic of machine learning to psychologists. It presents the basic machine learning concepts and
   terminology in an accessible manner, providing a comprehensive overview of the field. To complement the conceptual discussion, an example is constructed using data
   from a study on the effectiveness of internet-based positive psychology interventions for depression. Finally, implications of using this type of technology in
   psychotherapy research and practice are discussed.
   \vspace{\onelineskip}
 
   \noindent 
   \textbf{Keywords}: machine learning, psychotherapy, positive psychology, depression.
 \end{otherlanguage*}
\end{resumo}


% ---
% LISTA DE ILUSTRAÇÕES
% ---
\pdfbookmark[0]{\listfigurename}{lof}
\listoffigures*
\cleardoublepage
% ---


% ---
% LISTA DE QUADROS
% ---
% \pdfbookmark[0]{\listofquadrosname}{loq}
% \listofquadros*
% \cleardoublepage
% ---


% ---
% LISTA DE TABELAS
% ---
% \pdfbookmark[0]{\listtablename}{lot}
% \listoftables*
% \cleardoublepage
% ---


% ---
% LISTA DE ABREVIATURAS E SIGLAS
% ---
% \begin{siglas}
%   \item[ABNT] Associação Brasileira de Normas Técnicas
%   \item[abnTeX] ABsurdas Normas para TeX
% \end{siglas}
% ---


% ---
% LISTA DE SÍMBOLOS
% ---
% \begin{simbolos}
%   \item[$ \Gamma $] Letra grega Gama
%   \item[$ \Lambda $] Lambda
%   \item[$ \zeta $] Letra grega minúscula zeta
%   \item[$ \in $] Pertence
% \end{simbolos}
% ---


% ---
% Sumário
% ---
\pdfbookmark[0]{\contentsname}{toc}
\tableofcontents*
\cleardoublepage
% ---


% ----------------------------------------------------------
% ELEMENTOS TEXTUAIS
% ----------------------------------------------------------
\textual

\include{./01-introducao/main.tex}

\include{./02-machine-learning/main.tex}

\include{./03-exemplo/main.tex}

\include{./04-conclusao/main.tex}

% ----------------------------------------------------------
% ELEMENTOS PÓS-TEXTUAIS
% ----------------------------------------------------------
\postextual
% ----------------------------------------------------------


% ----------------------------------------------------------
% REFERÊNCIAS
% ----------------------------------------------------------
\bibliography{references.bib}


% ----------------------------------------------------------
% GLOSSÁRIO
% ----------------------------------------------------------
%
% Consulte o manual da classe abntex2 para orientações sobre o glossário.
%
%\glossary

% ----------------------------------------------------------
% APÊNDICES
% ----------------------------------------------------------

% ---
% Inicia os apêndices
% ---
% \begin{apendicesenv}

% % Imprime uma página indicando o início dos apêndices
% \partapendices

% % ----------------------------------------------------------
% \chapter{Quisque libero justo}
% % ----------------------------------------------------------

% % ----------------------------------------------------------
% \chapter{Nullam elementum urna}
% % ----------------------------------------------------------

% \end{apendicesenv}
% ---


% ----------------------------------------------------------
% ANEXOS
% ----------------------------------------------------------

% ---
% Inicia os anexos
% ---
% \begin{anexosenv}

% % Imprime uma página indicando o início dos anexos
% \partanexos

% % ---
% \chapter{Morbi ultrices rutrum lorem}
% % ---

% % ---
% \chapter{Cras non urna sed feugiat cum}
% % ---


% % ---
% \chapter{Fusce facilisis lacinia dui}
% % ---

% \end{anexosenv}

%---------------------------------------------------------------------
% INDICE REMISSIVO
%---------------------------------------------------------------------
% \phantompart
% \printindex
%---------------------------------------------------------------------
\includepdf{./05-extras/contracapa.pdf}
\end{document}


\documentclass[
	% Opções da classe memoir
	12pt,				        % tamanho da fonte
	openany,			        % capítulos começam em pág ímpar (insere página vazia caso preciso)
	oneside,			        % para impressão em recto e verso. Oposto a oneside
	a4paper,			        % tamanho do papel. 
	% Opções da classe abntex2
	%chapter=TITLE,		        % títulos de capítulos convertidos em letras maiúsculas
	%section=TITLE,		        % títulos de seções convertidos em letras maiúsculas
	%subsection=TITLE,	        % títulos de subseções convertidos em letras maiúsculas
	%subsubsection=TITLE,       % títulos de subsubseções convertidos em letras maiúsculas
	% Opções do pacote babel
	english,			        % idioma adicional para hifenização
	brazil				        % o último idioma é o principal do documento
	]{abntex2}

% ---
% Pacotes básicos 
% ---
\usepackage{lmodern}			% Usa a fonte Latin Modern			
\usepackage[T1]{fontenc}		% Selecao de codigos de fonte.
\usepackage[utf8]{inputenc}		% Codificacao do documento (conversão automática dos acentos)
\usepackage{indentfirst}		% Indenta o primeiro parágrafo de cada seção.
\usepackage{color}				% Controle das cores
\usepackage{graphicx}			% Inclusão de gráficos
\usepackage{microtype} 			% para melhorias de justificação
% ---
		

% ---
% Pacotes de citações
% ---
% \usepackage[alf]{abntex2cite}						% Citações padrão ABNT
\usepackage[
  alf,
  abnt-nbr10520=1988,                     % Citações em sentence case
  abnt-emphasize=bf
  ]{abntex2cite}


% ---
% Pacotes adicionais
% ---
\usepackage{pdfpages}	% Inclusão de páginas em PDF


% ---
% Informações para CAPA e FOLHA DE ROSTO
% ---
\titulo{APRENDIZAGEM DE MÁQUINA PARA PREDIÇÃO DE DESFECHOS EM SAÚDE MENTAL: \textmd{UMA INTRODUÇÃO PARA PSICÓLOGOS}}
\autor{GUILHERME AUGUSTO DIAS}
\local{Porto Alegre}
\data{2024}
\orientador{Prof. Dr. Wagner de Lara Machado}
% \coorientador{}
% \instituicao{
%   Pontifícia Universidade Católica do Rio Grande do Sul
%   \par
%   Escola de Ciências da Saúde e da Vida
%   \par
%   Curso de Psicologia}
\tipotrabalho{TCC}
\preambulo{Trabalho de Conclusão de Curso apresentado como requisito para a obtenção do grau de Bacharel em Psicologia pela Escola de Ciências da Saúde e da Vida da Pontifícia Universidade Católica do Rio Grande do Sul.}
% ---

% ---
% Configurações de aparência do PDF final
% ---

% Alterando o aspecto da cor azul
\definecolor{blue}{RGB}{41,5,195}

% Informações do PDF
\makeatletter
\hypersetup{
		pdftitle={Aprendizagem de máquina para predição de desfechos em saúde mental: uma introdução para psicólogos}, 
		pdfauthor={Guilherme Augusto Dias},
    	pdfsubject={\imprimirpreambulo},
	    pdfcreator={LaTeX with abnTeX2},
		pdfkeywords={aprendizagem de máquina}{psicoterapia}{psicologia positiva}{depressão}, 
		colorlinks=true,		% true: colored links; false: boxed links
    	linkcolor=black,		% color of internal links
    	citecolor=black,		% color of links to bibliography
    	filecolor=black,		% color of file links
		urlcolor=black,
		bookmarksdepth=4
}
\makeatother
% --- 


% ---
% Posiciona figuras e tabelas no topo da página quando adicionadas sozinhas
% em um página em branco. Ver https://github.com/abntex/abntex2/issues/170
\makeatletter
\setlength{\@fptop}{5pt} % Set distance from top of page to first float
\makeatother
% ---

% ---
% Possibilita criação de Quadros e Lista de quadros.
% Ver https://github.com/abntex/abntex2/issues/176
%
\newcommand{\quadroname}{Quadro}
\newcommand{\listofquadrosname}{Lista de quadros}

\newfloat[chapter]{quadro}{loq}{\quadroname}
\newlistof{listofquadros}{loq}{\listofquadrosname}
\newlistentry{quadro}{loq}{0}

% Configurações para atender às regras da ABNT
\setfloatadjustment{quadro}{\centering}
\counterwithout{quadro}{chapter}
\renewcommand{\cftquadroname}{\quadroname\space} 
\renewcommand*{\cftquadroaftersnum}{\hfill--\hfill}

\setfloatlocations{quadro}{hbtp} % Ver https://github.com/abntex/abntex2/issues/176
% ---


% --- 
% Espaçamentos entre linhas e parágrafos 
% --- 

% O tamanho do parágrafo é dado por:
\setlength{\parindent}{1.3cm}

% Controle do espaçamento entre um parágrafo e outro:
\setlength{\parskip}{0.2cm}  % tente também \onelineskip

% ---
% Compila o índice
% ---
\makeindex
% ---


% ----
% Início do documento
% ----
\begin{document}

% Seleciona o idioma do documento (conforme pacotes do babel)
\selectlanguage{brazil}

% Retira espaço extra obsoleto entre as frases.
\frenchspacing 

% ----------------------------------------------------------
% ELEMENTOS PRÉ-TEXTUAIS
% ----------------------------------------------------------
\pretextual

% ---
% CAPA
% ---
\includepdf{./05-extras/capa.pdf}
% ---


% ---
% FOLHA DE ROSTO
% (o * indica que haverá a ficha bibliográfica)
% ---
\imprimirfolhaderosto
% ---


% ---
% Ficha bibliografica
% ---

% Isto é um exemplo de Ficha Catalográfica, ou ``Dados internacionais de
% catalogação-na-publicação''. Você pode utilizar este modelo como referência. 
% Porém, provavelmente a biblioteca da sua universidade lhe fornecerá um PDF
% com a ficha catalográfica definitiva após a defesa do trabalho. Quando estiver
% com o documento, salve-o como PDF no diretório do seu projeto e substitua todo
% o conteúdo de implementação deste arquivo pelo comando abaixo:
%
% \usepackage{pdfpages}		% necessário para comando \includepdf
% \begin{fichacatalografica}
%     \includepdf{fig_ficha_catalografica.pdf}
% \end{fichacatalografica}

% \begin{fichacatalografica}
% 	\sffamily
% 	\vspace*{\fill}                             % Posição vertical
% 	\begin{center}                              % Minipage Centralizado
% 	\fbox{\begin{minipage}[c][8cm]{13.5cm}		% Largura
% 	\small
% 	\imprimirautor
% 	%Sobrenome, Nome do autor
	
% 	\hspace{0.5cm} \imprimirtitulo  / \imprimirautor. --
% 	\imprimirlocal, \imprimirdata-
	
% 	\hspace{0.5cm} \thelastpage p. : il. (algumas color.) ; 30 cm.\\
	
% 	\hspace{0.5cm} \imprimirorientadorRotulo~\imprimirorientador\\
	
% 	\hspace{0.5cm}
% 	\parbox[t]{\textwidth}{\imprimirtipotrabalho~--~\imprimirinstituicao,
% 	\imprimirdata.}\\
	
% 	\hspace{0.5cm}
% 		1. Palavra-chave1.
% 		2. Palavra-chave2.
% 		2. Palavra-chave3.
% 		I. Orientador.
% 		II. Universidade xxx.
% 		III. Faculdade de xxx.
% 		IV. Título 			
% 	\end{minipage}}
% 	\end{center}
% \end{fichacatalografica}
% ---


% ---
% FOLHA DE APROVAÇÃO
% ---

% Isto é um exemplo de Folha de aprovação, elemento obrigatório da NBR
% 14724/2011 (seção 4.2.1.3). Você pode utilizar este modelo até a aprovação
% do trabalho. Após isso, substitua todo o conteúdo deste arquivo por uma
% imagem da página assinada pela banca com o comando abaixo:
%
% \begin{folhadeaprovacao}
% \includepdf{folhadeaprovacao_final.pdf}
% \end{folhadeaprovacao}
%
\setlength{\ABNTEXsignskip}{2cm}
\setlength{\ABNTEXsignwidth}{10cm}

\begin{folhadeaprovacao}
  \begin{center}
    {\ABNTEXchapterfont\large\imprimirautor}

    \vspace*{1cm}

    \begin{center}
      \ABNTEXchapterfont\bfseries\Large\imprimirtitulo
    \end{center}

    \vspace*{1cm}
    
    \hspace{.45\textwidth}
    \begin{minipage}{.5\textwidth}
        \imprimirpreambulo
    \end{minipage}%
   \end{center}

   \vspace*{1cm}
        
   \begin{center}
   Aprovado em 17 de julho de 2024.

   \vspace*{1cm}

   Banca examinadora:
   \end{center}

   \vspace*{-1cm}

   \assinatura{Prof. Ma. Caroline Santa Maria Rodrigues -- PUCRS}
   \assinatura{Prof. Dra. Clarissa Pinto Pizarro de Freitas -- PUCRS}
   \assinatura{\imprimirorientador -- PUCRS}
\end{folhadeaprovacao}
% ---


% ---
% Agradecimentos
% ---
\begin{agradecimentos}
Agradeço aos meus pais, Ademilson José Dias e Irinês Porfírio Dias, pelo apoio incondicional que me deram ao longo da vida,
incentivando mesmo as ideias questionáveis como aquela de fazer uma segunda graduação.

A minha irmã, Maria Eduarda Dias, cuja trajetória acadêmica é fonte de inspiração.

A minha namorada, Nathalia Fallavena Ceratti, que acompanhou a esta jornada de perto, oferecendo apoio e compreensão que
fizeram toda a diferença.

Aos meus amigos, Guilherme da Silva Petró, Norton Neves e Uilinton Luciano Marques da Silva, que foram compreensíveis nos
finais de semana em que não pude jogar RPG.

Ao professor Wagner de Lara Machado pela orientação e apoio no desenvolvimento deste trabalho.
\end{agradecimentos}
% ---


% ---
% RESUMOS
% ---
% Resumo em português
\setlength{\absparsep}{18pt} % ajusta o espaçamento dos parágrafos do resumo
\begin{resumo}
Desenvolvimentos na área de aprendizagem de máquina representam uma série de oportunidades para diversas áreas de pesquisa e de atuação profissional.
O uso de modelos de aprendizagem de máquina na predição de desfecho para tratamentos em saúde mental tem sido investigado e apresenta resultados promissores.
A aplicação desse tipo de modelo possibilitaria, por exemplo, maior assertividade na personalização de tratamentos psicoterápicos a nível individual.
Apesar dos benefícios potenciais, o tema ainda é pouco discutido nos espaços de pesquisa, formação e atuação em psicologia clínica. Este trabalho tem
por objetivo introduzir o tema de aprendizagem de máquina a psicólogos. Busca-se apresentar, de maneira acessível, os conceitos e terminologia básicos
em aprendizagem de máquina, proporcionando um panorama geral da área. Para complementar a exposição conceitual, é construído um exemplo a partir de dados
de um estudo sobre eficácia de intervenções em psicologia positiva para depressão via internet. Por fim, discutem-se algumas implicações do uso desse tipo
de tecnologia para a pesquisa e para a prática em psicoterapia.

 \textbf{Palavras-chave}: aprendizagem de máquina, psicoterapia, psicologia positiva, depressão.
\end{resumo}

% Resumo em inglês
\begin{resumo}[Abstract]
 \begin{otherlanguage*}{english}
   Advances in machine learning have opened up a multitude of opportunities for various research and professional fields. The use of machine learning models to
   predict treatment outcomes in mental health has been investigated showing promising results. The application of such models would allow greater assertiveness
   in tailoring psychotherapeutic treatments at the individual level. Despite its potential benefits, the topic remains unexplored in the fields of research, training
   and practice in clinical psychology. This work aims to introduce the topic of machine learning to psychologists. It presents the basic machine learning concepts and
   terminology in an accessible manner, providing a comprehensive overview of the field. To complement the conceptual discussion, an example is constructed using data
   from a study on the effectiveness of internet-based positive psychology interventions for depression. Finally, implications of using this type of technology in
   psychotherapy research and practice are discussed.
   \vspace{\onelineskip}
 
   \noindent 
   \textbf{Keywords}: machine learning, psychotherapy, positive psychology, depression.
 \end{otherlanguage*}
\end{resumo}


% ---
% LISTA DE ILUSTRAÇÕES
% ---
\pdfbookmark[0]{\listfigurename}{lof}
\listoffigures*
\cleardoublepage
% ---


% ---
% LISTA DE QUADROS
% ---
% \pdfbookmark[0]{\listofquadrosname}{loq}
% \listofquadros*
% \cleardoublepage
% ---


% ---
% LISTA DE TABELAS
% ---
% \pdfbookmark[0]{\listtablename}{lot}
% \listoftables*
% \cleardoublepage
% ---


% ---
% LISTA DE ABREVIATURAS E SIGLAS
% ---
% \begin{siglas}
%   \item[ABNT] Associação Brasileira de Normas Técnicas
%   \item[abnTeX] ABsurdas Normas para TeX
% \end{siglas}
% ---


% ---
% LISTA DE SÍMBOLOS
% ---
% \begin{simbolos}
%   \item[$ \Gamma $] Letra grega Gama
%   \item[$ \Lambda $] Lambda
%   \item[$ \zeta $] Letra grega minúscula zeta
%   \item[$ \in $] Pertence
% \end{simbolos}
% ---


% ---
% Sumário
% ---
\pdfbookmark[0]{\contentsname}{toc}
\tableofcontents*
\cleardoublepage
% ---


% ----------------------------------------------------------
% ELEMENTOS TEXTUAIS
% ----------------------------------------------------------
\textual

\include{./01-introducao/main.tex}

\include{./02-machine-learning/main.tex}

\include{./03-exemplo/main.tex}

\include{./04-conclusao/main.tex}

% ----------------------------------------------------------
% ELEMENTOS PÓS-TEXTUAIS
% ----------------------------------------------------------
\postextual
% ----------------------------------------------------------


% ----------------------------------------------------------
% REFERÊNCIAS
% ----------------------------------------------------------
\bibliography{references.bib}


% ----------------------------------------------------------
% GLOSSÁRIO
% ----------------------------------------------------------
%
% Consulte o manual da classe abntex2 para orientações sobre o glossário.
%
%\glossary

% ----------------------------------------------------------
% APÊNDICES
% ----------------------------------------------------------

% ---
% Inicia os apêndices
% ---
% \begin{apendicesenv}

% % Imprime uma página indicando o início dos apêndices
% \partapendices

% % ----------------------------------------------------------
% \chapter{Quisque libero justo}
% % ----------------------------------------------------------

% % ----------------------------------------------------------
% \chapter{Nullam elementum urna}
% % ----------------------------------------------------------

% \end{apendicesenv}
% ---


% ----------------------------------------------------------
% ANEXOS
% ----------------------------------------------------------

% ---
% Inicia os anexos
% ---
% \begin{anexosenv}

% % Imprime uma página indicando o início dos anexos
% \partanexos

% % ---
% \chapter{Morbi ultrices rutrum lorem}
% % ---

% % ---
% \chapter{Cras non urna sed feugiat cum}
% % ---


% % ---
% \chapter{Fusce facilisis lacinia dui}
% % ---

% \end{anexosenv}

%---------------------------------------------------------------------
% INDICE REMISSIVO
%---------------------------------------------------------------------
% \phantompart
% \printindex
%---------------------------------------------------------------------
\includepdf{./05-extras/contracapa.pdf}
\end{document}


\documentclass[
	% Opções da classe memoir
	12pt,				        % tamanho da fonte
	openany,			        % capítulos começam em pág ímpar (insere página vazia caso preciso)
	oneside,			        % para impressão em recto e verso. Oposto a oneside
	a4paper,			        % tamanho do papel. 
	% Opções da classe abntex2
	%chapter=TITLE,		        % títulos de capítulos convertidos em letras maiúsculas
	%section=TITLE,		        % títulos de seções convertidos em letras maiúsculas
	%subsection=TITLE,	        % títulos de subseções convertidos em letras maiúsculas
	%subsubsection=TITLE,       % títulos de subsubseções convertidos em letras maiúsculas
	% Opções do pacote babel
	english,			        % idioma adicional para hifenização
	brazil				        % o último idioma é o principal do documento
	]{abntex2}

% ---
% Pacotes básicos 
% ---
\usepackage{lmodern}			% Usa a fonte Latin Modern			
\usepackage[T1]{fontenc}		% Selecao de codigos de fonte.
\usepackage[utf8]{inputenc}		% Codificacao do documento (conversão automática dos acentos)
\usepackage{indentfirst}		% Indenta o primeiro parágrafo de cada seção.
\usepackage{color}				% Controle das cores
\usepackage{graphicx}			% Inclusão de gráficos
\usepackage{microtype} 			% para melhorias de justificação
% ---
		

% ---
% Pacotes de citações
% ---
% \usepackage[alf]{abntex2cite}						% Citações padrão ABNT
\usepackage[
  alf,
  abnt-nbr10520=1988,                     % Citações em sentence case
  abnt-emphasize=bf
  ]{abntex2cite}


% ---
% Pacotes adicionais
% ---
\usepackage{pdfpages}	% Inclusão de páginas em PDF


% ---
% Informações para CAPA e FOLHA DE ROSTO
% ---
\titulo{APRENDIZAGEM DE MÁQUINA PARA PREDIÇÃO DE DESFECHOS EM SAÚDE MENTAL: \textmd{UMA INTRODUÇÃO PARA PSICÓLOGOS}}
\autor{GUILHERME AUGUSTO DIAS}
\local{Porto Alegre}
\data{2024}
\orientador{Prof. Dr. Wagner de Lara Machado}
% \coorientador{}
% \instituicao{
%   Pontifícia Universidade Católica do Rio Grande do Sul
%   \par
%   Escola de Ciências da Saúde e da Vida
%   \par
%   Curso de Psicologia}
\tipotrabalho{TCC}
\preambulo{Trabalho de Conclusão de Curso apresentado como requisito para a obtenção do grau de Bacharel em Psicologia pela Escola de Ciências da Saúde e da Vida da Pontifícia Universidade Católica do Rio Grande do Sul.}
% ---

% ---
% Configurações de aparência do PDF final
% ---

% Alterando o aspecto da cor azul
\definecolor{blue}{RGB}{41,5,195}

% Informações do PDF
\makeatletter
\hypersetup{
		pdftitle={Aprendizagem de máquina para predição de desfechos em saúde mental: uma introdução para psicólogos}, 
		pdfauthor={Guilherme Augusto Dias},
    	pdfsubject={\imprimirpreambulo},
	    pdfcreator={LaTeX with abnTeX2},
		pdfkeywords={aprendizagem de máquina}{psicoterapia}{psicologia positiva}{depressão}, 
		colorlinks=true,		% true: colored links; false: boxed links
    	linkcolor=black,		% color of internal links
    	citecolor=black,		% color of links to bibliography
    	filecolor=black,		% color of file links
		urlcolor=black,
		bookmarksdepth=4
}
\makeatother
% --- 


% ---
% Posiciona figuras e tabelas no topo da página quando adicionadas sozinhas
% em um página em branco. Ver https://github.com/abntex/abntex2/issues/170
\makeatletter
\setlength{\@fptop}{5pt} % Set distance from top of page to first float
\makeatother
% ---

% ---
% Possibilita criação de Quadros e Lista de quadros.
% Ver https://github.com/abntex/abntex2/issues/176
%
\newcommand{\quadroname}{Quadro}
\newcommand{\listofquadrosname}{Lista de quadros}

\newfloat[chapter]{quadro}{loq}{\quadroname}
\newlistof{listofquadros}{loq}{\listofquadrosname}
\newlistentry{quadro}{loq}{0}

% Configurações para atender às regras da ABNT
\setfloatadjustment{quadro}{\centering}
\counterwithout{quadro}{chapter}
\renewcommand{\cftquadroname}{\quadroname\space} 
\renewcommand*{\cftquadroaftersnum}{\hfill--\hfill}

\setfloatlocations{quadro}{hbtp} % Ver https://github.com/abntex/abntex2/issues/176
% ---


% --- 
% Espaçamentos entre linhas e parágrafos 
% --- 

% O tamanho do parágrafo é dado por:
\setlength{\parindent}{1.3cm}

% Controle do espaçamento entre um parágrafo e outro:
\setlength{\parskip}{0.2cm}  % tente também \onelineskip

% ---
% Compila o índice
% ---
\makeindex
% ---


% ----
% Início do documento
% ----
\begin{document}

% Seleciona o idioma do documento (conforme pacotes do babel)
\selectlanguage{brazil}

% Retira espaço extra obsoleto entre as frases.
\frenchspacing 

% ----------------------------------------------------------
% ELEMENTOS PRÉ-TEXTUAIS
% ----------------------------------------------------------
\pretextual

% ---
% CAPA
% ---
\includepdf{./05-extras/capa.pdf}
% ---


% ---
% FOLHA DE ROSTO
% (o * indica que haverá a ficha bibliográfica)
% ---
\imprimirfolhaderosto
% ---


% ---
% Ficha bibliografica
% ---

% Isto é um exemplo de Ficha Catalográfica, ou ``Dados internacionais de
% catalogação-na-publicação''. Você pode utilizar este modelo como referência. 
% Porém, provavelmente a biblioteca da sua universidade lhe fornecerá um PDF
% com a ficha catalográfica definitiva após a defesa do trabalho. Quando estiver
% com o documento, salve-o como PDF no diretório do seu projeto e substitua todo
% o conteúdo de implementação deste arquivo pelo comando abaixo:
%
% \usepackage{pdfpages}		% necessário para comando \includepdf
% \begin{fichacatalografica}
%     \includepdf{fig_ficha_catalografica.pdf}
% \end{fichacatalografica}

% \begin{fichacatalografica}
% 	\sffamily
% 	\vspace*{\fill}                             % Posição vertical
% 	\begin{center}                              % Minipage Centralizado
% 	\fbox{\begin{minipage}[c][8cm]{13.5cm}		% Largura
% 	\small
% 	\imprimirautor
% 	%Sobrenome, Nome do autor
	
% 	\hspace{0.5cm} \imprimirtitulo  / \imprimirautor. --
% 	\imprimirlocal, \imprimirdata-
	
% 	\hspace{0.5cm} \thelastpage p. : il. (algumas color.) ; 30 cm.\\
	
% 	\hspace{0.5cm} \imprimirorientadorRotulo~\imprimirorientador\\
	
% 	\hspace{0.5cm}
% 	\parbox[t]{\textwidth}{\imprimirtipotrabalho~--~\imprimirinstituicao,
% 	\imprimirdata.}\\
	
% 	\hspace{0.5cm}
% 		1. Palavra-chave1.
% 		2. Palavra-chave2.
% 		2. Palavra-chave3.
% 		I. Orientador.
% 		II. Universidade xxx.
% 		III. Faculdade de xxx.
% 		IV. Título 			
% 	\end{minipage}}
% 	\end{center}
% \end{fichacatalografica}
% ---


% ---
% FOLHA DE APROVAÇÃO
% ---

% Isto é um exemplo de Folha de aprovação, elemento obrigatório da NBR
% 14724/2011 (seção 4.2.1.3). Você pode utilizar este modelo até a aprovação
% do trabalho. Após isso, substitua todo o conteúdo deste arquivo por uma
% imagem da página assinada pela banca com o comando abaixo:
%
% \begin{folhadeaprovacao}
% \includepdf{folhadeaprovacao_final.pdf}
% \end{folhadeaprovacao}
%
\setlength{\ABNTEXsignskip}{2cm}
\setlength{\ABNTEXsignwidth}{10cm}

\begin{folhadeaprovacao}
  \begin{center}
    {\ABNTEXchapterfont\large\imprimirautor}

    \vspace*{1cm}

    \begin{center}
      \ABNTEXchapterfont\bfseries\Large\imprimirtitulo
    \end{center}

    \vspace*{1cm}
    
    \hspace{.45\textwidth}
    \begin{minipage}{.5\textwidth}
        \imprimirpreambulo
    \end{minipage}%
   \end{center}

   \vspace*{1cm}
        
   \begin{center}
   Aprovado em 17 de julho de 2024.

   \vspace*{1cm}

   Banca examinadora:
   \end{center}

   \vspace*{-1cm}

   \assinatura{Prof. Ma. Caroline Santa Maria Rodrigues -- PUCRS}
   \assinatura{Prof. Dra. Clarissa Pinto Pizarro de Freitas -- PUCRS}
   \assinatura{\imprimirorientador -- PUCRS}
\end{folhadeaprovacao}
% ---


% ---
% Agradecimentos
% ---
\begin{agradecimentos}
Agradeço aos meus pais, Ademilson José Dias e Irinês Porfírio Dias, pelo apoio incondicional que me deram ao longo da vida,
incentivando mesmo as ideias questionáveis como aquela de fazer uma segunda graduação.

A minha irmã, Maria Eduarda Dias, cuja trajetória acadêmica é fonte de inspiração.

A minha namorada, Nathalia Fallavena Ceratti, que acompanhou a esta jornada de perto, oferecendo apoio e compreensão que
fizeram toda a diferença.

Aos meus amigos, Guilherme da Silva Petró, Norton Neves e Uilinton Luciano Marques da Silva, que foram compreensíveis nos
finais de semana em que não pude jogar RPG.

Ao professor Wagner de Lara Machado pela orientação e apoio no desenvolvimento deste trabalho.
\end{agradecimentos}
% ---


% ---
% RESUMOS
% ---
% Resumo em português
\setlength{\absparsep}{18pt} % ajusta o espaçamento dos parágrafos do resumo
\begin{resumo}
Desenvolvimentos na área de aprendizagem de máquina representam uma série de oportunidades para diversas áreas de pesquisa e de atuação profissional.
O uso de modelos de aprendizagem de máquina na predição de desfecho para tratamentos em saúde mental tem sido investigado e apresenta resultados promissores.
A aplicação desse tipo de modelo possibilitaria, por exemplo, maior assertividade na personalização de tratamentos psicoterápicos a nível individual.
Apesar dos benefícios potenciais, o tema ainda é pouco discutido nos espaços de pesquisa, formação e atuação em psicologia clínica. Este trabalho tem
por objetivo introduzir o tema de aprendizagem de máquina a psicólogos. Busca-se apresentar, de maneira acessível, os conceitos e terminologia básicos
em aprendizagem de máquina, proporcionando um panorama geral da área. Para complementar a exposição conceitual, é construído um exemplo a partir de dados
de um estudo sobre eficácia de intervenções em psicologia positiva para depressão via internet. Por fim, discutem-se algumas implicações do uso desse tipo
de tecnologia para a pesquisa e para a prática em psicoterapia.

 \textbf{Palavras-chave}: aprendizagem de máquina, psicoterapia, psicologia positiva, depressão.
\end{resumo}

% Resumo em inglês
\begin{resumo}[Abstract]
 \begin{otherlanguage*}{english}
   Advances in machine learning have opened up a multitude of opportunities for various research and professional fields. The use of machine learning models to
   predict treatment outcomes in mental health has been investigated showing promising results. The application of such models would allow greater assertiveness
   in tailoring psychotherapeutic treatments at the individual level. Despite its potential benefits, the topic remains unexplored in the fields of research, training
   and practice in clinical psychology. This work aims to introduce the topic of machine learning to psychologists. It presents the basic machine learning concepts and
   terminology in an accessible manner, providing a comprehensive overview of the field. To complement the conceptual discussion, an example is constructed using data
   from a study on the effectiveness of internet-based positive psychology interventions for depression. Finally, implications of using this type of technology in
   psychotherapy research and practice are discussed.
   \vspace{\onelineskip}
 
   \noindent 
   \textbf{Keywords}: machine learning, psychotherapy, positive psychology, depression.
 \end{otherlanguage*}
\end{resumo}


% ---
% LISTA DE ILUSTRAÇÕES
% ---
\pdfbookmark[0]{\listfigurename}{lof}
\listoffigures*
\cleardoublepage
% ---


% ---
% LISTA DE QUADROS
% ---
% \pdfbookmark[0]{\listofquadrosname}{loq}
% \listofquadros*
% \cleardoublepage
% ---


% ---
% LISTA DE TABELAS
% ---
% \pdfbookmark[0]{\listtablename}{lot}
% \listoftables*
% \cleardoublepage
% ---


% ---
% LISTA DE ABREVIATURAS E SIGLAS
% ---
% \begin{siglas}
%   \item[ABNT] Associação Brasileira de Normas Técnicas
%   \item[abnTeX] ABsurdas Normas para TeX
% \end{siglas}
% ---


% ---
% LISTA DE SÍMBOLOS
% ---
% \begin{simbolos}
%   \item[$ \Gamma $] Letra grega Gama
%   \item[$ \Lambda $] Lambda
%   \item[$ \zeta $] Letra grega minúscula zeta
%   \item[$ \in $] Pertence
% \end{simbolos}
% ---


% ---
% Sumário
% ---
\pdfbookmark[0]{\contentsname}{toc}
\tableofcontents*
\cleardoublepage
% ---


% ----------------------------------------------------------
% ELEMENTOS TEXTUAIS
% ----------------------------------------------------------
\textual

\include{./01-introducao/main.tex}

\include{./02-machine-learning/main.tex}

\include{./03-exemplo/main.tex}

\include{./04-conclusao/main.tex}

% ----------------------------------------------------------
% ELEMENTOS PÓS-TEXTUAIS
% ----------------------------------------------------------
\postextual
% ----------------------------------------------------------


% ----------------------------------------------------------
% REFERÊNCIAS
% ----------------------------------------------------------
\bibliography{references.bib}


% ----------------------------------------------------------
% GLOSSÁRIO
% ----------------------------------------------------------
%
% Consulte o manual da classe abntex2 para orientações sobre o glossário.
%
%\glossary

% ----------------------------------------------------------
% APÊNDICES
% ----------------------------------------------------------

% ---
% Inicia os apêndices
% ---
% \begin{apendicesenv}

% % Imprime uma página indicando o início dos apêndices
% \partapendices

% % ----------------------------------------------------------
% \chapter{Quisque libero justo}
% % ----------------------------------------------------------

% % ----------------------------------------------------------
% \chapter{Nullam elementum urna}
% % ----------------------------------------------------------

% \end{apendicesenv}
% ---


% ----------------------------------------------------------
% ANEXOS
% ----------------------------------------------------------

% ---
% Inicia os anexos
% ---
% \begin{anexosenv}

% % Imprime uma página indicando o início dos anexos
% \partanexos

% % ---
% \chapter{Morbi ultrices rutrum lorem}
% % ---

% % ---
% \chapter{Cras non urna sed feugiat cum}
% % ---


% % ---
% \chapter{Fusce facilisis lacinia dui}
% % ---

% \end{anexosenv}

%---------------------------------------------------------------------
% INDICE REMISSIVO
%---------------------------------------------------------------------
% \phantompart
% \printindex
%---------------------------------------------------------------------
\includepdf{./05-extras/contracapa.pdf}
\end{document}


% ----------------------------------------------------------
% ELEMENTOS PÓS-TEXTUAIS
% ----------------------------------------------------------
\postextual
% ----------------------------------------------------------


% ----------------------------------------------------------
% REFERÊNCIAS
% ----------------------------------------------------------
\bibliography{references.bib}


% ----------------------------------------------------------
% GLOSSÁRIO
% ----------------------------------------------------------
%
% Consulte o manual da classe abntex2 para orientações sobre o glossário.
%
%\glossary

% ----------------------------------------------------------
% APÊNDICES
% ----------------------------------------------------------

% ---
% Inicia os apêndices
% ---
% \begin{apendicesenv}

% % Imprime uma página indicando o início dos apêndices
% \partapendices

% % ----------------------------------------------------------
% \chapter{Quisque libero justo}
% % ----------------------------------------------------------

% % ----------------------------------------------------------
% \chapter{Nullam elementum urna}
% % ----------------------------------------------------------

% \end{apendicesenv}
% ---


% ----------------------------------------------------------
% ANEXOS
% ----------------------------------------------------------

% ---
% Inicia os anexos
% ---
% \begin{anexosenv}

% % Imprime uma página indicando o início dos anexos
% \partanexos

% % ---
% \chapter{Morbi ultrices rutrum lorem}
% % ---

% % ---
% \chapter{Cras non urna sed feugiat cum}
% % ---


% % ---
% \chapter{Fusce facilisis lacinia dui}
% % ---

% \end{anexosenv}

%---------------------------------------------------------------------
% INDICE REMISSIVO
%---------------------------------------------------------------------
% \phantompart
% \printindex
%---------------------------------------------------------------------
\includepdf{./05-extras/contracapa.pdf}
\end{document}


\documentclass[
	% Opções da classe memoir
	12pt,				        % tamanho da fonte
	openany,			        % capítulos começam em pág ímpar (insere página vazia caso preciso)
	oneside,			        % para impressão em recto e verso. Oposto a oneside
	a4paper,			        % tamanho do papel. 
	% Opções da classe abntex2
	%chapter=TITLE,		        % títulos de capítulos convertidos em letras maiúsculas
	%section=TITLE,		        % títulos de seções convertidos em letras maiúsculas
	%subsection=TITLE,	        % títulos de subseções convertidos em letras maiúsculas
	%subsubsection=TITLE,       % títulos de subsubseções convertidos em letras maiúsculas
	% Opções do pacote babel
	english,			        % idioma adicional para hifenização
	brazil				        % o último idioma é o principal do documento
	]{abntex2}

% ---
% Pacotes básicos 
% ---
\usepackage{lmodern}			% Usa a fonte Latin Modern			
\usepackage[T1]{fontenc}		% Selecao de codigos de fonte.
\usepackage[utf8]{inputenc}		% Codificacao do documento (conversão automática dos acentos)
\usepackage{indentfirst}		% Indenta o primeiro parágrafo de cada seção.
\usepackage{color}				% Controle das cores
\usepackage{graphicx}			% Inclusão de gráficos
\usepackage{microtype} 			% para melhorias de justificação
% ---
		

% ---
% Pacotes de citações
% ---
% \usepackage[alf]{abntex2cite}						% Citações padrão ABNT
\usepackage[
  alf,
  abnt-nbr10520=1988,                     % Citações em sentence case
  abnt-emphasize=bf
  ]{abntex2cite}


% ---
% Pacotes adicionais
% ---
\usepackage{pdfpages}	% Inclusão de páginas em PDF


% ---
% Informações para CAPA e FOLHA DE ROSTO
% ---
\titulo{APRENDIZAGEM DE MÁQUINA PARA PREDIÇÃO DE DESFECHOS EM SAÚDE MENTAL: \textmd{UMA INTRODUÇÃO PARA PSICÓLOGOS}}
\autor{GUILHERME AUGUSTO DIAS}
\local{Porto Alegre}
\data{2024}
\orientador{Prof. Dr. Wagner de Lara Machado}
% \coorientador{}
% \instituicao{
%   Pontifícia Universidade Católica do Rio Grande do Sul
%   \par
%   Escola de Ciências da Saúde e da Vida
%   \par
%   Curso de Psicologia}
\tipotrabalho{TCC}
\preambulo{Trabalho de Conclusão de Curso apresentado como requisito para a obtenção do grau de Bacharel em Psicologia pela Escola de Ciências da Saúde e da Vida da Pontifícia Universidade Católica do Rio Grande do Sul.}
% ---

% ---
% Configurações de aparência do PDF final
% ---

% Alterando o aspecto da cor azul
\definecolor{blue}{RGB}{41,5,195}

% Informações do PDF
\makeatletter
\hypersetup{
		pdftitle={Aprendizagem de máquina para predição de desfechos em saúde mental: uma introdução para psicólogos}, 
		pdfauthor={Guilherme Augusto Dias},
    	pdfsubject={\imprimirpreambulo},
	    pdfcreator={LaTeX with abnTeX2},
		pdfkeywords={aprendizagem de máquina}{psicoterapia}{psicologia positiva}{depressão}, 
		colorlinks=true,		% true: colored links; false: boxed links
    	linkcolor=black,		% color of internal links
    	citecolor=black,		% color of links to bibliography
    	filecolor=black,		% color of file links
		urlcolor=black,
		bookmarksdepth=4
}
\makeatother
% --- 


% ---
% Posiciona figuras e tabelas no topo da página quando adicionadas sozinhas
% em um página em branco. Ver https://github.com/abntex/abntex2/issues/170
\makeatletter
\setlength{\@fptop}{5pt} % Set distance from top of page to first float
\makeatother
% ---

% ---
% Possibilita criação de Quadros e Lista de quadros.
% Ver https://github.com/abntex/abntex2/issues/176
%
\newcommand{\quadroname}{Quadro}
\newcommand{\listofquadrosname}{Lista de quadros}

\newfloat[chapter]{quadro}{loq}{\quadroname}
\newlistof{listofquadros}{loq}{\listofquadrosname}
\newlistentry{quadro}{loq}{0}

% Configurações para atender às regras da ABNT
\setfloatadjustment{quadro}{\centering}
\counterwithout{quadro}{chapter}
\renewcommand{\cftquadroname}{\quadroname\space} 
\renewcommand*{\cftquadroaftersnum}{\hfill--\hfill}

\setfloatlocations{quadro}{hbtp} % Ver https://github.com/abntex/abntex2/issues/176
% ---


% --- 
% Espaçamentos entre linhas e parágrafos 
% --- 

% O tamanho do parágrafo é dado por:
\setlength{\parindent}{1.3cm}

% Controle do espaçamento entre um parágrafo e outro:
\setlength{\parskip}{0.2cm}  % tente também \onelineskip

% ---
% Compila o índice
% ---
\makeindex
% ---


% ----
% Início do documento
% ----
\begin{document}

% Seleciona o idioma do documento (conforme pacotes do babel)
\selectlanguage{brazil}

% Retira espaço extra obsoleto entre as frases.
\frenchspacing 

% ----------------------------------------------------------
% ELEMENTOS PRÉ-TEXTUAIS
% ----------------------------------------------------------
\pretextual

% ---
% CAPA
% ---
\includepdf{./05-extras/capa.pdf}
% ---


% ---
% FOLHA DE ROSTO
% (o * indica que haverá a ficha bibliográfica)
% ---
\imprimirfolhaderosto
% ---


% ---
% Ficha bibliografica
% ---

% Isto é um exemplo de Ficha Catalográfica, ou ``Dados internacionais de
% catalogação-na-publicação''. Você pode utilizar este modelo como referência. 
% Porém, provavelmente a biblioteca da sua universidade lhe fornecerá um PDF
% com a ficha catalográfica definitiva após a defesa do trabalho. Quando estiver
% com o documento, salve-o como PDF no diretório do seu projeto e substitua todo
% o conteúdo de implementação deste arquivo pelo comando abaixo:
%
% \usepackage{pdfpages}		% necessário para comando \includepdf
% \begin{fichacatalografica}
%     \includepdf{fig_ficha_catalografica.pdf}
% \end{fichacatalografica}

% \begin{fichacatalografica}
% 	\sffamily
% 	\vspace*{\fill}                             % Posição vertical
% 	\begin{center}                              % Minipage Centralizado
% 	\fbox{\begin{minipage}[c][8cm]{13.5cm}		% Largura
% 	\small
% 	\imprimirautor
% 	%Sobrenome, Nome do autor
	
% 	\hspace{0.5cm} \imprimirtitulo  / \imprimirautor. --
% 	\imprimirlocal, \imprimirdata-
	
% 	\hspace{0.5cm} \thelastpage p. : il. (algumas color.) ; 30 cm.\\
	
% 	\hspace{0.5cm} \imprimirorientadorRotulo~\imprimirorientador\\
	
% 	\hspace{0.5cm}
% 	\parbox[t]{\textwidth}{\imprimirtipotrabalho~--~\imprimirinstituicao,
% 	\imprimirdata.}\\
	
% 	\hspace{0.5cm}
% 		1. Palavra-chave1.
% 		2. Palavra-chave2.
% 		2. Palavra-chave3.
% 		I. Orientador.
% 		II. Universidade xxx.
% 		III. Faculdade de xxx.
% 		IV. Título 			
% 	\end{minipage}}
% 	\end{center}
% \end{fichacatalografica}
% ---


% ---
% FOLHA DE APROVAÇÃO
% ---

% Isto é um exemplo de Folha de aprovação, elemento obrigatório da NBR
% 14724/2011 (seção 4.2.1.3). Você pode utilizar este modelo até a aprovação
% do trabalho. Após isso, substitua todo o conteúdo deste arquivo por uma
% imagem da página assinada pela banca com o comando abaixo:
%
% \begin{folhadeaprovacao}
% \includepdf{folhadeaprovacao_final.pdf}
% \end{folhadeaprovacao}
%
\setlength{\ABNTEXsignskip}{2cm}
\setlength{\ABNTEXsignwidth}{10cm}

\begin{folhadeaprovacao}
  \begin{center}
    {\ABNTEXchapterfont\large\imprimirautor}

    \vspace*{1cm}

    \begin{center}
      \ABNTEXchapterfont\bfseries\Large\imprimirtitulo
    \end{center}

    \vspace*{1cm}
    
    \hspace{.45\textwidth}
    \begin{minipage}{.5\textwidth}
        \imprimirpreambulo
    \end{minipage}%
   \end{center}

   \vspace*{1cm}
        
   \begin{center}
   Aprovado em 17 de julho de 2024.

   \vspace*{1cm}

   Banca examinadora:
   \end{center}

   \vspace*{-1cm}

   \assinatura{Prof. Ma. Caroline Santa Maria Rodrigues -- PUCRS}
   \assinatura{Prof. Dra. Clarissa Pinto Pizarro de Freitas -- PUCRS}
   \assinatura{\imprimirorientador -- PUCRS}
\end{folhadeaprovacao}
% ---


% ---
% Agradecimentos
% ---
\begin{agradecimentos}
Agradeço aos meus pais, Ademilson José Dias e Irinês Porfírio Dias, pelo apoio incondicional que me deram ao longo da vida,
incentivando mesmo as ideias questionáveis como aquela de fazer uma segunda graduação.

A minha irmã, Maria Eduarda Dias, cuja trajetória acadêmica é fonte de inspiração.

A minha namorada, Nathalia Fallavena Ceratti, que acompanhou a esta jornada de perto, oferecendo apoio e compreensão que
fizeram toda a diferença.

Aos meus amigos, Guilherme da Silva Petró, Norton Neves e Uilinton Luciano Marques da Silva, que foram compreensíveis nos
finais de semana em que não pude jogar RPG.

Ao professor Wagner de Lara Machado pela orientação e apoio no desenvolvimento deste trabalho.
\end{agradecimentos}
% ---


% ---
% RESUMOS
% ---
% Resumo em português
\setlength{\absparsep}{18pt} % ajusta o espaçamento dos parágrafos do resumo
\begin{resumo}
Desenvolvimentos na área de aprendizagem de máquina representam uma série de oportunidades para diversas áreas de pesquisa e de atuação profissional.
O uso de modelos de aprendizagem de máquina na predição de desfecho para tratamentos em saúde mental tem sido investigado e apresenta resultados promissores.
A aplicação desse tipo de modelo possibilitaria, por exemplo, maior assertividade na personalização de tratamentos psicoterápicos a nível individual.
Apesar dos benefícios potenciais, o tema ainda é pouco discutido nos espaços de pesquisa, formação e atuação em psicologia clínica. Este trabalho tem
por objetivo introduzir o tema de aprendizagem de máquina a psicólogos. Busca-se apresentar, de maneira acessível, os conceitos e terminologia básicos
em aprendizagem de máquina, proporcionando um panorama geral da área. Para complementar a exposição conceitual, é construído um exemplo a partir de dados
de um estudo sobre eficácia de intervenções em psicologia positiva para depressão via internet. Por fim, discutem-se algumas implicações do uso desse tipo
de tecnologia para a pesquisa e para a prática em psicoterapia.

 \textbf{Palavras-chave}: aprendizagem de máquina, psicoterapia, psicologia positiva, depressão.
\end{resumo}

% Resumo em inglês
\begin{resumo}[Abstract]
 \begin{otherlanguage*}{english}
   Advances in machine learning have opened up a multitude of opportunities for various research and professional fields. The use of machine learning models to
   predict treatment outcomes in mental health has been investigated showing promising results. The application of such models would allow greater assertiveness
   in tailoring psychotherapeutic treatments at the individual level. Despite its potential benefits, the topic remains unexplored in the fields of research, training
   and practice in clinical psychology. This work aims to introduce the topic of machine learning to psychologists. It presents the basic machine learning concepts and
   terminology in an accessible manner, providing a comprehensive overview of the field. To complement the conceptual discussion, an example is constructed using data
   from a study on the effectiveness of internet-based positive psychology interventions for depression. Finally, implications of using this type of technology in
   psychotherapy research and practice are discussed.
   \vspace{\onelineskip}
 
   \noindent 
   \textbf{Keywords}: machine learning, psychotherapy, positive psychology, depression.
 \end{otherlanguage*}
\end{resumo}


% ---
% LISTA DE ILUSTRAÇÕES
% ---
\pdfbookmark[0]{\listfigurename}{lof}
\listoffigures*
\cleardoublepage
% ---


% ---
% LISTA DE QUADROS
% ---
% \pdfbookmark[0]{\listofquadrosname}{loq}
% \listofquadros*
% \cleardoublepage
% ---


% ---
% LISTA DE TABELAS
% ---
% \pdfbookmark[0]{\listtablename}{lot}
% \listoftables*
% \cleardoublepage
% ---


% ---
% LISTA DE ABREVIATURAS E SIGLAS
% ---
% \begin{siglas}
%   \item[ABNT] Associação Brasileira de Normas Técnicas
%   \item[abnTeX] ABsurdas Normas para TeX
% \end{siglas}
% ---


% ---
% LISTA DE SÍMBOLOS
% ---
% \begin{simbolos}
%   \item[$ \Gamma $] Letra grega Gama
%   \item[$ \Lambda $] Lambda
%   \item[$ \zeta $] Letra grega minúscula zeta
%   \item[$ \in $] Pertence
% \end{simbolos}
% ---


% ---
% Sumário
% ---
\pdfbookmark[0]{\contentsname}{toc}
\tableofcontents*
\cleardoublepage
% ---


% ----------------------------------------------------------
% ELEMENTOS TEXTUAIS
% ----------------------------------------------------------
\textual

\documentclass[
	% Opções da classe memoir
	12pt,				        % tamanho da fonte
	openany,			        % capítulos começam em pág ímpar (insere página vazia caso preciso)
	oneside,			        % para impressão em recto e verso. Oposto a oneside
	a4paper,			        % tamanho do papel. 
	% Opções da classe abntex2
	%chapter=TITLE,		        % títulos de capítulos convertidos em letras maiúsculas
	%section=TITLE,		        % títulos de seções convertidos em letras maiúsculas
	%subsection=TITLE,	        % títulos de subseções convertidos em letras maiúsculas
	%subsubsection=TITLE,       % títulos de subsubseções convertidos em letras maiúsculas
	% Opções do pacote babel
	english,			        % idioma adicional para hifenização
	brazil				        % o último idioma é o principal do documento
	]{abntex2}

% ---
% Pacotes básicos 
% ---
\usepackage{lmodern}			% Usa a fonte Latin Modern			
\usepackage[T1]{fontenc}		% Selecao de codigos de fonte.
\usepackage[utf8]{inputenc}		% Codificacao do documento (conversão automática dos acentos)
\usepackage{indentfirst}		% Indenta o primeiro parágrafo de cada seção.
\usepackage{color}				% Controle das cores
\usepackage{graphicx}			% Inclusão de gráficos
\usepackage{microtype} 			% para melhorias de justificação
% ---
		

% ---
% Pacotes de citações
% ---
% \usepackage[alf]{abntex2cite}						% Citações padrão ABNT
\usepackage[
  alf,
  abnt-nbr10520=1988,                     % Citações em sentence case
  abnt-emphasize=bf
  ]{abntex2cite}


% ---
% Pacotes adicionais
% ---
\usepackage{pdfpages}	% Inclusão de páginas em PDF


% ---
% Informações para CAPA e FOLHA DE ROSTO
% ---
\titulo{APRENDIZAGEM DE MÁQUINA PARA PREDIÇÃO DE DESFECHOS EM SAÚDE MENTAL: \textmd{UMA INTRODUÇÃO PARA PSICÓLOGOS}}
\autor{GUILHERME AUGUSTO DIAS}
\local{Porto Alegre}
\data{2024}
\orientador{Prof. Dr. Wagner de Lara Machado}
% \coorientador{}
% \instituicao{
%   Pontifícia Universidade Católica do Rio Grande do Sul
%   \par
%   Escola de Ciências da Saúde e da Vida
%   \par
%   Curso de Psicologia}
\tipotrabalho{TCC}
\preambulo{Trabalho de Conclusão de Curso apresentado como requisito para a obtenção do grau de Bacharel em Psicologia pela Escola de Ciências da Saúde e da Vida da Pontifícia Universidade Católica do Rio Grande do Sul.}
% ---

% ---
% Configurações de aparência do PDF final
% ---

% Alterando o aspecto da cor azul
\definecolor{blue}{RGB}{41,5,195}

% Informações do PDF
\makeatletter
\hypersetup{
		pdftitle={Aprendizagem de máquina para predição de desfechos em saúde mental: uma introdução para psicólogos}, 
		pdfauthor={Guilherme Augusto Dias},
    	pdfsubject={\imprimirpreambulo},
	    pdfcreator={LaTeX with abnTeX2},
		pdfkeywords={aprendizagem de máquina}{psicoterapia}{psicologia positiva}{depressão}, 
		colorlinks=true,		% true: colored links; false: boxed links
    	linkcolor=black,		% color of internal links
    	citecolor=black,		% color of links to bibliography
    	filecolor=black,		% color of file links
		urlcolor=black,
		bookmarksdepth=4
}
\makeatother
% --- 


% ---
% Posiciona figuras e tabelas no topo da página quando adicionadas sozinhas
% em um página em branco. Ver https://github.com/abntex/abntex2/issues/170
\makeatletter
\setlength{\@fptop}{5pt} % Set distance from top of page to first float
\makeatother
% ---

% ---
% Possibilita criação de Quadros e Lista de quadros.
% Ver https://github.com/abntex/abntex2/issues/176
%
\newcommand{\quadroname}{Quadro}
\newcommand{\listofquadrosname}{Lista de quadros}

\newfloat[chapter]{quadro}{loq}{\quadroname}
\newlistof{listofquadros}{loq}{\listofquadrosname}
\newlistentry{quadro}{loq}{0}

% Configurações para atender às regras da ABNT
\setfloatadjustment{quadro}{\centering}
\counterwithout{quadro}{chapter}
\renewcommand{\cftquadroname}{\quadroname\space} 
\renewcommand*{\cftquadroaftersnum}{\hfill--\hfill}

\setfloatlocations{quadro}{hbtp} % Ver https://github.com/abntex/abntex2/issues/176
% ---


% --- 
% Espaçamentos entre linhas e parágrafos 
% --- 

% O tamanho do parágrafo é dado por:
\setlength{\parindent}{1.3cm}

% Controle do espaçamento entre um parágrafo e outro:
\setlength{\parskip}{0.2cm}  % tente também \onelineskip

% ---
% Compila o índice
% ---
\makeindex
% ---


% ----
% Início do documento
% ----
\begin{document}

% Seleciona o idioma do documento (conforme pacotes do babel)
\selectlanguage{brazil}

% Retira espaço extra obsoleto entre as frases.
\frenchspacing 

% ----------------------------------------------------------
% ELEMENTOS PRÉ-TEXTUAIS
% ----------------------------------------------------------
\pretextual

% ---
% CAPA
% ---
\includepdf{./05-extras/capa.pdf}
% ---


% ---
% FOLHA DE ROSTO
% (o * indica que haverá a ficha bibliográfica)
% ---
\imprimirfolhaderosto
% ---


% ---
% Ficha bibliografica
% ---

% Isto é um exemplo de Ficha Catalográfica, ou ``Dados internacionais de
% catalogação-na-publicação''. Você pode utilizar este modelo como referência. 
% Porém, provavelmente a biblioteca da sua universidade lhe fornecerá um PDF
% com a ficha catalográfica definitiva após a defesa do trabalho. Quando estiver
% com o documento, salve-o como PDF no diretório do seu projeto e substitua todo
% o conteúdo de implementação deste arquivo pelo comando abaixo:
%
% \usepackage{pdfpages}		% necessário para comando \includepdf
% \begin{fichacatalografica}
%     \includepdf{fig_ficha_catalografica.pdf}
% \end{fichacatalografica}

% \begin{fichacatalografica}
% 	\sffamily
% 	\vspace*{\fill}                             % Posição vertical
% 	\begin{center}                              % Minipage Centralizado
% 	\fbox{\begin{minipage}[c][8cm]{13.5cm}		% Largura
% 	\small
% 	\imprimirautor
% 	%Sobrenome, Nome do autor
	
% 	\hspace{0.5cm} \imprimirtitulo  / \imprimirautor. --
% 	\imprimirlocal, \imprimirdata-
	
% 	\hspace{0.5cm} \thelastpage p. : il. (algumas color.) ; 30 cm.\\
	
% 	\hspace{0.5cm} \imprimirorientadorRotulo~\imprimirorientador\\
	
% 	\hspace{0.5cm}
% 	\parbox[t]{\textwidth}{\imprimirtipotrabalho~--~\imprimirinstituicao,
% 	\imprimirdata.}\\
	
% 	\hspace{0.5cm}
% 		1. Palavra-chave1.
% 		2. Palavra-chave2.
% 		2. Palavra-chave3.
% 		I. Orientador.
% 		II. Universidade xxx.
% 		III. Faculdade de xxx.
% 		IV. Título 			
% 	\end{minipage}}
% 	\end{center}
% \end{fichacatalografica}
% ---


% ---
% FOLHA DE APROVAÇÃO
% ---

% Isto é um exemplo de Folha de aprovação, elemento obrigatório da NBR
% 14724/2011 (seção 4.2.1.3). Você pode utilizar este modelo até a aprovação
% do trabalho. Após isso, substitua todo o conteúdo deste arquivo por uma
% imagem da página assinada pela banca com o comando abaixo:
%
% \begin{folhadeaprovacao}
% \includepdf{folhadeaprovacao_final.pdf}
% \end{folhadeaprovacao}
%
\setlength{\ABNTEXsignskip}{2cm}
\setlength{\ABNTEXsignwidth}{10cm}

\begin{folhadeaprovacao}
  \begin{center}
    {\ABNTEXchapterfont\large\imprimirautor}

    \vspace*{1cm}

    \begin{center}
      \ABNTEXchapterfont\bfseries\Large\imprimirtitulo
    \end{center}

    \vspace*{1cm}
    
    \hspace{.45\textwidth}
    \begin{minipage}{.5\textwidth}
        \imprimirpreambulo
    \end{minipage}%
   \end{center}

   \vspace*{1cm}
        
   \begin{center}
   Aprovado em 17 de julho de 2024.

   \vspace*{1cm}

   Banca examinadora:
   \end{center}

   \vspace*{-1cm}

   \assinatura{Prof. Ma. Caroline Santa Maria Rodrigues -- PUCRS}
   \assinatura{Prof. Dra. Clarissa Pinto Pizarro de Freitas -- PUCRS}
   \assinatura{\imprimirorientador -- PUCRS}
\end{folhadeaprovacao}
% ---


% ---
% Agradecimentos
% ---
\begin{agradecimentos}
Agradeço aos meus pais, Ademilson José Dias e Irinês Porfírio Dias, pelo apoio incondicional que me deram ao longo da vida,
incentivando mesmo as ideias questionáveis como aquela de fazer uma segunda graduação.

A minha irmã, Maria Eduarda Dias, cuja trajetória acadêmica é fonte de inspiração.

A minha namorada, Nathalia Fallavena Ceratti, que acompanhou a esta jornada de perto, oferecendo apoio e compreensão que
fizeram toda a diferença.

Aos meus amigos, Guilherme da Silva Petró, Norton Neves e Uilinton Luciano Marques da Silva, que foram compreensíveis nos
finais de semana em que não pude jogar RPG.

Ao professor Wagner de Lara Machado pela orientação e apoio no desenvolvimento deste trabalho.
\end{agradecimentos}
% ---


% ---
% RESUMOS
% ---
% Resumo em português
\setlength{\absparsep}{18pt} % ajusta o espaçamento dos parágrafos do resumo
\begin{resumo}
Desenvolvimentos na área de aprendizagem de máquina representam uma série de oportunidades para diversas áreas de pesquisa e de atuação profissional.
O uso de modelos de aprendizagem de máquina na predição de desfecho para tratamentos em saúde mental tem sido investigado e apresenta resultados promissores.
A aplicação desse tipo de modelo possibilitaria, por exemplo, maior assertividade na personalização de tratamentos psicoterápicos a nível individual.
Apesar dos benefícios potenciais, o tema ainda é pouco discutido nos espaços de pesquisa, formação e atuação em psicologia clínica. Este trabalho tem
por objetivo introduzir o tema de aprendizagem de máquina a psicólogos. Busca-se apresentar, de maneira acessível, os conceitos e terminologia básicos
em aprendizagem de máquina, proporcionando um panorama geral da área. Para complementar a exposição conceitual, é construído um exemplo a partir de dados
de um estudo sobre eficácia de intervenções em psicologia positiva para depressão via internet. Por fim, discutem-se algumas implicações do uso desse tipo
de tecnologia para a pesquisa e para a prática em psicoterapia.

 \textbf{Palavras-chave}: aprendizagem de máquina, psicoterapia, psicologia positiva, depressão.
\end{resumo}

% Resumo em inglês
\begin{resumo}[Abstract]
 \begin{otherlanguage*}{english}
   Advances in machine learning have opened up a multitude of opportunities for various research and professional fields. The use of machine learning models to
   predict treatment outcomes in mental health has been investigated showing promising results. The application of such models would allow greater assertiveness
   in tailoring psychotherapeutic treatments at the individual level. Despite its potential benefits, the topic remains unexplored in the fields of research, training
   and practice in clinical psychology. This work aims to introduce the topic of machine learning to psychologists. It presents the basic machine learning concepts and
   terminology in an accessible manner, providing a comprehensive overview of the field. To complement the conceptual discussion, an example is constructed using data
   from a study on the effectiveness of internet-based positive psychology interventions for depression. Finally, implications of using this type of technology in
   psychotherapy research and practice are discussed.
   \vspace{\onelineskip}
 
   \noindent 
   \textbf{Keywords}: machine learning, psychotherapy, positive psychology, depression.
 \end{otherlanguage*}
\end{resumo}


% ---
% LISTA DE ILUSTRAÇÕES
% ---
\pdfbookmark[0]{\listfigurename}{lof}
\listoffigures*
\cleardoublepage
% ---


% ---
% LISTA DE QUADROS
% ---
% \pdfbookmark[0]{\listofquadrosname}{loq}
% \listofquadros*
% \cleardoublepage
% ---


% ---
% LISTA DE TABELAS
% ---
% \pdfbookmark[0]{\listtablename}{lot}
% \listoftables*
% \cleardoublepage
% ---


% ---
% LISTA DE ABREVIATURAS E SIGLAS
% ---
% \begin{siglas}
%   \item[ABNT] Associação Brasileira de Normas Técnicas
%   \item[abnTeX] ABsurdas Normas para TeX
% \end{siglas}
% ---


% ---
% LISTA DE SÍMBOLOS
% ---
% \begin{simbolos}
%   \item[$ \Gamma $] Letra grega Gama
%   \item[$ \Lambda $] Lambda
%   \item[$ \zeta $] Letra grega minúscula zeta
%   \item[$ \in $] Pertence
% \end{simbolos}
% ---


% ---
% Sumário
% ---
\pdfbookmark[0]{\contentsname}{toc}
\tableofcontents*
\cleardoublepage
% ---


% ----------------------------------------------------------
% ELEMENTOS TEXTUAIS
% ----------------------------------------------------------
\textual

\include{./01-introducao/main.tex}

\include{./02-machine-learning/main.tex}

\include{./03-exemplo/main.tex}

\include{./04-conclusao/main.tex}

% ----------------------------------------------------------
% ELEMENTOS PÓS-TEXTUAIS
% ----------------------------------------------------------
\postextual
% ----------------------------------------------------------


% ----------------------------------------------------------
% REFERÊNCIAS
% ----------------------------------------------------------
\bibliography{references.bib}


% ----------------------------------------------------------
% GLOSSÁRIO
% ----------------------------------------------------------
%
% Consulte o manual da classe abntex2 para orientações sobre o glossário.
%
%\glossary

% ----------------------------------------------------------
% APÊNDICES
% ----------------------------------------------------------

% ---
% Inicia os apêndices
% ---
% \begin{apendicesenv}

% % Imprime uma página indicando o início dos apêndices
% \partapendices

% % ----------------------------------------------------------
% \chapter{Quisque libero justo}
% % ----------------------------------------------------------

% % ----------------------------------------------------------
% \chapter{Nullam elementum urna}
% % ----------------------------------------------------------

% \end{apendicesenv}
% ---


% ----------------------------------------------------------
% ANEXOS
% ----------------------------------------------------------

% ---
% Inicia os anexos
% ---
% \begin{anexosenv}

% % Imprime uma página indicando o início dos anexos
% \partanexos

% % ---
% \chapter{Morbi ultrices rutrum lorem}
% % ---

% % ---
% \chapter{Cras non urna sed feugiat cum}
% % ---


% % ---
% \chapter{Fusce facilisis lacinia dui}
% % ---

% \end{anexosenv}

%---------------------------------------------------------------------
% INDICE REMISSIVO
%---------------------------------------------------------------------
% \phantompart
% \printindex
%---------------------------------------------------------------------
\includepdf{./05-extras/contracapa.pdf}
\end{document}


\documentclass[
	% Opções da classe memoir
	12pt,				        % tamanho da fonte
	openany,			        % capítulos começam em pág ímpar (insere página vazia caso preciso)
	oneside,			        % para impressão em recto e verso. Oposto a oneside
	a4paper,			        % tamanho do papel. 
	% Opções da classe abntex2
	%chapter=TITLE,		        % títulos de capítulos convertidos em letras maiúsculas
	%section=TITLE,		        % títulos de seções convertidos em letras maiúsculas
	%subsection=TITLE,	        % títulos de subseções convertidos em letras maiúsculas
	%subsubsection=TITLE,       % títulos de subsubseções convertidos em letras maiúsculas
	% Opções do pacote babel
	english,			        % idioma adicional para hifenização
	brazil				        % o último idioma é o principal do documento
	]{abntex2}

% ---
% Pacotes básicos 
% ---
\usepackage{lmodern}			% Usa a fonte Latin Modern			
\usepackage[T1]{fontenc}		% Selecao de codigos de fonte.
\usepackage[utf8]{inputenc}		% Codificacao do documento (conversão automática dos acentos)
\usepackage{indentfirst}		% Indenta o primeiro parágrafo de cada seção.
\usepackage{color}				% Controle das cores
\usepackage{graphicx}			% Inclusão de gráficos
\usepackage{microtype} 			% para melhorias de justificação
% ---
		

% ---
% Pacotes de citações
% ---
% \usepackage[alf]{abntex2cite}						% Citações padrão ABNT
\usepackage[
  alf,
  abnt-nbr10520=1988,                     % Citações em sentence case
  abnt-emphasize=bf
  ]{abntex2cite}


% ---
% Pacotes adicionais
% ---
\usepackage{pdfpages}	% Inclusão de páginas em PDF


% ---
% Informações para CAPA e FOLHA DE ROSTO
% ---
\titulo{APRENDIZAGEM DE MÁQUINA PARA PREDIÇÃO DE DESFECHOS EM SAÚDE MENTAL: \textmd{UMA INTRODUÇÃO PARA PSICÓLOGOS}}
\autor{GUILHERME AUGUSTO DIAS}
\local{Porto Alegre}
\data{2024}
\orientador{Prof. Dr. Wagner de Lara Machado}
% \coorientador{}
% \instituicao{
%   Pontifícia Universidade Católica do Rio Grande do Sul
%   \par
%   Escola de Ciências da Saúde e da Vida
%   \par
%   Curso de Psicologia}
\tipotrabalho{TCC}
\preambulo{Trabalho de Conclusão de Curso apresentado como requisito para a obtenção do grau de Bacharel em Psicologia pela Escola de Ciências da Saúde e da Vida da Pontifícia Universidade Católica do Rio Grande do Sul.}
% ---

% ---
% Configurações de aparência do PDF final
% ---

% Alterando o aspecto da cor azul
\definecolor{blue}{RGB}{41,5,195}

% Informações do PDF
\makeatletter
\hypersetup{
		pdftitle={Aprendizagem de máquina para predição de desfechos em saúde mental: uma introdução para psicólogos}, 
		pdfauthor={Guilherme Augusto Dias},
    	pdfsubject={\imprimirpreambulo},
	    pdfcreator={LaTeX with abnTeX2},
		pdfkeywords={aprendizagem de máquina}{psicoterapia}{psicologia positiva}{depressão}, 
		colorlinks=true,		% true: colored links; false: boxed links
    	linkcolor=black,		% color of internal links
    	citecolor=black,		% color of links to bibliography
    	filecolor=black,		% color of file links
		urlcolor=black,
		bookmarksdepth=4
}
\makeatother
% --- 


% ---
% Posiciona figuras e tabelas no topo da página quando adicionadas sozinhas
% em um página em branco. Ver https://github.com/abntex/abntex2/issues/170
\makeatletter
\setlength{\@fptop}{5pt} % Set distance from top of page to first float
\makeatother
% ---

% ---
% Possibilita criação de Quadros e Lista de quadros.
% Ver https://github.com/abntex/abntex2/issues/176
%
\newcommand{\quadroname}{Quadro}
\newcommand{\listofquadrosname}{Lista de quadros}

\newfloat[chapter]{quadro}{loq}{\quadroname}
\newlistof{listofquadros}{loq}{\listofquadrosname}
\newlistentry{quadro}{loq}{0}

% Configurações para atender às regras da ABNT
\setfloatadjustment{quadro}{\centering}
\counterwithout{quadro}{chapter}
\renewcommand{\cftquadroname}{\quadroname\space} 
\renewcommand*{\cftquadroaftersnum}{\hfill--\hfill}

\setfloatlocations{quadro}{hbtp} % Ver https://github.com/abntex/abntex2/issues/176
% ---


% --- 
% Espaçamentos entre linhas e parágrafos 
% --- 

% O tamanho do parágrafo é dado por:
\setlength{\parindent}{1.3cm}

% Controle do espaçamento entre um parágrafo e outro:
\setlength{\parskip}{0.2cm}  % tente também \onelineskip

% ---
% Compila o índice
% ---
\makeindex
% ---


% ----
% Início do documento
% ----
\begin{document}

% Seleciona o idioma do documento (conforme pacotes do babel)
\selectlanguage{brazil}

% Retira espaço extra obsoleto entre as frases.
\frenchspacing 

% ----------------------------------------------------------
% ELEMENTOS PRÉ-TEXTUAIS
% ----------------------------------------------------------
\pretextual

% ---
% CAPA
% ---
\includepdf{./05-extras/capa.pdf}
% ---


% ---
% FOLHA DE ROSTO
% (o * indica que haverá a ficha bibliográfica)
% ---
\imprimirfolhaderosto
% ---


% ---
% Ficha bibliografica
% ---

% Isto é um exemplo de Ficha Catalográfica, ou ``Dados internacionais de
% catalogação-na-publicação''. Você pode utilizar este modelo como referência. 
% Porém, provavelmente a biblioteca da sua universidade lhe fornecerá um PDF
% com a ficha catalográfica definitiva após a defesa do trabalho. Quando estiver
% com o documento, salve-o como PDF no diretório do seu projeto e substitua todo
% o conteúdo de implementação deste arquivo pelo comando abaixo:
%
% \usepackage{pdfpages}		% necessário para comando \includepdf
% \begin{fichacatalografica}
%     \includepdf{fig_ficha_catalografica.pdf}
% \end{fichacatalografica}

% \begin{fichacatalografica}
% 	\sffamily
% 	\vspace*{\fill}                             % Posição vertical
% 	\begin{center}                              % Minipage Centralizado
% 	\fbox{\begin{minipage}[c][8cm]{13.5cm}		% Largura
% 	\small
% 	\imprimirautor
% 	%Sobrenome, Nome do autor
	
% 	\hspace{0.5cm} \imprimirtitulo  / \imprimirautor. --
% 	\imprimirlocal, \imprimirdata-
	
% 	\hspace{0.5cm} \thelastpage p. : il. (algumas color.) ; 30 cm.\\
	
% 	\hspace{0.5cm} \imprimirorientadorRotulo~\imprimirorientador\\
	
% 	\hspace{0.5cm}
% 	\parbox[t]{\textwidth}{\imprimirtipotrabalho~--~\imprimirinstituicao,
% 	\imprimirdata.}\\
	
% 	\hspace{0.5cm}
% 		1. Palavra-chave1.
% 		2. Palavra-chave2.
% 		2. Palavra-chave3.
% 		I. Orientador.
% 		II. Universidade xxx.
% 		III. Faculdade de xxx.
% 		IV. Título 			
% 	\end{minipage}}
% 	\end{center}
% \end{fichacatalografica}
% ---


% ---
% FOLHA DE APROVAÇÃO
% ---

% Isto é um exemplo de Folha de aprovação, elemento obrigatório da NBR
% 14724/2011 (seção 4.2.1.3). Você pode utilizar este modelo até a aprovação
% do trabalho. Após isso, substitua todo o conteúdo deste arquivo por uma
% imagem da página assinada pela banca com o comando abaixo:
%
% \begin{folhadeaprovacao}
% \includepdf{folhadeaprovacao_final.pdf}
% \end{folhadeaprovacao}
%
\setlength{\ABNTEXsignskip}{2cm}
\setlength{\ABNTEXsignwidth}{10cm}

\begin{folhadeaprovacao}
  \begin{center}
    {\ABNTEXchapterfont\large\imprimirautor}

    \vspace*{1cm}

    \begin{center}
      \ABNTEXchapterfont\bfseries\Large\imprimirtitulo
    \end{center}

    \vspace*{1cm}
    
    \hspace{.45\textwidth}
    \begin{minipage}{.5\textwidth}
        \imprimirpreambulo
    \end{minipage}%
   \end{center}

   \vspace*{1cm}
        
   \begin{center}
   Aprovado em 17 de julho de 2024.

   \vspace*{1cm}

   Banca examinadora:
   \end{center}

   \vspace*{-1cm}

   \assinatura{Prof. Ma. Caroline Santa Maria Rodrigues -- PUCRS}
   \assinatura{Prof. Dra. Clarissa Pinto Pizarro de Freitas -- PUCRS}
   \assinatura{\imprimirorientador -- PUCRS}
\end{folhadeaprovacao}
% ---


% ---
% Agradecimentos
% ---
\begin{agradecimentos}
Agradeço aos meus pais, Ademilson José Dias e Irinês Porfírio Dias, pelo apoio incondicional que me deram ao longo da vida,
incentivando mesmo as ideias questionáveis como aquela de fazer uma segunda graduação.

A minha irmã, Maria Eduarda Dias, cuja trajetória acadêmica é fonte de inspiração.

A minha namorada, Nathalia Fallavena Ceratti, que acompanhou a esta jornada de perto, oferecendo apoio e compreensão que
fizeram toda a diferença.

Aos meus amigos, Guilherme da Silva Petró, Norton Neves e Uilinton Luciano Marques da Silva, que foram compreensíveis nos
finais de semana em que não pude jogar RPG.

Ao professor Wagner de Lara Machado pela orientação e apoio no desenvolvimento deste trabalho.
\end{agradecimentos}
% ---


% ---
% RESUMOS
% ---
% Resumo em português
\setlength{\absparsep}{18pt} % ajusta o espaçamento dos parágrafos do resumo
\begin{resumo}
Desenvolvimentos na área de aprendizagem de máquina representam uma série de oportunidades para diversas áreas de pesquisa e de atuação profissional.
O uso de modelos de aprendizagem de máquina na predição de desfecho para tratamentos em saúde mental tem sido investigado e apresenta resultados promissores.
A aplicação desse tipo de modelo possibilitaria, por exemplo, maior assertividade na personalização de tratamentos psicoterápicos a nível individual.
Apesar dos benefícios potenciais, o tema ainda é pouco discutido nos espaços de pesquisa, formação e atuação em psicologia clínica. Este trabalho tem
por objetivo introduzir o tema de aprendizagem de máquina a psicólogos. Busca-se apresentar, de maneira acessível, os conceitos e terminologia básicos
em aprendizagem de máquina, proporcionando um panorama geral da área. Para complementar a exposição conceitual, é construído um exemplo a partir de dados
de um estudo sobre eficácia de intervenções em psicologia positiva para depressão via internet. Por fim, discutem-se algumas implicações do uso desse tipo
de tecnologia para a pesquisa e para a prática em psicoterapia.

 \textbf{Palavras-chave}: aprendizagem de máquina, psicoterapia, psicologia positiva, depressão.
\end{resumo}

% Resumo em inglês
\begin{resumo}[Abstract]
 \begin{otherlanguage*}{english}
   Advances in machine learning have opened up a multitude of opportunities for various research and professional fields. The use of machine learning models to
   predict treatment outcomes in mental health has been investigated showing promising results. The application of such models would allow greater assertiveness
   in tailoring psychotherapeutic treatments at the individual level. Despite its potential benefits, the topic remains unexplored in the fields of research, training
   and practice in clinical psychology. This work aims to introduce the topic of machine learning to psychologists. It presents the basic machine learning concepts and
   terminology in an accessible manner, providing a comprehensive overview of the field. To complement the conceptual discussion, an example is constructed using data
   from a study on the effectiveness of internet-based positive psychology interventions for depression. Finally, implications of using this type of technology in
   psychotherapy research and practice are discussed.
   \vspace{\onelineskip}
 
   \noindent 
   \textbf{Keywords}: machine learning, psychotherapy, positive psychology, depression.
 \end{otherlanguage*}
\end{resumo}


% ---
% LISTA DE ILUSTRAÇÕES
% ---
\pdfbookmark[0]{\listfigurename}{lof}
\listoffigures*
\cleardoublepage
% ---


% ---
% LISTA DE QUADROS
% ---
% \pdfbookmark[0]{\listofquadrosname}{loq}
% \listofquadros*
% \cleardoublepage
% ---


% ---
% LISTA DE TABELAS
% ---
% \pdfbookmark[0]{\listtablename}{lot}
% \listoftables*
% \cleardoublepage
% ---


% ---
% LISTA DE ABREVIATURAS E SIGLAS
% ---
% \begin{siglas}
%   \item[ABNT] Associação Brasileira de Normas Técnicas
%   \item[abnTeX] ABsurdas Normas para TeX
% \end{siglas}
% ---


% ---
% LISTA DE SÍMBOLOS
% ---
% \begin{simbolos}
%   \item[$ \Gamma $] Letra grega Gama
%   \item[$ \Lambda $] Lambda
%   \item[$ \zeta $] Letra grega minúscula zeta
%   \item[$ \in $] Pertence
% \end{simbolos}
% ---


% ---
% Sumário
% ---
\pdfbookmark[0]{\contentsname}{toc}
\tableofcontents*
\cleardoublepage
% ---


% ----------------------------------------------------------
% ELEMENTOS TEXTUAIS
% ----------------------------------------------------------
\textual

\include{./01-introducao/main.tex}

\include{./02-machine-learning/main.tex}

\include{./03-exemplo/main.tex}

\include{./04-conclusao/main.tex}

% ----------------------------------------------------------
% ELEMENTOS PÓS-TEXTUAIS
% ----------------------------------------------------------
\postextual
% ----------------------------------------------------------


% ----------------------------------------------------------
% REFERÊNCIAS
% ----------------------------------------------------------
\bibliography{references.bib}


% ----------------------------------------------------------
% GLOSSÁRIO
% ----------------------------------------------------------
%
% Consulte o manual da classe abntex2 para orientações sobre o glossário.
%
%\glossary

% ----------------------------------------------------------
% APÊNDICES
% ----------------------------------------------------------

% ---
% Inicia os apêndices
% ---
% \begin{apendicesenv}

% % Imprime uma página indicando o início dos apêndices
% \partapendices

% % ----------------------------------------------------------
% \chapter{Quisque libero justo}
% % ----------------------------------------------------------

% % ----------------------------------------------------------
% \chapter{Nullam elementum urna}
% % ----------------------------------------------------------

% \end{apendicesenv}
% ---


% ----------------------------------------------------------
% ANEXOS
% ----------------------------------------------------------

% ---
% Inicia os anexos
% ---
% \begin{anexosenv}

% % Imprime uma página indicando o início dos anexos
% \partanexos

% % ---
% \chapter{Morbi ultrices rutrum lorem}
% % ---

% % ---
% \chapter{Cras non urna sed feugiat cum}
% % ---


% % ---
% \chapter{Fusce facilisis lacinia dui}
% % ---

% \end{anexosenv}

%---------------------------------------------------------------------
% INDICE REMISSIVO
%---------------------------------------------------------------------
% \phantompart
% \printindex
%---------------------------------------------------------------------
\includepdf{./05-extras/contracapa.pdf}
\end{document}


\documentclass[
	% Opções da classe memoir
	12pt,				        % tamanho da fonte
	openany,			        % capítulos começam em pág ímpar (insere página vazia caso preciso)
	oneside,			        % para impressão em recto e verso. Oposto a oneside
	a4paper,			        % tamanho do papel. 
	% Opções da classe abntex2
	%chapter=TITLE,		        % títulos de capítulos convertidos em letras maiúsculas
	%section=TITLE,		        % títulos de seções convertidos em letras maiúsculas
	%subsection=TITLE,	        % títulos de subseções convertidos em letras maiúsculas
	%subsubsection=TITLE,       % títulos de subsubseções convertidos em letras maiúsculas
	% Opções do pacote babel
	english,			        % idioma adicional para hifenização
	brazil				        % o último idioma é o principal do documento
	]{abntex2}

% ---
% Pacotes básicos 
% ---
\usepackage{lmodern}			% Usa a fonte Latin Modern			
\usepackage[T1]{fontenc}		% Selecao de codigos de fonte.
\usepackage[utf8]{inputenc}		% Codificacao do documento (conversão automática dos acentos)
\usepackage{indentfirst}		% Indenta o primeiro parágrafo de cada seção.
\usepackage{color}				% Controle das cores
\usepackage{graphicx}			% Inclusão de gráficos
\usepackage{microtype} 			% para melhorias de justificação
% ---
		

% ---
% Pacotes de citações
% ---
% \usepackage[alf]{abntex2cite}						% Citações padrão ABNT
\usepackage[
  alf,
  abnt-nbr10520=1988,                     % Citações em sentence case
  abnt-emphasize=bf
  ]{abntex2cite}


% ---
% Pacotes adicionais
% ---
\usepackage{pdfpages}	% Inclusão de páginas em PDF


% ---
% Informações para CAPA e FOLHA DE ROSTO
% ---
\titulo{APRENDIZAGEM DE MÁQUINA PARA PREDIÇÃO DE DESFECHOS EM SAÚDE MENTAL: \textmd{UMA INTRODUÇÃO PARA PSICÓLOGOS}}
\autor{GUILHERME AUGUSTO DIAS}
\local{Porto Alegre}
\data{2024}
\orientador{Prof. Dr. Wagner de Lara Machado}
% \coorientador{}
% \instituicao{
%   Pontifícia Universidade Católica do Rio Grande do Sul
%   \par
%   Escola de Ciências da Saúde e da Vida
%   \par
%   Curso de Psicologia}
\tipotrabalho{TCC}
\preambulo{Trabalho de Conclusão de Curso apresentado como requisito para a obtenção do grau de Bacharel em Psicologia pela Escola de Ciências da Saúde e da Vida da Pontifícia Universidade Católica do Rio Grande do Sul.}
% ---

% ---
% Configurações de aparência do PDF final
% ---

% Alterando o aspecto da cor azul
\definecolor{blue}{RGB}{41,5,195}

% Informações do PDF
\makeatletter
\hypersetup{
		pdftitle={Aprendizagem de máquina para predição de desfechos em saúde mental: uma introdução para psicólogos}, 
		pdfauthor={Guilherme Augusto Dias},
    	pdfsubject={\imprimirpreambulo},
	    pdfcreator={LaTeX with abnTeX2},
		pdfkeywords={aprendizagem de máquina}{psicoterapia}{psicologia positiva}{depressão}, 
		colorlinks=true,		% true: colored links; false: boxed links
    	linkcolor=black,		% color of internal links
    	citecolor=black,		% color of links to bibliography
    	filecolor=black,		% color of file links
		urlcolor=black,
		bookmarksdepth=4
}
\makeatother
% --- 


% ---
% Posiciona figuras e tabelas no topo da página quando adicionadas sozinhas
% em um página em branco. Ver https://github.com/abntex/abntex2/issues/170
\makeatletter
\setlength{\@fptop}{5pt} % Set distance from top of page to first float
\makeatother
% ---

% ---
% Possibilita criação de Quadros e Lista de quadros.
% Ver https://github.com/abntex/abntex2/issues/176
%
\newcommand{\quadroname}{Quadro}
\newcommand{\listofquadrosname}{Lista de quadros}

\newfloat[chapter]{quadro}{loq}{\quadroname}
\newlistof{listofquadros}{loq}{\listofquadrosname}
\newlistentry{quadro}{loq}{0}

% Configurações para atender às regras da ABNT
\setfloatadjustment{quadro}{\centering}
\counterwithout{quadro}{chapter}
\renewcommand{\cftquadroname}{\quadroname\space} 
\renewcommand*{\cftquadroaftersnum}{\hfill--\hfill}

\setfloatlocations{quadro}{hbtp} % Ver https://github.com/abntex/abntex2/issues/176
% ---


% --- 
% Espaçamentos entre linhas e parágrafos 
% --- 

% O tamanho do parágrafo é dado por:
\setlength{\parindent}{1.3cm}

% Controle do espaçamento entre um parágrafo e outro:
\setlength{\parskip}{0.2cm}  % tente também \onelineskip

% ---
% Compila o índice
% ---
\makeindex
% ---


% ----
% Início do documento
% ----
\begin{document}

% Seleciona o idioma do documento (conforme pacotes do babel)
\selectlanguage{brazil}

% Retira espaço extra obsoleto entre as frases.
\frenchspacing 

% ----------------------------------------------------------
% ELEMENTOS PRÉ-TEXTUAIS
% ----------------------------------------------------------
\pretextual

% ---
% CAPA
% ---
\includepdf{./05-extras/capa.pdf}
% ---


% ---
% FOLHA DE ROSTO
% (o * indica que haverá a ficha bibliográfica)
% ---
\imprimirfolhaderosto
% ---


% ---
% Ficha bibliografica
% ---

% Isto é um exemplo de Ficha Catalográfica, ou ``Dados internacionais de
% catalogação-na-publicação''. Você pode utilizar este modelo como referência. 
% Porém, provavelmente a biblioteca da sua universidade lhe fornecerá um PDF
% com a ficha catalográfica definitiva após a defesa do trabalho. Quando estiver
% com o documento, salve-o como PDF no diretório do seu projeto e substitua todo
% o conteúdo de implementação deste arquivo pelo comando abaixo:
%
% \usepackage{pdfpages}		% necessário para comando \includepdf
% \begin{fichacatalografica}
%     \includepdf{fig_ficha_catalografica.pdf}
% \end{fichacatalografica}

% \begin{fichacatalografica}
% 	\sffamily
% 	\vspace*{\fill}                             % Posição vertical
% 	\begin{center}                              % Minipage Centralizado
% 	\fbox{\begin{minipage}[c][8cm]{13.5cm}		% Largura
% 	\small
% 	\imprimirautor
% 	%Sobrenome, Nome do autor
	
% 	\hspace{0.5cm} \imprimirtitulo  / \imprimirautor. --
% 	\imprimirlocal, \imprimirdata-
	
% 	\hspace{0.5cm} \thelastpage p. : il. (algumas color.) ; 30 cm.\\
	
% 	\hspace{0.5cm} \imprimirorientadorRotulo~\imprimirorientador\\
	
% 	\hspace{0.5cm}
% 	\parbox[t]{\textwidth}{\imprimirtipotrabalho~--~\imprimirinstituicao,
% 	\imprimirdata.}\\
	
% 	\hspace{0.5cm}
% 		1. Palavra-chave1.
% 		2. Palavra-chave2.
% 		2. Palavra-chave3.
% 		I. Orientador.
% 		II. Universidade xxx.
% 		III. Faculdade de xxx.
% 		IV. Título 			
% 	\end{minipage}}
% 	\end{center}
% \end{fichacatalografica}
% ---


% ---
% FOLHA DE APROVAÇÃO
% ---

% Isto é um exemplo de Folha de aprovação, elemento obrigatório da NBR
% 14724/2011 (seção 4.2.1.3). Você pode utilizar este modelo até a aprovação
% do trabalho. Após isso, substitua todo o conteúdo deste arquivo por uma
% imagem da página assinada pela banca com o comando abaixo:
%
% \begin{folhadeaprovacao}
% \includepdf{folhadeaprovacao_final.pdf}
% \end{folhadeaprovacao}
%
\setlength{\ABNTEXsignskip}{2cm}
\setlength{\ABNTEXsignwidth}{10cm}

\begin{folhadeaprovacao}
  \begin{center}
    {\ABNTEXchapterfont\large\imprimirautor}

    \vspace*{1cm}

    \begin{center}
      \ABNTEXchapterfont\bfseries\Large\imprimirtitulo
    \end{center}

    \vspace*{1cm}
    
    \hspace{.45\textwidth}
    \begin{minipage}{.5\textwidth}
        \imprimirpreambulo
    \end{minipage}%
   \end{center}

   \vspace*{1cm}
        
   \begin{center}
   Aprovado em 17 de julho de 2024.

   \vspace*{1cm}

   Banca examinadora:
   \end{center}

   \vspace*{-1cm}

   \assinatura{Prof. Ma. Caroline Santa Maria Rodrigues -- PUCRS}
   \assinatura{Prof. Dra. Clarissa Pinto Pizarro de Freitas -- PUCRS}
   \assinatura{\imprimirorientador -- PUCRS}
\end{folhadeaprovacao}
% ---


% ---
% Agradecimentos
% ---
\begin{agradecimentos}
Agradeço aos meus pais, Ademilson José Dias e Irinês Porfírio Dias, pelo apoio incondicional que me deram ao longo da vida,
incentivando mesmo as ideias questionáveis como aquela de fazer uma segunda graduação.

A minha irmã, Maria Eduarda Dias, cuja trajetória acadêmica é fonte de inspiração.

A minha namorada, Nathalia Fallavena Ceratti, que acompanhou a esta jornada de perto, oferecendo apoio e compreensão que
fizeram toda a diferença.

Aos meus amigos, Guilherme da Silva Petró, Norton Neves e Uilinton Luciano Marques da Silva, que foram compreensíveis nos
finais de semana em que não pude jogar RPG.

Ao professor Wagner de Lara Machado pela orientação e apoio no desenvolvimento deste trabalho.
\end{agradecimentos}
% ---


% ---
% RESUMOS
% ---
% Resumo em português
\setlength{\absparsep}{18pt} % ajusta o espaçamento dos parágrafos do resumo
\begin{resumo}
Desenvolvimentos na área de aprendizagem de máquina representam uma série de oportunidades para diversas áreas de pesquisa e de atuação profissional.
O uso de modelos de aprendizagem de máquina na predição de desfecho para tratamentos em saúde mental tem sido investigado e apresenta resultados promissores.
A aplicação desse tipo de modelo possibilitaria, por exemplo, maior assertividade na personalização de tratamentos psicoterápicos a nível individual.
Apesar dos benefícios potenciais, o tema ainda é pouco discutido nos espaços de pesquisa, formação e atuação em psicologia clínica. Este trabalho tem
por objetivo introduzir o tema de aprendizagem de máquina a psicólogos. Busca-se apresentar, de maneira acessível, os conceitos e terminologia básicos
em aprendizagem de máquina, proporcionando um panorama geral da área. Para complementar a exposição conceitual, é construído um exemplo a partir de dados
de um estudo sobre eficácia de intervenções em psicologia positiva para depressão via internet. Por fim, discutem-se algumas implicações do uso desse tipo
de tecnologia para a pesquisa e para a prática em psicoterapia.

 \textbf{Palavras-chave}: aprendizagem de máquina, psicoterapia, psicologia positiva, depressão.
\end{resumo}

% Resumo em inglês
\begin{resumo}[Abstract]
 \begin{otherlanguage*}{english}
   Advances in machine learning have opened up a multitude of opportunities for various research and professional fields. The use of machine learning models to
   predict treatment outcomes in mental health has been investigated showing promising results. The application of such models would allow greater assertiveness
   in tailoring psychotherapeutic treatments at the individual level. Despite its potential benefits, the topic remains unexplored in the fields of research, training
   and practice in clinical psychology. This work aims to introduce the topic of machine learning to psychologists. It presents the basic machine learning concepts and
   terminology in an accessible manner, providing a comprehensive overview of the field. To complement the conceptual discussion, an example is constructed using data
   from a study on the effectiveness of internet-based positive psychology interventions for depression. Finally, implications of using this type of technology in
   psychotherapy research and practice are discussed.
   \vspace{\onelineskip}
 
   \noindent 
   \textbf{Keywords}: machine learning, psychotherapy, positive psychology, depression.
 \end{otherlanguage*}
\end{resumo}


% ---
% LISTA DE ILUSTRAÇÕES
% ---
\pdfbookmark[0]{\listfigurename}{lof}
\listoffigures*
\cleardoublepage
% ---


% ---
% LISTA DE QUADROS
% ---
% \pdfbookmark[0]{\listofquadrosname}{loq}
% \listofquadros*
% \cleardoublepage
% ---


% ---
% LISTA DE TABELAS
% ---
% \pdfbookmark[0]{\listtablename}{lot}
% \listoftables*
% \cleardoublepage
% ---


% ---
% LISTA DE ABREVIATURAS E SIGLAS
% ---
% \begin{siglas}
%   \item[ABNT] Associação Brasileira de Normas Técnicas
%   \item[abnTeX] ABsurdas Normas para TeX
% \end{siglas}
% ---


% ---
% LISTA DE SÍMBOLOS
% ---
% \begin{simbolos}
%   \item[$ \Gamma $] Letra grega Gama
%   \item[$ \Lambda $] Lambda
%   \item[$ \zeta $] Letra grega minúscula zeta
%   \item[$ \in $] Pertence
% \end{simbolos}
% ---


% ---
% Sumário
% ---
\pdfbookmark[0]{\contentsname}{toc}
\tableofcontents*
\cleardoublepage
% ---


% ----------------------------------------------------------
% ELEMENTOS TEXTUAIS
% ----------------------------------------------------------
\textual

\include{./01-introducao/main.tex}

\include{./02-machine-learning/main.tex}

\include{./03-exemplo/main.tex}

\include{./04-conclusao/main.tex}

% ----------------------------------------------------------
% ELEMENTOS PÓS-TEXTUAIS
% ----------------------------------------------------------
\postextual
% ----------------------------------------------------------


% ----------------------------------------------------------
% REFERÊNCIAS
% ----------------------------------------------------------
\bibliography{references.bib}


% ----------------------------------------------------------
% GLOSSÁRIO
% ----------------------------------------------------------
%
% Consulte o manual da classe abntex2 para orientações sobre o glossário.
%
%\glossary

% ----------------------------------------------------------
% APÊNDICES
% ----------------------------------------------------------

% ---
% Inicia os apêndices
% ---
% \begin{apendicesenv}

% % Imprime uma página indicando o início dos apêndices
% \partapendices

% % ----------------------------------------------------------
% \chapter{Quisque libero justo}
% % ----------------------------------------------------------

% % ----------------------------------------------------------
% \chapter{Nullam elementum urna}
% % ----------------------------------------------------------

% \end{apendicesenv}
% ---


% ----------------------------------------------------------
% ANEXOS
% ----------------------------------------------------------

% ---
% Inicia os anexos
% ---
% \begin{anexosenv}

% % Imprime uma página indicando o início dos anexos
% \partanexos

% % ---
% \chapter{Morbi ultrices rutrum lorem}
% % ---

% % ---
% \chapter{Cras non urna sed feugiat cum}
% % ---


% % ---
% \chapter{Fusce facilisis lacinia dui}
% % ---

% \end{anexosenv}

%---------------------------------------------------------------------
% INDICE REMISSIVO
%---------------------------------------------------------------------
% \phantompart
% \printindex
%---------------------------------------------------------------------
\includepdf{./05-extras/contracapa.pdf}
\end{document}


\documentclass[
	% Opções da classe memoir
	12pt,				        % tamanho da fonte
	openany,			        % capítulos começam em pág ímpar (insere página vazia caso preciso)
	oneside,			        % para impressão em recto e verso. Oposto a oneside
	a4paper,			        % tamanho do papel. 
	% Opções da classe abntex2
	%chapter=TITLE,		        % títulos de capítulos convertidos em letras maiúsculas
	%section=TITLE,		        % títulos de seções convertidos em letras maiúsculas
	%subsection=TITLE,	        % títulos de subseções convertidos em letras maiúsculas
	%subsubsection=TITLE,       % títulos de subsubseções convertidos em letras maiúsculas
	% Opções do pacote babel
	english,			        % idioma adicional para hifenização
	brazil				        % o último idioma é o principal do documento
	]{abntex2}

% ---
% Pacotes básicos 
% ---
\usepackage{lmodern}			% Usa a fonte Latin Modern			
\usepackage[T1]{fontenc}		% Selecao de codigos de fonte.
\usepackage[utf8]{inputenc}		% Codificacao do documento (conversão automática dos acentos)
\usepackage{indentfirst}		% Indenta o primeiro parágrafo de cada seção.
\usepackage{color}				% Controle das cores
\usepackage{graphicx}			% Inclusão de gráficos
\usepackage{microtype} 			% para melhorias de justificação
% ---
		

% ---
% Pacotes de citações
% ---
% \usepackage[alf]{abntex2cite}						% Citações padrão ABNT
\usepackage[
  alf,
  abnt-nbr10520=1988,                     % Citações em sentence case
  abnt-emphasize=bf
  ]{abntex2cite}


% ---
% Pacotes adicionais
% ---
\usepackage{pdfpages}	% Inclusão de páginas em PDF


% ---
% Informações para CAPA e FOLHA DE ROSTO
% ---
\titulo{APRENDIZAGEM DE MÁQUINA PARA PREDIÇÃO DE DESFECHOS EM SAÚDE MENTAL: \textmd{UMA INTRODUÇÃO PARA PSICÓLOGOS}}
\autor{GUILHERME AUGUSTO DIAS}
\local{Porto Alegre}
\data{2024}
\orientador{Prof. Dr. Wagner de Lara Machado}
% \coorientador{}
% \instituicao{
%   Pontifícia Universidade Católica do Rio Grande do Sul
%   \par
%   Escola de Ciências da Saúde e da Vida
%   \par
%   Curso de Psicologia}
\tipotrabalho{TCC}
\preambulo{Trabalho de Conclusão de Curso apresentado como requisito para a obtenção do grau de Bacharel em Psicologia pela Escola de Ciências da Saúde e da Vida da Pontifícia Universidade Católica do Rio Grande do Sul.}
% ---

% ---
% Configurações de aparência do PDF final
% ---

% Alterando o aspecto da cor azul
\definecolor{blue}{RGB}{41,5,195}

% Informações do PDF
\makeatletter
\hypersetup{
		pdftitle={Aprendizagem de máquina para predição de desfechos em saúde mental: uma introdução para psicólogos}, 
		pdfauthor={Guilherme Augusto Dias},
    	pdfsubject={\imprimirpreambulo},
	    pdfcreator={LaTeX with abnTeX2},
		pdfkeywords={aprendizagem de máquina}{psicoterapia}{psicologia positiva}{depressão}, 
		colorlinks=true,		% true: colored links; false: boxed links
    	linkcolor=black,		% color of internal links
    	citecolor=black,		% color of links to bibliography
    	filecolor=black,		% color of file links
		urlcolor=black,
		bookmarksdepth=4
}
\makeatother
% --- 


% ---
% Posiciona figuras e tabelas no topo da página quando adicionadas sozinhas
% em um página em branco. Ver https://github.com/abntex/abntex2/issues/170
\makeatletter
\setlength{\@fptop}{5pt} % Set distance from top of page to first float
\makeatother
% ---

% ---
% Possibilita criação de Quadros e Lista de quadros.
% Ver https://github.com/abntex/abntex2/issues/176
%
\newcommand{\quadroname}{Quadro}
\newcommand{\listofquadrosname}{Lista de quadros}

\newfloat[chapter]{quadro}{loq}{\quadroname}
\newlistof{listofquadros}{loq}{\listofquadrosname}
\newlistentry{quadro}{loq}{0}

% Configurações para atender às regras da ABNT
\setfloatadjustment{quadro}{\centering}
\counterwithout{quadro}{chapter}
\renewcommand{\cftquadroname}{\quadroname\space} 
\renewcommand*{\cftquadroaftersnum}{\hfill--\hfill}

\setfloatlocations{quadro}{hbtp} % Ver https://github.com/abntex/abntex2/issues/176
% ---


% --- 
% Espaçamentos entre linhas e parágrafos 
% --- 

% O tamanho do parágrafo é dado por:
\setlength{\parindent}{1.3cm}

% Controle do espaçamento entre um parágrafo e outro:
\setlength{\parskip}{0.2cm}  % tente também \onelineskip

% ---
% Compila o índice
% ---
\makeindex
% ---


% ----
% Início do documento
% ----
\begin{document}

% Seleciona o idioma do documento (conforme pacotes do babel)
\selectlanguage{brazil}

% Retira espaço extra obsoleto entre as frases.
\frenchspacing 

% ----------------------------------------------------------
% ELEMENTOS PRÉ-TEXTUAIS
% ----------------------------------------------------------
\pretextual

% ---
% CAPA
% ---
\includepdf{./05-extras/capa.pdf}
% ---


% ---
% FOLHA DE ROSTO
% (o * indica que haverá a ficha bibliográfica)
% ---
\imprimirfolhaderosto
% ---


% ---
% Ficha bibliografica
% ---

% Isto é um exemplo de Ficha Catalográfica, ou ``Dados internacionais de
% catalogação-na-publicação''. Você pode utilizar este modelo como referência. 
% Porém, provavelmente a biblioteca da sua universidade lhe fornecerá um PDF
% com a ficha catalográfica definitiva após a defesa do trabalho. Quando estiver
% com o documento, salve-o como PDF no diretório do seu projeto e substitua todo
% o conteúdo de implementação deste arquivo pelo comando abaixo:
%
% \usepackage{pdfpages}		% necessário para comando \includepdf
% \begin{fichacatalografica}
%     \includepdf{fig_ficha_catalografica.pdf}
% \end{fichacatalografica}

% \begin{fichacatalografica}
% 	\sffamily
% 	\vspace*{\fill}                             % Posição vertical
% 	\begin{center}                              % Minipage Centralizado
% 	\fbox{\begin{minipage}[c][8cm]{13.5cm}		% Largura
% 	\small
% 	\imprimirautor
% 	%Sobrenome, Nome do autor
	
% 	\hspace{0.5cm} \imprimirtitulo  / \imprimirautor. --
% 	\imprimirlocal, \imprimirdata-
	
% 	\hspace{0.5cm} \thelastpage p. : il. (algumas color.) ; 30 cm.\\
	
% 	\hspace{0.5cm} \imprimirorientadorRotulo~\imprimirorientador\\
	
% 	\hspace{0.5cm}
% 	\parbox[t]{\textwidth}{\imprimirtipotrabalho~--~\imprimirinstituicao,
% 	\imprimirdata.}\\
	
% 	\hspace{0.5cm}
% 		1. Palavra-chave1.
% 		2. Palavra-chave2.
% 		2. Palavra-chave3.
% 		I. Orientador.
% 		II. Universidade xxx.
% 		III. Faculdade de xxx.
% 		IV. Título 			
% 	\end{minipage}}
% 	\end{center}
% \end{fichacatalografica}
% ---


% ---
% FOLHA DE APROVAÇÃO
% ---

% Isto é um exemplo de Folha de aprovação, elemento obrigatório da NBR
% 14724/2011 (seção 4.2.1.3). Você pode utilizar este modelo até a aprovação
% do trabalho. Após isso, substitua todo o conteúdo deste arquivo por uma
% imagem da página assinada pela banca com o comando abaixo:
%
% \begin{folhadeaprovacao}
% \includepdf{folhadeaprovacao_final.pdf}
% \end{folhadeaprovacao}
%
\setlength{\ABNTEXsignskip}{2cm}
\setlength{\ABNTEXsignwidth}{10cm}

\begin{folhadeaprovacao}
  \begin{center}
    {\ABNTEXchapterfont\large\imprimirautor}

    \vspace*{1cm}

    \begin{center}
      \ABNTEXchapterfont\bfseries\Large\imprimirtitulo
    \end{center}

    \vspace*{1cm}
    
    \hspace{.45\textwidth}
    \begin{minipage}{.5\textwidth}
        \imprimirpreambulo
    \end{minipage}%
   \end{center}

   \vspace*{1cm}
        
   \begin{center}
   Aprovado em 17 de julho de 2024.

   \vspace*{1cm}

   Banca examinadora:
   \end{center}

   \vspace*{-1cm}

   \assinatura{Prof. Ma. Caroline Santa Maria Rodrigues -- PUCRS}
   \assinatura{Prof. Dra. Clarissa Pinto Pizarro de Freitas -- PUCRS}
   \assinatura{\imprimirorientador -- PUCRS}
\end{folhadeaprovacao}
% ---


% ---
% Agradecimentos
% ---
\begin{agradecimentos}
Agradeço aos meus pais, Ademilson José Dias e Irinês Porfírio Dias, pelo apoio incondicional que me deram ao longo da vida,
incentivando mesmo as ideias questionáveis como aquela de fazer uma segunda graduação.

A minha irmã, Maria Eduarda Dias, cuja trajetória acadêmica é fonte de inspiração.

A minha namorada, Nathalia Fallavena Ceratti, que acompanhou a esta jornada de perto, oferecendo apoio e compreensão que
fizeram toda a diferença.

Aos meus amigos, Guilherme da Silva Petró, Norton Neves e Uilinton Luciano Marques da Silva, que foram compreensíveis nos
finais de semana em que não pude jogar RPG.

Ao professor Wagner de Lara Machado pela orientação e apoio no desenvolvimento deste trabalho.
\end{agradecimentos}
% ---


% ---
% RESUMOS
% ---
% Resumo em português
\setlength{\absparsep}{18pt} % ajusta o espaçamento dos parágrafos do resumo
\begin{resumo}
Desenvolvimentos na área de aprendizagem de máquina representam uma série de oportunidades para diversas áreas de pesquisa e de atuação profissional.
O uso de modelos de aprendizagem de máquina na predição de desfecho para tratamentos em saúde mental tem sido investigado e apresenta resultados promissores.
A aplicação desse tipo de modelo possibilitaria, por exemplo, maior assertividade na personalização de tratamentos psicoterápicos a nível individual.
Apesar dos benefícios potenciais, o tema ainda é pouco discutido nos espaços de pesquisa, formação e atuação em psicologia clínica. Este trabalho tem
por objetivo introduzir o tema de aprendizagem de máquina a psicólogos. Busca-se apresentar, de maneira acessível, os conceitos e terminologia básicos
em aprendizagem de máquina, proporcionando um panorama geral da área. Para complementar a exposição conceitual, é construído um exemplo a partir de dados
de um estudo sobre eficácia de intervenções em psicologia positiva para depressão via internet. Por fim, discutem-se algumas implicações do uso desse tipo
de tecnologia para a pesquisa e para a prática em psicoterapia.

 \textbf{Palavras-chave}: aprendizagem de máquina, psicoterapia, psicologia positiva, depressão.
\end{resumo}

% Resumo em inglês
\begin{resumo}[Abstract]
 \begin{otherlanguage*}{english}
   Advances in machine learning have opened up a multitude of opportunities for various research and professional fields. The use of machine learning models to
   predict treatment outcomes in mental health has been investigated showing promising results. The application of such models would allow greater assertiveness
   in tailoring psychotherapeutic treatments at the individual level. Despite its potential benefits, the topic remains unexplored in the fields of research, training
   and practice in clinical psychology. This work aims to introduce the topic of machine learning to psychologists. It presents the basic machine learning concepts and
   terminology in an accessible manner, providing a comprehensive overview of the field. To complement the conceptual discussion, an example is constructed using data
   from a study on the effectiveness of internet-based positive psychology interventions for depression. Finally, implications of using this type of technology in
   psychotherapy research and practice are discussed.
   \vspace{\onelineskip}
 
   \noindent 
   \textbf{Keywords}: machine learning, psychotherapy, positive psychology, depression.
 \end{otherlanguage*}
\end{resumo}


% ---
% LISTA DE ILUSTRAÇÕES
% ---
\pdfbookmark[0]{\listfigurename}{lof}
\listoffigures*
\cleardoublepage
% ---


% ---
% LISTA DE QUADROS
% ---
% \pdfbookmark[0]{\listofquadrosname}{loq}
% \listofquadros*
% \cleardoublepage
% ---


% ---
% LISTA DE TABELAS
% ---
% \pdfbookmark[0]{\listtablename}{lot}
% \listoftables*
% \cleardoublepage
% ---


% ---
% LISTA DE ABREVIATURAS E SIGLAS
% ---
% \begin{siglas}
%   \item[ABNT] Associação Brasileira de Normas Técnicas
%   \item[abnTeX] ABsurdas Normas para TeX
% \end{siglas}
% ---


% ---
% LISTA DE SÍMBOLOS
% ---
% \begin{simbolos}
%   \item[$ \Gamma $] Letra grega Gama
%   \item[$ \Lambda $] Lambda
%   \item[$ \zeta $] Letra grega minúscula zeta
%   \item[$ \in $] Pertence
% \end{simbolos}
% ---


% ---
% Sumário
% ---
\pdfbookmark[0]{\contentsname}{toc}
\tableofcontents*
\cleardoublepage
% ---


% ----------------------------------------------------------
% ELEMENTOS TEXTUAIS
% ----------------------------------------------------------
\textual

\include{./01-introducao/main.tex}

\include{./02-machine-learning/main.tex}

\include{./03-exemplo/main.tex}

\include{./04-conclusao/main.tex}

% ----------------------------------------------------------
% ELEMENTOS PÓS-TEXTUAIS
% ----------------------------------------------------------
\postextual
% ----------------------------------------------------------


% ----------------------------------------------------------
% REFERÊNCIAS
% ----------------------------------------------------------
\bibliography{references.bib}


% ----------------------------------------------------------
% GLOSSÁRIO
% ----------------------------------------------------------
%
% Consulte o manual da classe abntex2 para orientações sobre o glossário.
%
%\glossary

% ----------------------------------------------------------
% APÊNDICES
% ----------------------------------------------------------

% ---
% Inicia os apêndices
% ---
% \begin{apendicesenv}

% % Imprime uma página indicando o início dos apêndices
% \partapendices

% % ----------------------------------------------------------
% \chapter{Quisque libero justo}
% % ----------------------------------------------------------

% % ----------------------------------------------------------
% \chapter{Nullam elementum urna}
% % ----------------------------------------------------------

% \end{apendicesenv}
% ---


% ----------------------------------------------------------
% ANEXOS
% ----------------------------------------------------------

% ---
% Inicia os anexos
% ---
% \begin{anexosenv}

% % Imprime uma página indicando o início dos anexos
% \partanexos

% % ---
% \chapter{Morbi ultrices rutrum lorem}
% % ---

% % ---
% \chapter{Cras non urna sed feugiat cum}
% % ---


% % ---
% \chapter{Fusce facilisis lacinia dui}
% % ---

% \end{anexosenv}

%---------------------------------------------------------------------
% INDICE REMISSIVO
%---------------------------------------------------------------------
% \phantompart
% \printindex
%---------------------------------------------------------------------
\includepdf{./05-extras/contracapa.pdf}
\end{document}


% ----------------------------------------------------------
% ELEMENTOS PÓS-TEXTUAIS
% ----------------------------------------------------------
\postextual
% ----------------------------------------------------------


% ----------------------------------------------------------
% REFERÊNCIAS
% ----------------------------------------------------------
\bibliography{references.bib}


% ----------------------------------------------------------
% GLOSSÁRIO
% ----------------------------------------------------------
%
% Consulte o manual da classe abntex2 para orientações sobre o glossário.
%
%\glossary

% ----------------------------------------------------------
% APÊNDICES
% ----------------------------------------------------------

% ---
% Inicia os apêndices
% ---
% \begin{apendicesenv}

% % Imprime uma página indicando o início dos apêndices
% \partapendices

% % ----------------------------------------------------------
% \chapter{Quisque libero justo}
% % ----------------------------------------------------------

% % ----------------------------------------------------------
% \chapter{Nullam elementum urna}
% % ----------------------------------------------------------

% \end{apendicesenv}
% ---


% ----------------------------------------------------------
% ANEXOS
% ----------------------------------------------------------

% ---
% Inicia os anexos
% ---
% \begin{anexosenv}

% % Imprime uma página indicando o início dos anexos
% \partanexos

% % ---
% \chapter{Morbi ultrices rutrum lorem}
% % ---

% % ---
% \chapter{Cras non urna sed feugiat cum}
% % ---


% % ---
% \chapter{Fusce facilisis lacinia dui}
% % ---

% \end{anexosenv}

%---------------------------------------------------------------------
% INDICE REMISSIVO
%---------------------------------------------------------------------
% \phantompart
% \printindex
%---------------------------------------------------------------------
\includepdf{./05-extras/contracapa.pdf}
\end{document}


\documentclass[
	% Opções da classe memoir
	12pt,				        % tamanho da fonte
	openany,			        % capítulos começam em pág ímpar (insere página vazia caso preciso)
	oneside,			        % para impressão em recto e verso. Oposto a oneside
	a4paper,			        % tamanho do papel. 
	% Opções da classe abntex2
	%chapter=TITLE,		        % títulos de capítulos convertidos em letras maiúsculas
	%section=TITLE,		        % títulos de seções convertidos em letras maiúsculas
	%subsection=TITLE,	        % títulos de subseções convertidos em letras maiúsculas
	%subsubsection=TITLE,       % títulos de subsubseções convertidos em letras maiúsculas
	% Opções do pacote babel
	english,			        % idioma adicional para hifenização
	brazil				        % o último idioma é o principal do documento
	]{abntex2}

% ---
% Pacotes básicos 
% ---
\usepackage{lmodern}			% Usa a fonte Latin Modern			
\usepackage[T1]{fontenc}		% Selecao de codigos de fonte.
\usepackage[utf8]{inputenc}		% Codificacao do documento (conversão automática dos acentos)
\usepackage{indentfirst}		% Indenta o primeiro parágrafo de cada seção.
\usepackage{color}				% Controle das cores
\usepackage{graphicx}			% Inclusão de gráficos
\usepackage{microtype} 			% para melhorias de justificação
% ---
		

% ---
% Pacotes de citações
% ---
% \usepackage[alf]{abntex2cite}						% Citações padrão ABNT
\usepackage[
  alf,
  abnt-nbr10520=1988,                     % Citações em sentence case
  abnt-emphasize=bf
  ]{abntex2cite}


% ---
% Pacotes adicionais
% ---
\usepackage{pdfpages}	% Inclusão de páginas em PDF


% ---
% Informações para CAPA e FOLHA DE ROSTO
% ---
\titulo{APRENDIZAGEM DE MÁQUINA PARA PREDIÇÃO DE DESFECHOS EM SAÚDE MENTAL: \textmd{UMA INTRODUÇÃO PARA PSICÓLOGOS}}
\autor{GUILHERME AUGUSTO DIAS}
\local{Porto Alegre}
\data{2024}
\orientador{Prof. Dr. Wagner de Lara Machado}
% \coorientador{}
% \instituicao{
%   Pontifícia Universidade Católica do Rio Grande do Sul
%   \par
%   Escola de Ciências da Saúde e da Vida
%   \par
%   Curso de Psicologia}
\tipotrabalho{TCC}
\preambulo{Trabalho de Conclusão de Curso apresentado como requisito para a obtenção do grau de Bacharel em Psicologia pela Escola de Ciências da Saúde e da Vida da Pontifícia Universidade Católica do Rio Grande do Sul.}
% ---

% ---
% Configurações de aparência do PDF final
% ---

% Alterando o aspecto da cor azul
\definecolor{blue}{RGB}{41,5,195}

% Informações do PDF
\makeatletter
\hypersetup{
		pdftitle={Aprendizagem de máquina para predição de desfechos em saúde mental: uma introdução para psicólogos}, 
		pdfauthor={Guilherme Augusto Dias},
    	pdfsubject={\imprimirpreambulo},
	    pdfcreator={LaTeX with abnTeX2},
		pdfkeywords={aprendizagem de máquina}{psicoterapia}{psicologia positiva}{depressão}, 
		colorlinks=true,		% true: colored links; false: boxed links
    	linkcolor=black,		% color of internal links
    	citecolor=black,		% color of links to bibliography
    	filecolor=black,		% color of file links
		urlcolor=black,
		bookmarksdepth=4
}
\makeatother
% --- 


% ---
% Posiciona figuras e tabelas no topo da página quando adicionadas sozinhas
% em um página em branco. Ver https://github.com/abntex/abntex2/issues/170
\makeatletter
\setlength{\@fptop}{5pt} % Set distance from top of page to first float
\makeatother
% ---

% ---
% Possibilita criação de Quadros e Lista de quadros.
% Ver https://github.com/abntex/abntex2/issues/176
%
\newcommand{\quadroname}{Quadro}
\newcommand{\listofquadrosname}{Lista de quadros}

\newfloat[chapter]{quadro}{loq}{\quadroname}
\newlistof{listofquadros}{loq}{\listofquadrosname}
\newlistentry{quadro}{loq}{0}

% Configurações para atender às regras da ABNT
\setfloatadjustment{quadro}{\centering}
\counterwithout{quadro}{chapter}
\renewcommand{\cftquadroname}{\quadroname\space} 
\renewcommand*{\cftquadroaftersnum}{\hfill--\hfill}

\setfloatlocations{quadro}{hbtp} % Ver https://github.com/abntex/abntex2/issues/176
% ---


% --- 
% Espaçamentos entre linhas e parágrafos 
% --- 

% O tamanho do parágrafo é dado por:
\setlength{\parindent}{1.3cm}

% Controle do espaçamento entre um parágrafo e outro:
\setlength{\parskip}{0.2cm}  % tente também \onelineskip

% ---
% Compila o índice
% ---
\makeindex
% ---


% ----
% Início do documento
% ----
\begin{document}

% Seleciona o idioma do documento (conforme pacotes do babel)
\selectlanguage{brazil}

% Retira espaço extra obsoleto entre as frases.
\frenchspacing 

% ----------------------------------------------------------
% ELEMENTOS PRÉ-TEXTUAIS
% ----------------------------------------------------------
\pretextual

% ---
% CAPA
% ---
\includepdf{./05-extras/capa.pdf}
% ---


% ---
% FOLHA DE ROSTO
% (o * indica que haverá a ficha bibliográfica)
% ---
\imprimirfolhaderosto
% ---


% ---
% Ficha bibliografica
% ---

% Isto é um exemplo de Ficha Catalográfica, ou ``Dados internacionais de
% catalogação-na-publicação''. Você pode utilizar este modelo como referência. 
% Porém, provavelmente a biblioteca da sua universidade lhe fornecerá um PDF
% com a ficha catalográfica definitiva após a defesa do trabalho. Quando estiver
% com o documento, salve-o como PDF no diretório do seu projeto e substitua todo
% o conteúdo de implementação deste arquivo pelo comando abaixo:
%
% \usepackage{pdfpages}		% necessário para comando \includepdf
% \begin{fichacatalografica}
%     \includepdf{fig_ficha_catalografica.pdf}
% \end{fichacatalografica}

% \begin{fichacatalografica}
% 	\sffamily
% 	\vspace*{\fill}                             % Posição vertical
% 	\begin{center}                              % Minipage Centralizado
% 	\fbox{\begin{minipage}[c][8cm]{13.5cm}		% Largura
% 	\small
% 	\imprimirautor
% 	%Sobrenome, Nome do autor
	
% 	\hspace{0.5cm} \imprimirtitulo  / \imprimirautor. --
% 	\imprimirlocal, \imprimirdata-
	
% 	\hspace{0.5cm} \thelastpage p. : il. (algumas color.) ; 30 cm.\\
	
% 	\hspace{0.5cm} \imprimirorientadorRotulo~\imprimirorientador\\
	
% 	\hspace{0.5cm}
% 	\parbox[t]{\textwidth}{\imprimirtipotrabalho~--~\imprimirinstituicao,
% 	\imprimirdata.}\\
	
% 	\hspace{0.5cm}
% 		1. Palavra-chave1.
% 		2. Palavra-chave2.
% 		2. Palavra-chave3.
% 		I. Orientador.
% 		II. Universidade xxx.
% 		III. Faculdade de xxx.
% 		IV. Título 			
% 	\end{minipage}}
% 	\end{center}
% \end{fichacatalografica}
% ---


% ---
% FOLHA DE APROVAÇÃO
% ---

% Isto é um exemplo de Folha de aprovação, elemento obrigatório da NBR
% 14724/2011 (seção 4.2.1.3). Você pode utilizar este modelo até a aprovação
% do trabalho. Após isso, substitua todo o conteúdo deste arquivo por uma
% imagem da página assinada pela banca com o comando abaixo:
%
% \begin{folhadeaprovacao}
% \includepdf{folhadeaprovacao_final.pdf}
% \end{folhadeaprovacao}
%
\setlength{\ABNTEXsignskip}{2cm}
\setlength{\ABNTEXsignwidth}{10cm}

\begin{folhadeaprovacao}
  \begin{center}
    {\ABNTEXchapterfont\large\imprimirautor}

    \vspace*{1cm}

    \begin{center}
      \ABNTEXchapterfont\bfseries\Large\imprimirtitulo
    \end{center}

    \vspace*{1cm}
    
    \hspace{.45\textwidth}
    \begin{minipage}{.5\textwidth}
        \imprimirpreambulo
    \end{minipage}%
   \end{center}

   \vspace*{1cm}
        
   \begin{center}
   Aprovado em 17 de julho de 2024.

   \vspace*{1cm}

   Banca examinadora:
   \end{center}

   \vspace*{-1cm}

   \assinatura{Prof. Ma. Caroline Santa Maria Rodrigues -- PUCRS}
   \assinatura{Prof. Dra. Clarissa Pinto Pizarro de Freitas -- PUCRS}
   \assinatura{\imprimirorientador -- PUCRS}
\end{folhadeaprovacao}
% ---


% ---
% Agradecimentos
% ---
\begin{agradecimentos}
Agradeço aos meus pais, Ademilson José Dias e Irinês Porfírio Dias, pelo apoio incondicional que me deram ao longo da vida,
incentivando mesmo as ideias questionáveis como aquela de fazer uma segunda graduação.

A minha irmã, Maria Eduarda Dias, cuja trajetória acadêmica é fonte de inspiração.

A minha namorada, Nathalia Fallavena Ceratti, que acompanhou a esta jornada de perto, oferecendo apoio e compreensão que
fizeram toda a diferença.

Aos meus amigos, Guilherme da Silva Petró, Norton Neves e Uilinton Luciano Marques da Silva, que foram compreensíveis nos
finais de semana em que não pude jogar RPG.

Ao professor Wagner de Lara Machado pela orientação e apoio no desenvolvimento deste trabalho.
\end{agradecimentos}
% ---


% ---
% RESUMOS
% ---
% Resumo em português
\setlength{\absparsep}{18pt} % ajusta o espaçamento dos parágrafos do resumo
\begin{resumo}
Desenvolvimentos na área de aprendizagem de máquina representam uma série de oportunidades para diversas áreas de pesquisa e de atuação profissional.
O uso de modelos de aprendizagem de máquina na predição de desfecho para tratamentos em saúde mental tem sido investigado e apresenta resultados promissores.
A aplicação desse tipo de modelo possibilitaria, por exemplo, maior assertividade na personalização de tratamentos psicoterápicos a nível individual.
Apesar dos benefícios potenciais, o tema ainda é pouco discutido nos espaços de pesquisa, formação e atuação em psicologia clínica. Este trabalho tem
por objetivo introduzir o tema de aprendizagem de máquina a psicólogos. Busca-se apresentar, de maneira acessível, os conceitos e terminologia básicos
em aprendizagem de máquina, proporcionando um panorama geral da área. Para complementar a exposição conceitual, é construído um exemplo a partir de dados
de um estudo sobre eficácia de intervenções em psicologia positiva para depressão via internet. Por fim, discutem-se algumas implicações do uso desse tipo
de tecnologia para a pesquisa e para a prática em psicoterapia.

 \textbf{Palavras-chave}: aprendizagem de máquina, psicoterapia, psicologia positiva, depressão.
\end{resumo}

% Resumo em inglês
\begin{resumo}[Abstract]
 \begin{otherlanguage*}{english}
   Advances in machine learning have opened up a multitude of opportunities for various research and professional fields. The use of machine learning models to
   predict treatment outcomes in mental health has been investigated showing promising results. The application of such models would allow greater assertiveness
   in tailoring psychotherapeutic treatments at the individual level. Despite its potential benefits, the topic remains unexplored in the fields of research, training
   and practice in clinical psychology. This work aims to introduce the topic of machine learning to psychologists. It presents the basic machine learning concepts and
   terminology in an accessible manner, providing a comprehensive overview of the field. To complement the conceptual discussion, an example is constructed using data
   from a study on the effectiveness of internet-based positive psychology interventions for depression. Finally, implications of using this type of technology in
   psychotherapy research and practice are discussed.
   \vspace{\onelineskip}
 
   \noindent 
   \textbf{Keywords}: machine learning, psychotherapy, positive psychology, depression.
 \end{otherlanguage*}
\end{resumo}


% ---
% LISTA DE ILUSTRAÇÕES
% ---
\pdfbookmark[0]{\listfigurename}{lof}
\listoffigures*
\cleardoublepage
% ---


% ---
% LISTA DE QUADROS
% ---
% \pdfbookmark[0]{\listofquadrosname}{loq}
% \listofquadros*
% \cleardoublepage
% ---


% ---
% LISTA DE TABELAS
% ---
% \pdfbookmark[0]{\listtablename}{lot}
% \listoftables*
% \cleardoublepage
% ---


% ---
% LISTA DE ABREVIATURAS E SIGLAS
% ---
% \begin{siglas}
%   \item[ABNT] Associação Brasileira de Normas Técnicas
%   \item[abnTeX] ABsurdas Normas para TeX
% \end{siglas}
% ---


% ---
% LISTA DE SÍMBOLOS
% ---
% \begin{simbolos}
%   \item[$ \Gamma $] Letra grega Gama
%   \item[$ \Lambda $] Lambda
%   \item[$ \zeta $] Letra grega minúscula zeta
%   \item[$ \in $] Pertence
% \end{simbolos}
% ---


% ---
% Sumário
% ---
\pdfbookmark[0]{\contentsname}{toc}
\tableofcontents*
\cleardoublepage
% ---


% ----------------------------------------------------------
% ELEMENTOS TEXTUAIS
% ----------------------------------------------------------
\textual

\documentclass[
	% Opções da classe memoir
	12pt,				        % tamanho da fonte
	openany,			        % capítulos começam em pág ímpar (insere página vazia caso preciso)
	oneside,			        % para impressão em recto e verso. Oposto a oneside
	a4paper,			        % tamanho do papel. 
	% Opções da classe abntex2
	%chapter=TITLE,		        % títulos de capítulos convertidos em letras maiúsculas
	%section=TITLE,		        % títulos de seções convertidos em letras maiúsculas
	%subsection=TITLE,	        % títulos de subseções convertidos em letras maiúsculas
	%subsubsection=TITLE,       % títulos de subsubseções convertidos em letras maiúsculas
	% Opções do pacote babel
	english,			        % idioma adicional para hifenização
	brazil				        % o último idioma é o principal do documento
	]{abntex2}

% ---
% Pacotes básicos 
% ---
\usepackage{lmodern}			% Usa a fonte Latin Modern			
\usepackage[T1]{fontenc}		% Selecao de codigos de fonte.
\usepackage[utf8]{inputenc}		% Codificacao do documento (conversão automática dos acentos)
\usepackage{indentfirst}		% Indenta o primeiro parágrafo de cada seção.
\usepackage{color}				% Controle das cores
\usepackage{graphicx}			% Inclusão de gráficos
\usepackage{microtype} 			% para melhorias de justificação
% ---
		

% ---
% Pacotes de citações
% ---
% \usepackage[alf]{abntex2cite}						% Citações padrão ABNT
\usepackage[
  alf,
  abnt-nbr10520=1988,                     % Citações em sentence case
  abnt-emphasize=bf
  ]{abntex2cite}


% ---
% Pacotes adicionais
% ---
\usepackage{pdfpages}	% Inclusão de páginas em PDF


% ---
% Informações para CAPA e FOLHA DE ROSTO
% ---
\titulo{APRENDIZAGEM DE MÁQUINA PARA PREDIÇÃO DE DESFECHOS EM SAÚDE MENTAL: \textmd{UMA INTRODUÇÃO PARA PSICÓLOGOS}}
\autor{GUILHERME AUGUSTO DIAS}
\local{Porto Alegre}
\data{2024}
\orientador{Prof. Dr. Wagner de Lara Machado}
% \coorientador{}
% \instituicao{
%   Pontifícia Universidade Católica do Rio Grande do Sul
%   \par
%   Escola de Ciências da Saúde e da Vida
%   \par
%   Curso de Psicologia}
\tipotrabalho{TCC}
\preambulo{Trabalho de Conclusão de Curso apresentado como requisito para a obtenção do grau de Bacharel em Psicologia pela Escola de Ciências da Saúde e da Vida da Pontifícia Universidade Católica do Rio Grande do Sul.}
% ---

% ---
% Configurações de aparência do PDF final
% ---

% Alterando o aspecto da cor azul
\definecolor{blue}{RGB}{41,5,195}

% Informações do PDF
\makeatletter
\hypersetup{
		pdftitle={Aprendizagem de máquina para predição de desfechos em saúde mental: uma introdução para psicólogos}, 
		pdfauthor={Guilherme Augusto Dias},
    	pdfsubject={\imprimirpreambulo},
	    pdfcreator={LaTeX with abnTeX2},
		pdfkeywords={aprendizagem de máquina}{psicoterapia}{psicologia positiva}{depressão}, 
		colorlinks=true,		% true: colored links; false: boxed links
    	linkcolor=black,		% color of internal links
    	citecolor=black,		% color of links to bibliography
    	filecolor=black,		% color of file links
		urlcolor=black,
		bookmarksdepth=4
}
\makeatother
% --- 


% ---
% Posiciona figuras e tabelas no topo da página quando adicionadas sozinhas
% em um página em branco. Ver https://github.com/abntex/abntex2/issues/170
\makeatletter
\setlength{\@fptop}{5pt} % Set distance from top of page to first float
\makeatother
% ---

% ---
% Possibilita criação de Quadros e Lista de quadros.
% Ver https://github.com/abntex/abntex2/issues/176
%
\newcommand{\quadroname}{Quadro}
\newcommand{\listofquadrosname}{Lista de quadros}

\newfloat[chapter]{quadro}{loq}{\quadroname}
\newlistof{listofquadros}{loq}{\listofquadrosname}
\newlistentry{quadro}{loq}{0}

% Configurações para atender às regras da ABNT
\setfloatadjustment{quadro}{\centering}
\counterwithout{quadro}{chapter}
\renewcommand{\cftquadroname}{\quadroname\space} 
\renewcommand*{\cftquadroaftersnum}{\hfill--\hfill}

\setfloatlocations{quadro}{hbtp} % Ver https://github.com/abntex/abntex2/issues/176
% ---


% --- 
% Espaçamentos entre linhas e parágrafos 
% --- 

% O tamanho do parágrafo é dado por:
\setlength{\parindent}{1.3cm}

% Controle do espaçamento entre um parágrafo e outro:
\setlength{\parskip}{0.2cm}  % tente também \onelineskip

% ---
% Compila o índice
% ---
\makeindex
% ---


% ----
% Início do documento
% ----
\begin{document}

% Seleciona o idioma do documento (conforme pacotes do babel)
\selectlanguage{brazil}

% Retira espaço extra obsoleto entre as frases.
\frenchspacing 

% ----------------------------------------------------------
% ELEMENTOS PRÉ-TEXTUAIS
% ----------------------------------------------------------
\pretextual

% ---
% CAPA
% ---
\includepdf{./05-extras/capa.pdf}
% ---


% ---
% FOLHA DE ROSTO
% (o * indica que haverá a ficha bibliográfica)
% ---
\imprimirfolhaderosto
% ---


% ---
% Ficha bibliografica
% ---

% Isto é um exemplo de Ficha Catalográfica, ou ``Dados internacionais de
% catalogação-na-publicação''. Você pode utilizar este modelo como referência. 
% Porém, provavelmente a biblioteca da sua universidade lhe fornecerá um PDF
% com a ficha catalográfica definitiva após a defesa do trabalho. Quando estiver
% com o documento, salve-o como PDF no diretório do seu projeto e substitua todo
% o conteúdo de implementação deste arquivo pelo comando abaixo:
%
% \usepackage{pdfpages}		% necessário para comando \includepdf
% \begin{fichacatalografica}
%     \includepdf{fig_ficha_catalografica.pdf}
% \end{fichacatalografica}

% \begin{fichacatalografica}
% 	\sffamily
% 	\vspace*{\fill}                             % Posição vertical
% 	\begin{center}                              % Minipage Centralizado
% 	\fbox{\begin{minipage}[c][8cm]{13.5cm}		% Largura
% 	\small
% 	\imprimirautor
% 	%Sobrenome, Nome do autor
	
% 	\hspace{0.5cm} \imprimirtitulo  / \imprimirautor. --
% 	\imprimirlocal, \imprimirdata-
	
% 	\hspace{0.5cm} \thelastpage p. : il. (algumas color.) ; 30 cm.\\
	
% 	\hspace{0.5cm} \imprimirorientadorRotulo~\imprimirorientador\\
	
% 	\hspace{0.5cm}
% 	\parbox[t]{\textwidth}{\imprimirtipotrabalho~--~\imprimirinstituicao,
% 	\imprimirdata.}\\
	
% 	\hspace{0.5cm}
% 		1. Palavra-chave1.
% 		2. Palavra-chave2.
% 		2. Palavra-chave3.
% 		I. Orientador.
% 		II. Universidade xxx.
% 		III. Faculdade de xxx.
% 		IV. Título 			
% 	\end{minipage}}
% 	\end{center}
% \end{fichacatalografica}
% ---


% ---
% FOLHA DE APROVAÇÃO
% ---

% Isto é um exemplo de Folha de aprovação, elemento obrigatório da NBR
% 14724/2011 (seção 4.2.1.3). Você pode utilizar este modelo até a aprovação
% do trabalho. Após isso, substitua todo o conteúdo deste arquivo por uma
% imagem da página assinada pela banca com o comando abaixo:
%
% \begin{folhadeaprovacao}
% \includepdf{folhadeaprovacao_final.pdf}
% \end{folhadeaprovacao}
%
\setlength{\ABNTEXsignskip}{2cm}
\setlength{\ABNTEXsignwidth}{10cm}

\begin{folhadeaprovacao}
  \begin{center}
    {\ABNTEXchapterfont\large\imprimirautor}

    \vspace*{1cm}

    \begin{center}
      \ABNTEXchapterfont\bfseries\Large\imprimirtitulo
    \end{center}

    \vspace*{1cm}
    
    \hspace{.45\textwidth}
    \begin{minipage}{.5\textwidth}
        \imprimirpreambulo
    \end{minipage}%
   \end{center}

   \vspace*{1cm}
        
   \begin{center}
   Aprovado em 17 de julho de 2024.

   \vspace*{1cm}

   Banca examinadora:
   \end{center}

   \vspace*{-1cm}

   \assinatura{Prof. Ma. Caroline Santa Maria Rodrigues -- PUCRS}
   \assinatura{Prof. Dra. Clarissa Pinto Pizarro de Freitas -- PUCRS}
   \assinatura{\imprimirorientador -- PUCRS}
\end{folhadeaprovacao}
% ---


% ---
% Agradecimentos
% ---
\begin{agradecimentos}
Agradeço aos meus pais, Ademilson José Dias e Irinês Porfírio Dias, pelo apoio incondicional que me deram ao longo da vida,
incentivando mesmo as ideias questionáveis como aquela de fazer uma segunda graduação.

A minha irmã, Maria Eduarda Dias, cuja trajetória acadêmica é fonte de inspiração.

A minha namorada, Nathalia Fallavena Ceratti, que acompanhou a esta jornada de perto, oferecendo apoio e compreensão que
fizeram toda a diferença.

Aos meus amigos, Guilherme da Silva Petró, Norton Neves e Uilinton Luciano Marques da Silva, que foram compreensíveis nos
finais de semana em que não pude jogar RPG.

Ao professor Wagner de Lara Machado pela orientação e apoio no desenvolvimento deste trabalho.
\end{agradecimentos}
% ---


% ---
% RESUMOS
% ---
% Resumo em português
\setlength{\absparsep}{18pt} % ajusta o espaçamento dos parágrafos do resumo
\begin{resumo}
Desenvolvimentos na área de aprendizagem de máquina representam uma série de oportunidades para diversas áreas de pesquisa e de atuação profissional.
O uso de modelos de aprendizagem de máquina na predição de desfecho para tratamentos em saúde mental tem sido investigado e apresenta resultados promissores.
A aplicação desse tipo de modelo possibilitaria, por exemplo, maior assertividade na personalização de tratamentos psicoterápicos a nível individual.
Apesar dos benefícios potenciais, o tema ainda é pouco discutido nos espaços de pesquisa, formação e atuação em psicologia clínica. Este trabalho tem
por objetivo introduzir o tema de aprendizagem de máquina a psicólogos. Busca-se apresentar, de maneira acessível, os conceitos e terminologia básicos
em aprendizagem de máquina, proporcionando um panorama geral da área. Para complementar a exposição conceitual, é construído um exemplo a partir de dados
de um estudo sobre eficácia de intervenções em psicologia positiva para depressão via internet. Por fim, discutem-se algumas implicações do uso desse tipo
de tecnologia para a pesquisa e para a prática em psicoterapia.

 \textbf{Palavras-chave}: aprendizagem de máquina, psicoterapia, psicologia positiva, depressão.
\end{resumo}

% Resumo em inglês
\begin{resumo}[Abstract]
 \begin{otherlanguage*}{english}
   Advances in machine learning have opened up a multitude of opportunities for various research and professional fields. The use of machine learning models to
   predict treatment outcomes in mental health has been investigated showing promising results. The application of such models would allow greater assertiveness
   in tailoring psychotherapeutic treatments at the individual level. Despite its potential benefits, the topic remains unexplored in the fields of research, training
   and practice in clinical psychology. This work aims to introduce the topic of machine learning to psychologists. It presents the basic machine learning concepts and
   terminology in an accessible manner, providing a comprehensive overview of the field. To complement the conceptual discussion, an example is constructed using data
   from a study on the effectiveness of internet-based positive psychology interventions for depression. Finally, implications of using this type of technology in
   psychotherapy research and practice are discussed.
   \vspace{\onelineskip}
 
   \noindent 
   \textbf{Keywords}: machine learning, psychotherapy, positive psychology, depression.
 \end{otherlanguage*}
\end{resumo}


% ---
% LISTA DE ILUSTRAÇÕES
% ---
\pdfbookmark[0]{\listfigurename}{lof}
\listoffigures*
\cleardoublepage
% ---


% ---
% LISTA DE QUADROS
% ---
% \pdfbookmark[0]{\listofquadrosname}{loq}
% \listofquadros*
% \cleardoublepage
% ---


% ---
% LISTA DE TABELAS
% ---
% \pdfbookmark[0]{\listtablename}{lot}
% \listoftables*
% \cleardoublepage
% ---


% ---
% LISTA DE ABREVIATURAS E SIGLAS
% ---
% \begin{siglas}
%   \item[ABNT] Associação Brasileira de Normas Técnicas
%   \item[abnTeX] ABsurdas Normas para TeX
% \end{siglas}
% ---


% ---
% LISTA DE SÍMBOLOS
% ---
% \begin{simbolos}
%   \item[$ \Gamma $] Letra grega Gama
%   \item[$ \Lambda $] Lambda
%   \item[$ \zeta $] Letra grega minúscula zeta
%   \item[$ \in $] Pertence
% \end{simbolos}
% ---


% ---
% Sumário
% ---
\pdfbookmark[0]{\contentsname}{toc}
\tableofcontents*
\cleardoublepage
% ---


% ----------------------------------------------------------
% ELEMENTOS TEXTUAIS
% ----------------------------------------------------------
\textual

\include{./01-introducao/main.tex}

\include{./02-machine-learning/main.tex}

\include{./03-exemplo/main.tex}

\include{./04-conclusao/main.tex}

% ----------------------------------------------------------
% ELEMENTOS PÓS-TEXTUAIS
% ----------------------------------------------------------
\postextual
% ----------------------------------------------------------


% ----------------------------------------------------------
% REFERÊNCIAS
% ----------------------------------------------------------
\bibliography{references.bib}


% ----------------------------------------------------------
% GLOSSÁRIO
% ----------------------------------------------------------
%
% Consulte o manual da classe abntex2 para orientações sobre o glossário.
%
%\glossary

% ----------------------------------------------------------
% APÊNDICES
% ----------------------------------------------------------

% ---
% Inicia os apêndices
% ---
% \begin{apendicesenv}

% % Imprime uma página indicando o início dos apêndices
% \partapendices

% % ----------------------------------------------------------
% \chapter{Quisque libero justo}
% % ----------------------------------------------------------

% % ----------------------------------------------------------
% \chapter{Nullam elementum urna}
% % ----------------------------------------------------------

% \end{apendicesenv}
% ---


% ----------------------------------------------------------
% ANEXOS
% ----------------------------------------------------------

% ---
% Inicia os anexos
% ---
% \begin{anexosenv}

% % Imprime uma página indicando o início dos anexos
% \partanexos

% % ---
% \chapter{Morbi ultrices rutrum lorem}
% % ---

% % ---
% \chapter{Cras non urna sed feugiat cum}
% % ---


% % ---
% \chapter{Fusce facilisis lacinia dui}
% % ---

% \end{anexosenv}

%---------------------------------------------------------------------
% INDICE REMISSIVO
%---------------------------------------------------------------------
% \phantompart
% \printindex
%---------------------------------------------------------------------
\includepdf{./05-extras/contracapa.pdf}
\end{document}


\documentclass[
	% Opções da classe memoir
	12pt,				        % tamanho da fonte
	openany,			        % capítulos começam em pág ímpar (insere página vazia caso preciso)
	oneside,			        % para impressão em recto e verso. Oposto a oneside
	a4paper,			        % tamanho do papel. 
	% Opções da classe abntex2
	%chapter=TITLE,		        % títulos de capítulos convertidos em letras maiúsculas
	%section=TITLE,		        % títulos de seções convertidos em letras maiúsculas
	%subsection=TITLE,	        % títulos de subseções convertidos em letras maiúsculas
	%subsubsection=TITLE,       % títulos de subsubseções convertidos em letras maiúsculas
	% Opções do pacote babel
	english,			        % idioma adicional para hifenização
	brazil				        % o último idioma é o principal do documento
	]{abntex2}

% ---
% Pacotes básicos 
% ---
\usepackage{lmodern}			% Usa a fonte Latin Modern			
\usepackage[T1]{fontenc}		% Selecao de codigos de fonte.
\usepackage[utf8]{inputenc}		% Codificacao do documento (conversão automática dos acentos)
\usepackage{indentfirst}		% Indenta o primeiro parágrafo de cada seção.
\usepackage{color}				% Controle das cores
\usepackage{graphicx}			% Inclusão de gráficos
\usepackage{microtype} 			% para melhorias de justificação
% ---
		

% ---
% Pacotes de citações
% ---
% \usepackage[alf]{abntex2cite}						% Citações padrão ABNT
\usepackage[
  alf,
  abnt-nbr10520=1988,                     % Citações em sentence case
  abnt-emphasize=bf
  ]{abntex2cite}


% ---
% Pacotes adicionais
% ---
\usepackage{pdfpages}	% Inclusão de páginas em PDF


% ---
% Informações para CAPA e FOLHA DE ROSTO
% ---
\titulo{APRENDIZAGEM DE MÁQUINA PARA PREDIÇÃO DE DESFECHOS EM SAÚDE MENTAL: \textmd{UMA INTRODUÇÃO PARA PSICÓLOGOS}}
\autor{GUILHERME AUGUSTO DIAS}
\local{Porto Alegre}
\data{2024}
\orientador{Prof. Dr. Wagner de Lara Machado}
% \coorientador{}
% \instituicao{
%   Pontifícia Universidade Católica do Rio Grande do Sul
%   \par
%   Escola de Ciências da Saúde e da Vida
%   \par
%   Curso de Psicologia}
\tipotrabalho{TCC}
\preambulo{Trabalho de Conclusão de Curso apresentado como requisito para a obtenção do grau de Bacharel em Psicologia pela Escola de Ciências da Saúde e da Vida da Pontifícia Universidade Católica do Rio Grande do Sul.}
% ---

% ---
% Configurações de aparência do PDF final
% ---

% Alterando o aspecto da cor azul
\definecolor{blue}{RGB}{41,5,195}

% Informações do PDF
\makeatletter
\hypersetup{
		pdftitle={Aprendizagem de máquina para predição de desfechos em saúde mental: uma introdução para psicólogos}, 
		pdfauthor={Guilherme Augusto Dias},
    	pdfsubject={\imprimirpreambulo},
	    pdfcreator={LaTeX with abnTeX2},
		pdfkeywords={aprendizagem de máquina}{psicoterapia}{psicologia positiva}{depressão}, 
		colorlinks=true,		% true: colored links; false: boxed links
    	linkcolor=black,		% color of internal links
    	citecolor=black,		% color of links to bibliography
    	filecolor=black,		% color of file links
		urlcolor=black,
		bookmarksdepth=4
}
\makeatother
% --- 


% ---
% Posiciona figuras e tabelas no topo da página quando adicionadas sozinhas
% em um página em branco. Ver https://github.com/abntex/abntex2/issues/170
\makeatletter
\setlength{\@fptop}{5pt} % Set distance from top of page to first float
\makeatother
% ---

% ---
% Possibilita criação de Quadros e Lista de quadros.
% Ver https://github.com/abntex/abntex2/issues/176
%
\newcommand{\quadroname}{Quadro}
\newcommand{\listofquadrosname}{Lista de quadros}

\newfloat[chapter]{quadro}{loq}{\quadroname}
\newlistof{listofquadros}{loq}{\listofquadrosname}
\newlistentry{quadro}{loq}{0}

% Configurações para atender às regras da ABNT
\setfloatadjustment{quadro}{\centering}
\counterwithout{quadro}{chapter}
\renewcommand{\cftquadroname}{\quadroname\space} 
\renewcommand*{\cftquadroaftersnum}{\hfill--\hfill}

\setfloatlocations{quadro}{hbtp} % Ver https://github.com/abntex/abntex2/issues/176
% ---


% --- 
% Espaçamentos entre linhas e parágrafos 
% --- 

% O tamanho do parágrafo é dado por:
\setlength{\parindent}{1.3cm}

% Controle do espaçamento entre um parágrafo e outro:
\setlength{\parskip}{0.2cm}  % tente também \onelineskip

% ---
% Compila o índice
% ---
\makeindex
% ---


% ----
% Início do documento
% ----
\begin{document}

% Seleciona o idioma do documento (conforme pacotes do babel)
\selectlanguage{brazil}

% Retira espaço extra obsoleto entre as frases.
\frenchspacing 

% ----------------------------------------------------------
% ELEMENTOS PRÉ-TEXTUAIS
% ----------------------------------------------------------
\pretextual

% ---
% CAPA
% ---
\includepdf{./05-extras/capa.pdf}
% ---


% ---
% FOLHA DE ROSTO
% (o * indica que haverá a ficha bibliográfica)
% ---
\imprimirfolhaderosto
% ---


% ---
% Ficha bibliografica
% ---

% Isto é um exemplo de Ficha Catalográfica, ou ``Dados internacionais de
% catalogação-na-publicação''. Você pode utilizar este modelo como referência. 
% Porém, provavelmente a biblioteca da sua universidade lhe fornecerá um PDF
% com a ficha catalográfica definitiva após a defesa do trabalho. Quando estiver
% com o documento, salve-o como PDF no diretório do seu projeto e substitua todo
% o conteúdo de implementação deste arquivo pelo comando abaixo:
%
% \usepackage{pdfpages}		% necessário para comando \includepdf
% \begin{fichacatalografica}
%     \includepdf{fig_ficha_catalografica.pdf}
% \end{fichacatalografica}

% \begin{fichacatalografica}
% 	\sffamily
% 	\vspace*{\fill}                             % Posição vertical
% 	\begin{center}                              % Minipage Centralizado
% 	\fbox{\begin{minipage}[c][8cm]{13.5cm}		% Largura
% 	\small
% 	\imprimirautor
% 	%Sobrenome, Nome do autor
	
% 	\hspace{0.5cm} \imprimirtitulo  / \imprimirautor. --
% 	\imprimirlocal, \imprimirdata-
	
% 	\hspace{0.5cm} \thelastpage p. : il. (algumas color.) ; 30 cm.\\
	
% 	\hspace{0.5cm} \imprimirorientadorRotulo~\imprimirorientador\\
	
% 	\hspace{0.5cm}
% 	\parbox[t]{\textwidth}{\imprimirtipotrabalho~--~\imprimirinstituicao,
% 	\imprimirdata.}\\
	
% 	\hspace{0.5cm}
% 		1. Palavra-chave1.
% 		2. Palavra-chave2.
% 		2. Palavra-chave3.
% 		I. Orientador.
% 		II. Universidade xxx.
% 		III. Faculdade de xxx.
% 		IV. Título 			
% 	\end{minipage}}
% 	\end{center}
% \end{fichacatalografica}
% ---


% ---
% FOLHA DE APROVAÇÃO
% ---

% Isto é um exemplo de Folha de aprovação, elemento obrigatório da NBR
% 14724/2011 (seção 4.2.1.3). Você pode utilizar este modelo até a aprovação
% do trabalho. Após isso, substitua todo o conteúdo deste arquivo por uma
% imagem da página assinada pela banca com o comando abaixo:
%
% \begin{folhadeaprovacao}
% \includepdf{folhadeaprovacao_final.pdf}
% \end{folhadeaprovacao}
%
\setlength{\ABNTEXsignskip}{2cm}
\setlength{\ABNTEXsignwidth}{10cm}

\begin{folhadeaprovacao}
  \begin{center}
    {\ABNTEXchapterfont\large\imprimirautor}

    \vspace*{1cm}

    \begin{center}
      \ABNTEXchapterfont\bfseries\Large\imprimirtitulo
    \end{center}

    \vspace*{1cm}
    
    \hspace{.45\textwidth}
    \begin{minipage}{.5\textwidth}
        \imprimirpreambulo
    \end{minipage}%
   \end{center}

   \vspace*{1cm}
        
   \begin{center}
   Aprovado em 17 de julho de 2024.

   \vspace*{1cm}

   Banca examinadora:
   \end{center}

   \vspace*{-1cm}

   \assinatura{Prof. Ma. Caroline Santa Maria Rodrigues -- PUCRS}
   \assinatura{Prof. Dra. Clarissa Pinto Pizarro de Freitas -- PUCRS}
   \assinatura{\imprimirorientador -- PUCRS}
\end{folhadeaprovacao}
% ---


% ---
% Agradecimentos
% ---
\begin{agradecimentos}
Agradeço aos meus pais, Ademilson José Dias e Irinês Porfírio Dias, pelo apoio incondicional que me deram ao longo da vida,
incentivando mesmo as ideias questionáveis como aquela de fazer uma segunda graduação.

A minha irmã, Maria Eduarda Dias, cuja trajetória acadêmica é fonte de inspiração.

A minha namorada, Nathalia Fallavena Ceratti, que acompanhou a esta jornada de perto, oferecendo apoio e compreensão que
fizeram toda a diferença.

Aos meus amigos, Guilherme da Silva Petró, Norton Neves e Uilinton Luciano Marques da Silva, que foram compreensíveis nos
finais de semana em que não pude jogar RPG.

Ao professor Wagner de Lara Machado pela orientação e apoio no desenvolvimento deste trabalho.
\end{agradecimentos}
% ---


% ---
% RESUMOS
% ---
% Resumo em português
\setlength{\absparsep}{18pt} % ajusta o espaçamento dos parágrafos do resumo
\begin{resumo}
Desenvolvimentos na área de aprendizagem de máquina representam uma série de oportunidades para diversas áreas de pesquisa e de atuação profissional.
O uso de modelos de aprendizagem de máquina na predição de desfecho para tratamentos em saúde mental tem sido investigado e apresenta resultados promissores.
A aplicação desse tipo de modelo possibilitaria, por exemplo, maior assertividade na personalização de tratamentos psicoterápicos a nível individual.
Apesar dos benefícios potenciais, o tema ainda é pouco discutido nos espaços de pesquisa, formação e atuação em psicologia clínica. Este trabalho tem
por objetivo introduzir o tema de aprendizagem de máquina a psicólogos. Busca-se apresentar, de maneira acessível, os conceitos e terminologia básicos
em aprendizagem de máquina, proporcionando um panorama geral da área. Para complementar a exposição conceitual, é construído um exemplo a partir de dados
de um estudo sobre eficácia de intervenções em psicologia positiva para depressão via internet. Por fim, discutem-se algumas implicações do uso desse tipo
de tecnologia para a pesquisa e para a prática em psicoterapia.

 \textbf{Palavras-chave}: aprendizagem de máquina, psicoterapia, psicologia positiva, depressão.
\end{resumo}

% Resumo em inglês
\begin{resumo}[Abstract]
 \begin{otherlanguage*}{english}
   Advances in machine learning have opened up a multitude of opportunities for various research and professional fields. The use of machine learning models to
   predict treatment outcomes in mental health has been investigated showing promising results. The application of such models would allow greater assertiveness
   in tailoring psychotherapeutic treatments at the individual level. Despite its potential benefits, the topic remains unexplored in the fields of research, training
   and practice in clinical psychology. This work aims to introduce the topic of machine learning to psychologists. It presents the basic machine learning concepts and
   terminology in an accessible manner, providing a comprehensive overview of the field. To complement the conceptual discussion, an example is constructed using data
   from a study on the effectiveness of internet-based positive psychology interventions for depression. Finally, implications of using this type of technology in
   psychotherapy research and practice are discussed.
   \vspace{\onelineskip}
 
   \noindent 
   \textbf{Keywords}: machine learning, psychotherapy, positive psychology, depression.
 \end{otherlanguage*}
\end{resumo}


% ---
% LISTA DE ILUSTRAÇÕES
% ---
\pdfbookmark[0]{\listfigurename}{lof}
\listoffigures*
\cleardoublepage
% ---


% ---
% LISTA DE QUADROS
% ---
% \pdfbookmark[0]{\listofquadrosname}{loq}
% \listofquadros*
% \cleardoublepage
% ---


% ---
% LISTA DE TABELAS
% ---
% \pdfbookmark[0]{\listtablename}{lot}
% \listoftables*
% \cleardoublepage
% ---


% ---
% LISTA DE ABREVIATURAS E SIGLAS
% ---
% \begin{siglas}
%   \item[ABNT] Associação Brasileira de Normas Técnicas
%   \item[abnTeX] ABsurdas Normas para TeX
% \end{siglas}
% ---


% ---
% LISTA DE SÍMBOLOS
% ---
% \begin{simbolos}
%   \item[$ \Gamma $] Letra grega Gama
%   \item[$ \Lambda $] Lambda
%   \item[$ \zeta $] Letra grega minúscula zeta
%   \item[$ \in $] Pertence
% \end{simbolos}
% ---


% ---
% Sumário
% ---
\pdfbookmark[0]{\contentsname}{toc}
\tableofcontents*
\cleardoublepage
% ---


% ----------------------------------------------------------
% ELEMENTOS TEXTUAIS
% ----------------------------------------------------------
\textual

\include{./01-introducao/main.tex}

\include{./02-machine-learning/main.tex}

\include{./03-exemplo/main.tex}

\include{./04-conclusao/main.tex}

% ----------------------------------------------------------
% ELEMENTOS PÓS-TEXTUAIS
% ----------------------------------------------------------
\postextual
% ----------------------------------------------------------


% ----------------------------------------------------------
% REFERÊNCIAS
% ----------------------------------------------------------
\bibliography{references.bib}


% ----------------------------------------------------------
% GLOSSÁRIO
% ----------------------------------------------------------
%
% Consulte o manual da classe abntex2 para orientações sobre o glossário.
%
%\glossary

% ----------------------------------------------------------
% APÊNDICES
% ----------------------------------------------------------

% ---
% Inicia os apêndices
% ---
% \begin{apendicesenv}

% % Imprime uma página indicando o início dos apêndices
% \partapendices

% % ----------------------------------------------------------
% \chapter{Quisque libero justo}
% % ----------------------------------------------------------

% % ----------------------------------------------------------
% \chapter{Nullam elementum urna}
% % ----------------------------------------------------------

% \end{apendicesenv}
% ---


% ----------------------------------------------------------
% ANEXOS
% ----------------------------------------------------------

% ---
% Inicia os anexos
% ---
% \begin{anexosenv}

% % Imprime uma página indicando o início dos anexos
% \partanexos

% % ---
% \chapter{Morbi ultrices rutrum lorem}
% % ---

% % ---
% \chapter{Cras non urna sed feugiat cum}
% % ---


% % ---
% \chapter{Fusce facilisis lacinia dui}
% % ---

% \end{anexosenv}

%---------------------------------------------------------------------
% INDICE REMISSIVO
%---------------------------------------------------------------------
% \phantompart
% \printindex
%---------------------------------------------------------------------
\includepdf{./05-extras/contracapa.pdf}
\end{document}


\documentclass[
	% Opções da classe memoir
	12pt,				        % tamanho da fonte
	openany,			        % capítulos começam em pág ímpar (insere página vazia caso preciso)
	oneside,			        % para impressão em recto e verso. Oposto a oneside
	a4paper,			        % tamanho do papel. 
	% Opções da classe abntex2
	%chapter=TITLE,		        % títulos de capítulos convertidos em letras maiúsculas
	%section=TITLE,		        % títulos de seções convertidos em letras maiúsculas
	%subsection=TITLE,	        % títulos de subseções convertidos em letras maiúsculas
	%subsubsection=TITLE,       % títulos de subsubseções convertidos em letras maiúsculas
	% Opções do pacote babel
	english,			        % idioma adicional para hifenização
	brazil				        % o último idioma é o principal do documento
	]{abntex2}

% ---
% Pacotes básicos 
% ---
\usepackage{lmodern}			% Usa a fonte Latin Modern			
\usepackage[T1]{fontenc}		% Selecao de codigos de fonte.
\usepackage[utf8]{inputenc}		% Codificacao do documento (conversão automática dos acentos)
\usepackage{indentfirst}		% Indenta o primeiro parágrafo de cada seção.
\usepackage{color}				% Controle das cores
\usepackage{graphicx}			% Inclusão de gráficos
\usepackage{microtype} 			% para melhorias de justificação
% ---
		

% ---
% Pacotes de citações
% ---
% \usepackage[alf]{abntex2cite}						% Citações padrão ABNT
\usepackage[
  alf,
  abnt-nbr10520=1988,                     % Citações em sentence case
  abnt-emphasize=bf
  ]{abntex2cite}


% ---
% Pacotes adicionais
% ---
\usepackage{pdfpages}	% Inclusão de páginas em PDF


% ---
% Informações para CAPA e FOLHA DE ROSTO
% ---
\titulo{APRENDIZAGEM DE MÁQUINA PARA PREDIÇÃO DE DESFECHOS EM SAÚDE MENTAL: \textmd{UMA INTRODUÇÃO PARA PSICÓLOGOS}}
\autor{GUILHERME AUGUSTO DIAS}
\local{Porto Alegre}
\data{2024}
\orientador{Prof. Dr. Wagner de Lara Machado}
% \coorientador{}
% \instituicao{
%   Pontifícia Universidade Católica do Rio Grande do Sul
%   \par
%   Escola de Ciências da Saúde e da Vida
%   \par
%   Curso de Psicologia}
\tipotrabalho{TCC}
\preambulo{Trabalho de Conclusão de Curso apresentado como requisito para a obtenção do grau de Bacharel em Psicologia pela Escola de Ciências da Saúde e da Vida da Pontifícia Universidade Católica do Rio Grande do Sul.}
% ---

% ---
% Configurações de aparência do PDF final
% ---

% Alterando o aspecto da cor azul
\definecolor{blue}{RGB}{41,5,195}

% Informações do PDF
\makeatletter
\hypersetup{
		pdftitle={Aprendizagem de máquina para predição de desfechos em saúde mental: uma introdução para psicólogos}, 
		pdfauthor={Guilherme Augusto Dias},
    	pdfsubject={\imprimirpreambulo},
	    pdfcreator={LaTeX with abnTeX2},
		pdfkeywords={aprendizagem de máquina}{psicoterapia}{psicologia positiva}{depressão}, 
		colorlinks=true,		% true: colored links; false: boxed links
    	linkcolor=black,		% color of internal links
    	citecolor=black,		% color of links to bibliography
    	filecolor=black,		% color of file links
		urlcolor=black,
		bookmarksdepth=4
}
\makeatother
% --- 


% ---
% Posiciona figuras e tabelas no topo da página quando adicionadas sozinhas
% em um página em branco. Ver https://github.com/abntex/abntex2/issues/170
\makeatletter
\setlength{\@fptop}{5pt} % Set distance from top of page to first float
\makeatother
% ---

% ---
% Possibilita criação de Quadros e Lista de quadros.
% Ver https://github.com/abntex/abntex2/issues/176
%
\newcommand{\quadroname}{Quadro}
\newcommand{\listofquadrosname}{Lista de quadros}

\newfloat[chapter]{quadro}{loq}{\quadroname}
\newlistof{listofquadros}{loq}{\listofquadrosname}
\newlistentry{quadro}{loq}{0}

% Configurações para atender às regras da ABNT
\setfloatadjustment{quadro}{\centering}
\counterwithout{quadro}{chapter}
\renewcommand{\cftquadroname}{\quadroname\space} 
\renewcommand*{\cftquadroaftersnum}{\hfill--\hfill}

\setfloatlocations{quadro}{hbtp} % Ver https://github.com/abntex/abntex2/issues/176
% ---


% --- 
% Espaçamentos entre linhas e parágrafos 
% --- 

% O tamanho do parágrafo é dado por:
\setlength{\parindent}{1.3cm}

% Controle do espaçamento entre um parágrafo e outro:
\setlength{\parskip}{0.2cm}  % tente também \onelineskip

% ---
% Compila o índice
% ---
\makeindex
% ---


% ----
% Início do documento
% ----
\begin{document}

% Seleciona o idioma do documento (conforme pacotes do babel)
\selectlanguage{brazil}

% Retira espaço extra obsoleto entre as frases.
\frenchspacing 

% ----------------------------------------------------------
% ELEMENTOS PRÉ-TEXTUAIS
% ----------------------------------------------------------
\pretextual

% ---
% CAPA
% ---
\includepdf{./05-extras/capa.pdf}
% ---


% ---
% FOLHA DE ROSTO
% (o * indica que haverá a ficha bibliográfica)
% ---
\imprimirfolhaderosto
% ---


% ---
% Ficha bibliografica
% ---

% Isto é um exemplo de Ficha Catalográfica, ou ``Dados internacionais de
% catalogação-na-publicação''. Você pode utilizar este modelo como referência. 
% Porém, provavelmente a biblioteca da sua universidade lhe fornecerá um PDF
% com a ficha catalográfica definitiva após a defesa do trabalho. Quando estiver
% com o documento, salve-o como PDF no diretório do seu projeto e substitua todo
% o conteúdo de implementação deste arquivo pelo comando abaixo:
%
% \usepackage{pdfpages}		% necessário para comando \includepdf
% \begin{fichacatalografica}
%     \includepdf{fig_ficha_catalografica.pdf}
% \end{fichacatalografica}

% \begin{fichacatalografica}
% 	\sffamily
% 	\vspace*{\fill}                             % Posição vertical
% 	\begin{center}                              % Minipage Centralizado
% 	\fbox{\begin{minipage}[c][8cm]{13.5cm}		% Largura
% 	\small
% 	\imprimirautor
% 	%Sobrenome, Nome do autor
	
% 	\hspace{0.5cm} \imprimirtitulo  / \imprimirautor. --
% 	\imprimirlocal, \imprimirdata-
	
% 	\hspace{0.5cm} \thelastpage p. : il. (algumas color.) ; 30 cm.\\
	
% 	\hspace{0.5cm} \imprimirorientadorRotulo~\imprimirorientador\\
	
% 	\hspace{0.5cm}
% 	\parbox[t]{\textwidth}{\imprimirtipotrabalho~--~\imprimirinstituicao,
% 	\imprimirdata.}\\
	
% 	\hspace{0.5cm}
% 		1. Palavra-chave1.
% 		2. Palavra-chave2.
% 		2. Palavra-chave3.
% 		I. Orientador.
% 		II. Universidade xxx.
% 		III. Faculdade de xxx.
% 		IV. Título 			
% 	\end{minipage}}
% 	\end{center}
% \end{fichacatalografica}
% ---


% ---
% FOLHA DE APROVAÇÃO
% ---

% Isto é um exemplo de Folha de aprovação, elemento obrigatório da NBR
% 14724/2011 (seção 4.2.1.3). Você pode utilizar este modelo até a aprovação
% do trabalho. Após isso, substitua todo o conteúdo deste arquivo por uma
% imagem da página assinada pela banca com o comando abaixo:
%
% \begin{folhadeaprovacao}
% \includepdf{folhadeaprovacao_final.pdf}
% \end{folhadeaprovacao}
%
\setlength{\ABNTEXsignskip}{2cm}
\setlength{\ABNTEXsignwidth}{10cm}

\begin{folhadeaprovacao}
  \begin{center}
    {\ABNTEXchapterfont\large\imprimirautor}

    \vspace*{1cm}

    \begin{center}
      \ABNTEXchapterfont\bfseries\Large\imprimirtitulo
    \end{center}

    \vspace*{1cm}
    
    \hspace{.45\textwidth}
    \begin{minipage}{.5\textwidth}
        \imprimirpreambulo
    \end{minipage}%
   \end{center}

   \vspace*{1cm}
        
   \begin{center}
   Aprovado em 17 de julho de 2024.

   \vspace*{1cm}

   Banca examinadora:
   \end{center}

   \vspace*{-1cm}

   \assinatura{Prof. Ma. Caroline Santa Maria Rodrigues -- PUCRS}
   \assinatura{Prof. Dra. Clarissa Pinto Pizarro de Freitas -- PUCRS}
   \assinatura{\imprimirorientador -- PUCRS}
\end{folhadeaprovacao}
% ---


% ---
% Agradecimentos
% ---
\begin{agradecimentos}
Agradeço aos meus pais, Ademilson José Dias e Irinês Porfírio Dias, pelo apoio incondicional que me deram ao longo da vida,
incentivando mesmo as ideias questionáveis como aquela de fazer uma segunda graduação.

A minha irmã, Maria Eduarda Dias, cuja trajetória acadêmica é fonte de inspiração.

A minha namorada, Nathalia Fallavena Ceratti, que acompanhou a esta jornada de perto, oferecendo apoio e compreensão que
fizeram toda a diferença.

Aos meus amigos, Guilherme da Silva Petró, Norton Neves e Uilinton Luciano Marques da Silva, que foram compreensíveis nos
finais de semana em que não pude jogar RPG.

Ao professor Wagner de Lara Machado pela orientação e apoio no desenvolvimento deste trabalho.
\end{agradecimentos}
% ---


% ---
% RESUMOS
% ---
% Resumo em português
\setlength{\absparsep}{18pt} % ajusta o espaçamento dos parágrafos do resumo
\begin{resumo}
Desenvolvimentos na área de aprendizagem de máquina representam uma série de oportunidades para diversas áreas de pesquisa e de atuação profissional.
O uso de modelos de aprendizagem de máquina na predição de desfecho para tratamentos em saúde mental tem sido investigado e apresenta resultados promissores.
A aplicação desse tipo de modelo possibilitaria, por exemplo, maior assertividade na personalização de tratamentos psicoterápicos a nível individual.
Apesar dos benefícios potenciais, o tema ainda é pouco discutido nos espaços de pesquisa, formação e atuação em psicologia clínica. Este trabalho tem
por objetivo introduzir o tema de aprendizagem de máquina a psicólogos. Busca-se apresentar, de maneira acessível, os conceitos e terminologia básicos
em aprendizagem de máquina, proporcionando um panorama geral da área. Para complementar a exposição conceitual, é construído um exemplo a partir de dados
de um estudo sobre eficácia de intervenções em psicologia positiva para depressão via internet. Por fim, discutem-se algumas implicações do uso desse tipo
de tecnologia para a pesquisa e para a prática em psicoterapia.

 \textbf{Palavras-chave}: aprendizagem de máquina, psicoterapia, psicologia positiva, depressão.
\end{resumo}

% Resumo em inglês
\begin{resumo}[Abstract]
 \begin{otherlanguage*}{english}
   Advances in machine learning have opened up a multitude of opportunities for various research and professional fields. The use of machine learning models to
   predict treatment outcomes in mental health has been investigated showing promising results. The application of such models would allow greater assertiveness
   in tailoring psychotherapeutic treatments at the individual level. Despite its potential benefits, the topic remains unexplored in the fields of research, training
   and practice in clinical psychology. This work aims to introduce the topic of machine learning to psychologists. It presents the basic machine learning concepts and
   terminology in an accessible manner, providing a comprehensive overview of the field. To complement the conceptual discussion, an example is constructed using data
   from a study on the effectiveness of internet-based positive psychology interventions for depression. Finally, implications of using this type of technology in
   psychotherapy research and practice are discussed.
   \vspace{\onelineskip}
 
   \noindent 
   \textbf{Keywords}: machine learning, psychotherapy, positive psychology, depression.
 \end{otherlanguage*}
\end{resumo}


% ---
% LISTA DE ILUSTRAÇÕES
% ---
\pdfbookmark[0]{\listfigurename}{lof}
\listoffigures*
\cleardoublepage
% ---


% ---
% LISTA DE QUADROS
% ---
% \pdfbookmark[0]{\listofquadrosname}{loq}
% \listofquadros*
% \cleardoublepage
% ---


% ---
% LISTA DE TABELAS
% ---
% \pdfbookmark[0]{\listtablename}{lot}
% \listoftables*
% \cleardoublepage
% ---


% ---
% LISTA DE ABREVIATURAS E SIGLAS
% ---
% \begin{siglas}
%   \item[ABNT] Associação Brasileira de Normas Técnicas
%   \item[abnTeX] ABsurdas Normas para TeX
% \end{siglas}
% ---


% ---
% LISTA DE SÍMBOLOS
% ---
% \begin{simbolos}
%   \item[$ \Gamma $] Letra grega Gama
%   \item[$ \Lambda $] Lambda
%   \item[$ \zeta $] Letra grega minúscula zeta
%   \item[$ \in $] Pertence
% \end{simbolos}
% ---


% ---
% Sumário
% ---
\pdfbookmark[0]{\contentsname}{toc}
\tableofcontents*
\cleardoublepage
% ---


% ----------------------------------------------------------
% ELEMENTOS TEXTUAIS
% ----------------------------------------------------------
\textual

\include{./01-introducao/main.tex}

\include{./02-machine-learning/main.tex}

\include{./03-exemplo/main.tex}

\include{./04-conclusao/main.tex}

% ----------------------------------------------------------
% ELEMENTOS PÓS-TEXTUAIS
% ----------------------------------------------------------
\postextual
% ----------------------------------------------------------


% ----------------------------------------------------------
% REFERÊNCIAS
% ----------------------------------------------------------
\bibliography{references.bib}


% ----------------------------------------------------------
% GLOSSÁRIO
% ----------------------------------------------------------
%
% Consulte o manual da classe abntex2 para orientações sobre o glossário.
%
%\glossary

% ----------------------------------------------------------
% APÊNDICES
% ----------------------------------------------------------

% ---
% Inicia os apêndices
% ---
% \begin{apendicesenv}

% % Imprime uma página indicando o início dos apêndices
% \partapendices

% % ----------------------------------------------------------
% \chapter{Quisque libero justo}
% % ----------------------------------------------------------

% % ----------------------------------------------------------
% \chapter{Nullam elementum urna}
% % ----------------------------------------------------------

% \end{apendicesenv}
% ---


% ----------------------------------------------------------
% ANEXOS
% ----------------------------------------------------------

% ---
% Inicia os anexos
% ---
% \begin{anexosenv}

% % Imprime uma página indicando o início dos anexos
% \partanexos

% % ---
% \chapter{Morbi ultrices rutrum lorem}
% % ---

% % ---
% \chapter{Cras non urna sed feugiat cum}
% % ---


% % ---
% \chapter{Fusce facilisis lacinia dui}
% % ---

% \end{anexosenv}

%---------------------------------------------------------------------
% INDICE REMISSIVO
%---------------------------------------------------------------------
% \phantompart
% \printindex
%---------------------------------------------------------------------
\includepdf{./05-extras/contracapa.pdf}
\end{document}


\documentclass[
	% Opções da classe memoir
	12pt,				        % tamanho da fonte
	openany,			        % capítulos começam em pág ímpar (insere página vazia caso preciso)
	oneside,			        % para impressão em recto e verso. Oposto a oneside
	a4paper,			        % tamanho do papel. 
	% Opções da classe abntex2
	%chapter=TITLE,		        % títulos de capítulos convertidos em letras maiúsculas
	%section=TITLE,		        % títulos de seções convertidos em letras maiúsculas
	%subsection=TITLE,	        % títulos de subseções convertidos em letras maiúsculas
	%subsubsection=TITLE,       % títulos de subsubseções convertidos em letras maiúsculas
	% Opções do pacote babel
	english,			        % idioma adicional para hifenização
	brazil				        % o último idioma é o principal do documento
	]{abntex2}

% ---
% Pacotes básicos 
% ---
\usepackage{lmodern}			% Usa a fonte Latin Modern			
\usepackage[T1]{fontenc}		% Selecao de codigos de fonte.
\usepackage[utf8]{inputenc}		% Codificacao do documento (conversão automática dos acentos)
\usepackage{indentfirst}		% Indenta o primeiro parágrafo de cada seção.
\usepackage{color}				% Controle das cores
\usepackage{graphicx}			% Inclusão de gráficos
\usepackage{microtype} 			% para melhorias de justificação
% ---
		

% ---
% Pacotes de citações
% ---
% \usepackage[alf]{abntex2cite}						% Citações padrão ABNT
\usepackage[
  alf,
  abnt-nbr10520=1988,                     % Citações em sentence case
  abnt-emphasize=bf
  ]{abntex2cite}


% ---
% Pacotes adicionais
% ---
\usepackage{pdfpages}	% Inclusão de páginas em PDF


% ---
% Informações para CAPA e FOLHA DE ROSTO
% ---
\titulo{APRENDIZAGEM DE MÁQUINA PARA PREDIÇÃO DE DESFECHOS EM SAÚDE MENTAL: \textmd{UMA INTRODUÇÃO PARA PSICÓLOGOS}}
\autor{GUILHERME AUGUSTO DIAS}
\local{Porto Alegre}
\data{2024}
\orientador{Prof. Dr. Wagner de Lara Machado}
% \coorientador{}
% \instituicao{
%   Pontifícia Universidade Católica do Rio Grande do Sul
%   \par
%   Escola de Ciências da Saúde e da Vida
%   \par
%   Curso de Psicologia}
\tipotrabalho{TCC}
\preambulo{Trabalho de Conclusão de Curso apresentado como requisito para a obtenção do grau de Bacharel em Psicologia pela Escola de Ciências da Saúde e da Vida da Pontifícia Universidade Católica do Rio Grande do Sul.}
% ---

% ---
% Configurações de aparência do PDF final
% ---

% Alterando o aspecto da cor azul
\definecolor{blue}{RGB}{41,5,195}

% Informações do PDF
\makeatletter
\hypersetup{
		pdftitle={Aprendizagem de máquina para predição de desfechos em saúde mental: uma introdução para psicólogos}, 
		pdfauthor={Guilherme Augusto Dias},
    	pdfsubject={\imprimirpreambulo},
	    pdfcreator={LaTeX with abnTeX2},
		pdfkeywords={aprendizagem de máquina}{psicoterapia}{psicologia positiva}{depressão}, 
		colorlinks=true,		% true: colored links; false: boxed links
    	linkcolor=black,		% color of internal links
    	citecolor=black,		% color of links to bibliography
    	filecolor=black,		% color of file links
		urlcolor=black,
		bookmarksdepth=4
}
\makeatother
% --- 


% ---
% Posiciona figuras e tabelas no topo da página quando adicionadas sozinhas
% em um página em branco. Ver https://github.com/abntex/abntex2/issues/170
\makeatletter
\setlength{\@fptop}{5pt} % Set distance from top of page to first float
\makeatother
% ---

% ---
% Possibilita criação de Quadros e Lista de quadros.
% Ver https://github.com/abntex/abntex2/issues/176
%
\newcommand{\quadroname}{Quadro}
\newcommand{\listofquadrosname}{Lista de quadros}

\newfloat[chapter]{quadro}{loq}{\quadroname}
\newlistof{listofquadros}{loq}{\listofquadrosname}
\newlistentry{quadro}{loq}{0}

% Configurações para atender às regras da ABNT
\setfloatadjustment{quadro}{\centering}
\counterwithout{quadro}{chapter}
\renewcommand{\cftquadroname}{\quadroname\space} 
\renewcommand*{\cftquadroaftersnum}{\hfill--\hfill}

\setfloatlocations{quadro}{hbtp} % Ver https://github.com/abntex/abntex2/issues/176
% ---


% --- 
% Espaçamentos entre linhas e parágrafos 
% --- 

% O tamanho do parágrafo é dado por:
\setlength{\parindent}{1.3cm}

% Controle do espaçamento entre um parágrafo e outro:
\setlength{\parskip}{0.2cm}  % tente também \onelineskip

% ---
% Compila o índice
% ---
\makeindex
% ---


% ----
% Início do documento
% ----
\begin{document}

% Seleciona o idioma do documento (conforme pacotes do babel)
\selectlanguage{brazil}

% Retira espaço extra obsoleto entre as frases.
\frenchspacing 

% ----------------------------------------------------------
% ELEMENTOS PRÉ-TEXTUAIS
% ----------------------------------------------------------
\pretextual

% ---
% CAPA
% ---
\includepdf{./05-extras/capa.pdf}
% ---


% ---
% FOLHA DE ROSTO
% (o * indica que haverá a ficha bibliográfica)
% ---
\imprimirfolhaderosto
% ---


% ---
% Ficha bibliografica
% ---

% Isto é um exemplo de Ficha Catalográfica, ou ``Dados internacionais de
% catalogação-na-publicação''. Você pode utilizar este modelo como referência. 
% Porém, provavelmente a biblioteca da sua universidade lhe fornecerá um PDF
% com a ficha catalográfica definitiva após a defesa do trabalho. Quando estiver
% com o documento, salve-o como PDF no diretório do seu projeto e substitua todo
% o conteúdo de implementação deste arquivo pelo comando abaixo:
%
% \usepackage{pdfpages}		% necessário para comando \includepdf
% \begin{fichacatalografica}
%     \includepdf{fig_ficha_catalografica.pdf}
% \end{fichacatalografica}

% \begin{fichacatalografica}
% 	\sffamily
% 	\vspace*{\fill}                             % Posição vertical
% 	\begin{center}                              % Minipage Centralizado
% 	\fbox{\begin{minipage}[c][8cm]{13.5cm}		% Largura
% 	\small
% 	\imprimirautor
% 	%Sobrenome, Nome do autor
	
% 	\hspace{0.5cm} \imprimirtitulo  / \imprimirautor. --
% 	\imprimirlocal, \imprimirdata-
	
% 	\hspace{0.5cm} \thelastpage p. : il. (algumas color.) ; 30 cm.\\
	
% 	\hspace{0.5cm} \imprimirorientadorRotulo~\imprimirorientador\\
	
% 	\hspace{0.5cm}
% 	\parbox[t]{\textwidth}{\imprimirtipotrabalho~--~\imprimirinstituicao,
% 	\imprimirdata.}\\
	
% 	\hspace{0.5cm}
% 		1. Palavra-chave1.
% 		2. Palavra-chave2.
% 		2. Palavra-chave3.
% 		I. Orientador.
% 		II. Universidade xxx.
% 		III. Faculdade de xxx.
% 		IV. Título 			
% 	\end{minipage}}
% 	\end{center}
% \end{fichacatalografica}
% ---


% ---
% FOLHA DE APROVAÇÃO
% ---

% Isto é um exemplo de Folha de aprovação, elemento obrigatório da NBR
% 14724/2011 (seção 4.2.1.3). Você pode utilizar este modelo até a aprovação
% do trabalho. Após isso, substitua todo o conteúdo deste arquivo por uma
% imagem da página assinada pela banca com o comando abaixo:
%
% \begin{folhadeaprovacao}
% \includepdf{folhadeaprovacao_final.pdf}
% \end{folhadeaprovacao}
%
\setlength{\ABNTEXsignskip}{2cm}
\setlength{\ABNTEXsignwidth}{10cm}

\begin{folhadeaprovacao}
  \begin{center}
    {\ABNTEXchapterfont\large\imprimirautor}

    \vspace*{1cm}

    \begin{center}
      \ABNTEXchapterfont\bfseries\Large\imprimirtitulo
    \end{center}

    \vspace*{1cm}
    
    \hspace{.45\textwidth}
    \begin{minipage}{.5\textwidth}
        \imprimirpreambulo
    \end{minipage}%
   \end{center}

   \vspace*{1cm}
        
   \begin{center}
   Aprovado em 17 de julho de 2024.

   \vspace*{1cm}

   Banca examinadora:
   \end{center}

   \vspace*{-1cm}

   \assinatura{Prof. Ma. Caroline Santa Maria Rodrigues -- PUCRS}
   \assinatura{Prof. Dra. Clarissa Pinto Pizarro de Freitas -- PUCRS}
   \assinatura{\imprimirorientador -- PUCRS}
\end{folhadeaprovacao}
% ---


% ---
% Agradecimentos
% ---
\begin{agradecimentos}
Agradeço aos meus pais, Ademilson José Dias e Irinês Porfírio Dias, pelo apoio incondicional que me deram ao longo da vida,
incentivando mesmo as ideias questionáveis como aquela de fazer uma segunda graduação.

A minha irmã, Maria Eduarda Dias, cuja trajetória acadêmica é fonte de inspiração.

A minha namorada, Nathalia Fallavena Ceratti, que acompanhou a esta jornada de perto, oferecendo apoio e compreensão que
fizeram toda a diferença.

Aos meus amigos, Guilherme da Silva Petró, Norton Neves e Uilinton Luciano Marques da Silva, que foram compreensíveis nos
finais de semana em que não pude jogar RPG.

Ao professor Wagner de Lara Machado pela orientação e apoio no desenvolvimento deste trabalho.
\end{agradecimentos}
% ---


% ---
% RESUMOS
% ---
% Resumo em português
\setlength{\absparsep}{18pt} % ajusta o espaçamento dos parágrafos do resumo
\begin{resumo}
Desenvolvimentos na área de aprendizagem de máquina representam uma série de oportunidades para diversas áreas de pesquisa e de atuação profissional.
O uso de modelos de aprendizagem de máquina na predição de desfecho para tratamentos em saúde mental tem sido investigado e apresenta resultados promissores.
A aplicação desse tipo de modelo possibilitaria, por exemplo, maior assertividade na personalização de tratamentos psicoterápicos a nível individual.
Apesar dos benefícios potenciais, o tema ainda é pouco discutido nos espaços de pesquisa, formação e atuação em psicologia clínica. Este trabalho tem
por objetivo introduzir o tema de aprendizagem de máquina a psicólogos. Busca-se apresentar, de maneira acessível, os conceitos e terminologia básicos
em aprendizagem de máquina, proporcionando um panorama geral da área. Para complementar a exposição conceitual, é construído um exemplo a partir de dados
de um estudo sobre eficácia de intervenções em psicologia positiva para depressão via internet. Por fim, discutem-se algumas implicações do uso desse tipo
de tecnologia para a pesquisa e para a prática em psicoterapia.

 \textbf{Palavras-chave}: aprendizagem de máquina, psicoterapia, psicologia positiva, depressão.
\end{resumo}

% Resumo em inglês
\begin{resumo}[Abstract]
 \begin{otherlanguage*}{english}
   Advances in machine learning have opened up a multitude of opportunities for various research and professional fields. The use of machine learning models to
   predict treatment outcomes in mental health has been investigated showing promising results. The application of such models would allow greater assertiveness
   in tailoring psychotherapeutic treatments at the individual level. Despite its potential benefits, the topic remains unexplored in the fields of research, training
   and practice in clinical psychology. This work aims to introduce the topic of machine learning to psychologists. It presents the basic machine learning concepts and
   terminology in an accessible manner, providing a comprehensive overview of the field. To complement the conceptual discussion, an example is constructed using data
   from a study on the effectiveness of internet-based positive psychology interventions for depression. Finally, implications of using this type of technology in
   psychotherapy research and practice are discussed.
   \vspace{\onelineskip}
 
   \noindent 
   \textbf{Keywords}: machine learning, psychotherapy, positive psychology, depression.
 \end{otherlanguage*}
\end{resumo}


% ---
% LISTA DE ILUSTRAÇÕES
% ---
\pdfbookmark[0]{\listfigurename}{lof}
\listoffigures*
\cleardoublepage
% ---


% ---
% LISTA DE QUADROS
% ---
% \pdfbookmark[0]{\listofquadrosname}{loq}
% \listofquadros*
% \cleardoublepage
% ---


% ---
% LISTA DE TABELAS
% ---
% \pdfbookmark[0]{\listtablename}{lot}
% \listoftables*
% \cleardoublepage
% ---


% ---
% LISTA DE ABREVIATURAS E SIGLAS
% ---
% \begin{siglas}
%   \item[ABNT] Associação Brasileira de Normas Técnicas
%   \item[abnTeX] ABsurdas Normas para TeX
% \end{siglas}
% ---


% ---
% LISTA DE SÍMBOLOS
% ---
% \begin{simbolos}
%   \item[$ \Gamma $] Letra grega Gama
%   \item[$ \Lambda $] Lambda
%   \item[$ \zeta $] Letra grega minúscula zeta
%   \item[$ \in $] Pertence
% \end{simbolos}
% ---


% ---
% Sumário
% ---
\pdfbookmark[0]{\contentsname}{toc}
\tableofcontents*
\cleardoublepage
% ---


% ----------------------------------------------------------
% ELEMENTOS TEXTUAIS
% ----------------------------------------------------------
\textual

\include{./01-introducao/main.tex}

\include{./02-machine-learning/main.tex}

\include{./03-exemplo/main.tex}

\include{./04-conclusao/main.tex}

% ----------------------------------------------------------
% ELEMENTOS PÓS-TEXTUAIS
% ----------------------------------------------------------
\postextual
% ----------------------------------------------------------


% ----------------------------------------------------------
% REFERÊNCIAS
% ----------------------------------------------------------
\bibliography{references.bib}


% ----------------------------------------------------------
% GLOSSÁRIO
% ----------------------------------------------------------
%
% Consulte o manual da classe abntex2 para orientações sobre o glossário.
%
%\glossary

% ----------------------------------------------------------
% APÊNDICES
% ----------------------------------------------------------

% ---
% Inicia os apêndices
% ---
% \begin{apendicesenv}

% % Imprime uma página indicando o início dos apêndices
% \partapendices

% % ----------------------------------------------------------
% \chapter{Quisque libero justo}
% % ----------------------------------------------------------

% % ----------------------------------------------------------
% \chapter{Nullam elementum urna}
% % ----------------------------------------------------------

% \end{apendicesenv}
% ---


% ----------------------------------------------------------
% ANEXOS
% ----------------------------------------------------------

% ---
% Inicia os anexos
% ---
% \begin{anexosenv}

% % Imprime uma página indicando o início dos anexos
% \partanexos

% % ---
% \chapter{Morbi ultrices rutrum lorem}
% % ---

% % ---
% \chapter{Cras non urna sed feugiat cum}
% % ---


% % ---
% \chapter{Fusce facilisis lacinia dui}
% % ---

% \end{anexosenv}

%---------------------------------------------------------------------
% INDICE REMISSIVO
%---------------------------------------------------------------------
% \phantompart
% \printindex
%---------------------------------------------------------------------
\includepdf{./05-extras/contracapa.pdf}
\end{document}


% ----------------------------------------------------------
% ELEMENTOS PÓS-TEXTUAIS
% ----------------------------------------------------------
\postextual
% ----------------------------------------------------------


% ----------------------------------------------------------
% REFERÊNCIAS
% ----------------------------------------------------------
\bibliography{references.bib}


% ----------------------------------------------------------
% GLOSSÁRIO
% ----------------------------------------------------------
%
% Consulte o manual da classe abntex2 para orientações sobre o glossário.
%
%\glossary

% ----------------------------------------------------------
% APÊNDICES
% ----------------------------------------------------------

% ---
% Inicia os apêndices
% ---
% \begin{apendicesenv}

% % Imprime uma página indicando o início dos apêndices
% \partapendices

% % ----------------------------------------------------------
% \chapter{Quisque libero justo}
% % ----------------------------------------------------------

% % ----------------------------------------------------------
% \chapter{Nullam elementum urna}
% % ----------------------------------------------------------

% \end{apendicesenv}
% ---


% ----------------------------------------------------------
% ANEXOS
% ----------------------------------------------------------

% ---
% Inicia os anexos
% ---
% \begin{anexosenv}

% % Imprime uma página indicando o início dos anexos
% \partanexos

% % ---
% \chapter{Morbi ultrices rutrum lorem}
% % ---

% % ---
% \chapter{Cras non urna sed feugiat cum}
% % ---


% % ---
% \chapter{Fusce facilisis lacinia dui}
% % ---

% \end{anexosenv}

%---------------------------------------------------------------------
% INDICE REMISSIVO
%---------------------------------------------------------------------
% \phantompart
% \printindex
%---------------------------------------------------------------------
\includepdf{./05-extras/contracapa.pdf}
\end{document}


\documentclass[
	% Opções da classe memoir
	12pt,				        % tamanho da fonte
	openany,			        % capítulos começam em pág ímpar (insere página vazia caso preciso)
	oneside,			        % para impressão em recto e verso. Oposto a oneside
	a4paper,			        % tamanho do papel. 
	% Opções da classe abntex2
	%chapter=TITLE,		        % títulos de capítulos convertidos em letras maiúsculas
	%section=TITLE,		        % títulos de seções convertidos em letras maiúsculas
	%subsection=TITLE,	        % títulos de subseções convertidos em letras maiúsculas
	%subsubsection=TITLE,       % títulos de subsubseções convertidos em letras maiúsculas
	% Opções do pacote babel
	english,			        % idioma adicional para hifenização
	brazil				        % o último idioma é o principal do documento
	]{abntex2}

% ---
% Pacotes básicos 
% ---
\usepackage{lmodern}			% Usa a fonte Latin Modern			
\usepackage[T1]{fontenc}		% Selecao de codigos de fonte.
\usepackage[utf8]{inputenc}		% Codificacao do documento (conversão automática dos acentos)
\usepackage{indentfirst}		% Indenta o primeiro parágrafo de cada seção.
\usepackage{color}				% Controle das cores
\usepackage{graphicx}			% Inclusão de gráficos
\usepackage{microtype} 			% para melhorias de justificação
% ---
		

% ---
% Pacotes de citações
% ---
% \usepackage[alf]{abntex2cite}						% Citações padrão ABNT
\usepackage[
  alf,
  abnt-nbr10520=1988,                     % Citações em sentence case
  abnt-emphasize=bf
  ]{abntex2cite}


% ---
% Pacotes adicionais
% ---
\usepackage{pdfpages}	% Inclusão de páginas em PDF


% ---
% Informações para CAPA e FOLHA DE ROSTO
% ---
\titulo{APRENDIZAGEM DE MÁQUINA PARA PREDIÇÃO DE DESFECHOS EM SAÚDE MENTAL: \textmd{UMA INTRODUÇÃO PARA PSICÓLOGOS}}
\autor{GUILHERME AUGUSTO DIAS}
\local{Porto Alegre}
\data{2024}
\orientador{Prof. Dr. Wagner de Lara Machado}
% \coorientador{}
% \instituicao{
%   Pontifícia Universidade Católica do Rio Grande do Sul
%   \par
%   Escola de Ciências da Saúde e da Vida
%   \par
%   Curso de Psicologia}
\tipotrabalho{TCC}
\preambulo{Trabalho de Conclusão de Curso apresentado como requisito para a obtenção do grau de Bacharel em Psicologia pela Escola de Ciências da Saúde e da Vida da Pontifícia Universidade Católica do Rio Grande do Sul.}
% ---

% ---
% Configurações de aparência do PDF final
% ---

% Alterando o aspecto da cor azul
\definecolor{blue}{RGB}{41,5,195}

% Informações do PDF
\makeatletter
\hypersetup{
		pdftitle={Aprendizagem de máquina para predição de desfechos em saúde mental: uma introdução para psicólogos}, 
		pdfauthor={Guilherme Augusto Dias},
    	pdfsubject={\imprimirpreambulo},
	    pdfcreator={LaTeX with abnTeX2},
		pdfkeywords={aprendizagem de máquina}{psicoterapia}{psicologia positiva}{depressão}, 
		colorlinks=true,		% true: colored links; false: boxed links
    	linkcolor=black,		% color of internal links
    	citecolor=black,		% color of links to bibliography
    	filecolor=black,		% color of file links
		urlcolor=black,
		bookmarksdepth=4
}
\makeatother
% --- 


% ---
% Posiciona figuras e tabelas no topo da página quando adicionadas sozinhas
% em um página em branco. Ver https://github.com/abntex/abntex2/issues/170
\makeatletter
\setlength{\@fptop}{5pt} % Set distance from top of page to first float
\makeatother
% ---

% ---
% Possibilita criação de Quadros e Lista de quadros.
% Ver https://github.com/abntex/abntex2/issues/176
%
\newcommand{\quadroname}{Quadro}
\newcommand{\listofquadrosname}{Lista de quadros}

\newfloat[chapter]{quadro}{loq}{\quadroname}
\newlistof{listofquadros}{loq}{\listofquadrosname}
\newlistentry{quadro}{loq}{0}

% Configurações para atender às regras da ABNT
\setfloatadjustment{quadro}{\centering}
\counterwithout{quadro}{chapter}
\renewcommand{\cftquadroname}{\quadroname\space} 
\renewcommand*{\cftquadroaftersnum}{\hfill--\hfill}

\setfloatlocations{quadro}{hbtp} % Ver https://github.com/abntex/abntex2/issues/176
% ---


% --- 
% Espaçamentos entre linhas e parágrafos 
% --- 

% O tamanho do parágrafo é dado por:
\setlength{\parindent}{1.3cm}

% Controle do espaçamento entre um parágrafo e outro:
\setlength{\parskip}{0.2cm}  % tente também \onelineskip

% ---
% Compila o índice
% ---
\makeindex
% ---


% ----
% Início do documento
% ----
\begin{document}

% Seleciona o idioma do documento (conforme pacotes do babel)
\selectlanguage{brazil}

% Retira espaço extra obsoleto entre as frases.
\frenchspacing 

% ----------------------------------------------------------
% ELEMENTOS PRÉ-TEXTUAIS
% ----------------------------------------------------------
\pretextual

% ---
% CAPA
% ---
\includepdf{./05-extras/capa.pdf}
% ---


% ---
% FOLHA DE ROSTO
% (o * indica que haverá a ficha bibliográfica)
% ---
\imprimirfolhaderosto
% ---


% ---
% Ficha bibliografica
% ---

% Isto é um exemplo de Ficha Catalográfica, ou ``Dados internacionais de
% catalogação-na-publicação''. Você pode utilizar este modelo como referência. 
% Porém, provavelmente a biblioteca da sua universidade lhe fornecerá um PDF
% com a ficha catalográfica definitiva após a defesa do trabalho. Quando estiver
% com o documento, salve-o como PDF no diretório do seu projeto e substitua todo
% o conteúdo de implementação deste arquivo pelo comando abaixo:
%
% \usepackage{pdfpages}		% necessário para comando \includepdf
% \begin{fichacatalografica}
%     \includepdf{fig_ficha_catalografica.pdf}
% \end{fichacatalografica}

% \begin{fichacatalografica}
% 	\sffamily
% 	\vspace*{\fill}                             % Posição vertical
% 	\begin{center}                              % Minipage Centralizado
% 	\fbox{\begin{minipage}[c][8cm]{13.5cm}		% Largura
% 	\small
% 	\imprimirautor
% 	%Sobrenome, Nome do autor
	
% 	\hspace{0.5cm} \imprimirtitulo  / \imprimirautor. --
% 	\imprimirlocal, \imprimirdata-
	
% 	\hspace{0.5cm} \thelastpage p. : il. (algumas color.) ; 30 cm.\\
	
% 	\hspace{0.5cm} \imprimirorientadorRotulo~\imprimirorientador\\
	
% 	\hspace{0.5cm}
% 	\parbox[t]{\textwidth}{\imprimirtipotrabalho~--~\imprimirinstituicao,
% 	\imprimirdata.}\\
	
% 	\hspace{0.5cm}
% 		1. Palavra-chave1.
% 		2. Palavra-chave2.
% 		2. Palavra-chave3.
% 		I. Orientador.
% 		II. Universidade xxx.
% 		III. Faculdade de xxx.
% 		IV. Título 			
% 	\end{minipage}}
% 	\end{center}
% \end{fichacatalografica}
% ---


% ---
% FOLHA DE APROVAÇÃO
% ---

% Isto é um exemplo de Folha de aprovação, elemento obrigatório da NBR
% 14724/2011 (seção 4.2.1.3). Você pode utilizar este modelo até a aprovação
% do trabalho. Após isso, substitua todo o conteúdo deste arquivo por uma
% imagem da página assinada pela banca com o comando abaixo:
%
% \begin{folhadeaprovacao}
% \includepdf{folhadeaprovacao_final.pdf}
% \end{folhadeaprovacao}
%
\setlength{\ABNTEXsignskip}{2cm}
\setlength{\ABNTEXsignwidth}{10cm}

\begin{folhadeaprovacao}
  \begin{center}
    {\ABNTEXchapterfont\large\imprimirautor}

    \vspace*{1cm}

    \begin{center}
      \ABNTEXchapterfont\bfseries\Large\imprimirtitulo
    \end{center}

    \vspace*{1cm}
    
    \hspace{.45\textwidth}
    \begin{minipage}{.5\textwidth}
        \imprimirpreambulo
    \end{minipage}%
   \end{center}

   \vspace*{1cm}
        
   \begin{center}
   Aprovado em 17 de julho de 2024.

   \vspace*{1cm}

   Banca examinadora:
   \end{center}

   \vspace*{-1cm}

   \assinatura{Prof. Ma. Caroline Santa Maria Rodrigues -- PUCRS}
   \assinatura{Prof. Dra. Clarissa Pinto Pizarro de Freitas -- PUCRS}
   \assinatura{\imprimirorientador -- PUCRS}
\end{folhadeaprovacao}
% ---


% ---
% Agradecimentos
% ---
\begin{agradecimentos}
Agradeço aos meus pais, Ademilson José Dias e Irinês Porfírio Dias, pelo apoio incondicional que me deram ao longo da vida,
incentivando mesmo as ideias questionáveis como aquela de fazer uma segunda graduação.

A minha irmã, Maria Eduarda Dias, cuja trajetória acadêmica é fonte de inspiração.

A minha namorada, Nathalia Fallavena Ceratti, que acompanhou a esta jornada de perto, oferecendo apoio e compreensão que
fizeram toda a diferença.

Aos meus amigos, Guilherme da Silva Petró, Norton Neves e Uilinton Luciano Marques da Silva, que foram compreensíveis nos
finais de semana em que não pude jogar RPG.

Ao professor Wagner de Lara Machado pela orientação e apoio no desenvolvimento deste trabalho.
\end{agradecimentos}
% ---


% ---
% RESUMOS
% ---
% Resumo em português
\setlength{\absparsep}{18pt} % ajusta o espaçamento dos parágrafos do resumo
\begin{resumo}
Desenvolvimentos na área de aprendizagem de máquina representam uma série de oportunidades para diversas áreas de pesquisa e de atuação profissional.
O uso de modelos de aprendizagem de máquina na predição de desfecho para tratamentos em saúde mental tem sido investigado e apresenta resultados promissores.
A aplicação desse tipo de modelo possibilitaria, por exemplo, maior assertividade na personalização de tratamentos psicoterápicos a nível individual.
Apesar dos benefícios potenciais, o tema ainda é pouco discutido nos espaços de pesquisa, formação e atuação em psicologia clínica. Este trabalho tem
por objetivo introduzir o tema de aprendizagem de máquina a psicólogos. Busca-se apresentar, de maneira acessível, os conceitos e terminologia básicos
em aprendizagem de máquina, proporcionando um panorama geral da área. Para complementar a exposição conceitual, é construído um exemplo a partir de dados
de um estudo sobre eficácia de intervenções em psicologia positiva para depressão via internet. Por fim, discutem-se algumas implicações do uso desse tipo
de tecnologia para a pesquisa e para a prática em psicoterapia.

 \textbf{Palavras-chave}: aprendizagem de máquina, psicoterapia, psicologia positiva, depressão.
\end{resumo}

% Resumo em inglês
\begin{resumo}[Abstract]
 \begin{otherlanguage*}{english}
   Advances in machine learning have opened up a multitude of opportunities for various research and professional fields. The use of machine learning models to
   predict treatment outcomes in mental health has been investigated showing promising results. The application of such models would allow greater assertiveness
   in tailoring psychotherapeutic treatments at the individual level. Despite its potential benefits, the topic remains unexplored in the fields of research, training
   and practice in clinical psychology. This work aims to introduce the topic of machine learning to psychologists. It presents the basic machine learning concepts and
   terminology in an accessible manner, providing a comprehensive overview of the field. To complement the conceptual discussion, an example is constructed using data
   from a study on the effectiveness of internet-based positive psychology interventions for depression. Finally, implications of using this type of technology in
   psychotherapy research and practice are discussed.
   \vspace{\onelineskip}
 
   \noindent 
   \textbf{Keywords}: machine learning, psychotherapy, positive psychology, depression.
 \end{otherlanguage*}
\end{resumo}


% ---
% LISTA DE ILUSTRAÇÕES
% ---
\pdfbookmark[0]{\listfigurename}{lof}
\listoffigures*
\cleardoublepage
% ---


% ---
% LISTA DE QUADROS
% ---
% \pdfbookmark[0]{\listofquadrosname}{loq}
% \listofquadros*
% \cleardoublepage
% ---


% ---
% LISTA DE TABELAS
% ---
% \pdfbookmark[0]{\listtablename}{lot}
% \listoftables*
% \cleardoublepage
% ---


% ---
% LISTA DE ABREVIATURAS E SIGLAS
% ---
% \begin{siglas}
%   \item[ABNT] Associação Brasileira de Normas Técnicas
%   \item[abnTeX] ABsurdas Normas para TeX
% \end{siglas}
% ---


% ---
% LISTA DE SÍMBOLOS
% ---
% \begin{simbolos}
%   \item[$ \Gamma $] Letra grega Gama
%   \item[$ \Lambda $] Lambda
%   \item[$ \zeta $] Letra grega minúscula zeta
%   \item[$ \in $] Pertence
% \end{simbolos}
% ---


% ---
% Sumário
% ---
\pdfbookmark[0]{\contentsname}{toc}
\tableofcontents*
\cleardoublepage
% ---


% ----------------------------------------------------------
% ELEMENTOS TEXTUAIS
% ----------------------------------------------------------
\textual

\documentclass[
	% Opções da classe memoir
	12pt,				        % tamanho da fonte
	openany,			        % capítulos começam em pág ímpar (insere página vazia caso preciso)
	oneside,			        % para impressão em recto e verso. Oposto a oneside
	a4paper,			        % tamanho do papel. 
	% Opções da classe abntex2
	%chapter=TITLE,		        % títulos de capítulos convertidos em letras maiúsculas
	%section=TITLE,		        % títulos de seções convertidos em letras maiúsculas
	%subsection=TITLE,	        % títulos de subseções convertidos em letras maiúsculas
	%subsubsection=TITLE,       % títulos de subsubseções convertidos em letras maiúsculas
	% Opções do pacote babel
	english,			        % idioma adicional para hifenização
	brazil				        % o último idioma é o principal do documento
	]{abntex2}

% ---
% Pacotes básicos 
% ---
\usepackage{lmodern}			% Usa a fonte Latin Modern			
\usepackage[T1]{fontenc}		% Selecao de codigos de fonte.
\usepackage[utf8]{inputenc}		% Codificacao do documento (conversão automática dos acentos)
\usepackage{indentfirst}		% Indenta o primeiro parágrafo de cada seção.
\usepackage{color}				% Controle das cores
\usepackage{graphicx}			% Inclusão de gráficos
\usepackage{microtype} 			% para melhorias de justificação
% ---
		

% ---
% Pacotes de citações
% ---
% \usepackage[alf]{abntex2cite}						% Citações padrão ABNT
\usepackage[
  alf,
  abnt-nbr10520=1988,                     % Citações em sentence case
  abnt-emphasize=bf
  ]{abntex2cite}


% ---
% Pacotes adicionais
% ---
\usepackage{pdfpages}	% Inclusão de páginas em PDF


% ---
% Informações para CAPA e FOLHA DE ROSTO
% ---
\titulo{APRENDIZAGEM DE MÁQUINA PARA PREDIÇÃO DE DESFECHOS EM SAÚDE MENTAL: \textmd{UMA INTRODUÇÃO PARA PSICÓLOGOS}}
\autor{GUILHERME AUGUSTO DIAS}
\local{Porto Alegre}
\data{2024}
\orientador{Prof. Dr. Wagner de Lara Machado}
% \coorientador{}
% \instituicao{
%   Pontifícia Universidade Católica do Rio Grande do Sul
%   \par
%   Escola de Ciências da Saúde e da Vida
%   \par
%   Curso de Psicologia}
\tipotrabalho{TCC}
\preambulo{Trabalho de Conclusão de Curso apresentado como requisito para a obtenção do grau de Bacharel em Psicologia pela Escola de Ciências da Saúde e da Vida da Pontifícia Universidade Católica do Rio Grande do Sul.}
% ---

% ---
% Configurações de aparência do PDF final
% ---

% Alterando o aspecto da cor azul
\definecolor{blue}{RGB}{41,5,195}

% Informações do PDF
\makeatletter
\hypersetup{
		pdftitle={Aprendizagem de máquina para predição de desfechos em saúde mental: uma introdução para psicólogos}, 
		pdfauthor={Guilherme Augusto Dias},
    	pdfsubject={\imprimirpreambulo},
	    pdfcreator={LaTeX with abnTeX2},
		pdfkeywords={aprendizagem de máquina}{psicoterapia}{psicologia positiva}{depressão}, 
		colorlinks=true,		% true: colored links; false: boxed links
    	linkcolor=black,		% color of internal links
    	citecolor=black,		% color of links to bibliography
    	filecolor=black,		% color of file links
		urlcolor=black,
		bookmarksdepth=4
}
\makeatother
% --- 


% ---
% Posiciona figuras e tabelas no topo da página quando adicionadas sozinhas
% em um página em branco. Ver https://github.com/abntex/abntex2/issues/170
\makeatletter
\setlength{\@fptop}{5pt} % Set distance from top of page to first float
\makeatother
% ---

% ---
% Possibilita criação de Quadros e Lista de quadros.
% Ver https://github.com/abntex/abntex2/issues/176
%
\newcommand{\quadroname}{Quadro}
\newcommand{\listofquadrosname}{Lista de quadros}

\newfloat[chapter]{quadro}{loq}{\quadroname}
\newlistof{listofquadros}{loq}{\listofquadrosname}
\newlistentry{quadro}{loq}{0}

% Configurações para atender às regras da ABNT
\setfloatadjustment{quadro}{\centering}
\counterwithout{quadro}{chapter}
\renewcommand{\cftquadroname}{\quadroname\space} 
\renewcommand*{\cftquadroaftersnum}{\hfill--\hfill}

\setfloatlocations{quadro}{hbtp} % Ver https://github.com/abntex/abntex2/issues/176
% ---


% --- 
% Espaçamentos entre linhas e parágrafos 
% --- 

% O tamanho do parágrafo é dado por:
\setlength{\parindent}{1.3cm}

% Controle do espaçamento entre um parágrafo e outro:
\setlength{\parskip}{0.2cm}  % tente também \onelineskip

% ---
% Compila o índice
% ---
\makeindex
% ---


% ----
% Início do documento
% ----
\begin{document}

% Seleciona o idioma do documento (conforme pacotes do babel)
\selectlanguage{brazil}

% Retira espaço extra obsoleto entre as frases.
\frenchspacing 

% ----------------------------------------------------------
% ELEMENTOS PRÉ-TEXTUAIS
% ----------------------------------------------------------
\pretextual

% ---
% CAPA
% ---
\includepdf{./05-extras/capa.pdf}
% ---


% ---
% FOLHA DE ROSTO
% (o * indica que haverá a ficha bibliográfica)
% ---
\imprimirfolhaderosto
% ---


% ---
% Ficha bibliografica
% ---

% Isto é um exemplo de Ficha Catalográfica, ou ``Dados internacionais de
% catalogação-na-publicação''. Você pode utilizar este modelo como referência. 
% Porém, provavelmente a biblioteca da sua universidade lhe fornecerá um PDF
% com a ficha catalográfica definitiva após a defesa do trabalho. Quando estiver
% com o documento, salve-o como PDF no diretório do seu projeto e substitua todo
% o conteúdo de implementação deste arquivo pelo comando abaixo:
%
% \usepackage{pdfpages}		% necessário para comando \includepdf
% \begin{fichacatalografica}
%     \includepdf{fig_ficha_catalografica.pdf}
% \end{fichacatalografica}

% \begin{fichacatalografica}
% 	\sffamily
% 	\vspace*{\fill}                             % Posição vertical
% 	\begin{center}                              % Minipage Centralizado
% 	\fbox{\begin{minipage}[c][8cm]{13.5cm}		% Largura
% 	\small
% 	\imprimirautor
% 	%Sobrenome, Nome do autor
	
% 	\hspace{0.5cm} \imprimirtitulo  / \imprimirautor. --
% 	\imprimirlocal, \imprimirdata-
	
% 	\hspace{0.5cm} \thelastpage p. : il. (algumas color.) ; 30 cm.\\
	
% 	\hspace{0.5cm} \imprimirorientadorRotulo~\imprimirorientador\\
	
% 	\hspace{0.5cm}
% 	\parbox[t]{\textwidth}{\imprimirtipotrabalho~--~\imprimirinstituicao,
% 	\imprimirdata.}\\
	
% 	\hspace{0.5cm}
% 		1. Palavra-chave1.
% 		2. Palavra-chave2.
% 		2. Palavra-chave3.
% 		I. Orientador.
% 		II. Universidade xxx.
% 		III. Faculdade de xxx.
% 		IV. Título 			
% 	\end{minipage}}
% 	\end{center}
% \end{fichacatalografica}
% ---


% ---
% FOLHA DE APROVAÇÃO
% ---

% Isto é um exemplo de Folha de aprovação, elemento obrigatório da NBR
% 14724/2011 (seção 4.2.1.3). Você pode utilizar este modelo até a aprovação
% do trabalho. Após isso, substitua todo o conteúdo deste arquivo por uma
% imagem da página assinada pela banca com o comando abaixo:
%
% \begin{folhadeaprovacao}
% \includepdf{folhadeaprovacao_final.pdf}
% \end{folhadeaprovacao}
%
\setlength{\ABNTEXsignskip}{2cm}
\setlength{\ABNTEXsignwidth}{10cm}

\begin{folhadeaprovacao}
  \begin{center}
    {\ABNTEXchapterfont\large\imprimirautor}

    \vspace*{1cm}

    \begin{center}
      \ABNTEXchapterfont\bfseries\Large\imprimirtitulo
    \end{center}

    \vspace*{1cm}
    
    \hspace{.45\textwidth}
    \begin{minipage}{.5\textwidth}
        \imprimirpreambulo
    \end{minipage}%
   \end{center}

   \vspace*{1cm}
        
   \begin{center}
   Aprovado em 17 de julho de 2024.

   \vspace*{1cm}

   Banca examinadora:
   \end{center}

   \vspace*{-1cm}

   \assinatura{Prof. Ma. Caroline Santa Maria Rodrigues -- PUCRS}
   \assinatura{Prof. Dra. Clarissa Pinto Pizarro de Freitas -- PUCRS}
   \assinatura{\imprimirorientador -- PUCRS}
\end{folhadeaprovacao}
% ---


% ---
% Agradecimentos
% ---
\begin{agradecimentos}
Agradeço aos meus pais, Ademilson José Dias e Irinês Porfírio Dias, pelo apoio incondicional que me deram ao longo da vida,
incentivando mesmo as ideias questionáveis como aquela de fazer uma segunda graduação.

A minha irmã, Maria Eduarda Dias, cuja trajetória acadêmica é fonte de inspiração.

A minha namorada, Nathalia Fallavena Ceratti, que acompanhou a esta jornada de perto, oferecendo apoio e compreensão que
fizeram toda a diferença.

Aos meus amigos, Guilherme da Silva Petró, Norton Neves e Uilinton Luciano Marques da Silva, que foram compreensíveis nos
finais de semana em que não pude jogar RPG.

Ao professor Wagner de Lara Machado pela orientação e apoio no desenvolvimento deste trabalho.
\end{agradecimentos}
% ---


% ---
% RESUMOS
% ---
% Resumo em português
\setlength{\absparsep}{18pt} % ajusta o espaçamento dos parágrafos do resumo
\begin{resumo}
Desenvolvimentos na área de aprendizagem de máquina representam uma série de oportunidades para diversas áreas de pesquisa e de atuação profissional.
O uso de modelos de aprendizagem de máquina na predição de desfecho para tratamentos em saúde mental tem sido investigado e apresenta resultados promissores.
A aplicação desse tipo de modelo possibilitaria, por exemplo, maior assertividade na personalização de tratamentos psicoterápicos a nível individual.
Apesar dos benefícios potenciais, o tema ainda é pouco discutido nos espaços de pesquisa, formação e atuação em psicologia clínica. Este trabalho tem
por objetivo introduzir o tema de aprendizagem de máquina a psicólogos. Busca-se apresentar, de maneira acessível, os conceitos e terminologia básicos
em aprendizagem de máquina, proporcionando um panorama geral da área. Para complementar a exposição conceitual, é construído um exemplo a partir de dados
de um estudo sobre eficácia de intervenções em psicologia positiva para depressão via internet. Por fim, discutem-se algumas implicações do uso desse tipo
de tecnologia para a pesquisa e para a prática em psicoterapia.

 \textbf{Palavras-chave}: aprendizagem de máquina, psicoterapia, psicologia positiva, depressão.
\end{resumo}

% Resumo em inglês
\begin{resumo}[Abstract]
 \begin{otherlanguage*}{english}
   Advances in machine learning have opened up a multitude of opportunities for various research and professional fields. The use of machine learning models to
   predict treatment outcomes in mental health has been investigated showing promising results. The application of such models would allow greater assertiveness
   in tailoring psychotherapeutic treatments at the individual level. Despite its potential benefits, the topic remains unexplored in the fields of research, training
   and practice in clinical psychology. This work aims to introduce the topic of machine learning to psychologists. It presents the basic machine learning concepts and
   terminology in an accessible manner, providing a comprehensive overview of the field. To complement the conceptual discussion, an example is constructed using data
   from a study on the effectiveness of internet-based positive psychology interventions for depression. Finally, implications of using this type of technology in
   psychotherapy research and practice are discussed.
   \vspace{\onelineskip}
 
   \noindent 
   \textbf{Keywords}: machine learning, psychotherapy, positive psychology, depression.
 \end{otherlanguage*}
\end{resumo}


% ---
% LISTA DE ILUSTRAÇÕES
% ---
\pdfbookmark[0]{\listfigurename}{lof}
\listoffigures*
\cleardoublepage
% ---


% ---
% LISTA DE QUADROS
% ---
% \pdfbookmark[0]{\listofquadrosname}{loq}
% \listofquadros*
% \cleardoublepage
% ---


% ---
% LISTA DE TABELAS
% ---
% \pdfbookmark[0]{\listtablename}{lot}
% \listoftables*
% \cleardoublepage
% ---


% ---
% LISTA DE ABREVIATURAS E SIGLAS
% ---
% \begin{siglas}
%   \item[ABNT] Associação Brasileira de Normas Técnicas
%   \item[abnTeX] ABsurdas Normas para TeX
% \end{siglas}
% ---


% ---
% LISTA DE SÍMBOLOS
% ---
% \begin{simbolos}
%   \item[$ \Gamma $] Letra grega Gama
%   \item[$ \Lambda $] Lambda
%   \item[$ \zeta $] Letra grega minúscula zeta
%   \item[$ \in $] Pertence
% \end{simbolos}
% ---


% ---
% Sumário
% ---
\pdfbookmark[0]{\contentsname}{toc}
\tableofcontents*
\cleardoublepage
% ---


% ----------------------------------------------------------
% ELEMENTOS TEXTUAIS
% ----------------------------------------------------------
\textual

\include{./01-introducao/main.tex}

\include{./02-machine-learning/main.tex}

\include{./03-exemplo/main.tex}

\include{./04-conclusao/main.tex}

% ----------------------------------------------------------
% ELEMENTOS PÓS-TEXTUAIS
% ----------------------------------------------------------
\postextual
% ----------------------------------------------------------


% ----------------------------------------------------------
% REFERÊNCIAS
% ----------------------------------------------------------
\bibliography{references.bib}


% ----------------------------------------------------------
% GLOSSÁRIO
% ----------------------------------------------------------
%
% Consulte o manual da classe abntex2 para orientações sobre o glossário.
%
%\glossary

% ----------------------------------------------------------
% APÊNDICES
% ----------------------------------------------------------

% ---
% Inicia os apêndices
% ---
% \begin{apendicesenv}

% % Imprime uma página indicando o início dos apêndices
% \partapendices

% % ----------------------------------------------------------
% \chapter{Quisque libero justo}
% % ----------------------------------------------------------

% % ----------------------------------------------------------
% \chapter{Nullam elementum urna}
% % ----------------------------------------------------------

% \end{apendicesenv}
% ---


% ----------------------------------------------------------
% ANEXOS
% ----------------------------------------------------------

% ---
% Inicia os anexos
% ---
% \begin{anexosenv}

% % Imprime uma página indicando o início dos anexos
% \partanexos

% % ---
% \chapter{Morbi ultrices rutrum lorem}
% % ---

% % ---
% \chapter{Cras non urna sed feugiat cum}
% % ---


% % ---
% \chapter{Fusce facilisis lacinia dui}
% % ---

% \end{anexosenv}

%---------------------------------------------------------------------
% INDICE REMISSIVO
%---------------------------------------------------------------------
% \phantompart
% \printindex
%---------------------------------------------------------------------
\includepdf{./05-extras/contracapa.pdf}
\end{document}


\documentclass[
	% Opções da classe memoir
	12pt,				        % tamanho da fonte
	openany,			        % capítulos começam em pág ímpar (insere página vazia caso preciso)
	oneside,			        % para impressão em recto e verso. Oposto a oneside
	a4paper,			        % tamanho do papel. 
	% Opções da classe abntex2
	%chapter=TITLE,		        % títulos de capítulos convertidos em letras maiúsculas
	%section=TITLE,		        % títulos de seções convertidos em letras maiúsculas
	%subsection=TITLE,	        % títulos de subseções convertidos em letras maiúsculas
	%subsubsection=TITLE,       % títulos de subsubseções convertidos em letras maiúsculas
	% Opções do pacote babel
	english,			        % idioma adicional para hifenização
	brazil				        % o último idioma é o principal do documento
	]{abntex2}

% ---
% Pacotes básicos 
% ---
\usepackage{lmodern}			% Usa a fonte Latin Modern			
\usepackage[T1]{fontenc}		% Selecao de codigos de fonte.
\usepackage[utf8]{inputenc}		% Codificacao do documento (conversão automática dos acentos)
\usepackage{indentfirst}		% Indenta o primeiro parágrafo de cada seção.
\usepackage{color}				% Controle das cores
\usepackage{graphicx}			% Inclusão de gráficos
\usepackage{microtype} 			% para melhorias de justificação
% ---
		

% ---
% Pacotes de citações
% ---
% \usepackage[alf]{abntex2cite}						% Citações padrão ABNT
\usepackage[
  alf,
  abnt-nbr10520=1988,                     % Citações em sentence case
  abnt-emphasize=bf
  ]{abntex2cite}


% ---
% Pacotes adicionais
% ---
\usepackage{pdfpages}	% Inclusão de páginas em PDF


% ---
% Informações para CAPA e FOLHA DE ROSTO
% ---
\titulo{APRENDIZAGEM DE MÁQUINA PARA PREDIÇÃO DE DESFECHOS EM SAÚDE MENTAL: \textmd{UMA INTRODUÇÃO PARA PSICÓLOGOS}}
\autor{GUILHERME AUGUSTO DIAS}
\local{Porto Alegre}
\data{2024}
\orientador{Prof. Dr. Wagner de Lara Machado}
% \coorientador{}
% \instituicao{
%   Pontifícia Universidade Católica do Rio Grande do Sul
%   \par
%   Escola de Ciências da Saúde e da Vida
%   \par
%   Curso de Psicologia}
\tipotrabalho{TCC}
\preambulo{Trabalho de Conclusão de Curso apresentado como requisito para a obtenção do grau de Bacharel em Psicologia pela Escola de Ciências da Saúde e da Vida da Pontifícia Universidade Católica do Rio Grande do Sul.}
% ---

% ---
% Configurações de aparência do PDF final
% ---

% Alterando o aspecto da cor azul
\definecolor{blue}{RGB}{41,5,195}

% Informações do PDF
\makeatletter
\hypersetup{
		pdftitle={Aprendizagem de máquina para predição de desfechos em saúde mental: uma introdução para psicólogos}, 
		pdfauthor={Guilherme Augusto Dias},
    	pdfsubject={\imprimirpreambulo},
	    pdfcreator={LaTeX with abnTeX2},
		pdfkeywords={aprendizagem de máquina}{psicoterapia}{psicologia positiva}{depressão}, 
		colorlinks=true,		% true: colored links; false: boxed links
    	linkcolor=black,		% color of internal links
    	citecolor=black,		% color of links to bibliography
    	filecolor=black,		% color of file links
		urlcolor=black,
		bookmarksdepth=4
}
\makeatother
% --- 


% ---
% Posiciona figuras e tabelas no topo da página quando adicionadas sozinhas
% em um página em branco. Ver https://github.com/abntex/abntex2/issues/170
\makeatletter
\setlength{\@fptop}{5pt} % Set distance from top of page to first float
\makeatother
% ---

% ---
% Possibilita criação de Quadros e Lista de quadros.
% Ver https://github.com/abntex/abntex2/issues/176
%
\newcommand{\quadroname}{Quadro}
\newcommand{\listofquadrosname}{Lista de quadros}

\newfloat[chapter]{quadro}{loq}{\quadroname}
\newlistof{listofquadros}{loq}{\listofquadrosname}
\newlistentry{quadro}{loq}{0}

% Configurações para atender às regras da ABNT
\setfloatadjustment{quadro}{\centering}
\counterwithout{quadro}{chapter}
\renewcommand{\cftquadroname}{\quadroname\space} 
\renewcommand*{\cftquadroaftersnum}{\hfill--\hfill}

\setfloatlocations{quadro}{hbtp} % Ver https://github.com/abntex/abntex2/issues/176
% ---


% --- 
% Espaçamentos entre linhas e parágrafos 
% --- 

% O tamanho do parágrafo é dado por:
\setlength{\parindent}{1.3cm}

% Controle do espaçamento entre um parágrafo e outro:
\setlength{\parskip}{0.2cm}  % tente também \onelineskip

% ---
% Compila o índice
% ---
\makeindex
% ---


% ----
% Início do documento
% ----
\begin{document}

% Seleciona o idioma do documento (conforme pacotes do babel)
\selectlanguage{brazil}

% Retira espaço extra obsoleto entre as frases.
\frenchspacing 

% ----------------------------------------------------------
% ELEMENTOS PRÉ-TEXTUAIS
% ----------------------------------------------------------
\pretextual

% ---
% CAPA
% ---
\includepdf{./05-extras/capa.pdf}
% ---


% ---
% FOLHA DE ROSTO
% (o * indica que haverá a ficha bibliográfica)
% ---
\imprimirfolhaderosto
% ---


% ---
% Ficha bibliografica
% ---

% Isto é um exemplo de Ficha Catalográfica, ou ``Dados internacionais de
% catalogação-na-publicação''. Você pode utilizar este modelo como referência. 
% Porém, provavelmente a biblioteca da sua universidade lhe fornecerá um PDF
% com a ficha catalográfica definitiva após a defesa do trabalho. Quando estiver
% com o documento, salve-o como PDF no diretório do seu projeto e substitua todo
% o conteúdo de implementação deste arquivo pelo comando abaixo:
%
% \usepackage{pdfpages}		% necessário para comando \includepdf
% \begin{fichacatalografica}
%     \includepdf{fig_ficha_catalografica.pdf}
% \end{fichacatalografica}

% \begin{fichacatalografica}
% 	\sffamily
% 	\vspace*{\fill}                             % Posição vertical
% 	\begin{center}                              % Minipage Centralizado
% 	\fbox{\begin{minipage}[c][8cm]{13.5cm}		% Largura
% 	\small
% 	\imprimirautor
% 	%Sobrenome, Nome do autor
	
% 	\hspace{0.5cm} \imprimirtitulo  / \imprimirautor. --
% 	\imprimirlocal, \imprimirdata-
	
% 	\hspace{0.5cm} \thelastpage p. : il. (algumas color.) ; 30 cm.\\
	
% 	\hspace{0.5cm} \imprimirorientadorRotulo~\imprimirorientador\\
	
% 	\hspace{0.5cm}
% 	\parbox[t]{\textwidth}{\imprimirtipotrabalho~--~\imprimirinstituicao,
% 	\imprimirdata.}\\
	
% 	\hspace{0.5cm}
% 		1. Palavra-chave1.
% 		2. Palavra-chave2.
% 		2. Palavra-chave3.
% 		I. Orientador.
% 		II. Universidade xxx.
% 		III. Faculdade de xxx.
% 		IV. Título 			
% 	\end{minipage}}
% 	\end{center}
% \end{fichacatalografica}
% ---


% ---
% FOLHA DE APROVAÇÃO
% ---

% Isto é um exemplo de Folha de aprovação, elemento obrigatório da NBR
% 14724/2011 (seção 4.2.1.3). Você pode utilizar este modelo até a aprovação
% do trabalho. Após isso, substitua todo o conteúdo deste arquivo por uma
% imagem da página assinada pela banca com o comando abaixo:
%
% \begin{folhadeaprovacao}
% \includepdf{folhadeaprovacao_final.pdf}
% \end{folhadeaprovacao}
%
\setlength{\ABNTEXsignskip}{2cm}
\setlength{\ABNTEXsignwidth}{10cm}

\begin{folhadeaprovacao}
  \begin{center}
    {\ABNTEXchapterfont\large\imprimirautor}

    \vspace*{1cm}

    \begin{center}
      \ABNTEXchapterfont\bfseries\Large\imprimirtitulo
    \end{center}

    \vspace*{1cm}
    
    \hspace{.45\textwidth}
    \begin{minipage}{.5\textwidth}
        \imprimirpreambulo
    \end{minipage}%
   \end{center}

   \vspace*{1cm}
        
   \begin{center}
   Aprovado em 17 de julho de 2024.

   \vspace*{1cm}

   Banca examinadora:
   \end{center}

   \vspace*{-1cm}

   \assinatura{Prof. Ma. Caroline Santa Maria Rodrigues -- PUCRS}
   \assinatura{Prof. Dra. Clarissa Pinto Pizarro de Freitas -- PUCRS}
   \assinatura{\imprimirorientador -- PUCRS}
\end{folhadeaprovacao}
% ---


% ---
% Agradecimentos
% ---
\begin{agradecimentos}
Agradeço aos meus pais, Ademilson José Dias e Irinês Porfírio Dias, pelo apoio incondicional que me deram ao longo da vida,
incentivando mesmo as ideias questionáveis como aquela de fazer uma segunda graduação.

A minha irmã, Maria Eduarda Dias, cuja trajetória acadêmica é fonte de inspiração.

A minha namorada, Nathalia Fallavena Ceratti, que acompanhou a esta jornada de perto, oferecendo apoio e compreensão que
fizeram toda a diferença.

Aos meus amigos, Guilherme da Silva Petró, Norton Neves e Uilinton Luciano Marques da Silva, que foram compreensíveis nos
finais de semana em que não pude jogar RPG.

Ao professor Wagner de Lara Machado pela orientação e apoio no desenvolvimento deste trabalho.
\end{agradecimentos}
% ---


% ---
% RESUMOS
% ---
% Resumo em português
\setlength{\absparsep}{18pt} % ajusta o espaçamento dos parágrafos do resumo
\begin{resumo}
Desenvolvimentos na área de aprendizagem de máquina representam uma série de oportunidades para diversas áreas de pesquisa e de atuação profissional.
O uso de modelos de aprendizagem de máquina na predição de desfecho para tratamentos em saúde mental tem sido investigado e apresenta resultados promissores.
A aplicação desse tipo de modelo possibilitaria, por exemplo, maior assertividade na personalização de tratamentos psicoterápicos a nível individual.
Apesar dos benefícios potenciais, o tema ainda é pouco discutido nos espaços de pesquisa, formação e atuação em psicologia clínica. Este trabalho tem
por objetivo introduzir o tema de aprendizagem de máquina a psicólogos. Busca-se apresentar, de maneira acessível, os conceitos e terminologia básicos
em aprendizagem de máquina, proporcionando um panorama geral da área. Para complementar a exposição conceitual, é construído um exemplo a partir de dados
de um estudo sobre eficácia de intervenções em psicologia positiva para depressão via internet. Por fim, discutem-se algumas implicações do uso desse tipo
de tecnologia para a pesquisa e para a prática em psicoterapia.

 \textbf{Palavras-chave}: aprendizagem de máquina, psicoterapia, psicologia positiva, depressão.
\end{resumo}

% Resumo em inglês
\begin{resumo}[Abstract]
 \begin{otherlanguage*}{english}
   Advances in machine learning have opened up a multitude of opportunities for various research and professional fields. The use of machine learning models to
   predict treatment outcomes in mental health has been investigated showing promising results. The application of such models would allow greater assertiveness
   in tailoring psychotherapeutic treatments at the individual level. Despite its potential benefits, the topic remains unexplored in the fields of research, training
   and practice in clinical psychology. This work aims to introduce the topic of machine learning to psychologists. It presents the basic machine learning concepts and
   terminology in an accessible manner, providing a comprehensive overview of the field. To complement the conceptual discussion, an example is constructed using data
   from a study on the effectiveness of internet-based positive psychology interventions for depression. Finally, implications of using this type of technology in
   psychotherapy research and practice are discussed.
   \vspace{\onelineskip}
 
   \noindent 
   \textbf{Keywords}: machine learning, psychotherapy, positive psychology, depression.
 \end{otherlanguage*}
\end{resumo}


% ---
% LISTA DE ILUSTRAÇÕES
% ---
\pdfbookmark[0]{\listfigurename}{lof}
\listoffigures*
\cleardoublepage
% ---


% ---
% LISTA DE QUADROS
% ---
% \pdfbookmark[0]{\listofquadrosname}{loq}
% \listofquadros*
% \cleardoublepage
% ---


% ---
% LISTA DE TABELAS
% ---
% \pdfbookmark[0]{\listtablename}{lot}
% \listoftables*
% \cleardoublepage
% ---


% ---
% LISTA DE ABREVIATURAS E SIGLAS
% ---
% \begin{siglas}
%   \item[ABNT] Associação Brasileira de Normas Técnicas
%   \item[abnTeX] ABsurdas Normas para TeX
% \end{siglas}
% ---


% ---
% LISTA DE SÍMBOLOS
% ---
% \begin{simbolos}
%   \item[$ \Gamma $] Letra grega Gama
%   \item[$ \Lambda $] Lambda
%   \item[$ \zeta $] Letra grega minúscula zeta
%   \item[$ \in $] Pertence
% \end{simbolos}
% ---


% ---
% Sumário
% ---
\pdfbookmark[0]{\contentsname}{toc}
\tableofcontents*
\cleardoublepage
% ---


% ----------------------------------------------------------
% ELEMENTOS TEXTUAIS
% ----------------------------------------------------------
\textual

\include{./01-introducao/main.tex}

\include{./02-machine-learning/main.tex}

\include{./03-exemplo/main.tex}

\include{./04-conclusao/main.tex}

% ----------------------------------------------------------
% ELEMENTOS PÓS-TEXTUAIS
% ----------------------------------------------------------
\postextual
% ----------------------------------------------------------


% ----------------------------------------------------------
% REFERÊNCIAS
% ----------------------------------------------------------
\bibliography{references.bib}


% ----------------------------------------------------------
% GLOSSÁRIO
% ----------------------------------------------------------
%
% Consulte o manual da classe abntex2 para orientações sobre o glossário.
%
%\glossary

% ----------------------------------------------------------
% APÊNDICES
% ----------------------------------------------------------

% ---
% Inicia os apêndices
% ---
% \begin{apendicesenv}

% % Imprime uma página indicando o início dos apêndices
% \partapendices

% % ----------------------------------------------------------
% \chapter{Quisque libero justo}
% % ----------------------------------------------------------

% % ----------------------------------------------------------
% \chapter{Nullam elementum urna}
% % ----------------------------------------------------------

% \end{apendicesenv}
% ---


% ----------------------------------------------------------
% ANEXOS
% ----------------------------------------------------------

% ---
% Inicia os anexos
% ---
% \begin{anexosenv}

% % Imprime uma página indicando o início dos anexos
% \partanexos

% % ---
% \chapter{Morbi ultrices rutrum lorem}
% % ---

% % ---
% \chapter{Cras non urna sed feugiat cum}
% % ---


% % ---
% \chapter{Fusce facilisis lacinia dui}
% % ---

% \end{anexosenv}

%---------------------------------------------------------------------
% INDICE REMISSIVO
%---------------------------------------------------------------------
% \phantompart
% \printindex
%---------------------------------------------------------------------
\includepdf{./05-extras/contracapa.pdf}
\end{document}


\documentclass[
	% Opções da classe memoir
	12pt,				        % tamanho da fonte
	openany,			        % capítulos começam em pág ímpar (insere página vazia caso preciso)
	oneside,			        % para impressão em recto e verso. Oposto a oneside
	a4paper,			        % tamanho do papel. 
	% Opções da classe abntex2
	%chapter=TITLE,		        % títulos de capítulos convertidos em letras maiúsculas
	%section=TITLE,		        % títulos de seções convertidos em letras maiúsculas
	%subsection=TITLE,	        % títulos de subseções convertidos em letras maiúsculas
	%subsubsection=TITLE,       % títulos de subsubseções convertidos em letras maiúsculas
	% Opções do pacote babel
	english,			        % idioma adicional para hifenização
	brazil				        % o último idioma é o principal do documento
	]{abntex2}

% ---
% Pacotes básicos 
% ---
\usepackage{lmodern}			% Usa a fonte Latin Modern			
\usepackage[T1]{fontenc}		% Selecao de codigos de fonte.
\usepackage[utf8]{inputenc}		% Codificacao do documento (conversão automática dos acentos)
\usepackage{indentfirst}		% Indenta o primeiro parágrafo de cada seção.
\usepackage{color}				% Controle das cores
\usepackage{graphicx}			% Inclusão de gráficos
\usepackage{microtype} 			% para melhorias de justificação
% ---
		

% ---
% Pacotes de citações
% ---
% \usepackage[alf]{abntex2cite}						% Citações padrão ABNT
\usepackage[
  alf,
  abnt-nbr10520=1988,                     % Citações em sentence case
  abnt-emphasize=bf
  ]{abntex2cite}


% ---
% Pacotes adicionais
% ---
\usepackage{pdfpages}	% Inclusão de páginas em PDF


% ---
% Informações para CAPA e FOLHA DE ROSTO
% ---
\titulo{APRENDIZAGEM DE MÁQUINA PARA PREDIÇÃO DE DESFECHOS EM SAÚDE MENTAL: \textmd{UMA INTRODUÇÃO PARA PSICÓLOGOS}}
\autor{GUILHERME AUGUSTO DIAS}
\local{Porto Alegre}
\data{2024}
\orientador{Prof. Dr. Wagner de Lara Machado}
% \coorientador{}
% \instituicao{
%   Pontifícia Universidade Católica do Rio Grande do Sul
%   \par
%   Escola de Ciências da Saúde e da Vida
%   \par
%   Curso de Psicologia}
\tipotrabalho{TCC}
\preambulo{Trabalho de Conclusão de Curso apresentado como requisito para a obtenção do grau de Bacharel em Psicologia pela Escola de Ciências da Saúde e da Vida da Pontifícia Universidade Católica do Rio Grande do Sul.}
% ---

% ---
% Configurações de aparência do PDF final
% ---

% Alterando o aspecto da cor azul
\definecolor{blue}{RGB}{41,5,195}

% Informações do PDF
\makeatletter
\hypersetup{
		pdftitle={Aprendizagem de máquina para predição de desfechos em saúde mental: uma introdução para psicólogos}, 
		pdfauthor={Guilherme Augusto Dias},
    	pdfsubject={\imprimirpreambulo},
	    pdfcreator={LaTeX with abnTeX2},
		pdfkeywords={aprendizagem de máquina}{psicoterapia}{psicologia positiva}{depressão}, 
		colorlinks=true,		% true: colored links; false: boxed links
    	linkcolor=black,		% color of internal links
    	citecolor=black,		% color of links to bibliography
    	filecolor=black,		% color of file links
		urlcolor=black,
		bookmarksdepth=4
}
\makeatother
% --- 


% ---
% Posiciona figuras e tabelas no topo da página quando adicionadas sozinhas
% em um página em branco. Ver https://github.com/abntex/abntex2/issues/170
\makeatletter
\setlength{\@fptop}{5pt} % Set distance from top of page to first float
\makeatother
% ---

% ---
% Possibilita criação de Quadros e Lista de quadros.
% Ver https://github.com/abntex/abntex2/issues/176
%
\newcommand{\quadroname}{Quadro}
\newcommand{\listofquadrosname}{Lista de quadros}

\newfloat[chapter]{quadro}{loq}{\quadroname}
\newlistof{listofquadros}{loq}{\listofquadrosname}
\newlistentry{quadro}{loq}{0}

% Configurações para atender às regras da ABNT
\setfloatadjustment{quadro}{\centering}
\counterwithout{quadro}{chapter}
\renewcommand{\cftquadroname}{\quadroname\space} 
\renewcommand*{\cftquadroaftersnum}{\hfill--\hfill}

\setfloatlocations{quadro}{hbtp} % Ver https://github.com/abntex/abntex2/issues/176
% ---


% --- 
% Espaçamentos entre linhas e parágrafos 
% --- 

% O tamanho do parágrafo é dado por:
\setlength{\parindent}{1.3cm}

% Controle do espaçamento entre um parágrafo e outro:
\setlength{\parskip}{0.2cm}  % tente também \onelineskip

% ---
% Compila o índice
% ---
\makeindex
% ---


% ----
% Início do documento
% ----
\begin{document}

% Seleciona o idioma do documento (conforme pacotes do babel)
\selectlanguage{brazil}

% Retira espaço extra obsoleto entre as frases.
\frenchspacing 

% ----------------------------------------------------------
% ELEMENTOS PRÉ-TEXTUAIS
% ----------------------------------------------------------
\pretextual

% ---
% CAPA
% ---
\includepdf{./05-extras/capa.pdf}
% ---


% ---
% FOLHA DE ROSTO
% (o * indica que haverá a ficha bibliográfica)
% ---
\imprimirfolhaderosto
% ---


% ---
% Ficha bibliografica
% ---

% Isto é um exemplo de Ficha Catalográfica, ou ``Dados internacionais de
% catalogação-na-publicação''. Você pode utilizar este modelo como referência. 
% Porém, provavelmente a biblioteca da sua universidade lhe fornecerá um PDF
% com a ficha catalográfica definitiva após a defesa do trabalho. Quando estiver
% com o documento, salve-o como PDF no diretório do seu projeto e substitua todo
% o conteúdo de implementação deste arquivo pelo comando abaixo:
%
% \usepackage{pdfpages}		% necessário para comando \includepdf
% \begin{fichacatalografica}
%     \includepdf{fig_ficha_catalografica.pdf}
% \end{fichacatalografica}

% \begin{fichacatalografica}
% 	\sffamily
% 	\vspace*{\fill}                             % Posição vertical
% 	\begin{center}                              % Minipage Centralizado
% 	\fbox{\begin{minipage}[c][8cm]{13.5cm}		% Largura
% 	\small
% 	\imprimirautor
% 	%Sobrenome, Nome do autor
	
% 	\hspace{0.5cm} \imprimirtitulo  / \imprimirautor. --
% 	\imprimirlocal, \imprimirdata-
	
% 	\hspace{0.5cm} \thelastpage p. : il. (algumas color.) ; 30 cm.\\
	
% 	\hspace{0.5cm} \imprimirorientadorRotulo~\imprimirorientador\\
	
% 	\hspace{0.5cm}
% 	\parbox[t]{\textwidth}{\imprimirtipotrabalho~--~\imprimirinstituicao,
% 	\imprimirdata.}\\
	
% 	\hspace{0.5cm}
% 		1. Palavra-chave1.
% 		2. Palavra-chave2.
% 		2. Palavra-chave3.
% 		I. Orientador.
% 		II. Universidade xxx.
% 		III. Faculdade de xxx.
% 		IV. Título 			
% 	\end{minipage}}
% 	\end{center}
% \end{fichacatalografica}
% ---


% ---
% FOLHA DE APROVAÇÃO
% ---

% Isto é um exemplo de Folha de aprovação, elemento obrigatório da NBR
% 14724/2011 (seção 4.2.1.3). Você pode utilizar este modelo até a aprovação
% do trabalho. Após isso, substitua todo o conteúdo deste arquivo por uma
% imagem da página assinada pela banca com o comando abaixo:
%
% \begin{folhadeaprovacao}
% \includepdf{folhadeaprovacao_final.pdf}
% \end{folhadeaprovacao}
%
\setlength{\ABNTEXsignskip}{2cm}
\setlength{\ABNTEXsignwidth}{10cm}

\begin{folhadeaprovacao}
  \begin{center}
    {\ABNTEXchapterfont\large\imprimirautor}

    \vspace*{1cm}

    \begin{center}
      \ABNTEXchapterfont\bfseries\Large\imprimirtitulo
    \end{center}

    \vspace*{1cm}
    
    \hspace{.45\textwidth}
    \begin{minipage}{.5\textwidth}
        \imprimirpreambulo
    \end{minipage}%
   \end{center}

   \vspace*{1cm}
        
   \begin{center}
   Aprovado em 17 de julho de 2024.

   \vspace*{1cm}

   Banca examinadora:
   \end{center}

   \vspace*{-1cm}

   \assinatura{Prof. Ma. Caroline Santa Maria Rodrigues -- PUCRS}
   \assinatura{Prof. Dra. Clarissa Pinto Pizarro de Freitas -- PUCRS}
   \assinatura{\imprimirorientador -- PUCRS}
\end{folhadeaprovacao}
% ---


% ---
% Agradecimentos
% ---
\begin{agradecimentos}
Agradeço aos meus pais, Ademilson José Dias e Irinês Porfírio Dias, pelo apoio incondicional que me deram ao longo da vida,
incentivando mesmo as ideias questionáveis como aquela de fazer uma segunda graduação.

A minha irmã, Maria Eduarda Dias, cuja trajetória acadêmica é fonte de inspiração.

A minha namorada, Nathalia Fallavena Ceratti, que acompanhou a esta jornada de perto, oferecendo apoio e compreensão que
fizeram toda a diferença.

Aos meus amigos, Guilherme da Silva Petró, Norton Neves e Uilinton Luciano Marques da Silva, que foram compreensíveis nos
finais de semana em que não pude jogar RPG.

Ao professor Wagner de Lara Machado pela orientação e apoio no desenvolvimento deste trabalho.
\end{agradecimentos}
% ---


% ---
% RESUMOS
% ---
% Resumo em português
\setlength{\absparsep}{18pt} % ajusta o espaçamento dos parágrafos do resumo
\begin{resumo}
Desenvolvimentos na área de aprendizagem de máquina representam uma série de oportunidades para diversas áreas de pesquisa e de atuação profissional.
O uso de modelos de aprendizagem de máquina na predição de desfecho para tratamentos em saúde mental tem sido investigado e apresenta resultados promissores.
A aplicação desse tipo de modelo possibilitaria, por exemplo, maior assertividade na personalização de tratamentos psicoterápicos a nível individual.
Apesar dos benefícios potenciais, o tema ainda é pouco discutido nos espaços de pesquisa, formação e atuação em psicologia clínica. Este trabalho tem
por objetivo introduzir o tema de aprendizagem de máquina a psicólogos. Busca-se apresentar, de maneira acessível, os conceitos e terminologia básicos
em aprendizagem de máquina, proporcionando um panorama geral da área. Para complementar a exposição conceitual, é construído um exemplo a partir de dados
de um estudo sobre eficácia de intervenções em psicologia positiva para depressão via internet. Por fim, discutem-se algumas implicações do uso desse tipo
de tecnologia para a pesquisa e para a prática em psicoterapia.

 \textbf{Palavras-chave}: aprendizagem de máquina, psicoterapia, psicologia positiva, depressão.
\end{resumo}

% Resumo em inglês
\begin{resumo}[Abstract]
 \begin{otherlanguage*}{english}
   Advances in machine learning have opened up a multitude of opportunities for various research and professional fields. The use of machine learning models to
   predict treatment outcomes in mental health has been investigated showing promising results. The application of such models would allow greater assertiveness
   in tailoring psychotherapeutic treatments at the individual level. Despite its potential benefits, the topic remains unexplored in the fields of research, training
   and practice in clinical psychology. This work aims to introduce the topic of machine learning to psychologists. It presents the basic machine learning concepts and
   terminology in an accessible manner, providing a comprehensive overview of the field. To complement the conceptual discussion, an example is constructed using data
   from a study on the effectiveness of internet-based positive psychology interventions for depression. Finally, implications of using this type of technology in
   psychotherapy research and practice are discussed.
   \vspace{\onelineskip}
 
   \noindent 
   \textbf{Keywords}: machine learning, psychotherapy, positive psychology, depression.
 \end{otherlanguage*}
\end{resumo}


% ---
% LISTA DE ILUSTRAÇÕES
% ---
\pdfbookmark[0]{\listfigurename}{lof}
\listoffigures*
\cleardoublepage
% ---


% ---
% LISTA DE QUADROS
% ---
% \pdfbookmark[0]{\listofquadrosname}{loq}
% \listofquadros*
% \cleardoublepage
% ---


% ---
% LISTA DE TABELAS
% ---
% \pdfbookmark[0]{\listtablename}{lot}
% \listoftables*
% \cleardoublepage
% ---


% ---
% LISTA DE ABREVIATURAS E SIGLAS
% ---
% \begin{siglas}
%   \item[ABNT] Associação Brasileira de Normas Técnicas
%   \item[abnTeX] ABsurdas Normas para TeX
% \end{siglas}
% ---


% ---
% LISTA DE SÍMBOLOS
% ---
% \begin{simbolos}
%   \item[$ \Gamma $] Letra grega Gama
%   \item[$ \Lambda $] Lambda
%   \item[$ \zeta $] Letra grega minúscula zeta
%   \item[$ \in $] Pertence
% \end{simbolos}
% ---


% ---
% Sumário
% ---
\pdfbookmark[0]{\contentsname}{toc}
\tableofcontents*
\cleardoublepage
% ---


% ----------------------------------------------------------
% ELEMENTOS TEXTUAIS
% ----------------------------------------------------------
\textual

\include{./01-introducao/main.tex}

\include{./02-machine-learning/main.tex}

\include{./03-exemplo/main.tex}

\include{./04-conclusao/main.tex}

% ----------------------------------------------------------
% ELEMENTOS PÓS-TEXTUAIS
% ----------------------------------------------------------
\postextual
% ----------------------------------------------------------


% ----------------------------------------------------------
% REFERÊNCIAS
% ----------------------------------------------------------
\bibliography{references.bib}


% ----------------------------------------------------------
% GLOSSÁRIO
% ----------------------------------------------------------
%
% Consulte o manual da classe abntex2 para orientações sobre o glossário.
%
%\glossary

% ----------------------------------------------------------
% APÊNDICES
% ----------------------------------------------------------

% ---
% Inicia os apêndices
% ---
% \begin{apendicesenv}

% % Imprime uma página indicando o início dos apêndices
% \partapendices

% % ----------------------------------------------------------
% \chapter{Quisque libero justo}
% % ----------------------------------------------------------

% % ----------------------------------------------------------
% \chapter{Nullam elementum urna}
% % ----------------------------------------------------------

% \end{apendicesenv}
% ---


% ----------------------------------------------------------
% ANEXOS
% ----------------------------------------------------------

% ---
% Inicia os anexos
% ---
% \begin{anexosenv}

% % Imprime uma página indicando o início dos anexos
% \partanexos

% % ---
% \chapter{Morbi ultrices rutrum lorem}
% % ---

% % ---
% \chapter{Cras non urna sed feugiat cum}
% % ---


% % ---
% \chapter{Fusce facilisis lacinia dui}
% % ---

% \end{anexosenv}

%---------------------------------------------------------------------
% INDICE REMISSIVO
%---------------------------------------------------------------------
% \phantompart
% \printindex
%---------------------------------------------------------------------
\includepdf{./05-extras/contracapa.pdf}
\end{document}


\documentclass[
	% Opções da classe memoir
	12pt,				        % tamanho da fonte
	openany,			        % capítulos começam em pág ímpar (insere página vazia caso preciso)
	oneside,			        % para impressão em recto e verso. Oposto a oneside
	a4paper,			        % tamanho do papel. 
	% Opções da classe abntex2
	%chapter=TITLE,		        % títulos de capítulos convertidos em letras maiúsculas
	%section=TITLE,		        % títulos de seções convertidos em letras maiúsculas
	%subsection=TITLE,	        % títulos de subseções convertidos em letras maiúsculas
	%subsubsection=TITLE,       % títulos de subsubseções convertidos em letras maiúsculas
	% Opções do pacote babel
	english,			        % idioma adicional para hifenização
	brazil				        % o último idioma é o principal do documento
	]{abntex2}

% ---
% Pacotes básicos 
% ---
\usepackage{lmodern}			% Usa a fonte Latin Modern			
\usepackage[T1]{fontenc}		% Selecao de codigos de fonte.
\usepackage[utf8]{inputenc}		% Codificacao do documento (conversão automática dos acentos)
\usepackage{indentfirst}		% Indenta o primeiro parágrafo de cada seção.
\usepackage{color}				% Controle das cores
\usepackage{graphicx}			% Inclusão de gráficos
\usepackage{microtype} 			% para melhorias de justificação
% ---
		

% ---
% Pacotes de citações
% ---
% \usepackage[alf]{abntex2cite}						% Citações padrão ABNT
\usepackage[
  alf,
  abnt-nbr10520=1988,                     % Citações em sentence case
  abnt-emphasize=bf
  ]{abntex2cite}


% ---
% Pacotes adicionais
% ---
\usepackage{pdfpages}	% Inclusão de páginas em PDF


% ---
% Informações para CAPA e FOLHA DE ROSTO
% ---
\titulo{APRENDIZAGEM DE MÁQUINA PARA PREDIÇÃO DE DESFECHOS EM SAÚDE MENTAL: \textmd{UMA INTRODUÇÃO PARA PSICÓLOGOS}}
\autor{GUILHERME AUGUSTO DIAS}
\local{Porto Alegre}
\data{2024}
\orientador{Prof. Dr. Wagner de Lara Machado}
% \coorientador{}
% \instituicao{
%   Pontifícia Universidade Católica do Rio Grande do Sul
%   \par
%   Escola de Ciências da Saúde e da Vida
%   \par
%   Curso de Psicologia}
\tipotrabalho{TCC}
\preambulo{Trabalho de Conclusão de Curso apresentado como requisito para a obtenção do grau de Bacharel em Psicologia pela Escola de Ciências da Saúde e da Vida da Pontifícia Universidade Católica do Rio Grande do Sul.}
% ---

% ---
% Configurações de aparência do PDF final
% ---

% Alterando o aspecto da cor azul
\definecolor{blue}{RGB}{41,5,195}

% Informações do PDF
\makeatletter
\hypersetup{
		pdftitle={Aprendizagem de máquina para predição de desfechos em saúde mental: uma introdução para psicólogos}, 
		pdfauthor={Guilherme Augusto Dias},
    	pdfsubject={\imprimirpreambulo},
	    pdfcreator={LaTeX with abnTeX2},
		pdfkeywords={aprendizagem de máquina}{psicoterapia}{psicologia positiva}{depressão}, 
		colorlinks=true,		% true: colored links; false: boxed links
    	linkcolor=black,		% color of internal links
    	citecolor=black,		% color of links to bibliography
    	filecolor=black,		% color of file links
		urlcolor=black,
		bookmarksdepth=4
}
\makeatother
% --- 


% ---
% Posiciona figuras e tabelas no topo da página quando adicionadas sozinhas
% em um página em branco. Ver https://github.com/abntex/abntex2/issues/170
\makeatletter
\setlength{\@fptop}{5pt} % Set distance from top of page to first float
\makeatother
% ---

% ---
% Possibilita criação de Quadros e Lista de quadros.
% Ver https://github.com/abntex/abntex2/issues/176
%
\newcommand{\quadroname}{Quadro}
\newcommand{\listofquadrosname}{Lista de quadros}

\newfloat[chapter]{quadro}{loq}{\quadroname}
\newlistof{listofquadros}{loq}{\listofquadrosname}
\newlistentry{quadro}{loq}{0}

% Configurações para atender às regras da ABNT
\setfloatadjustment{quadro}{\centering}
\counterwithout{quadro}{chapter}
\renewcommand{\cftquadroname}{\quadroname\space} 
\renewcommand*{\cftquadroaftersnum}{\hfill--\hfill}

\setfloatlocations{quadro}{hbtp} % Ver https://github.com/abntex/abntex2/issues/176
% ---


% --- 
% Espaçamentos entre linhas e parágrafos 
% --- 

% O tamanho do parágrafo é dado por:
\setlength{\parindent}{1.3cm}

% Controle do espaçamento entre um parágrafo e outro:
\setlength{\parskip}{0.2cm}  % tente também \onelineskip

% ---
% Compila o índice
% ---
\makeindex
% ---


% ----
% Início do documento
% ----
\begin{document}

% Seleciona o idioma do documento (conforme pacotes do babel)
\selectlanguage{brazil}

% Retira espaço extra obsoleto entre as frases.
\frenchspacing 

% ----------------------------------------------------------
% ELEMENTOS PRÉ-TEXTUAIS
% ----------------------------------------------------------
\pretextual

% ---
% CAPA
% ---
\includepdf{./05-extras/capa.pdf}
% ---


% ---
% FOLHA DE ROSTO
% (o * indica que haverá a ficha bibliográfica)
% ---
\imprimirfolhaderosto
% ---


% ---
% Ficha bibliografica
% ---

% Isto é um exemplo de Ficha Catalográfica, ou ``Dados internacionais de
% catalogação-na-publicação''. Você pode utilizar este modelo como referência. 
% Porém, provavelmente a biblioteca da sua universidade lhe fornecerá um PDF
% com a ficha catalográfica definitiva após a defesa do trabalho. Quando estiver
% com o documento, salve-o como PDF no diretório do seu projeto e substitua todo
% o conteúdo de implementação deste arquivo pelo comando abaixo:
%
% \usepackage{pdfpages}		% necessário para comando \includepdf
% \begin{fichacatalografica}
%     \includepdf{fig_ficha_catalografica.pdf}
% \end{fichacatalografica}

% \begin{fichacatalografica}
% 	\sffamily
% 	\vspace*{\fill}                             % Posição vertical
% 	\begin{center}                              % Minipage Centralizado
% 	\fbox{\begin{minipage}[c][8cm]{13.5cm}		% Largura
% 	\small
% 	\imprimirautor
% 	%Sobrenome, Nome do autor
	
% 	\hspace{0.5cm} \imprimirtitulo  / \imprimirautor. --
% 	\imprimirlocal, \imprimirdata-
	
% 	\hspace{0.5cm} \thelastpage p. : il. (algumas color.) ; 30 cm.\\
	
% 	\hspace{0.5cm} \imprimirorientadorRotulo~\imprimirorientador\\
	
% 	\hspace{0.5cm}
% 	\parbox[t]{\textwidth}{\imprimirtipotrabalho~--~\imprimirinstituicao,
% 	\imprimirdata.}\\
	
% 	\hspace{0.5cm}
% 		1. Palavra-chave1.
% 		2. Palavra-chave2.
% 		2. Palavra-chave3.
% 		I. Orientador.
% 		II. Universidade xxx.
% 		III. Faculdade de xxx.
% 		IV. Título 			
% 	\end{minipage}}
% 	\end{center}
% \end{fichacatalografica}
% ---


% ---
% FOLHA DE APROVAÇÃO
% ---

% Isto é um exemplo de Folha de aprovação, elemento obrigatório da NBR
% 14724/2011 (seção 4.2.1.3). Você pode utilizar este modelo até a aprovação
% do trabalho. Após isso, substitua todo o conteúdo deste arquivo por uma
% imagem da página assinada pela banca com o comando abaixo:
%
% \begin{folhadeaprovacao}
% \includepdf{folhadeaprovacao_final.pdf}
% \end{folhadeaprovacao}
%
\setlength{\ABNTEXsignskip}{2cm}
\setlength{\ABNTEXsignwidth}{10cm}

\begin{folhadeaprovacao}
  \begin{center}
    {\ABNTEXchapterfont\large\imprimirautor}

    \vspace*{1cm}

    \begin{center}
      \ABNTEXchapterfont\bfseries\Large\imprimirtitulo
    \end{center}

    \vspace*{1cm}
    
    \hspace{.45\textwidth}
    \begin{minipage}{.5\textwidth}
        \imprimirpreambulo
    \end{minipage}%
   \end{center}

   \vspace*{1cm}
        
   \begin{center}
   Aprovado em 17 de julho de 2024.

   \vspace*{1cm}

   Banca examinadora:
   \end{center}

   \vspace*{-1cm}

   \assinatura{Prof. Ma. Caroline Santa Maria Rodrigues -- PUCRS}
   \assinatura{Prof. Dra. Clarissa Pinto Pizarro de Freitas -- PUCRS}
   \assinatura{\imprimirorientador -- PUCRS}
\end{folhadeaprovacao}
% ---


% ---
% Agradecimentos
% ---
\begin{agradecimentos}
Agradeço aos meus pais, Ademilson José Dias e Irinês Porfírio Dias, pelo apoio incondicional que me deram ao longo da vida,
incentivando mesmo as ideias questionáveis como aquela de fazer uma segunda graduação.

A minha irmã, Maria Eduarda Dias, cuja trajetória acadêmica é fonte de inspiração.

A minha namorada, Nathalia Fallavena Ceratti, que acompanhou a esta jornada de perto, oferecendo apoio e compreensão que
fizeram toda a diferença.

Aos meus amigos, Guilherme da Silva Petró, Norton Neves e Uilinton Luciano Marques da Silva, que foram compreensíveis nos
finais de semana em que não pude jogar RPG.

Ao professor Wagner de Lara Machado pela orientação e apoio no desenvolvimento deste trabalho.
\end{agradecimentos}
% ---


% ---
% RESUMOS
% ---
% Resumo em português
\setlength{\absparsep}{18pt} % ajusta o espaçamento dos parágrafos do resumo
\begin{resumo}
Desenvolvimentos na área de aprendizagem de máquina representam uma série de oportunidades para diversas áreas de pesquisa e de atuação profissional.
O uso de modelos de aprendizagem de máquina na predição de desfecho para tratamentos em saúde mental tem sido investigado e apresenta resultados promissores.
A aplicação desse tipo de modelo possibilitaria, por exemplo, maior assertividade na personalização de tratamentos psicoterápicos a nível individual.
Apesar dos benefícios potenciais, o tema ainda é pouco discutido nos espaços de pesquisa, formação e atuação em psicologia clínica. Este trabalho tem
por objetivo introduzir o tema de aprendizagem de máquina a psicólogos. Busca-se apresentar, de maneira acessível, os conceitos e terminologia básicos
em aprendizagem de máquina, proporcionando um panorama geral da área. Para complementar a exposição conceitual, é construído um exemplo a partir de dados
de um estudo sobre eficácia de intervenções em psicologia positiva para depressão via internet. Por fim, discutem-se algumas implicações do uso desse tipo
de tecnologia para a pesquisa e para a prática em psicoterapia.

 \textbf{Palavras-chave}: aprendizagem de máquina, psicoterapia, psicologia positiva, depressão.
\end{resumo}

% Resumo em inglês
\begin{resumo}[Abstract]
 \begin{otherlanguage*}{english}
   Advances in machine learning have opened up a multitude of opportunities for various research and professional fields. The use of machine learning models to
   predict treatment outcomes in mental health has been investigated showing promising results. The application of such models would allow greater assertiveness
   in tailoring psychotherapeutic treatments at the individual level. Despite its potential benefits, the topic remains unexplored in the fields of research, training
   and practice in clinical psychology. This work aims to introduce the topic of machine learning to psychologists. It presents the basic machine learning concepts and
   terminology in an accessible manner, providing a comprehensive overview of the field. To complement the conceptual discussion, an example is constructed using data
   from a study on the effectiveness of internet-based positive psychology interventions for depression. Finally, implications of using this type of technology in
   psychotherapy research and practice are discussed.
   \vspace{\onelineskip}
 
   \noindent 
   \textbf{Keywords}: machine learning, psychotherapy, positive psychology, depression.
 \end{otherlanguage*}
\end{resumo}


% ---
% LISTA DE ILUSTRAÇÕES
% ---
\pdfbookmark[0]{\listfigurename}{lof}
\listoffigures*
\cleardoublepage
% ---


% ---
% LISTA DE QUADROS
% ---
% \pdfbookmark[0]{\listofquadrosname}{loq}
% \listofquadros*
% \cleardoublepage
% ---


% ---
% LISTA DE TABELAS
% ---
% \pdfbookmark[0]{\listtablename}{lot}
% \listoftables*
% \cleardoublepage
% ---


% ---
% LISTA DE ABREVIATURAS E SIGLAS
% ---
% \begin{siglas}
%   \item[ABNT] Associação Brasileira de Normas Técnicas
%   \item[abnTeX] ABsurdas Normas para TeX
% \end{siglas}
% ---


% ---
% LISTA DE SÍMBOLOS
% ---
% \begin{simbolos}
%   \item[$ \Gamma $] Letra grega Gama
%   \item[$ \Lambda $] Lambda
%   \item[$ \zeta $] Letra grega minúscula zeta
%   \item[$ \in $] Pertence
% \end{simbolos}
% ---


% ---
% Sumário
% ---
\pdfbookmark[0]{\contentsname}{toc}
\tableofcontents*
\cleardoublepage
% ---


% ----------------------------------------------------------
% ELEMENTOS TEXTUAIS
% ----------------------------------------------------------
\textual

\include{./01-introducao/main.tex}

\include{./02-machine-learning/main.tex}

\include{./03-exemplo/main.tex}

\include{./04-conclusao/main.tex}

% ----------------------------------------------------------
% ELEMENTOS PÓS-TEXTUAIS
% ----------------------------------------------------------
\postextual
% ----------------------------------------------------------


% ----------------------------------------------------------
% REFERÊNCIAS
% ----------------------------------------------------------
\bibliography{references.bib}


% ----------------------------------------------------------
% GLOSSÁRIO
% ----------------------------------------------------------
%
% Consulte o manual da classe abntex2 para orientações sobre o glossário.
%
%\glossary

% ----------------------------------------------------------
% APÊNDICES
% ----------------------------------------------------------

% ---
% Inicia os apêndices
% ---
% \begin{apendicesenv}

% % Imprime uma página indicando o início dos apêndices
% \partapendices

% % ----------------------------------------------------------
% \chapter{Quisque libero justo}
% % ----------------------------------------------------------

% % ----------------------------------------------------------
% \chapter{Nullam elementum urna}
% % ----------------------------------------------------------

% \end{apendicesenv}
% ---


% ----------------------------------------------------------
% ANEXOS
% ----------------------------------------------------------

% ---
% Inicia os anexos
% ---
% \begin{anexosenv}

% % Imprime uma página indicando o início dos anexos
% \partanexos

% % ---
% \chapter{Morbi ultrices rutrum lorem}
% % ---

% % ---
% \chapter{Cras non urna sed feugiat cum}
% % ---


% % ---
% \chapter{Fusce facilisis lacinia dui}
% % ---

% \end{anexosenv}

%---------------------------------------------------------------------
% INDICE REMISSIVO
%---------------------------------------------------------------------
% \phantompart
% \printindex
%---------------------------------------------------------------------
\includepdf{./05-extras/contracapa.pdf}
\end{document}


% ----------------------------------------------------------
% ELEMENTOS PÓS-TEXTUAIS
% ----------------------------------------------------------
\postextual
% ----------------------------------------------------------


% ----------------------------------------------------------
% REFERÊNCIAS
% ----------------------------------------------------------
\bibliography{references.bib}


% ----------------------------------------------------------
% GLOSSÁRIO
% ----------------------------------------------------------
%
% Consulte o manual da classe abntex2 para orientações sobre o glossário.
%
%\glossary

% ----------------------------------------------------------
% APÊNDICES
% ----------------------------------------------------------

% ---
% Inicia os apêndices
% ---
% \begin{apendicesenv}

% % Imprime uma página indicando o início dos apêndices
% \partapendices

% % ----------------------------------------------------------
% \chapter{Quisque libero justo}
% % ----------------------------------------------------------

% % ----------------------------------------------------------
% \chapter{Nullam elementum urna}
% % ----------------------------------------------------------

% \end{apendicesenv}
% ---


% ----------------------------------------------------------
% ANEXOS
% ----------------------------------------------------------

% ---
% Inicia os anexos
% ---
% \begin{anexosenv}

% % Imprime uma página indicando o início dos anexos
% \partanexos

% % ---
% \chapter{Morbi ultrices rutrum lorem}
% % ---

% % ---
% \chapter{Cras non urna sed feugiat cum}
% % ---


% % ---
% \chapter{Fusce facilisis lacinia dui}
% % ---

% \end{anexosenv}

%---------------------------------------------------------------------
% INDICE REMISSIVO
%---------------------------------------------------------------------
% \phantompart
% \printindex
%---------------------------------------------------------------------
\includepdf{./05-extras/contracapa.pdf}
\end{document}


% ----------------------------------------------------------
% ELEMENTOS PÓS-TEXTUAIS
% ----------------------------------------------------------
\postextual
% ----------------------------------------------------------


% ----------------------------------------------------------
% REFERÊNCIAS
% ----------------------------------------------------------
\bibliography{references.bib}


% ----------------------------------------------------------
% GLOSSÁRIO
% ----------------------------------------------------------
%
% Consulte o manual da classe abntex2 para orientações sobre o glossário.
%
%\glossary

% ----------------------------------------------------------
% APÊNDICES
% ----------------------------------------------------------

% ---
% Inicia os apêndices
% ---
% \begin{apendicesenv}

% % Imprime uma página indicando o início dos apêndices
% \partapendices

% % ----------------------------------------------------------
% \chapter{Quisque libero justo}
% % ----------------------------------------------------------

% % ----------------------------------------------------------
% \chapter{Nullam elementum urna}
% % ----------------------------------------------------------

% \end{apendicesenv}
% ---


% ----------------------------------------------------------
% ANEXOS
% ----------------------------------------------------------

% ---
% Inicia os anexos
% ---
% \begin{anexosenv}

% % Imprime uma página indicando o início dos anexos
% \partanexos

% % ---
% \chapter{Morbi ultrices rutrum lorem}
% % ---

% % ---
% \chapter{Cras non urna sed feugiat cum}
% % ---


% % ---
% \chapter{Fusce facilisis lacinia dui}
% % ---

% \end{anexosenv}

%---------------------------------------------------------------------
% INDICE REMISSIVO
%---------------------------------------------------------------------
% \phantompart
% \printindex
%---------------------------------------------------------------------
\includepdf{./05-extras/contracapa.pdf}
\end{document}


\documentclass[
	% Opções da classe memoir
	12pt,				        % tamanho da fonte
	openany,			        % capítulos começam em pág ímpar (insere página vazia caso preciso)
	oneside,			        % para impressão em recto e verso. Oposto a oneside
	a4paper,			        % tamanho do papel. 
	% Opções da classe abntex2
	%chapter=TITLE,		        % títulos de capítulos convertidos em letras maiúsculas
	%section=TITLE,		        % títulos de seções convertidos em letras maiúsculas
	%subsection=TITLE,	        % títulos de subseções convertidos em letras maiúsculas
	%subsubsection=TITLE,       % títulos de subsubseções convertidos em letras maiúsculas
	% Opções do pacote babel
	english,			        % idioma adicional para hifenização
	brazil				        % o último idioma é o principal do documento
	]{abntex2}

% ---
% Pacotes básicos 
% ---
\usepackage{lmodern}			% Usa a fonte Latin Modern			
\usepackage[T1]{fontenc}		% Selecao de codigos de fonte.
\usepackage[utf8]{inputenc}		% Codificacao do documento (conversão automática dos acentos)
\usepackage{indentfirst}		% Indenta o primeiro parágrafo de cada seção.
\usepackage{color}				% Controle das cores
\usepackage{graphicx}			% Inclusão de gráficos
\usepackage{microtype} 			% para melhorias de justificação
% ---
		

% ---
% Pacotes de citações
% ---
% \usepackage[alf]{abntex2cite}						% Citações padrão ABNT
\usepackage[
  alf,
  abnt-nbr10520=1988,                     % Citações em sentence case
  abnt-emphasize=bf
  ]{abntex2cite}


% ---
% Pacotes adicionais
% ---
\usepackage{pdfpages}	% Inclusão de páginas em PDF


% ---
% Informações para CAPA e FOLHA DE ROSTO
% ---
\titulo{APRENDIZAGEM DE MÁQUINA PARA PREDIÇÃO DE DESFECHOS EM SAÚDE MENTAL: \textmd{UMA INTRODUÇÃO PARA PSICÓLOGOS}}
\autor{GUILHERME AUGUSTO DIAS}
\local{Porto Alegre}
\data{2024}
\orientador{Prof. Dr. Wagner de Lara Machado}
% \coorientador{}
% \instituicao{
%   Pontifícia Universidade Católica do Rio Grande do Sul
%   \par
%   Escola de Ciências da Saúde e da Vida
%   \par
%   Curso de Psicologia}
\tipotrabalho{TCC}
\preambulo{Trabalho de Conclusão de Curso apresentado como requisito para a obtenção do grau de Bacharel em Psicologia pela Escola de Ciências da Saúde e da Vida da Pontifícia Universidade Católica do Rio Grande do Sul.}
% ---

% ---
% Configurações de aparência do PDF final
% ---

% Alterando o aspecto da cor azul
\definecolor{blue}{RGB}{41,5,195}

% Informações do PDF
\makeatletter
\hypersetup{
		pdftitle={Aprendizagem de máquina para predição de desfechos em saúde mental: uma introdução para psicólogos}, 
		pdfauthor={Guilherme Augusto Dias},
    	pdfsubject={\imprimirpreambulo},
	    pdfcreator={LaTeX with abnTeX2},
		pdfkeywords={aprendizagem de máquina}{psicoterapia}{psicologia positiva}{depressão}, 
		colorlinks=true,		% true: colored links; false: boxed links
    	linkcolor=black,		% color of internal links
    	citecolor=black,		% color of links to bibliography
    	filecolor=black,		% color of file links
		urlcolor=black,
		bookmarksdepth=4
}
\makeatother
% --- 


% ---
% Posiciona figuras e tabelas no topo da página quando adicionadas sozinhas
% em um página em branco. Ver https://github.com/abntex/abntex2/issues/170
\makeatletter
\setlength{\@fptop}{5pt} % Set distance from top of page to first float
\makeatother
% ---

% ---
% Possibilita criação de Quadros e Lista de quadros.
% Ver https://github.com/abntex/abntex2/issues/176
%
\newcommand{\quadroname}{Quadro}
\newcommand{\listofquadrosname}{Lista de quadros}

\newfloat[chapter]{quadro}{loq}{\quadroname}
\newlistof{listofquadros}{loq}{\listofquadrosname}
\newlistentry{quadro}{loq}{0}

% Configurações para atender às regras da ABNT
\setfloatadjustment{quadro}{\centering}
\counterwithout{quadro}{chapter}
\renewcommand{\cftquadroname}{\quadroname\space} 
\renewcommand*{\cftquadroaftersnum}{\hfill--\hfill}

\setfloatlocations{quadro}{hbtp} % Ver https://github.com/abntex/abntex2/issues/176
% ---


% --- 
% Espaçamentos entre linhas e parágrafos 
% --- 

% O tamanho do parágrafo é dado por:
\setlength{\parindent}{1.3cm}

% Controle do espaçamento entre um parágrafo e outro:
\setlength{\parskip}{0.2cm}  % tente também \onelineskip

% ---
% Compila o índice
% ---
\makeindex
% ---


% ----
% Início do documento
% ----
\begin{document}

% Seleciona o idioma do documento (conforme pacotes do babel)
\selectlanguage{brazil}

% Retira espaço extra obsoleto entre as frases.
\frenchspacing 

% ----------------------------------------------------------
% ELEMENTOS PRÉ-TEXTUAIS
% ----------------------------------------------------------
\pretextual

% ---
% CAPA
% ---
\includepdf{./05-extras/capa.pdf}
% ---


% ---
% FOLHA DE ROSTO
% (o * indica que haverá a ficha bibliográfica)
% ---
\imprimirfolhaderosto
% ---


% ---
% Ficha bibliografica
% ---

% Isto é um exemplo de Ficha Catalográfica, ou ``Dados internacionais de
% catalogação-na-publicação''. Você pode utilizar este modelo como referência. 
% Porém, provavelmente a biblioteca da sua universidade lhe fornecerá um PDF
% com a ficha catalográfica definitiva após a defesa do trabalho. Quando estiver
% com o documento, salve-o como PDF no diretório do seu projeto e substitua todo
% o conteúdo de implementação deste arquivo pelo comando abaixo:
%
% \usepackage{pdfpages}		% necessário para comando \includepdf
% \begin{fichacatalografica}
%     \includepdf{fig_ficha_catalografica.pdf}
% \end{fichacatalografica}

% \begin{fichacatalografica}
% 	\sffamily
% 	\vspace*{\fill}                             % Posição vertical
% 	\begin{center}                              % Minipage Centralizado
% 	\fbox{\begin{minipage}[c][8cm]{13.5cm}		% Largura
% 	\small
% 	\imprimirautor
% 	%Sobrenome, Nome do autor
	
% 	\hspace{0.5cm} \imprimirtitulo  / \imprimirautor. --
% 	\imprimirlocal, \imprimirdata-
	
% 	\hspace{0.5cm} \thelastpage p. : il. (algumas color.) ; 30 cm.\\
	
% 	\hspace{0.5cm} \imprimirorientadorRotulo~\imprimirorientador\\
	
% 	\hspace{0.5cm}
% 	\parbox[t]{\textwidth}{\imprimirtipotrabalho~--~\imprimirinstituicao,
% 	\imprimirdata.}\\
	
% 	\hspace{0.5cm}
% 		1. Palavra-chave1.
% 		2. Palavra-chave2.
% 		2. Palavra-chave3.
% 		I. Orientador.
% 		II. Universidade xxx.
% 		III. Faculdade de xxx.
% 		IV. Título 			
% 	\end{minipage}}
% 	\end{center}
% \end{fichacatalografica}
% ---


% ---
% FOLHA DE APROVAÇÃO
% ---

% Isto é um exemplo de Folha de aprovação, elemento obrigatório da NBR
% 14724/2011 (seção 4.2.1.3). Você pode utilizar este modelo até a aprovação
% do trabalho. Após isso, substitua todo o conteúdo deste arquivo por uma
% imagem da página assinada pela banca com o comando abaixo:
%
% \begin{folhadeaprovacao}
% \includepdf{folhadeaprovacao_final.pdf}
% \end{folhadeaprovacao}
%
\setlength{\ABNTEXsignskip}{2cm}
\setlength{\ABNTEXsignwidth}{10cm}

\begin{folhadeaprovacao}
  \begin{center}
    {\ABNTEXchapterfont\large\imprimirautor}

    \vspace*{1cm}

    \begin{center}
      \ABNTEXchapterfont\bfseries\Large\imprimirtitulo
    \end{center}

    \vspace*{1cm}
    
    \hspace{.45\textwidth}
    \begin{minipage}{.5\textwidth}
        \imprimirpreambulo
    \end{minipage}%
   \end{center}

   \vspace*{1cm}
        
   \begin{center}
   Aprovado em 17 de julho de 2024.

   \vspace*{1cm}

   Banca examinadora:
   \end{center}

   \vspace*{-1cm}

   \assinatura{Prof. Ma. Caroline Santa Maria Rodrigues -- PUCRS}
   \assinatura{Prof. Dra. Clarissa Pinto Pizarro de Freitas -- PUCRS}
   \assinatura{\imprimirorientador -- PUCRS}
\end{folhadeaprovacao}
% ---


% ---
% Agradecimentos
% ---
\begin{agradecimentos}
Agradeço aos meus pais, Ademilson José Dias e Irinês Porfírio Dias, pelo apoio incondicional que me deram ao longo da vida,
incentivando mesmo as ideias questionáveis como aquela de fazer uma segunda graduação.

A minha irmã, Maria Eduarda Dias, cuja trajetória acadêmica é fonte de inspiração.

A minha namorada, Nathalia Fallavena Ceratti, que acompanhou a esta jornada de perto, oferecendo apoio e compreensão que
fizeram toda a diferença.

Aos meus amigos, Guilherme da Silva Petró, Norton Neves e Uilinton Luciano Marques da Silva, que foram compreensíveis nos
finais de semana em que não pude jogar RPG.

Ao professor Wagner de Lara Machado pela orientação e apoio no desenvolvimento deste trabalho.
\end{agradecimentos}
% ---


% ---
% RESUMOS
% ---
% Resumo em português
\setlength{\absparsep}{18pt} % ajusta o espaçamento dos parágrafos do resumo
\begin{resumo}
Desenvolvimentos na área de aprendizagem de máquina representam uma série de oportunidades para diversas áreas de pesquisa e de atuação profissional.
O uso de modelos de aprendizagem de máquina na predição de desfecho para tratamentos em saúde mental tem sido investigado e apresenta resultados promissores.
A aplicação desse tipo de modelo possibilitaria, por exemplo, maior assertividade na personalização de tratamentos psicoterápicos a nível individual.
Apesar dos benefícios potenciais, o tema ainda é pouco discutido nos espaços de pesquisa, formação e atuação em psicologia clínica. Este trabalho tem
por objetivo introduzir o tema de aprendizagem de máquina a psicólogos. Busca-se apresentar, de maneira acessível, os conceitos e terminologia básicos
em aprendizagem de máquina, proporcionando um panorama geral da área. Para complementar a exposição conceitual, é construído um exemplo a partir de dados
de um estudo sobre eficácia de intervenções em psicologia positiva para depressão via internet. Por fim, discutem-se algumas implicações do uso desse tipo
de tecnologia para a pesquisa e para a prática em psicoterapia.

 \textbf{Palavras-chave}: aprendizagem de máquina, psicoterapia, psicologia positiva, depressão.
\end{resumo}

% Resumo em inglês
\begin{resumo}[Abstract]
 \begin{otherlanguage*}{english}
   Advances in machine learning have opened up a multitude of opportunities for various research and professional fields. The use of machine learning models to
   predict treatment outcomes in mental health has been investigated showing promising results. The application of such models would allow greater assertiveness
   in tailoring psychotherapeutic treatments at the individual level. Despite its potential benefits, the topic remains unexplored in the fields of research, training
   and practice in clinical psychology. This work aims to introduce the topic of machine learning to psychologists. It presents the basic machine learning concepts and
   terminology in an accessible manner, providing a comprehensive overview of the field. To complement the conceptual discussion, an example is constructed using data
   from a study on the effectiveness of internet-based positive psychology interventions for depression. Finally, implications of using this type of technology in
   psychotherapy research and practice are discussed.
   \vspace{\onelineskip}
 
   \noindent 
   \textbf{Keywords}: machine learning, psychotherapy, positive psychology, depression.
 \end{otherlanguage*}
\end{resumo}


% ---
% LISTA DE ILUSTRAÇÕES
% ---
\pdfbookmark[0]{\listfigurename}{lof}
\listoffigures*
\cleardoublepage
% ---


% ---
% LISTA DE QUADROS
% ---
% \pdfbookmark[0]{\listofquadrosname}{loq}
% \listofquadros*
% \cleardoublepage
% ---


% ---
% LISTA DE TABELAS
% ---
% \pdfbookmark[0]{\listtablename}{lot}
% \listoftables*
% \cleardoublepage
% ---


% ---
% LISTA DE ABREVIATURAS E SIGLAS
% ---
% \begin{siglas}
%   \item[ABNT] Associação Brasileira de Normas Técnicas
%   \item[abnTeX] ABsurdas Normas para TeX
% \end{siglas}
% ---


% ---
% LISTA DE SÍMBOLOS
% ---
% \begin{simbolos}
%   \item[$ \Gamma $] Letra grega Gama
%   \item[$ \Lambda $] Lambda
%   \item[$ \zeta $] Letra grega minúscula zeta
%   \item[$ \in $] Pertence
% \end{simbolos}
% ---


% ---
% Sumário
% ---
\pdfbookmark[0]{\contentsname}{toc}
\tableofcontents*
\cleardoublepage
% ---


% ----------------------------------------------------------
% ELEMENTOS TEXTUAIS
% ----------------------------------------------------------
\textual

\documentclass[
	% Opções da classe memoir
	12pt,				        % tamanho da fonte
	openany,			        % capítulos começam em pág ímpar (insere página vazia caso preciso)
	oneside,			        % para impressão em recto e verso. Oposto a oneside
	a4paper,			        % tamanho do papel. 
	% Opções da classe abntex2
	%chapter=TITLE,		        % títulos de capítulos convertidos em letras maiúsculas
	%section=TITLE,		        % títulos de seções convertidos em letras maiúsculas
	%subsection=TITLE,	        % títulos de subseções convertidos em letras maiúsculas
	%subsubsection=TITLE,       % títulos de subsubseções convertidos em letras maiúsculas
	% Opções do pacote babel
	english,			        % idioma adicional para hifenização
	brazil				        % o último idioma é o principal do documento
	]{abntex2}

% ---
% Pacotes básicos 
% ---
\usepackage{lmodern}			% Usa a fonte Latin Modern			
\usepackage[T1]{fontenc}		% Selecao de codigos de fonte.
\usepackage[utf8]{inputenc}		% Codificacao do documento (conversão automática dos acentos)
\usepackage{indentfirst}		% Indenta o primeiro parágrafo de cada seção.
\usepackage{color}				% Controle das cores
\usepackage{graphicx}			% Inclusão de gráficos
\usepackage{microtype} 			% para melhorias de justificação
% ---
		

% ---
% Pacotes de citações
% ---
% \usepackage[alf]{abntex2cite}						% Citações padrão ABNT
\usepackage[
  alf,
  abnt-nbr10520=1988,                     % Citações em sentence case
  abnt-emphasize=bf
  ]{abntex2cite}


% ---
% Pacotes adicionais
% ---
\usepackage{pdfpages}	% Inclusão de páginas em PDF


% ---
% Informações para CAPA e FOLHA DE ROSTO
% ---
\titulo{APRENDIZAGEM DE MÁQUINA PARA PREDIÇÃO DE DESFECHOS EM SAÚDE MENTAL: \textmd{UMA INTRODUÇÃO PARA PSICÓLOGOS}}
\autor{GUILHERME AUGUSTO DIAS}
\local{Porto Alegre}
\data{2024}
\orientador{Prof. Dr. Wagner de Lara Machado}
% \coorientador{}
% \instituicao{
%   Pontifícia Universidade Católica do Rio Grande do Sul
%   \par
%   Escola de Ciências da Saúde e da Vida
%   \par
%   Curso de Psicologia}
\tipotrabalho{TCC}
\preambulo{Trabalho de Conclusão de Curso apresentado como requisito para a obtenção do grau de Bacharel em Psicologia pela Escola de Ciências da Saúde e da Vida da Pontifícia Universidade Católica do Rio Grande do Sul.}
% ---

% ---
% Configurações de aparência do PDF final
% ---

% Alterando o aspecto da cor azul
\definecolor{blue}{RGB}{41,5,195}

% Informações do PDF
\makeatletter
\hypersetup{
		pdftitle={Aprendizagem de máquina para predição de desfechos em saúde mental: uma introdução para psicólogos}, 
		pdfauthor={Guilherme Augusto Dias},
    	pdfsubject={\imprimirpreambulo},
	    pdfcreator={LaTeX with abnTeX2},
		pdfkeywords={aprendizagem de máquina}{psicoterapia}{psicologia positiva}{depressão}, 
		colorlinks=true,		% true: colored links; false: boxed links
    	linkcolor=black,		% color of internal links
    	citecolor=black,		% color of links to bibliography
    	filecolor=black,		% color of file links
		urlcolor=black,
		bookmarksdepth=4
}
\makeatother
% --- 


% ---
% Posiciona figuras e tabelas no topo da página quando adicionadas sozinhas
% em um página em branco. Ver https://github.com/abntex/abntex2/issues/170
\makeatletter
\setlength{\@fptop}{5pt} % Set distance from top of page to first float
\makeatother
% ---

% ---
% Possibilita criação de Quadros e Lista de quadros.
% Ver https://github.com/abntex/abntex2/issues/176
%
\newcommand{\quadroname}{Quadro}
\newcommand{\listofquadrosname}{Lista de quadros}

\newfloat[chapter]{quadro}{loq}{\quadroname}
\newlistof{listofquadros}{loq}{\listofquadrosname}
\newlistentry{quadro}{loq}{0}

% Configurações para atender às regras da ABNT
\setfloatadjustment{quadro}{\centering}
\counterwithout{quadro}{chapter}
\renewcommand{\cftquadroname}{\quadroname\space} 
\renewcommand*{\cftquadroaftersnum}{\hfill--\hfill}

\setfloatlocations{quadro}{hbtp} % Ver https://github.com/abntex/abntex2/issues/176
% ---


% --- 
% Espaçamentos entre linhas e parágrafos 
% --- 

% O tamanho do parágrafo é dado por:
\setlength{\parindent}{1.3cm}

% Controle do espaçamento entre um parágrafo e outro:
\setlength{\parskip}{0.2cm}  % tente também \onelineskip

% ---
% Compila o índice
% ---
\makeindex
% ---


% ----
% Início do documento
% ----
\begin{document}

% Seleciona o idioma do documento (conforme pacotes do babel)
\selectlanguage{brazil}

% Retira espaço extra obsoleto entre as frases.
\frenchspacing 

% ----------------------------------------------------------
% ELEMENTOS PRÉ-TEXTUAIS
% ----------------------------------------------------------
\pretextual

% ---
% CAPA
% ---
\includepdf{./05-extras/capa.pdf}
% ---


% ---
% FOLHA DE ROSTO
% (o * indica que haverá a ficha bibliográfica)
% ---
\imprimirfolhaderosto
% ---


% ---
% Ficha bibliografica
% ---

% Isto é um exemplo de Ficha Catalográfica, ou ``Dados internacionais de
% catalogação-na-publicação''. Você pode utilizar este modelo como referência. 
% Porém, provavelmente a biblioteca da sua universidade lhe fornecerá um PDF
% com a ficha catalográfica definitiva após a defesa do trabalho. Quando estiver
% com o documento, salve-o como PDF no diretório do seu projeto e substitua todo
% o conteúdo de implementação deste arquivo pelo comando abaixo:
%
% \usepackage{pdfpages}		% necessário para comando \includepdf
% \begin{fichacatalografica}
%     \includepdf{fig_ficha_catalografica.pdf}
% \end{fichacatalografica}

% \begin{fichacatalografica}
% 	\sffamily
% 	\vspace*{\fill}                             % Posição vertical
% 	\begin{center}                              % Minipage Centralizado
% 	\fbox{\begin{minipage}[c][8cm]{13.5cm}		% Largura
% 	\small
% 	\imprimirautor
% 	%Sobrenome, Nome do autor
	
% 	\hspace{0.5cm} \imprimirtitulo  / \imprimirautor. --
% 	\imprimirlocal, \imprimirdata-
	
% 	\hspace{0.5cm} \thelastpage p. : il. (algumas color.) ; 30 cm.\\
	
% 	\hspace{0.5cm} \imprimirorientadorRotulo~\imprimirorientador\\
	
% 	\hspace{0.5cm}
% 	\parbox[t]{\textwidth}{\imprimirtipotrabalho~--~\imprimirinstituicao,
% 	\imprimirdata.}\\
	
% 	\hspace{0.5cm}
% 		1. Palavra-chave1.
% 		2. Palavra-chave2.
% 		2. Palavra-chave3.
% 		I. Orientador.
% 		II. Universidade xxx.
% 		III. Faculdade de xxx.
% 		IV. Título 			
% 	\end{minipage}}
% 	\end{center}
% \end{fichacatalografica}
% ---


% ---
% FOLHA DE APROVAÇÃO
% ---

% Isto é um exemplo de Folha de aprovação, elemento obrigatório da NBR
% 14724/2011 (seção 4.2.1.3). Você pode utilizar este modelo até a aprovação
% do trabalho. Após isso, substitua todo o conteúdo deste arquivo por uma
% imagem da página assinada pela banca com o comando abaixo:
%
% \begin{folhadeaprovacao}
% \includepdf{folhadeaprovacao_final.pdf}
% \end{folhadeaprovacao}
%
\setlength{\ABNTEXsignskip}{2cm}
\setlength{\ABNTEXsignwidth}{10cm}

\begin{folhadeaprovacao}
  \begin{center}
    {\ABNTEXchapterfont\large\imprimirautor}

    \vspace*{1cm}

    \begin{center}
      \ABNTEXchapterfont\bfseries\Large\imprimirtitulo
    \end{center}

    \vspace*{1cm}
    
    \hspace{.45\textwidth}
    \begin{minipage}{.5\textwidth}
        \imprimirpreambulo
    \end{minipage}%
   \end{center}

   \vspace*{1cm}
        
   \begin{center}
   Aprovado em 17 de julho de 2024.

   \vspace*{1cm}

   Banca examinadora:
   \end{center}

   \vspace*{-1cm}

   \assinatura{Prof. Ma. Caroline Santa Maria Rodrigues -- PUCRS}
   \assinatura{Prof. Dra. Clarissa Pinto Pizarro de Freitas -- PUCRS}
   \assinatura{\imprimirorientador -- PUCRS}
\end{folhadeaprovacao}
% ---


% ---
% Agradecimentos
% ---
\begin{agradecimentos}
Agradeço aos meus pais, Ademilson José Dias e Irinês Porfírio Dias, pelo apoio incondicional que me deram ao longo da vida,
incentivando mesmo as ideias questionáveis como aquela de fazer uma segunda graduação.

A minha irmã, Maria Eduarda Dias, cuja trajetória acadêmica é fonte de inspiração.

A minha namorada, Nathalia Fallavena Ceratti, que acompanhou a esta jornada de perto, oferecendo apoio e compreensão que
fizeram toda a diferença.

Aos meus amigos, Guilherme da Silva Petró, Norton Neves e Uilinton Luciano Marques da Silva, que foram compreensíveis nos
finais de semana em que não pude jogar RPG.

Ao professor Wagner de Lara Machado pela orientação e apoio no desenvolvimento deste trabalho.
\end{agradecimentos}
% ---


% ---
% RESUMOS
% ---
% Resumo em português
\setlength{\absparsep}{18pt} % ajusta o espaçamento dos parágrafos do resumo
\begin{resumo}
Desenvolvimentos na área de aprendizagem de máquina representam uma série de oportunidades para diversas áreas de pesquisa e de atuação profissional.
O uso de modelos de aprendizagem de máquina na predição de desfecho para tratamentos em saúde mental tem sido investigado e apresenta resultados promissores.
A aplicação desse tipo de modelo possibilitaria, por exemplo, maior assertividade na personalização de tratamentos psicoterápicos a nível individual.
Apesar dos benefícios potenciais, o tema ainda é pouco discutido nos espaços de pesquisa, formação e atuação em psicologia clínica. Este trabalho tem
por objetivo introduzir o tema de aprendizagem de máquina a psicólogos. Busca-se apresentar, de maneira acessível, os conceitos e terminologia básicos
em aprendizagem de máquina, proporcionando um panorama geral da área. Para complementar a exposição conceitual, é construído um exemplo a partir de dados
de um estudo sobre eficácia de intervenções em psicologia positiva para depressão via internet. Por fim, discutem-se algumas implicações do uso desse tipo
de tecnologia para a pesquisa e para a prática em psicoterapia.

 \textbf{Palavras-chave}: aprendizagem de máquina, psicoterapia, psicologia positiva, depressão.
\end{resumo}

% Resumo em inglês
\begin{resumo}[Abstract]
 \begin{otherlanguage*}{english}
   Advances in machine learning have opened up a multitude of opportunities for various research and professional fields. The use of machine learning models to
   predict treatment outcomes in mental health has been investigated showing promising results. The application of such models would allow greater assertiveness
   in tailoring psychotherapeutic treatments at the individual level. Despite its potential benefits, the topic remains unexplored in the fields of research, training
   and practice in clinical psychology. This work aims to introduce the topic of machine learning to psychologists. It presents the basic machine learning concepts and
   terminology in an accessible manner, providing a comprehensive overview of the field. To complement the conceptual discussion, an example is constructed using data
   from a study on the effectiveness of internet-based positive psychology interventions for depression. Finally, implications of using this type of technology in
   psychotherapy research and practice are discussed.
   \vspace{\onelineskip}
 
   \noindent 
   \textbf{Keywords}: machine learning, psychotherapy, positive psychology, depression.
 \end{otherlanguage*}
\end{resumo}


% ---
% LISTA DE ILUSTRAÇÕES
% ---
\pdfbookmark[0]{\listfigurename}{lof}
\listoffigures*
\cleardoublepage
% ---


% ---
% LISTA DE QUADROS
% ---
% \pdfbookmark[0]{\listofquadrosname}{loq}
% \listofquadros*
% \cleardoublepage
% ---


% ---
% LISTA DE TABELAS
% ---
% \pdfbookmark[0]{\listtablename}{lot}
% \listoftables*
% \cleardoublepage
% ---


% ---
% LISTA DE ABREVIATURAS E SIGLAS
% ---
% \begin{siglas}
%   \item[ABNT] Associação Brasileira de Normas Técnicas
%   \item[abnTeX] ABsurdas Normas para TeX
% \end{siglas}
% ---


% ---
% LISTA DE SÍMBOLOS
% ---
% \begin{simbolos}
%   \item[$ \Gamma $] Letra grega Gama
%   \item[$ \Lambda $] Lambda
%   \item[$ \zeta $] Letra grega minúscula zeta
%   \item[$ \in $] Pertence
% \end{simbolos}
% ---


% ---
% Sumário
% ---
\pdfbookmark[0]{\contentsname}{toc}
\tableofcontents*
\cleardoublepage
% ---


% ----------------------------------------------------------
% ELEMENTOS TEXTUAIS
% ----------------------------------------------------------
\textual

\documentclass[
	% Opções da classe memoir
	12pt,				        % tamanho da fonte
	openany,			        % capítulos começam em pág ímpar (insere página vazia caso preciso)
	oneside,			        % para impressão em recto e verso. Oposto a oneside
	a4paper,			        % tamanho do papel. 
	% Opções da classe abntex2
	%chapter=TITLE,		        % títulos de capítulos convertidos em letras maiúsculas
	%section=TITLE,		        % títulos de seções convertidos em letras maiúsculas
	%subsection=TITLE,	        % títulos de subseções convertidos em letras maiúsculas
	%subsubsection=TITLE,       % títulos de subsubseções convertidos em letras maiúsculas
	% Opções do pacote babel
	english,			        % idioma adicional para hifenização
	brazil				        % o último idioma é o principal do documento
	]{abntex2}

% ---
% Pacotes básicos 
% ---
\usepackage{lmodern}			% Usa a fonte Latin Modern			
\usepackage[T1]{fontenc}		% Selecao de codigos de fonte.
\usepackage[utf8]{inputenc}		% Codificacao do documento (conversão automática dos acentos)
\usepackage{indentfirst}		% Indenta o primeiro parágrafo de cada seção.
\usepackage{color}				% Controle das cores
\usepackage{graphicx}			% Inclusão de gráficos
\usepackage{microtype} 			% para melhorias de justificação
% ---
		

% ---
% Pacotes de citações
% ---
% \usepackage[alf]{abntex2cite}						% Citações padrão ABNT
\usepackage[
  alf,
  abnt-nbr10520=1988,                     % Citações em sentence case
  abnt-emphasize=bf
  ]{abntex2cite}


% ---
% Pacotes adicionais
% ---
\usepackage{pdfpages}	% Inclusão de páginas em PDF


% ---
% Informações para CAPA e FOLHA DE ROSTO
% ---
\titulo{APRENDIZAGEM DE MÁQUINA PARA PREDIÇÃO DE DESFECHOS EM SAÚDE MENTAL: \textmd{UMA INTRODUÇÃO PARA PSICÓLOGOS}}
\autor{GUILHERME AUGUSTO DIAS}
\local{Porto Alegre}
\data{2024}
\orientador{Prof. Dr. Wagner de Lara Machado}
% \coorientador{}
% \instituicao{
%   Pontifícia Universidade Católica do Rio Grande do Sul
%   \par
%   Escola de Ciências da Saúde e da Vida
%   \par
%   Curso de Psicologia}
\tipotrabalho{TCC}
\preambulo{Trabalho de Conclusão de Curso apresentado como requisito para a obtenção do grau de Bacharel em Psicologia pela Escola de Ciências da Saúde e da Vida da Pontifícia Universidade Católica do Rio Grande do Sul.}
% ---

% ---
% Configurações de aparência do PDF final
% ---

% Alterando o aspecto da cor azul
\definecolor{blue}{RGB}{41,5,195}

% Informações do PDF
\makeatletter
\hypersetup{
		pdftitle={Aprendizagem de máquina para predição de desfechos em saúde mental: uma introdução para psicólogos}, 
		pdfauthor={Guilherme Augusto Dias},
    	pdfsubject={\imprimirpreambulo},
	    pdfcreator={LaTeX with abnTeX2},
		pdfkeywords={aprendizagem de máquina}{psicoterapia}{psicologia positiva}{depressão}, 
		colorlinks=true,		% true: colored links; false: boxed links
    	linkcolor=black,		% color of internal links
    	citecolor=black,		% color of links to bibliography
    	filecolor=black,		% color of file links
		urlcolor=black,
		bookmarksdepth=4
}
\makeatother
% --- 


% ---
% Posiciona figuras e tabelas no topo da página quando adicionadas sozinhas
% em um página em branco. Ver https://github.com/abntex/abntex2/issues/170
\makeatletter
\setlength{\@fptop}{5pt} % Set distance from top of page to first float
\makeatother
% ---

% ---
% Possibilita criação de Quadros e Lista de quadros.
% Ver https://github.com/abntex/abntex2/issues/176
%
\newcommand{\quadroname}{Quadro}
\newcommand{\listofquadrosname}{Lista de quadros}

\newfloat[chapter]{quadro}{loq}{\quadroname}
\newlistof{listofquadros}{loq}{\listofquadrosname}
\newlistentry{quadro}{loq}{0}

% Configurações para atender às regras da ABNT
\setfloatadjustment{quadro}{\centering}
\counterwithout{quadro}{chapter}
\renewcommand{\cftquadroname}{\quadroname\space} 
\renewcommand*{\cftquadroaftersnum}{\hfill--\hfill}

\setfloatlocations{quadro}{hbtp} % Ver https://github.com/abntex/abntex2/issues/176
% ---


% --- 
% Espaçamentos entre linhas e parágrafos 
% --- 

% O tamanho do parágrafo é dado por:
\setlength{\parindent}{1.3cm}

% Controle do espaçamento entre um parágrafo e outro:
\setlength{\parskip}{0.2cm}  % tente também \onelineskip

% ---
% Compila o índice
% ---
\makeindex
% ---


% ----
% Início do documento
% ----
\begin{document}

% Seleciona o idioma do documento (conforme pacotes do babel)
\selectlanguage{brazil}

% Retira espaço extra obsoleto entre as frases.
\frenchspacing 

% ----------------------------------------------------------
% ELEMENTOS PRÉ-TEXTUAIS
% ----------------------------------------------------------
\pretextual

% ---
% CAPA
% ---
\includepdf{./05-extras/capa.pdf}
% ---


% ---
% FOLHA DE ROSTO
% (o * indica que haverá a ficha bibliográfica)
% ---
\imprimirfolhaderosto
% ---


% ---
% Ficha bibliografica
% ---

% Isto é um exemplo de Ficha Catalográfica, ou ``Dados internacionais de
% catalogação-na-publicação''. Você pode utilizar este modelo como referência. 
% Porém, provavelmente a biblioteca da sua universidade lhe fornecerá um PDF
% com a ficha catalográfica definitiva após a defesa do trabalho. Quando estiver
% com o documento, salve-o como PDF no diretório do seu projeto e substitua todo
% o conteúdo de implementação deste arquivo pelo comando abaixo:
%
% \usepackage{pdfpages}		% necessário para comando \includepdf
% \begin{fichacatalografica}
%     \includepdf{fig_ficha_catalografica.pdf}
% \end{fichacatalografica}

% \begin{fichacatalografica}
% 	\sffamily
% 	\vspace*{\fill}                             % Posição vertical
% 	\begin{center}                              % Minipage Centralizado
% 	\fbox{\begin{minipage}[c][8cm]{13.5cm}		% Largura
% 	\small
% 	\imprimirautor
% 	%Sobrenome, Nome do autor
	
% 	\hspace{0.5cm} \imprimirtitulo  / \imprimirautor. --
% 	\imprimirlocal, \imprimirdata-
	
% 	\hspace{0.5cm} \thelastpage p. : il. (algumas color.) ; 30 cm.\\
	
% 	\hspace{0.5cm} \imprimirorientadorRotulo~\imprimirorientador\\
	
% 	\hspace{0.5cm}
% 	\parbox[t]{\textwidth}{\imprimirtipotrabalho~--~\imprimirinstituicao,
% 	\imprimirdata.}\\
	
% 	\hspace{0.5cm}
% 		1. Palavra-chave1.
% 		2. Palavra-chave2.
% 		2. Palavra-chave3.
% 		I. Orientador.
% 		II. Universidade xxx.
% 		III. Faculdade de xxx.
% 		IV. Título 			
% 	\end{minipage}}
% 	\end{center}
% \end{fichacatalografica}
% ---


% ---
% FOLHA DE APROVAÇÃO
% ---

% Isto é um exemplo de Folha de aprovação, elemento obrigatório da NBR
% 14724/2011 (seção 4.2.1.3). Você pode utilizar este modelo até a aprovação
% do trabalho. Após isso, substitua todo o conteúdo deste arquivo por uma
% imagem da página assinada pela banca com o comando abaixo:
%
% \begin{folhadeaprovacao}
% \includepdf{folhadeaprovacao_final.pdf}
% \end{folhadeaprovacao}
%
\setlength{\ABNTEXsignskip}{2cm}
\setlength{\ABNTEXsignwidth}{10cm}

\begin{folhadeaprovacao}
  \begin{center}
    {\ABNTEXchapterfont\large\imprimirautor}

    \vspace*{1cm}

    \begin{center}
      \ABNTEXchapterfont\bfseries\Large\imprimirtitulo
    \end{center}

    \vspace*{1cm}
    
    \hspace{.45\textwidth}
    \begin{minipage}{.5\textwidth}
        \imprimirpreambulo
    \end{minipage}%
   \end{center}

   \vspace*{1cm}
        
   \begin{center}
   Aprovado em 17 de julho de 2024.

   \vspace*{1cm}

   Banca examinadora:
   \end{center}

   \vspace*{-1cm}

   \assinatura{Prof. Ma. Caroline Santa Maria Rodrigues -- PUCRS}
   \assinatura{Prof. Dra. Clarissa Pinto Pizarro de Freitas -- PUCRS}
   \assinatura{\imprimirorientador -- PUCRS}
\end{folhadeaprovacao}
% ---


% ---
% Agradecimentos
% ---
\begin{agradecimentos}
Agradeço aos meus pais, Ademilson José Dias e Irinês Porfírio Dias, pelo apoio incondicional que me deram ao longo da vida,
incentivando mesmo as ideias questionáveis como aquela de fazer uma segunda graduação.

A minha irmã, Maria Eduarda Dias, cuja trajetória acadêmica é fonte de inspiração.

A minha namorada, Nathalia Fallavena Ceratti, que acompanhou a esta jornada de perto, oferecendo apoio e compreensão que
fizeram toda a diferença.

Aos meus amigos, Guilherme da Silva Petró, Norton Neves e Uilinton Luciano Marques da Silva, que foram compreensíveis nos
finais de semana em que não pude jogar RPG.

Ao professor Wagner de Lara Machado pela orientação e apoio no desenvolvimento deste trabalho.
\end{agradecimentos}
% ---


% ---
% RESUMOS
% ---
% Resumo em português
\setlength{\absparsep}{18pt} % ajusta o espaçamento dos parágrafos do resumo
\begin{resumo}
Desenvolvimentos na área de aprendizagem de máquina representam uma série de oportunidades para diversas áreas de pesquisa e de atuação profissional.
O uso de modelos de aprendizagem de máquina na predição de desfecho para tratamentos em saúde mental tem sido investigado e apresenta resultados promissores.
A aplicação desse tipo de modelo possibilitaria, por exemplo, maior assertividade na personalização de tratamentos psicoterápicos a nível individual.
Apesar dos benefícios potenciais, o tema ainda é pouco discutido nos espaços de pesquisa, formação e atuação em psicologia clínica. Este trabalho tem
por objetivo introduzir o tema de aprendizagem de máquina a psicólogos. Busca-se apresentar, de maneira acessível, os conceitos e terminologia básicos
em aprendizagem de máquina, proporcionando um panorama geral da área. Para complementar a exposição conceitual, é construído um exemplo a partir de dados
de um estudo sobre eficácia de intervenções em psicologia positiva para depressão via internet. Por fim, discutem-se algumas implicações do uso desse tipo
de tecnologia para a pesquisa e para a prática em psicoterapia.

 \textbf{Palavras-chave}: aprendizagem de máquina, psicoterapia, psicologia positiva, depressão.
\end{resumo}

% Resumo em inglês
\begin{resumo}[Abstract]
 \begin{otherlanguage*}{english}
   Advances in machine learning have opened up a multitude of opportunities for various research and professional fields. The use of machine learning models to
   predict treatment outcomes in mental health has been investigated showing promising results. The application of such models would allow greater assertiveness
   in tailoring psychotherapeutic treatments at the individual level. Despite its potential benefits, the topic remains unexplored in the fields of research, training
   and practice in clinical psychology. This work aims to introduce the topic of machine learning to psychologists. It presents the basic machine learning concepts and
   terminology in an accessible manner, providing a comprehensive overview of the field. To complement the conceptual discussion, an example is constructed using data
   from a study on the effectiveness of internet-based positive psychology interventions for depression. Finally, implications of using this type of technology in
   psychotherapy research and practice are discussed.
   \vspace{\onelineskip}
 
   \noindent 
   \textbf{Keywords}: machine learning, psychotherapy, positive psychology, depression.
 \end{otherlanguage*}
\end{resumo}


% ---
% LISTA DE ILUSTRAÇÕES
% ---
\pdfbookmark[0]{\listfigurename}{lof}
\listoffigures*
\cleardoublepage
% ---


% ---
% LISTA DE QUADROS
% ---
% \pdfbookmark[0]{\listofquadrosname}{loq}
% \listofquadros*
% \cleardoublepage
% ---


% ---
% LISTA DE TABELAS
% ---
% \pdfbookmark[0]{\listtablename}{lot}
% \listoftables*
% \cleardoublepage
% ---


% ---
% LISTA DE ABREVIATURAS E SIGLAS
% ---
% \begin{siglas}
%   \item[ABNT] Associação Brasileira de Normas Técnicas
%   \item[abnTeX] ABsurdas Normas para TeX
% \end{siglas}
% ---


% ---
% LISTA DE SÍMBOLOS
% ---
% \begin{simbolos}
%   \item[$ \Gamma $] Letra grega Gama
%   \item[$ \Lambda $] Lambda
%   \item[$ \zeta $] Letra grega minúscula zeta
%   \item[$ \in $] Pertence
% \end{simbolos}
% ---


% ---
% Sumário
% ---
\pdfbookmark[0]{\contentsname}{toc}
\tableofcontents*
\cleardoublepage
% ---


% ----------------------------------------------------------
% ELEMENTOS TEXTUAIS
% ----------------------------------------------------------
\textual

\include{./01-introducao/main.tex}

\include{./02-machine-learning/main.tex}

\include{./03-exemplo/main.tex}

\include{./04-conclusao/main.tex}

% ----------------------------------------------------------
% ELEMENTOS PÓS-TEXTUAIS
% ----------------------------------------------------------
\postextual
% ----------------------------------------------------------


% ----------------------------------------------------------
% REFERÊNCIAS
% ----------------------------------------------------------
\bibliography{references.bib}


% ----------------------------------------------------------
% GLOSSÁRIO
% ----------------------------------------------------------
%
% Consulte o manual da classe abntex2 para orientações sobre o glossário.
%
%\glossary

% ----------------------------------------------------------
% APÊNDICES
% ----------------------------------------------------------

% ---
% Inicia os apêndices
% ---
% \begin{apendicesenv}

% % Imprime uma página indicando o início dos apêndices
% \partapendices

% % ----------------------------------------------------------
% \chapter{Quisque libero justo}
% % ----------------------------------------------------------

% % ----------------------------------------------------------
% \chapter{Nullam elementum urna}
% % ----------------------------------------------------------

% \end{apendicesenv}
% ---


% ----------------------------------------------------------
% ANEXOS
% ----------------------------------------------------------

% ---
% Inicia os anexos
% ---
% \begin{anexosenv}

% % Imprime uma página indicando o início dos anexos
% \partanexos

% % ---
% \chapter{Morbi ultrices rutrum lorem}
% % ---

% % ---
% \chapter{Cras non urna sed feugiat cum}
% % ---


% % ---
% \chapter{Fusce facilisis lacinia dui}
% % ---

% \end{anexosenv}

%---------------------------------------------------------------------
% INDICE REMISSIVO
%---------------------------------------------------------------------
% \phantompart
% \printindex
%---------------------------------------------------------------------
\includepdf{./05-extras/contracapa.pdf}
\end{document}


\documentclass[
	% Opções da classe memoir
	12pt,				        % tamanho da fonte
	openany,			        % capítulos começam em pág ímpar (insere página vazia caso preciso)
	oneside,			        % para impressão em recto e verso. Oposto a oneside
	a4paper,			        % tamanho do papel. 
	% Opções da classe abntex2
	%chapter=TITLE,		        % títulos de capítulos convertidos em letras maiúsculas
	%section=TITLE,		        % títulos de seções convertidos em letras maiúsculas
	%subsection=TITLE,	        % títulos de subseções convertidos em letras maiúsculas
	%subsubsection=TITLE,       % títulos de subsubseções convertidos em letras maiúsculas
	% Opções do pacote babel
	english,			        % idioma adicional para hifenização
	brazil				        % o último idioma é o principal do documento
	]{abntex2}

% ---
% Pacotes básicos 
% ---
\usepackage{lmodern}			% Usa a fonte Latin Modern			
\usepackage[T1]{fontenc}		% Selecao de codigos de fonte.
\usepackage[utf8]{inputenc}		% Codificacao do documento (conversão automática dos acentos)
\usepackage{indentfirst}		% Indenta o primeiro parágrafo de cada seção.
\usepackage{color}				% Controle das cores
\usepackage{graphicx}			% Inclusão de gráficos
\usepackage{microtype} 			% para melhorias de justificação
% ---
		

% ---
% Pacotes de citações
% ---
% \usepackage[alf]{abntex2cite}						% Citações padrão ABNT
\usepackage[
  alf,
  abnt-nbr10520=1988,                     % Citações em sentence case
  abnt-emphasize=bf
  ]{abntex2cite}


% ---
% Pacotes adicionais
% ---
\usepackage{pdfpages}	% Inclusão de páginas em PDF


% ---
% Informações para CAPA e FOLHA DE ROSTO
% ---
\titulo{APRENDIZAGEM DE MÁQUINA PARA PREDIÇÃO DE DESFECHOS EM SAÚDE MENTAL: \textmd{UMA INTRODUÇÃO PARA PSICÓLOGOS}}
\autor{GUILHERME AUGUSTO DIAS}
\local{Porto Alegre}
\data{2024}
\orientador{Prof. Dr. Wagner de Lara Machado}
% \coorientador{}
% \instituicao{
%   Pontifícia Universidade Católica do Rio Grande do Sul
%   \par
%   Escola de Ciências da Saúde e da Vida
%   \par
%   Curso de Psicologia}
\tipotrabalho{TCC}
\preambulo{Trabalho de Conclusão de Curso apresentado como requisito para a obtenção do grau de Bacharel em Psicologia pela Escola de Ciências da Saúde e da Vida da Pontifícia Universidade Católica do Rio Grande do Sul.}
% ---

% ---
% Configurações de aparência do PDF final
% ---

% Alterando o aspecto da cor azul
\definecolor{blue}{RGB}{41,5,195}

% Informações do PDF
\makeatletter
\hypersetup{
		pdftitle={Aprendizagem de máquina para predição de desfechos em saúde mental: uma introdução para psicólogos}, 
		pdfauthor={Guilherme Augusto Dias},
    	pdfsubject={\imprimirpreambulo},
	    pdfcreator={LaTeX with abnTeX2},
		pdfkeywords={aprendizagem de máquina}{psicoterapia}{psicologia positiva}{depressão}, 
		colorlinks=true,		% true: colored links; false: boxed links
    	linkcolor=black,		% color of internal links
    	citecolor=black,		% color of links to bibliography
    	filecolor=black,		% color of file links
		urlcolor=black,
		bookmarksdepth=4
}
\makeatother
% --- 


% ---
% Posiciona figuras e tabelas no topo da página quando adicionadas sozinhas
% em um página em branco. Ver https://github.com/abntex/abntex2/issues/170
\makeatletter
\setlength{\@fptop}{5pt} % Set distance from top of page to first float
\makeatother
% ---

% ---
% Possibilita criação de Quadros e Lista de quadros.
% Ver https://github.com/abntex/abntex2/issues/176
%
\newcommand{\quadroname}{Quadro}
\newcommand{\listofquadrosname}{Lista de quadros}

\newfloat[chapter]{quadro}{loq}{\quadroname}
\newlistof{listofquadros}{loq}{\listofquadrosname}
\newlistentry{quadro}{loq}{0}

% Configurações para atender às regras da ABNT
\setfloatadjustment{quadro}{\centering}
\counterwithout{quadro}{chapter}
\renewcommand{\cftquadroname}{\quadroname\space} 
\renewcommand*{\cftquadroaftersnum}{\hfill--\hfill}

\setfloatlocations{quadro}{hbtp} % Ver https://github.com/abntex/abntex2/issues/176
% ---


% --- 
% Espaçamentos entre linhas e parágrafos 
% --- 

% O tamanho do parágrafo é dado por:
\setlength{\parindent}{1.3cm}

% Controle do espaçamento entre um parágrafo e outro:
\setlength{\parskip}{0.2cm}  % tente também \onelineskip

% ---
% Compila o índice
% ---
\makeindex
% ---


% ----
% Início do documento
% ----
\begin{document}

% Seleciona o idioma do documento (conforme pacotes do babel)
\selectlanguage{brazil}

% Retira espaço extra obsoleto entre as frases.
\frenchspacing 

% ----------------------------------------------------------
% ELEMENTOS PRÉ-TEXTUAIS
% ----------------------------------------------------------
\pretextual

% ---
% CAPA
% ---
\includepdf{./05-extras/capa.pdf}
% ---


% ---
% FOLHA DE ROSTO
% (o * indica que haverá a ficha bibliográfica)
% ---
\imprimirfolhaderosto
% ---


% ---
% Ficha bibliografica
% ---

% Isto é um exemplo de Ficha Catalográfica, ou ``Dados internacionais de
% catalogação-na-publicação''. Você pode utilizar este modelo como referência. 
% Porém, provavelmente a biblioteca da sua universidade lhe fornecerá um PDF
% com a ficha catalográfica definitiva após a defesa do trabalho. Quando estiver
% com o documento, salve-o como PDF no diretório do seu projeto e substitua todo
% o conteúdo de implementação deste arquivo pelo comando abaixo:
%
% \usepackage{pdfpages}		% necessário para comando \includepdf
% \begin{fichacatalografica}
%     \includepdf{fig_ficha_catalografica.pdf}
% \end{fichacatalografica}

% \begin{fichacatalografica}
% 	\sffamily
% 	\vspace*{\fill}                             % Posição vertical
% 	\begin{center}                              % Minipage Centralizado
% 	\fbox{\begin{minipage}[c][8cm]{13.5cm}		% Largura
% 	\small
% 	\imprimirautor
% 	%Sobrenome, Nome do autor
	
% 	\hspace{0.5cm} \imprimirtitulo  / \imprimirautor. --
% 	\imprimirlocal, \imprimirdata-
	
% 	\hspace{0.5cm} \thelastpage p. : il. (algumas color.) ; 30 cm.\\
	
% 	\hspace{0.5cm} \imprimirorientadorRotulo~\imprimirorientador\\
	
% 	\hspace{0.5cm}
% 	\parbox[t]{\textwidth}{\imprimirtipotrabalho~--~\imprimirinstituicao,
% 	\imprimirdata.}\\
	
% 	\hspace{0.5cm}
% 		1. Palavra-chave1.
% 		2. Palavra-chave2.
% 		2. Palavra-chave3.
% 		I. Orientador.
% 		II. Universidade xxx.
% 		III. Faculdade de xxx.
% 		IV. Título 			
% 	\end{minipage}}
% 	\end{center}
% \end{fichacatalografica}
% ---


% ---
% FOLHA DE APROVAÇÃO
% ---

% Isto é um exemplo de Folha de aprovação, elemento obrigatório da NBR
% 14724/2011 (seção 4.2.1.3). Você pode utilizar este modelo até a aprovação
% do trabalho. Após isso, substitua todo o conteúdo deste arquivo por uma
% imagem da página assinada pela banca com o comando abaixo:
%
% \begin{folhadeaprovacao}
% \includepdf{folhadeaprovacao_final.pdf}
% \end{folhadeaprovacao}
%
\setlength{\ABNTEXsignskip}{2cm}
\setlength{\ABNTEXsignwidth}{10cm}

\begin{folhadeaprovacao}
  \begin{center}
    {\ABNTEXchapterfont\large\imprimirautor}

    \vspace*{1cm}

    \begin{center}
      \ABNTEXchapterfont\bfseries\Large\imprimirtitulo
    \end{center}

    \vspace*{1cm}
    
    \hspace{.45\textwidth}
    \begin{minipage}{.5\textwidth}
        \imprimirpreambulo
    \end{minipage}%
   \end{center}

   \vspace*{1cm}
        
   \begin{center}
   Aprovado em 17 de julho de 2024.

   \vspace*{1cm}

   Banca examinadora:
   \end{center}

   \vspace*{-1cm}

   \assinatura{Prof. Ma. Caroline Santa Maria Rodrigues -- PUCRS}
   \assinatura{Prof. Dra. Clarissa Pinto Pizarro de Freitas -- PUCRS}
   \assinatura{\imprimirorientador -- PUCRS}
\end{folhadeaprovacao}
% ---


% ---
% Agradecimentos
% ---
\begin{agradecimentos}
Agradeço aos meus pais, Ademilson José Dias e Irinês Porfírio Dias, pelo apoio incondicional que me deram ao longo da vida,
incentivando mesmo as ideias questionáveis como aquela de fazer uma segunda graduação.

A minha irmã, Maria Eduarda Dias, cuja trajetória acadêmica é fonte de inspiração.

A minha namorada, Nathalia Fallavena Ceratti, que acompanhou a esta jornada de perto, oferecendo apoio e compreensão que
fizeram toda a diferença.

Aos meus amigos, Guilherme da Silva Petró, Norton Neves e Uilinton Luciano Marques da Silva, que foram compreensíveis nos
finais de semana em que não pude jogar RPG.

Ao professor Wagner de Lara Machado pela orientação e apoio no desenvolvimento deste trabalho.
\end{agradecimentos}
% ---


% ---
% RESUMOS
% ---
% Resumo em português
\setlength{\absparsep}{18pt} % ajusta o espaçamento dos parágrafos do resumo
\begin{resumo}
Desenvolvimentos na área de aprendizagem de máquina representam uma série de oportunidades para diversas áreas de pesquisa e de atuação profissional.
O uso de modelos de aprendizagem de máquina na predição de desfecho para tratamentos em saúde mental tem sido investigado e apresenta resultados promissores.
A aplicação desse tipo de modelo possibilitaria, por exemplo, maior assertividade na personalização de tratamentos psicoterápicos a nível individual.
Apesar dos benefícios potenciais, o tema ainda é pouco discutido nos espaços de pesquisa, formação e atuação em psicologia clínica. Este trabalho tem
por objetivo introduzir o tema de aprendizagem de máquina a psicólogos. Busca-se apresentar, de maneira acessível, os conceitos e terminologia básicos
em aprendizagem de máquina, proporcionando um panorama geral da área. Para complementar a exposição conceitual, é construído um exemplo a partir de dados
de um estudo sobre eficácia de intervenções em psicologia positiva para depressão via internet. Por fim, discutem-se algumas implicações do uso desse tipo
de tecnologia para a pesquisa e para a prática em psicoterapia.

 \textbf{Palavras-chave}: aprendizagem de máquina, psicoterapia, psicologia positiva, depressão.
\end{resumo}

% Resumo em inglês
\begin{resumo}[Abstract]
 \begin{otherlanguage*}{english}
   Advances in machine learning have opened up a multitude of opportunities for various research and professional fields. The use of machine learning models to
   predict treatment outcomes in mental health has been investigated showing promising results. The application of such models would allow greater assertiveness
   in tailoring psychotherapeutic treatments at the individual level. Despite its potential benefits, the topic remains unexplored in the fields of research, training
   and practice in clinical psychology. This work aims to introduce the topic of machine learning to psychologists. It presents the basic machine learning concepts and
   terminology in an accessible manner, providing a comprehensive overview of the field. To complement the conceptual discussion, an example is constructed using data
   from a study on the effectiveness of internet-based positive psychology interventions for depression. Finally, implications of using this type of technology in
   psychotherapy research and practice are discussed.
   \vspace{\onelineskip}
 
   \noindent 
   \textbf{Keywords}: machine learning, psychotherapy, positive psychology, depression.
 \end{otherlanguage*}
\end{resumo}


% ---
% LISTA DE ILUSTRAÇÕES
% ---
\pdfbookmark[0]{\listfigurename}{lof}
\listoffigures*
\cleardoublepage
% ---


% ---
% LISTA DE QUADROS
% ---
% \pdfbookmark[0]{\listofquadrosname}{loq}
% \listofquadros*
% \cleardoublepage
% ---


% ---
% LISTA DE TABELAS
% ---
% \pdfbookmark[0]{\listtablename}{lot}
% \listoftables*
% \cleardoublepage
% ---


% ---
% LISTA DE ABREVIATURAS E SIGLAS
% ---
% \begin{siglas}
%   \item[ABNT] Associação Brasileira de Normas Técnicas
%   \item[abnTeX] ABsurdas Normas para TeX
% \end{siglas}
% ---


% ---
% LISTA DE SÍMBOLOS
% ---
% \begin{simbolos}
%   \item[$ \Gamma $] Letra grega Gama
%   \item[$ \Lambda $] Lambda
%   \item[$ \zeta $] Letra grega minúscula zeta
%   \item[$ \in $] Pertence
% \end{simbolos}
% ---


% ---
% Sumário
% ---
\pdfbookmark[0]{\contentsname}{toc}
\tableofcontents*
\cleardoublepage
% ---


% ----------------------------------------------------------
% ELEMENTOS TEXTUAIS
% ----------------------------------------------------------
\textual

\include{./01-introducao/main.tex}

\include{./02-machine-learning/main.tex}

\include{./03-exemplo/main.tex}

\include{./04-conclusao/main.tex}

% ----------------------------------------------------------
% ELEMENTOS PÓS-TEXTUAIS
% ----------------------------------------------------------
\postextual
% ----------------------------------------------------------


% ----------------------------------------------------------
% REFERÊNCIAS
% ----------------------------------------------------------
\bibliography{references.bib}


% ----------------------------------------------------------
% GLOSSÁRIO
% ----------------------------------------------------------
%
% Consulte o manual da classe abntex2 para orientações sobre o glossário.
%
%\glossary

% ----------------------------------------------------------
% APÊNDICES
% ----------------------------------------------------------

% ---
% Inicia os apêndices
% ---
% \begin{apendicesenv}

% % Imprime uma página indicando o início dos apêndices
% \partapendices

% % ----------------------------------------------------------
% \chapter{Quisque libero justo}
% % ----------------------------------------------------------

% % ----------------------------------------------------------
% \chapter{Nullam elementum urna}
% % ----------------------------------------------------------

% \end{apendicesenv}
% ---


% ----------------------------------------------------------
% ANEXOS
% ----------------------------------------------------------

% ---
% Inicia os anexos
% ---
% \begin{anexosenv}

% % Imprime uma página indicando o início dos anexos
% \partanexos

% % ---
% \chapter{Morbi ultrices rutrum lorem}
% % ---

% % ---
% \chapter{Cras non urna sed feugiat cum}
% % ---


% % ---
% \chapter{Fusce facilisis lacinia dui}
% % ---

% \end{anexosenv}

%---------------------------------------------------------------------
% INDICE REMISSIVO
%---------------------------------------------------------------------
% \phantompart
% \printindex
%---------------------------------------------------------------------
\includepdf{./05-extras/contracapa.pdf}
\end{document}


\documentclass[
	% Opções da classe memoir
	12pt,				        % tamanho da fonte
	openany,			        % capítulos começam em pág ímpar (insere página vazia caso preciso)
	oneside,			        % para impressão em recto e verso. Oposto a oneside
	a4paper,			        % tamanho do papel. 
	% Opções da classe abntex2
	%chapter=TITLE,		        % títulos de capítulos convertidos em letras maiúsculas
	%section=TITLE,		        % títulos de seções convertidos em letras maiúsculas
	%subsection=TITLE,	        % títulos de subseções convertidos em letras maiúsculas
	%subsubsection=TITLE,       % títulos de subsubseções convertidos em letras maiúsculas
	% Opções do pacote babel
	english,			        % idioma adicional para hifenização
	brazil				        % o último idioma é o principal do documento
	]{abntex2}

% ---
% Pacotes básicos 
% ---
\usepackage{lmodern}			% Usa a fonte Latin Modern			
\usepackage[T1]{fontenc}		% Selecao de codigos de fonte.
\usepackage[utf8]{inputenc}		% Codificacao do documento (conversão automática dos acentos)
\usepackage{indentfirst}		% Indenta o primeiro parágrafo de cada seção.
\usepackage{color}				% Controle das cores
\usepackage{graphicx}			% Inclusão de gráficos
\usepackage{microtype} 			% para melhorias de justificação
% ---
		

% ---
% Pacotes de citações
% ---
% \usepackage[alf]{abntex2cite}						% Citações padrão ABNT
\usepackage[
  alf,
  abnt-nbr10520=1988,                     % Citações em sentence case
  abnt-emphasize=bf
  ]{abntex2cite}


% ---
% Pacotes adicionais
% ---
\usepackage{pdfpages}	% Inclusão de páginas em PDF


% ---
% Informações para CAPA e FOLHA DE ROSTO
% ---
\titulo{APRENDIZAGEM DE MÁQUINA PARA PREDIÇÃO DE DESFECHOS EM SAÚDE MENTAL: \textmd{UMA INTRODUÇÃO PARA PSICÓLOGOS}}
\autor{GUILHERME AUGUSTO DIAS}
\local{Porto Alegre}
\data{2024}
\orientador{Prof. Dr. Wagner de Lara Machado}
% \coorientador{}
% \instituicao{
%   Pontifícia Universidade Católica do Rio Grande do Sul
%   \par
%   Escola de Ciências da Saúde e da Vida
%   \par
%   Curso de Psicologia}
\tipotrabalho{TCC}
\preambulo{Trabalho de Conclusão de Curso apresentado como requisito para a obtenção do grau de Bacharel em Psicologia pela Escola de Ciências da Saúde e da Vida da Pontifícia Universidade Católica do Rio Grande do Sul.}
% ---

% ---
% Configurações de aparência do PDF final
% ---

% Alterando o aspecto da cor azul
\definecolor{blue}{RGB}{41,5,195}

% Informações do PDF
\makeatletter
\hypersetup{
		pdftitle={Aprendizagem de máquina para predição de desfechos em saúde mental: uma introdução para psicólogos}, 
		pdfauthor={Guilherme Augusto Dias},
    	pdfsubject={\imprimirpreambulo},
	    pdfcreator={LaTeX with abnTeX2},
		pdfkeywords={aprendizagem de máquina}{psicoterapia}{psicologia positiva}{depressão}, 
		colorlinks=true,		% true: colored links; false: boxed links
    	linkcolor=black,		% color of internal links
    	citecolor=black,		% color of links to bibliography
    	filecolor=black,		% color of file links
		urlcolor=black,
		bookmarksdepth=4
}
\makeatother
% --- 


% ---
% Posiciona figuras e tabelas no topo da página quando adicionadas sozinhas
% em um página em branco. Ver https://github.com/abntex/abntex2/issues/170
\makeatletter
\setlength{\@fptop}{5pt} % Set distance from top of page to first float
\makeatother
% ---

% ---
% Possibilita criação de Quadros e Lista de quadros.
% Ver https://github.com/abntex/abntex2/issues/176
%
\newcommand{\quadroname}{Quadro}
\newcommand{\listofquadrosname}{Lista de quadros}

\newfloat[chapter]{quadro}{loq}{\quadroname}
\newlistof{listofquadros}{loq}{\listofquadrosname}
\newlistentry{quadro}{loq}{0}

% Configurações para atender às regras da ABNT
\setfloatadjustment{quadro}{\centering}
\counterwithout{quadro}{chapter}
\renewcommand{\cftquadroname}{\quadroname\space} 
\renewcommand*{\cftquadroaftersnum}{\hfill--\hfill}

\setfloatlocations{quadro}{hbtp} % Ver https://github.com/abntex/abntex2/issues/176
% ---


% --- 
% Espaçamentos entre linhas e parágrafos 
% --- 

% O tamanho do parágrafo é dado por:
\setlength{\parindent}{1.3cm}

% Controle do espaçamento entre um parágrafo e outro:
\setlength{\parskip}{0.2cm}  % tente também \onelineskip

% ---
% Compila o índice
% ---
\makeindex
% ---


% ----
% Início do documento
% ----
\begin{document}

% Seleciona o idioma do documento (conforme pacotes do babel)
\selectlanguage{brazil}

% Retira espaço extra obsoleto entre as frases.
\frenchspacing 

% ----------------------------------------------------------
% ELEMENTOS PRÉ-TEXTUAIS
% ----------------------------------------------------------
\pretextual

% ---
% CAPA
% ---
\includepdf{./05-extras/capa.pdf}
% ---


% ---
% FOLHA DE ROSTO
% (o * indica que haverá a ficha bibliográfica)
% ---
\imprimirfolhaderosto
% ---


% ---
% Ficha bibliografica
% ---

% Isto é um exemplo de Ficha Catalográfica, ou ``Dados internacionais de
% catalogação-na-publicação''. Você pode utilizar este modelo como referência. 
% Porém, provavelmente a biblioteca da sua universidade lhe fornecerá um PDF
% com a ficha catalográfica definitiva após a defesa do trabalho. Quando estiver
% com o documento, salve-o como PDF no diretório do seu projeto e substitua todo
% o conteúdo de implementação deste arquivo pelo comando abaixo:
%
% \usepackage{pdfpages}		% necessário para comando \includepdf
% \begin{fichacatalografica}
%     \includepdf{fig_ficha_catalografica.pdf}
% \end{fichacatalografica}

% \begin{fichacatalografica}
% 	\sffamily
% 	\vspace*{\fill}                             % Posição vertical
% 	\begin{center}                              % Minipage Centralizado
% 	\fbox{\begin{minipage}[c][8cm]{13.5cm}		% Largura
% 	\small
% 	\imprimirautor
% 	%Sobrenome, Nome do autor
	
% 	\hspace{0.5cm} \imprimirtitulo  / \imprimirautor. --
% 	\imprimirlocal, \imprimirdata-
	
% 	\hspace{0.5cm} \thelastpage p. : il. (algumas color.) ; 30 cm.\\
	
% 	\hspace{0.5cm} \imprimirorientadorRotulo~\imprimirorientador\\
	
% 	\hspace{0.5cm}
% 	\parbox[t]{\textwidth}{\imprimirtipotrabalho~--~\imprimirinstituicao,
% 	\imprimirdata.}\\
	
% 	\hspace{0.5cm}
% 		1. Palavra-chave1.
% 		2. Palavra-chave2.
% 		2. Palavra-chave3.
% 		I. Orientador.
% 		II. Universidade xxx.
% 		III. Faculdade de xxx.
% 		IV. Título 			
% 	\end{minipage}}
% 	\end{center}
% \end{fichacatalografica}
% ---


% ---
% FOLHA DE APROVAÇÃO
% ---

% Isto é um exemplo de Folha de aprovação, elemento obrigatório da NBR
% 14724/2011 (seção 4.2.1.3). Você pode utilizar este modelo até a aprovação
% do trabalho. Após isso, substitua todo o conteúdo deste arquivo por uma
% imagem da página assinada pela banca com o comando abaixo:
%
% \begin{folhadeaprovacao}
% \includepdf{folhadeaprovacao_final.pdf}
% \end{folhadeaprovacao}
%
\setlength{\ABNTEXsignskip}{2cm}
\setlength{\ABNTEXsignwidth}{10cm}

\begin{folhadeaprovacao}
  \begin{center}
    {\ABNTEXchapterfont\large\imprimirautor}

    \vspace*{1cm}

    \begin{center}
      \ABNTEXchapterfont\bfseries\Large\imprimirtitulo
    \end{center}

    \vspace*{1cm}
    
    \hspace{.45\textwidth}
    \begin{minipage}{.5\textwidth}
        \imprimirpreambulo
    \end{minipage}%
   \end{center}

   \vspace*{1cm}
        
   \begin{center}
   Aprovado em 17 de julho de 2024.

   \vspace*{1cm}

   Banca examinadora:
   \end{center}

   \vspace*{-1cm}

   \assinatura{Prof. Ma. Caroline Santa Maria Rodrigues -- PUCRS}
   \assinatura{Prof. Dra. Clarissa Pinto Pizarro de Freitas -- PUCRS}
   \assinatura{\imprimirorientador -- PUCRS}
\end{folhadeaprovacao}
% ---


% ---
% Agradecimentos
% ---
\begin{agradecimentos}
Agradeço aos meus pais, Ademilson José Dias e Irinês Porfírio Dias, pelo apoio incondicional que me deram ao longo da vida,
incentivando mesmo as ideias questionáveis como aquela de fazer uma segunda graduação.

A minha irmã, Maria Eduarda Dias, cuja trajetória acadêmica é fonte de inspiração.

A minha namorada, Nathalia Fallavena Ceratti, que acompanhou a esta jornada de perto, oferecendo apoio e compreensão que
fizeram toda a diferença.

Aos meus amigos, Guilherme da Silva Petró, Norton Neves e Uilinton Luciano Marques da Silva, que foram compreensíveis nos
finais de semana em que não pude jogar RPG.

Ao professor Wagner de Lara Machado pela orientação e apoio no desenvolvimento deste trabalho.
\end{agradecimentos}
% ---


% ---
% RESUMOS
% ---
% Resumo em português
\setlength{\absparsep}{18pt} % ajusta o espaçamento dos parágrafos do resumo
\begin{resumo}
Desenvolvimentos na área de aprendizagem de máquina representam uma série de oportunidades para diversas áreas de pesquisa e de atuação profissional.
O uso de modelos de aprendizagem de máquina na predição de desfecho para tratamentos em saúde mental tem sido investigado e apresenta resultados promissores.
A aplicação desse tipo de modelo possibilitaria, por exemplo, maior assertividade na personalização de tratamentos psicoterápicos a nível individual.
Apesar dos benefícios potenciais, o tema ainda é pouco discutido nos espaços de pesquisa, formação e atuação em psicologia clínica. Este trabalho tem
por objetivo introduzir o tema de aprendizagem de máquina a psicólogos. Busca-se apresentar, de maneira acessível, os conceitos e terminologia básicos
em aprendizagem de máquina, proporcionando um panorama geral da área. Para complementar a exposição conceitual, é construído um exemplo a partir de dados
de um estudo sobre eficácia de intervenções em psicologia positiva para depressão via internet. Por fim, discutem-se algumas implicações do uso desse tipo
de tecnologia para a pesquisa e para a prática em psicoterapia.

 \textbf{Palavras-chave}: aprendizagem de máquina, psicoterapia, psicologia positiva, depressão.
\end{resumo}

% Resumo em inglês
\begin{resumo}[Abstract]
 \begin{otherlanguage*}{english}
   Advances in machine learning have opened up a multitude of opportunities for various research and professional fields. The use of machine learning models to
   predict treatment outcomes in mental health has been investigated showing promising results. The application of such models would allow greater assertiveness
   in tailoring psychotherapeutic treatments at the individual level. Despite its potential benefits, the topic remains unexplored in the fields of research, training
   and practice in clinical psychology. This work aims to introduce the topic of machine learning to psychologists. It presents the basic machine learning concepts and
   terminology in an accessible manner, providing a comprehensive overview of the field. To complement the conceptual discussion, an example is constructed using data
   from a study on the effectiveness of internet-based positive psychology interventions for depression. Finally, implications of using this type of technology in
   psychotherapy research and practice are discussed.
   \vspace{\onelineskip}
 
   \noindent 
   \textbf{Keywords}: machine learning, psychotherapy, positive psychology, depression.
 \end{otherlanguage*}
\end{resumo}


% ---
% LISTA DE ILUSTRAÇÕES
% ---
\pdfbookmark[0]{\listfigurename}{lof}
\listoffigures*
\cleardoublepage
% ---


% ---
% LISTA DE QUADROS
% ---
% \pdfbookmark[0]{\listofquadrosname}{loq}
% \listofquadros*
% \cleardoublepage
% ---


% ---
% LISTA DE TABELAS
% ---
% \pdfbookmark[0]{\listtablename}{lot}
% \listoftables*
% \cleardoublepage
% ---


% ---
% LISTA DE ABREVIATURAS E SIGLAS
% ---
% \begin{siglas}
%   \item[ABNT] Associação Brasileira de Normas Técnicas
%   \item[abnTeX] ABsurdas Normas para TeX
% \end{siglas}
% ---


% ---
% LISTA DE SÍMBOLOS
% ---
% \begin{simbolos}
%   \item[$ \Gamma $] Letra grega Gama
%   \item[$ \Lambda $] Lambda
%   \item[$ \zeta $] Letra grega minúscula zeta
%   \item[$ \in $] Pertence
% \end{simbolos}
% ---


% ---
% Sumário
% ---
\pdfbookmark[0]{\contentsname}{toc}
\tableofcontents*
\cleardoublepage
% ---


% ----------------------------------------------------------
% ELEMENTOS TEXTUAIS
% ----------------------------------------------------------
\textual

\include{./01-introducao/main.tex}

\include{./02-machine-learning/main.tex}

\include{./03-exemplo/main.tex}

\include{./04-conclusao/main.tex}

% ----------------------------------------------------------
% ELEMENTOS PÓS-TEXTUAIS
% ----------------------------------------------------------
\postextual
% ----------------------------------------------------------


% ----------------------------------------------------------
% REFERÊNCIAS
% ----------------------------------------------------------
\bibliography{references.bib}


% ----------------------------------------------------------
% GLOSSÁRIO
% ----------------------------------------------------------
%
% Consulte o manual da classe abntex2 para orientações sobre o glossário.
%
%\glossary

% ----------------------------------------------------------
% APÊNDICES
% ----------------------------------------------------------

% ---
% Inicia os apêndices
% ---
% \begin{apendicesenv}

% % Imprime uma página indicando o início dos apêndices
% \partapendices

% % ----------------------------------------------------------
% \chapter{Quisque libero justo}
% % ----------------------------------------------------------

% % ----------------------------------------------------------
% \chapter{Nullam elementum urna}
% % ----------------------------------------------------------

% \end{apendicesenv}
% ---


% ----------------------------------------------------------
% ANEXOS
% ----------------------------------------------------------

% ---
% Inicia os anexos
% ---
% \begin{anexosenv}

% % Imprime uma página indicando o início dos anexos
% \partanexos

% % ---
% \chapter{Morbi ultrices rutrum lorem}
% % ---

% % ---
% \chapter{Cras non urna sed feugiat cum}
% % ---


% % ---
% \chapter{Fusce facilisis lacinia dui}
% % ---

% \end{anexosenv}

%---------------------------------------------------------------------
% INDICE REMISSIVO
%---------------------------------------------------------------------
% \phantompart
% \printindex
%---------------------------------------------------------------------
\includepdf{./05-extras/contracapa.pdf}
\end{document}


\documentclass[
	% Opções da classe memoir
	12pt,				        % tamanho da fonte
	openany,			        % capítulos começam em pág ímpar (insere página vazia caso preciso)
	oneside,			        % para impressão em recto e verso. Oposto a oneside
	a4paper,			        % tamanho do papel. 
	% Opções da classe abntex2
	%chapter=TITLE,		        % títulos de capítulos convertidos em letras maiúsculas
	%section=TITLE,		        % títulos de seções convertidos em letras maiúsculas
	%subsection=TITLE,	        % títulos de subseções convertidos em letras maiúsculas
	%subsubsection=TITLE,       % títulos de subsubseções convertidos em letras maiúsculas
	% Opções do pacote babel
	english,			        % idioma adicional para hifenização
	brazil				        % o último idioma é o principal do documento
	]{abntex2}

% ---
% Pacotes básicos 
% ---
\usepackage{lmodern}			% Usa a fonte Latin Modern			
\usepackage[T1]{fontenc}		% Selecao de codigos de fonte.
\usepackage[utf8]{inputenc}		% Codificacao do documento (conversão automática dos acentos)
\usepackage{indentfirst}		% Indenta o primeiro parágrafo de cada seção.
\usepackage{color}				% Controle das cores
\usepackage{graphicx}			% Inclusão de gráficos
\usepackage{microtype} 			% para melhorias de justificação
% ---
		

% ---
% Pacotes de citações
% ---
% \usepackage[alf]{abntex2cite}						% Citações padrão ABNT
\usepackage[
  alf,
  abnt-nbr10520=1988,                     % Citações em sentence case
  abnt-emphasize=bf
  ]{abntex2cite}


% ---
% Pacotes adicionais
% ---
\usepackage{pdfpages}	% Inclusão de páginas em PDF


% ---
% Informações para CAPA e FOLHA DE ROSTO
% ---
\titulo{APRENDIZAGEM DE MÁQUINA PARA PREDIÇÃO DE DESFECHOS EM SAÚDE MENTAL: \textmd{UMA INTRODUÇÃO PARA PSICÓLOGOS}}
\autor{GUILHERME AUGUSTO DIAS}
\local{Porto Alegre}
\data{2024}
\orientador{Prof. Dr. Wagner de Lara Machado}
% \coorientador{}
% \instituicao{
%   Pontifícia Universidade Católica do Rio Grande do Sul
%   \par
%   Escola de Ciências da Saúde e da Vida
%   \par
%   Curso de Psicologia}
\tipotrabalho{TCC}
\preambulo{Trabalho de Conclusão de Curso apresentado como requisito para a obtenção do grau de Bacharel em Psicologia pela Escola de Ciências da Saúde e da Vida da Pontifícia Universidade Católica do Rio Grande do Sul.}
% ---

% ---
% Configurações de aparência do PDF final
% ---

% Alterando o aspecto da cor azul
\definecolor{blue}{RGB}{41,5,195}

% Informações do PDF
\makeatletter
\hypersetup{
		pdftitle={Aprendizagem de máquina para predição de desfechos em saúde mental: uma introdução para psicólogos}, 
		pdfauthor={Guilherme Augusto Dias},
    	pdfsubject={\imprimirpreambulo},
	    pdfcreator={LaTeX with abnTeX2},
		pdfkeywords={aprendizagem de máquina}{psicoterapia}{psicologia positiva}{depressão}, 
		colorlinks=true,		% true: colored links; false: boxed links
    	linkcolor=black,		% color of internal links
    	citecolor=black,		% color of links to bibliography
    	filecolor=black,		% color of file links
		urlcolor=black,
		bookmarksdepth=4
}
\makeatother
% --- 


% ---
% Posiciona figuras e tabelas no topo da página quando adicionadas sozinhas
% em um página em branco. Ver https://github.com/abntex/abntex2/issues/170
\makeatletter
\setlength{\@fptop}{5pt} % Set distance from top of page to first float
\makeatother
% ---

% ---
% Possibilita criação de Quadros e Lista de quadros.
% Ver https://github.com/abntex/abntex2/issues/176
%
\newcommand{\quadroname}{Quadro}
\newcommand{\listofquadrosname}{Lista de quadros}

\newfloat[chapter]{quadro}{loq}{\quadroname}
\newlistof{listofquadros}{loq}{\listofquadrosname}
\newlistentry{quadro}{loq}{0}

% Configurações para atender às regras da ABNT
\setfloatadjustment{quadro}{\centering}
\counterwithout{quadro}{chapter}
\renewcommand{\cftquadroname}{\quadroname\space} 
\renewcommand*{\cftquadroaftersnum}{\hfill--\hfill}

\setfloatlocations{quadro}{hbtp} % Ver https://github.com/abntex/abntex2/issues/176
% ---


% --- 
% Espaçamentos entre linhas e parágrafos 
% --- 

% O tamanho do parágrafo é dado por:
\setlength{\parindent}{1.3cm}

% Controle do espaçamento entre um parágrafo e outro:
\setlength{\parskip}{0.2cm}  % tente também \onelineskip

% ---
% Compila o índice
% ---
\makeindex
% ---


% ----
% Início do documento
% ----
\begin{document}

% Seleciona o idioma do documento (conforme pacotes do babel)
\selectlanguage{brazil}

% Retira espaço extra obsoleto entre as frases.
\frenchspacing 

% ----------------------------------------------------------
% ELEMENTOS PRÉ-TEXTUAIS
% ----------------------------------------------------------
\pretextual

% ---
% CAPA
% ---
\includepdf{./05-extras/capa.pdf}
% ---


% ---
% FOLHA DE ROSTO
% (o * indica que haverá a ficha bibliográfica)
% ---
\imprimirfolhaderosto
% ---


% ---
% Ficha bibliografica
% ---

% Isto é um exemplo de Ficha Catalográfica, ou ``Dados internacionais de
% catalogação-na-publicação''. Você pode utilizar este modelo como referência. 
% Porém, provavelmente a biblioteca da sua universidade lhe fornecerá um PDF
% com a ficha catalográfica definitiva após a defesa do trabalho. Quando estiver
% com o documento, salve-o como PDF no diretório do seu projeto e substitua todo
% o conteúdo de implementação deste arquivo pelo comando abaixo:
%
% \usepackage{pdfpages}		% necessário para comando \includepdf
% \begin{fichacatalografica}
%     \includepdf{fig_ficha_catalografica.pdf}
% \end{fichacatalografica}

% \begin{fichacatalografica}
% 	\sffamily
% 	\vspace*{\fill}                             % Posição vertical
% 	\begin{center}                              % Minipage Centralizado
% 	\fbox{\begin{minipage}[c][8cm]{13.5cm}		% Largura
% 	\small
% 	\imprimirautor
% 	%Sobrenome, Nome do autor
	
% 	\hspace{0.5cm} \imprimirtitulo  / \imprimirautor. --
% 	\imprimirlocal, \imprimirdata-
	
% 	\hspace{0.5cm} \thelastpage p. : il. (algumas color.) ; 30 cm.\\
	
% 	\hspace{0.5cm} \imprimirorientadorRotulo~\imprimirorientador\\
	
% 	\hspace{0.5cm}
% 	\parbox[t]{\textwidth}{\imprimirtipotrabalho~--~\imprimirinstituicao,
% 	\imprimirdata.}\\
	
% 	\hspace{0.5cm}
% 		1. Palavra-chave1.
% 		2. Palavra-chave2.
% 		2. Palavra-chave3.
% 		I. Orientador.
% 		II. Universidade xxx.
% 		III. Faculdade de xxx.
% 		IV. Título 			
% 	\end{minipage}}
% 	\end{center}
% \end{fichacatalografica}
% ---


% ---
% FOLHA DE APROVAÇÃO
% ---

% Isto é um exemplo de Folha de aprovação, elemento obrigatório da NBR
% 14724/2011 (seção 4.2.1.3). Você pode utilizar este modelo até a aprovação
% do trabalho. Após isso, substitua todo o conteúdo deste arquivo por uma
% imagem da página assinada pela banca com o comando abaixo:
%
% \begin{folhadeaprovacao}
% \includepdf{folhadeaprovacao_final.pdf}
% \end{folhadeaprovacao}
%
\setlength{\ABNTEXsignskip}{2cm}
\setlength{\ABNTEXsignwidth}{10cm}

\begin{folhadeaprovacao}
  \begin{center}
    {\ABNTEXchapterfont\large\imprimirautor}

    \vspace*{1cm}

    \begin{center}
      \ABNTEXchapterfont\bfseries\Large\imprimirtitulo
    \end{center}

    \vspace*{1cm}
    
    \hspace{.45\textwidth}
    \begin{minipage}{.5\textwidth}
        \imprimirpreambulo
    \end{minipage}%
   \end{center}

   \vspace*{1cm}
        
   \begin{center}
   Aprovado em 17 de julho de 2024.

   \vspace*{1cm}

   Banca examinadora:
   \end{center}

   \vspace*{-1cm}

   \assinatura{Prof. Ma. Caroline Santa Maria Rodrigues -- PUCRS}
   \assinatura{Prof. Dra. Clarissa Pinto Pizarro de Freitas -- PUCRS}
   \assinatura{\imprimirorientador -- PUCRS}
\end{folhadeaprovacao}
% ---


% ---
% Agradecimentos
% ---
\begin{agradecimentos}
Agradeço aos meus pais, Ademilson José Dias e Irinês Porfírio Dias, pelo apoio incondicional que me deram ao longo da vida,
incentivando mesmo as ideias questionáveis como aquela de fazer uma segunda graduação.

A minha irmã, Maria Eduarda Dias, cuja trajetória acadêmica é fonte de inspiração.

A minha namorada, Nathalia Fallavena Ceratti, que acompanhou a esta jornada de perto, oferecendo apoio e compreensão que
fizeram toda a diferença.

Aos meus amigos, Guilherme da Silva Petró, Norton Neves e Uilinton Luciano Marques da Silva, que foram compreensíveis nos
finais de semana em que não pude jogar RPG.

Ao professor Wagner de Lara Machado pela orientação e apoio no desenvolvimento deste trabalho.
\end{agradecimentos}
% ---


% ---
% RESUMOS
% ---
% Resumo em português
\setlength{\absparsep}{18pt} % ajusta o espaçamento dos parágrafos do resumo
\begin{resumo}
Desenvolvimentos na área de aprendizagem de máquina representam uma série de oportunidades para diversas áreas de pesquisa e de atuação profissional.
O uso de modelos de aprendizagem de máquina na predição de desfecho para tratamentos em saúde mental tem sido investigado e apresenta resultados promissores.
A aplicação desse tipo de modelo possibilitaria, por exemplo, maior assertividade na personalização de tratamentos psicoterápicos a nível individual.
Apesar dos benefícios potenciais, o tema ainda é pouco discutido nos espaços de pesquisa, formação e atuação em psicologia clínica. Este trabalho tem
por objetivo introduzir o tema de aprendizagem de máquina a psicólogos. Busca-se apresentar, de maneira acessível, os conceitos e terminologia básicos
em aprendizagem de máquina, proporcionando um panorama geral da área. Para complementar a exposição conceitual, é construído um exemplo a partir de dados
de um estudo sobre eficácia de intervenções em psicologia positiva para depressão via internet. Por fim, discutem-se algumas implicações do uso desse tipo
de tecnologia para a pesquisa e para a prática em psicoterapia.

 \textbf{Palavras-chave}: aprendizagem de máquina, psicoterapia, psicologia positiva, depressão.
\end{resumo}

% Resumo em inglês
\begin{resumo}[Abstract]
 \begin{otherlanguage*}{english}
   Advances in machine learning have opened up a multitude of opportunities for various research and professional fields. The use of machine learning models to
   predict treatment outcomes in mental health has been investigated showing promising results. The application of such models would allow greater assertiveness
   in tailoring psychotherapeutic treatments at the individual level. Despite its potential benefits, the topic remains unexplored in the fields of research, training
   and practice in clinical psychology. This work aims to introduce the topic of machine learning to psychologists. It presents the basic machine learning concepts and
   terminology in an accessible manner, providing a comprehensive overview of the field. To complement the conceptual discussion, an example is constructed using data
   from a study on the effectiveness of internet-based positive psychology interventions for depression. Finally, implications of using this type of technology in
   psychotherapy research and practice are discussed.
   \vspace{\onelineskip}
 
   \noindent 
   \textbf{Keywords}: machine learning, psychotherapy, positive psychology, depression.
 \end{otherlanguage*}
\end{resumo}


% ---
% LISTA DE ILUSTRAÇÕES
% ---
\pdfbookmark[0]{\listfigurename}{lof}
\listoffigures*
\cleardoublepage
% ---


% ---
% LISTA DE QUADROS
% ---
% \pdfbookmark[0]{\listofquadrosname}{loq}
% \listofquadros*
% \cleardoublepage
% ---


% ---
% LISTA DE TABELAS
% ---
% \pdfbookmark[0]{\listtablename}{lot}
% \listoftables*
% \cleardoublepage
% ---


% ---
% LISTA DE ABREVIATURAS E SIGLAS
% ---
% \begin{siglas}
%   \item[ABNT] Associação Brasileira de Normas Técnicas
%   \item[abnTeX] ABsurdas Normas para TeX
% \end{siglas}
% ---


% ---
% LISTA DE SÍMBOLOS
% ---
% \begin{simbolos}
%   \item[$ \Gamma $] Letra grega Gama
%   \item[$ \Lambda $] Lambda
%   \item[$ \zeta $] Letra grega minúscula zeta
%   \item[$ \in $] Pertence
% \end{simbolos}
% ---


% ---
% Sumário
% ---
\pdfbookmark[0]{\contentsname}{toc}
\tableofcontents*
\cleardoublepage
% ---


% ----------------------------------------------------------
% ELEMENTOS TEXTUAIS
% ----------------------------------------------------------
\textual

\include{./01-introducao/main.tex}

\include{./02-machine-learning/main.tex}

\include{./03-exemplo/main.tex}

\include{./04-conclusao/main.tex}

% ----------------------------------------------------------
% ELEMENTOS PÓS-TEXTUAIS
% ----------------------------------------------------------
\postextual
% ----------------------------------------------------------


% ----------------------------------------------------------
% REFERÊNCIAS
% ----------------------------------------------------------
\bibliography{references.bib}


% ----------------------------------------------------------
% GLOSSÁRIO
% ----------------------------------------------------------
%
% Consulte o manual da classe abntex2 para orientações sobre o glossário.
%
%\glossary

% ----------------------------------------------------------
% APÊNDICES
% ----------------------------------------------------------

% ---
% Inicia os apêndices
% ---
% \begin{apendicesenv}

% % Imprime uma página indicando o início dos apêndices
% \partapendices

% % ----------------------------------------------------------
% \chapter{Quisque libero justo}
% % ----------------------------------------------------------

% % ----------------------------------------------------------
% \chapter{Nullam elementum urna}
% % ----------------------------------------------------------

% \end{apendicesenv}
% ---


% ----------------------------------------------------------
% ANEXOS
% ----------------------------------------------------------

% ---
% Inicia os anexos
% ---
% \begin{anexosenv}

% % Imprime uma página indicando o início dos anexos
% \partanexos

% % ---
% \chapter{Morbi ultrices rutrum lorem}
% % ---

% % ---
% \chapter{Cras non urna sed feugiat cum}
% % ---


% % ---
% \chapter{Fusce facilisis lacinia dui}
% % ---

% \end{anexosenv}

%---------------------------------------------------------------------
% INDICE REMISSIVO
%---------------------------------------------------------------------
% \phantompart
% \printindex
%---------------------------------------------------------------------
\includepdf{./05-extras/contracapa.pdf}
\end{document}


% ----------------------------------------------------------
% ELEMENTOS PÓS-TEXTUAIS
% ----------------------------------------------------------
\postextual
% ----------------------------------------------------------


% ----------------------------------------------------------
% REFERÊNCIAS
% ----------------------------------------------------------
\bibliography{references.bib}


% ----------------------------------------------------------
% GLOSSÁRIO
% ----------------------------------------------------------
%
% Consulte o manual da classe abntex2 para orientações sobre o glossário.
%
%\glossary

% ----------------------------------------------------------
% APÊNDICES
% ----------------------------------------------------------

% ---
% Inicia os apêndices
% ---
% \begin{apendicesenv}

% % Imprime uma página indicando o início dos apêndices
% \partapendices

% % ----------------------------------------------------------
% \chapter{Quisque libero justo}
% % ----------------------------------------------------------

% % ----------------------------------------------------------
% \chapter{Nullam elementum urna}
% % ----------------------------------------------------------

% \end{apendicesenv}
% ---


% ----------------------------------------------------------
% ANEXOS
% ----------------------------------------------------------

% ---
% Inicia os anexos
% ---
% \begin{anexosenv}

% % Imprime uma página indicando o início dos anexos
% \partanexos

% % ---
% \chapter{Morbi ultrices rutrum lorem}
% % ---

% % ---
% \chapter{Cras non urna sed feugiat cum}
% % ---


% % ---
% \chapter{Fusce facilisis lacinia dui}
% % ---

% \end{anexosenv}

%---------------------------------------------------------------------
% INDICE REMISSIVO
%---------------------------------------------------------------------
% \phantompart
% \printindex
%---------------------------------------------------------------------
\includepdf{./05-extras/contracapa.pdf}
\end{document}


\documentclass[
	% Opções da classe memoir
	12pt,				        % tamanho da fonte
	openany,			        % capítulos começam em pág ímpar (insere página vazia caso preciso)
	oneside,			        % para impressão em recto e verso. Oposto a oneside
	a4paper,			        % tamanho do papel. 
	% Opções da classe abntex2
	%chapter=TITLE,		        % títulos de capítulos convertidos em letras maiúsculas
	%section=TITLE,		        % títulos de seções convertidos em letras maiúsculas
	%subsection=TITLE,	        % títulos de subseções convertidos em letras maiúsculas
	%subsubsection=TITLE,       % títulos de subsubseções convertidos em letras maiúsculas
	% Opções do pacote babel
	english,			        % idioma adicional para hifenização
	brazil				        % o último idioma é o principal do documento
	]{abntex2}

% ---
% Pacotes básicos 
% ---
\usepackage{lmodern}			% Usa a fonte Latin Modern			
\usepackage[T1]{fontenc}		% Selecao de codigos de fonte.
\usepackage[utf8]{inputenc}		% Codificacao do documento (conversão automática dos acentos)
\usepackage{indentfirst}		% Indenta o primeiro parágrafo de cada seção.
\usepackage{color}				% Controle das cores
\usepackage{graphicx}			% Inclusão de gráficos
\usepackage{microtype} 			% para melhorias de justificação
% ---
		

% ---
% Pacotes de citações
% ---
% \usepackage[alf]{abntex2cite}						% Citações padrão ABNT
\usepackage[
  alf,
  abnt-nbr10520=1988,                     % Citações em sentence case
  abnt-emphasize=bf
  ]{abntex2cite}


% ---
% Pacotes adicionais
% ---
\usepackage{pdfpages}	% Inclusão de páginas em PDF


% ---
% Informações para CAPA e FOLHA DE ROSTO
% ---
\titulo{APRENDIZAGEM DE MÁQUINA PARA PREDIÇÃO DE DESFECHOS EM SAÚDE MENTAL: \textmd{UMA INTRODUÇÃO PARA PSICÓLOGOS}}
\autor{GUILHERME AUGUSTO DIAS}
\local{Porto Alegre}
\data{2024}
\orientador{Prof. Dr. Wagner de Lara Machado}
% \coorientador{}
% \instituicao{
%   Pontifícia Universidade Católica do Rio Grande do Sul
%   \par
%   Escola de Ciências da Saúde e da Vida
%   \par
%   Curso de Psicologia}
\tipotrabalho{TCC}
\preambulo{Trabalho de Conclusão de Curso apresentado como requisito para a obtenção do grau de Bacharel em Psicologia pela Escola de Ciências da Saúde e da Vida da Pontifícia Universidade Católica do Rio Grande do Sul.}
% ---

% ---
% Configurações de aparência do PDF final
% ---

% Alterando o aspecto da cor azul
\definecolor{blue}{RGB}{41,5,195}

% Informações do PDF
\makeatletter
\hypersetup{
		pdftitle={Aprendizagem de máquina para predição de desfechos em saúde mental: uma introdução para psicólogos}, 
		pdfauthor={Guilherme Augusto Dias},
    	pdfsubject={\imprimirpreambulo},
	    pdfcreator={LaTeX with abnTeX2},
		pdfkeywords={aprendizagem de máquina}{psicoterapia}{psicologia positiva}{depressão}, 
		colorlinks=true,		% true: colored links; false: boxed links
    	linkcolor=black,		% color of internal links
    	citecolor=black,		% color of links to bibliography
    	filecolor=black,		% color of file links
		urlcolor=black,
		bookmarksdepth=4
}
\makeatother
% --- 


% ---
% Posiciona figuras e tabelas no topo da página quando adicionadas sozinhas
% em um página em branco. Ver https://github.com/abntex/abntex2/issues/170
\makeatletter
\setlength{\@fptop}{5pt} % Set distance from top of page to first float
\makeatother
% ---

% ---
% Possibilita criação de Quadros e Lista de quadros.
% Ver https://github.com/abntex/abntex2/issues/176
%
\newcommand{\quadroname}{Quadro}
\newcommand{\listofquadrosname}{Lista de quadros}

\newfloat[chapter]{quadro}{loq}{\quadroname}
\newlistof{listofquadros}{loq}{\listofquadrosname}
\newlistentry{quadro}{loq}{0}

% Configurações para atender às regras da ABNT
\setfloatadjustment{quadro}{\centering}
\counterwithout{quadro}{chapter}
\renewcommand{\cftquadroname}{\quadroname\space} 
\renewcommand*{\cftquadroaftersnum}{\hfill--\hfill}

\setfloatlocations{quadro}{hbtp} % Ver https://github.com/abntex/abntex2/issues/176
% ---


% --- 
% Espaçamentos entre linhas e parágrafos 
% --- 

% O tamanho do parágrafo é dado por:
\setlength{\parindent}{1.3cm}

% Controle do espaçamento entre um parágrafo e outro:
\setlength{\parskip}{0.2cm}  % tente também \onelineskip

% ---
% Compila o índice
% ---
\makeindex
% ---


% ----
% Início do documento
% ----
\begin{document}

% Seleciona o idioma do documento (conforme pacotes do babel)
\selectlanguage{brazil}

% Retira espaço extra obsoleto entre as frases.
\frenchspacing 

% ----------------------------------------------------------
% ELEMENTOS PRÉ-TEXTUAIS
% ----------------------------------------------------------
\pretextual

% ---
% CAPA
% ---
\includepdf{./05-extras/capa.pdf}
% ---


% ---
% FOLHA DE ROSTO
% (o * indica que haverá a ficha bibliográfica)
% ---
\imprimirfolhaderosto
% ---


% ---
% Ficha bibliografica
% ---

% Isto é um exemplo de Ficha Catalográfica, ou ``Dados internacionais de
% catalogação-na-publicação''. Você pode utilizar este modelo como referência. 
% Porém, provavelmente a biblioteca da sua universidade lhe fornecerá um PDF
% com a ficha catalográfica definitiva após a defesa do trabalho. Quando estiver
% com o documento, salve-o como PDF no diretório do seu projeto e substitua todo
% o conteúdo de implementação deste arquivo pelo comando abaixo:
%
% \usepackage{pdfpages}		% necessário para comando \includepdf
% \begin{fichacatalografica}
%     \includepdf{fig_ficha_catalografica.pdf}
% \end{fichacatalografica}

% \begin{fichacatalografica}
% 	\sffamily
% 	\vspace*{\fill}                             % Posição vertical
% 	\begin{center}                              % Minipage Centralizado
% 	\fbox{\begin{minipage}[c][8cm]{13.5cm}		% Largura
% 	\small
% 	\imprimirautor
% 	%Sobrenome, Nome do autor
	
% 	\hspace{0.5cm} \imprimirtitulo  / \imprimirautor. --
% 	\imprimirlocal, \imprimirdata-
	
% 	\hspace{0.5cm} \thelastpage p. : il. (algumas color.) ; 30 cm.\\
	
% 	\hspace{0.5cm} \imprimirorientadorRotulo~\imprimirorientador\\
	
% 	\hspace{0.5cm}
% 	\parbox[t]{\textwidth}{\imprimirtipotrabalho~--~\imprimirinstituicao,
% 	\imprimirdata.}\\
	
% 	\hspace{0.5cm}
% 		1. Palavra-chave1.
% 		2. Palavra-chave2.
% 		2. Palavra-chave3.
% 		I. Orientador.
% 		II. Universidade xxx.
% 		III. Faculdade de xxx.
% 		IV. Título 			
% 	\end{minipage}}
% 	\end{center}
% \end{fichacatalografica}
% ---


% ---
% FOLHA DE APROVAÇÃO
% ---

% Isto é um exemplo de Folha de aprovação, elemento obrigatório da NBR
% 14724/2011 (seção 4.2.1.3). Você pode utilizar este modelo até a aprovação
% do trabalho. Após isso, substitua todo o conteúdo deste arquivo por uma
% imagem da página assinada pela banca com o comando abaixo:
%
% \begin{folhadeaprovacao}
% \includepdf{folhadeaprovacao_final.pdf}
% \end{folhadeaprovacao}
%
\setlength{\ABNTEXsignskip}{2cm}
\setlength{\ABNTEXsignwidth}{10cm}

\begin{folhadeaprovacao}
  \begin{center}
    {\ABNTEXchapterfont\large\imprimirautor}

    \vspace*{1cm}

    \begin{center}
      \ABNTEXchapterfont\bfseries\Large\imprimirtitulo
    \end{center}

    \vspace*{1cm}
    
    \hspace{.45\textwidth}
    \begin{minipage}{.5\textwidth}
        \imprimirpreambulo
    \end{minipage}%
   \end{center}

   \vspace*{1cm}
        
   \begin{center}
   Aprovado em 17 de julho de 2024.

   \vspace*{1cm}

   Banca examinadora:
   \end{center}

   \vspace*{-1cm}

   \assinatura{Prof. Ma. Caroline Santa Maria Rodrigues -- PUCRS}
   \assinatura{Prof. Dra. Clarissa Pinto Pizarro de Freitas -- PUCRS}
   \assinatura{\imprimirorientador -- PUCRS}
\end{folhadeaprovacao}
% ---


% ---
% Agradecimentos
% ---
\begin{agradecimentos}
Agradeço aos meus pais, Ademilson José Dias e Irinês Porfírio Dias, pelo apoio incondicional que me deram ao longo da vida,
incentivando mesmo as ideias questionáveis como aquela de fazer uma segunda graduação.

A minha irmã, Maria Eduarda Dias, cuja trajetória acadêmica é fonte de inspiração.

A minha namorada, Nathalia Fallavena Ceratti, que acompanhou a esta jornada de perto, oferecendo apoio e compreensão que
fizeram toda a diferença.

Aos meus amigos, Guilherme da Silva Petró, Norton Neves e Uilinton Luciano Marques da Silva, que foram compreensíveis nos
finais de semana em que não pude jogar RPG.

Ao professor Wagner de Lara Machado pela orientação e apoio no desenvolvimento deste trabalho.
\end{agradecimentos}
% ---


% ---
% RESUMOS
% ---
% Resumo em português
\setlength{\absparsep}{18pt} % ajusta o espaçamento dos parágrafos do resumo
\begin{resumo}
Desenvolvimentos na área de aprendizagem de máquina representam uma série de oportunidades para diversas áreas de pesquisa e de atuação profissional.
O uso de modelos de aprendizagem de máquina na predição de desfecho para tratamentos em saúde mental tem sido investigado e apresenta resultados promissores.
A aplicação desse tipo de modelo possibilitaria, por exemplo, maior assertividade na personalização de tratamentos psicoterápicos a nível individual.
Apesar dos benefícios potenciais, o tema ainda é pouco discutido nos espaços de pesquisa, formação e atuação em psicologia clínica. Este trabalho tem
por objetivo introduzir o tema de aprendizagem de máquina a psicólogos. Busca-se apresentar, de maneira acessível, os conceitos e terminologia básicos
em aprendizagem de máquina, proporcionando um panorama geral da área. Para complementar a exposição conceitual, é construído um exemplo a partir de dados
de um estudo sobre eficácia de intervenções em psicologia positiva para depressão via internet. Por fim, discutem-se algumas implicações do uso desse tipo
de tecnologia para a pesquisa e para a prática em psicoterapia.

 \textbf{Palavras-chave}: aprendizagem de máquina, psicoterapia, psicologia positiva, depressão.
\end{resumo}

% Resumo em inglês
\begin{resumo}[Abstract]
 \begin{otherlanguage*}{english}
   Advances in machine learning have opened up a multitude of opportunities for various research and professional fields. The use of machine learning models to
   predict treatment outcomes in mental health has been investigated showing promising results. The application of such models would allow greater assertiveness
   in tailoring psychotherapeutic treatments at the individual level. Despite its potential benefits, the topic remains unexplored in the fields of research, training
   and practice in clinical psychology. This work aims to introduce the topic of machine learning to psychologists. It presents the basic machine learning concepts and
   terminology in an accessible manner, providing a comprehensive overview of the field. To complement the conceptual discussion, an example is constructed using data
   from a study on the effectiveness of internet-based positive psychology interventions for depression. Finally, implications of using this type of technology in
   psychotherapy research and practice are discussed.
   \vspace{\onelineskip}
 
   \noindent 
   \textbf{Keywords}: machine learning, psychotherapy, positive psychology, depression.
 \end{otherlanguage*}
\end{resumo}


% ---
% LISTA DE ILUSTRAÇÕES
% ---
\pdfbookmark[0]{\listfigurename}{lof}
\listoffigures*
\cleardoublepage
% ---


% ---
% LISTA DE QUADROS
% ---
% \pdfbookmark[0]{\listofquadrosname}{loq}
% \listofquadros*
% \cleardoublepage
% ---


% ---
% LISTA DE TABELAS
% ---
% \pdfbookmark[0]{\listtablename}{lot}
% \listoftables*
% \cleardoublepage
% ---


% ---
% LISTA DE ABREVIATURAS E SIGLAS
% ---
% \begin{siglas}
%   \item[ABNT] Associação Brasileira de Normas Técnicas
%   \item[abnTeX] ABsurdas Normas para TeX
% \end{siglas}
% ---


% ---
% LISTA DE SÍMBOLOS
% ---
% \begin{simbolos}
%   \item[$ \Gamma $] Letra grega Gama
%   \item[$ \Lambda $] Lambda
%   \item[$ \zeta $] Letra grega minúscula zeta
%   \item[$ \in $] Pertence
% \end{simbolos}
% ---


% ---
% Sumário
% ---
\pdfbookmark[0]{\contentsname}{toc}
\tableofcontents*
\cleardoublepage
% ---


% ----------------------------------------------------------
% ELEMENTOS TEXTUAIS
% ----------------------------------------------------------
\textual

\documentclass[
	% Opções da classe memoir
	12pt,				        % tamanho da fonte
	openany,			        % capítulos começam em pág ímpar (insere página vazia caso preciso)
	oneside,			        % para impressão em recto e verso. Oposto a oneside
	a4paper,			        % tamanho do papel. 
	% Opções da classe abntex2
	%chapter=TITLE,		        % títulos de capítulos convertidos em letras maiúsculas
	%section=TITLE,		        % títulos de seções convertidos em letras maiúsculas
	%subsection=TITLE,	        % títulos de subseções convertidos em letras maiúsculas
	%subsubsection=TITLE,       % títulos de subsubseções convertidos em letras maiúsculas
	% Opções do pacote babel
	english,			        % idioma adicional para hifenização
	brazil				        % o último idioma é o principal do documento
	]{abntex2}

% ---
% Pacotes básicos 
% ---
\usepackage{lmodern}			% Usa a fonte Latin Modern			
\usepackage[T1]{fontenc}		% Selecao de codigos de fonte.
\usepackage[utf8]{inputenc}		% Codificacao do documento (conversão automática dos acentos)
\usepackage{indentfirst}		% Indenta o primeiro parágrafo de cada seção.
\usepackage{color}				% Controle das cores
\usepackage{graphicx}			% Inclusão de gráficos
\usepackage{microtype} 			% para melhorias de justificação
% ---
		

% ---
% Pacotes de citações
% ---
% \usepackage[alf]{abntex2cite}						% Citações padrão ABNT
\usepackage[
  alf,
  abnt-nbr10520=1988,                     % Citações em sentence case
  abnt-emphasize=bf
  ]{abntex2cite}


% ---
% Pacotes adicionais
% ---
\usepackage{pdfpages}	% Inclusão de páginas em PDF


% ---
% Informações para CAPA e FOLHA DE ROSTO
% ---
\titulo{APRENDIZAGEM DE MÁQUINA PARA PREDIÇÃO DE DESFECHOS EM SAÚDE MENTAL: \textmd{UMA INTRODUÇÃO PARA PSICÓLOGOS}}
\autor{GUILHERME AUGUSTO DIAS}
\local{Porto Alegre}
\data{2024}
\orientador{Prof. Dr. Wagner de Lara Machado}
% \coorientador{}
% \instituicao{
%   Pontifícia Universidade Católica do Rio Grande do Sul
%   \par
%   Escola de Ciências da Saúde e da Vida
%   \par
%   Curso de Psicologia}
\tipotrabalho{TCC}
\preambulo{Trabalho de Conclusão de Curso apresentado como requisito para a obtenção do grau de Bacharel em Psicologia pela Escola de Ciências da Saúde e da Vida da Pontifícia Universidade Católica do Rio Grande do Sul.}
% ---

% ---
% Configurações de aparência do PDF final
% ---

% Alterando o aspecto da cor azul
\definecolor{blue}{RGB}{41,5,195}

% Informações do PDF
\makeatletter
\hypersetup{
		pdftitle={Aprendizagem de máquina para predição de desfechos em saúde mental: uma introdução para psicólogos}, 
		pdfauthor={Guilherme Augusto Dias},
    	pdfsubject={\imprimirpreambulo},
	    pdfcreator={LaTeX with abnTeX2},
		pdfkeywords={aprendizagem de máquina}{psicoterapia}{psicologia positiva}{depressão}, 
		colorlinks=true,		% true: colored links; false: boxed links
    	linkcolor=black,		% color of internal links
    	citecolor=black,		% color of links to bibliography
    	filecolor=black,		% color of file links
		urlcolor=black,
		bookmarksdepth=4
}
\makeatother
% --- 


% ---
% Posiciona figuras e tabelas no topo da página quando adicionadas sozinhas
% em um página em branco. Ver https://github.com/abntex/abntex2/issues/170
\makeatletter
\setlength{\@fptop}{5pt} % Set distance from top of page to first float
\makeatother
% ---

% ---
% Possibilita criação de Quadros e Lista de quadros.
% Ver https://github.com/abntex/abntex2/issues/176
%
\newcommand{\quadroname}{Quadro}
\newcommand{\listofquadrosname}{Lista de quadros}

\newfloat[chapter]{quadro}{loq}{\quadroname}
\newlistof{listofquadros}{loq}{\listofquadrosname}
\newlistentry{quadro}{loq}{0}

% Configurações para atender às regras da ABNT
\setfloatadjustment{quadro}{\centering}
\counterwithout{quadro}{chapter}
\renewcommand{\cftquadroname}{\quadroname\space} 
\renewcommand*{\cftquadroaftersnum}{\hfill--\hfill}

\setfloatlocations{quadro}{hbtp} % Ver https://github.com/abntex/abntex2/issues/176
% ---


% --- 
% Espaçamentos entre linhas e parágrafos 
% --- 

% O tamanho do parágrafo é dado por:
\setlength{\parindent}{1.3cm}

% Controle do espaçamento entre um parágrafo e outro:
\setlength{\parskip}{0.2cm}  % tente também \onelineskip

% ---
% Compila o índice
% ---
\makeindex
% ---


% ----
% Início do documento
% ----
\begin{document}

% Seleciona o idioma do documento (conforme pacotes do babel)
\selectlanguage{brazil}

% Retira espaço extra obsoleto entre as frases.
\frenchspacing 

% ----------------------------------------------------------
% ELEMENTOS PRÉ-TEXTUAIS
% ----------------------------------------------------------
\pretextual

% ---
% CAPA
% ---
\includepdf{./05-extras/capa.pdf}
% ---


% ---
% FOLHA DE ROSTO
% (o * indica que haverá a ficha bibliográfica)
% ---
\imprimirfolhaderosto
% ---


% ---
% Ficha bibliografica
% ---

% Isto é um exemplo de Ficha Catalográfica, ou ``Dados internacionais de
% catalogação-na-publicação''. Você pode utilizar este modelo como referência. 
% Porém, provavelmente a biblioteca da sua universidade lhe fornecerá um PDF
% com a ficha catalográfica definitiva após a defesa do trabalho. Quando estiver
% com o documento, salve-o como PDF no diretório do seu projeto e substitua todo
% o conteúdo de implementação deste arquivo pelo comando abaixo:
%
% \usepackage{pdfpages}		% necessário para comando \includepdf
% \begin{fichacatalografica}
%     \includepdf{fig_ficha_catalografica.pdf}
% \end{fichacatalografica}

% \begin{fichacatalografica}
% 	\sffamily
% 	\vspace*{\fill}                             % Posição vertical
% 	\begin{center}                              % Minipage Centralizado
% 	\fbox{\begin{minipage}[c][8cm]{13.5cm}		% Largura
% 	\small
% 	\imprimirautor
% 	%Sobrenome, Nome do autor
	
% 	\hspace{0.5cm} \imprimirtitulo  / \imprimirautor. --
% 	\imprimirlocal, \imprimirdata-
	
% 	\hspace{0.5cm} \thelastpage p. : il. (algumas color.) ; 30 cm.\\
	
% 	\hspace{0.5cm} \imprimirorientadorRotulo~\imprimirorientador\\
	
% 	\hspace{0.5cm}
% 	\parbox[t]{\textwidth}{\imprimirtipotrabalho~--~\imprimirinstituicao,
% 	\imprimirdata.}\\
	
% 	\hspace{0.5cm}
% 		1. Palavra-chave1.
% 		2. Palavra-chave2.
% 		2. Palavra-chave3.
% 		I. Orientador.
% 		II. Universidade xxx.
% 		III. Faculdade de xxx.
% 		IV. Título 			
% 	\end{minipage}}
% 	\end{center}
% \end{fichacatalografica}
% ---


% ---
% FOLHA DE APROVAÇÃO
% ---

% Isto é um exemplo de Folha de aprovação, elemento obrigatório da NBR
% 14724/2011 (seção 4.2.1.3). Você pode utilizar este modelo até a aprovação
% do trabalho. Após isso, substitua todo o conteúdo deste arquivo por uma
% imagem da página assinada pela banca com o comando abaixo:
%
% \begin{folhadeaprovacao}
% \includepdf{folhadeaprovacao_final.pdf}
% \end{folhadeaprovacao}
%
\setlength{\ABNTEXsignskip}{2cm}
\setlength{\ABNTEXsignwidth}{10cm}

\begin{folhadeaprovacao}
  \begin{center}
    {\ABNTEXchapterfont\large\imprimirautor}

    \vspace*{1cm}

    \begin{center}
      \ABNTEXchapterfont\bfseries\Large\imprimirtitulo
    \end{center}

    \vspace*{1cm}
    
    \hspace{.45\textwidth}
    \begin{minipage}{.5\textwidth}
        \imprimirpreambulo
    \end{minipage}%
   \end{center}

   \vspace*{1cm}
        
   \begin{center}
   Aprovado em 17 de julho de 2024.

   \vspace*{1cm}

   Banca examinadora:
   \end{center}

   \vspace*{-1cm}

   \assinatura{Prof. Ma. Caroline Santa Maria Rodrigues -- PUCRS}
   \assinatura{Prof. Dra. Clarissa Pinto Pizarro de Freitas -- PUCRS}
   \assinatura{\imprimirorientador -- PUCRS}
\end{folhadeaprovacao}
% ---


% ---
% Agradecimentos
% ---
\begin{agradecimentos}
Agradeço aos meus pais, Ademilson José Dias e Irinês Porfírio Dias, pelo apoio incondicional que me deram ao longo da vida,
incentivando mesmo as ideias questionáveis como aquela de fazer uma segunda graduação.

A minha irmã, Maria Eduarda Dias, cuja trajetória acadêmica é fonte de inspiração.

A minha namorada, Nathalia Fallavena Ceratti, que acompanhou a esta jornada de perto, oferecendo apoio e compreensão que
fizeram toda a diferença.

Aos meus amigos, Guilherme da Silva Petró, Norton Neves e Uilinton Luciano Marques da Silva, que foram compreensíveis nos
finais de semana em que não pude jogar RPG.

Ao professor Wagner de Lara Machado pela orientação e apoio no desenvolvimento deste trabalho.
\end{agradecimentos}
% ---


% ---
% RESUMOS
% ---
% Resumo em português
\setlength{\absparsep}{18pt} % ajusta o espaçamento dos parágrafos do resumo
\begin{resumo}
Desenvolvimentos na área de aprendizagem de máquina representam uma série de oportunidades para diversas áreas de pesquisa e de atuação profissional.
O uso de modelos de aprendizagem de máquina na predição de desfecho para tratamentos em saúde mental tem sido investigado e apresenta resultados promissores.
A aplicação desse tipo de modelo possibilitaria, por exemplo, maior assertividade na personalização de tratamentos psicoterápicos a nível individual.
Apesar dos benefícios potenciais, o tema ainda é pouco discutido nos espaços de pesquisa, formação e atuação em psicologia clínica. Este trabalho tem
por objetivo introduzir o tema de aprendizagem de máquina a psicólogos. Busca-se apresentar, de maneira acessível, os conceitos e terminologia básicos
em aprendizagem de máquina, proporcionando um panorama geral da área. Para complementar a exposição conceitual, é construído um exemplo a partir de dados
de um estudo sobre eficácia de intervenções em psicologia positiva para depressão via internet. Por fim, discutem-se algumas implicações do uso desse tipo
de tecnologia para a pesquisa e para a prática em psicoterapia.

 \textbf{Palavras-chave}: aprendizagem de máquina, psicoterapia, psicologia positiva, depressão.
\end{resumo}

% Resumo em inglês
\begin{resumo}[Abstract]
 \begin{otherlanguage*}{english}
   Advances in machine learning have opened up a multitude of opportunities for various research and professional fields. The use of machine learning models to
   predict treatment outcomes in mental health has been investigated showing promising results. The application of such models would allow greater assertiveness
   in tailoring psychotherapeutic treatments at the individual level. Despite its potential benefits, the topic remains unexplored in the fields of research, training
   and practice in clinical psychology. This work aims to introduce the topic of machine learning to psychologists. It presents the basic machine learning concepts and
   terminology in an accessible manner, providing a comprehensive overview of the field. To complement the conceptual discussion, an example is constructed using data
   from a study on the effectiveness of internet-based positive psychology interventions for depression. Finally, implications of using this type of technology in
   psychotherapy research and practice are discussed.
   \vspace{\onelineskip}
 
   \noindent 
   \textbf{Keywords}: machine learning, psychotherapy, positive psychology, depression.
 \end{otherlanguage*}
\end{resumo}


% ---
% LISTA DE ILUSTRAÇÕES
% ---
\pdfbookmark[0]{\listfigurename}{lof}
\listoffigures*
\cleardoublepage
% ---


% ---
% LISTA DE QUADROS
% ---
% \pdfbookmark[0]{\listofquadrosname}{loq}
% \listofquadros*
% \cleardoublepage
% ---


% ---
% LISTA DE TABELAS
% ---
% \pdfbookmark[0]{\listtablename}{lot}
% \listoftables*
% \cleardoublepage
% ---


% ---
% LISTA DE ABREVIATURAS E SIGLAS
% ---
% \begin{siglas}
%   \item[ABNT] Associação Brasileira de Normas Técnicas
%   \item[abnTeX] ABsurdas Normas para TeX
% \end{siglas}
% ---


% ---
% LISTA DE SÍMBOLOS
% ---
% \begin{simbolos}
%   \item[$ \Gamma $] Letra grega Gama
%   \item[$ \Lambda $] Lambda
%   \item[$ \zeta $] Letra grega minúscula zeta
%   \item[$ \in $] Pertence
% \end{simbolos}
% ---


% ---
% Sumário
% ---
\pdfbookmark[0]{\contentsname}{toc}
\tableofcontents*
\cleardoublepage
% ---


% ----------------------------------------------------------
% ELEMENTOS TEXTUAIS
% ----------------------------------------------------------
\textual

\include{./01-introducao/main.tex}

\include{./02-machine-learning/main.tex}

\include{./03-exemplo/main.tex}

\include{./04-conclusao/main.tex}

% ----------------------------------------------------------
% ELEMENTOS PÓS-TEXTUAIS
% ----------------------------------------------------------
\postextual
% ----------------------------------------------------------


% ----------------------------------------------------------
% REFERÊNCIAS
% ----------------------------------------------------------
\bibliography{references.bib}


% ----------------------------------------------------------
% GLOSSÁRIO
% ----------------------------------------------------------
%
% Consulte o manual da classe abntex2 para orientações sobre o glossário.
%
%\glossary

% ----------------------------------------------------------
% APÊNDICES
% ----------------------------------------------------------

% ---
% Inicia os apêndices
% ---
% \begin{apendicesenv}

% % Imprime uma página indicando o início dos apêndices
% \partapendices

% % ----------------------------------------------------------
% \chapter{Quisque libero justo}
% % ----------------------------------------------------------

% % ----------------------------------------------------------
% \chapter{Nullam elementum urna}
% % ----------------------------------------------------------

% \end{apendicesenv}
% ---


% ----------------------------------------------------------
% ANEXOS
% ----------------------------------------------------------

% ---
% Inicia os anexos
% ---
% \begin{anexosenv}

% % Imprime uma página indicando o início dos anexos
% \partanexos

% % ---
% \chapter{Morbi ultrices rutrum lorem}
% % ---

% % ---
% \chapter{Cras non urna sed feugiat cum}
% % ---


% % ---
% \chapter{Fusce facilisis lacinia dui}
% % ---

% \end{anexosenv}

%---------------------------------------------------------------------
% INDICE REMISSIVO
%---------------------------------------------------------------------
% \phantompart
% \printindex
%---------------------------------------------------------------------
\includepdf{./05-extras/contracapa.pdf}
\end{document}


\documentclass[
	% Opções da classe memoir
	12pt,				        % tamanho da fonte
	openany,			        % capítulos começam em pág ímpar (insere página vazia caso preciso)
	oneside,			        % para impressão em recto e verso. Oposto a oneside
	a4paper,			        % tamanho do papel. 
	% Opções da classe abntex2
	%chapter=TITLE,		        % títulos de capítulos convertidos em letras maiúsculas
	%section=TITLE,		        % títulos de seções convertidos em letras maiúsculas
	%subsection=TITLE,	        % títulos de subseções convertidos em letras maiúsculas
	%subsubsection=TITLE,       % títulos de subsubseções convertidos em letras maiúsculas
	% Opções do pacote babel
	english,			        % idioma adicional para hifenização
	brazil				        % o último idioma é o principal do documento
	]{abntex2}

% ---
% Pacotes básicos 
% ---
\usepackage{lmodern}			% Usa a fonte Latin Modern			
\usepackage[T1]{fontenc}		% Selecao de codigos de fonte.
\usepackage[utf8]{inputenc}		% Codificacao do documento (conversão automática dos acentos)
\usepackage{indentfirst}		% Indenta o primeiro parágrafo de cada seção.
\usepackage{color}				% Controle das cores
\usepackage{graphicx}			% Inclusão de gráficos
\usepackage{microtype} 			% para melhorias de justificação
% ---
		

% ---
% Pacotes de citações
% ---
% \usepackage[alf]{abntex2cite}						% Citações padrão ABNT
\usepackage[
  alf,
  abnt-nbr10520=1988,                     % Citações em sentence case
  abnt-emphasize=bf
  ]{abntex2cite}


% ---
% Pacotes adicionais
% ---
\usepackage{pdfpages}	% Inclusão de páginas em PDF


% ---
% Informações para CAPA e FOLHA DE ROSTO
% ---
\titulo{APRENDIZAGEM DE MÁQUINA PARA PREDIÇÃO DE DESFECHOS EM SAÚDE MENTAL: \textmd{UMA INTRODUÇÃO PARA PSICÓLOGOS}}
\autor{GUILHERME AUGUSTO DIAS}
\local{Porto Alegre}
\data{2024}
\orientador{Prof. Dr. Wagner de Lara Machado}
% \coorientador{}
% \instituicao{
%   Pontifícia Universidade Católica do Rio Grande do Sul
%   \par
%   Escola de Ciências da Saúde e da Vida
%   \par
%   Curso de Psicologia}
\tipotrabalho{TCC}
\preambulo{Trabalho de Conclusão de Curso apresentado como requisito para a obtenção do grau de Bacharel em Psicologia pela Escola de Ciências da Saúde e da Vida da Pontifícia Universidade Católica do Rio Grande do Sul.}
% ---

% ---
% Configurações de aparência do PDF final
% ---

% Alterando o aspecto da cor azul
\definecolor{blue}{RGB}{41,5,195}

% Informações do PDF
\makeatletter
\hypersetup{
		pdftitle={Aprendizagem de máquina para predição de desfechos em saúde mental: uma introdução para psicólogos}, 
		pdfauthor={Guilherme Augusto Dias},
    	pdfsubject={\imprimirpreambulo},
	    pdfcreator={LaTeX with abnTeX2},
		pdfkeywords={aprendizagem de máquina}{psicoterapia}{psicologia positiva}{depressão}, 
		colorlinks=true,		% true: colored links; false: boxed links
    	linkcolor=black,		% color of internal links
    	citecolor=black,		% color of links to bibliography
    	filecolor=black,		% color of file links
		urlcolor=black,
		bookmarksdepth=4
}
\makeatother
% --- 


% ---
% Posiciona figuras e tabelas no topo da página quando adicionadas sozinhas
% em um página em branco. Ver https://github.com/abntex/abntex2/issues/170
\makeatletter
\setlength{\@fptop}{5pt} % Set distance from top of page to first float
\makeatother
% ---

% ---
% Possibilita criação de Quadros e Lista de quadros.
% Ver https://github.com/abntex/abntex2/issues/176
%
\newcommand{\quadroname}{Quadro}
\newcommand{\listofquadrosname}{Lista de quadros}

\newfloat[chapter]{quadro}{loq}{\quadroname}
\newlistof{listofquadros}{loq}{\listofquadrosname}
\newlistentry{quadro}{loq}{0}

% Configurações para atender às regras da ABNT
\setfloatadjustment{quadro}{\centering}
\counterwithout{quadro}{chapter}
\renewcommand{\cftquadroname}{\quadroname\space} 
\renewcommand*{\cftquadroaftersnum}{\hfill--\hfill}

\setfloatlocations{quadro}{hbtp} % Ver https://github.com/abntex/abntex2/issues/176
% ---


% --- 
% Espaçamentos entre linhas e parágrafos 
% --- 

% O tamanho do parágrafo é dado por:
\setlength{\parindent}{1.3cm}

% Controle do espaçamento entre um parágrafo e outro:
\setlength{\parskip}{0.2cm}  % tente também \onelineskip

% ---
% Compila o índice
% ---
\makeindex
% ---


% ----
% Início do documento
% ----
\begin{document}

% Seleciona o idioma do documento (conforme pacotes do babel)
\selectlanguage{brazil}

% Retira espaço extra obsoleto entre as frases.
\frenchspacing 

% ----------------------------------------------------------
% ELEMENTOS PRÉ-TEXTUAIS
% ----------------------------------------------------------
\pretextual

% ---
% CAPA
% ---
\includepdf{./05-extras/capa.pdf}
% ---


% ---
% FOLHA DE ROSTO
% (o * indica que haverá a ficha bibliográfica)
% ---
\imprimirfolhaderosto
% ---


% ---
% Ficha bibliografica
% ---

% Isto é um exemplo de Ficha Catalográfica, ou ``Dados internacionais de
% catalogação-na-publicação''. Você pode utilizar este modelo como referência. 
% Porém, provavelmente a biblioteca da sua universidade lhe fornecerá um PDF
% com a ficha catalográfica definitiva após a defesa do trabalho. Quando estiver
% com o documento, salve-o como PDF no diretório do seu projeto e substitua todo
% o conteúdo de implementação deste arquivo pelo comando abaixo:
%
% \usepackage{pdfpages}		% necessário para comando \includepdf
% \begin{fichacatalografica}
%     \includepdf{fig_ficha_catalografica.pdf}
% \end{fichacatalografica}

% \begin{fichacatalografica}
% 	\sffamily
% 	\vspace*{\fill}                             % Posição vertical
% 	\begin{center}                              % Minipage Centralizado
% 	\fbox{\begin{minipage}[c][8cm]{13.5cm}		% Largura
% 	\small
% 	\imprimirautor
% 	%Sobrenome, Nome do autor
	
% 	\hspace{0.5cm} \imprimirtitulo  / \imprimirautor. --
% 	\imprimirlocal, \imprimirdata-
	
% 	\hspace{0.5cm} \thelastpage p. : il. (algumas color.) ; 30 cm.\\
	
% 	\hspace{0.5cm} \imprimirorientadorRotulo~\imprimirorientador\\
	
% 	\hspace{0.5cm}
% 	\parbox[t]{\textwidth}{\imprimirtipotrabalho~--~\imprimirinstituicao,
% 	\imprimirdata.}\\
	
% 	\hspace{0.5cm}
% 		1. Palavra-chave1.
% 		2. Palavra-chave2.
% 		2. Palavra-chave3.
% 		I. Orientador.
% 		II. Universidade xxx.
% 		III. Faculdade de xxx.
% 		IV. Título 			
% 	\end{minipage}}
% 	\end{center}
% \end{fichacatalografica}
% ---


% ---
% FOLHA DE APROVAÇÃO
% ---

% Isto é um exemplo de Folha de aprovação, elemento obrigatório da NBR
% 14724/2011 (seção 4.2.1.3). Você pode utilizar este modelo até a aprovação
% do trabalho. Após isso, substitua todo o conteúdo deste arquivo por uma
% imagem da página assinada pela banca com o comando abaixo:
%
% \begin{folhadeaprovacao}
% \includepdf{folhadeaprovacao_final.pdf}
% \end{folhadeaprovacao}
%
\setlength{\ABNTEXsignskip}{2cm}
\setlength{\ABNTEXsignwidth}{10cm}

\begin{folhadeaprovacao}
  \begin{center}
    {\ABNTEXchapterfont\large\imprimirautor}

    \vspace*{1cm}

    \begin{center}
      \ABNTEXchapterfont\bfseries\Large\imprimirtitulo
    \end{center}

    \vspace*{1cm}
    
    \hspace{.45\textwidth}
    \begin{minipage}{.5\textwidth}
        \imprimirpreambulo
    \end{minipage}%
   \end{center}

   \vspace*{1cm}
        
   \begin{center}
   Aprovado em 17 de julho de 2024.

   \vspace*{1cm}

   Banca examinadora:
   \end{center}

   \vspace*{-1cm}

   \assinatura{Prof. Ma. Caroline Santa Maria Rodrigues -- PUCRS}
   \assinatura{Prof. Dra. Clarissa Pinto Pizarro de Freitas -- PUCRS}
   \assinatura{\imprimirorientador -- PUCRS}
\end{folhadeaprovacao}
% ---


% ---
% Agradecimentos
% ---
\begin{agradecimentos}
Agradeço aos meus pais, Ademilson José Dias e Irinês Porfírio Dias, pelo apoio incondicional que me deram ao longo da vida,
incentivando mesmo as ideias questionáveis como aquela de fazer uma segunda graduação.

A minha irmã, Maria Eduarda Dias, cuja trajetória acadêmica é fonte de inspiração.

A minha namorada, Nathalia Fallavena Ceratti, que acompanhou a esta jornada de perto, oferecendo apoio e compreensão que
fizeram toda a diferença.

Aos meus amigos, Guilherme da Silva Petró, Norton Neves e Uilinton Luciano Marques da Silva, que foram compreensíveis nos
finais de semana em que não pude jogar RPG.

Ao professor Wagner de Lara Machado pela orientação e apoio no desenvolvimento deste trabalho.
\end{agradecimentos}
% ---


% ---
% RESUMOS
% ---
% Resumo em português
\setlength{\absparsep}{18pt} % ajusta o espaçamento dos parágrafos do resumo
\begin{resumo}
Desenvolvimentos na área de aprendizagem de máquina representam uma série de oportunidades para diversas áreas de pesquisa e de atuação profissional.
O uso de modelos de aprendizagem de máquina na predição de desfecho para tratamentos em saúde mental tem sido investigado e apresenta resultados promissores.
A aplicação desse tipo de modelo possibilitaria, por exemplo, maior assertividade na personalização de tratamentos psicoterápicos a nível individual.
Apesar dos benefícios potenciais, o tema ainda é pouco discutido nos espaços de pesquisa, formação e atuação em psicologia clínica. Este trabalho tem
por objetivo introduzir o tema de aprendizagem de máquina a psicólogos. Busca-se apresentar, de maneira acessível, os conceitos e terminologia básicos
em aprendizagem de máquina, proporcionando um panorama geral da área. Para complementar a exposição conceitual, é construído um exemplo a partir de dados
de um estudo sobre eficácia de intervenções em psicologia positiva para depressão via internet. Por fim, discutem-se algumas implicações do uso desse tipo
de tecnologia para a pesquisa e para a prática em psicoterapia.

 \textbf{Palavras-chave}: aprendizagem de máquina, psicoterapia, psicologia positiva, depressão.
\end{resumo}

% Resumo em inglês
\begin{resumo}[Abstract]
 \begin{otherlanguage*}{english}
   Advances in machine learning have opened up a multitude of opportunities for various research and professional fields. The use of machine learning models to
   predict treatment outcomes in mental health has been investigated showing promising results. The application of such models would allow greater assertiveness
   in tailoring psychotherapeutic treatments at the individual level. Despite its potential benefits, the topic remains unexplored in the fields of research, training
   and practice in clinical psychology. This work aims to introduce the topic of machine learning to psychologists. It presents the basic machine learning concepts and
   terminology in an accessible manner, providing a comprehensive overview of the field. To complement the conceptual discussion, an example is constructed using data
   from a study on the effectiveness of internet-based positive psychology interventions for depression. Finally, implications of using this type of technology in
   psychotherapy research and practice are discussed.
   \vspace{\onelineskip}
 
   \noindent 
   \textbf{Keywords}: machine learning, psychotherapy, positive psychology, depression.
 \end{otherlanguage*}
\end{resumo}


% ---
% LISTA DE ILUSTRAÇÕES
% ---
\pdfbookmark[0]{\listfigurename}{lof}
\listoffigures*
\cleardoublepage
% ---


% ---
% LISTA DE QUADROS
% ---
% \pdfbookmark[0]{\listofquadrosname}{loq}
% \listofquadros*
% \cleardoublepage
% ---


% ---
% LISTA DE TABELAS
% ---
% \pdfbookmark[0]{\listtablename}{lot}
% \listoftables*
% \cleardoublepage
% ---


% ---
% LISTA DE ABREVIATURAS E SIGLAS
% ---
% \begin{siglas}
%   \item[ABNT] Associação Brasileira de Normas Técnicas
%   \item[abnTeX] ABsurdas Normas para TeX
% \end{siglas}
% ---


% ---
% LISTA DE SÍMBOLOS
% ---
% \begin{simbolos}
%   \item[$ \Gamma $] Letra grega Gama
%   \item[$ \Lambda $] Lambda
%   \item[$ \zeta $] Letra grega minúscula zeta
%   \item[$ \in $] Pertence
% \end{simbolos}
% ---


% ---
% Sumário
% ---
\pdfbookmark[0]{\contentsname}{toc}
\tableofcontents*
\cleardoublepage
% ---


% ----------------------------------------------------------
% ELEMENTOS TEXTUAIS
% ----------------------------------------------------------
\textual

\include{./01-introducao/main.tex}

\include{./02-machine-learning/main.tex}

\include{./03-exemplo/main.tex}

\include{./04-conclusao/main.tex}

% ----------------------------------------------------------
% ELEMENTOS PÓS-TEXTUAIS
% ----------------------------------------------------------
\postextual
% ----------------------------------------------------------


% ----------------------------------------------------------
% REFERÊNCIAS
% ----------------------------------------------------------
\bibliography{references.bib}


% ----------------------------------------------------------
% GLOSSÁRIO
% ----------------------------------------------------------
%
% Consulte o manual da classe abntex2 para orientações sobre o glossário.
%
%\glossary

% ----------------------------------------------------------
% APÊNDICES
% ----------------------------------------------------------

% ---
% Inicia os apêndices
% ---
% \begin{apendicesenv}

% % Imprime uma página indicando o início dos apêndices
% \partapendices

% % ----------------------------------------------------------
% \chapter{Quisque libero justo}
% % ----------------------------------------------------------

% % ----------------------------------------------------------
% \chapter{Nullam elementum urna}
% % ----------------------------------------------------------

% \end{apendicesenv}
% ---


% ----------------------------------------------------------
% ANEXOS
% ----------------------------------------------------------

% ---
% Inicia os anexos
% ---
% \begin{anexosenv}

% % Imprime uma página indicando o início dos anexos
% \partanexos

% % ---
% \chapter{Morbi ultrices rutrum lorem}
% % ---

% % ---
% \chapter{Cras non urna sed feugiat cum}
% % ---


% % ---
% \chapter{Fusce facilisis lacinia dui}
% % ---

% \end{anexosenv}

%---------------------------------------------------------------------
% INDICE REMISSIVO
%---------------------------------------------------------------------
% \phantompart
% \printindex
%---------------------------------------------------------------------
\includepdf{./05-extras/contracapa.pdf}
\end{document}


\documentclass[
	% Opções da classe memoir
	12pt,				        % tamanho da fonte
	openany,			        % capítulos começam em pág ímpar (insere página vazia caso preciso)
	oneside,			        % para impressão em recto e verso. Oposto a oneside
	a4paper,			        % tamanho do papel. 
	% Opções da classe abntex2
	%chapter=TITLE,		        % títulos de capítulos convertidos em letras maiúsculas
	%section=TITLE,		        % títulos de seções convertidos em letras maiúsculas
	%subsection=TITLE,	        % títulos de subseções convertidos em letras maiúsculas
	%subsubsection=TITLE,       % títulos de subsubseções convertidos em letras maiúsculas
	% Opções do pacote babel
	english,			        % idioma adicional para hifenização
	brazil				        % o último idioma é o principal do documento
	]{abntex2}

% ---
% Pacotes básicos 
% ---
\usepackage{lmodern}			% Usa a fonte Latin Modern			
\usepackage[T1]{fontenc}		% Selecao de codigos de fonte.
\usepackage[utf8]{inputenc}		% Codificacao do documento (conversão automática dos acentos)
\usepackage{indentfirst}		% Indenta o primeiro parágrafo de cada seção.
\usepackage{color}				% Controle das cores
\usepackage{graphicx}			% Inclusão de gráficos
\usepackage{microtype} 			% para melhorias de justificação
% ---
		

% ---
% Pacotes de citações
% ---
% \usepackage[alf]{abntex2cite}						% Citações padrão ABNT
\usepackage[
  alf,
  abnt-nbr10520=1988,                     % Citações em sentence case
  abnt-emphasize=bf
  ]{abntex2cite}


% ---
% Pacotes adicionais
% ---
\usepackage{pdfpages}	% Inclusão de páginas em PDF


% ---
% Informações para CAPA e FOLHA DE ROSTO
% ---
\titulo{APRENDIZAGEM DE MÁQUINA PARA PREDIÇÃO DE DESFECHOS EM SAÚDE MENTAL: \textmd{UMA INTRODUÇÃO PARA PSICÓLOGOS}}
\autor{GUILHERME AUGUSTO DIAS}
\local{Porto Alegre}
\data{2024}
\orientador{Prof. Dr. Wagner de Lara Machado}
% \coorientador{}
% \instituicao{
%   Pontifícia Universidade Católica do Rio Grande do Sul
%   \par
%   Escola de Ciências da Saúde e da Vida
%   \par
%   Curso de Psicologia}
\tipotrabalho{TCC}
\preambulo{Trabalho de Conclusão de Curso apresentado como requisito para a obtenção do grau de Bacharel em Psicologia pela Escola de Ciências da Saúde e da Vida da Pontifícia Universidade Católica do Rio Grande do Sul.}
% ---

% ---
% Configurações de aparência do PDF final
% ---

% Alterando o aspecto da cor azul
\definecolor{blue}{RGB}{41,5,195}

% Informações do PDF
\makeatletter
\hypersetup{
		pdftitle={Aprendizagem de máquina para predição de desfechos em saúde mental: uma introdução para psicólogos}, 
		pdfauthor={Guilherme Augusto Dias},
    	pdfsubject={\imprimirpreambulo},
	    pdfcreator={LaTeX with abnTeX2},
		pdfkeywords={aprendizagem de máquina}{psicoterapia}{psicologia positiva}{depressão}, 
		colorlinks=true,		% true: colored links; false: boxed links
    	linkcolor=black,		% color of internal links
    	citecolor=black,		% color of links to bibliography
    	filecolor=black,		% color of file links
		urlcolor=black,
		bookmarksdepth=4
}
\makeatother
% --- 


% ---
% Posiciona figuras e tabelas no topo da página quando adicionadas sozinhas
% em um página em branco. Ver https://github.com/abntex/abntex2/issues/170
\makeatletter
\setlength{\@fptop}{5pt} % Set distance from top of page to first float
\makeatother
% ---

% ---
% Possibilita criação de Quadros e Lista de quadros.
% Ver https://github.com/abntex/abntex2/issues/176
%
\newcommand{\quadroname}{Quadro}
\newcommand{\listofquadrosname}{Lista de quadros}

\newfloat[chapter]{quadro}{loq}{\quadroname}
\newlistof{listofquadros}{loq}{\listofquadrosname}
\newlistentry{quadro}{loq}{0}

% Configurações para atender às regras da ABNT
\setfloatadjustment{quadro}{\centering}
\counterwithout{quadro}{chapter}
\renewcommand{\cftquadroname}{\quadroname\space} 
\renewcommand*{\cftquadroaftersnum}{\hfill--\hfill}

\setfloatlocations{quadro}{hbtp} % Ver https://github.com/abntex/abntex2/issues/176
% ---


% --- 
% Espaçamentos entre linhas e parágrafos 
% --- 

% O tamanho do parágrafo é dado por:
\setlength{\parindent}{1.3cm}

% Controle do espaçamento entre um parágrafo e outro:
\setlength{\parskip}{0.2cm}  % tente também \onelineskip

% ---
% Compila o índice
% ---
\makeindex
% ---


% ----
% Início do documento
% ----
\begin{document}

% Seleciona o idioma do documento (conforme pacotes do babel)
\selectlanguage{brazil}

% Retira espaço extra obsoleto entre as frases.
\frenchspacing 

% ----------------------------------------------------------
% ELEMENTOS PRÉ-TEXTUAIS
% ----------------------------------------------------------
\pretextual

% ---
% CAPA
% ---
\includepdf{./05-extras/capa.pdf}
% ---


% ---
% FOLHA DE ROSTO
% (o * indica que haverá a ficha bibliográfica)
% ---
\imprimirfolhaderosto
% ---


% ---
% Ficha bibliografica
% ---

% Isto é um exemplo de Ficha Catalográfica, ou ``Dados internacionais de
% catalogação-na-publicação''. Você pode utilizar este modelo como referência. 
% Porém, provavelmente a biblioteca da sua universidade lhe fornecerá um PDF
% com a ficha catalográfica definitiva após a defesa do trabalho. Quando estiver
% com o documento, salve-o como PDF no diretório do seu projeto e substitua todo
% o conteúdo de implementação deste arquivo pelo comando abaixo:
%
% \usepackage{pdfpages}		% necessário para comando \includepdf
% \begin{fichacatalografica}
%     \includepdf{fig_ficha_catalografica.pdf}
% \end{fichacatalografica}

% \begin{fichacatalografica}
% 	\sffamily
% 	\vspace*{\fill}                             % Posição vertical
% 	\begin{center}                              % Minipage Centralizado
% 	\fbox{\begin{minipage}[c][8cm]{13.5cm}		% Largura
% 	\small
% 	\imprimirautor
% 	%Sobrenome, Nome do autor
	
% 	\hspace{0.5cm} \imprimirtitulo  / \imprimirautor. --
% 	\imprimirlocal, \imprimirdata-
	
% 	\hspace{0.5cm} \thelastpage p. : il. (algumas color.) ; 30 cm.\\
	
% 	\hspace{0.5cm} \imprimirorientadorRotulo~\imprimirorientador\\
	
% 	\hspace{0.5cm}
% 	\parbox[t]{\textwidth}{\imprimirtipotrabalho~--~\imprimirinstituicao,
% 	\imprimirdata.}\\
	
% 	\hspace{0.5cm}
% 		1. Palavra-chave1.
% 		2. Palavra-chave2.
% 		2. Palavra-chave3.
% 		I. Orientador.
% 		II. Universidade xxx.
% 		III. Faculdade de xxx.
% 		IV. Título 			
% 	\end{minipage}}
% 	\end{center}
% \end{fichacatalografica}
% ---


% ---
% FOLHA DE APROVAÇÃO
% ---

% Isto é um exemplo de Folha de aprovação, elemento obrigatório da NBR
% 14724/2011 (seção 4.2.1.3). Você pode utilizar este modelo até a aprovação
% do trabalho. Após isso, substitua todo o conteúdo deste arquivo por uma
% imagem da página assinada pela banca com o comando abaixo:
%
% \begin{folhadeaprovacao}
% \includepdf{folhadeaprovacao_final.pdf}
% \end{folhadeaprovacao}
%
\setlength{\ABNTEXsignskip}{2cm}
\setlength{\ABNTEXsignwidth}{10cm}

\begin{folhadeaprovacao}
  \begin{center}
    {\ABNTEXchapterfont\large\imprimirautor}

    \vspace*{1cm}

    \begin{center}
      \ABNTEXchapterfont\bfseries\Large\imprimirtitulo
    \end{center}

    \vspace*{1cm}
    
    \hspace{.45\textwidth}
    \begin{minipage}{.5\textwidth}
        \imprimirpreambulo
    \end{minipage}%
   \end{center}

   \vspace*{1cm}
        
   \begin{center}
   Aprovado em 17 de julho de 2024.

   \vspace*{1cm}

   Banca examinadora:
   \end{center}

   \vspace*{-1cm}

   \assinatura{Prof. Ma. Caroline Santa Maria Rodrigues -- PUCRS}
   \assinatura{Prof. Dra. Clarissa Pinto Pizarro de Freitas -- PUCRS}
   \assinatura{\imprimirorientador -- PUCRS}
\end{folhadeaprovacao}
% ---


% ---
% Agradecimentos
% ---
\begin{agradecimentos}
Agradeço aos meus pais, Ademilson José Dias e Irinês Porfírio Dias, pelo apoio incondicional que me deram ao longo da vida,
incentivando mesmo as ideias questionáveis como aquela de fazer uma segunda graduação.

A minha irmã, Maria Eduarda Dias, cuja trajetória acadêmica é fonte de inspiração.

A minha namorada, Nathalia Fallavena Ceratti, que acompanhou a esta jornada de perto, oferecendo apoio e compreensão que
fizeram toda a diferença.

Aos meus amigos, Guilherme da Silva Petró, Norton Neves e Uilinton Luciano Marques da Silva, que foram compreensíveis nos
finais de semana em que não pude jogar RPG.

Ao professor Wagner de Lara Machado pela orientação e apoio no desenvolvimento deste trabalho.
\end{agradecimentos}
% ---


% ---
% RESUMOS
% ---
% Resumo em português
\setlength{\absparsep}{18pt} % ajusta o espaçamento dos parágrafos do resumo
\begin{resumo}
Desenvolvimentos na área de aprendizagem de máquina representam uma série de oportunidades para diversas áreas de pesquisa e de atuação profissional.
O uso de modelos de aprendizagem de máquina na predição de desfecho para tratamentos em saúde mental tem sido investigado e apresenta resultados promissores.
A aplicação desse tipo de modelo possibilitaria, por exemplo, maior assertividade na personalização de tratamentos psicoterápicos a nível individual.
Apesar dos benefícios potenciais, o tema ainda é pouco discutido nos espaços de pesquisa, formação e atuação em psicologia clínica. Este trabalho tem
por objetivo introduzir o tema de aprendizagem de máquina a psicólogos. Busca-se apresentar, de maneira acessível, os conceitos e terminologia básicos
em aprendizagem de máquina, proporcionando um panorama geral da área. Para complementar a exposição conceitual, é construído um exemplo a partir de dados
de um estudo sobre eficácia de intervenções em psicologia positiva para depressão via internet. Por fim, discutem-se algumas implicações do uso desse tipo
de tecnologia para a pesquisa e para a prática em psicoterapia.

 \textbf{Palavras-chave}: aprendizagem de máquina, psicoterapia, psicologia positiva, depressão.
\end{resumo}

% Resumo em inglês
\begin{resumo}[Abstract]
 \begin{otherlanguage*}{english}
   Advances in machine learning have opened up a multitude of opportunities for various research and professional fields. The use of machine learning models to
   predict treatment outcomes in mental health has been investigated showing promising results. The application of such models would allow greater assertiveness
   in tailoring psychotherapeutic treatments at the individual level. Despite its potential benefits, the topic remains unexplored in the fields of research, training
   and practice in clinical psychology. This work aims to introduce the topic of machine learning to psychologists. It presents the basic machine learning concepts and
   terminology in an accessible manner, providing a comprehensive overview of the field. To complement the conceptual discussion, an example is constructed using data
   from a study on the effectiveness of internet-based positive psychology interventions for depression. Finally, implications of using this type of technology in
   psychotherapy research and practice are discussed.
   \vspace{\onelineskip}
 
   \noindent 
   \textbf{Keywords}: machine learning, psychotherapy, positive psychology, depression.
 \end{otherlanguage*}
\end{resumo}


% ---
% LISTA DE ILUSTRAÇÕES
% ---
\pdfbookmark[0]{\listfigurename}{lof}
\listoffigures*
\cleardoublepage
% ---


% ---
% LISTA DE QUADROS
% ---
% \pdfbookmark[0]{\listofquadrosname}{loq}
% \listofquadros*
% \cleardoublepage
% ---


% ---
% LISTA DE TABELAS
% ---
% \pdfbookmark[0]{\listtablename}{lot}
% \listoftables*
% \cleardoublepage
% ---


% ---
% LISTA DE ABREVIATURAS E SIGLAS
% ---
% \begin{siglas}
%   \item[ABNT] Associação Brasileira de Normas Técnicas
%   \item[abnTeX] ABsurdas Normas para TeX
% \end{siglas}
% ---


% ---
% LISTA DE SÍMBOLOS
% ---
% \begin{simbolos}
%   \item[$ \Gamma $] Letra grega Gama
%   \item[$ \Lambda $] Lambda
%   \item[$ \zeta $] Letra grega minúscula zeta
%   \item[$ \in $] Pertence
% \end{simbolos}
% ---


% ---
% Sumário
% ---
\pdfbookmark[0]{\contentsname}{toc}
\tableofcontents*
\cleardoublepage
% ---


% ----------------------------------------------------------
% ELEMENTOS TEXTUAIS
% ----------------------------------------------------------
\textual

\include{./01-introducao/main.tex}

\include{./02-machine-learning/main.tex}

\include{./03-exemplo/main.tex}

\include{./04-conclusao/main.tex}

% ----------------------------------------------------------
% ELEMENTOS PÓS-TEXTUAIS
% ----------------------------------------------------------
\postextual
% ----------------------------------------------------------


% ----------------------------------------------------------
% REFERÊNCIAS
% ----------------------------------------------------------
\bibliography{references.bib}


% ----------------------------------------------------------
% GLOSSÁRIO
% ----------------------------------------------------------
%
% Consulte o manual da classe abntex2 para orientações sobre o glossário.
%
%\glossary

% ----------------------------------------------------------
% APÊNDICES
% ----------------------------------------------------------

% ---
% Inicia os apêndices
% ---
% \begin{apendicesenv}

% % Imprime uma página indicando o início dos apêndices
% \partapendices

% % ----------------------------------------------------------
% \chapter{Quisque libero justo}
% % ----------------------------------------------------------

% % ----------------------------------------------------------
% \chapter{Nullam elementum urna}
% % ----------------------------------------------------------

% \end{apendicesenv}
% ---


% ----------------------------------------------------------
% ANEXOS
% ----------------------------------------------------------

% ---
% Inicia os anexos
% ---
% \begin{anexosenv}

% % Imprime uma página indicando o início dos anexos
% \partanexos

% % ---
% \chapter{Morbi ultrices rutrum lorem}
% % ---

% % ---
% \chapter{Cras non urna sed feugiat cum}
% % ---


% % ---
% \chapter{Fusce facilisis lacinia dui}
% % ---

% \end{anexosenv}

%---------------------------------------------------------------------
% INDICE REMISSIVO
%---------------------------------------------------------------------
% \phantompart
% \printindex
%---------------------------------------------------------------------
\includepdf{./05-extras/contracapa.pdf}
\end{document}


\documentclass[
	% Opções da classe memoir
	12pt,				        % tamanho da fonte
	openany,			        % capítulos começam em pág ímpar (insere página vazia caso preciso)
	oneside,			        % para impressão em recto e verso. Oposto a oneside
	a4paper,			        % tamanho do papel. 
	% Opções da classe abntex2
	%chapter=TITLE,		        % títulos de capítulos convertidos em letras maiúsculas
	%section=TITLE,		        % títulos de seções convertidos em letras maiúsculas
	%subsection=TITLE,	        % títulos de subseções convertidos em letras maiúsculas
	%subsubsection=TITLE,       % títulos de subsubseções convertidos em letras maiúsculas
	% Opções do pacote babel
	english,			        % idioma adicional para hifenização
	brazil				        % o último idioma é o principal do documento
	]{abntex2}

% ---
% Pacotes básicos 
% ---
\usepackage{lmodern}			% Usa a fonte Latin Modern			
\usepackage[T1]{fontenc}		% Selecao de codigos de fonte.
\usepackage[utf8]{inputenc}		% Codificacao do documento (conversão automática dos acentos)
\usepackage{indentfirst}		% Indenta o primeiro parágrafo de cada seção.
\usepackage{color}				% Controle das cores
\usepackage{graphicx}			% Inclusão de gráficos
\usepackage{microtype} 			% para melhorias de justificação
% ---
		

% ---
% Pacotes de citações
% ---
% \usepackage[alf]{abntex2cite}						% Citações padrão ABNT
\usepackage[
  alf,
  abnt-nbr10520=1988,                     % Citações em sentence case
  abnt-emphasize=bf
  ]{abntex2cite}


% ---
% Pacotes adicionais
% ---
\usepackage{pdfpages}	% Inclusão de páginas em PDF


% ---
% Informações para CAPA e FOLHA DE ROSTO
% ---
\titulo{APRENDIZAGEM DE MÁQUINA PARA PREDIÇÃO DE DESFECHOS EM SAÚDE MENTAL: \textmd{UMA INTRODUÇÃO PARA PSICÓLOGOS}}
\autor{GUILHERME AUGUSTO DIAS}
\local{Porto Alegre}
\data{2024}
\orientador{Prof. Dr. Wagner de Lara Machado}
% \coorientador{}
% \instituicao{
%   Pontifícia Universidade Católica do Rio Grande do Sul
%   \par
%   Escola de Ciências da Saúde e da Vida
%   \par
%   Curso de Psicologia}
\tipotrabalho{TCC}
\preambulo{Trabalho de Conclusão de Curso apresentado como requisito para a obtenção do grau de Bacharel em Psicologia pela Escola de Ciências da Saúde e da Vida da Pontifícia Universidade Católica do Rio Grande do Sul.}
% ---

% ---
% Configurações de aparência do PDF final
% ---

% Alterando o aspecto da cor azul
\definecolor{blue}{RGB}{41,5,195}

% Informações do PDF
\makeatletter
\hypersetup{
		pdftitle={Aprendizagem de máquina para predição de desfechos em saúde mental: uma introdução para psicólogos}, 
		pdfauthor={Guilherme Augusto Dias},
    	pdfsubject={\imprimirpreambulo},
	    pdfcreator={LaTeX with abnTeX2},
		pdfkeywords={aprendizagem de máquina}{psicoterapia}{psicologia positiva}{depressão}, 
		colorlinks=true,		% true: colored links; false: boxed links
    	linkcolor=black,		% color of internal links
    	citecolor=black,		% color of links to bibliography
    	filecolor=black,		% color of file links
		urlcolor=black,
		bookmarksdepth=4
}
\makeatother
% --- 


% ---
% Posiciona figuras e tabelas no topo da página quando adicionadas sozinhas
% em um página em branco. Ver https://github.com/abntex/abntex2/issues/170
\makeatletter
\setlength{\@fptop}{5pt} % Set distance from top of page to first float
\makeatother
% ---

% ---
% Possibilita criação de Quadros e Lista de quadros.
% Ver https://github.com/abntex/abntex2/issues/176
%
\newcommand{\quadroname}{Quadro}
\newcommand{\listofquadrosname}{Lista de quadros}

\newfloat[chapter]{quadro}{loq}{\quadroname}
\newlistof{listofquadros}{loq}{\listofquadrosname}
\newlistentry{quadro}{loq}{0}

% Configurações para atender às regras da ABNT
\setfloatadjustment{quadro}{\centering}
\counterwithout{quadro}{chapter}
\renewcommand{\cftquadroname}{\quadroname\space} 
\renewcommand*{\cftquadroaftersnum}{\hfill--\hfill}

\setfloatlocations{quadro}{hbtp} % Ver https://github.com/abntex/abntex2/issues/176
% ---


% --- 
% Espaçamentos entre linhas e parágrafos 
% --- 

% O tamanho do parágrafo é dado por:
\setlength{\parindent}{1.3cm}

% Controle do espaçamento entre um parágrafo e outro:
\setlength{\parskip}{0.2cm}  % tente também \onelineskip

% ---
% Compila o índice
% ---
\makeindex
% ---


% ----
% Início do documento
% ----
\begin{document}

% Seleciona o idioma do documento (conforme pacotes do babel)
\selectlanguage{brazil}

% Retira espaço extra obsoleto entre as frases.
\frenchspacing 

% ----------------------------------------------------------
% ELEMENTOS PRÉ-TEXTUAIS
% ----------------------------------------------------------
\pretextual

% ---
% CAPA
% ---
\includepdf{./05-extras/capa.pdf}
% ---


% ---
% FOLHA DE ROSTO
% (o * indica que haverá a ficha bibliográfica)
% ---
\imprimirfolhaderosto
% ---


% ---
% Ficha bibliografica
% ---

% Isto é um exemplo de Ficha Catalográfica, ou ``Dados internacionais de
% catalogação-na-publicação''. Você pode utilizar este modelo como referência. 
% Porém, provavelmente a biblioteca da sua universidade lhe fornecerá um PDF
% com a ficha catalográfica definitiva após a defesa do trabalho. Quando estiver
% com o documento, salve-o como PDF no diretório do seu projeto e substitua todo
% o conteúdo de implementação deste arquivo pelo comando abaixo:
%
% \usepackage{pdfpages}		% necessário para comando \includepdf
% \begin{fichacatalografica}
%     \includepdf{fig_ficha_catalografica.pdf}
% \end{fichacatalografica}

% \begin{fichacatalografica}
% 	\sffamily
% 	\vspace*{\fill}                             % Posição vertical
% 	\begin{center}                              % Minipage Centralizado
% 	\fbox{\begin{minipage}[c][8cm]{13.5cm}		% Largura
% 	\small
% 	\imprimirautor
% 	%Sobrenome, Nome do autor
	
% 	\hspace{0.5cm} \imprimirtitulo  / \imprimirautor. --
% 	\imprimirlocal, \imprimirdata-
	
% 	\hspace{0.5cm} \thelastpage p. : il. (algumas color.) ; 30 cm.\\
	
% 	\hspace{0.5cm} \imprimirorientadorRotulo~\imprimirorientador\\
	
% 	\hspace{0.5cm}
% 	\parbox[t]{\textwidth}{\imprimirtipotrabalho~--~\imprimirinstituicao,
% 	\imprimirdata.}\\
	
% 	\hspace{0.5cm}
% 		1. Palavra-chave1.
% 		2. Palavra-chave2.
% 		2. Palavra-chave3.
% 		I. Orientador.
% 		II. Universidade xxx.
% 		III. Faculdade de xxx.
% 		IV. Título 			
% 	\end{minipage}}
% 	\end{center}
% \end{fichacatalografica}
% ---


% ---
% FOLHA DE APROVAÇÃO
% ---

% Isto é um exemplo de Folha de aprovação, elemento obrigatório da NBR
% 14724/2011 (seção 4.2.1.3). Você pode utilizar este modelo até a aprovação
% do trabalho. Após isso, substitua todo o conteúdo deste arquivo por uma
% imagem da página assinada pela banca com o comando abaixo:
%
% \begin{folhadeaprovacao}
% \includepdf{folhadeaprovacao_final.pdf}
% \end{folhadeaprovacao}
%
\setlength{\ABNTEXsignskip}{2cm}
\setlength{\ABNTEXsignwidth}{10cm}

\begin{folhadeaprovacao}
  \begin{center}
    {\ABNTEXchapterfont\large\imprimirautor}

    \vspace*{1cm}

    \begin{center}
      \ABNTEXchapterfont\bfseries\Large\imprimirtitulo
    \end{center}

    \vspace*{1cm}
    
    \hspace{.45\textwidth}
    \begin{minipage}{.5\textwidth}
        \imprimirpreambulo
    \end{minipage}%
   \end{center}

   \vspace*{1cm}
        
   \begin{center}
   Aprovado em 17 de julho de 2024.

   \vspace*{1cm}

   Banca examinadora:
   \end{center}

   \vspace*{-1cm}

   \assinatura{Prof. Ma. Caroline Santa Maria Rodrigues -- PUCRS}
   \assinatura{Prof. Dra. Clarissa Pinto Pizarro de Freitas -- PUCRS}
   \assinatura{\imprimirorientador -- PUCRS}
\end{folhadeaprovacao}
% ---


% ---
% Agradecimentos
% ---
\begin{agradecimentos}
Agradeço aos meus pais, Ademilson José Dias e Irinês Porfírio Dias, pelo apoio incondicional que me deram ao longo da vida,
incentivando mesmo as ideias questionáveis como aquela de fazer uma segunda graduação.

A minha irmã, Maria Eduarda Dias, cuja trajetória acadêmica é fonte de inspiração.

A minha namorada, Nathalia Fallavena Ceratti, que acompanhou a esta jornada de perto, oferecendo apoio e compreensão que
fizeram toda a diferença.

Aos meus amigos, Guilherme da Silva Petró, Norton Neves e Uilinton Luciano Marques da Silva, que foram compreensíveis nos
finais de semana em que não pude jogar RPG.

Ao professor Wagner de Lara Machado pela orientação e apoio no desenvolvimento deste trabalho.
\end{agradecimentos}
% ---


% ---
% RESUMOS
% ---
% Resumo em português
\setlength{\absparsep}{18pt} % ajusta o espaçamento dos parágrafos do resumo
\begin{resumo}
Desenvolvimentos na área de aprendizagem de máquina representam uma série de oportunidades para diversas áreas de pesquisa e de atuação profissional.
O uso de modelos de aprendizagem de máquina na predição de desfecho para tratamentos em saúde mental tem sido investigado e apresenta resultados promissores.
A aplicação desse tipo de modelo possibilitaria, por exemplo, maior assertividade na personalização de tratamentos psicoterápicos a nível individual.
Apesar dos benefícios potenciais, o tema ainda é pouco discutido nos espaços de pesquisa, formação e atuação em psicologia clínica. Este trabalho tem
por objetivo introduzir o tema de aprendizagem de máquina a psicólogos. Busca-se apresentar, de maneira acessível, os conceitos e terminologia básicos
em aprendizagem de máquina, proporcionando um panorama geral da área. Para complementar a exposição conceitual, é construído um exemplo a partir de dados
de um estudo sobre eficácia de intervenções em psicologia positiva para depressão via internet. Por fim, discutem-se algumas implicações do uso desse tipo
de tecnologia para a pesquisa e para a prática em psicoterapia.

 \textbf{Palavras-chave}: aprendizagem de máquina, psicoterapia, psicologia positiva, depressão.
\end{resumo}

% Resumo em inglês
\begin{resumo}[Abstract]
 \begin{otherlanguage*}{english}
   Advances in machine learning have opened up a multitude of opportunities for various research and professional fields. The use of machine learning models to
   predict treatment outcomes in mental health has been investigated showing promising results. The application of such models would allow greater assertiveness
   in tailoring psychotherapeutic treatments at the individual level. Despite its potential benefits, the topic remains unexplored in the fields of research, training
   and practice in clinical psychology. This work aims to introduce the topic of machine learning to psychologists. It presents the basic machine learning concepts and
   terminology in an accessible manner, providing a comprehensive overview of the field. To complement the conceptual discussion, an example is constructed using data
   from a study on the effectiveness of internet-based positive psychology interventions for depression. Finally, implications of using this type of technology in
   psychotherapy research and practice are discussed.
   \vspace{\onelineskip}
 
   \noindent 
   \textbf{Keywords}: machine learning, psychotherapy, positive psychology, depression.
 \end{otherlanguage*}
\end{resumo}


% ---
% LISTA DE ILUSTRAÇÕES
% ---
\pdfbookmark[0]{\listfigurename}{lof}
\listoffigures*
\cleardoublepage
% ---


% ---
% LISTA DE QUADROS
% ---
% \pdfbookmark[0]{\listofquadrosname}{loq}
% \listofquadros*
% \cleardoublepage
% ---


% ---
% LISTA DE TABELAS
% ---
% \pdfbookmark[0]{\listtablename}{lot}
% \listoftables*
% \cleardoublepage
% ---


% ---
% LISTA DE ABREVIATURAS E SIGLAS
% ---
% \begin{siglas}
%   \item[ABNT] Associação Brasileira de Normas Técnicas
%   \item[abnTeX] ABsurdas Normas para TeX
% \end{siglas}
% ---


% ---
% LISTA DE SÍMBOLOS
% ---
% \begin{simbolos}
%   \item[$ \Gamma $] Letra grega Gama
%   \item[$ \Lambda $] Lambda
%   \item[$ \zeta $] Letra grega minúscula zeta
%   \item[$ \in $] Pertence
% \end{simbolos}
% ---


% ---
% Sumário
% ---
\pdfbookmark[0]{\contentsname}{toc}
\tableofcontents*
\cleardoublepage
% ---


% ----------------------------------------------------------
% ELEMENTOS TEXTUAIS
% ----------------------------------------------------------
\textual

\include{./01-introducao/main.tex}

\include{./02-machine-learning/main.tex}

\include{./03-exemplo/main.tex}

\include{./04-conclusao/main.tex}

% ----------------------------------------------------------
% ELEMENTOS PÓS-TEXTUAIS
% ----------------------------------------------------------
\postextual
% ----------------------------------------------------------


% ----------------------------------------------------------
% REFERÊNCIAS
% ----------------------------------------------------------
\bibliography{references.bib}


% ----------------------------------------------------------
% GLOSSÁRIO
% ----------------------------------------------------------
%
% Consulte o manual da classe abntex2 para orientações sobre o glossário.
%
%\glossary

% ----------------------------------------------------------
% APÊNDICES
% ----------------------------------------------------------

% ---
% Inicia os apêndices
% ---
% \begin{apendicesenv}

% % Imprime uma página indicando o início dos apêndices
% \partapendices

% % ----------------------------------------------------------
% \chapter{Quisque libero justo}
% % ----------------------------------------------------------

% % ----------------------------------------------------------
% \chapter{Nullam elementum urna}
% % ----------------------------------------------------------

% \end{apendicesenv}
% ---


% ----------------------------------------------------------
% ANEXOS
% ----------------------------------------------------------

% ---
% Inicia os anexos
% ---
% \begin{anexosenv}

% % Imprime uma página indicando o início dos anexos
% \partanexos

% % ---
% \chapter{Morbi ultrices rutrum lorem}
% % ---

% % ---
% \chapter{Cras non urna sed feugiat cum}
% % ---


% % ---
% \chapter{Fusce facilisis lacinia dui}
% % ---

% \end{anexosenv}

%---------------------------------------------------------------------
% INDICE REMISSIVO
%---------------------------------------------------------------------
% \phantompart
% \printindex
%---------------------------------------------------------------------
\includepdf{./05-extras/contracapa.pdf}
\end{document}


% ----------------------------------------------------------
% ELEMENTOS PÓS-TEXTUAIS
% ----------------------------------------------------------
\postextual
% ----------------------------------------------------------


% ----------------------------------------------------------
% REFERÊNCIAS
% ----------------------------------------------------------
\bibliography{references.bib}


% ----------------------------------------------------------
% GLOSSÁRIO
% ----------------------------------------------------------
%
% Consulte o manual da classe abntex2 para orientações sobre o glossário.
%
%\glossary

% ----------------------------------------------------------
% APÊNDICES
% ----------------------------------------------------------

% ---
% Inicia os apêndices
% ---
% \begin{apendicesenv}

% % Imprime uma página indicando o início dos apêndices
% \partapendices

% % ----------------------------------------------------------
% \chapter{Quisque libero justo}
% % ----------------------------------------------------------

% % ----------------------------------------------------------
% \chapter{Nullam elementum urna}
% % ----------------------------------------------------------

% \end{apendicesenv}
% ---


% ----------------------------------------------------------
% ANEXOS
% ----------------------------------------------------------

% ---
% Inicia os anexos
% ---
% \begin{anexosenv}

% % Imprime uma página indicando o início dos anexos
% \partanexos

% % ---
% \chapter{Morbi ultrices rutrum lorem}
% % ---

% % ---
% \chapter{Cras non urna sed feugiat cum}
% % ---


% % ---
% \chapter{Fusce facilisis lacinia dui}
% % ---

% \end{anexosenv}

%---------------------------------------------------------------------
% INDICE REMISSIVO
%---------------------------------------------------------------------
% \phantompart
% \printindex
%---------------------------------------------------------------------
\includepdf{./05-extras/contracapa.pdf}
\end{document}


\documentclass[
	% Opções da classe memoir
	12pt,				        % tamanho da fonte
	openany,			        % capítulos começam em pág ímpar (insere página vazia caso preciso)
	oneside,			        % para impressão em recto e verso. Oposto a oneside
	a4paper,			        % tamanho do papel. 
	% Opções da classe abntex2
	%chapter=TITLE,		        % títulos de capítulos convertidos em letras maiúsculas
	%section=TITLE,		        % títulos de seções convertidos em letras maiúsculas
	%subsection=TITLE,	        % títulos de subseções convertidos em letras maiúsculas
	%subsubsection=TITLE,       % títulos de subsubseções convertidos em letras maiúsculas
	% Opções do pacote babel
	english,			        % idioma adicional para hifenização
	brazil				        % o último idioma é o principal do documento
	]{abntex2}

% ---
% Pacotes básicos 
% ---
\usepackage{lmodern}			% Usa a fonte Latin Modern			
\usepackage[T1]{fontenc}		% Selecao de codigos de fonte.
\usepackage[utf8]{inputenc}		% Codificacao do documento (conversão automática dos acentos)
\usepackage{indentfirst}		% Indenta o primeiro parágrafo de cada seção.
\usepackage{color}				% Controle das cores
\usepackage{graphicx}			% Inclusão de gráficos
\usepackage{microtype} 			% para melhorias de justificação
% ---
		

% ---
% Pacotes de citações
% ---
% \usepackage[alf]{abntex2cite}						% Citações padrão ABNT
\usepackage[
  alf,
  abnt-nbr10520=1988,                     % Citações em sentence case
  abnt-emphasize=bf
  ]{abntex2cite}


% ---
% Pacotes adicionais
% ---
\usepackage{pdfpages}	% Inclusão de páginas em PDF


% ---
% Informações para CAPA e FOLHA DE ROSTO
% ---
\titulo{APRENDIZAGEM DE MÁQUINA PARA PREDIÇÃO DE DESFECHOS EM SAÚDE MENTAL: \textmd{UMA INTRODUÇÃO PARA PSICÓLOGOS}}
\autor{GUILHERME AUGUSTO DIAS}
\local{Porto Alegre}
\data{2024}
\orientador{Prof. Dr. Wagner de Lara Machado}
% \coorientador{}
% \instituicao{
%   Pontifícia Universidade Católica do Rio Grande do Sul
%   \par
%   Escola de Ciências da Saúde e da Vida
%   \par
%   Curso de Psicologia}
\tipotrabalho{TCC}
\preambulo{Trabalho de Conclusão de Curso apresentado como requisito para a obtenção do grau de Bacharel em Psicologia pela Escola de Ciências da Saúde e da Vida da Pontifícia Universidade Católica do Rio Grande do Sul.}
% ---

% ---
% Configurações de aparência do PDF final
% ---

% Alterando o aspecto da cor azul
\definecolor{blue}{RGB}{41,5,195}

% Informações do PDF
\makeatletter
\hypersetup{
		pdftitle={Aprendizagem de máquina para predição de desfechos em saúde mental: uma introdução para psicólogos}, 
		pdfauthor={Guilherme Augusto Dias},
    	pdfsubject={\imprimirpreambulo},
	    pdfcreator={LaTeX with abnTeX2},
		pdfkeywords={aprendizagem de máquina}{psicoterapia}{psicologia positiva}{depressão}, 
		colorlinks=true,		% true: colored links; false: boxed links
    	linkcolor=black,		% color of internal links
    	citecolor=black,		% color of links to bibliography
    	filecolor=black,		% color of file links
		urlcolor=black,
		bookmarksdepth=4
}
\makeatother
% --- 


% ---
% Posiciona figuras e tabelas no topo da página quando adicionadas sozinhas
% em um página em branco. Ver https://github.com/abntex/abntex2/issues/170
\makeatletter
\setlength{\@fptop}{5pt} % Set distance from top of page to first float
\makeatother
% ---

% ---
% Possibilita criação de Quadros e Lista de quadros.
% Ver https://github.com/abntex/abntex2/issues/176
%
\newcommand{\quadroname}{Quadro}
\newcommand{\listofquadrosname}{Lista de quadros}

\newfloat[chapter]{quadro}{loq}{\quadroname}
\newlistof{listofquadros}{loq}{\listofquadrosname}
\newlistentry{quadro}{loq}{0}

% Configurações para atender às regras da ABNT
\setfloatadjustment{quadro}{\centering}
\counterwithout{quadro}{chapter}
\renewcommand{\cftquadroname}{\quadroname\space} 
\renewcommand*{\cftquadroaftersnum}{\hfill--\hfill}

\setfloatlocations{quadro}{hbtp} % Ver https://github.com/abntex/abntex2/issues/176
% ---


% --- 
% Espaçamentos entre linhas e parágrafos 
% --- 

% O tamanho do parágrafo é dado por:
\setlength{\parindent}{1.3cm}

% Controle do espaçamento entre um parágrafo e outro:
\setlength{\parskip}{0.2cm}  % tente também \onelineskip

% ---
% Compila o índice
% ---
\makeindex
% ---


% ----
% Início do documento
% ----
\begin{document}

% Seleciona o idioma do documento (conforme pacotes do babel)
\selectlanguage{brazil}

% Retira espaço extra obsoleto entre as frases.
\frenchspacing 

% ----------------------------------------------------------
% ELEMENTOS PRÉ-TEXTUAIS
% ----------------------------------------------------------
\pretextual

% ---
% CAPA
% ---
\includepdf{./05-extras/capa.pdf}
% ---


% ---
% FOLHA DE ROSTO
% (o * indica que haverá a ficha bibliográfica)
% ---
\imprimirfolhaderosto
% ---


% ---
% Ficha bibliografica
% ---

% Isto é um exemplo de Ficha Catalográfica, ou ``Dados internacionais de
% catalogação-na-publicação''. Você pode utilizar este modelo como referência. 
% Porém, provavelmente a biblioteca da sua universidade lhe fornecerá um PDF
% com a ficha catalográfica definitiva após a defesa do trabalho. Quando estiver
% com o documento, salve-o como PDF no diretório do seu projeto e substitua todo
% o conteúdo de implementação deste arquivo pelo comando abaixo:
%
% \usepackage{pdfpages}		% necessário para comando \includepdf
% \begin{fichacatalografica}
%     \includepdf{fig_ficha_catalografica.pdf}
% \end{fichacatalografica}

% \begin{fichacatalografica}
% 	\sffamily
% 	\vspace*{\fill}                             % Posição vertical
% 	\begin{center}                              % Minipage Centralizado
% 	\fbox{\begin{minipage}[c][8cm]{13.5cm}		% Largura
% 	\small
% 	\imprimirautor
% 	%Sobrenome, Nome do autor
	
% 	\hspace{0.5cm} \imprimirtitulo  / \imprimirautor. --
% 	\imprimirlocal, \imprimirdata-
	
% 	\hspace{0.5cm} \thelastpage p. : il. (algumas color.) ; 30 cm.\\
	
% 	\hspace{0.5cm} \imprimirorientadorRotulo~\imprimirorientador\\
	
% 	\hspace{0.5cm}
% 	\parbox[t]{\textwidth}{\imprimirtipotrabalho~--~\imprimirinstituicao,
% 	\imprimirdata.}\\
	
% 	\hspace{0.5cm}
% 		1. Palavra-chave1.
% 		2. Palavra-chave2.
% 		2. Palavra-chave3.
% 		I. Orientador.
% 		II. Universidade xxx.
% 		III. Faculdade de xxx.
% 		IV. Título 			
% 	\end{minipage}}
% 	\end{center}
% \end{fichacatalografica}
% ---


% ---
% FOLHA DE APROVAÇÃO
% ---

% Isto é um exemplo de Folha de aprovação, elemento obrigatório da NBR
% 14724/2011 (seção 4.2.1.3). Você pode utilizar este modelo até a aprovação
% do trabalho. Após isso, substitua todo o conteúdo deste arquivo por uma
% imagem da página assinada pela banca com o comando abaixo:
%
% \begin{folhadeaprovacao}
% \includepdf{folhadeaprovacao_final.pdf}
% \end{folhadeaprovacao}
%
\setlength{\ABNTEXsignskip}{2cm}
\setlength{\ABNTEXsignwidth}{10cm}

\begin{folhadeaprovacao}
  \begin{center}
    {\ABNTEXchapterfont\large\imprimirautor}

    \vspace*{1cm}

    \begin{center}
      \ABNTEXchapterfont\bfseries\Large\imprimirtitulo
    \end{center}

    \vspace*{1cm}
    
    \hspace{.45\textwidth}
    \begin{minipage}{.5\textwidth}
        \imprimirpreambulo
    \end{minipage}%
   \end{center}

   \vspace*{1cm}
        
   \begin{center}
   Aprovado em 17 de julho de 2024.

   \vspace*{1cm}

   Banca examinadora:
   \end{center}

   \vspace*{-1cm}

   \assinatura{Prof. Ma. Caroline Santa Maria Rodrigues -- PUCRS}
   \assinatura{Prof. Dra. Clarissa Pinto Pizarro de Freitas -- PUCRS}
   \assinatura{\imprimirorientador -- PUCRS}
\end{folhadeaprovacao}
% ---


% ---
% Agradecimentos
% ---
\begin{agradecimentos}
Agradeço aos meus pais, Ademilson José Dias e Irinês Porfírio Dias, pelo apoio incondicional que me deram ao longo da vida,
incentivando mesmo as ideias questionáveis como aquela de fazer uma segunda graduação.

A minha irmã, Maria Eduarda Dias, cuja trajetória acadêmica é fonte de inspiração.

A minha namorada, Nathalia Fallavena Ceratti, que acompanhou a esta jornada de perto, oferecendo apoio e compreensão que
fizeram toda a diferença.

Aos meus amigos, Guilherme da Silva Petró, Norton Neves e Uilinton Luciano Marques da Silva, que foram compreensíveis nos
finais de semana em que não pude jogar RPG.

Ao professor Wagner de Lara Machado pela orientação e apoio no desenvolvimento deste trabalho.
\end{agradecimentos}
% ---


% ---
% RESUMOS
% ---
% Resumo em português
\setlength{\absparsep}{18pt} % ajusta o espaçamento dos parágrafos do resumo
\begin{resumo}
Desenvolvimentos na área de aprendizagem de máquina representam uma série de oportunidades para diversas áreas de pesquisa e de atuação profissional.
O uso de modelos de aprendizagem de máquina na predição de desfecho para tratamentos em saúde mental tem sido investigado e apresenta resultados promissores.
A aplicação desse tipo de modelo possibilitaria, por exemplo, maior assertividade na personalização de tratamentos psicoterápicos a nível individual.
Apesar dos benefícios potenciais, o tema ainda é pouco discutido nos espaços de pesquisa, formação e atuação em psicologia clínica. Este trabalho tem
por objetivo introduzir o tema de aprendizagem de máquina a psicólogos. Busca-se apresentar, de maneira acessível, os conceitos e terminologia básicos
em aprendizagem de máquina, proporcionando um panorama geral da área. Para complementar a exposição conceitual, é construído um exemplo a partir de dados
de um estudo sobre eficácia de intervenções em psicologia positiva para depressão via internet. Por fim, discutem-se algumas implicações do uso desse tipo
de tecnologia para a pesquisa e para a prática em psicoterapia.

 \textbf{Palavras-chave}: aprendizagem de máquina, psicoterapia, psicologia positiva, depressão.
\end{resumo}

% Resumo em inglês
\begin{resumo}[Abstract]
 \begin{otherlanguage*}{english}
   Advances in machine learning have opened up a multitude of opportunities for various research and professional fields. The use of machine learning models to
   predict treatment outcomes in mental health has been investigated showing promising results. The application of such models would allow greater assertiveness
   in tailoring psychotherapeutic treatments at the individual level. Despite its potential benefits, the topic remains unexplored in the fields of research, training
   and practice in clinical psychology. This work aims to introduce the topic of machine learning to psychologists. It presents the basic machine learning concepts and
   terminology in an accessible manner, providing a comprehensive overview of the field. To complement the conceptual discussion, an example is constructed using data
   from a study on the effectiveness of internet-based positive psychology interventions for depression. Finally, implications of using this type of technology in
   psychotherapy research and practice are discussed.
   \vspace{\onelineskip}
 
   \noindent 
   \textbf{Keywords}: machine learning, psychotherapy, positive psychology, depression.
 \end{otherlanguage*}
\end{resumo}


% ---
% LISTA DE ILUSTRAÇÕES
% ---
\pdfbookmark[0]{\listfigurename}{lof}
\listoffigures*
\cleardoublepage
% ---


% ---
% LISTA DE QUADROS
% ---
% \pdfbookmark[0]{\listofquadrosname}{loq}
% \listofquadros*
% \cleardoublepage
% ---


% ---
% LISTA DE TABELAS
% ---
% \pdfbookmark[0]{\listtablename}{lot}
% \listoftables*
% \cleardoublepage
% ---


% ---
% LISTA DE ABREVIATURAS E SIGLAS
% ---
% \begin{siglas}
%   \item[ABNT] Associação Brasileira de Normas Técnicas
%   \item[abnTeX] ABsurdas Normas para TeX
% \end{siglas}
% ---


% ---
% LISTA DE SÍMBOLOS
% ---
% \begin{simbolos}
%   \item[$ \Gamma $] Letra grega Gama
%   \item[$ \Lambda $] Lambda
%   \item[$ \zeta $] Letra grega minúscula zeta
%   \item[$ \in $] Pertence
% \end{simbolos}
% ---


% ---
% Sumário
% ---
\pdfbookmark[0]{\contentsname}{toc}
\tableofcontents*
\cleardoublepage
% ---


% ----------------------------------------------------------
% ELEMENTOS TEXTUAIS
% ----------------------------------------------------------
\textual

\documentclass[
	% Opções da classe memoir
	12pt,				        % tamanho da fonte
	openany,			        % capítulos começam em pág ímpar (insere página vazia caso preciso)
	oneside,			        % para impressão em recto e verso. Oposto a oneside
	a4paper,			        % tamanho do papel. 
	% Opções da classe abntex2
	%chapter=TITLE,		        % títulos de capítulos convertidos em letras maiúsculas
	%section=TITLE,		        % títulos de seções convertidos em letras maiúsculas
	%subsection=TITLE,	        % títulos de subseções convertidos em letras maiúsculas
	%subsubsection=TITLE,       % títulos de subsubseções convertidos em letras maiúsculas
	% Opções do pacote babel
	english,			        % idioma adicional para hifenização
	brazil				        % o último idioma é o principal do documento
	]{abntex2}

% ---
% Pacotes básicos 
% ---
\usepackage{lmodern}			% Usa a fonte Latin Modern			
\usepackage[T1]{fontenc}		% Selecao de codigos de fonte.
\usepackage[utf8]{inputenc}		% Codificacao do documento (conversão automática dos acentos)
\usepackage{indentfirst}		% Indenta o primeiro parágrafo de cada seção.
\usepackage{color}				% Controle das cores
\usepackage{graphicx}			% Inclusão de gráficos
\usepackage{microtype} 			% para melhorias de justificação
% ---
		

% ---
% Pacotes de citações
% ---
% \usepackage[alf]{abntex2cite}						% Citações padrão ABNT
\usepackage[
  alf,
  abnt-nbr10520=1988,                     % Citações em sentence case
  abnt-emphasize=bf
  ]{abntex2cite}


% ---
% Pacotes adicionais
% ---
\usepackage{pdfpages}	% Inclusão de páginas em PDF


% ---
% Informações para CAPA e FOLHA DE ROSTO
% ---
\titulo{APRENDIZAGEM DE MÁQUINA PARA PREDIÇÃO DE DESFECHOS EM SAÚDE MENTAL: \textmd{UMA INTRODUÇÃO PARA PSICÓLOGOS}}
\autor{GUILHERME AUGUSTO DIAS}
\local{Porto Alegre}
\data{2024}
\orientador{Prof. Dr. Wagner de Lara Machado}
% \coorientador{}
% \instituicao{
%   Pontifícia Universidade Católica do Rio Grande do Sul
%   \par
%   Escola de Ciências da Saúde e da Vida
%   \par
%   Curso de Psicologia}
\tipotrabalho{TCC}
\preambulo{Trabalho de Conclusão de Curso apresentado como requisito para a obtenção do grau de Bacharel em Psicologia pela Escola de Ciências da Saúde e da Vida da Pontifícia Universidade Católica do Rio Grande do Sul.}
% ---

% ---
% Configurações de aparência do PDF final
% ---

% Alterando o aspecto da cor azul
\definecolor{blue}{RGB}{41,5,195}

% Informações do PDF
\makeatletter
\hypersetup{
		pdftitle={Aprendizagem de máquina para predição de desfechos em saúde mental: uma introdução para psicólogos}, 
		pdfauthor={Guilherme Augusto Dias},
    	pdfsubject={\imprimirpreambulo},
	    pdfcreator={LaTeX with abnTeX2},
		pdfkeywords={aprendizagem de máquina}{psicoterapia}{psicologia positiva}{depressão}, 
		colorlinks=true,		% true: colored links; false: boxed links
    	linkcolor=black,		% color of internal links
    	citecolor=black,		% color of links to bibliography
    	filecolor=black,		% color of file links
		urlcolor=black,
		bookmarksdepth=4
}
\makeatother
% --- 


% ---
% Posiciona figuras e tabelas no topo da página quando adicionadas sozinhas
% em um página em branco. Ver https://github.com/abntex/abntex2/issues/170
\makeatletter
\setlength{\@fptop}{5pt} % Set distance from top of page to first float
\makeatother
% ---

% ---
% Possibilita criação de Quadros e Lista de quadros.
% Ver https://github.com/abntex/abntex2/issues/176
%
\newcommand{\quadroname}{Quadro}
\newcommand{\listofquadrosname}{Lista de quadros}

\newfloat[chapter]{quadro}{loq}{\quadroname}
\newlistof{listofquadros}{loq}{\listofquadrosname}
\newlistentry{quadro}{loq}{0}

% Configurações para atender às regras da ABNT
\setfloatadjustment{quadro}{\centering}
\counterwithout{quadro}{chapter}
\renewcommand{\cftquadroname}{\quadroname\space} 
\renewcommand*{\cftquadroaftersnum}{\hfill--\hfill}

\setfloatlocations{quadro}{hbtp} % Ver https://github.com/abntex/abntex2/issues/176
% ---


% --- 
% Espaçamentos entre linhas e parágrafos 
% --- 

% O tamanho do parágrafo é dado por:
\setlength{\parindent}{1.3cm}

% Controle do espaçamento entre um parágrafo e outro:
\setlength{\parskip}{0.2cm}  % tente também \onelineskip

% ---
% Compila o índice
% ---
\makeindex
% ---


% ----
% Início do documento
% ----
\begin{document}

% Seleciona o idioma do documento (conforme pacotes do babel)
\selectlanguage{brazil}

% Retira espaço extra obsoleto entre as frases.
\frenchspacing 

% ----------------------------------------------------------
% ELEMENTOS PRÉ-TEXTUAIS
% ----------------------------------------------------------
\pretextual

% ---
% CAPA
% ---
\includepdf{./05-extras/capa.pdf}
% ---


% ---
% FOLHA DE ROSTO
% (o * indica que haverá a ficha bibliográfica)
% ---
\imprimirfolhaderosto
% ---


% ---
% Ficha bibliografica
% ---

% Isto é um exemplo de Ficha Catalográfica, ou ``Dados internacionais de
% catalogação-na-publicação''. Você pode utilizar este modelo como referência. 
% Porém, provavelmente a biblioteca da sua universidade lhe fornecerá um PDF
% com a ficha catalográfica definitiva após a defesa do trabalho. Quando estiver
% com o documento, salve-o como PDF no diretório do seu projeto e substitua todo
% o conteúdo de implementação deste arquivo pelo comando abaixo:
%
% \usepackage{pdfpages}		% necessário para comando \includepdf
% \begin{fichacatalografica}
%     \includepdf{fig_ficha_catalografica.pdf}
% \end{fichacatalografica}

% \begin{fichacatalografica}
% 	\sffamily
% 	\vspace*{\fill}                             % Posição vertical
% 	\begin{center}                              % Minipage Centralizado
% 	\fbox{\begin{minipage}[c][8cm]{13.5cm}		% Largura
% 	\small
% 	\imprimirautor
% 	%Sobrenome, Nome do autor
	
% 	\hspace{0.5cm} \imprimirtitulo  / \imprimirautor. --
% 	\imprimirlocal, \imprimirdata-
	
% 	\hspace{0.5cm} \thelastpage p. : il. (algumas color.) ; 30 cm.\\
	
% 	\hspace{0.5cm} \imprimirorientadorRotulo~\imprimirorientador\\
	
% 	\hspace{0.5cm}
% 	\parbox[t]{\textwidth}{\imprimirtipotrabalho~--~\imprimirinstituicao,
% 	\imprimirdata.}\\
	
% 	\hspace{0.5cm}
% 		1. Palavra-chave1.
% 		2. Palavra-chave2.
% 		2. Palavra-chave3.
% 		I. Orientador.
% 		II. Universidade xxx.
% 		III. Faculdade de xxx.
% 		IV. Título 			
% 	\end{minipage}}
% 	\end{center}
% \end{fichacatalografica}
% ---


% ---
% FOLHA DE APROVAÇÃO
% ---

% Isto é um exemplo de Folha de aprovação, elemento obrigatório da NBR
% 14724/2011 (seção 4.2.1.3). Você pode utilizar este modelo até a aprovação
% do trabalho. Após isso, substitua todo o conteúdo deste arquivo por uma
% imagem da página assinada pela banca com o comando abaixo:
%
% \begin{folhadeaprovacao}
% \includepdf{folhadeaprovacao_final.pdf}
% \end{folhadeaprovacao}
%
\setlength{\ABNTEXsignskip}{2cm}
\setlength{\ABNTEXsignwidth}{10cm}

\begin{folhadeaprovacao}
  \begin{center}
    {\ABNTEXchapterfont\large\imprimirautor}

    \vspace*{1cm}

    \begin{center}
      \ABNTEXchapterfont\bfseries\Large\imprimirtitulo
    \end{center}

    \vspace*{1cm}
    
    \hspace{.45\textwidth}
    \begin{minipage}{.5\textwidth}
        \imprimirpreambulo
    \end{minipage}%
   \end{center}

   \vspace*{1cm}
        
   \begin{center}
   Aprovado em 17 de julho de 2024.

   \vspace*{1cm}

   Banca examinadora:
   \end{center}

   \vspace*{-1cm}

   \assinatura{Prof. Ma. Caroline Santa Maria Rodrigues -- PUCRS}
   \assinatura{Prof. Dra. Clarissa Pinto Pizarro de Freitas -- PUCRS}
   \assinatura{\imprimirorientador -- PUCRS}
\end{folhadeaprovacao}
% ---


% ---
% Agradecimentos
% ---
\begin{agradecimentos}
Agradeço aos meus pais, Ademilson José Dias e Irinês Porfírio Dias, pelo apoio incondicional que me deram ao longo da vida,
incentivando mesmo as ideias questionáveis como aquela de fazer uma segunda graduação.

A minha irmã, Maria Eduarda Dias, cuja trajetória acadêmica é fonte de inspiração.

A minha namorada, Nathalia Fallavena Ceratti, que acompanhou a esta jornada de perto, oferecendo apoio e compreensão que
fizeram toda a diferença.

Aos meus amigos, Guilherme da Silva Petró, Norton Neves e Uilinton Luciano Marques da Silva, que foram compreensíveis nos
finais de semana em que não pude jogar RPG.

Ao professor Wagner de Lara Machado pela orientação e apoio no desenvolvimento deste trabalho.
\end{agradecimentos}
% ---


% ---
% RESUMOS
% ---
% Resumo em português
\setlength{\absparsep}{18pt} % ajusta o espaçamento dos parágrafos do resumo
\begin{resumo}
Desenvolvimentos na área de aprendizagem de máquina representam uma série de oportunidades para diversas áreas de pesquisa e de atuação profissional.
O uso de modelos de aprendizagem de máquina na predição de desfecho para tratamentos em saúde mental tem sido investigado e apresenta resultados promissores.
A aplicação desse tipo de modelo possibilitaria, por exemplo, maior assertividade na personalização de tratamentos psicoterápicos a nível individual.
Apesar dos benefícios potenciais, o tema ainda é pouco discutido nos espaços de pesquisa, formação e atuação em psicologia clínica. Este trabalho tem
por objetivo introduzir o tema de aprendizagem de máquina a psicólogos. Busca-se apresentar, de maneira acessível, os conceitos e terminologia básicos
em aprendizagem de máquina, proporcionando um panorama geral da área. Para complementar a exposição conceitual, é construído um exemplo a partir de dados
de um estudo sobre eficácia de intervenções em psicologia positiva para depressão via internet. Por fim, discutem-se algumas implicações do uso desse tipo
de tecnologia para a pesquisa e para a prática em psicoterapia.

 \textbf{Palavras-chave}: aprendizagem de máquina, psicoterapia, psicologia positiva, depressão.
\end{resumo}

% Resumo em inglês
\begin{resumo}[Abstract]
 \begin{otherlanguage*}{english}
   Advances in machine learning have opened up a multitude of opportunities for various research and professional fields. The use of machine learning models to
   predict treatment outcomes in mental health has been investigated showing promising results. The application of such models would allow greater assertiveness
   in tailoring psychotherapeutic treatments at the individual level. Despite its potential benefits, the topic remains unexplored in the fields of research, training
   and practice in clinical psychology. This work aims to introduce the topic of machine learning to psychologists. It presents the basic machine learning concepts and
   terminology in an accessible manner, providing a comprehensive overview of the field. To complement the conceptual discussion, an example is constructed using data
   from a study on the effectiveness of internet-based positive psychology interventions for depression. Finally, implications of using this type of technology in
   psychotherapy research and practice are discussed.
   \vspace{\onelineskip}
 
   \noindent 
   \textbf{Keywords}: machine learning, psychotherapy, positive psychology, depression.
 \end{otherlanguage*}
\end{resumo}


% ---
% LISTA DE ILUSTRAÇÕES
% ---
\pdfbookmark[0]{\listfigurename}{lof}
\listoffigures*
\cleardoublepage
% ---


% ---
% LISTA DE QUADROS
% ---
% \pdfbookmark[0]{\listofquadrosname}{loq}
% \listofquadros*
% \cleardoublepage
% ---


% ---
% LISTA DE TABELAS
% ---
% \pdfbookmark[0]{\listtablename}{lot}
% \listoftables*
% \cleardoublepage
% ---


% ---
% LISTA DE ABREVIATURAS E SIGLAS
% ---
% \begin{siglas}
%   \item[ABNT] Associação Brasileira de Normas Técnicas
%   \item[abnTeX] ABsurdas Normas para TeX
% \end{siglas}
% ---


% ---
% LISTA DE SÍMBOLOS
% ---
% \begin{simbolos}
%   \item[$ \Gamma $] Letra grega Gama
%   \item[$ \Lambda $] Lambda
%   \item[$ \zeta $] Letra grega minúscula zeta
%   \item[$ \in $] Pertence
% \end{simbolos}
% ---


% ---
% Sumário
% ---
\pdfbookmark[0]{\contentsname}{toc}
\tableofcontents*
\cleardoublepage
% ---


% ----------------------------------------------------------
% ELEMENTOS TEXTUAIS
% ----------------------------------------------------------
\textual

\include{./01-introducao/main.tex}

\include{./02-machine-learning/main.tex}

\include{./03-exemplo/main.tex}

\include{./04-conclusao/main.tex}

% ----------------------------------------------------------
% ELEMENTOS PÓS-TEXTUAIS
% ----------------------------------------------------------
\postextual
% ----------------------------------------------------------


% ----------------------------------------------------------
% REFERÊNCIAS
% ----------------------------------------------------------
\bibliography{references.bib}


% ----------------------------------------------------------
% GLOSSÁRIO
% ----------------------------------------------------------
%
% Consulte o manual da classe abntex2 para orientações sobre o glossário.
%
%\glossary

% ----------------------------------------------------------
% APÊNDICES
% ----------------------------------------------------------

% ---
% Inicia os apêndices
% ---
% \begin{apendicesenv}

% % Imprime uma página indicando o início dos apêndices
% \partapendices

% % ----------------------------------------------------------
% \chapter{Quisque libero justo}
% % ----------------------------------------------------------

% % ----------------------------------------------------------
% \chapter{Nullam elementum urna}
% % ----------------------------------------------------------

% \end{apendicesenv}
% ---


% ----------------------------------------------------------
% ANEXOS
% ----------------------------------------------------------

% ---
% Inicia os anexos
% ---
% \begin{anexosenv}

% % Imprime uma página indicando o início dos anexos
% \partanexos

% % ---
% \chapter{Morbi ultrices rutrum lorem}
% % ---

% % ---
% \chapter{Cras non urna sed feugiat cum}
% % ---


% % ---
% \chapter{Fusce facilisis lacinia dui}
% % ---

% \end{anexosenv}

%---------------------------------------------------------------------
% INDICE REMISSIVO
%---------------------------------------------------------------------
% \phantompart
% \printindex
%---------------------------------------------------------------------
\includepdf{./05-extras/contracapa.pdf}
\end{document}


\documentclass[
	% Opções da classe memoir
	12pt,				        % tamanho da fonte
	openany,			        % capítulos começam em pág ímpar (insere página vazia caso preciso)
	oneside,			        % para impressão em recto e verso. Oposto a oneside
	a4paper,			        % tamanho do papel. 
	% Opções da classe abntex2
	%chapter=TITLE,		        % títulos de capítulos convertidos em letras maiúsculas
	%section=TITLE,		        % títulos de seções convertidos em letras maiúsculas
	%subsection=TITLE,	        % títulos de subseções convertidos em letras maiúsculas
	%subsubsection=TITLE,       % títulos de subsubseções convertidos em letras maiúsculas
	% Opções do pacote babel
	english,			        % idioma adicional para hifenização
	brazil				        % o último idioma é o principal do documento
	]{abntex2}

% ---
% Pacotes básicos 
% ---
\usepackage{lmodern}			% Usa a fonte Latin Modern			
\usepackage[T1]{fontenc}		% Selecao de codigos de fonte.
\usepackage[utf8]{inputenc}		% Codificacao do documento (conversão automática dos acentos)
\usepackage{indentfirst}		% Indenta o primeiro parágrafo de cada seção.
\usepackage{color}				% Controle das cores
\usepackage{graphicx}			% Inclusão de gráficos
\usepackage{microtype} 			% para melhorias de justificação
% ---
		

% ---
% Pacotes de citações
% ---
% \usepackage[alf]{abntex2cite}						% Citações padrão ABNT
\usepackage[
  alf,
  abnt-nbr10520=1988,                     % Citações em sentence case
  abnt-emphasize=bf
  ]{abntex2cite}


% ---
% Pacotes adicionais
% ---
\usepackage{pdfpages}	% Inclusão de páginas em PDF


% ---
% Informações para CAPA e FOLHA DE ROSTO
% ---
\titulo{APRENDIZAGEM DE MÁQUINA PARA PREDIÇÃO DE DESFECHOS EM SAÚDE MENTAL: \textmd{UMA INTRODUÇÃO PARA PSICÓLOGOS}}
\autor{GUILHERME AUGUSTO DIAS}
\local{Porto Alegre}
\data{2024}
\orientador{Prof. Dr. Wagner de Lara Machado}
% \coorientador{}
% \instituicao{
%   Pontifícia Universidade Católica do Rio Grande do Sul
%   \par
%   Escola de Ciências da Saúde e da Vida
%   \par
%   Curso de Psicologia}
\tipotrabalho{TCC}
\preambulo{Trabalho de Conclusão de Curso apresentado como requisito para a obtenção do grau de Bacharel em Psicologia pela Escola de Ciências da Saúde e da Vida da Pontifícia Universidade Católica do Rio Grande do Sul.}
% ---

% ---
% Configurações de aparência do PDF final
% ---

% Alterando o aspecto da cor azul
\definecolor{blue}{RGB}{41,5,195}

% Informações do PDF
\makeatletter
\hypersetup{
		pdftitle={Aprendizagem de máquina para predição de desfechos em saúde mental: uma introdução para psicólogos}, 
		pdfauthor={Guilherme Augusto Dias},
    	pdfsubject={\imprimirpreambulo},
	    pdfcreator={LaTeX with abnTeX2},
		pdfkeywords={aprendizagem de máquina}{psicoterapia}{psicologia positiva}{depressão}, 
		colorlinks=true,		% true: colored links; false: boxed links
    	linkcolor=black,		% color of internal links
    	citecolor=black,		% color of links to bibliography
    	filecolor=black,		% color of file links
		urlcolor=black,
		bookmarksdepth=4
}
\makeatother
% --- 


% ---
% Posiciona figuras e tabelas no topo da página quando adicionadas sozinhas
% em um página em branco. Ver https://github.com/abntex/abntex2/issues/170
\makeatletter
\setlength{\@fptop}{5pt} % Set distance from top of page to first float
\makeatother
% ---

% ---
% Possibilita criação de Quadros e Lista de quadros.
% Ver https://github.com/abntex/abntex2/issues/176
%
\newcommand{\quadroname}{Quadro}
\newcommand{\listofquadrosname}{Lista de quadros}

\newfloat[chapter]{quadro}{loq}{\quadroname}
\newlistof{listofquadros}{loq}{\listofquadrosname}
\newlistentry{quadro}{loq}{0}

% Configurações para atender às regras da ABNT
\setfloatadjustment{quadro}{\centering}
\counterwithout{quadro}{chapter}
\renewcommand{\cftquadroname}{\quadroname\space} 
\renewcommand*{\cftquadroaftersnum}{\hfill--\hfill}

\setfloatlocations{quadro}{hbtp} % Ver https://github.com/abntex/abntex2/issues/176
% ---


% --- 
% Espaçamentos entre linhas e parágrafos 
% --- 

% O tamanho do parágrafo é dado por:
\setlength{\parindent}{1.3cm}

% Controle do espaçamento entre um parágrafo e outro:
\setlength{\parskip}{0.2cm}  % tente também \onelineskip

% ---
% Compila o índice
% ---
\makeindex
% ---


% ----
% Início do documento
% ----
\begin{document}

% Seleciona o idioma do documento (conforme pacotes do babel)
\selectlanguage{brazil}

% Retira espaço extra obsoleto entre as frases.
\frenchspacing 

% ----------------------------------------------------------
% ELEMENTOS PRÉ-TEXTUAIS
% ----------------------------------------------------------
\pretextual

% ---
% CAPA
% ---
\includepdf{./05-extras/capa.pdf}
% ---


% ---
% FOLHA DE ROSTO
% (o * indica que haverá a ficha bibliográfica)
% ---
\imprimirfolhaderosto
% ---


% ---
% Ficha bibliografica
% ---

% Isto é um exemplo de Ficha Catalográfica, ou ``Dados internacionais de
% catalogação-na-publicação''. Você pode utilizar este modelo como referência. 
% Porém, provavelmente a biblioteca da sua universidade lhe fornecerá um PDF
% com a ficha catalográfica definitiva após a defesa do trabalho. Quando estiver
% com o documento, salve-o como PDF no diretório do seu projeto e substitua todo
% o conteúdo de implementação deste arquivo pelo comando abaixo:
%
% \usepackage{pdfpages}		% necessário para comando \includepdf
% \begin{fichacatalografica}
%     \includepdf{fig_ficha_catalografica.pdf}
% \end{fichacatalografica}

% \begin{fichacatalografica}
% 	\sffamily
% 	\vspace*{\fill}                             % Posição vertical
% 	\begin{center}                              % Minipage Centralizado
% 	\fbox{\begin{minipage}[c][8cm]{13.5cm}		% Largura
% 	\small
% 	\imprimirautor
% 	%Sobrenome, Nome do autor
	
% 	\hspace{0.5cm} \imprimirtitulo  / \imprimirautor. --
% 	\imprimirlocal, \imprimirdata-
	
% 	\hspace{0.5cm} \thelastpage p. : il. (algumas color.) ; 30 cm.\\
	
% 	\hspace{0.5cm} \imprimirorientadorRotulo~\imprimirorientador\\
	
% 	\hspace{0.5cm}
% 	\parbox[t]{\textwidth}{\imprimirtipotrabalho~--~\imprimirinstituicao,
% 	\imprimirdata.}\\
	
% 	\hspace{0.5cm}
% 		1. Palavra-chave1.
% 		2. Palavra-chave2.
% 		2. Palavra-chave3.
% 		I. Orientador.
% 		II. Universidade xxx.
% 		III. Faculdade de xxx.
% 		IV. Título 			
% 	\end{minipage}}
% 	\end{center}
% \end{fichacatalografica}
% ---


% ---
% FOLHA DE APROVAÇÃO
% ---

% Isto é um exemplo de Folha de aprovação, elemento obrigatório da NBR
% 14724/2011 (seção 4.2.1.3). Você pode utilizar este modelo até a aprovação
% do trabalho. Após isso, substitua todo o conteúdo deste arquivo por uma
% imagem da página assinada pela banca com o comando abaixo:
%
% \begin{folhadeaprovacao}
% \includepdf{folhadeaprovacao_final.pdf}
% \end{folhadeaprovacao}
%
\setlength{\ABNTEXsignskip}{2cm}
\setlength{\ABNTEXsignwidth}{10cm}

\begin{folhadeaprovacao}
  \begin{center}
    {\ABNTEXchapterfont\large\imprimirautor}

    \vspace*{1cm}

    \begin{center}
      \ABNTEXchapterfont\bfseries\Large\imprimirtitulo
    \end{center}

    \vspace*{1cm}
    
    \hspace{.45\textwidth}
    \begin{minipage}{.5\textwidth}
        \imprimirpreambulo
    \end{minipage}%
   \end{center}

   \vspace*{1cm}
        
   \begin{center}
   Aprovado em 17 de julho de 2024.

   \vspace*{1cm}

   Banca examinadora:
   \end{center}

   \vspace*{-1cm}

   \assinatura{Prof. Ma. Caroline Santa Maria Rodrigues -- PUCRS}
   \assinatura{Prof. Dra. Clarissa Pinto Pizarro de Freitas -- PUCRS}
   \assinatura{\imprimirorientador -- PUCRS}
\end{folhadeaprovacao}
% ---


% ---
% Agradecimentos
% ---
\begin{agradecimentos}
Agradeço aos meus pais, Ademilson José Dias e Irinês Porfírio Dias, pelo apoio incondicional que me deram ao longo da vida,
incentivando mesmo as ideias questionáveis como aquela de fazer uma segunda graduação.

A minha irmã, Maria Eduarda Dias, cuja trajetória acadêmica é fonte de inspiração.

A minha namorada, Nathalia Fallavena Ceratti, que acompanhou a esta jornada de perto, oferecendo apoio e compreensão que
fizeram toda a diferença.

Aos meus amigos, Guilherme da Silva Petró, Norton Neves e Uilinton Luciano Marques da Silva, que foram compreensíveis nos
finais de semana em que não pude jogar RPG.

Ao professor Wagner de Lara Machado pela orientação e apoio no desenvolvimento deste trabalho.
\end{agradecimentos}
% ---


% ---
% RESUMOS
% ---
% Resumo em português
\setlength{\absparsep}{18pt} % ajusta o espaçamento dos parágrafos do resumo
\begin{resumo}
Desenvolvimentos na área de aprendizagem de máquina representam uma série de oportunidades para diversas áreas de pesquisa e de atuação profissional.
O uso de modelos de aprendizagem de máquina na predição de desfecho para tratamentos em saúde mental tem sido investigado e apresenta resultados promissores.
A aplicação desse tipo de modelo possibilitaria, por exemplo, maior assertividade na personalização de tratamentos psicoterápicos a nível individual.
Apesar dos benefícios potenciais, o tema ainda é pouco discutido nos espaços de pesquisa, formação e atuação em psicologia clínica. Este trabalho tem
por objetivo introduzir o tema de aprendizagem de máquina a psicólogos. Busca-se apresentar, de maneira acessível, os conceitos e terminologia básicos
em aprendizagem de máquina, proporcionando um panorama geral da área. Para complementar a exposição conceitual, é construído um exemplo a partir de dados
de um estudo sobre eficácia de intervenções em psicologia positiva para depressão via internet. Por fim, discutem-se algumas implicações do uso desse tipo
de tecnologia para a pesquisa e para a prática em psicoterapia.

 \textbf{Palavras-chave}: aprendizagem de máquina, psicoterapia, psicologia positiva, depressão.
\end{resumo}

% Resumo em inglês
\begin{resumo}[Abstract]
 \begin{otherlanguage*}{english}
   Advances in machine learning have opened up a multitude of opportunities for various research and professional fields. The use of machine learning models to
   predict treatment outcomes in mental health has been investigated showing promising results. The application of such models would allow greater assertiveness
   in tailoring psychotherapeutic treatments at the individual level. Despite its potential benefits, the topic remains unexplored in the fields of research, training
   and practice in clinical psychology. This work aims to introduce the topic of machine learning to psychologists. It presents the basic machine learning concepts and
   terminology in an accessible manner, providing a comprehensive overview of the field. To complement the conceptual discussion, an example is constructed using data
   from a study on the effectiveness of internet-based positive psychology interventions for depression. Finally, implications of using this type of technology in
   psychotherapy research and practice are discussed.
   \vspace{\onelineskip}
 
   \noindent 
   \textbf{Keywords}: machine learning, psychotherapy, positive psychology, depression.
 \end{otherlanguage*}
\end{resumo}


% ---
% LISTA DE ILUSTRAÇÕES
% ---
\pdfbookmark[0]{\listfigurename}{lof}
\listoffigures*
\cleardoublepage
% ---


% ---
% LISTA DE QUADROS
% ---
% \pdfbookmark[0]{\listofquadrosname}{loq}
% \listofquadros*
% \cleardoublepage
% ---


% ---
% LISTA DE TABELAS
% ---
% \pdfbookmark[0]{\listtablename}{lot}
% \listoftables*
% \cleardoublepage
% ---


% ---
% LISTA DE ABREVIATURAS E SIGLAS
% ---
% \begin{siglas}
%   \item[ABNT] Associação Brasileira de Normas Técnicas
%   \item[abnTeX] ABsurdas Normas para TeX
% \end{siglas}
% ---


% ---
% LISTA DE SÍMBOLOS
% ---
% \begin{simbolos}
%   \item[$ \Gamma $] Letra grega Gama
%   \item[$ \Lambda $] Lambda
%   \item[$ \zeta $] Letra grega minúscula zeta
%   \item[$ \in $] Pertence
% \end{simbolos}
% ---


% ---
% Sumário
% ---
\pdfbookmark[0]{\contentsname}{toc}
\tableofcontents*
\cleardoublepage
% ---


% ----------------------------------------------------------
% ELEMENTOS TEXTUAIS
% ----------------------------------------------------------
\textual

\include{./01-introducao/main.tex}

\include{./02-machine-learning/main.tex}

\include{./03-exemplo/main.tex}

\include{./04-conclusao/main.tex}

% ----------------------------------------------------------
% ELEMENTOS PÓS-TEXTUAIS
% ----------------------------------------------------------
\postextual
% ----------------------------------------------------------


% ----------------------------------------------------------
% REFERÊNCIAS
% ----------------------------------------------------------
\bibliography{references.bib}


% ----------------------------------------------------------
% GLOSSÁRIO
% ----------------------------------------------------------
%
% Consulte o manual da classe abntex2 para orientações sobre o glossário.
%
%\glossary

% ----------------------------------------------------------
% APÊNDICES
% ----------------------------------------------------------

% ---
% Inicia os apêndices
% ---
% \begin{apendicesenv}

% % Imprime uma página indicando o início dos apêndices
% \partapendices

% % ----------------------------------------------------------
% \chapter{Quisque libero justo}
% % ----------------------------------------------------------

% % ----------------------------------------------------------
% \chapter{Nullam elementum urna}
% % ----------------------------------------------------------

% \end{apendicesenv}
% ---


% ----------------------------------------------------------
% ANEXOS
% ----------------------------------------------------------

% ---
% Inicia os anexos
% ---
% \begin{anexosenv}

% % Imprime uma página indicando o início dos anexos
% \partanexos

% % ---
% \chapter{Morbi ultrices rutrum lorem}
% % ---

% % ---
% \chapter{Cras non urna sed feugiat cum}
% % ---


% % ---
% \chapter{Fusce facilisis lacinia dui}
% % ---

% \end{anexosenv}

%---------------------------------------------------------------------
% INDICE REMISSIVO
%---------------------------------------------------------------------
% \phantompart
% \printindex
%---------------------------------------------------------------------
\includepdf{./05-extras/contracapa.pdf}
\end{document}


\documentclass[
	% Opções da classe memoir
	12pt,				        % tamanho da fonte
	openany,			        % capítulos começam em pág ímpar (insere página vazia caso preciso)
	oneside,			        % para impressão em recto e verso. Oposto a oneside
	a4paper,			        % tamanho do papel. 
	% Opções da classe abntex2
	%chapter=TITLE,		        % títulos de capítulos convertidos em letras maiúsculas
	%section=TITLE,		        % títulos de seções convertidos em letras maiúsculas
	%subsection=TITLE,	        % títulos de subseções convertidos em letras maiúsculas
	%subsubsection=TITLE,       % títulos de subsubseções convertidos em letras maiúsculas
	% Opções do pacote babel
	english,			        % idioma adicional para hifenização
	brazil				        % o último idioma é o principal do documento
	]{abntex2}

% ---
% Pacotes básicos 
% ---
\usepackage{lmodern}			% Usa a fonte Latin Modern			
\usepackage[T1]{fontenc}		% Selecao de codigos de fonte.
\usepackage[utf8]{inputenc}		% Codificacao do documento (conversão automática dos acentos)
\usepackage{indentfirst}		% Indenta o primeiro parágrafo de cada seção.
\usepackage{color}				% Controle das cores
\usepackage{graphicx}			% Inclusão de gráficos
\usepackage{microtype} 			% para melhorias de justificação
% ---
		

% ---
% Pacotes de citações
% ---
% \usepackage[alf]{abntex2cite}						% Citações padrão ABNT
\usepackage[
  alf,
  abnt-nbr10520=1988,                     % Citações em sentence case
  abnt-emphasize=bf
  ]{abntex2cite}


% ---
% Pacotes adicionais
% ---
\usepackage{pdfpages}	% Inclusão de páginas em PDF


% ---
% Informações para CAPA e FOLHA DE ROSTO
% ---
\titulo{APRENDIZAGEM DE MÁQUINA PARA PREDIÇÃO DE DESFECHOS EM SAÚDE MENTAL: \textmd{UMA INTRODUÇÃO PARA PSICÓLOGOS}}
\autor{GUILHERME AUGUSTO DIAS}
\local{Porto Alegre}
\data{2024}
\orientador{Prof. Dr. Wagner de Lara Machado}
% \coorientador{}
% \instituicao{
%   Pontifícia Universidade Católica do Rio Grande do Sul
%   \par
%   Escola de Ciências da Saúde e da Vida
%   \par
%   Curso de Psicologia}
\tipotrabalho{TCC}
\preambulo{Trabalho de Conclusão de Curso apresentado como requisito para a obtenção do grau de Bacharel em Psicologia pela Escola de Ciências da Saúde e da Vida da Pontifícia Universidade Católica do Rio Grande do Sul.}
% ---

% ---
% Configurações de aparência do PDF final
% ---

% Alterando o aspecto da cor azul
\definecolor{blue}{RGB}{41,5,195}

% Informações do PDF
\makeatletter
\hypersetup{
		pdftitle={Aprendizagem de máquina para predição de desfechos em saúde mental: uma introdução para psicólogos}, 
		pdfauthor={Guilherme Augusto Dias},
    	pdfsubject={\imprimirpreambulo},
	    pdfcreator={LaTeX with abnTeX2},
		pdfkeywords={aprendizagem de máquina}{psicoterapia}{psicologia positiva}{depressão}, 
		colorlinks=true,		% true: colored links; false: boxed links
    	linkcolor=black,		% color of internal links
    	citecolor=black,		% color of links to bibliography
    	filecolor=black,		% color of file links
		urlcolor=black,
		bookmarksdepth=4
}
\makeatother
% --- 


% ---
% Posiciona figuras e tabelas no topo da página quando adicionadas sozinhas
% em um página em branco. Ver https://github.com/abntex/abntex2/issues/170
\makeatletter
\setlength{\@fptop}{5pt} % Set distance from top of page to first float
\makeatother
% ---

% ---
% Possibilita criação de Quadros e Lista de quadros.
% Ver https://github.com/abntex/abntex2/issues/176
%
\newcommand{\quadroname}{Quadro}
\newcommand{\listofquadrosname}{Lista de quadros}

\newfloat[chapter]{quadro}{loq}{\quadroname}
\newlistof{listofquadros}{loq}{\listofquadrosname}
\newlistentry{quadro}{loq}{0}

% Configurações para atender às regras da ABNT
\setfloatadjustment{quadro}{\centering}
\counterwithout{quadro}{chapter}
\renewcommand{\cftquadroname}{\quadroname\space} 
\renewcommand*{\cftquadroaftersnum}{\hfill--\hfill}

\setfloatlocations{quadro}{hbtp} % Ver https://github.com/abntex/abntex2/issues/176
% ---


% --- 
% Espaçamentos entre linhas e parágrafos 
% --- 

% O tamanho do parágrafo é dado por:
\setlength{\parindent}{1.3cm}

% Controle do espaçamento entre um parágrafo e outro:
\setlength{\parskip}{0.2cm}  % tente também \onelineskip

% ---
% Compila o índice
% ---
\makeindex
% ---


% ----
% Início do documento
% ----
\begin{document}

% Seleciona o idioma do documento (conforme pacotes do babel)
\selectlanguage{brazil}

% Retira espaço extra obsoleto entre as frases.
\frenchspacing 

% ----------------------------------------------------------
% ELEMENTOS PRÉ-TEXTUAIS
% ----------------------------------------------------------
\pretextual

% ---
% CAPA
% ---
\includepdf{./05-extras/capa.pdf}
% ---


% ---
% FOLHA DE ROSTO
% (o * indica que haverá a ficha bibliográfica)
% ---
\imprimirfolhaderosto
% ---


% ---
% Ficha bibliografica
% ---

% Isto é um exemplo de Ficha Catalográfica, ou ``Dados internacionais de
% catalogação-na-publicação''. Você pode utilizar este modelo como referência. 
% Porém, provavelmente a biblioteca da sua universidade lhe fornecerá um PDF
% com a ficha catalográfica definitiva após a defesa do trabalho. Quando estiver
% com o documento, salve-o como PDF no diretório do seu projeto e substitua todo
% o conteúdo de implementação deste arquivo pelo comando abaixo:
%
% \usepackage{pdfpages}		% necessário para comando \includepdf
% \begin{fichacatalografica}
%     \includepdf{fig_ficha_catalografica.pdf}
% \end{fichacatalografica}

% \begin{fichacatalografica}
% 	\sffamily
% 	\vspace*{\fill}                             % Posição vertical
% 	\begin{center}                              % Minipage Centralizado
% 	\fbox{\begin{minipage}[c][8cm]{13.5cm}		% Largura
% 	\small
% 	\imprimirautor
% 	%Sobrenome, Nome do autor
	
% 	\hspace{0.5cm} \imprimirtitulo  / \imprimirautor. --
% 	\imprimirlocal, \imprimirdata-
	
% 	\hspace{0.5cm} \thelastpage p. : il. (algumas color.) ; 30 cm.\\
	
% 	\hspace{0.5cm} \imprimirorientadorRotulo~\imprimirorientador\\
	
% 	\hspace{0.5cm}
% 	\parbox[t]{\textwidth}{\imprimirtipotrabalho~--~\imprimirinstituicao,
% 	\imprimirdata.}\\
	
% 	\hspace{0.5cm}
% 		1. Palavra-chave1.
% 		2. Palavra-chave2.
% 		2. Palavra-chave3.
% 		I. Orientador.
% 		II. Universidade xxx.
% 		III. Faculdade de xxx.
% 		IV. Título 			
% 	\end{minipage}}
% 	\end{center}
% \end{fichacatalografica}
% ---


% ---
% FOLHA DE APROVAÇÃO
% ---

% Isto é um exemplo de Folha de aprovação, elemento obrigatório da NBR
% 14724/2011 (seção 4.2.1.3). Você pode utilizar este modelo até a aprovação
% do trabalho. Após isso, substitua todo o conteúdo deste arquivo por uma
% imagem da página assinada pela banca com o comando abaixo:
%
% \begin{folhadeaprovacao}
% \includepdf{folhadeaprovacao_final.pdf}
% \end{folhadeaprovacao}
%
\setlength{\ABNTEXsignskip}{2cm}
\setlength{\ABNTEXsignwidth}{10cm}

\begin{folhadeaprovacao}
  \begin{center}
    {\ABNTEXchapterfont\large\imprimirautor}

    \vspace*{1cm}

    \begin{center}
      \ABNTEXchapterfont\bfseries\Large\imprimirtitulo
    \end{center}

    \vspace*{1cm}
    
    \hspace{.45\textwidth}
    \begin{minipage}{.5\textwidth}
        \imprimirpreambulo
    \end{minipage}%
   \end{center}

   \vspace*{1cm}
        
   \begin{center}
   Aprovado em 17 de julho de 2024.

   \vspace*{1cm}

   Banca examinadora:
   \end{center}

   \vspace*{-1cm}

   \assinatura{Prof. Ma. Caroline Santa Maria Rodrigues -- PUCRS}
   \assinatura{Prof. Dra. Clarissa Pinto Pizarro de Freitas -- PUCRS}
   \assinatura{\imprimirorientador -- PUCRS}
\end{folhadeaprovacao}
% ---


% ---
% Agradecimentos
% ---
\begin{agradecimentos}
Agradeço aos meus pais, Ademilson José Dias e Irinês Porfírio Dias, pelo apoio incondicional que me deram ao longo da vida,
incentivando mesmo as ideias questionáveis como aquela de fazer uma segunda graduação.

A minha irmã, Maria Eduarda Dias, cuja trajetória acadêmica é fonte de inspiração.

A minha namorada, Nathalia Fallavena Ceratti, que acompanhou a esta jornada de perto, oferecendo apoio e compreensão que
fizeram toda a diferença.

Aos meus amigos, Guilherme da Silva Petró, Norton Neves e Uilinton Luciano Marques da Silva, que foram compreensíveis nos
finais de semana em que não pude jogar RPG.

Ao professor Wagner de Lara Machado pela orientação e apoio no desenvolvimento deste trabalho.
\end{agradecimentos}
% ---


% ---
% RESUMOS
% ---
% Resumo em português
\setlength{\absparsep}{18pt} % ajusta o espaçamento dos parágrafos do resumo
\begin{resumo}
Desenvolvimentos na área de aprendizagem de máquina representam uma série de oportunidades para diversas áreas de pesquisa e de atuação profissional.
O uso de modelos de aprendizagem de máquina na predição de desfecho para tratamentos em saúde mental tem sido investigado e apresenta resultados promissores.
A aplicação desse tipo de modelo possibilitaria, por exemplo, maior assertividade na personalização de tratamentos psicoterápicos a nível individual.
Apesar dos benefícios potenciais, o tema ainda é pouco discutido nos espaços de pesquisa, formação e atuação em psicologia clínica. Este trabalho tem
por objetivo introduzir o tema de aprendizagem de máquina a psicólogos. Busca-se apresentar, de maneira acessível, os conceitos e terminologia básicos
em aprendizagem de máquina, proporcionando um panorama geral da área. Para complementar a exposição conceitual, é construído um exemplo a partir de dados
de um estudo sobre eficácia de intervenções em psicologia positiva para depressão via internet. Por fim, discutem-se algumas implicações do uso desse tipo
de tecnologia para a pesquisa e para a prática em psicoterapia.

 \textbf{Palavras-chave}: aprendizagem de máquina, psicoterapia, psicologia positiva, depressão.
\end{resumo}

% Resumo em inglês
\begin{resumo}[Abstract]
 \begin{otherlanguage*}{english}
   Advances in machine learning have opened up a multitude of opportunities for various research and professional fields. The use of machine learning models to
   predict treatment outcomes in mental health has been investigated showing promising results. The application of such models would allow greater assertiveness
   in tailoring psychotherapeutic treatments at the individual level. Despite its potential benefits, the topic remains unexplored in the fields of research, training
   and practice in clinical psychology. This work aims to introduce the topic of machine learning to psychologists. It presents the basic machine learning concepts and
   terminology in an accessible manner, providing a comprehensive overview of the field. To complement the conceptual discussion, an example is constructed using data
   from a study on the effectiveness of internet-based positive psychology interventions for depression. Finally, implications of using this type of technology in
   psychotherapy research and practice are discussed.
   \vspace{\onelineskip}
 
   \noindent 
   \textbf{Keywords}: machine learning, psychotherapy, positive psychology, depression.
 \end{otherlanguage*}
\end{resumo}


% ---
% LISTA DE ILUSTRAÇÕES
% ---
\pdfbookmark[0]{\listfigurename}{lof}
\listoffigures*
\cleardoublepage
% ---


% ---
% LISTA DE QUADROS
% ---
% \pdfbookmark[0]{\listofquadrosname}{loq}
% \listofquadros*
% \cleardoublepage
% ---


% ---
% LISTA DE TABELAS
% ---
% \pdfbookmark[0]{\listtablename}{lot}
% \listoftables*
% \cleardoublepage
% ---


% ---
% LISTA DE ABREVIATURAS E SIGLAS
% ---
% \begin{siglas}
%   \item[ABNT] Associação Brasileira de Normas Técnicas
%   \item[abnTeX] ABsurdas Normas para TeX
% \end{siglas}
% ---


% ---
% LISTA DE SÍMBOLOS
% ---
% \begin{simbolos}
%   \item[$ \Gamma $] Letra grega Gama
%   \item[$ \Lambda $] Lambda
%   \item[$ \zeta $] Letra grega minúscula zeta
%   \item[$ \in $] Pertence
% \end{simbolos}
% ---


% ---
% Sumário
% ---
\pdfbookmark[0]{\contentsname}{toc}
\tableofcontents*
\cleardoublepage
% ---


% ----------------------------------------------------------
% ELEMENTOS TEXTUAIS
% ----------------------------------------------------------
\textual

\include{./01-introducao/main.tex}

\include{./02-machine-learning/main.tex}

\include{./03-exemplo/main.tex}

\include{./04-conclusao/main.tex}

% ----------------------------------------------------------
% ELEMENTOS PÓS-TEXTUAIS
% ----------------------------------------------------------
\postextual
% ----------------------------------------------------------


% ----------------------------------------------------------
% REFERÊNCIAS
% ----------------------------------------------------------
\bibliography{references.bib}


% ----------------------------------------------------------
% GLOSSÁRIO
% ----------------------------------------------------------
%
% Consulte o manual da classe abntex2 para orientações sobre o glossário.
%
%\glossary

% ----------------------------------------------------------
% APÊNDICES
% ----------------------------------------------------------

% ---
% Inicia os apêndices
% ---
% \begin{apendicesenv}

% % Imprime uma página indicando o início dos apêndices
% \partapendices

% % ----------------------------------------------------------
% \chapter{Quisque libero justo}
% % ----------------------------------------------------------

% % ----------------------------------------------------------
% \chapter{Nullam elementum urna}
% % ----------------------------------------------------------

% \end{apendicesenv}
% ---


% ----------------------------------------------------------
% ANEXOS
% ----------------------------------------------------------

% ---
% Inicia os anexos
% ---
% \begin{anexosenv}

% % Imprime uma página indicando o início dos anexos
% \partanexos

% % ---
% \chapter{Morbi ultrices rutrum lorem}
% % ---

% % ---
% \chapter{Cras non urna sed feugiat cum}
% % ---


% % ---
% \chapter{Fusce facilisis lacinia dui}
% % ---

% \end{anexosenv}

%---------------------------------------------------------------------
% INDICE REMISSIVO
%---------------------------------------------------------------------
% \phantompart
% \printindex
%---------------------------------------------------------------------
\includepdf{./05-extras/contracapa.pdf}
\end{document}


\documentclass[
	% Opções da classe memoir
	12pt,				        % tamanho da fonte
	openany,			        % capítulos começam em pág ímpar (insere página vazia caso preciso)
	oneside,			        % para impressão em recto e verso. Oposto a oneside
	a4paper,			        % tamanho do papel. 
	% Opções da classe abntex2
	%chapter=TITLE,		        % títulos de capítulos convertidos em letras maiúsculas
	%section=TITLE,		        % títulos de seções convertidos em letras maiúsculas
	%subsection=TITLE,	        % títulos de subseções convertidos em letras maiúsculas
	%subsubsection=TITLE,       % títulos de subsubseções convertidos em letras maiúsculas
	% Opções do pacote babel
	english,			        % idioma adicional para hifenização
	brazil				        % o último idioma é o principal do documento
	]{abntex2}

% ---
% Pacotes básicos 
% ---
\usepackage{lmodern}			% Usa a fonte Latin Modern			
\usepackage[T1]{fontenc}		% Selecao de codigos de fonte.
\usepackage[utf8]{inputenc}		% Codificacao do documento (conversão automática dos acentos)
\usepackage{indentfirst}		% Indenta o primeiro parágrafo de cada seção.
\usepackage{color}				% Controle das cores
\usepackage{graphicx}			% Inclusão de gráficos
\usepackage{microtype} 			% para melhorias de justificação
% ---
		

% ---
% Pacotes de citações
% ---
% \usepackage[alf]{abntex2cite}						% Citações padrão ABNT
\usepackage[
  alf,
  abnt-nbr10520=1988,                     % Citações em sentence case
  abnt-emphasize=bf
  ]{abntex2cite}


% ---
% Pacotes adicionais
% ---
\usepackage{pdfpages}	% Inclusão de páginas em PDF


% ---
% Informações para CAPA e FOLHA DE ROSTO
% ---
\titulo{APRENDIZAGEM DE MÁQUINA PARA PREDIÇÃO DE DESFECHOS EM SAÚDE MENTAL: \textmd{UMA INTRODUÇÃO PARA PSICÓLOGOS}}
\autor{GUILHERME AUGUSTO DIAS}
\local{Porto Alegre}
\data{2024}
\orientador{Prof. Dr. Wagner de Lara Machado}
% \coorientador{}
% \instituicao{
%   Pontifícia Universidade Católica do Rio Grande do Sul
%   \par
%   Escola de Ciências da Saúde e da Vida
%   \par
%   Curso de Psicologia}
\tipotrabalho{TCC}
\preambulo{Trabalho de Conclusão de Curso apresentado como requisito para a obtenção do grau de Bacharel em Psicologia pela Escola de Ciências da Saúde e da Vida da Pontifícia Universidade Católica do Rio Grande do Sul.}
% ---

% ---
% Configurações de aparência do PDF final
% ---

% Alterando o aspecto da cor azul
\definecolor{blue}{RGB}{41,5,195}

% Informações do PDF
\makeatletter
\hypersetup{
		pdftitle={Aprendizagem de máquina para predição de desfechos em saúde mental: uma introdução para psicólogos}, 
		pdfauthor={Guilherme Augusto Dias},
    	pdfsubject={\imprimirpreambulo},
	    pdfcreator={LaTeX with abnTeX2},
		pdfkeywords={aprendizagem de máquina}{psicoterapia}{psicologia positiva}{depressão}, 
		colorlinks=true,		% true: colored links; false: boxed links
    	linkcolor=black,		% color of internal links
    	citecolor=black,		% color of links to bibliography
    	filecolor=black,		% color of file links
		urlcolor=black,
		bookmarksdepth=4
}
\makeatother
% --- 


% ---
% Posiciona figuras e tabelas no topo da página quando adicionadas sozinhas
% em um página em branco. Ver https://github.com/abntex/abntex2/issues/170
\makeatletter
\setlength{\@fptop}{5pt} % Set distance from top of page to first float
\makeatother
% ---

% ---
% Possibilita criação de Quadros e Lista de quadros.
% Ver https://github.com/abntex/abntex2/issues/176
%
\newcommand{\quadroname}{Quadro}
\newcommand{\listofquadrosname}{Lista de quadros}

\newfloat[chapter]{quadro}{loq}{\quadroname}
\newlistof{listofquadros}{loq}{\listofquadrosname}
\newlistentry{quadro}{loq}{0}

% Configurações para atender às regras da ABNT
\setfloatadjustment{quadro}{\centering}
\counterwithout{quadro}{chapter}
\renewcommand{\cftquadroname}{\quadroname\space} 
\renewcommand*{\cftquadroaftersnum}{\hfill--\hfill}

\setfloatlocations{quadro}{hbtp} % Ver https://github.com/abntex/abntex2/issues/176
% ---


% --- 
% Espaçamentos entre linhas e parágrafos 
% --- 

% O tamanho do parágrafo é dado por:
\setlength{\parindent}{1.3cm}

% Controle do espaçamento entre um parágrafo e outro:
\setlength{\parskip}{0.2cm}  % tente também \onelineskip

% ---
% Compila o índice
% ---
\makeindex
% ---


% ----
% Início do documento
% ----
\begin{document}

% Seleciona o idioma do documento (conforme pacotes do babel)
\selectlanguage{brazil}

% Retira espaço extra obsoleto entre as frases.
\frenchspacing 

% ----------------------------------------------------------
% ELEMENTOS PRÉ-TEXTUAIS
% ----------------------------------------------------------
\pretextual

% ---
% CAPA
% ---
\includepdf{./05-extras/capa.pdf}
% ---


% ---
% FOLHA DE ROSTO
% (o * indica que haverá a ficha bibliográfica)
% ---
\imprimirfolhaderosto
% ---


% ---
% Ficha bibliografica
% ---

% Isto é um exemplo de Ficha Catalográfica, ou ``Dados internacionais de
% catalogação-na-publicação''. Você pode utilizar este modelo como referência. 
% Porém, provavelmente a biblioteca da sua universidade lhe fornecerá um PDF
% com a ficha catalográfica definitiva após a defesa do trabalho. Quando estiver
% com o documento, salve-o como PDF no diretório do seu projeto e substitua todo
% o conteúdo de implementação deste arquivo pelo comando abaixo:
%
% \usepackage{pdfpages}		% necessário para comando \includepdf
% \begin{fichacatalografica}
%     \includepdf{fig_ficha_catalografica.pdf}
% \end{fichacatalografica}

% \begin{fichacatalografica}
% 	\sffamily
% 	\vspace*{\fill}                             % Posição vertical
% 	\begin{center}                              % Minipage Centralizado
% 	\fbox{\begin{minipage}[c][8cm]{13.5cm}		% Largura
% 	\small
% 	\imprimirautor
% 	%Sobrenome, Nome do autor
	
% 	\hspace{0.5cm} \imprimirtitulo  / \imprimirautor. --
% 	\imprimirlocal, \imprimirdata-
	
% 	\hspace{0.5cm} \thelastpage p. : il. (algumas color.) ; 30 cm.\\
	
% 	\hspace{0.5cm} \imprimirorientadorRotulo~\imprimirorientador\\
	
% 	\hspace{0.5cm}
% 	\parbox[t]{\textwidth}{\imprimirtipotrabalho~--~\imprimirinstituicao,
% 	\imprimirdata.}\\
	
% 	\hspace{0.5cm}
% 		1. Palavra-chave1.
% 		2. Palavra-chave2.
% 		2. Palavra-chave3.
% 		I. Orientador.
% 		II. Universidade xxx.
% 		III. Faculdade de xxx.
% 		IV. Título 			
% 	\end{minipage}}
% 	\end{center}
% \end{fichacatalografica}
% ---


% ---
% FOLHA DE APROVAÇÃO
% ---

% Isto é um exemplo de Folha de aprovação, elemento obrigatório da NBR
% 14724/2011 (seção 4.2.1.3). Você pode utilizar este modelo até a aprovação
% do trabalho. Após isso, substitua todo o conteúdo deste arquivo por uma
% imagem da página assinada pela banca com o comando abaixo:
%
% \begin{folhadeaprovacao}
% \includepdf{folhadeaprovacao_final.pdf}
% \end{folhadeaprovacao}
%
\setlength{\ABNTEXsignskip}{2cm}
\setlength{\ABNTEXsignwidth}{10cm}

\begin{folhadeaprovacao}
  \begin{center}
    {\ABNTEXchapterfont\large\imprimirautor}

    \vspace*{1cm}

    \begin{center}
      \ABNTEXchapterfont\bfseries\Large\imprimirtitulo
    \end{center}

    \vspace*{1cm}
    
    \hspace{.45\textwidth}
    \begin{minipage}{.5\textwidth}
        \imprimirpreambulo
    \end{minipage}%
   \end{center}

   \vspace*{1cm}
        
   \begin{center}
   Aprovado em 17 de julho de 2024.

   \vspace*{1cm}

   Banca examinadora:
   \end{center}

   \vspace*{-1cm}

   \assinatura{Prof. Ma. Caroline Santa Maria Rodrigues -- PUCRS}
   \assinatura{Prof. Dra. Clarissa Pinto Pizarro de Freitas -- PUCRS}
   \assinatura{\imprimirorientador -- PUCRS}
\end{folhadeaprovacao}
% ---


% ---
% Agradecimentos
% ---
\begin{agradecimentos}
Agradeço aos meus pais, Ademilson José Dias e Irinês Porfírio Dias, pelo apoio incondicional que me deram ao longo da vida,
incentivando mesmo as ideias questionáveis como aquela de fazer uma segunda graduação.

A minha irmã, Maria Eduarda Dias, cuja trajetória acadêmica é fonte de inspiração.

A minha namorada, Nathalia Fallavena Ceratti, que acompanhou a esta jornada de perto, oferecendo apoio e compreensão que
fizeram toda a diferença.

Aos meus amigos, Guilherme da Silva Petró, Norton Neves e Uilinton Luciano Marques da Silva, que foram compreensíveis nos
finais de semana em que não pude jogar RPG.

Ao professor Wagner de Lara Machado pela orientação e apoio no desenvolvimento deste trabalho.
\end{agradecimentos}
% ---


% ---
% RESUMOS
% ---
% Resumo em português
\setlength{\absparsep}{18pt} % ajusta o espaçamento dos parágrafos do resumo
\begin{resumo}
Desenvolvimentos na área de aprendizagem de máquina representam uma série de oportunidades para diversas áreas de pesquisa e de atuação profissional.
O uso de modelos de aprendizagem de máquina na predição de desfecho para tratamentos em saúde mental tem sido investigado e apresenta resultados promissores.
A aplicação desse tipo de modelo possibilitaria, por exemplo, maior assertividade na personalização de tratamentos psicoterápicos a nível individual.
Apesar dos benefícios potenciais, o tema ainda é pouco discutido nos espaços de pesquisa, formação e atuação em psicologia clínica. Este trabalho tem
por objetivo introduzir o tema de aprendizagem de máquina a psicólogos. Busca-se apresentar, de maneira acessível, os conceitos e terminologia básicos
em aprendizagem de máquina, proporcionando um panorama geral da área. Para complementar a exposição conceitual, é construído um exemplo a partir de dados
de um estudo sobre eficácia de intervenções em psicologia positiva para depressão via internet. Por fim, discutem-se algumas implicações do uso desse tipo
de tecnologia para a pesquisa e para a prática em psicoterapia.

 \textbf{Palavras-chave}: aprendizagem de máquina, psicoterapia, psicologia positiva, depressão.
\end{resumo}

% Resumo em inglês
\begin{resumo}[Abstract]
 \begin{otherlanguage*}{english}
   Advances in machine learning have opened up a multitude of opportunities for various research and professional fields. The use of machine learning models to
   predict treatment outcomes in mental health has been investigated showing promising results. The application of such models would allow greater assertiveness
   in tailoring psychotherapeutic treatments at the individual level. Despite its potential benefits, the topic remains unexplored in the fields of research, training
   and practice in clinical psychology. This work aims to introduce the topic of machine learning to psychologists. It presents the basic machine learning concepts and
   terminology in an accessible manner, providing a comprehensive overview of the field. To complement the conceptual discussion, an example is constructed using data
   from a study on the effectiveness of internet-based positive psychology interventions for depression. Finally, implications of using this type of technology in
   psychotherapy research and practice are discussed.
   \vspace{\onelineskip}
 
   \noindent 
   \textbf{Keywords}: machine learning, psychotherapy, positive psychology, depression.
 \end{otherlanguage*}
\end{resumo}


% ---
% LISTA DE ILUSTRAÇÕES
% ---
\pdfbookmark[0]{\listfigurename}{lof}
\listoffigures*
\cleardoublepage
% ---


% ---
% LISTA DE QUADROS
% ---
% \pdfbookmark[0]{\listofquadrosname}{loq}
% \listofquadros*
% \cleardoublepage
% ---


% ---
% LISTA DE TABELAS
% ---
% \pdfbookmark[0]{\listtablename}{lot}
% \listoftables*
% \cleardoublepage
% ---


% ---
% LISTA DE ABREVIATURAS E SIGLAS
% ---
% \begin{siglas}
%   \item[ABNT] Associação Brasileira de Normas Técnicas
%   \item[abnTeX] ABsurdas Normas para TeX
% \end{siglas}
% ---


% ---
% LISTA DE SÍMBOLOS
% ---
% \begin{simbolos}
%   \item[$ \Gamma $] Letra grega Gama
%   \item[$ \Lambda $] Lambda
%   \item[$ \zeta $] Letra grega minúscula zeta
%   \item[$ \in $] Pertence
% \end{simbolos}
% ---


% ---
% Sumário
% ---
\pdfbookmark[0]{\contentsname}{toc}
\tableofcontents*
\cleardoublepage
% ---


% ----------------------------------------------------------
% ELEMENTOS TEXTUAIS
% ----------------------------------------------------------
\textual

\include{./01-introducao/main.tex}

\include{./02-machine-learning/main.tex}

\include{./03-exemplo/main.tex}

\include{./04-conclusao/main.tex}

% ----------------------------------------------------------
% ELEMENTOS PÓS-TEXTUAIS
% ----------------------------------------------------------
\postextual
% ----------------------------------------------------------


% ----------------------------------------------------------
% REFERÊNCIAS
% ----------------------------------------------------------
\bibliography{references.bib}


% ----------------------------------------------------------
% GLOSSÁRIO
% ----------------------------------------------------------
%
% Consulte o manual da classe abntex2 para orientações sobre o glossário.
%
%\glossary

% ----------------------------------------------------------
% APÊNDICES
% ----------------------------------------------------------

% ---
% Inicia os apêndices
% ---
% \begin{apendicesenv}

% % Imprime uma página indicando o início dos apêndices
% \partapendices

% % ----------------------------------------------------------
% \chapter{Quisque libero justo}
% % ----------------------------------------------------------

% % ----------------------------------------------------------
% \chapter{Nullam elementum urna}
% % ----------------------------------------------------------

% \end{apendicesenv}
% ---


% ----------------------------------------------------------
% ANEXOS
% ----------------------------------------------------------

% ---
% Inicia os anexos
% ---
% \begin{anexosenv}

% % Imprime uma página indicando o início dos anexos
% \partanexos

% % ---
% \chapter{Morbi ultrices rutrum lorem}
% % ---

% % ---
% \chapter{Cras non urna sed feugiat cum}
% % ---


% % ---
% \chapter{Fusce facilisis lacinia dui}
% % ---

% \end{anexosenv}

%---------------------------------------------------------------------
% INDICE REMISSIVO
%---------------------------------------------------------------------
% \phantompart
% \printindex
%---------------------------------------------------------------------
\includepdf{./05-extras/contracapa.pdf}
\end{document}


% ----------------------------------------------------------
% ELEMENTOS PÓS-TEXTUAIS
% ----------------------------------------------------------
\postextual
% ----------------------------------------------------------


% ----------------------------------------------------------
% REFERÊNCIAS
% ----------------------------------------------------------
\bibliography{references.bib}


% ----------------------------------------------------------
% GLOSSÁRIO
% ----------------------------------------------------------
%
% Consulte o manual da classe abntex2 para orientações sobre o glossário.
%
%\glossary

% ----------------------------------------------------------
% APÊNDICES
% ----------------------------------------------------------

% ---
% Inicia os apêndices
% ---
% \begin{apendicesenv}

% % Imprime uma página indicando o início dos apêndices
% \partapendices

% % ----------------------------------------------------------
% \chapter{Quisque libero justo}
% % ----------------------------------------------------------

% % ----------------------------------------------------------
% \chapter{Nullam elementum urna}
% % ----------------------------------------------------------

% \end{apendicesenv}
% ---


% ----------------------------------------------------------
% ANEXOS
% ----------------------------------------------------------

% ---
% Inicia os anexos
% ---
% \begin{anexosenv}

% % Imprime uma página indicando o início dos anexos
% \partanexos

% % ---
% \chapter{Morbi ultrices rutrum lorem}
% % ---

% % ---
% \chapter{Cras non urna sed feugiat cum}
% % ---


% % ---
% \chapter{Fusce facilisis lacinia dui}
% % ---

% \end{anexosenv}

%---------------------------------------------------------------------
% INDICE REMISSIVO
%---------------------------------------------------------------------
% \phantompart
% \printindex
%---------------------------------------------------------------------
\includepdf{./05-extras/contracapa.pdf}
\end{document}


\documentclass[
	% Opções da classe memoir
	12pt,				        % tamanho da fonte
	openany,			        % capítulos começam em pág ímpar (insere página vazia caso preciso)
	oneside,			        % para impressão em recto e verso. Oposto a oneside
	a4paper,			        % tamanho do papel. 
	% Opções da classe abntex2
	%chapter=TITLE,		        % títulos de capítulos convertidos em letras maiúsculas
	%section=TITLE,		        % títulos de seções convertidos em letras maiúsculas
	%subsection=TITLE,	        % títulos de subseções convertidos em letras maiúsculas
	%subsubsection=TITLE,       % títulos de subsubseções convertidos em letras maiúsculas
	% Opções do pacote babel
	english,			        % idioma adicional para hifenização
	brazil				        % o último idioma é o principal do documento
	]{abntex2}

% ---
% Pacotes básicos 
% ---
\usepackage{lmodern}			% Usa a fonte Latin Modern			
\usepackage[T1]{fontenc}		% Selecao de codigos de fonte.
\usepackage[utf8]{inputenc}		% Codificacao do documento (conversão automática dos acentos)
\usepackage{indentfirst}		% Indenta o primeiro parágrafo de cada seção.
\usepackage{color}				% Controle das cores
\usepackage{graphicx}			% Inclusão de gráficos
\usepackage{microtype} 			% para melhorias de justificação
% ---
		

% ---
% Pacotes de citações
% ---
% \usepackage[alf]{abntex2cite}						% Citações padrão ABNT
\usepackage[
  alf,
  abnt-nbr10520=1988,                     % Citações em sentence case
  abnt-emphasize=bf
  ]{abntex2cite}


% ---
% Pacotes adicionais
% ---
\usepackage{pdfpages}	% Inclusão de páginas em PDF


% ---
% Informações para CAPA e FOLHA DE ROSTO
% ---
\titulo{APRENDIZAGEM DE MÁQUINA PARA PREDIÇÃO DE DESFECHOS EM SAÚDE MENTAL: \textmd{UMA INTRODUÇÃO PARA PSICÓLOGOS}}
\autor{GUILHERME AUGUSTO DIAS}
\local{Porto Alegre}
\data{2024}
\orientador{Prof. Dr. Wagner de Lara Machado}
% \coorientador{}
% \instituicao{
%   Pontifícia Universidade Católica do Rio Grande do Sul
%   \par
%   Escola de Ciências da Saúde e da Vida
%   \par
%   Curso de Psicologia}
\tipotrabalho{TCC}
\preambulo{Trabalho de Conclusão de Curso apresentado como requisito para a obtenção do grau de Bacharel em Psicologia pela Escola de Ciências da Saúde e da Vida da Pontifícia Universidade Católica do Rio Grande do Sul.}
% ---

% ---
% Configurações de aparência do PDF final
% ---

% Alterando o aspecto da cor azul
\definecolor{blue}{RGB}{41,5,195}

% Informações do PDF
\makeatletter
\hypersetup{
		pdftitle={Aprendizagem de máquina para predição de desfechos em saúde mental: uma introdução para psicólogos}, 
		pdfauthor={Guilherme Augusto Dias},
    	pdfsubject={\imprimirpreambulo},
	    pdfcreator={LaTeX with abnTeX2},
		pdfkeywords={aprendizagem de máquina}{psicoterapia}{psicologia positiva}{depressão}, 
		colorlinks=true,		% true: colored links; false: boxed links
    	linkcolor=black,		% color of internal links
    	citecolor=black,		% color of links to bibliography
    	filecolor=black,		% color of file links
		urlcolor=black,
		bookmarksdepth=4
}
\makeatother
% --- 


% ---
% Posiciona figuras e tabelas no topo da página quando adicionadas sozinhas
% em um página em branco. Ver https://github.com/abntex/abntex2/issues/170
\makeatletter
\setlength{\@fptop}{5pt} % Set distance from top of page to first float
\makeatother
% ---

% ---
% Possibilita criação de Quadros e Lista de quadros.
% Ver https://github.com/abntex/abntex2/issues/176
%
\newcommand{\quadroname}{Quadro}
\newcommand{\listofquadrosname}{Lista de quadros}

\newfloat[chapter]{quadro}{loq}{\quadroname}
\newlistof{listofquadros}{loq}{\listofquadrosname}
\newlistentry{quadro}{loq}{0}

% Configurações para atender às regras da ABNT
\setfloatadjustment{quadro}{\centering}
\counterwithout{quadro}{chapter}
\renewcommand{\cftquadroname}{\quadroname\space} 
\renewcommand*{\cftquadroaftersnum}{\hfill--\hfill}

\setfloatlocations{quadro}{hbtp} % Ver https://github.com/abntex/abntex2/issues/176
% ---


% --- 
% Espaçamentos entre linhas e parágrafos 
% --- 

% O tamanho do parágrafo é dado por:
\setlength{\parindent}{1.3cm}

% Controle do espaçamento entre um parágrafo e outro:
\setlength{\parskip}{0.2cm}  % tente também \onelineskip

% ---
% Compila o índice
% ---
\makeindex
% ---


% ----
% Início do documento
% ----
\begin{document}

% Seleciona o idioma do documento (conforme pacotes do babel)
\selectlanguage{brazil}

% Retira espaço extra obsoleto entre as frases.
\frenchspacing 

% ----------------------------------------------------------
% ELEMENTOS PRÉ-TEXTUAIS
% ----------------------------------------------------------
\pretextual

% ---
% CAPA
% ---
\includepdf{./05-extras/capa.pdf}
% ---


% ---
% FOLHA DE ROSTO
% (o * indica que haverá a ficha bibliográfica)
% ---
\imprimirfolhaderosto
% ---


% ---
% Ficha bibliografica
% ---

% Isto é um exemplo de Ficha Catalográfica, ou ``Dados internacionais de
% catalogação-na-publicação''. Você pode utilizar este modelo como referência. 
% Porém, provavelmente a biblioteca da sua universidade lhe fornecerá um PDF
% com a ficha catalográfica definitiva após a defesa do trabalho. Quando estiver
% com o documento, salve-o como PDF no diretório do seu projeto e substitua todo
% o conteúdo de implementação deste arquivo pelo comando abaixo:
%
% \usepackage{pdfpages}		% necessário para comando \includepdf
% \begin{fichacatalografica}
%     \includepdf{fig_ficha_catalografica.pdf}
% \end{fichacatalografica}

% \begin{fichacatalografica}
% 	\sffamily
% 	\vspace*{\fill}                             % Posição vertical
% 	\begin{center}                              % Minipage Centralizado
% 	\fbox{\begin{minipage}[c][8cm]{13.5cm}		% Largura
% 	\small
% 	\imprimirautor
% 	%Sobrenome, Nome do autor
	
% 	\hspace{0.5cm} \imprimirtitulo  / \imprimirautor. --
% 	\imprimirlocal, \imprimirdata-
	
% 	\hspace{0.5cm} \thelastpage p. : il. (algumas color.) ; 30 cm.\\
	
% 	\hspace{0.5cm} \imprimirorientadorRotulo~\imprimirorientador\\
	
% 	\hspace{0.5cm}
% 	\parbox[t]{\textwidth}{\imprimirtipotrabalho~--~\imprimirinstituicao,
% 	\imprimirdata.}\\
	
% 	\hspace{0.5cm}
% 		1. Palavra-chave1.
% 		2. Palavra-chave2.
% 		2. Palavra-chave3.
% 		I. Orientador.
% 		II. Universidade xxx.
% 		III. Faculdade de xxx.
% 		IV. Título 			
% 	\end{minipage}}
% 	\end{center}
% \end{fichacatalografica}
% ---


% ---
% FOLHA DE APROVAÇÃO
% ---

% Isto é um exemplo de Folha de aprovação, elemento obrigatório da NBR
% 14724/2011 (seção 4.2.1.3). Você pode utilizar este modelo até a aprovação
% do trabalho. Após isso, substitua todo o conteúdo deste arquivo por uma
% imagem da página assinada pela banca com o comando abaixo:
%
% \begin{folhadeaprovacao}
% \includepdf{folhadeaprovacao_final.pdf}
% \end{folhadeaprovacao}
%
\setlength{\ABNTEXsignskip}{2cm}
\setlength{\ABNTEXsignwidth}{10cm}

\begin{folhadeaprovacao}
  \begin{center}
    {\ABNTEXchapterfont\large\imprimirautor}

    \vspace*{1cm}

    \begin{center}
      \ABNTEXchapterfont\bfseries\Large\imprimirtitulo
    \end{center}

    \vspace*{1cm}
    
    \hspace{.45\textwidth}
    \begin{minipage}{.5\textwidth}
        \imprimirpreambulo
    \end{minipage}%
   \end{center}

   \vspace*{1cm}
        
   \begin{center}
   Aprovado em 17 de julho de 2024.

   \vspace*{1cm}

   Banca examinadora:
   \end{center}

   \vspace*{-1cm}

   \assinatura{Prof. Ma. Caroline Santa Maria Rodrigues -- PUCRS}
   \assinatura{Prof. Dra. Clarissa Pinto Pizarro de Freitas -- PUCRS}
   \assinatura{\imprimirorientador -- PUCRS}
\end{folhadeaprovacao}
% ---


% ---
% Agradecimentos
% ---
\begin{agradecimentos}
Agradeço aos meus pais, Ademilson José Dias e Irinês Porfírio Dias, pelo apoio incondicional que me deram ao longo da vida,
incentivando mesmo as ideias questionáveis como aquela de fazer uma segunda graduação.

A minha irmã, Maria Eduarda Dias, cuja trajetória acadêmica é fonte de inspiração.

A minha namorada, Nathalia Fallavena Ceratti, que acompanhou a esta jornada de perto, oferecendo apoio e compreensão que
fizeram toda a diferença.

Aos meus amigos, Guilherme da Silva Petró, Norton Neves e Uilinton Luciano Marques da Silva, que foram compreensíveis nos
finais de semana em que não pude jogar RPG.

Ao professor Wagner de Lara Machado pela orientação e apoio no desenvolvimento deste trabalho.
\end{agradecimentos}
% ---


% ---
% RESUMOS
% ---
% Resumo em português
\setlength{\absparsep}{18pt} % ajusta o espaçamento dos parágrafos do resumo
\begin{resumo}
Desenvolvimentos na área de aprendizagem de máquina representam uma série de oportunidades para diversas áreas de pesquisa e de atuação profissional.
O uso de modelos de aprendizagem de máquina na predição de desfecho para tratamentos em saúde mental tem sido investigado e apresenta resultados promissores.
A aplicação desse tipo de modelo possibilitaria, por exemplo, maior assertividade na personalização de tratamentos psicoterápicos a nível individual.
Apesar dos benefícios potenciais, o tema ainda é pouco discutido nos espaços de pesquisa, formação e atuação em psicologia clínica. Este trabalho tem
por objetivo introduzir o tema de aprendizagem de máquina a psicólogos. Busca-se apresentar, de maneira acessível, os conceitos e terminologia básicos
em aprendizagem de máquina, proporcionando um panorama geral da área. Para complementar a exposição conceitual, é construído um exemplo a partir de dados
de um estudo sobre eficácia de intervenções em psicologia positiva para depressão via internet. Por fim, discutem-se algumas implicações do uso desse tipo
de tecnologia para a pesquisa e para a prática em psicoterapia.

 \textbf{Palavras-chave}: aprendizagem de máquina, psicoterapia, psicologia positiva, depressão.
\end{resumo}

% Resumo em inglês
\begin{resumo}[Abstract]
 \begin{otherlanguage*}{english}
   Advances in machine learning have opened up a multitude of opportunities for various research and professional fields. The use of machine learning models to
   predict treatment outcomes in mental health has been investigated showing promising results. The application of such models would allow greater assertiveness
   in tailoring psychotherapeutic treatments at the individual level. Despite its potential benefits, the topic remains unexplored in the fields of research, training
   and practice in clinical psychology. This work aims to introduce the topic of machine learning to psychologists. It presents the basic machine learning concepts and
   terminology in an accessible manner, providing a comprehensive overview of the field. To complement the conceptual discussion, an example is constructed using data
   from a study on the effectiveness of internet-based positive psychology interventions for depression. Finally, implications of using this type of technology in
   psychotherapy research and practice are discussed.
   \vspace{\onelineskip}
 
   \noindent 
   \textbf{Keywords}: machine learning, psychotherapy, positive psychology, depression.
 \end{otherlanguage*}
\end{resumo}


% ---
% LISTA DE ILUSTRAÇÕES
% ---
\pdfbookmark[0]{\listfigurename}{lof}
\listoffigures*
\cleardoublepage
% ---


% ---
% LISTA DE QUADROS
% ---
% \pdfbookmark[0]{\listofquadrosname}{loq}
% \listofquadros*
% \cleardoublepage
% ---


% ---
% LISTA DE TABELAS
% ---
% \pdfbookmark[0]{\listtablename}{lot}
% \listoftables*
% \cleardoublepage
% ---


% ---
% LISTA DE ABREVIATURAS E SIGLAS
% ---
% \begin{siglas}
%   \item[ABNT] Associação Brasileira de Normas Técnicas
%   \item[abnTeX] ABsurdas Normas para TeX
% \end{siglas}
% ---


% ---
% LISTA DE SÍMBOLOS
% ---
% \begin{simbolos}
%   \item[$ \Gamma $] Letra grega Gama
%   \item[$ \Lambda $] Lambda
%   \item[$ \zeta $] Letra grega minúscula zeta
%   \item[$ \in $] Pertence
% \end{simbolos}
% ---


% ---
% Sumário
% ---
\pdfbookmark[0]{\contentsname}{toc}
\tableofcontents*
\cleardoublepage
% ---


% ----------------------------------------------------------
% ELEMENTOS TEXTUAIS
% ----------------------------------------------------------
\textual

\documentclass[
	% Opções da classe memoir
	12pt,				        % tamanho da fonte
	openany,			        % capítulos começam em pág ímpar (insere página vazia caso preciso)
	oneside,			        % para impressão em recto e verso. Oposto a oneside
	a4paper,			        % tamanho do papel. 
	% Opções da classe abntex2
	%chapter=TITLE,		        % títulos de capítulos convertidos em letras maiúsculas
	%section=TITLE,		        % títulos de seções convertidos em letras maiúsculas
	%subsection=TITLE,	        % títulos de subseções convertidos em letras maiúsculas
	%subsubsection=TITLE,       % títulos de subsubseções convertidos em letras maiúsculas
	% Opções do pacote babel
	english,			        % idioma adicional para hifenização
	brazil				        % o último idioma é o principal do documento
	]{abntex2}

% ---
% Pacotes básicos 
% ---
\usepackage{lmodern}			% Usa a fonte Latin Modern			
\usepackage[T1]{fontenc}		% Selecao de codigos de fonte.
\usepackage[utf8]{inputenc}		% Codificacao do documento (conversão automática dos acentos)
\usepackage{indentfirst}		% Indenta o primeiro parágrafo de cada seção.
\usepackage{color}				% Controle das cores
\usepackage{graphicx}			% Inclusão de gráficos
\usepackage{microtype} 			% para melhorias de justificação
% ---
		

% ---
% Pacotes de citações
% ---
% \usepackage[alf]{abntex2cite}						% Citações padrão ABNT
\usepackage[
  alf,
  abnt-nbr10520=1988,                     % Citações em sentence case
  abnt-emphasize=bf
  ]{abntex2cite}


% ---
% Pacotes adicionais
% ---
\usepackage{pdfpages}	% Inclusão de páginas em PDF


% ---
% Informações para CAPA e FOLHA DE ROSTO
% ---
\titulo{APRENDIZAGEM DE MÁQUINA PARA PREDIÇÃO DE DESFECHOS EM SAÚDE MENTAL: \textmd{UMA INTRODUÇÃO PARA PSICÓLOGOS}}
\autor{GUILHERME AUGUSTO DIAS}
\local{Porto Alegre}
\data{2024}
\orientador{Prof. Dr. Wagner de Lara Machado}
% \coorientador{}
% \instituicao{
%   Pontifícia Universidade Católica do Rio Grande do Sul
%   \par
%   Escola de Ciências da Saúde e da Vida
%   \par
%   Curso de Psicologia}
\tipotrabalho{TCC}
\preambulo{Trabalho de Conclusão de Curso apresentado como requisito para a obtenção do grau de Bacharel em Psicologia pela Escola de Ciências da Saúde e da Vida da Pontifícia Universidade Católica do Rio Grande do Sul.}
% ---

% ---
% Configurações de aparência do PDF final
% ---

% Alterando o aspecto da cor azul
\definecolor{blue}{RGB}{41,5,195}

% Informações do PDF
\makeatletter
\hypersetup{
		pdftitle={Aprendizagem de máquina para predição de desfechos em saúde mental: uma introdução para psicólogos}, 
		pdfauthor={Guilherme Augusto Dias},
    	pdfsubject={\imprimirpreambulo},
	    pdfcreator={LaTeX with abnTeX2},
		pdfkeywords={aprendizagem de máquina}{psicoterapia}{psicologia positiva}{depressão}, 
		colorlinks=true,		% true: colored links; false: boxed links
    	linkcolor=black,		% color of internal links
    	citecolor=black,		% color of links to bibliography
    	filecolor=black,		% color of file links
		urlcolor=black,
		bookmarksdepth=4
}
\makeatother
% --- 


% ---
% Posiciona figuras e tabelas no topo da página quando adicionadas sozinhas
% em um página em branco. Ver https://github.com/abntex/abntex2/issues/170
\makeatletter
\setlength{\@fptop}{5pt} % Set distance from top of page to first float
\makeatother
% ---

% ---
% Possibilita criação de Quadros e Lista de quadros.
% Ver https://github.com/abntex/abntex2/issues/176
%
\newcommand{\quadroname}{Quadro}
\newcommand{\listofquadrosname}{Lista de quadros}

\newfloat[chapter]{quadro}{loq}{\quadroname}
\newlistof{listofquadros}{loq}{\listofquadrosname}
\newlistentry{quadro}{loq}{0}

% Configurações para atender às regras da ABNT
\setfloatadjustment{quadro}{\centering}
\counterwithout{quadro}{chapter}
\renewcommand{\cftquadroname}{\quadroname\space} 
\renewcommand*{\cftquadroaftersnum}{\hfill--\hfill}

\setfloatlocations{quadro}{hbtp} % Ver https://github.com/abntex/abntex2/issues/176
% ---


% --- 
% Espaçamentos entre linhas e parágrafos 
% --- 

% O tamanho do parágrafo é dado por:
\setlength{\parindent}{1.3cm}

% Controle do espaçamento entre um parágrafo e outro:
\setlength{\parskip}{0.2cm}  % tente também \onelineskip

% ---
% Compila o índice
% ---
\makeindex
% ---


% ----
% Início do documento
% ----
\begin{document}

% Seleciona o idioma do documento (conforme pacotes do babel)
\selectlanguage{brazil}

% Retira espaço extra obsoleto entre as frases.
\frenchspacing 

% ----------------------------------------------------------
% ELEMENTOS PRÉ-TEXTUAIS
% ----------------------------------------------------------
\pretextual

% ---
% CAPA
% ---
\includepdf{./05-extras/capa.pdf}
% ---


% ---
% FOLHA DE ROSTO
% (o * indica que haverá a ficha bibliográfica)
% ---
\imprimirfolhaderosto
% ---


% ---
% Ficha bibliografica
% ---

% Isto é um exemplo de Ficha Catalográfica, ou ``Dados internacionais de
% catalogação-na-publicação''. Você pode utilizar este modelo como referência. 
% Porém, provavelmente a biblioteca da sua universidade lhe fornecerá um PDF
% com a ficha catalográfica definitiva após a defesa do trabalho. Quando estiver
% com o documento, salve-o como PDF no diretório do seu projeto e substitua todo
% o conteúdo de implementação deste arquivo pelo comando abaixo:
%
% \usepackage{pdfpages}		% necessário para comando \includepdf
% \begin{fichacatalografica}
%     \includepdf{fig_ficha_catalografica.pdf}
% \end{fichacatalografica}

% \begin{fichacatalografica}
% 	\sffamily
% 	\vspace*{\fill}                             % Posição vertical
% 	\begin{center}                              % Minipage Centralizado
% 	\fbox{\begin{minipage}[c][8cm]{13.5cm}		% Largura
% 	\small
% 	\imprimirautor
% 	%Sobrenome, Nome do autor
	
% 	\hspace{0.5cm} \imprimirtitulo  / \imprimirautor. --
% 	\imprimirlocal, \imprimirdata-
	
% 	\hspace{0.5cm} \thelastpage p. : il. (algumas color.) ; 30 cm.\\
	
% 	\hspace{0.5cm} \imprimirorientadorRotulo~\imprimirorientador\\
	
% 	\hspace{0.5cm}
% 	\parbox[t]{\textwidth}{\imprimirtipotrabalho~--~\imprimirinstituicao,
% 	\imprimirdata.}\\
	
% 	\hspace{0.5cm}
% 		1. Palavra-chave1.
% 		2. Palavra-chave2.
% 		2. Palavra-chave3.
% 		I. Orientador.
% 		II. Universidade xxx.
% 		III. Faculdade de xxx.
% 		IV. Título 			
% 	\end{minipage}}
% 	\end{center}
% \end{fichacatalografica}
% ---


% ---
% FOLHA DE APROVAÇÃO
% ---

% Isto é um exemplo de Folha de aprovação, elemento obrigatório da NBR
% 14724/2011 (seção 4.2.1.3). Você pode utilizar este modelo até a aprovação
% do trabalho. Após isso, substitua todo o conteúdo deste arquivo por uma
% imagem da página assinada pela banca com o comando abaixo:
%
% \begin{folhadeaprovacao}
% \includepdf{folhadeaprovacao_final.pdf}
% \end{folhadeaprovacao}
%
\setlength{\ABNTEXsignskip}{2cm}
\setlength{\ABNTEXsignwidth}{10cm}

\begin{folhadeaprovacao}
  \begin{center}
    {\ABNTEXchapterfont\large\imprimirautor}

    \vspace*{1cm}

    \begin{center}
      \ABNTEXchapterfont\bfseries\Large\imprimirtitulo
    \end{center}

    \vspace*{1cm}
    
    \hspace{.45\textwidth}
    \begin{minipage}{.5\textwidth}
        \imprimirpreambulo
    \end{minipage}%
   \end{center}

   \vspace*{1cm}
        
   \begin{center}
   Aprovado em 17 de julho de 2024.

   \vspace*{1cm}

   Banca examinadora:
   \end{center}

   \vspace*{-1cm}

   \assinatura{Prof. Ma. Caroline Santa Maria Rodrigues -- PUCRS}
   \assinatura{Prof. Dra. Clarissa Pinto Pizarro de Freitas -- PUCRS}
   \assinatura{\imprimirorientador -- PUCRS}
\end{folhadeaprovacao}
% ---


% ---
% Agradecimentos
% ---
\begin{agradecimentos}
Agradeço aos meus pais, Ademilson José Dias e Irinês Porfírio Dias, pelo apoio incondicional que me deram ao longo da vida,
incentivando mesmo as ideias questionáveis como aquela de fazer uma segunda graduação.

A minha irmã, Maria Eduarda Dias, cuja trajetória acadêmica é fonte de inspiração.

A minha namorada, Nathalia Fallavena Ceratti, que acompanhou a esta jornada de perto, oferecendo apoio e compreensão que
fizeram toda a diferença.

Aos meus amigos, Guilherme da Silva Petró, Norton Neves e Uilinton Luciano Marques da Silva, que foram compreensíveis nos
finais de semana em que não pude jogar RPG.

Ao professor Wagner de Lara Machado pela orientação e apoio no desenvolvimento deste trabalho.
\end{agradecimentos}
% ---


% ---
% RESUMOS
% ---
% Resumo em português
\setlength{\absparsep}{18pt} % ajusta o espaçamento dos parágrafos do resumo
\begin{resumo}
Desenvolvimentos na área de aprendizagem de máquina representam uma série de oportunidades para diversas áreas de pesquisa e de atuação profissional.
O uso de modelos de aprendizagem de máquina na predição de desfecho para tratamentos em saúde mental tem sido investigado e apresenta resultados promissores.
A aplicação desse tipo de modelo possibilitaria, por exemplo, maior assertividade na personalização de tratamentos psicoterápicos a nível individual.
Apesar dos benefícios potenciais, o tema ainda é pouco discutido nos espaços de pesquisa, formação e atuação em psicologia clínica. Este trabalho tem
por objetivo introduzir o tema de aprendizagem de máquina a psicólogos. Busca-se apresentar, de maneira acessível, os conceitos e terminologia básicos
em aprendizagem de máquina, proporcionando um panorama geral da área. Para complementar a exposição conceitual, é construído um exemplo a partir de dados
de um estudo sobre eficácia de intervenções em psicologia positiva para depressão via internet. Por fim, discutem-se algumas implicações do uso desse tipo
de tecnologia para a pesquisa e para a prática em psicoterapia.

 \textbf{Palavras-chave}: aprendizagem de máquina, psicoterapia, psicologia positiva, depressão.
\end{resumo}

% Resumo em inglês
\begin{resumo}[Abstract]
 \begin{otherlanguage*}{english}
   Advances in machine learning have opened up a multitude of opportunities for various research and professional fields. The use of machine learning models to
   predict treatment outcomes in mental health has been investigated showing promising results. The application of such models would allow greater assertiveness
   in tailoring psychotherapeutic treatments at the individual level. Despite its potential benefits, the topic remains unexplored in the fields of research, training
   and practice in clinical psychology. This work aims to introduce the topic of machine learning to psychologists. It presents the basic machine learning concepts and
   terminology in an accessible manner, providing a comprehensive overview of the field. To complement the conceptual discussion, an example is constructed using data
   from a study on the effectiveness of internet-based positive psychology interventions for depression. Finally, implications of using this type of technology in
   psychotherapy research and practice are discussed.
   \vspace{\onelineskip}
 
   \noindent 
   \textbf{Keywords}: machine learning, psychotherapy, positive psychology, depression.
 \end{otherlanguage*}
\end{resumo}


% ---
% LISTA DE ILUSTRAÇÕES
% ---
\pdfbookmark[0]{\listfigurename}{lof}
\listoffigures*
\cleardoublepage
% ---


% ---
% LISTA DE QUADROS
% ---
% \pdfbookmark[0]{\listofquadrosname}{loq}
% \listofquadros*
% \cleardoublepage
% ---


% ---
% LISTA DE TABELAS
% ---
% \pdfbookmark[0]{\listtablename}{lot}
% \listoftables*
% \cleardoublepage
% ---


% ---
% LISTA DE ABREVIATURAS E SIGLAS
% ---
% \begin{siglas}
%   \item[ABNT] Associação Brasileira de Normas Técnicas
%   \item[abnTeX] ABsurdas Normas para TeX
% \end{siglas}
% ---


% ---
% LISTA DE SÍMBOLOS
% ---
% \begin{simbolos}
%   \item[$ \Gamma $] Letra grega Gama
%   \item[$ \Lambda $] Lambda
%   \item[$ \zeta $] Letra grega minúscula zeta
%   \item[$ \in $] Pertence
% \end{simbolos}
% ---


% ---
% Sumário
% ---
\pdfbookmark[0]{\contentsname}{toc}
\tableofcontents*
\cleardoublepage
% ---


% ----------------------------------------------------------
% ELEMENTOS TEXTUAIS
% ----------------------------------------------------------
\textual

\include{./01-introducao/main.tex}

\include{./02-machine-learning/main.tex}

\include{./03-exemplo/main.tex}

\include{./04-conclusao/main.tex}

% ----------------------------------------------------------
% ELEMENTOS PÓS-TEXTUAIS
% ----------------------------------------------------------
\postextual
% ----------------------------------------------------------


% ----------------------------------------------------------
% REFERÊNCIAS
% ----------------------------------------------------------
\bibliography{references.bib}


% ----------------------------------------------------------
% GLOSSÁRIO
% ----------------------------------------------------------
%
% Consulte o manual da classe abntex2 para orientações sobre o glossário.
%
%\glossary

% ----------------------------------------------------------
% APÊNDICES
% ----------------------------------------------------------

% ---
% Inicia os apêndices
% ---
% \begin{apendicesenv}

% % Imprime uma página indicando o início dos apêndices
% \partapendices

% % ----------------------------------------------------------
% \chapter{Quisque libero justo}
% % ----------------------------------------------------------

% % ----------------------------------------------------------
% \chapter{Nullam elementum urna}
% % ----------------------------------------------------------

% \end{apendicesenv}
% ---


% ----------------------------------------------------------
% ANEXOS
% ----------------------------------------------------------

% ---
% Inicia os anexos
% ---
% \begin{anexosenv}

% % Imprime uma página indicando o início dos anexos
% \partanexos

% % ---
% \chapter{Morbi ultrices rutrum lorem}
% % ---

% % ---
% \chapter{Cras non urna sed feugiat cum}
% % ---


% % ---
% \chapter{Fusce facilisis lacinia dui}
% % ---

% \end{anexosenv}

%---------------------------------------------------------------------
% INDICE REMISSIVO
%---------------------------------------------------------------------
% \phantompart
% \printindex
%---------------------------------------------------------------------
\includepdf{./05-extras/contracapa.pdf}
\end{document}


\documentclass[
	% Opções da classe memoir
	12pt,				        % tamanho da fonte
	openany,			        % capítulos começam em pág ímpar (insere página vazia caso preciso)
	oneside,			        % para impressão em recto e verso. Oposto a oneside
	a4paper,			        % tamanho do papel. 
	% Opções da classe abntex2
	%chapter=TITLE,		        % títulos de capítulos convertidos em letras maiúsculas
	%section=TITLE,		        % títulos de seções convertidos em letras maiúsculas
	%subsection=TITLE,	        % títulos de subseções convertidos em letras maiúsculas
	%subsubsection=TITLE,       % títulos de subsubseções convertidos em letras maiúsculas
	% Opções do pacote babel
	english,			        % idioma adicional para hifenização
	brazil				        % o último idioma é o principal do documento
	]{abntex2}

% ---
% Pacotes básicos 
% ---
\usepackage{lmodern}			% Usa a fonte Latin Modern			
\usepackage[T1]{fontenc}		% Selecao de codigos de fonte.
\usepackage[utf8]{inputenc}		% Codificacao do documento (conversão automática dos acentos)
\usepackage{indentfirst}		% Indenta o primeiro parágrafo de cada seção.
\usepackage{color}				% Controle das cores
\usepackage{graphicx}			% Inclusão de gráficos
\usepackage{microtype} 			% para melhorias de justificação
% ---
		

% ---
% Pacotes de citações
% ---
% \usepackage[alf]{abntex2cite}						% Citações padrão ABNT
\usepackage[
  alf,
  abnt-nbr10520=1988,                     % Citações em sentence case
  abnt-emphasize=bf
  ]{abntex2cite}


% ---
% Pacotes adicionais
% ---
\usepackage{pdfpages}	% Inclusão de páginas em PDF


% ---
% Informações para CAPA e FOLHA DE ROSTO
% ---
\titulo{APRENDIZAGEM DE MÁQUINA PARA PREDIÇÃO DE DESFECHOS EM SAÚDE MENTAL: \textmd{UMA INTRODUÇÃO PARA PSICÓLOGOS}}
\autor{GUILHERME AUGUSTO DIAS}
\local{Porto Alegre}
\data{2024}
\orientador{Prof. Dr. Wagner de Lara Machado}
% \coorientador{}
% \instituicao{
%   Pontifícia Universidade Católica do Rio Grande do Sul
%   \par
%   Escola de Ciências da Saúde e da Vida
%   \par
%   Curso de Psicologia}
\tipotrabalho{TCC}
\preambulo{Trabalho de Conclusão de Curso apresentado como requisito para a obtenção do grau de Bacharel em Psicologia pela Escola de Ciências da Saúde e da Vida da Pontifícia Universidade Católica do Rio Grande do Sul.}
% ---

% ---
% Configurações de aparência do PDF final
% ---

% Alterando o aspecto da cor azul
\definecolor{blue}{RGB}{41,5,195}

% Informações do PDF
\makeatletter
\hypersetup{
		pdftitle={Aprendizagem de máquina para predição de desfechos em saúde mental: uma introdução para psicólogos}, 
		pdfauthor={Guilherme Augusto Dias},
    	pdfsubject={\imprimirpreambulo},
	    pdfcreator={LaTeX with abnTeX2},
		pdfkeywords={aprendizagem de máquina}{psicoterapia}{psicologia positiva}{depressão}, 
		colorlinks=true,		% true: colored links; false: boxed links
    	linkcolor=black,		% color of internal links
    	citecolor=black,		% color of links to bibliography
    	filecolor=black,		% color of file links
		urlcolor=black,
		bookmarksdepth=4
}
\makeatother
% --- 


% ---
% Posiciona figuras e tabelas no topo da página quando adicionadas sozinhas
% em um página em branco. Ver https://github.com/abntex/abntex2/issues/170
\makeatletter
\setlength{\@fptop}{5pt} % Set distance from top of page to first float
\makeatother
% ---

% ---
% Possibilita criação de Quadros e Lista de quadros.
% Ver https://github.com/abntex/abntex2/issues/176
%
\newcommand{\quadroname}{Quadro}
\newcommand{\listofquadrosname}{Lista de quadros}

\newfloat[chapter]{quadro}{loq}{\quadroname}
\newlistof{listofquadros}{loq}{\listofquadrosname}
\newlistentry{quadro}{loq}{0}

% Configurações para atender às regras da ABNT
\setfloatadjustment{quadro}{\centering}
\counterwithout{quadro}{chapter}
\renewcommand{\cftquadroname}{\quadroname\space} 
\renewcommand*{\cftquadroaftersnum}{\hfill--\hfill}

\setfloatlocations{quadro}{hbtp} % Ver https://github.com/abntex/abntex2/issues/176
% ---


% --- 
% Espaçamentos entre linhas e parágrafos 
% --- 

% O tamanho do parágrafo é dado por:
\setlength{\parindent}{1.3cm}

% Controle do espaçamento entre um parágrafo e outro:
\setlength{\parskip}{0.2cm}  % tente também \onelineskip

% ---
% Compila o índice
% ---
\makeindex
% ---


% ----
% Início do documento
% ----
\begin{document}

% Seleciona o idioma do documento (conforme pacotes do babel)
\selectlanguage{brazil}

% Retira espaço extra obsoleto entre as frases.
\frenchspacing 

% ----------------------------------------------------------
% ELEMENTOS PRÉ-TEXTUAIS
% ----------------------------------------------------------
\pretextual

% ---
% CAPA
% ---
\includepdf{./05-extras/capa.pdf}
% ---


% ---
% FOLHA DE ROSTO
% (o * indica que haverá a ficha bibliográfica)
% ---
\imprimirfolhaderosto
% ---


% ---
% Ficha bibliografica
% ---

% Isto é um exemplo de Ficha Catalográfica, ou ``Dados internacionais de
% catalogação-na-publicação''. Você pode utilizar este modelo como referência. 
% Porém, provavelmente a biblioteca da sua universidade lhe fornecerá um PDF
% com a ficha catalográfica definitiva após a defesa do trabalho. Quando estiver
% com o documento, salve-o como PDF no diretório do seu projeto e substitua todo
% o conteúdo de implementação deste arquivo pelo comando abaixo:
%
% \usepackage{pdfpages}		% necessário para comando \includepdf
% \begin{fichacatalografica}
%     \includepdf{fig_ficha_catalografica.pdf}
% \end{fichacatalografica}

% \begin{fichacatalografica}
% 	\sffamily
% 	\vspace*{\fill}                             % Posição vertical
% 	\begin{center}                              % Minipage Centralizado
% 	\fbox{\begin{minipage}[c][8cm]{13.5cm}		% Largura
% 	\small
% 	\imprimirautor
% 	%Sobrenome, Nome do autor
	
% 	\hspace{0.5cm} \imprimirtitulo  / \imprimirautor. --
% 	\imprimirlocal, \imprimirdata-
	
% 	\hspace{0.5cm} \thelastpage p. : il. (algumas color.) ; 30 cm.\\
	
% 	\hspace{0.5cm} \imprimirorientadorRotulo~\imprimirorientador\\
	
% 	\hspace{0.5cm}
% 	\parbox[t]{\textwidth}{\imprimirtipotrabalho~--~\imprimirinstituicao,
% 	\imprimirdata.}\\
	
% 	\hspace{0.5cm}
% 		1. Palavra-chave1.
% 		2. Palavra-chave2.
% 		2. Palavra-chave3.
% 		I. Orientador.
% 		II. Universidade xxx.
% 		III. Faculdade de xxx.
% 		IV. Título 			
% 	\end{minipage}}
% 	\end{center}
% \end{fichacatalografica}
% ---


% ---
% FOLHA DE APROVAÇÃO
% ---

% Isto é um exemplo de Folha de aprovação, elemento obrigatório da NBR
% 14724/2011 (seção 4.2.1.3). Você pode utilizar este modelo até a aprovação
% do trabalho. Após isso, substitua todo o conteúdo deste arquivo por uma
% imagem da página assinada pela banca com o comando abaixo:
%
% \begin{folhadeaprovacao}
% \includepdf{folhadeaprovacao_final.pdf}
% \end{folhadeaprovacao}
%
\setlength{\ABNTEXsignskip}{2cm}
\setlength{\ABNTEXsignwidth}{10cm}

\begin{folhadeaprovacao}
  \begin{center}
    {\ABNTEXchapterfont\large\imprimirautor}

    \vspace*{1cm}

    \begin{center}
      \ABNTEXchapterfont\bfseries\Large\imprimirtitulo
    \end{center}

    \vspace*{1cm}
    
    \hspace{.45\textwidth}
    \begin{minipage}{.5\textwidth}
        \imprimirpreambulo
    \end{minipage}%
   \end{center}

   \vspace*{1cm}
        
   \begin{center}
   Aprovado em 17 de julho de 2024.

   \vspace*{1cm}

   Banca examinadora:
   \end{center}

   \vspace*{-1cm}

   \assinatura{Prof. Ma. Caroline Santa Maria Rodrigues -- PUCRS}
   \assinatura{Prof. Dra. Clarissa Pinto Pizarro de Freitas -- PUCRS}
   \assinatura{\imprimirorientador -- PUCRS}
\end{folhadeaprovacao}
% ---


% ---
% Agradecimentos
% ---
\begin{agradecimentos}
Agradeço aos meus pais, Ademilson José Dias e Irinês Porfírio Dias, pelo apoio incondicional que me deram ao longo da vida,
incentivando mesmo as ideias questionáveis como aquela de fazer uma segunda graduação.

A minha irmã, Maria Eduarda Dias, cuja trajetória acadêmica é fonte de inspiração.

A minha namorada, Nathalia Fallavena Ceratti, que acompanhou a esta jornada de perto, oferecendo apoio e compreensão que
fizeram toda a diferença.

Aos meus amigos, Guilherme da Silva Petró, Norton Neves e Uilinton Luciano Marques da Silva, que foram compreensíveis nos
finais de semana em que não pude jogar RPG.

Ao professor Wagner de Lara Machado pela orientação e apoio no desenvolvimento deste trabalho.
\end{agradecimentos}
% ---


% ---
% RESUMOS
% ---
% Resumo em português
\setlength{\absparsep}{18pt} % ajusta o espaçamento dos parágrafos do resumo
\begin{resumo}
Desenvolvimentos na área de aprendizagem de máquina representam uma série de oportunidades para diversas áreas de pesquisa e de atuação profissional.
O uso de modelos de aprendizagem de máquina na predição de desfecho para tratamentos em saúde mental tem sido investigado e apresenta resultados promissores.
A aplicação desse tipo de modelo possibilitaria, por exemplo, maior assertividade na personalização de tratamentos psicoterápicos a nível individual.
Apesar dos benefícios potenciais, o tema ainda é pouco discutido nos espaços de pesquisa, formação e atuação em psicologia clínica. Este trabalho tem
por objetivo introduzir o tema de aprendizagem de máquina a psicólogos. Busca-se apresentar, de maneira acessível, os conceitos e terminologia básicos
em aprendizagem de máquina, proporcionando um panorama geral da área. Para complementar a exposição conceitual, é construído um exemplo a partir de dados
de um estudo sobre eficácia de intervenções em psicologia positiva para depressão via internet. Por fim, discutem-se algumas implicações do uso desse tipo
de tecnologia para a pesquisa e para a prática em psicoterapia.

 \textbf{Palavras-chave}: aprendizagem de máquina, psicoterapia, psicologia positiva, depressão.
\end{resumo}

% Resumo em inglês
\begin{resumo}[Abstract]
 \begin{otherlanguage*}{english}
   Advances in machine learning have opened up a multitude of opportunities for various research and professional fields. The use of machine learning models to
   predict treatment outcomes in mental health has been investigated showing promising results. The application of such models would allow greater assertiveness
   in tailoring psychotherapeutic treatments at the individual level. Despite its potential benefits, the topic remains unexplored in the fields of research, training
   and practice in clinical psychology. This work aims to introduce the topic of machine learning to psychologists. It presents the basic machine learning concepts and
   terminology in an accessible manner, providing a comprehensive overview of the field. To complement the conceptual discussion, an example is constructed using data
   from a study on the effectiveness of internet-based positive psychology interventions for depression. Finally, implications of using this type of technology in
   psychotherapy research and practice are discussed.
   \vspace{\onelineskip}
 
   \noindent 
   \textbf{Keywords}: machine learning, psychotherapy, positive psychology, depression.
 \end{otherlanguage*}
\end{resumo}


% ---
% LISTA DE ILUSTRAÇÕES
% ---
\pdfbookmark[0]{\listfigurename}{lof}
\listoffigures*
\cleardoublepage
% ---


% ---
% LISTA DE QUADROS
% ---
% \pdfbookmark[0]{\listofquadrosname}{loq}
% \listofquadros*
% \cleardoublepage
% ---


% ---
% LISTA DE TABELAS
% ---
% \pdfbookmark[0]{\listtablename}{lot}
% \listoftables*
% \cleardoublepage
% ---


% ---
% LISTA DE ABREVIATURAS E SIGLAS
% ---
% \begin{siglas}
%   \item[ABNT] Associação Brasileira de Normas Técnicas
%   \item[abnTeX] ABsurdas Normas para TeX
% \end{siglas}
% ---


% ---
% LISTA DE SÍMBOLOS
% ---
% \begin{simbolos}
%   \item[$ \Gamma $] Letra grega Gama
%   \item[$ \Lambda $] Lambda
%   \item[$ \zeta $] Letra grega minúscula zeta
%   \item[$ \in $] Pertence
% \end{simbolos}
% ---


% ---
% Sumário
% ---
\pdfbookmark[0]{\contentsname}{toc}
\tableofcontents*
\cleardoublepage
% ---


% ----------------------------------------------------------
% ELEMENTOS TEXTUAIS
% ----------------------------------------------------------
\textual

\include{./01-introducao/main.tex}

\include{./02-machine-learning/main.tex}

\include{./03-exemplo/main.tex}

\include{./04-conclusao/main.tex}

% ----------------------------------------------------------
% ELEMENTOS PÓS-TEXTUAIS
% ----------------------------------------------------------
\postextual
% ----------------------------------------------------------


% ----------------------------------------------------------
% REFERÊNCIAS
% ----------------------------------------------------------
\bibliography{references.bib}


% ----------------------------------------------------------
% GLOSSÁRIO
% ----------------------------------------------------------
%
% Consulte o manual da classe abntex2 para orientações sobre o glossário.
%
%\glossary

% ----------------------------------------------------------
% APÊNDICES
% ----------------------------------------------------------

% ---
% Inicia os apêndices
% ---
% \begin{apendicesenv}

% % Imprime uma página indicando o início dos apêndices
% \partapendices

% % ----------------------------------------------------------
% \chapter{Quisque libero justo}
% % ----------------------------------------------------------

% % ----------------------------------------------------------
% \chapter{Nullam elementum urna}
% % ----------------------------------------------------------

% \end{apendicesenv}
% ---


% ----------------------------------------------------------
% ANEXOS
% ----------------------------------------------------------

% ---
% Inicia os anexos
% ---
% \begin{anexosenv}

% % Imprime uma página indicando o início dos anexos
% \partanexos

% % ---
% \chapter{Morbi ultrices rutrum lorem}
% % ---

% % ---
% \chapter{Cras non urna sed feugiat cum}
% % ---


% % ---
% \chapter{Fusce facilisis lacinia dui}
% % ---

% \end{anexosenv}

%---------------------------------------------------------------------
% INDICE REMISSIVO
%---------------------------------------------------------------------
% \phantompart
% \printindex
%---------------------------------------------------------------------
\includepdf{./05-extras/contracapa.pdf}
\end{document}


\documentclass[
	% Opções da classe memoir
	12pt,				        % tamanho da fonte
	openany,			        % capítulos começam em pág ímpar (insere página vazia caso preciso)
	oneside,			        % para impressão em recto e verso. Oposto a oneside
	a4paper,			        % tamanho do papel. 
	% Opções da classe abntex2
	%chapter=TITLE,		        % títulos de capítulos convertidos em letras maiúsculas
	%section=TITLE,		        % títulos de seções convertidos em letras maiúsculas
	%subsection=TITLE,	        % títulos de subseções convertidos em letras maiúsculas
	%subsubsection=TITLE,       % títulos de subsubseções convertidos em letras maiúsculas
	% Opções do pacote babel
	english,			        % idioma adicional para hifenização
	brazil				        % o último idioma é o principal do documento
	]{abntex2}

% ---
% Pacotes básicos 
% ---
\usepackage{lmodern}			% Usa a fonte Latin Modern			
\usepackage[T1]{fontenc}		% Selecao de codigos de fonte.
\usepackage[utf8]{inputenc}		% Codificacao do documento (conversão automática dos acentos)
\usepackage{indentfirst}		% Indenta o primeiro parágrafo de cada seção.
\usepackage{color}				% Controle das cores
\usepackage{graphicx}			% Inclusão de gráficos
\usepackage{microtype} 			% para melhorias de justificação
% ---
		

% ---
% Pacotes de citações
% ---
% \usepackage[alf]{abntex2cite}						% Citações padrão ABNT
\usepackage[
  alf,
  abnt-nbr10520=1988,                     % Citações em sentence case
  abnt-emphasize=bf
  ]{abntex2cite}


% ---
% Pacotes adicionais
% ---
\usepackage{pdfpages}	% Inclusão de páginas em PDF


% ---
% Informações para CAPA e FOLHA DE ROSTO
% ---
\titulo{APRENDIZAGEM DE MÁQUINA PARA PREDIÇÃO DE DESFECHOS EM SAÚDE MENTAL: \textmd{UMA INTRODUÇÃO PARA PSICÓLOGOS}}
\autor{GUILHERME AUGUSTO DIAS}
\local{Porto Alegre}
\data{2024}
\orientador{Prof. Dr. Wagner de Lara Machado}
% \coorientador{}
% \instituicao{
%   Pontifícia Universidade Católica do Rio Grande do Sul
%   \par
%   Escola de Ciências da Saúde e da Vida
%   \par
%   Curso de Psicologia}
\tipotrabalho{TCC}
\preambulo{Trabalho de Conclusão de Curso apresentado como requisito para a obtenção do grau de Bacharel em Psicologia pela Escola de Ciências da Saúde e da Vida da Pontifícia Universidade Católica do Rio Grande do Sul.}
% ---

% ---
% Configurações de aparência do PDF final
% ---

% Alterando o aspecto da cor azul
\definecolor{blue}{RGB}{41,5,195}

% Informações do PDF
\makeatletter
\hypersetup{
		pdftitle={Aprendizagem de máquina para predição de desfechos em saúde mental: uma introdução para psicólogos}, 
		pdfauthor={Guilherme Augusto Dias},
    	pdfsubject={\imprimirpreambulo},
	    pdfcreator={LaTeX with abnTeX2},
		pdfkeywords={aprendizagem de máquina}{psicoterapia}{psicologia positiva}{depressão}, 
		colorlinks=true,		% true: colored links; false: boxed links
    	linkcolor=black,		% color of internal links
    	citecolor=black,		% color of links to bibliography
    	filecolor=black,		% color of file links
		urlcolor=black,
		bookmarksdepth=4
}
\makeatother
% --- 


% ---
% Posiciona figuras e tabelas no topo da página quando adicionadas sozinhas
% em um página em branco. Ver https://github.com/abntex/abntex2/issues/170
\makeatletter
\setlength{\@fptop}{5pt} % Set distance from top of page to first float
\makeatother
% ---

% ---
% Possibilita criação de Quadros e Lista de quadros.
% Ver https://github.com/abntex/abntex2/issues/176
%
\newcommand{\quadroname}{Quadro}
\newcommand{\listofquadrosname}{Lista de quadros}

\newfloat[chapter]{quadro}{loq}{\quadroname}
\newlistof{listofquadros}{loq}{\listofquadrosname}
\newlistentry{quadro}{loq}{0}

% Configurações para atender às regras da ABNT
\setfloatadjustment{quadro}{\centering}
\counterwithout{quadro}{chapter}
\renewcommand{\cftquadroname}{\quadroname\space} 
\renewcommand*{\cftquadroaftersnum}{\hfill--\hfill}

\setfloatlocations{quadro}{hbtp} % Ver https://github.com/abntex/abntex2/issues/176
% ---


% --- 
% Espaçamentos entre linhas e parágrafos 
% --- 

% O tamanho do parágrafo é dado por:
\setlength{\parindent}{1.3cm}

% Controle do espaçamento entre um parágrafo e outro:
\setlength{\parskip}{0.2cm}  % tente também \onelineskip

% ---
% Compila o índice
% ---
\makeindex
% ---


% ----
% Início do documento
% ----
\begin{document}

% Seleciona o idioma do documento (conforme pacotes do babel)
\selectlanguage{brazil}

% Retira espaço extra obsoleto entre as frases.
\frenchspacing 

% ----------------------------------------------------------
% ELEMENTOS PRÉ-TEXTUAIS
% ----------------------------------------------------------
\pretextual

% ---
% CAPA
% ---
\includepdf{./05-extras/capa.pdf}
% ---


% ---
% FOLHA DE ROSTO
% (o * indica que haverá a ficha bibliográfica)
% ---
\imprimirfolhaderosto
% ---


% ---
% Ficha bibliografica
% ---

% Isto é um exemplo de Ficha Catalográfica, ou ``Dados internacionais de
% catalogação-na-publicação''. Você pode utilizar este modelo como referência. 
% Porém, provavelmente a biblioteca da sua universidade lhe fornecerá um PDF
% com a ficha catalográfica definitiva após a defesa do trabalho. Quando estiver
% com o documento, salve-o como PDF no diretório do seu projeto e substitua todo
% o conteúdo de implementação deste arquivo pelo comando abaixo:
%
% \usepackage{pdfpages}		% necessário para comando \includepdf
% \begin{fichacatalografica}
%     \includepdf{fig_ficha_catalografica.pdf}
% \end{fichacatalografica}

% \begin{fichacatalografica}
% 	\sffamily
% 	\vspace*{\fill}                             % Posição vertical
% 	\begin{center}                              % Minipage Centralizado
% 	\fbox{\begin{minipage}[c][8cm]{13.5cm}		% Largura
% 	\small
% 	\imprimirautor
% 	%Sobrenome, Nome do autor
	
% 	\hspace{0.5cm} \imprimirtitulo  / \imprimirautor. --
% 	\imprimirlocal, \imprimirdata-
	
% 	\hspace{0.5cm} \thelastpage p. : il. (algumas color.) ; 30 cm.\\
	
% 	\hspace{0.5cm} \imprimirorientadorRotulo~\imprimirorientador\\
	
% 	\hspace{0.5cm}
% 	\parbox[t]{\textwidth}{\imprimirtipotrabalho~--~\imprimirinstituicao,
% 	\imprimirdata.}\\
	
% 	\hspace{0.5cm}
% 		1. Palavra-chave1.
% 		2. Palavra-chave2.
% 		2. Palavra-chave3.
% 		I. Orientador.
% 		II. Universidade xxx.
% 		III. Faculdade de xxx.
% 		IV. Título 			
% 	\end{minipage}}
% 	\end{center}
% \end{fichacatalografica}
% ---


% ---
% FOLHA DE APROVAÇÃO
% ---

% Isto é um exemplo de Folha de aprovação, elemento obrigatório da NBR
% 14724/2011 (seção 4.2.1.3). Você pode utilizar este modelo até a aprovação
% do trabalho. Após isso, substitua todo o conteúdo deste arquivo por uma
% imagem da página assinada pela banca com o comando abaixo:
%
% \begin{folhadeaprovacao}
% \includepdf{folhadeaprovacao_final.pdf}
% \end{folhadeaprovacao}
%
\setlength{\ABNTEXsignskip}{2cm}
\setlength{\ABNTEXsignwidth}{10cm}

\begin{folhadeaprovacao}
  \begin{center}
    {\ABNTEXchapterfont\large\imprimirautor}

    \vspace*{1cm}

    \begin{center}
      \ABNTEXchapterfont\bfseries\Large\imprimirtitulo
    \end{center}

    \vspace*{1cm}
    
    \hspace{.45\textwidth}
    \begin{minipage}{.5\textwidth}
        \imprimirpreambulo
    \end{minipage}%
   \end{center}

   \vspace*{1cm}
        
   \begin{center}
   Aprovado em 17 de julho de 2024.

   \vspace*{1cm}

   Banca examinadora:
   \end{center}

   \vspace*{-1cm}

   \assinatura{Prof. Ma. Caroline Santa Maria Rodrigues -- PUCRS}
   \assinatura{Prof. Dra. Clarissa Pinto Pizarro de Freitas -- PUCRS}
   \assinatura{\imprimirorientador -- PUCRS}
\end{folhadeaprovacao}
% ---


% ---
% Agradecimentos
% ---
\begin{agradecimentos}
Agradeço aos meus pais, Ademilson José Dias e Irinês Porfírio Dias, pelo apoio incondicional que me deram ao longo da vida,
incentivando mesmo as ideias questionáveis como aquela de fazer uma segunda graduação.

A minha irmã, Maria Eduarda Dias, cuja trajetória acadêmica é fonte de inspiração.

A minha namorada, Nathalia Fallavena Ceratti, que acompanhou a esta jornada de perto, oferecendo apoio e compreensão que
fizeram toda a diferença.

Aos meus amigos, Guilherme da Silva Petró, Norton Neves e Uilinton Luciano Marques da Silva, que foram compreensíveis nos
finais de semana em que não pude jogar RPG.

Ao professor Wagner de Lara Machado pela orientação e apoio no desenvolvimento deste trabalho.
\end{agradecimentos}
% ---


% ---
% RESUMOS
% ---
% Resumo em português
\setlength{\absparsep}{18pt} % ajusta o espaçamento dos parágrafos do resumo
\begin{resumo}
Desenvolvimentos na área de aprendizagem de máquina representam uma série de oportunidades para diversas áreas de pesquisa e de atuação profissional.
O uso de modelos de aprendizagem de máquina na predição de desfecho para tratamentos em saúde mental tem sido investigado e apresenta resultados promissores.
A aplicação desse tipo de modelo possibilitaria, por exemplo, maior assertividade na personalização de tratamentos psicoterápicos a nível individual.
Apesar dos benefícios potenciais, o tema ainda é pouco discutido nos espaços de pesquisa, formação e atuação em psicologia clínica. Este trabalho tem
por objetivo introduzir o tema de aprendizagem de máquina a psicólogos. Busca-se apresentar, de maneira acessível, os conceitos e terminologia básicos
em aprendizagem de máquina, proporcionando um panorama geral da área. Para complementar a exposição conceitual, é construído um exemplo a partir de dados
de um estudo sobre eficácia de intervenções em psicologia positiva para depressão via internet. Por fim, discutem-se algumas implicações do uso desse tipo
de tecnologia para a pesquisa e para a prática em psicoterapia.

 \textbf{Palavras-chave}: aprendizagem de máquina, psicoterapia, psicologia positiva, depressão.
\end{resumo}

% Resumo em inglês
\begin{resumo}[Abstract]
 \begin{otherlanguage*}{english}
   Advances in machine learning have opened up a multitude of opportunities for various research and professional fields. The use of machine learning models to
   predict treatment outcomes in mental health has been investigated showing promising results. The application of such models would allow greater assertiveness
   in tailoring psychotherapeutic treatments at the individual level. Despite its potential benefits, the topic remains unexplored in the fields of research, training
   and practice in clinical psychology. This work aims to introduce the topic of machine learning to psychologists. It presents the basic machine learning concepts and
   terminology in an accessible manner, providing a comprehensive overview of the field. To complement the conceptual discussion, an example is constructed using data
   from a study on the effectiveness of internet-based positive psychology interventions for depression. Finally, implications of using this type of technology in
   psychotherapy research and practice are discussed.
   \vspace{\onelineskip}
 
   \noindent 
   \textbf{Keywords}: machine learning, psychotherapy, positive psychology, depression.
 \end{otherlanguage*}
\end{resumo}


% ---
% LISTA DE ILUSTRAÇÕES
% ---
\pdfbookmark[0]{\listfigurename}{lof}
\listoffigures*
\cleardoublepage
% ---


% ---
% LISTA DE QUADROS
% ---
% \pdfbookmark[0]{\listofquadrosname}{loq}
% \listofquadros*
% \cleardoublepage
% ---


% ---
% LISTA DE TABELAS
% ---
% \pdfbookmark[0]{\listtablename}{lot}
% \listoftables*
% \cleardoublepage
% ---


% ---
% LISTA DE ABREVIATURAS E SIGLAS
% ---
% \begin{siglas}
%   \item[ABNT] Associação Brasileira de Normas Técnicas
%   \item[abnTeX] ABsurdas Normas para TeX
% \end{siglas}
% ---


% ---
% LISTA DE SÍMBOLOS
% ---
% \begin{simbolos}
%   \item[$ \Gamma $] Letra grega Gama
%   \item[$ \Lambda $] Lambda
%   \item[$ \zeta $] Letra grega minúscula zeta
%   \item[$ \in $] Pertence
% \end{simbolos}
% ---


% ---
% Sumário
% ---
\pdfbookmark[0]{\contentsname}{toc}
\tableofcontents*
\cleardoublepage
% ---


% ----------------------------------------------------------
% ELEMENTOS TEXTUAIS
% ----------------------------------------------------------
\textual

\include{./01-introducao/main.tex}

\include{./02-machine-learning/main.tex}

\include{./03-exemplo/main.tex}

\include{./04-conclusao/main.tex}

% ----------------------------------------------------------
% ELEMENTOS PÓS-TEXTUAIS
% ----------------------------------------------------------
\postextual
% ----------------------------------------------------------


% ----------------------------------------------------------
% REFERÊNCIAS
% ----------------------------------------------------------
\bibliography{references.bib}


% ----------------------------------------------------------
% GLOSSÁRIO
% ----------------------------------------------------------
%
% Consulte o manual da classe abntex2 para orientações sobre o glossário.
%
%\glossary

% ----------------------------------------------------------
% APÊNDICES
% ----------------------------------------------------------

% ---
% Inicia os apêndices
% ---
% \begin{apendicesenv}

% % Imprime uma página indicando o início dos apêndices
% \partapendices

% % ----------------------------------------------------------
% \chapter{Quisque libero justo}
% % ----------------------------------------------------------

% % ----------------------------------------------------------
% \chapter{Nullam elementum urna}
% % ----------------------------------------------------------

% \end{apendicesenv}
% ---


% ----------------------------------------------------------
% ANEXOS
% ----------------------------------------------------------

% ---
% Inicia os anexos
% ---
% \begin{anexosenv}

% % Imprime uma página indicando o início dos anexos
% \partanexos

% % ---
% \chapter{Morbi ultrices rutrum lorem}
% % ---

% % ---
% \chapter{Cras non urna sed feugiat cum}
% % ---


% % ---
% \chapter{Fusce facilisis lacinia dui}
% % ---

% \end{anexosenv}

%---------------------------------------------------------------------
% INDICE REMISSIVO
%---------------------------------------------------------------------
% \phantompart
% \printindex
%---------------------------------------------------------------------
\includepdf{./05-extras/contracapa.pdf}
\end{document}


\documentclass[
	% Opções da classe memoir
	12pt,				        % tamanho da fonte
	openany,			        % capítulos começam em pág ímpar (insere página vazia caso preciso)
	oneside,			        % para impressão em recto e verso. Oposto a oneside
	a4paper,			        % tamanho do papel. 
	% Opções da classe abntex2
	%chapter=TITLE,		        % títulos de capítulos convertidos em letras maiúsculas
	%section=TITLE,		        % títulos de seções convertidos em letras maiúsculas
	%subsection=TITLE,	        % títulos de subseções convertidos em letras maiúsculas
	%subsubsection=TITLE,       % títulos de subsubseções convertidos em letras maiúsculas
	% Opções do pacote babel
	english,			        % idioma adicional para hifenização
	brazil				        % o último idioma é o principal do documento
	]{abntex2}

% ---
% Pacotes básicos 
% ---
\usepackage{lmodern}			% Usa a fonte Latin Modern			
\usepackage[T1]{fontenc}		% Selecao de codigos de fonte.
\usepackage[utf8]{inputenc}		% Codificacao do documento (conversão automática dos acentos)
\usepackage{indentfirst}		% Indenta o primeiro parágrafo de cada seção.
\usepackage{color}				% Controle das cores
\usepackage{graphicx}			% Inclusão de gráficos
\usepackage{microtype} 			% para melhorias de justificação
% ---
		

% ---
% Pacotes de citações
% ---
% \usepackage[alf]{abntex2cite}						% Citações padrão ABNT
\usepackage[
  alf,
  abnt-nbr10520=1988,                     % Citações em sentence case
  abnt-emphasize=bf
  ]{abntex2cite}


% ---
% Pacotes adicionais
% ---
\usepackage{pdfpages}	% Inclusão de páginas em PDF


% ---
% Informações para CAPA e FOLHA DE ROSTO
% ---
\titulo{APRENDIZAGEM DE MÁQUINA PARA PREDIÇÃO DE DESFECHOS EM SAÚDE MENTAL: \textmd{UMA INTRODUÇÃO PARA PSICÓLOGOS}}
\autor{GUILHERME AUGUSTO DIAS}
\local{Porto Alegre}
\data{2024}
\orientador{Prof. Dr. Wagner de Lara Machado}
% \coorientador{}
% \instituicao{
%   Pontifícia Universidade Católica do Rio Grande do Sul
%   \par
%   Escola de Ciências da Saúde e da Vida
%   \par
%   Curso de Psicologia}
\tipotrabalho{TCC}
\preambulo{Trabalho de Conclusão de Curso apresentado como requisito para a obtenção do grau de Bacharel em Psicologia pela Escola de Ciências da Saúde e da Vida da Pontifícia Universidade Católica do Rio Grande do Sul.}
% ---

% ---
% Configurações de aparência do PDF final
% ---

% Alterando o aspecto da cor azul
\definecolor{blue}{RGB}{41,5,195}

% Informações do PDF
\makeatletter
\hypersetup{
		pdftitle={Aprendizagem de máquina para predição de desfechos em saúde mental: uma introdução para psicólogos}, 
		pdfauthor={Guilherme Augusto Dias},
    	pdfsubject={\imprimirpreambulo},
	    pdfcreator={LaTeX with abnTeX2},
		pdfkeywords={aprendizagem de máquina}{psicoterapia}{psicologia positiva}{depressão}, 
		colorlinks=true,		% true: colored links; false: boxed links
    	linkcolor=black,		% color of internal links
    	citecolor=black,		% color of links to bibliography
    	filecolor=black,		% color of file links
		urlcolor=black,
		bookmarksdepth=4
}
\makeatother
% --- 


% ---
% Posiciona figuras e tabelas no topo da página quando adicionadas sozinhas
% em um página em branco. Ver https://github.com/abntex/abntex2/issues/170
\makeatletter
\setlength{\@fptop}{5pt} % Set distance from top of page to first float
\makeatother
% ---

% ---
% Possibilita criação de Quadros e Lista de quadros.
% Ver https://github.com/abntex/abntex2/issues/176
%
\newcommand{\quadroname}{Quadro}
\newcommand{\listofquadrosname}{Lista de quadros}

\newfloat[chapter]{quadro}{loq}{\quadroname}
\newlistof{listofquadros}{loq}{\listofquadrosname}
\newlistentry{quadro}{loq}{0}

% Configurações para atender às regras da ABNT
\setfloatadjustment{quadro}{\centering}
\counterwithout{quadro}{chapter}
\renewcommand{\cftquadroname}{\quadroname\space} 
\renewcommand*{\cftquadroaftersnum}{\hfill--\hfill}

\setfloatlocations{quadro}{hbtp} % Ver https://github.com/abntex/abntex2/issues/176
% ---


% --- 
% Espaçamentos entre linhas e parágrafos 
% --- 

% O tamanho do parágrafo é dado por:
\setlength{\parindent}{1.3cm}

% Controle do espaçamento entre um parágrafo e outro:
\setlength{\parskip}{0.2cm}  % tente também \onelineskip

% ---
% Compila o índice
% ---
\makeindex
% ---


% ----
% Início do documento
% ----
\begin{document}

% Seleciona o idioma do documento (conforme pacotes do babel)
\selectlanguage{brazil}

% Retira espaço extra obsoleto entre as frases.
\frenchspacing 

% ----------------------------------------------------------
% ELEMENTOS PRÉ-TEXTUAIS
% ----------------------------------------------------------
\pretextual

% ---
% CAPA
% ---
\includepdf{./05-extras/capa.pdf}
% ---


% ---
% FOLHA DE ROSTO
% (o * indica que haverá a ficha bibliográfica)
% ---
\imprimirfolhaderosto
% ---


% ---
% Ficha bibliografica
% ---

% Isto é um exemplo de Ficha Catalográfica, ou ``Dados internacionais de
% catalogação-na-publicação''. Você pode utilizar este modelo como referência. 
% Porém, provavelmente a biblioteca da sua universidade lhe fornecerá um PDF
% com a ficha catalográfica definitiva após a defesa do trabalho. Quando estiver
% com o documento, salve-o como PDF no diretório do seu projeto e substitua todo
% o conteúdo de implementação deste arquivo pelo comando abaixo:
%
% \usepackage{pdfpages}		% necessário para comando \includepdf
% \begin{fichacatalografica}
%     \includepdf{fig_ficha_catalografica.pdf}
% \end{fichacatalografica}

% \begin{fichacatalografica}
% 	\sffamily
% 	\vspace*{\fill}                             % Posição vertical
% 	\begin{center}                              % Minipage Centralizado
% 	\fbox{\begin{minipage}[c][8cm]{13.5cm}		% Largura
% 	\small
% 	\imprimirautor
% 	%Sobrenome, Nome do autor
	
% 	\hspace{0.5cm} \imprimirtitulo  / \imprimirautor. --
% 	\imprimirlocal, \imprimirdata-
	
% 	\hspace{0.5cm} \thelastpage p. : il. (algumas color.) ; 30 cm.\\
	
% 	\hspace{0.5cm} \imprimirorientadorRotulo~\imprimirorientador\\
	
% 	\hspace{0.5cm}
% 	\parbox[t]{\textwidth}{\imprimirtipotrabalho~--~\imprimirinstituicao,
% 	\imprimirdata.}\\
	
% 	\hspace{0.5cm}
% 		1. Palavra-chave1.
% 		2. Palavra-chave2.
% 		2. Palavra-chave3.
% 		I. Orientador.
% 		II. Universidade xxx.
% 		III. Faculdade de xxx.
% 		IV. Título 			
% 	\end{minipage}}
% 	\end{center}
% \end{fichacatalografica}
% ---


% ---
% FOLHA DE APROVAÇÃO
% ---

% Isto é um exemplo de Folha de aprovação, elemento obrigatório da NBR
% 14724/2011 (seção 4.2.1.3). Você pode utilizar este modelo até a aprovação
% do trabalho. Após isso, substitua todo o conteúdo deste arquivo por uma
% imagem da página assinada pela banca com o comando abaixo:
%
% \begin{folhadeaprovacao}
% \includepdf{folhadeaprovacao_final.pdf}
% \end{folhadeaprovacao}
%
\setlength{\ABNTEXsignskip}{2cm}
\setlength{\ABNTEXsignwidth}{10cm}

\begin{folhadeaprovacao}
  \begin{center}
    {\ABNTEXchapterfont\large\imprimirautor}

    \vspace*{1cm}

    \begin{center}
      \ABNTEXchapterfont\bfseries\Large\imprimirtitulo
    \end{center}

    \vspace*{1cm}
    
    \hspace{.45\textwidth}
    \begin{minipage}{.5\textwidth}
        \imprimirpreambulo
    \end{minipage}%
   \end{center}

   \vspace*{1cm}
        
   \begin{center}
   Aprovado em 17 de julho de 2024.

   \vspace*{1cm}

   Banca examinadora:
   \end{center}

   \vspace*{-1cm}

   \assinatura{Prof. Ma. Caroline Santa Maria Rodrigues -- PUCRS}
   \assinatura{Prof. Dra. Clarissa Pinto Pizarro de Freitas -- PUCRS}
   \assinatura{\imprimirorientador -- PUCRS}
\end{folhadeaprovacao}
% ---


% ---
% Agradecimentos
% ---
\begin{agradecimentos}
Agradeço aos meus pais, Ademilson José Dias e Irinês Porfírio Dias, pelo apoio incondicional que me deram ao longo da vida,
incentivando mesmo as ideias questionáveis como aquela de fazer uma segunda graduação.

A minha irmã, Maria Eduarda Dias, cuja trajetória acadêmica é fonte de inspiração.

A minha namorada, Nathalia Fallavena Ceratti, que acompanhou a esta jornada de perto, oferecendo apoio e compreensão que
fizeram toda a diferença.

Aos meus amigos, Guilherme da Silva Petró, Norton Neves e Uilinton Luciano Marques da Silva, que foram compreensíveis nos
finais de semana em que não pude jogar RPG.

Ao professor Wagner de Lara Machado pela orientação e apoio no desenvolvimento deste trabalho.
\end{agradecimentos}
% ---


% ---
% RESUMOS
% ---
% Resumo em português
\setlength{\absparsep}{18pt} % ajusta o espaçamento dos parágrafos do resumo
\begin{resumo}
Desenvolvimentos na área de aprendizagem de máquina representam uma série de oportunidades para diversas áreas de pesquisa e de atuação profissional.
O uso de modelos de aprendizagem de máquina na predição de desfecho para tratamentos em saúde mental tem sido investigado e apresenta resultados promissores.
A aplicação desse tipo de modelo possibilitaria, por exemplo, maior assertividade na personalização de tratamentos psicoterápicos a nível individual.
Apesar dos benefícios potenciais, o tema ainda é pouco discutido nos espaços de pesquisa, formação e atuação em psicologia clínica. Este trabalho tem
por objetivo introduzir o tema de aprendizagem de máquina a psicólogos. Busca-se apresentar, de maneira acessível, os conceitos e terminologia básicos
em aprendizagem de máquina, proporcionando um panorama geral da área. Para complementar a exposição conceitual, é construído um exemplo a partir de dados
de um estudo sobre eficácia de intervenções em psicologia positiva para depressão via internet. Por fim, discutem-se algumas implicações do uso desse tipo
de tecnologia para a pesquisa e para a prática em psicoterapia.

 \textbf{Palavras-chave}: aprendizagem de máquina, psicoterapia, psicologia positiva, depressão.
\end{resumo}

% Resumo em inglês
\begin{resumo}[Abstract]
 \begin{otherlanguage*}{english}
   Advances in machine learning have opened up a multitude of opportunities for various research and professional fields. The use of machine learning models to
   predict treatment outcomes in mental health has been investigated showing promising results. The application of such models would allow greater assertiveness
   in tailoring psychotherapeutic treatments at the individual level. Despite its potential benefits, the topic remains unexplored in the fields of research, training
   and practice in clinical psychology. This work aims to introduce the topic of machine learning to psychologists. It presents the basic machine learning concepts and
   terminology in an accessible manner, providing a comprehensive overview of the field. To complement the conceptual discussion, an example is constructed using data
   from a study on the effectiveness of internet-based positive psychology interventions for depression. Finally, implications of using this type of technology in
   psychotherapy research and practice are discussed.
   \vspace{\onelineskip}
 
   \noindent 
   \textbf{Keywords}: machine learning, psychotherapy, positive psychology, depression.
 \end{otherlanguage*}
\end{resumo}


% ---
% LISTA DE ILUSTRAÇÕES
% ---
\pdfbookmark[0]{\listfigurename}{lof}
\listoffigures*
\cleardoublepage
% ---


% ---
% LISTA DE QUADROS
% ---
% \pdfbookmark[0]{\listofquadrosname}{loq}
% \listofquadros*
% \cleardoublepage
% ---


% ---
% LISTA DE TABELAS
% ---
% \pdfbookmark[0]{\listtablename}{lot}
% \listoftables*
% \cleardoublepage
% ---


% ---
% LISTA DE ABREVIATURAS E SIGLAS
% ---
% \begin{siglas}
%   \item[ABNT] Associação Brasileira de Normas Técnicas
%   \item[abnTeX] ABsurdas Normas para TeX
% \end{siglas}
% ---


% ---
% LISTA DE SÍMBOLOS
% ---
% \begin{simbolos}
%   \item[$ \Gamma $] Letra grega Gama
%   \item[$ \Lambda $] Lambda
%   \item[$ \zeta $] Letra grega minúscula zeta
%   \item[$ \in $] Pertence
% \end{simbolos}
% ---


% ---
% Sumário
% ---
\pdfbookmark[0]{\contentsname}{toc}
\tableofcontents*
\cleardoublepage
% ---


% ----------------------------------------------------------
% ELEMENTOS TEXTUAIS
% ----------------------------------------------------------
\textual

\include{./01-introducao/main.tex}

\include{./02-machine-learning/main.tex}

\include{./03-exemplo/main.tex}

\include{./04-conclusao/main.tex}

% ----------------------------------------------------------
% ELEMENTOS PÓS-TEXTUAIS
% ----------------------------------------------------------
\postextual
% ----------------------------------------------------------


% ----------------------------------------------------------
% REFERÊNCIAS
% ----------------------------------------------------------
\bibliography{references.bib}


% ----------------------------------------------------------
% GLOSSÁRIO
% ----------------------------------------------------------
%
% Consulte o manual da classe abntex2 para orientações sobre o glossário.
%
%\glossary

% ----------------------------------------------------------
% APÊNDICES
% ----------------------------------------------------------

% ---
% Inicia os apêndices
% ---
% \begin{apendicesenv}

% % Imprime uma página indicando o início dos apêndices
% \partapendices

% % ----------------------------------------------------------
% \chapter{Quisque libero justo}
% % ----------------------------------------------------------

% % ----------------------------------------------------------
% \chapter{Nullam elementum urna}
% % ----------------------------------------------------------

% \end{apendicesenv}
% ---


% ----------------------------------------------------------
% ANEXOS
% ----------------------------------------------------------

% ---
% Inicia os anexos
% ---
% \begin{anexosenv}

% % Imprime uma página indicando o início dos anexos
% \partanexos

% % ---
% \chapter{Morbi ultrices rutrum lorem}
% % ---

% % ---
% \chapter{Cras non urna sed feugiat cum}
% % ---


% % ---
% \chapter{Fusce facilisis lacinia dui}
% % ---

% \end{anexosenv}

%---------------------------------------------------------------------
% INDICE REMISSIVO
%---------------------------------------------------------------------
% \phantompart
% \printindex
%---------------------------------------------------------------------
\includepdf{./05-extras/contracapa.pdf}
\end{document}


% ----------------------------------------------------------
% ELEMENTOS PÓS-TEXTUAIS
% ----------------------------------------------------------
\postextual
% ----------------------------------------------------------


% ----------------------------------------------------------
% REFERÊNCIAS
% ----------------------------------------------------------
\bibliography{references.bib}


% ----------------------------------------------------------
% GLOSSÁRIO
% ----------------------------------------------------------
%
% Consulte o manual da classe abntex2 para orientações sobre o glossário.
%
%\glossary

% ----------------------------------------------------------
% APÊNDICES
% ----------------------------------------------------------

% ---
% Inicia os apêndices
% ---
% \begin{apendicesenv}

% % Imprime uma página indicando o início dos apêndices
% \partapendices

% % ----------------------------------------------------------
% \chapter{Quisque libero justo}
% % ----------------------------------------------------------

% % ----------------------------------------------------------
% \chapter{Nullam elementum urna}
% % ----------------------------------------------------------

% \end{apendicesenv}
% ---


% ----------------------------------------------------------
% ANEXOS
% ----------------------------------------------------------

% ---
% Inicia os anexos
% ---
% \begin{anexosenv}

% % Imprime uma página indicando o início dos anexos
% \partanexos

% % ---
% \chapter{Morbi ultrices rutrum lorem}
% % ---

% % ---
% \chapter{Cras non urna sed feugiat cum}
% % ---


% % ---
% \chapter{Fusce facilisis lacinia dui}
% % ---

% \end{anexosenv}

%---------------------------------------------------------------------
% INDICE REMISSIVO
%---------------------------------------------------------------------
% \phantompart
% \printindex
%---------------------------------------------------------------------
\includepdf{./05-extras/contracapa.pdf}
\end{document}


% ----------------------------------------------------------
% ELEMENTOS PÓS-TEXTUAIS
% ----------------------------------------------------------
\postextual
% ----------------------------------------------------------


% ----------------------------------------------------------
% REFERÊNCIAS
% ----------------------------------------------------------
\bibliography{references.bib}


% ----------------------------------------------------------
% GLOSSÁRIO
% ----------------------------------------------------------
%
% Consulte o manual da classe abntex2 para orientações sobre o glossário.
%
%\glossary

% ----------------------------------------------------------
% APÊNDICES
% ----------------------------------------------------------

% ---
% Inicia os apêndices
% ---
% \begin{apendicesenv}

% % Imprime uma página indicando o início dos apêndices
% \partapendices

% % ----------------------------------------------------------
% \chapter{Quisque libero justo}
% % ----------------------------------------------------------

% % ----------------------------------------------------------
% \chapter{Nullam elementum urna}
% % ----------------------------------------------------------

% \end{apendicesenv}
% ---


% ----------------------------------------------------------
% ANEXOS
% ----------------------------------------------------------

% ---
% Inicia os anexos
% ---
% \begin{anexosenv}

% % Imprime uma página indicando o início dos anexos
% \partanexos

% % ---
% \chapter{Morbi ultrices rutrum lorem}
% % ---

% % ---
% \chapter{Cras non urna sed feugiat cum}
% % ---


% % ---
% \chapter{Fusce facilisis lacinia dui}
% % ---

% \end{anexosenv}

%---------------------------------------------------------------------
% INDICE REMISSIVO
%---------------------------------------------------------------------
% \phantompart
% \printindex
%---------------------------------------------------------------------
\includepdf{./05-extras/contracapa.pdf}
\end{document}


\documentclass[
	% Opções da classe memoir
	12pt,				        % tamanho da fonte
	openany,			        % capítulos começam em pág ímpar (insere página vazia caso preciso)
	oneside,			        % para impressão em recto e verso. Oposto a oneside
	a4paper,			        % tamanho do papel. 
	% Opções da classe abntex2
	%chapter=TITLE,		        % títulos de capítulos convertidos em letras maiúsculas
	%section=TITLE,		        % títulos de seções convertidos em letras maiúsculas
	%subsection=TITLE,	        % títulos de subseções convertidos em letras maiúsculas
	%subsubsection=TITLE,       % títulos de subsubseções convertidos em letras maiúsculas
	% Opções do pacote babel
	english,			        % idioma adicional para hifenização
	brazil				        % o último idioma é o principal do documento
	]{abntex2}

% ---
% Pacotes básicos 
% ---
\usepackage{lmodern}			% Usa a fonte Latin Modern			
\usepackage[T1]{fontenc}		% Selecao de codigos de fonte.
\usepackage[utf8]{inputenc}		% Codificacao do documento (conversão automática dos acentos)
\usepackage{indentfirst}		% Indenta o primeiro parágrafo de cada seção.
\usepackage{color}				% Controle das cores
\usepackage{graphicx}			% Inclusão de gráficos
\usepackage{microtype} 			% para melhorias de justificação
% ---
		

% ---
% Pacotes de citações
% ---
% \usepackage[alf]{abntex2cite}						% Citações padrão ABNT
\usepackage[
  alf,
  abnt-nbr10520=1988,                     % Citações em sentence case
  abnt-emphasize=bf
  ]{abntex2cite}


% ---
% Pacotes adicionais
% ---
\usepackage{pdfpages}	% Inclusão de páginas em PDF


% ---
% Informações para CAPA e FOLHA DE ROSTO
% ---
\titulo{APRENDIZAGEM DE MÁQUINA PARA PREDIÇÃO DE DESFECHOS EM SAÚDE MENTAL: \textmd{UMA INTRODUÇÃO PARA PSICÓLOGOS}}
\autor{GUILHERME AUGUSTO DIAS}
\local{Porto Alegre}
\data{2024}
\orientador{Prof. Dr. Wagner de Lara Machado}
% \coorientador{}
% \instituicao{
%   Pontifícia Universidade Católica do Rio Grande do Sul
%   \par
%   Escola de Ciências da Saúde e da Vida
%   \par
%   Curso de Psicologia}
\tipotrabalho{TCC}
\preambulo{Trabalho de Conclusão de Curso apresentado como requisito para a obtenção do grau de Bacharel em Psicologia pela Escola de Ciências da Saúde e da Vida da Pontifícia Universidade Católica do Rio Grande do Sul.}
% ---

% ---
% Configurações de aparência do PDF final
% ---

% Alterando o aspecto da cor azul
\definecolor{blue}{RGB}{41,5,195}

% Informações do PDF
\makeatletter
\hypersetup{
		pdftitle={Aprendizagem de máquina para predição de desfechos em saúde mental: uma introdução para psicólogos}, 
		pdfauthor={Guilherme Augusto Dias},
    	pdfsubject={\imprimirpreambulo},
	    pdfcreator={LaTeX with abnTeX2},
		pdfkeywords={aprendizagem de máquina}{psicoterapia}{psicologia positiva}{depressão}, 
		colorlinks=true,		% true: colored links; false: boxed links
    	linkcolor=black,		% color of internal links
    	citecolor=black,		% color of links to bibliography
    	filecolor=black,		% color of file links
		urlcolor=black,
		bookmarksdepth=4
}
\makeatother
% --- 


% ---
% Posiciona figuras e tabelas no topo da página quando adicionadas sozinhas
% em um página em branco. Ver https://github.com/abntex/abntex2/issues/170
\makeatletter
\setlength{\@fptop}{5pt} % Set distance from top of page to first float
\makeatother
% ---

% ---
% Possibilita criação de Quadros e Lista de quadros.
% Ver https://github.com/abntex/abntex2/issues/176
%
\newcommand{\quadroname}{Quadro}
\newcommand{\listofquadrosname}{Lista de quadros}

\newfloat[chapter]{quadro}{loq}{\quadroname}
\newlistof{listofquadros}{loq}{\listofquadrosname}
\newlistentry{quadro}{loq}{0}

% Configurações para atender às regras da ABNT
\setfloatadjustment{quadro}{\centering}
\counterwithout{quadro}{chapter}
\renewcommand{\cftquadroname}{\quadroname\space} 
\renewcommand*{\cftquadroaftersnum}{\hfill--\hfill}

\setfloatlocations{quadro}{hbtp} % Ver https://github.com/abntex/abntex2/issues/176
% ---


% --- 
% Espaçamentos entre linhas e parágrafos 
% --- 

% O tamanho do parágrafo é dado por:
\setlength{\parindent}{1.3cm}

% Controle do espaçamento entre um parágrafo e outro:
\setlength{\parskip}{0.2cm}  % tente também \onelineskip

% ---
% Compila o índice
% ---
\makeindex
% ---


% ----
% Início do documento
% ----
\begin{document}

% Seleciona o idioma do documento (conforme pacotes do babel)
\selectlanguage{brazil}

% Retira espaço extra obsoleto entre as frases.
\frenchspacing 

% ----------------------------------------------------------
% ELEMENTOS PRÉ-TEXTUAIS
% ----------------------------------------------------------
\pretextual

% ---
% CAPA
% ---
\includepdf{./05-extras/capa.pdf}
% ---


% ---
% FOLHA DE ROSTO
% (o * indica que haverá a ficha bibliográfica)
% ---
\imprimirfolhaderosto
% ---


% ---
% Ficha bibliografica
% ---

% Isto é um exemplo de Ficha Catalográfica, ou ``Dados internacionais de
% catalogação-na-publicação''. Você pode utilizar este modelo como referência. 
% Porém, provavelmente a biblioteca da sua universidade lhe fornecerá um PDF
% com a ficha catalográfica definitiva após a defesa do trabalho. Quando estiver
% com o documento, salve-o como PDF no diretório do seu projeto e substitua todo
% o conteúdo de implementação deste arquivo pelo comando abaixo:
%
% \usepackage{pdfpages}		% necessário para comando \includepdf
% \begin{fichacatalografica}
%     \includepdf{fig_ficha_catalografica.pdf}
% \end{fichacatalografica}

% \begin{fichacatalografica}
% 	\sffamily
% 	\vspace*{\fill}                             % Posição vertical
% 	\begin{center}                              % Minipage Centralizado
% 	\fbox{\begin{minipage}[c][8cm]{13.5cm}		% Largura
% 	\small
% 	\imprimirautor
% 	%Sobrenome, Nome do autor
	
% 	\hspace{0.5cm} \imprimirtitulo  / \imprimirautor. --
% 	\imprimirlocal, \imprimirdata-
	
% 	\hspace{0.5cm} \thelastpage p. : il. (algumas color.) ; 30 cm.\\
	
% 	\hspace{0.5cm} \imprimirorientadorRotulo~\imprimirorientador\\
	
% 	\hspace{0.5cm}
% 	\parbox[t]{\textwidth}{\imprimirtipotrabalho~--~\imprimirinstituicao,
% 	\imprimirdata.}\\
	
% 	\hspace{0.5cm}
% 		1. Palavra-chave1.
% 		2. Palavra-chave2.
% 		2. Palavra-chave3.
% 		I. Orientador.
% 		II. Universidade xxx.
% 		III. Faculdade de xxx.
% 		IV. Título 			
% 	\end{minipage}}
% 	\end{center}
% \end{fichacatalografica}
% ---


% ---
% FOLHA DE APROVAÇÃO
% ---

% Isto é um exemplo de Folha de aprovação, elemento obrigatório da NBR
% 14724/2011 (seção 4.2.1.3). Você pode utilizar este modelo até a aprovação
% do trabalho. Após isso, substitua todo o conteúdo deste arquivo por uma
% imagem da página assinada pela banca com o comando abaixo:
%
% \begin{folhadeaprovacao}
% \includepdf{folhadeaprovacao_final.pdf}
% \end{folhadeaprovacao}
%
\setlength{\ABNTEXsignskip}{2cm}
\setlength{\ABNTEXsignwidth}{10cm}

\begin{folhadeaprovacao}
  \begin{center}
    {\ABNTEXchapterfont\large\imprimirautor}

    \vspace*{1cm}

    \begin{center}
      \ABNTEXchapterfont\bfseries\Large\imprimirtitulo
    \end{center}

    \vspace*{1cm}
    
    \hspace{.45\textwidth}
    \begin{minipage}{.5\textwidth}
        \imprimirpreambulo
    \end{minipage}%
   \end{center}

   \vspace*{1cm}
        
   \begin{center}
   Aprovado em 17 de julho de 2024.

   \vspace*{1cm}

   Banca examinadora:
   \end{center}

   \vspace*{-1cm}

   \assinatura{Prof. Ma. Caroline Santa Maria Rodrigues -- PUCRS}
   \assinatura{Prof. Dra. Clarissa Pinto Pizarro de Freitas -- PUCRS}
   \assinatura{\imprimirorientador -- PUCRS}
\end{folhadeaprovacao}
% ---


% ---
% Agradecimentos
% ---
\begin{agradecimentos}
Agradeço aos meus pais, Ademilson José Dias e Irinês Porfírio Dias, pelo apoio incondicional que me deram ao longo da vida,
incentivando mesmo as ideias questionáveis como aquela de fazer uma segunda graduação.

A minha irmã, Maria Eduarda Dias, cuja trajetória acadêmica é fonte de inspiração.

A minha namorada, Nathalia Fallavena Ceratti, que acompanhou a esta jornada de perto, oferecendo apoio e compreensão que
fizeram toda a diferença.

Aos meus amigos, Guilherme da Silva Petró, Norton Neves e Uilinton Luciano Marques da Silva, que foram compreensíveis nos
finais de semana em que não pude jogar RPG.

Ao professor Wagner de Lara Machado pela orientação e apoio no desenvolvimento deste trabalho.
\end{agradecimentos}
% ---


% ---
% RESUMOS
% ---
% Resumo em português
\setlength{\absparsep}{18pt} % ajusta o espaçamento dos parágrafos do resumo
\begin{resumo}
Desenvolvimentos na área de aprendizagem de máquina representam uma série de oportunidades para diversas áreas de pesquisa e de atuação profissional.
O uso de modelos de aprendizagem de máquina na predição de desfecho para tratamentos em saúde mental tem sido investigado e apresenta resultados promissores.
A aplicação desse tipo de modelo possibilitaria, por exemplo, maior assertividade na personalização de tratamentos psicoterápicos a nível individual.
Apesar dos benefícios potenciais, o tema ainda é pouco discutido nos espaços de pesquisa, formação e atuação em psicologia clínica. Este trabalho tem
por objetivo introduzir o tema de aprendizagem de máquina a psicólogos. Busca-se apresentar, de maneira acessível, os conceitos e terminologia básicos
em aprendizagem de máquina, proporcionando um panorama geral da área. Para complementar a exposição conceitual, é construído um exemplo a partir de dados
de um estudo sobre eficácia de intervenções em psicologia positiva para depressão via internet. Por fim, discutem-se algumas implicações do uso desse tipo
de tecnologia para a pesquisa e para a prática em psicoterapia.

 \textbf{Palavras-chave}: aprendizagem de máquina, psicoterapia, psicologia positiva, depressão.
\end{resumo}

% Resumo em inglês
\begin{resumo}[Abstract]
 \begin{otherlanguage*}{english}
   Advances in machine learning have opened up a multitude of opportunities for various research and professional fields. The use of machine learning models to
   predict treatment outcomes in mental health has been investigated showing promising results. The application of such models would allow greater assertiveness
   in tailoring psychotherapeutic treatments at the individual level. Despite its potential benefits, the topic remains unexplored in the fields of research, training
   and practice in clinical psychology. This work aims to introduce the topic of machine learning to psychologists. It presents the basic machine learning concepts and
   terminology in an accessible manner, providing a comprehensive overview of the field. To complement the conceptual discussion, an example is constructed using data
   from a study on the effectiveness of internet-based positive psychology interventions for depression. Finally, implications of using this type of technology in
   psychotherapy research and practice are discussed.
   \vspace{\onelineskip}
 
   \noindent 
   \textbf{Keywords}: machine learning, psychotherapy, positive psychology, depression.
 \end{otherlanguage*}
\end{resumo}


% ---
% LISTA DE ILUSTRAÇÕES
% ---
\pdfbookmark[0]{\listfigurename}{lof}
\listoffigures*
\cleardoublepage
% ---


% ---
% LISTA DE QUADROS
% ---
% \pdfbookmark[0]{\listofquadrosname}{loq}
% \listofquadros*
% \cleardoublepage
% ---


% ---
% LISTA DE TABELAS
% ---
% \pdfbookmark[0]{\listtablename}{lot}
% \listoftables*
% \cleardoublepage
% ---


% ---
% LISTA DE ABREVIATURAS E SIGLAS
% ---
% \begin{siglas}
%   \item[ABNT] Associação Brasileira de Normas Técnicas
%   \item[abnTeX] ABsurdas Normas para TeX
% \end{siglas}
% ---


% ---
% LISTA DE SÍMBOLOS
% ---
% \begin{simbolos}
%   \item[$ \Gamma $] Letra grega Gama
%   \item[$ \Lambda $] Lambda
%   \item[$ \zeta $] Letra grega minúscula zeta
%   \item[$ \in $] Pertence
% \end{simbolos}
% ---


% ---
% Sumário
% ---
\pdfbookmark[0]{\contentsname}{toc}
\tableofcontents*
\cleardoublepage
% ---


% ----------------------------------------------------------
% ELEMENTOS TEXTUAIS
% ----------------------------------------------------------
\textual

\documentclass[
	% Opções da classe memoir
	12pt,				        % tamanho da fonte
	openany,			        % capítulos começam em pág ímpar (insere página vazia caso preciso)
	oneside,			        % para impressão em recto e verso. Oposto a oneside
	a4paper,			        % tamanho do papel. 
	% Opções da classe abntex2
	%chapter=TITLE,		        % títulos de capítulos convertidos em letras maiúsculas
	%section=TITLE,		        % títulos de seções convertidos em letras maiúsculas
	%subsection=TITLE,	        % títulos de subseções convertidos em letras maiúsculas
	%subsubsection=TITLE,       % títulos de subsubseções convertidos em letras maiúsculas
	% Opções do pacote babel
	english,			        % idioma adicional para hifenização
	brazil				        % o último idioma é o principal do documento
	]{abntex2}

% ---
% Pacotes básicos 
% ---
\usepackage{lmodern}			% Usa a fonte Latin Modern			
\usepackage[T1]{fontenc}		% Selecao de codigos de fonte.
\usepackage[utf8]{inputenc}		% Codificacao do documento (conversão automática dos acentos)
\usepackage{indentfirst}		% Indenta o primeiro parágrafo de cada seção.
\usepackage{color}				% Controle das cores
\usepackage{graphicx}			% Inclusão de gráficos
\usepackage{microtype} 			% para melhorias de justificação
% ---
		

% ---
% Pacotes de citações
% ---
% \usepackage[alf]{abntex2cite}						% Citações padrão ABNT
\usepackage[
  alf,
  abnt-nbr10520=1988,                     % Citações em sentence case
  abnt-emphasize=bf
  ]{abntex2cite}


% ---
% Pacotes adicionais
% ---
\usepackage{pdfpages}	% Inclusão de páginas em PDF


% ---
% Informações para CAPA e FOLHA DE ROSTO
% ---
\titulo{APRENDIZAGEM DE MÁQUINA PARA PREDIÇÃO DE DESFECHOS EM SAÚDE MENTAL: \textmd{UMA INTRODUÇÃO PARA PSICÓLOGOS}}
\autor{GUILHERME AUGUSTO DIAS}
\local{Porto Alegre}
\data{2024}
\orientador{Prof. Dr. Wagner de Lara Machado}
% \coorientador{}
% \instituicao{
%   Pontifícia Universidade Católica do Rio Grande do Sul
%   \par
%   Escola de Ciências da Saúde e da Vida
%   \par
%   Curso de Psicologia}
\tipotrabalho{TCC}
\preambulo{Trabalho de Conclusão de Curso apresentado como requisito para a obtenção do grau de Bacharel em Psicologia pela Escola de Ciências da Saúde e da Vida da Pontifícia Universidade Católica do Rio Grande do Sul.}
% ---

% ---
% Configurações de aparência do PDF final
% ---

% Alterando o aspecto da cor azul
\definecolor{blue}{RGB}{41,5,195}

% Informações do PDF
\makeatletter
\hypersetup{
		pdftitle={Aprendizagem de máquina para predição de desfechos em saúde mental: uma introdução para psicólogos}, 
		pdfauthor={Guilherme Augusto Dias},
    	pdfsubject={\imprimirpreambulo},
	    pdfcreator={LaTeX with abnTeX2},
		pdfkeywords={aprendizagem de máquina}{psicoterapia}{psicologia positiva}{depressão}, 
		colorlinks=true,		% true: colored links; false: boxed links
    	linkcolor=black,		% color of internal links
    	citecolor=black,		% color of links to bibliography
    	filecolor=black,		% color of file links
		urlcolor=black,
		bookmarksdepth=4
}
\makeatother
% --- 


% ---
% Posiciona figuras e tabelas no topo da página quando adicionadas sozinhas
% em um página em branco. Ver https://github.com/abntex/abntex2/issues/170
\makeatletter
\setlength{\@fptop}{5pt} % Set distance from top of page to first float
\makeatother
% ---

% ---
% Possibilita criação de Quadros e Lista de quadros.
% Ver https://github.com/abntex/abntex2/issues/176
%
\newcommand{\quadroname}{Quadro}
\newcommand{\listofquadrosname}{Lista de quadros}

\newfloat[chapter]{quadro}{loq}{\quadroname}
\newlistof{listofquadros}{loq}{\listofquadrosname}
\newlistentry{quadro}{loq}{0}

% Configurações para atender às regras da ABNT
\setfloatadjustment{quadro}{\centering}
\counterwithout{quadro}{chapter}
\renewcommand{\cftquadroname}{\quadroname\space} 
\renewcommand*{\cftquadroaftersnum}{\hfill--\hfill}

\setfloatlocations{quadro}{hbtp} % Ver https://github.com/abntex/abntex2/issues/176
% ---


% --- 
% Espaçamentos entre linhas e parágrafos 
% --- 

% O tamanho do parágrafo é dado por:
\setlength{\parindent}{1.3cm}

% Controle do espaçamento entre um parágrafo e outro:
\setlength{\parskip}{0.2cm}  % tente também \onelineskip

% ---
% Compila o índice
% ---
\makeindex
% ---


% ----
% Início do documento
% ----
\begin{document}

% Seleciona o idioma do documento (conforme pacotes do babel)
\selectlanguage{brazil}

% Retira espaço extra obsoleto entre as frases.
\frenchspacing 

% ----------------------------------------------------------
% ELEMENTOS PRÉ-TEXTUAIS
% ----------------------------------------------------------
\pretextual

% ---
% CAPA
% ---
\includepdf{./05-extras/capa.pdf}
% ---


% ---
% FOLHA DE ROSTO
% (o * indica que haverá a ficha bibliográfica)
% ---
\imprimirfolhaderosto
% ---


% ---
% Ficha bibliografica
% ---

% Isto é um exemplo de Ficha Catalográfica, ou ``Dados internacionais de
% catalogação-na-publicação''. Você pode utilizar este modelo como referência. 
% Porém, provavelmente a biblioteca da sua universidade lhe fornecerá um PDF
% com a ficha catalográfica definitiva após a defesa do trabalho. Quando estiver
% com o documento, salve-o como PDF no diretório do seu projeto e substitua todo
% o conteúdo de implementação deste arquivo pelo comando abaixo:
%
% \usepackage{pdfpages}		% necessário para comando \includepdf
% \begin{fichacatalografica}
%     \includepdf{fig_ficha_catalografica.pdf}
% \end{fichacatalografica}

% \begin{fichacatalografica}
% 	\sffamily
% 	\vspace*{\fill}                             % Posição vertical
% 	\begin{center}                              % Minipage Centralizado
% 	\fbox{\begin{minipage}[c][8cm]{13.5cm}		% Largura
% 	\small
% 	\imprimirautor
% 	%Sobrenome, Nome do autor
	
% 	\hspace{0.5cm} \imprimirtitulo  / \imprimirautor. --
% 	\imprimirlocal, \imprimirdata-
	
% 	\hspace{0.5cm} \thelastpage p. : il. (algumas color.) ; 30 cm.\\
	
% 	\hspace{0.5cm} \imprimirorientadorRotulo~\imprimirorientador\\
	
% 	\hspace{0.5cm}
% 	\parbox[t]{\textwidth}{\imprimirtipotrabalho~--~\imprimirinstituicao,
% 	\imprimirdata.}\\
	
% 	\hspace{0.5cm}
% 		1. Palavra-chave1.
% 		2. Palavra-chave2.
% 		2. Palavra-chave3.
% 		I. Orientador.
% 		II. Universidade xxx.
% 		III. Faculdade de xxx.
% 		IV. Título 			
% 	\end{minipage}}
% 	\end{center}
% \end{fichacatalografica}
% ---


% ---
% FOLHA DE APROVAÇÃO
% ---

% Isto é um exemplo de Folha de aprovação, elemento obrigatório da NBR
% 14724/2011 (seção 4.2.1.3). Você pode utilizar este modelo até a aprovação
% do trabalho. Após isso, substitua todo o conteúdo deste arquivo por uma
% imagem da página assinada pela banca com o comando abaixo:
%
% \begin{folhadeaprovacao}
% \includepdf{folhadeaprovacao_final.pdf}
% \end{folhadeaprovacao}
%
\setlength{\ABNTEXsignskip}{2cm}
\setlength{\ABNTEXsignwidth}{10cm}

\begin{folhadeaprovacao}
  \begin{center}
    {\ABNTEXchapterfont\large\imprimirautor}

    \vspace*{1cm}

    \begin{center}
      \ABNTEXchapterfont\bfseries\Large\imprimirtitulo
    \end{center}

    \vspace*{1cm}
    
    \hspace{.45\textwidth}
    \begin{minipage}{.5\textwidth}
        \imprimirpreambulo
    \end{minipage}%
   \end{center}

   \vspace*{1cm}
        
   \begin{center}
   Aprovado em 17 de julho de 2024.

   \vspace*{1cm}

   Banca examinadora:
   \end{center}

   \vspace*{-1cm}

   \assinatura{Prof. Ma. Caroline Santa Maria Rodrigues -- PUCRS}
   \assinatura{Prof. Dra. Clarissa Pinto Pizarro de Freitas -- PUCRS}
   \assinatura{\imprimirorientador -- PUCRS}
\end{folhadeaprovacao}
% ---


% ---
% Agradecimentos
% ---
\begin{agradecimentos}
Agradeço aos meus pais, Ademilson José Dias e Irinês Porfírio Dias, pelo apoio incondicional que me deram ao longo da vida,
incentivando mesmo as ideias questionáveis como aquela de fazer uma segunda graduação.

A minha irmã, Maria Eduarda Dias, cuja trajetória acadêmica é fonte de inspiração.

A minha namorada, Nathalia Fallavena Ceratti, que acompanhou a esta jornada de perto, oferecendo apoio e compreensão que
fizeram toda a diferença.

Aos meus amigos, Guilherme da Silva Petró, Norton Neves e Uilinton Luciano Marques da Silva, que foram compreensíveis nos
finais de semana em que não pude jogar RPG.

Ao professor Wagner de Lara Machado pela orientação e apoio no desenvolvimento deste trabalho.
\end{agradecimentos}
% ---


% ---
% RESUMOS
% ---
% Resumo em português
\setlength{\absparsep}{18pt} % ajusta o espaçamento dos parágrafos do resumo
\begin{resumo}
Desenvolvimentos na área de aprendizagem de máquina representam uma série de oportunidades para diversas áreas de pesquisa e de atuação profissional.
O uso de modelos de aprendizagem de máquina na predição de desfecho para tratamentos em saúde mental tem sido investigado e apresenta resultados promissores.
A aplicação desse tipo de modelo possibilitaria, por exemplo, maior assertividade na personalização de tratamentos psicoterápicos a nível individual.
Apesar dos benefícios potenciais, o tema ainda é pouco discutido nos espaços de pesquisa, formação e atuação em psicologia clínica. Este trabalho tem
por objetivo introduzir o tema de aprendizagem de máquina a psicólogos. Busca-se apresentar, de maneira acessível, os conceitos e terminologia básicos
em aprendizagem de máquina, proporcionando um panorama geral da área. Para complementar a exposição conceitual, é construído um exemplo a partir de dados
de um estudo sobre eficácia de intervenções em psicologia positiva para depressão via internet. Por fim, discutem-se algumas implicações do uso desse tipo
de tecnologia para a pesquisa e para a prática em psicoterapia.

 \textbf{Palavras-chave}: aprendizagem de máquina, psicoterapia, psicologia positiva, depressão.
\end{resumo}

% Resumo em inglês
\begin{resumo}[Abstract]
 \begin{otherlanguage*}{english}
   Advances in machine learning have opened up a multitude of opportunities for various research and professional fields. The use of machine learning models to
   predict treatment outcomes in mental health has been investigated showing promising results. The application of such models would allow greater assertiveness
   in tailoring psychotherapeutic treatments at the individual level. Despite its potential benefits, the topic remains unexplored in the fields of research, training
   and practice in clinical psychology. This work aims to introduce the topic of machine learning to psychologists. It presents the basic machine learning concepts and
   terminology in an accessible manner, providing a comprehensive overview of the field. To complement the conceptual discussion, an example is constructed using data
   from a study on the effectiveness of internet-based positive psychology interventions for depression. Finally, implications of using this type of technology in
   psychotherapy research and practice are discussed.
   \vspace{\onelineskip}
 
   \noindent 
   \textbf{Keywords}: machine learning, psychotherapy, positive psychology, depression.
 \end{otherlanguage*}
\end{resumo}


% ---
% LISTA DE ILUSTRAÇÕES
% ---
\pdfbookmark[0]{\listfigurename}{lof}
\listoffigures*
\cleardoublepage
% ---


% ---
% LISTA DE QUADROS
% ---
% \pdfbookmark[0]{\listofquadrosname}{loq}
% \listofquadros*
% \cleardoublepage
% ---


% ---
% LISTA DE TABELAS
% ---
% \pdfbookmark[0]{\listtablename}{lot}
% \listoftables*
% \cleardoublepage
% ---


% ---
% LISTA DE ABREVIATURAS E SIGLAS
% ---
% \begin{siglas}
%   \item[ABNT] Associação Brasileira de Normas Técnicas
%   \item[abnTeX] ABsurdas Normas para TeX
% \end{siglas}
% ---


% ---
% LISTA DE SÍMBOLOS
% ---
% \begin{simbolos}
%   \item[$ \Gamma $] Letra grega Gama
%   \item[$ \Lambda $] Lambda
%   \item[$ \zeta $] Letra grega minúscula zeta
%   \item[$ \in $] Pertence
% \end{simbolos}
% ---


% ---
% Sumário
% ---
\pdfbookmark[0]{\contentsname}{toc}
\tableofcontents*
\cleardoublepage
% ---


% ----------------------------------------------------------
% ELEMENTOS TEXTUAIS
% ----------------------------------------------------------
\textual

\documentclass[
	% Opções da classe memoir
	12pt,				        % tamanho da fonte
	openany,			        % capítulos começam em pág ímpar (insere página vazia caso preciso)
	oneside,			        % para impressão em recto e verso. Oposto a oneside
	a4paper,			        % tamanho do papel. 
	% Opções da classe abntex2
	%chapter=TITLE,		        % títulos de capítulos convertidos em letras maiúsculas
	%section=TITLE,		        % títulos de seções convertidos em letras maiúsculas
	%subsection=TITLE,	        % títulos de subseções convertidos em letras maiúsculas
	%subsubsection=TITLE,       % títulos de subsubseções convertidos em letras maiúsculas
	% Opções do pacote babel
	english,			        % idioma adicional para hifenização
	brazil				        % o último idioma é o principal do documento
	]{abntex2}

% ---
% Pacotes básicos 
% ---
\usepackage{lmodern}			% Usa a fonte Latin Modern			
\usepackage[T1]{fontenc}		% Selecao de codigos de fonte.
\usepackage[utf8]{inputenc}		% Codificacao do documento (conversão automática dos acentos)
\usepackage{indentfirst}		% Indenta o primeiro parágrafo de cada seção.
\usepackage{color}				% Controle das cores
\usepackage{graphicx}			% Inclusão de gráficos
\usepackage{microtype} 			% para melhorias de justificação
% ---
		

% ---
% Pacotes de citações
% ---
% \usepackage[alf]{abntex2cite}						% Citações padrão ABNT
\usepackage[
  alf,
  abnt-nbr10520=1988,                     % Citações em sentence case
  abnt-emphasize=bf
  ]{abntex2cite}


% ---
% Pacotes adicionais
% ---
\usepackage{pdfpages}	% Inclusão de páginas em PDF


% ---
% Informações para CAPA e FOLHA DE ROSTO
% ---
\titulo{APRENDIZAGEM DE MÁQUINA PARA PREDIÇÃO DE DESFECHOS EM SAÚDE MENTAL: \textmd{UMA INTRODUÇÃO PARA PSICÓLOGOS}}
\autor{GUILHERME AUGUSTO DIAS}
\local{Porto Alegre}
\data{2024}
\orientador{Prof. Dr. Wagner de Lara Machado}
% \coorientador{}
% \instituicao{
%   Pontifícia Universidade Católica do Rio Grande do Sul
%   \par
%   Escola de Ciências da Saúde e da Vida
%   \par
%   Curso de Psicologia}
\tipotrabalho{TCC}
\preambulo{Trabalho de Conclusão de Curso apresentado como requisito para a obtenção do grau de Bacharel em Psicologia pela Escola de Ciências da Saúde e da Vida da Pontifícia Universidade Católica do Rio Grande do Sul.}
% ---

% ---
% Configurações de aparência do PDF final
% ---

% Alterando o aspecto da cor azul
\definecolor{blue}{RGB}{41,5,195}

% Informações do PDF
\makeatletter
\hypersetup{
		pdftitle={Aprendizagem de máquina para predição de desfechos em saúde mental: uma introdução para psicólogos}, 
		pdfauthor={Guilherme Augusto Dias},
    	pdfsubject={\imprimirpreambulo},
	    pdfcreator={LaTeX with abnTeX2},
		pdfkeywords={aprendizagem de máquina}{psicoterapia}{psicologia positiva}{depressão}, 
		colorlinks=true,		% true: colored links; false: boxed links
    	linkcolor=black,		% color of internal links
    	citecolor=black,		% color of links to bibliography
    	filecolor=black,		% color of file links
		urlcolor=black,
		bookmarksdepth=4
}
\makeatother
% --- 


% ---
% Posiciona figuras e tabelas no topo da página quando adicionadas sozinhas
% em um página em branco. Ver https://github.com/abntex/abntex2/issues/170
\makeatletter
\setlength{\@fptop}{5pt} % Set distance from top of page to first float
\makeatother
% ---

% ---
% Possibilita criação de Quadros e Lista de quadros.
% Ver https://github.com/abntex/abntex2/issues/176
%
\newcommand{\quadroname}{Quadro}
\newcommand{\listofquadrosname}{Lista de quadros}

\newfloat[chapter]{quadro}{loq}{\quadroname}
\newlistof{listofquadros}{loq}{\listofquadrosname}
\newlistentry{quadro}{loq}{0}

% Configurações para atender às regras da ABNT
\setfloatadjustment{quadro}{\centering}
\counterwithout{quadro}{chapter}
\renewcommand{\cftquadroname}{\quadroname\space} 
\renewcommand*{\cftquadroaftersnum}{\hfill--\hfill}

\setfloatlocations{quadro}{hbtp} % Ver https://github.com/abntex/abntex2/issues/176
% ---


% --- 
% Espaçamentos entre linhas e parágrafos 
% --- 

% O tamanho do parágrafo é dado por:
\setlength{\parindent}{1.3cm}

% Controle do espaçamento entre um parágrafo e outro:
\setlength{\parskip}{0.2cm}  % tente também \onelineskip

% ---
% Compila o índice
% ---
\makeindex
% ---


% ----
% Início do documento
% ----
\begin{document}

% Seleciona o idioma do documento (conforme pacotes do babel)
\selectlanguage{brazil}

% Retira espaço extra obsoleto entre as frases.
\frenchspacing 

% ----------------------------------------------------------
% ELEMENTOS PRÉ-TEXTUAIS
% ----------------------------------------------------------
\pretextual

% ---
% CAPA
% ---
\includepdf{./05-extras/capa.pdf}
% ---


% ---
% FOLHA DE ROSTO
% (o * indica que haverá a ficha bibliográfica)
% ---
\imprimirfolhaderosto
% ---


% ---
% Ficha bibliografica
% ---

% Isto é um exemplo de Ficha Catalográfica, ou ``Dados internacionais de
% catalogação-na-publicação''. Você pode utilizar este modelo como referência. 
% Porém, provavelmente a biblioteca da sua universidade lhe fornecerá um PDF
% com a ficha catalográfica definitiva após a defesa do trabalho. Quando estiver
% com o documento, salve-o como PDF no diretório do seu projeto e substitua todo
% o conteúdo de implementação deste arquivo pelo comando abaixo:
%
% \usepackage{pdfpages}		% necessário para comando \includepdf
% \begin{fichacatalografica}
%     \includepdf{fig_ficha_catalografica.pdf}
% \end{fichacatalografica}

% \begin{fichacatalografica}
% 	\sffamily
% 	\vspace*{\fill}                             % Posição vertical
% 	\begin{center}                              % Minipage Centralizado
% 	\fbox{\begin{minipage}[c][8cm]{13.5cm}		% Largura
% 	\small
% 	\imprimirautor
% 	%Sobrenome, Nome do autor
	
% 	\hspace{0.5cm} \imprimirtitulo  / \imprimirautor. --
% 	\imprimirlocal, \imprimirdata-
	
% 	\hspace{0.5cm} \thelastpage p. : il. (algumas color.) ; 30 cm.\\
	
% 	\hspace{0.5cm} \imprimirorientadorRotulo~\imprimirorientador\\
	
% 	\hspace{0.5cm}
% 	\parbox[t]{\textwidth}{\imprimirtipotrabalho~--~\imprimirinstituicao,
% 	\imprimirdata.}\\
	
% 	\hspace{0.5cm}
% 		1. Palavra-chave1.
% 		2. Palavra-chave2.
% 		2. Palavra-chave3.
% 		I. Orientador.
% 		II. Universidade xxx.
% 		III. Faculdade de xxx.
% 		IV. Título 			
% 	\end{minipage}}
% 	\end{center}
% \end{fichacatalografica}
% ---


% ---
% FOLHA DE APROVAÇÃO
% ---

% Isto é um exemplo de Folha de aprovação, elemento obrigatório da NBR
% 14724/2011 (seção 4.2.1.3). Você pode utilizar este modelo até a aprovação
% do trabalho. Após isso, substitua todo o conteúdo deste arquivo por uma
% imagem da página assinada pela banca com o comando abaixo:
%
% \begin{folhadeaprovacao}
% \includepdf{folhadeaprovacao_final.pdf}
% \end{folhadeaprovacao}
%
\setlength{\ABNTEXsignskip}{2cm}
\setlength{\ABNTEXsignwidth}{10cm}

\begin{folhadeaprovacao}
  \begin{center}
    {\ABNTEXchapterfont\large\imprimirautor}

    \vspace*{1cm}

    \begin{center}
      \ABNTEXchapterfont\bfseries\Large\imprimirtitulo
    \end{center}

    \vspace*{1cm}
    
    \hspace{.45\textwidth}
    \begin{minipage}{.5\textwidth}
        \imprimirpreambulo
    \end{minipage}%
   \end{center}

   \vspace*{1cm}
        
   \begin{center}
   Aprovado em 17 de julho de 2024.

   \vspace*{1cm}

   Banca examinadora:
   \end{center}

   \vspace*{-1cm}

   \assinatura{Prof. Ma. Caroline Santa Maria Rodrigues -- PUCRS}
   \assinatura{Prof. Dra. Clarissa Pinto Pizarro de Freitas -- PUCRS}
   \assinatura{\imprimirorientador -- PUCRS}
\end{folhadeaprovacao}
% ---


% ---
% Agradecimentos
% ---
\begin{agradecimentos}
Agradeço aos meus pais, Ademilson José Dias e Irinês Porfírio Dias, pelo apoio incondicional que me deram ao longo da vida,
incentivando mesmo as ideias questionáveis como aquela de fazer uma segunda graduação.

A minha irmã, Maria Eduarda Dias, cuja trajetória acadêmica é fonte de inspiração.

A minha namorada, Nathalia Fallavena Ceratti, que acompanhou a esta jornada de perto, oferecendo apoio e compreensão que
fizeram toda a diferença.

Aos meus amigos, Guilherme da Silva Petró, Norton Neves e Uilinton Luciano Marques da Silva, que foram compreensíveis nos
finais de semana em que não pude jogar RPG.

Ao professor Wagner de Lara Machado pela orientação e apoio no desenvolvimento deste trabalho.
\end{agradecimentos}
% ---


% ---
% RESUMOS
% ---
% Resumo em português
\setlength{\absparsep}{18pt} % ajusta o espaçamento dos parágrafos do resumo
\begin{resumo}
Desenvolvimentos na área de aprendizagem de máquina representam uma série de oportunidades para diversas áreas de pesquisa e de atuação profissional.
O uso de modelos de aprendizagem de máquina na predição de desfecho para tratamentos em saúde mental tem sido investigado e apresenta resultados promissores.
A aplicação desse tipo de modelo possibilitaria, por exemplo, maior assertividade na personalização de tratamentos psicoterápicos a nível individual.
Apesar dos benefícios potenciais, o tema ainda é pouco discutido nos espaços de pesquisa, formação e atuação em psicologia clínica. Este trabalho tem
por objetivo introduzir o tema de aprendizagem de máquina a psicólogos. Busca-se apresentar, de maneira acessível, os conceitos e terminologia básicos
em aprendizagem de máquina, proporcionando um panorama geral da área. Para complementar a exposição conceitual, é construído um exemplo a partir de dados
de um estudo sobre eficácia de intervenções em psicologia positiva para depressão via internet. Por fim, discutem-se algumas implicações do uso desse tipo
de tecnologia para a pesquisa e para a prática em psicoterapia.

 \textbf{Palavras-chave}: aprendizagem de máquina, psicoterapia, psicologia positiva, depressão.
\end{resumo}

% Resumo em inglês
\begin{resumo}[Abstract]
 \begin{otherlanguage*}{english}
   Advances in machine learning have opened up a multitude of opportunities for various research and professional fields. The use of machine learning models to
   predict treatment outcomes in mental health has been investigated showing promising results. The application of such models would allow greater assertiveness
   in tailoring psychotherapeutic treatments at the individual level. Despite its potential benefits, the topic remains unexplored in the fields of research, training
   and practice in clinical psychology. This work aims to introduce the topic of machine learning to psychologists. It presents the basic machine learning concepts and
   terminology in an accessible manner, providing a comprehensive overview of the field. To complement the conceptual discussion, an example is constructed using data
   from a study on the effectiveness of internet-based positive psychology interventions for depression. Finally, implications of using this type of technology in
   psychotherapy research and practice are discussed.
   \vspace{\onelineskip}
 
   \noindent 
   \textbf{Keywords}: machine learning, psychotherapy, positive psychology, depression.
 \end{otherlanguage*}
\end{resumo}


% ---
% LISTA DE ILUSTRAÇÕES
% ---
\pdfbookmark[0]{\listfigurename}{lof}
\listoffigures*
\cleardoublepage
% ---


% ---
% LISTA DE QUADROS
% ---
% \pdfbookmark[0]{\listofquadrosname}{loq}
% \listofquadros*
% \cleardoublepage
% ---


% ---
% LISTA DE TABELAS
% ---
% \pdfbookmark[0]{\listtablename}{lot}
% \listoftables*
% \cleardoublepage
% ---


% ---
% LISTA DE ABREVIATURAS E SIGLAS
% ---
% \begin{siglas}
%   \item[ABNT] Associação Brasileira de Normas Técnicas
%   \item[abnTeX] ABsurdas Normas para TeX
% \end{siglas}
% ---


% ---
% LISTA DE SÍMBOLOS
% ---
% \begin{simbolos}
%   \item[$ \Gamma $] Letra grega Gama
%   \item[$ \Lambda $] Lambda
%   \item[$ \zeta $] Letra grega minúscula zeta
%   \item[$ \in $] Pertence
% \end{simbolos}
% ---


% ---
% Sumário
% ---
\pdfbookmark[0]{\contentsname}{toc}
\tableofcontents*
\cleardoublepage
% ---


% ----------------------------------------------------------
% ELEMENTOS TEXTUAIS
% ----------------------------------------------------------
\textual

\include{./01-introducao/main.tex}

\include{./02-machine-learning/main.tex}

\include{./03-exemplo/main.tex}

\include{./04-conclusao/main.tex}

% ----------------------------------------------------------
% ELEMENTOS PÓS-TEXTUAIS
% ----------------------------------------------------------
\postextual
% ----------------------------------------------------------


% ----------------------------------------------------------
% REFERÊNCIAS
% ----------------------------------------------------------
\bibliography{references.bib}


% ----------------------------------------------------------
% GLOSSÁRIO
% ----------------------------------------------------------
%
% Consulte o manual da classe abntex2 para orientações sobre o glossário.
%
%\glossary

% ----------------------------------------------------------
% APÊNDICES
% ----------------------------------------------------------

% ---
% Inicia os apêndices
% ---
% \begin{apendicesenv}

% % Imprime uma página indicando o início dos apêndices
% \partapendices

% % ----------------------------------------------------------
% \chapter{Quisque libero justo}
% % ----------------------------------------------------------

% % ----------------------------------------------------------
% \chapter{Nullam elementum urna}
% % ----------------------------------------------------------

% \end{apendicesenv}
% ---


% ----------------------------------------------------------
% ANEXOS
% ----------------------------------------------------------

% ---
% Inicia os anexos
% ---
% \begin{anexosenv}

% % Imprime uma página indicando o início dos anexos
% \partanexos

% % ---
% \chapter{Morbi ultrices rutrum lorem}
% % ---

% % ---
% \chapter{Cras non urna sed feugiat cum}
% % ---


% % ---
% \chapter{Fusce facilisis lacinia dui}
% % ---

% \end{anexosenv}

%---------------------------------------------------------------------
% INDICE REMISSIVO
%---------------------------------------------------------------------
% \phantompart
% \printindex
%---------------------------------------------------------------------
\includepdf{./05-extras/contracapa.pdf}
\end{document}


\documentclass[
	% Opções da classe memoir
	12pt,				        % tamanho da fonte
	openany,			        % capítulos começam em pág ímpar (insere página vazia caso preciso)
	oneside,			        % para impressão em recto e verso. Oposto a oneside
	a4paper,			        % tamanho do papel. 
	% Opções da classe abntex2
	%chapter=TITLE,		        % títulos de capítulos convertidos em letras maiúsculas
	%section=TITLE,		        % títulos de seções convertidos em letras maiúsculas
	%subsection=TITLE,	        % títulos de subseções convertidos em letras maiúsculas
	%subsubsection=TITLE,       % títulos de subsubseções convertidos em letras maiúsculas
	% Opções do pacote babel
	english,			        % idioma adicional para hifenização
	brazil				        % o último idioma é o principal do documento
	]{abntex2}

% ---
% Pacotes básicos 
% ---
\usepackage{lmodern}			% Usa a fonte Latin Modern			
\usepackage[T1]{fontenc}		% Selecao de codigos de fonte.
\usepackage[utf8]{inputenc}		% Codificacao do documento (conversão automática dos acentos)
\usepackage{indentfirst}		% Indenta o primeiro parágrafo de cada seção.
\usepackage{color}				% Controle das cores
\usepackage{graphicx}			% Inclusão de gráficos
\usepackage{microtype} 			% para melhorias de justificação
% ---
		

% ---
% Pacotes de citações
% ---
% \usepackage[alf]{abntex2cite}						% Citações padrão ABNT
\usepackage[
  alf,
  abnt-nbr10520=1988,                     % Citações em sentence case
  abnt-emphasize=bf
  ]{abntex2cite}


% ---
% Pacotes adicionais
% ---
\usepackage{pdfpages}	% Inclusão de páginas em PDF


% ---
% Informações para CAPA e FOLHA DE ROSTO
% ---
\titulo{APRENDIZAGEM DE MÁQUINA PARA PREDIÇÃO DE DESFECHOS EM SAÚDE MENTAL: \textmd{UMA INTRODUÇÃO PARA PSICÓLOGOS}}
\autor{GUILHERME AUGUSTO DIAS}
\local{Porto Alegre}
\data{2024}
\orientador{Prof. Dr. Wagner de Lara Machado}
% \coorientador{}
% \instituicao{
%   Pontifícia Universidade Católica do Rio Grande do Sul
%   \par
%   Escola de Ciências da Saúde e da Vida
%   \par
%   Curso de Psicologia}
\tipotrabalho{TCC}
\preambulo{Trabalho de Conclusão de Curso apresentado como requisito para a obtenção do grau de Bacharel em Psicologia pela Escola de Ciências da Saúde e da Vida da Pontifícia Universidade Católica do Rio Grande do Sul.}
% ---

% ---
% Configurações de aparência do PDF final
% ---

% Alterando o aspecto da cor azul
\definecolor{blue}{RGB}{41,5,195}

% Informações do PDF
\makeatletter
\hypersetup{
		pdftitle={Aprendizagem de máquina para predição de desfechos em saúde mental: uma introdução para psicólogos}, 
		pdfauthor={Guilherme Augusto Dias},
    	pdfsubject={\imprimirpreambulo},
	    pdfcreator={LaTeX with abnTeX2},
		pdfkeywords={aprendizagem de máquina}{psicoterapia}{psicologia positiva}{depressão}, 
		colorlinks=true,		% true: colored links; false: boxed links
    	linkcolor=black,		% color of internal links
    	citecolor=black,		% color of links to bibliography
    	filecolor=black,		% color of file links
		urlcolor=black,
		bookmarksdepth=4
}
\makeatother
% --- 


% ---
% Posiciona figuras e tabelas no topo da página quando adicionadas sozinhas
% em um página em branco. Ver https://github.com/abntex/abntex2/issues/170
\makeatletter
\setlength{\@fptop}{5pt} % Set distance from top of page to first float
\makeatother
% ---

% ---
% Possibilita criação de Quadros e Lista de quadros.
% Ver https://github.com/abntex/abntex2/issues/176
%
\newcommand{\quadroname}{Quadro}
\newcommand{\listofquadrosname}{Lista de quadros}

\newfloat[chapter]{quadro}{loq}{\quadroname}
\newlistof{listofquadros}{loq}{\listofquadrosname}
\newlistentry{quadro}{loq}{0}

% Configurações para atender às regras da ABNT
\setfloatadjustment{quadro}{\centering}
\counterwithout{quadro}{chapter}
\renewcommand{\cftquadroname}{\quadroname\space} 
\renewcommand*{\cftquadroaftersnum}{\hfill--\hfill}

\setfloatlocations{quadro}{hbtp} % Ver https://github.com/abntex/abntex2/issues/176
% ---


% --- 
% Espaçamentos entre linhas e parágrafos 
% --- 

% O tamanho do parágrafo é dado por:
\setlength{\parindent}{1.3cm}

% Controle do espaçamento entre um parágrafo e outro:
\setlength{\parskip}{0.2cm}  % tente também \onelineskip

% ---
% Compila o índice
% ---
\makeindex
% ---


% ----
% Início do documento
% ----
\begin{document}

% Seleciona o idioma do documento (conforme pacotes do babel)
\selectlanguage{brazil}

% Retira espaço extra obsoleto entre as frases.
\frenchspacing 

% ----------------------------------------------------------
% ELEMENTOS PRÉ-TEXTUAIS
% ----------------------------------------------------------
\pretextual

% ---
% CAPA
% ---
\includepdf{./05-extras/capa.pdf}
% ---


% ---
% FOLHA DE ROSTO
% (o * indica que haverá a ficha bibliográfica)
% ---
\imprimirfolhaderosto
% ---


% ---
% Ficha bibliografica
% ---

% Isto é um exemplo de Ficha Catalográfica, ou ``Dados internacionais de
% catalogação-na-publicação''. Você pode utilizar este modelo como referência. 
% Porém, provavelmente a biblioteca da sua universidade lhe fornecerá um PDF
% com a ficha catalográfica definitiva após a defesa do trabalho. Quando estiver
% com o documento, salve-o como PDF no diretório do seu projeto e substitua todo
% o conteúdo de implementação deste arquivo pelo comando abaixo:
%
% \usepackage{pdfpages}		% necessário para comando \includepdf
% \begin{fichacatalografica}
%     \includepdf{fig_ficha_catalografica.pdf}
% \end{fichacatalografica}

% \begin{fichacatalografica}
% 	\sffamily
% 	\vspace*{\fill}                             % Posição vertical
% 	\begin{center}                              % Minipage Centralizado
% 	\fbox{\begin{minipage}[c][8cm]{13.5cm}		% Largura
% 	\small
% 	\imprimirautor
% 	%Sobrenome, Nome do autor
	
% 	\hspace{0.5cm} \imprimirtitulo  / \imprimirautor. --
% 	\imprimirlocal, \imprimirdata-
	
% 	\hspace{0.5cm} \thelastpage p. : il. (algumas color.) ; 30 cm.\\
	
% 	\hspace{0.5cm} \imprimirorientadorRotulo~\imprimirorientador\\
	
% 	\hspace{0.5cm}
% 	\parbox[t]{\textwidth}{\imprimirtipotrabalho~--~\imprimirinstituicao,
% 	\imprimirdata.}\\
	
% 	\hspace{0.5cm}
% 		1. Palavra-chave1.
% 		2. Palavra-chave2.
% 		2. Palavra-chave3.
% 		I. Orientador.
% 		II. Universidade xxx.
% 		III. Faculdade de xxx.
% 		IV. Título 			
% 	\end{minipage}}
% 	\end{center}
% \end{fichacatalografica}
% ---


% ---
% FOLHA DE APROVAÇÃO
% ---

% Isto é um exemplo de Folha de aprovação, elemento obrigatório da NBR
% 14724/2011 (seção 4.2.1.3). Você pode utilizar este modelo até a aprovação
% do trabalho. Após isso, substitua todo o conteúdo deste arquivo por uma
% imagem da página assinada pela banca com o comando abaixo:
%
% \begin{folhadeaprovacao}
% \includepdf{folhadeaprovacao_final.pdf}
% \end{folhadeaprovacao}
%
\setlength{\ABNTEXsignskip}{2cm}
\setlength{\ABNTEXsignwidth}{10cm}

\begin{folhadeaprovacao}
  \begin{center}
    {\ABNTEXchapterfont\large\imprimirautor}

    \vspace*{1cm}

    \begin{center}
      \ABNTEXchapterfont\bfseries\Large\imprimirtitulo
    \end{center}

    \vspace*{1cm}
    
    \hspace{.45\textwidth}
    \begin{minipage}{.5\textwidth}
        \imprimirpreambulo
    \end{minipage}%
   \end{center}

   \vspace*{1cm}
        
   \begin{center}
   Aprovado em 17 de julho de 2024.

   \vspace*{1cm}

   Banca examinadora:
   \end{center}

   \vspace*{-1cm}

   \assinatura{Prof. Ma. Caroline Santa Maria Rodrigues -- PUCRS}
   \assinatura{Prof. Dra. Clarissa Pinto Pizarro de Freitas -- PUCRS}
   \assinatura{\imprimirorientador -- PUCRS}
\end{folhadeaprovacao}
% ---


% ---
% Agradecimentos
% ---
\begin{agradecimentos}
Agradeço aos meus pais, Ademilson José Dias e Irinês Porfírio Dias, pelo apoio incondicional que me deram ao longo da vida,
incentivando mesmo as ideias questionáveis como aquela de fazer uma segunda graduação.

A minha irmã, Maria Eduarda Dias, cuja trajetória acadêmica é fonte de inspiração.

A minha namorada, Nathalia Fallavena Ceratti, que acompanhou a esta jornada de perto, oferecendo apoio e compreensão que
fizeram toda a diferença.

Aos meus amigos, Guilherme da Silva Petró, Norton Neves e Uilinton Luciano Marques da Silva, que foram compreensíveis nos
finais de semana em que não pude jogar RPG.

Ao professor Wagner de Lara Machado pela orientação e apoio no desenvolvimento deste trabalho.
\end{agradecimentos}
% ---


% ---
% RESUMOS
% ---
% Resumo em português
\setlength{\absparsep}{18pt} % ajusta o espaçamento dos parágrafos do resumo
\begin{resumo}
Desenvolvimentos na área de aprendizagem de máquina representam uma série de oportunidades para diversas áreas de pesquisa e de atuação profissional.
O uso de modelos de aprendizagem de máquina na predição de desfecho para tratamentos em saúde mental tem sido investigado e apresenta resultados promissores.
A aplicação desse tipo de modelo possibilitaria, por exemplo, maior assertividade na personalização de tratamentos psicoterápicos a nível individual.
Apesar dos benefícios potenciais, o tema ainda é pouco discutido nos espaços de pesquisa, formação e atuação em psicologia clínica. Este trabalho tem
por objetivo introduzir o tema de aprendizagem de máquina a psicólogos. Busca-se apresentar, de maneira acessível, os conceitos e terminologia básicos
em aprendizagem de máquina, proporcionando um panorama geral da área. Para complementar a exposição conceitual, é construído um exemplo a partir de dados
de um estudo sobre eficácia de intervenções em psicologia positiva para depressão via internet. Por fim, discutem-se algumas implicações do uso desse tipo
de tecnologia para a pesquisa e para a prática em psicoterapia.

 \textbf{Palavras-chave}: aprendizagem de máquina, psicoterapia, psicologia positiva, depressão.
\end{resumo}

% Resumo em inglês
\begin{resumo}[Abstract]
 \begin{otherlanguage*}{english}
   Advances in machine learning have opened up a multitude of opportunities for various research and professional fields. The use of machine learning models to
   predict treatment outcomes in mental health has been investigated showing promising results. The application of such models would allow greater assertiveness
   in tailoring psychotherapeutic treatments at the individual level. Despite its potential benefits, the topic remains unexplored in the fields of research, training
   and practice in clinical psychology. This work aims to introduce the topic of machine learning to psychologists. It presents the basic machine learning concepts and
   terminology in an accessible manner, providing a comprehensive overview of the field. To complement the conceptual discussion, an example is constructed using data
   from a study on the effectiveness of internet-based positive psychology interventions for depression. Finally, implications of using this type of technology in
   psychotherapy research and practice are discussed.
   \vspace{\onelineskip}
 
   \noindent 
   \textbf{Keywords}: machine learning, psychotherapy, positive psychology, depression.
 \end{otherlanguage*}
\end{resumo}


% ---
% LISTA DE ILUSTRAÇÕES
% ---
\pdfbookmark[0]{\listfigurename}{lof}
\listoffigures*
\cleardoublepage
% ---


% ---
% LISTA DE QUADROS
% ---
% \pdfbookmark[0]{\listofquadrosname}{loq}
% \listofquadros*
% \cleardoublepage
% ---


% ---
% LISTA DE TABELAS
% ---
% \pdfbookmark[0]{\listtablename}{lot}
% \listoftables*
% \cleardoublepage
% ---


% ---
% LISTA DE ABREVIATURAS E SIGLAS
% ---
% \begin{siglas}
%   \item[ABNT] Associação Brasileira de Normas Técnicas
%   \item[abnTeX] ABsurdas Normas para TeX
% \end{siglas}
% ---


% ---
% LISTA DE SÍMBOLOS
% ---
% \begin{simbolos}
%   \item[$ \Gamma $] Letra grega Gama
%   \item[$ \Lambda $] Lambda
%   \item[$ \zeta $] Letra grega minúscula zeta
%   \item[$ \in $] Pertence
% \end{simbolos}
% ---


% ---
% Sumário
% ---
\pdfbookmark[0]{\contentsname}{toc}
\tableofcontents*
\cleardoublepage
% ---


% ----------------------------------------------------------
% ELEMENTOS TEXTUAIS
% ----------------------------------------------------------
\textual

\include{./01-introducao/main.tex}

\include{./02-machine-learning/main.tex}

\include{./03-exemplo/main.tex}

\include{./04-conclusao/main.tex}

% ----------------------------------------------------------
% ELEMENTOS PÓS-TEXTUAIS
% ----------------------------------------------------------
\postextual
% ----------------------------------------------------------


% ----------------------------------------------------------
% REFERÊNCIAS
% ----------------------------------------------------------
\bibliography{references.bib}


% ----------------------------------------------------------
% GLOSSÁRIO
% ----------------------------------------------------------
%
% Consulte o manual da classe abntex2 para orientações sobre o glossário.
%
%\glossary

% ----------------------------------------------------------
% APÊNDICES
% ----------------------------------------------------------

% ---
% Inicia os apêndices
% ---
% \begin{apendicesenv}

% % Imprime uma página indicando o início dos apêndices
% \partapendices

% % ----------------------------------------------------------
% \chapter{Quisque libero justo}
% % ----------------------------------------------------------

% % ----------------------------------------------------------
% \chapter{Nullam elementum urna}
% % ----------------------------------------------------------

% \end{apendicesenv}
% ---


% ----------------------------------------------------------
% ANEXOS
% ----------------------------------------------------------

% ---
% Inicia os anexos
% ---
% \begin{anexosenv}

% % Imprime uma página indicando o início dos anexos
% \partanexos

% % ---
% \chapter{Morbi ultrices rutrum lorem}
% % ---

% % ---
% \chapter{Cras non urna sed feugiat cum}
% % ---


% % ---
% \chapter{Fusce facilisis lacinia dui}
% % ---

% \end{anexosenv}

%---------------------------------------------------------------------
% INDICE REMISSIVO
%---------------------------------------------------------------------
% \phantompart
% \printindex
%---------------------------------------------------------------------
\includepdf{./05-extras/contracapa.pdf}
\end{document}


\documentclass[
	% Opções da classe memoir
	12pt,				        % tamanho da fonte
	openany,			        % capítulos começam em pág ímpar (insere página vazia caso preciso)
	oneside,			        % para impressão em recto e verso. Oposto a oneside
	a4paper,			        % tamanho do papel. 
	% Opções da classe abntex2
	%chapter=TITLE,		        % títulos de capítulos convertidos em letras maiúsculas
	%section=TITLE,		        % títulos de seções convertidos em letras maiúsculas
	%subsection=TITLE,	        % títulos de subseções convertidos em letras maiúsculas
	%subsubsection=TITLE,       % títulos de subsubseções convertidos em letras maiúsculas
	% Opções do pacote babel
	english,			        % idioma adicional para hifenização
	brazil				        % o último idioma é o principal do documento
	]{abntex2}

% ---
% Pacotes básicos 
% ---
\usepackage{lmodern}			% Usa a fonte Latin Modern			
\usepackage[T1]{fontenc}		% Selecao de codigos de fonte.
\usepackage[utf8]{inputenc}		% Codificacao do documento (conversão automática dos acentos)
\usepackage{indentfirst}		% Indenta o primeiro parágrafo de cada seção.
\usepackage{color}				% Controle das cores
\usepackage{graphicx}			% Inclusão de gráficos
\usepackage{microtype} 			% para melhorias de justificação
% ---
		

% ---
% Pacotes de citações
% ---
% \usepackage[alf]{abntex2cite}						% Citações padrão ABNT
\usepackage[
  alf,
  abnt-nbr10520=1988,                     % Citações em sentence case
  abnt-emphasize=bf
  ]{abntex2cite}


% ---
% Pacotes adicionais
% ---
\usepackage{pdfpages}	% Inclusão de páginas em PDF


% ---
% Informações para CAPA e FOLHA DE ROSTO
% ---
\titulo{APRENDIZAGEM DE MÁQUINA PARA PREDIÇÃO DE DESFECHOS EM SAÚDE MENTAL: \textmd{UMA INTRODUÇÃO PARA PSICÓLOGOS}}
\autor{GUILHERME AUGUSTO DIAS}
\local{Porto Alegre}
\data{2024}
\orientador{Prof. Dr. Wagner de Lara Machado}
% \coorientador{}
% \instituicao{
%   Pontifícia Universidade Católica do Rio Grande do Sul
%   \par
%   Escola de Ciências da Saúde e da Vida
%   \par
%   Curso de Psicologia}
\tipotrabalho{TCC}
\preambulo{Trabalho de Conclusão de Curso apresentado como requisito para a obtenção do grau de Bacharel em Psicologia pela Escola de Ciências da Saúde e da Vida da Pontifícia Universidade Católica do Rio Grande do Sul.}
% ---

% ---
% Configurações de aparência do PDF final
% ---

% Alterando o aspecto da cor azul
\definecolor{blue}{RGB}{41,5,195}

% Informações do PDF
\makeatletter
\hypersetup{
		pdftitle={Aprendizagem de máquina para predição de desfechos em saúde mental: uma introdução para psicólogos}, 
		pdfauthor={Guilherme Augusto Dias},
    	pdfsubject={\imprimirpreambulo},
	    pdfcreator={LaTeX with abnTeX2},
		pdfkeywords={aprendizagem de máquina}{psicoterapia}{psicologia positiva}{depressão}, 
		colorlinks=true,		% true: colored links; false: boxed links
    	linkcolor=black,		% color of internal links
    	citecolor=black,		% color of links to bibliography
    	filecolor=black,		% color of file links
		urlcolor=black,
		bookmarksdepth=4
}
\makeatother
% --- 


% ---
% Posiciona figuras e tabelas no topo da página quando adicionadas sozinhas
% em um página em branco. Ver https://github.com/abntex/abntex2/issues/170
\makeatletter
\setlength{\@fptop}{5pt} % Set distance from top of page to first float
\makeatother
% ---

% ---
% Possibilita criação de Quadros e Lista de quadros.
% Ver https://github.com/abntex/abntex2/issues/176
%
\newcommand{\quadroname}{Quadro}
\newcommand{\listofquadrosname}{Lista de quadros}

\newfloat[chapter]{quadro}{loq}{\quadroname}
\newlistof{listofquadros}{loq}{\listofquadrosname}
\newlistentry{quadro}{loq}{0}

% Configurações para atender às regras da ABNT
\setfloatadjustment{quadro}{\centering}
\counterwithout{quadro}{chapter}
\renewcommand{\cftquadroname}{\quadroname\space} 
\renewcommand*{\cftquadroaftersnum}{\hfill--\hfill}

\setfloatlocations{quadro}{hbtp} % Ver https://github.com/abntex/abntex2/issues/176
% ---


% --- 
% Espaçamentos entre linhas e parágrafos 
% --- 

% O tamanho do parágrafo é dado por:
\setlength{\parindent}{1.3cm}

% Controle do espaçamento entre um parágrafo e outro:
\setlength{\parskip}{0.2cm}  % tente também \onelineskip

% ---
% Compila o índice
% ---
\makeindex
% ---


% ----
% Início do documento
% ----
\begin{document}

% Seleciona o idioma do documento (conforme pacotes do babel)
\selectlanguage{brazil}

% Retira espaço extra obsoleto entre as frases.
\frenchspacing 

% ----------------------------------------------------------
% ELEMENTOS PRÉ-TEXTUAIS
% ----------------------------------------------------------
\pretextual

% ---
% CAPA
% ---
\includepdf{./05-extras/capa.pdf}
% ---


% ---
% FOLHA DE ROSTO
% (o * indica que haverá a ficha bibliográfica)
% ---
\imprimirfolhaderosto
% ---


% ---
% Ficha bibliografica
% ---

% Isto é um exemplo de Ficha Catalográfica, ou ``Dados internacionais de
% catalogação-na-publicação''. Você pode utilizar este modelo como referência. 
% Porém, provavelmente a biblioteca da sua universidade lhe fornecerá um PDF
% com a ficha catalográfica definitiva após a defesa do trabalho. Quando estiver
% com o documento, salve-o como PDF no diretório do seu projeto e substitua todo
% o conteúdo de implementação deste arquivo pelo comando abaixo:
%
% \usepackage{pdfpages}		% necessário para comando \includepdf
% \begin{fichacatalografica}
%     \includepdf{fig_ficha_catalografica.pdf}
% \end{fichacatalografica}

% \begin{fichacatalografica}
% 	\sffamily
% 	\vspace*{\fill}                             % Posição vertical
% 	\begin{center}                              % Minipage Centralizado
% 	\fbox{\begin{minipage}[c][8cm]{13.5cm}		% Largura
% 	\small
% 	\imprimirautor
% 	%Sobrenome, Nome do autor
	
% 	\hspace{0.5cm} \imprimirtitulo  / \imprimirautor. --
% 	\imprimirlocal, \imprimirdata-
	
% 	\hspace{0.5cm} \thelastpage p. : il. (algumas color.) ; 30 cm.\\
	
% 	\hspace{0.5cm} \imprimirorientadorRotulo~\imprimirorientador\\
	
% 	\hspace{0.5cm}
% 	\parbox[t]{\textwidth}{\imprimirtipotrabalho~--~\imprimirinstituicao,
% 	\imprimirdata.}\\
	
% 	\hspace{0.5cm}
% 		1. Palavra-chave1.
% 		2. Palavra-chave2.
% 		2. Palavra-chave3.
% 		I. Orientador.
% 		II. Universidade xxx.
% 		III. Faculdade de xxx.
% 		IV. Título 			
% 	\end{minipage}}
% 	\end{center}
% \end{fichacatalografica}
% ---


% ---
% FOLHA DE APROVAÇÃO
% ---

% Isto é um exemplo de Folha de aprovação, elemento obrigatório da NBR
% 14724/2011 (seção 4.2.1.3). Você pode utilizar este modelo até a aprovação
% do trabalho. Após isso, substitua todo o conteúdo deste arquivo por uma
% imagem da página assinada pela banca com o comando abaixo:
%
% \begin{folhadeaprovacao}
% \includepdf{folhadeaprovacao_final.pdf}
% \end{folhadeaprovacao}
%
\setlength{\ABNTEXsignskip}{2cm}
\setlength{\ABNTEXsignwidth}{10cm}

\begin{folhadeaprovacao}
  \begin{center}
    {\ABNTEXchapterfont\large\imprimirautor}

    \vspace*{1cm}

    \begin{center}
      \ABNTEXchapterfont\bfseries\Large\imprimirtitulo
    \end{center}

    \vspace*{1cm}
    
    \hspace{.45\textwidth}
    \begin{minipage}{.5\textwidth}
        \imprimirpreambulo
    \end{minipage}%
   \end{center}

   \vspace*{1cm}
        
   \begin{center}
   Aprovado em 17 de julho de 2024.

   \vspace*{1cm}

   Banca examinadora:
   \end{center}

   \vspace*{-1cm}

   \assinatura{Prof. Ma. Caroline Santa Maria Rodrigues -- PUCRS}
   \assinatura{Prof. Dra. Clarissa Pinto Pizarro de Freitas -- PUCRS}
   \assinatura{\imprimirorientador -- PUCRS}
\end{folhadeaprovacao}
% ---


% ---
% Agradecimentos
% ---
\begin{agradecimentos}
Agradeço aos meus pais, Ademilson José Dias e Irinês Porfírio Dias, pelo apoio incondicional que me deram ao longo da vida,
incentivando mesmo as ideias questionáveis como aquela de fazer uma segunda graduação.

A minha irmã, Maria Eduarda Dias, cuja trajetória acadêmica é fonte de inspiração.

A minha namorada, Nathalia Fallavena Ceratti, que acompanhou a esta jornada de perto, oferecendo apoio e compreensão que
fizeram toda a diferença.

Aos meus amigos, Guilherme da Silva Petró, Norton Neves e Uilinton Luciano Marques da Silva, que foram compreensíveis nos
finais de semana em que não pude jogar RPG.

Ao professor Wagner de Lara Machado pela orientação e apoio no desenvolvimento deste trabalho.
\end{agradecimentos}
% ---


% ---
% RESUMOS
% ---
% Resumo em português
\setlength{\absparsep}{18pt} % ajusta o espaçamento dos parágrafos do resumo
\begin{resumo}
Desenvolvimentos na área de aprendizagem de máquina representam uma série de oportunidades para diversas áreas de pesquisa e de atuação profissional.
O uso de modelos de aprendizagem de máquina na predição de desfecho para tratamentos em saúde mental tem sido investigado e apresenta resultados promissores.
A aplicação desse tipo de modelo possibilitaria, por exemplo, maior assertividade na personalização de tratamentos psicoterápicos a nível individual.
Apesar dos benefícios potenciais, o tema ainda é pouco discutido nos espaços de pesquisa, formação e atuação em psicologia clínica. Este trabalho tem
por objetivo introduzir o tema de aprendizagem de máquina a psicólogos. Busca-se apresentar, de maneira acessível, os conceitos e terminologia básicos
em aprendizagem de máquina, proporcionando um panorama geral da área. Para complementar a exposição conceitual, é construído um exemplo a partir de dados
de um estudo sobre eficácia de intervenções em psicologia positiva para depressão via internet. Por fim, discutem-se algumas implicações do uso desse tipo
de tecnologia para a pesquisa e para a prática em psicoterapia.

 \textbf{Palavras-chave}: aprendizagem de máquina, psicoterapia, psicologia positiva, depressão.
\end{resumo}

% Resumo em inglês
\begin{resumo}[Abstract]
 \begin{otherlanguage*}{english}
   Advances in machine learning have opened up a multitude of opportunities for various research and professional fields. The use of machine learning models to
   predict treatment outcomes in mental health has been investigated showing promising results. The application of such models would allow greater assertiveness
   in tailoring psychotherapeutic treatments at the individual level. Despite its potential benefits, the topic remains unexplored in the fields of research, training
   and practice in clinical psychology. This work aims to introduce the topic of machine learning to psychologists. It presents the basic machine learning concepts and
   terminology in an accessible manner, providing a comprehensive overview of the field. To complement the conceptual discussion, an example is constructed using data
   from a study on the effectiveness of internet-based positive psychology interventions for depression. Finally, implications of using this type of technology in
   psychotherapy research and practice are discussed.
   \vspace{\onelineskip}
 
   \noindent 
   \textbf{Keywords}: machine learning, psychotherapy, positive psychology, depression.
 \end{otherlanguage*}
\end{resumo}


% ---
% LISTA DE ILUSTRAÇÕES
% ---
\pdfbookmark[0]{\listfigurename}{lof}
\listoffigures*
\cleardoublepage
% ---


% ---
% LISTA DE QUADROS
% ---
% \pdfbookmark[0]{\listofquadrosname}{loq}
% \listofquadros*
% \cleardoublepage
% ---


% ---
% LISTA DE TABELAS
% ---
% \pdfbookmark[0]{\listtablename}{lot}
% \listoftables*
% \cleardoublepage
% ---


% ---
% LISTA DE ABREVIATURAS E SIGLAS
% ---
% \begin{siglas}
%   \item[ABNT] Associação Brasileira de Normas Técnicas
%   \item[abnTeX] ABsurdas Normas para TeX
% \end{siglas}
% ---


% ---
% LISTA DE SÍMBOLOS
% ---
% \begin{simbolos}
%   \item[$ \Gamma $] Letra grega Gama
%   \item[$ \Lambda $] Lambda
%   \item[$ \zeta $] Letra grega minúscula zeta
%   \item[$ \in $] Pertence
% \end{simbolos}
% ---


% ---
% Sumário
% ---
\pdfbookmark[0]{\contentsname}{toc}
\tableofcontents*
\cleardoublepage
% ---


% ----------------------------------------------------------
% ELEMENTOS TEXTUAIS
% ----------------------------------------------------------
\textual

\include{./01-introducao/main.tex}

\include{./02-machine-learning/main.tex}

\include{./03-exemplo/main.tex}

\include{./04-conclusao/main.tex}

% ----------------------------------------------------------
% ELEMENTOS PÓS-TEXTUAIS
% ----------------------------------------------------------
\postextual
% ----------------------------------------------------------


% ----------------------------------------------------------
% REFERÊNCIAS
% ----------------------------------------------------------
\bibliography{references.bib}


% ----------------------------------------------------------
% GLOSSÁRIO
% ----------------------------------------------------------
%
% Consulte o manual da classe abntex2 para orientações sobre o glossário.
%
%\glossary

% ----------------------------------------------------------
% APÊNDICES
% ----------------------------------------------------------

% ---
% Inicia os apêndices
% ---
% \begin{apendicesenv}

% % Imprime uma página indicando o início dos apêndices
% \partapendices

% % ----------------------------------------------------------
% \chapter{Quisque libero justo}
% % ----------------------------------------------------------

% % ----------------------------------------------------------
% \chapter{Nullam elementum urna}
% % ----------------------------------------------------------

% \end{apendicesenv}
% ---


% ----------------------------------------------------------
% ANEXOS
% ----------------------------------------------------------

% ---
% Inicia os anexos
% ---
% \begin{anexosenv}

% % Imprime uma página indicando o início dos anexos
% \partanexos

% % ---
% \chapter{Morbi ultrices rutrum lorem}
% % ---

% % ---
% \chapter{Cras non urna sed feugiat cum}
% % ---


% % ---
% \chapter{Fusce facilisis lacinia dui}
% % ---

% \end{anexosenv}

%---------------------------------------------------------------------
% INDICE REMISSIVO
%---------------------------------------------------------------------
% \phantompart
% \printindex
%---------------------------------------------------------------------
\includepdf{./05-extras/contracapa.pdf}
\end{document}


\documentclass[
	% Opções da classe memoir
	12pt,				        % tamanho da fonte
	openany,			        % capítulos começam em pág ímpar (insere página vazia caso preciso)
	oneside,			        % para impressão em recto e verso. Oposto a oneside
	a4paper,			        % tamanho do papel. 
	% Opções da classe abntex2
	%chapter=TITLE,		        % títulos de capítulos convertidos em letras maiúsculas
	%section=TITLE,		        % títulos de seções convertidos em letras maiúsculas
	%subsection=TITLE,	        % títulos de subseções convertidos em letras maiúsculas
	%subsubsection=TITLE,       % títulos de subsubseções convertidos em letras maiúsculas
	% Opções do pacote babel
	english,			        % idioma adicional para hifenização
	brazil				        % o último idioma é o principal do documento
	]{abntex2}

% ---
% Pacotes básicos 
% ---
\usepackage{lmodern}			% Usa a fonte Latin Modern			
\usepackage[T1]{fontenc}		% Selecao de codigos de fonte.
\usepackage[utf8]{inputenc}		% Codificacao do documento (conversão automática dos acentos)
\usepackage{indentfirst}		% Indenta o primeiro parágrafo de cada seção.
\usepackage{color}				% Controle das cores
\usepackage{graphicx}			% Inclusão de gráficos
\usepackage{microtype} 			% para melhorias de justificação
% ---
		

% ---
% Pacotes de citações
% ---
% \usepackage[alf]{abntex2cite}						% Citações padrão ABNT
\usepackage[
  alf,
  abnt-nbr10520=1988,                     % Citações em sentence case
  abnt-emphasize=bf
  ]{abntex2cite}


% ---
% Pacotes adicionais
% ---
\usepackage{pdfpages}	% Inclusão de páginas em PDF


% ---
% Informações para CAPA e FOLHA DE ROSTO
% ---
\titulo{APRENDIZAGEM DE MÁQUINA PARA PREDIÇÃO DE DESFECHOS EM SAÚDE MENTAL: \textmd{UMA INTRODUÇÃO PARA PSICÓLOGOS}}
\autor{GUILHERME AUGUSTO DIAS}
\local{Porto Alegre}
\data{2024}
\orientador{Prof. Dr. Wagner de Lara Machado}
% \coorientador{}
% \instituicao{
%   Pontifícia Universidade Católica do Rio Grande do Sul
%   \par
%   Escola de Ciências da Saúde e da Vida
%   \par
%   Curso de Psicologia}
\tipotrabalho{TCC}
\preambulo{Trabalho de Conclusão de Curso apresentado como requisito para a obtenção do grau de Bacharel em Psicologia pela Escola de Ciências da Saúde e da Vida da Pontifícia Universidade Católica do Rio Grande do Sul.}
% ---

% ---
% Configurações de aparência do PDF final
% ---

% Alterando o aspecto da cor azul
\definecolor{blue}{RGB}{41,5,195}

% Informações do PDF
\makeatletter
\hypersetup{
		pdftitle={Aprendizagem de máquina para predição de desfechos em saúde mental: uma introdução para psicólogos}, 
		pdfauthor={Guilherme Augusto Dias},
    	pdfsubject={\imprimirpreambulo},
	    pdfcreator={LaTeX with abnTeX2},
		pdfkeywords={aprendizagem de máquina}{psicoterapia}{psicologia positiva}{depressão}, 
		colorlinks=true,		% true: colored links; false: boxed links
    	linkcolor=black,		% color of internal links
    	citecolor=black,		% color of links to bibliography
    	filecolor=black,		% color of file links
		urlcolor=black,
		bookmarksdepth=4
}
\makeatother
% --- 


% ---
% Posiciona figuras e tabelas no topo da página quando adicionadas sozinhas
% em um página em branco. Ver https://github.com/abntex/abntex2/issues/170
\makeatletter
\setlength{\@fptop}{5pt} % Set distance from top of page to first float
\makeatother
% ---

% ---
% Possibilita criação de Quadros e Lista de quadros.
% Ver https://github.com/abntex/abntex2/issues/176
%
\newcommand{\quadroname}{Quadro}
\newcommand{\listofquadrosname}{Lista de quadros}

\newfloat[chapter]{quadro}{loq}{\quadroname}
\newlistof{listofquadros}{loq}{\listofquadrosname}
\newlistentry{quadro}{loq}{0}

% Configurações para atender às regras da ABNT
\setfloatadjustment{quadro}{\centering}
\counterwithout{quadro}{chapter}
\renewcommand{\cftquadroname}{\quadroname\space} 
\renewcommand*{\cftquadroaftersnum}{\hfill--\hfill}

\setfloatlocations{quadro}{hbtp} % Ver https://github.com/abntex/abntex2/issues/176
% ---


% --- 
% Espaçamentos entre linhas e parágrafos 
% --- 

% O tamanho do parágrafo é dado por:
\setlength{\parindent}{1.3cm}

% Controle do espaçamento entre um parágrafo e outro:
\setlength{\parskip}{0.2cm}  % tente também \onelineskip

% ---
% Compila o índice
% ---
\makeindex
% ---


% ----
% Início do documento
% ----
\begin{document}

% Seleciona o idioma do documento (conforme pacotes do babel)
\selectlanguage{brazil}

% Retira espaço extra obsoleto entre as frases.
\frenchspacing 

% ----------------------------------------------------------
% ELEMENTOS PRÉ-TEXTUAIS
% ----------------------------------------------------------
\pretextual

% ---
% CAPA
% ---
\includepdf{./05-extras/capa.pdf}
% ---


% ---
% FOLHA DE ROSTO
% (o * indica que haverá a ficha bibliográfica)
% ---
\imprimirfolhaderosto
% ---


% ---
% Ficha bibliografica
% ---

% Isto é um exemplo de Ficha Catalográfica, ou ``Dados internacionais de
% catalogação-na-publicação''. Você pode utilizar este modelo como referência. 
% Porém, provavelmente a biblioteca da sua universidade lhe fornecerá um PDF
% com a ficha catalográfica definitiva após a defesa do trabalho. Quando estiver
% com o documento, salve-o como PDF no diretório do seu projeto e substitua todo
% o conteúdo de implementação deste arquivo pelo comando abaixo:
%
% \usepackage{pdfpages}		% necessário para comando \includepdf
% \begin{fichacatalografica}
%     \includepdf{fig_ficha_catalografica.pdf}
% \end{fichacatalografica}

% \begin{fichacatalografica}
% 	\sffamily
% 	\vspace*{\fill}                             % Posição vertical
% 	\begin{center}                              % Minipage Centralizado
% 	\fbox{\begin{minipage}[c][8cm]{13.5cm}		% Largura
% 	\small
% 	\imprimirautor
% 	%Sobrenome, Nome do autor
	
% 	\hspace{0.5cm} \imprimirtitulo  / \imprimirautor. --
% 	\imprimirlocal, \imprimirdata-
	
% 	\hspace{0.5cm} \thelastpage p. : il. (algumas color.) ; 30 cm.\\
	
% 	\hspace{0.5cm} \imprimirorientadorRotulo~\imprimirorientador\\
	
% 	\hspace{0.5cm}
% 	\parbox[t]{\textwidth}{\imprimirtipotrabalho~--~\imprimirinstituicao,
% 	\imprimirdata.}\\
	
% 	\hspace{0.5cm}
% 		1. Palavra-chave1.
% 		2. Palavra-chave2.
% 		2. Palavra-chave3.
% 		I. Orientador.
% 		II. Universidade xxx.
% 		III. Faculdade de xxx.
% 		IV. Título 			
% 	\end{minipage}}
% 	\end{center}
% \end{fichacatalografica}
% ---


% ---
% FOLHA DE APROVAÇÃO
% ---

% Isto é um exemplo de Folha de aprovação, elemento obrigatório da NBR
% 14724/2011 (seção 4.2.1.3). Você pode utilizar este modelo até a aprovação
% do trabalho. Após isso, substitua todo o conteúdo deste arquivo por uma
% imagem da página assinada pela banca com o comando abaixo:
%
% \begin{folhadeaprovacao}
% \includepdf{folhadeaprovacao_final.pdf}
% \end{folhadeaprovacao}
%
\setlength{\ABNTEXsignskip}{2cm}
\setlength{\ABNTEXsignwidth}{10cm}

\begin{folhadeaprovacao}
  \begin{center}
    {\ABNTEXchapterfont\large\imprimirautor}

    \vspace*{1cm}

    \begin{center}
      \ABNTEXchapterfont\bfseries\Large\imprimirtitulo
    \end{center}

    \vspace*{1cm}
    
    \hspace{.45\textwidth}
    \begin{minipage}{.5\textwidth}
        \imprimirpreambulo
    \end{minipage}%
   \end{center}

   \vspace*{1cm}
        
   \begin{center}
   Aprovado em 17 de julho de 2024.

   \vspace*{1cm}

   Banca examinadora:
   \end{center}

   \vspace*{-1cm}

   \assinatura{Prof. Ma. Caroline Santa Maria Rodrigues -- PUCRS}
   \assinatura{Prof. Dra. Clarissa Pinto Pizarro de Freitas -- PUCRS}
   \assinatura{\imprimirorientador -- PUCRS}
\end{folhadeaprovacao}
% ---


% ---
% Agradecimentos
% ---
\begin{agradecimentos}
Agradeço aos meus pais, Ademilson José Dias e Irinês Porfírio Dias, pelo apoio incondicional que me deram ao longo da vida,
incentivando mesmo as ideias questionáveis como aquela de fazer uma segunda graduação.

A minha irmã, Maria Eduarda Dias, cuja trajetória acadêmica é fonte de inspiração.

A minha namorada, Nathalia Fallavena Ceratti, que acompanhou a esta jornada de perto, oferecendo apoio e compreensão que
fizeram toda a diferença.

Aos meus amigos, Guilherme da Silva Petró, Norton Neves e Uilinton Luciano Marques da Silva, que foram compreensíveis nos
finais de semana em que não pude jogar RPG.

Ao professor Wagner de Lara Machado pela orientação e apoio no desenvolvimento deste trabalho.
\end{agradecimentos}
% ---


% ---
% RESUMOS
% ---
% Resumo em português
\setlength{\absparsep}{18pt} % ajusta o espaçamento dos parágrafos do resumo
\begin{resumo}
Desenvolvimentos na área de aprendizagem de máquina representam uma série de oportunidades para diversas áreas de pesquisa e de atuação profissional.
O uso de modelos de aprendizagem de máquina na predição de desfecho para tratamentos em saúde mental tem sido investigado e apresenta resultados promissores.
A aplicação desse tipo de modelo possibilitaria, por exemplo, maior assertividade na personalização de tratamentos psicoterápicos a nível individual.
Apesar dos benefícios potenciais, o tema ainda é pouco discutido nos espaços de pesquisa, formação e atuação em psicologia clínica. Este trabalho tem
por objetivo introduzir o tema de aprendizagem de máquina a psicólogos. Busca-se apresentar, de maneira acessível, os conceitos e terminologia básicos
em aprendizagem de máquina, proporcionando um panorama geral da área. Para complementar a exposição conceitual, é construído um exemplo a partir de dados
de um estudo sobre eficácia de intervenções em psicologia positiva para depressão via internet. Por fim, discutem-se algumas implicações do uso desse tipo
de tecnologia para a pesquisa e para a prática em psicoterapia.

 \textbf{Palavras-chave}: aprendizagem de máquina, psicoterapia, psicologia positiva, depressão.
\end{resumo}

% Resumo em inglês
\begin{resumo}[Abstract]
 \begin{otherlanguage*}{english}
   Advances in machine learning have opened up a multitude of opportunities for various research and professional fields. The use of machine learning models to
   predict treatment outcomes in mental health has been investigated showing promising results. The application of such models would allow greater assertiveness
   in tailoring psychotherapeutic treatments at the individual level. Despite its potential benefits, the topic remains unexplored in the fields of research, training
   and practice in clinical psychology. This work aims to introduce the topic of machine learning to psychologists. It presents the basic machine learning concepts and
   terminology in an accessible manner, providing a comprehensive overview of the field. To complement the conceptual discussion, an example is constructed using data
   from a study on the effectiveness of internet-based positive psychology interventions for depression. Finally, implications of using this type of technology in
   psychotherapy research and practice are discussed.
   \vspace{\onelineskip}
 
   \noindent 
   \textbf{Keywords}: machine learning, psychotherapy, positive psychology, depression.
 \end{otherlanguage*}
\end{resumo}


% ---
% LISTA DE ILUSTRAÇÕES
% ---
\pdfbookmark[0]{\listfigurename}{lof}
\listoffigures*
\cleardoublepage
% ---


% ---
% LISTA DE QUADROS
% ---
% \pdfbookmark[0]{\listofquadrosname}{loq}
% \listofquadros*
% \cleardoublepage
% ---


% ---
% LISTA DE TABELAS
% ---
% \pdfbookmark[0]{\listtablename}{lot}
% \listoftables*
% \cleardoublepage
% ---


% ---
% LISTA DE ABREVIATURAS E SIGLAS
% ---
% \begin{siglas}
%   \item[ABNT] Associação Brasileira de Normas Técnicas
%   \item[abnTeX] ABsurdas Normas para TeX
% \end{siglas}
% ---


% ---
% LISTA DE SÍMBOLOS
% ---
% \begin{simbolos}
%   \item[$ \Gamma $] Letra grega Gama
%   \item[$ \Lambda $] Lambda
%   \item[$ \zeta $] Letra grega minúscula zeta
%   \item[$ \in $] Pertence
% \end{simbolos}
% ---


% ---
% Sumário
% ---
\pdfbookmark[0]{\contentsname}{toc}
\tableofcontents*
\cleardoublepage
% ---


% ----------------------------------------------------------
% ELEMENTOS TEXTUAIS
% ----------------------------------------------------------
\textual

\include{./01-introducao/main.tex}

\include{./02-machine-learning/main.tex}

\include{./03-exemplo/main.tex}

\include{./04-conclusao/main.tex}

% ----------------------------------------------------------
% ELEMENTOS PÓS-TEXTUAIS
% ----------------------------------------------------------
\postextual
% ----------------------------------------------------------


% ----------------------------------------------------------
% REFERÊNCIAS
% ----------------------------------------------------------
\bibliography{references.bib}


% ----------------------------------------------------------
% GLOSSÁRIO
% ----------------------------------------------------------
%
% Consulte o manual da classe abntex2 para orientações sobre o glossário.
%
%\glossary

% ----------------------------------------------------------
% APÊNDICES
% ----------------------------------------------------------

% ---
% Inicia os apêndices
% ---
% \begin{apendicesenv}

% % Imprime uma página indicando o início dos apêndices
% \partapendices

% % ----------------------------------------------------------
% \chapter{Quisque libero justo}
% % ----------------------------------------------------------

% % ----------------------------------------------------------
% \chapter{Nullam elementum urna}
% % ----------------------------------------------------------

% \end{apendicesenv}
% ---


% ----------------------------------------------------------
% ANEXOS
% ----------------------------------------------------------

% ---
% Inicia os anexos
% ---
% \begin{anexosenv}

% % Imprime uma página indicando o início dos anexos
% \partanexos

% % ---
% \chapter{Morbi ultrices rutrum lorem}
% % ---

% % ---
% \chapter{Cras non urna sed feugiat cum}
% % ---


% % ---
% \chapter{Fusce facilisis lacinia dui}
% % ---

% \end{anexosenv}

%---------------------------------------------------------------------
% INDICE REMISSIVO
%---------------------------------------------------------------------
% \phantompart
% \printindex
%---------------------------------------------------------------------
\includepdf{./05-extras/contracapa.pdf}
\end{document}


% ----------------------------------------------------------
% ELEMENTOS PÓS-TEXTUAIS
% ----------------------------------------------------------
\postextual
% ----------------------------------------------------------


% ----------------------------------------------------------
% REFERÊNCIAS
% ----------------------------------------------------------
\bibliography{references.bib}


% ----------------------------------------------------------
% GLOSSÁRIO
% ----------------------------------------------------------
%
% Consulte o manual da classe abntex2 para orientações sobre o glossário.
%
%\glossary

% ----------------------------------------------------------
% APÊNDICES
% ----------------------------------------------------------

% ---
% Inicia os apêndices
% ---
% \begin{apendicesenv}

% % Imprime uma página indicando o início dos apêndices
% \partapendices

% % ----------------------------------------------------------
% \chapter{Quisque libero justo}
% % ----------------------------------------------------------

% % ----------------------------------------------------------
% \chapter{Nullam elementum urna}
% % ----------------------------------------------------------

% \end{apendicesenv}
% ---


% ----------------------------------------------------------
% ANEXOS
% ----------------------------------------------------------

% ---
% Inicia os anexos
% ---
% \begin{anexosenv}

% % Imprime uma página indicando o início dos anexos
% \partanexos

% % ---
% \chapter{Morbi ultrices rutrum lorem}
% % ---

% % ---
% \chapter{Cras non urna sed feugiat cum}
% % ---


% % ---
% \chapter{Fusce facilisis lacinia dui}
% % ---

% \end{anexosenv}

%---------------------------------------------------------------------
% INDICE REMISSIVO
%---------------------------------------------------------------------
% \phantompart
% \printindex
%---------------------------------------------------------------------
\includepdf{./05-extras/contracapa.pdf}
\end{document}


\documentclass[
	% Opções da classe memoir
	12pt,				        % tamanho da fonte
	openany,			        % capítulos começam em pág ímpar (insere página vazia caso preciso)
	oneside,			        % para impressão em recto e verso. Oposto a oneside
	a4paper,			        % tamanho do papel. 
	% Opções da classe abntex2
	%chapter=TITLE,		        % títulos de capítulos convertidos em letras maiúsculas
	%section=TITLE,		        % títulos de seções convertidos em letras maiúsculas
	%subsection=TITLE,	        % títulos de subseções convertidos em letras maiúsculas
	%subsubsection=TITLE,       % títulos de subsubseções convertidos em letras maiúsculas
	% Opções do pacote babel
	english,			        % idioma adicional para hifenização
	brazil				        % o último idioma é o principal do documento
	]{abntex2}

% ---
% Pacotes básicos 
% ---
\usepackage{lmodern}			% Usa a fonte Latin Modern			
\usepackage[T1]{fontenc}		% Selecao de codigos de fonte.
\usepackage[utf8]{inputenc}		% Codificacao do documento (conversão automática dos acentos)
\usepackage{indentfirst}		% Indenta o primeiro parágrafo de cada seção.
\usepackage{color}				% Controle das cores
\usepackage{graphicx}			% Inclusão de gráficos
\usepackage{microtype} 			% para melhorias de justificação
% ---
		

% ---
% Pacotes de citações
% ---
% \usepackage[alf]{abntex2cite}						% Citações padrão ABNT
\usepackage[
  alf,
  abnt-nbr10520=1988,                     % Citações em sentence case
  abnt-emphasize=bf
  ]{abntex2cite}


% ---
% Pacotes adicionais
% ---
\usepackage{pdfpages}	% Inclusão de páginas em PDF


% ---
% Informações para CAPA e FOLHA DE ROSTO
% ---
\titulo{APRENDIZAGEM DE MÁQUINA PARA PREDIÇÃO DE DESFECHOS EM SAÚDE MENTAL: \textmd{UMA INTRODUÇÃO PARA PSICÓLOGOS}}
\autor{GUILHERME AUGUSTO DIAS}
\local{Porto Alegre}
\data{2024}
\orientador{Prof. Dr. Wagner de Lara Machado}
% \coorientador{}
% \instituicao{
%   Pontifícia Universidade Católica do Rio Grande do Sul
%   \par
%   Escola de Ciências da Saúde e da Vida
%   \par
%   Curso de Psicologia}
\tipotrabalho{TCC}
\preambulo{Trabalho de Conclusão de Curso apresentado como requisito para a obtenção do grau de Bacharel em Psicologia pela Escola de Ciências da Saúde e da Vida da Pontifícia Universidade Católica do Rio Grande do Sul.}
% ---

% ---
% Configurações de aparência do PDF final
% ---

% Alterando o aspecto da cor azul
\definecolor{blue}{RGB}{41,5,195}

% Informações do PDF
\makeatletter
\hypersetup{
		pdftitle={Aprendizagem de máquina para predição de desfechos em saúde mental: uma introdução para psicólogos}, 
		pdfauthor={Guilherme Augusto Dias},
    	pdfsubject={\imprimirpreambulo},
	    pdfcreator={LaTeX with abnTeX2},
		pdfkeywords={aprendizagem de máquina}{psicoterapia}{psicologia positiva}{depressão}, 
		colorlinks=true,		% true: colored links; false: boxed links
    	linkcolor=black,		% color of internal links
    	citecolor=black,		% color of links to bibliography
    	filecolor=black,		% color of file links
		urlcolor=black,
		bookmarksdepth=4
}
\makeatother
% --- 


% ---
% Posiciona figuras e tabelas no topo da página quando adicionadas sozinhas
% em um página em branco. Ver https://github.com/abntex/abntex2/issues/170
\makeatletter
\setlength{\@fptop}{5pt} % Set distance from top of page to first float
\makeatother
% ---

% ---
% Possibilita criação de Quadros e Lista de quadros.
% Ver https://github.com/abntex/abntex2/issues/176
%
\newcommand{\quadroname}{Quadro}
\newcommand{\listofquadrosname}{Lista de quadros}

\newfloat[chapter]{quadro}{loq}{\quadroname}
\newlistof{listofquadros}{loq}{\listofquadrosname}
\newlistentry{quadro}{loq}{0}

% Configurações para atender às regras da ABNT
\setfloatadjustment{quadro}{\centering}
\counterwithout{quadro}{chapter}
\renewcommand{\cftquadroname}{\quadroname\space} 
\renewcommand*{\cftquadroaftersnum}{\hfill--\hfill}

\setfloatlocations{quadro}{hbtp} % Ver https://github.com/abntex/abntex2/issues/176
% ---


% --- 
% Espaçamentos entre linhas e parágrafos 
% --- 

% O tamanho do parágrafo é dado por:
\setlength{\parindent}{1.3cm}

% Controle do espaçamento entre um parágrafo e outro:
\setlength{\parskip}{0.2cm}  % tente também \onelineskip

% ---
% Compila o índice
% ---
\makeindex
% ---


% ----
% Início do documento
% ----
\begin{document}

% Seleciona o idioma do documento (conforme pacotes do babel)
\selectlanguage{brazil}

% Retira espaço extra obsoleto entre as frases.
\frenchspacing 

% ----------------------------------------------------------
% ELEMENTOS PRÉ-TEXTUAIS
% ----------------------------------------------------------
\pretextual

% ---
% CAPA
% ---
\includepdf{./05-extras/capa.pdf}
% ---


% ---
% FOLHA DE ROSTO
% (o * indica que haverá a ficha bibliográfica)
% ---
\imprimirfolhaderosto
% ---


% ---
% Ficha bibliografica
% ---

% Isto é um exemplo de Ficha Catalográfica, ou ``Dados internacionais de
% catalogação-na-publicação''. Você pode utilizar este modelo como referência. 
% Porém, provavelmente a biblioteca da sua universidade lhe fornecerá um PDF
% com a ficha catalográfica definitiva após a defesa do trabalho. Quando estiver
% com o documento, salve-o como PDF no diretório do seu projeto e substitua todo
% o conteúdo de implementação deste arquivo pelo comando abaixo:
%
% \usepackage{pdfpages}		% necessário para comando \includepdf
% \begin{fichacatalografica}
%     \includepdf{fig_ficha_catalografica.pdf}
% \end{fichacatalografica}

% \begin{fichacatalografica}
% 	\sffamily
% 	\vspace*{\fill}                             % Posição vertical
% 	\begin{center}                              % Minipage Centralizado
% 	\fbox{\begin{minipage}[c][8cm]{13.5cm}		% Largura
% 	\small
% 	\imprimirautor
% 	%Sobrenome, Nome do autor
	
% 	\hspace{0.5cm} \imprimirtitulo  / \imprimirautor. --
% 	\imprimirlocal, \imprimirdata-
	
% 	\hspace{0.5cm} \thelastpage p. : il. (algumas color.) ; 30 cm.\\
	
% 	\hspace{0.5cm} \imprimirorientadorRotulo~\imprimirorientador\\
	
% 	\hspace{0.5cm}
% 	\parbox[t]{\textwidth}{\imprimirtipotrabalho~--~\imprimirinstituicao,
% 	\imprimirdata.}\\
	
% 	\hspace{0.5cm}
% 		1. Palavra-chave1.
% 		2. Palavra-chave2.
% 		2. Palavra-chave3.
% 		I. Orientador.
% 		II. Universidade xxx.
% 		III. Faculdade de xxx.
% 		IV. Título 			
% 	\end{minipage}}
% 	\end{center}
% \end{fichacatalografica}
% ---


% ---
% FOLHA DE APROVAÇÃO
% ---

% Isto é um exemplo de Folha de aprovação, elemento obrigatório da NBR
% 14724/2011 (seção 4.2.1.3). Você pode utilizar este modelo até a aprovação
% do trabalho. Após isso, substitua todo o conteúdo deste arquivo por uma
% imagem da página assinada pela banca com o comando abaixo:
%
% \begin{folhadeaprovacao}
% \includepdf{folhadeaprovacao_final.pdf}
% \end{folhadeaprovacao}
%
\setlength{\ABNTEXsignskip}{2cm}
\setlength{\ABNTEXsignwidth}{10cm}

\begin{folhadeaprovacao}
  \begin{center}
    {\ABNTEXchapterfont\large\imprimirautor}

    \vspace*{1cm}

    \begin{center}
      \ABNTEXchapterfont\bfseries\Large\imprimirtitulo
    \end{center}

    \vspace*{1cm}
    
    \hspace{.45\textwidth}
    \begin{minipage}{.5\textwidth}
        \imprimirpreambulo
    \end{minipage}%
   \end{center}

   \vspace*{1cm}
        
   \begin{center}
   Aprovado em 17 de julho de 2024.

   \vspace*{1cm}

   Banca examinadora:
   \end{center}

   \vspace*{-1cm}

   \assinatura{Prof. Ma. Caroline Santa Maria Rodrigues -- PUCRS}
   \assinatura{Prof. Dra. Clarissa Pinto Pizarro de Freitas -- PUCRS}
   \assinatura{\imprimirorientador -- PUCRS}
\end{folhadeaprovacao}
% ---


% ---
% Agradecimentos
% ---
\begin{agradecimentos}
Agradeço aos meus pais, Ademilson José Dias e Irinês Porfírio Dias, pelo apoio incondicional que me deram ao longo da vida,
incentivando mesmo as ideias questionáveis como aquela de fazer uma segunda graduação.

A minha irmã, Maria Eduarda Dias, cuja trajetória acadêmica é fonte de inspiração.

A minha namorada, Nathalia Fallavena Ceratti, que acompanhou a esta jornada de perto, oferecendo apoio e compreensão que
fizeram toda a diferença.

Aos meus amigos, Guilherme da Silva Petró, Norton Neves e Uilinton Luciano Marques da Silva, que foram compreensíveis nos
finais de semana em que não pude jogar RPG.

Ao professor Wagner de Lara Machado pela orientação e apoio no desenvolvimento deste trabalho.
\end{agradecimentos}
% ---


% ---
% RESUMOS
% ---
% Resumo em português
\setlength{\absparsep}{18pt} % ajusta o espaçamento dos parágrafos do resumo
\begin{resumo}
Desenvolvimentos na área de aprendizagem de máquina representam uma série de oportunidades para diversas áreas de pesquisa e de atuação profissional.
O uso de modelos de aprendizagem de máquina na predição de desfecho para tratamentos em saúde mental tem sido investigado e apresenta resultados promissores.
A aplicação desse tipo de modelo possibilitaria, por exemplo, maior assertividade na personalização de tratamentos psicoterápicos a nível individual.
Apesar dos benefícios potenciais, o tema ainda é pouco discutido nos espaços de pesquisa, formação e atuação em psicologia clínica. Este trabalho tem
por objetivo introduzir o tema de aprendizagem de máquina a psicólogos. Busca-se apresentar, de maneira acessível, os conceitos e terminologia básicos
em aprendizagem de máquina, proporcionando um panorama geral da área. Para complementar a exposição conceitual, é construído um exemplo a partir de dados
de um estudo sobre eficácia de intervenções em psicologia positiva para depressão via internet. Por fim, discutem-se algumas implicações do uso desse tipo
de tecnologia para a pesquisa e para a prática em psicoterapia.

 \textbf{Palavras-chave}: aprendizagem de máquina, psicoterapia, psicologia positiva, depressão.
\end{resumo}

% Resumo em inglês
\begin{resumo}[Abstract]
 \begin{otherlanguage*}{english}
   Advances in machine learning have opened up a multitude of opportunities for various research and professional fields. The use of machine learning models to
   predict treatment outcomes in mental health has been investigated showing promising results. The application of such models would allow greater assertiveness
   in tailoring psychotherapeutic treatments at the individual level. Despite its potential benefits, the topic remains unexplored in the fields of research, training
   and practice in clinical psychology. This work aims to introduce the topic of machine learning to psychologists. It presents the basic machine learning concepts and
   terminology in an accessible manner, providing a comprehensive overview of the field. To complement the conceptual discussion, an example is constructed using data
   from a study on the effectiveness of internet-based positive psychology interventions for depression. Finally, implications of using this type of technology in
   psychotherapy research and practice are discussed.
   \vspace{\onelineskip}
 
   \noindent 
   \textbf{Keywords}: machine learning, psychotherapy, positive psychology, depression.
 \end{otherlanguage*}
\end{resumo}


% ---
% LISTA DE ILUSTRAÇÕES
% ---
\pdfbookmark[0]{\listfigurename}{lof}
\listoffigures*
\cleardoublepage
% ---


% ---
% LISTA DE QUADROS
% ---
% \pdfbookmark[0]{\listofquadrosname}{loq}
% \listofquadros*
% \cleardoublepage
% ---


% ---
% LISTA DE TABELAS
% ---
% \pdfbookmark[0]{\listtablename}{lot}
% \listoftables*
% \cleardoublepage
% ---


% ---
% LISTA DE ABREVIATURAS E SIGLAS
% ---
% \begin{siglas}
%   \item[ABNT] Associação Brasileira de Normas Técnicas
%   \item[abnTeX] ABsurdas Normas para TeX
% \end{siglas}
% ---


% ---
% LISTA DE SÍMBOLOS
% ---
% \begin{simbolos}
%   \item[$ \Gamma $] Letra grega Gama
%   \item[$ \Lambda $] Lambda
%   \item[$ \zeta $] Letra grega minúscula zeta
%   \item[$ \in $] Pertence
% \end{simbolos}
% ---


% ---
% Sumário
% ---
\pdfbookmark[0]{\contentsname}{toc}
\tableofcontents*
\cleardoublepage
% ---


% ----------------------------------------------------------
% ELEMENTOS TEXTUAIS
% ----------------------------------------------------------
\textual

\documentclass[
	% Opções da classe memoir
	12pt,				        % tamanho da fonte
	openany,			        % capítulos começam em pág ímpar (insere página vazia caso preciso)
	oneside,			        % para impressão em recto e verso. Oposto a oneside
	a4paper,			        % tamanho do papel. 
	% Opções da classe abntex2
	%chapter=TITLE,		        % títulos de capítulos convertidos em letras maiúsculas
	%section=TITLE,		        % títulos de seções convertidos em letras maiúsculas
	%subsection=TITLE,	        % títulos de subseções convertidos em letras maiúsculas
	%subsubsection=TITLE,       % títulos de subsubseções convertidos em letras maiúsculas
	% Opções do pacote babel
	english,			        % idioma adicional para hifenização
	brazil				        % o último idioma é o principal do documento
	]{abntex2}

% ---
% Pacotes básicos 
% ---
\usepackage{lmodern}			% Usa a fonte Latin Modern			
\usepackage[T1]{fontenc}		% Selecao de codigos de fonte.
\usepackage[utf8]{inputenc}		% Codificacao do documento (conversão automática dos acentos)
\usepackage{indentfirst}		% Indenta o primeiro parágrafo de cada seção.
\usepackage{color}				% Controle das cores
\usepackage{graphicx}			% Inclusão de gráficos
\usepackage{microtype} 			% para melhorias de justificação
% ---
		

% ---
% Pacotes de citações
% ---
% \usepackage[alf]{abntex2cite}						% Citações padrão ABNT
\usepackage[
  alf,
  abnt-nbr10520=1988,                     % Citações em sentence case
  abnt-emphasize=bf
  ]{abntex2cite}


% ---
% Pacotes adicionais
% ---
\usepackage{pdfpages}	% Inclusão de páginas em PDF


% ---
% Informações para CAPA e FOLHA DE ROSTO
% ---
\titulo{APRENDIZAGEM DE MÁQUINA PARA PREDIÇÃO DE DESFECHOS EM SAÚDE MENTAL: \textmd{UMA INTRODUÇÃO PARA PSICÓLOGOS}}
\autor{GUILHERME AUGUSTO DIAS}
\local{Porto Alegre}
\data{2024}
\orientador{Prof. Dr. Wagner de Lara Machado}
% \coorientador{}
% \instituicao{
%   Pontifícia Universidade Católica do Rio Grande do Sul
%   \par
%   Escola de Ciências da Saúde e da Vida
%   \par
%   Curso de Psicologia}
\tipotrabalho{TCC}
\preambulo{Trabalho de Conclusão de Curso apresentado como requisito para a obtenção do grau de Bacharel em Psicologia pela Escola de Ciências da Saúde e da Vida da Pontifícia Universidade Católica do Rio Grande do Sul.}
% ---

% ---
% Configurações de aparência do PDF final
% ---

% Alterando o aspecto da cor azul
\definecolor{blue}{RGB}{41,5,195}

% Informações do PDF
\makeatletter
\hypersetup{
		pdftitle={Aprendizagem de máquina para predição de desfechos em saúde mental: uma introdução para psicólogos}, 
		pdfauthor={Guilherme Augusto Dias},
    	pdfsubject={\imprimirpreambulo},
	    pdfcreator={LaTeX with abnTeX2},
		pdfkeywords={aprendizagem de máquina}{psicoterapia}{psicologia positiva}{depressão}, 
		colorlinks=true,		% true: colored links; false: boxed links
    	linkcolor=black,		% color of internal links
    	citecolor=black,		% color of links to bibliography
    	filecolor=black,		% color of file links
		urlcolor=black,
		bookmarksdepth=4
}
\makeatother
% --- 


% ---
% Posiciona figuras e tabelas no topo da página quando adicionadas sozinhas
% em um página em branco. Ver https://github.com/abntex/abntex2/issues/170
\makeatletter
\setlength{\@fptop}{5pt} % Set distance from top of page to first float
\makeatother
% ---

% ---
% Possibilita criação de Quadros e Lista de quadros.
% Ver https://github.com/abntex/abntex2/issues/176
%
\newcommand{\quadroname}{Quadro}
\newcommand{\listofquadrosname}{Lista de quadros}

\newfloat[chapter]{quadro}{loq}{\quadroname}
\newlistof{listofquadros}{loq}{\listofquadrosname}
\newlistentry{quadro}{loq}{0}

% Configurações para atender às regras da ABNT
\setfloatadjustment{quadro}{\centering}
\counterwithout{quadro}{chapter}
\renewcommand{\cftquadroname}{\quadroname\space} 
\renewcommand*{\cftquadroaftersnum}{\hfill--\hfill}

\setfloatlocations{quadro}{hbtp} % Ver https://github.com/abntex/abntex2/issues/176
% ---


% --- 
% Espaçamentos entre linhas e parágrafos 
% --- 

% O tamanho do parágrafo é dado por:
\setlength{\parindent}{1.3cm}

% Controle do espaçamento entre um parágrafo e outro:
\setlength{\parskip}{0.2cm}  % tente também \onelineskip

% ---
% Compila o índice
% ---
\makeindex
% ---


% ----
% Início do documento
% ----
\begin{document}

% Seleciona o idioma do documento (conforme pacotes do babel)
\selectlanguage{brazil}

% Retira espaço extra obsoleto entre as frases.
\frenchspacing 

% ----------------------------------------------------------
% ELEMENTOS PRÉ-TEXTUAIS
% ----------------------------------------------------------
\pretextual

% ---
% CAPA
% ---
\includepdf{./05-extras/capa.pdf}
% ---


% ---
% FOLHA DE ROSTO
% (o * indica que haverá a ficha bibliográfica)
% ---
\imprimirfolhaderosto
% ---


% ---
% Ficha bibliografica
% ---

% Isto é um exemplo de Ficha Catalográfica, ou ``Dados internacionais de
% catalogação-na-publicação''. Você pode utilizar este modelo como referência. 
% Porém, provavelmente a biblioteca da sua universidade lhe fornecerá um PDF
% com a ficha catalográfica definitiva após a defesa do trabalho. Quando estiver
% com o documento, salve-o como PDF no diretório do seu projeto e substitua todo
% o conteúdo de implementação deste arquivo pelo comando abaixo:
%
% \usepackage{pdfpages}		% necessário para comando \includepdf
% \begin{fichacatalografica}
%     \includepdf{fig_ficha_catalografica.pdf}
% \end{fichacatalografica}

% \begin{fichacatalografica}
% 	\sffamily
% 	\vspace*{\fill}                             % Posição vertical
% 	\begin{center}                              % Minipage Centralizado
% 	\fbox{\begin{minipage}[c][8cm]{13.5cm}		% Largura
% 	\small
% 	\imprimirautor
% 	%Sobrenome, Nome do autor
	
% 	\hspace{0.5cm} \imprimirtitulo  / \imprimirautor. --
% 	\imprimirlocal, \imprimirdata-
	
% 	\hspace{0.5cm} \thelastpage p. : il. (algumas color.) ; 30 cm.\\
	
% 	\hspace{0.5cm} \imprimirorientadorRotulo~\imprimirorientador\\
	
% 	\hspace{0.5cm}
% 	\parbox[t]{\textwidth}{\imprimirtipotrabalho~--~\imprimirinstituicao,
% 	\imprimirdata.}\\
	
% 	\hspace{0.5cm}
% 		1. Palavra-chave1.
% 		2. Palavra-chave2.
% 		2. Palavra-chave3.
% 		I. Orientador.
% 		II. Universidade xxx.
% 		III. Faculdade de xxx.
% 		IV. Título 			
% 	\end{minipage}}
% 	\end{center}
% \end{fichacatalografica}
% ---


% ---
% FOLHA DE APROVAÇÃO
% ---

% Isto é um exemplo de Folha de aprovação, elemento obrigatório da NBR
% 14724/2011 (seção 4.2.1.3). Você pode utilizar este modelo até a aprovação
% do trabalho. Após isso, substitua todo o conteúdo deste arquivo por uma
% imagem da página assinada pela banca com o comando abaixo:
%
% \begin{folhadeaprovacao}
% \includepdf{folhadeaprovacao_final.pdf}
% \end{folhadeaprovacao}
%
\setlength{\ABNTEXsignskip}{2cm}
\setlength{\ABNTEXsignwidth}{10cm}

\begin{folhadeaprovacao}
  \begin{center}
    {\ABNTEXchapterfont\large\imprimirautor}

    \vspace*{1cm}

    \begin{center}
      \ABNTEXchapterfont\bfseries\Large\imprimirtitulo
    \end{center}

    \vspace*{1cm}
    
    \hspace{.45\textwidth}
    \begin{minipage}{.5\textwidth}
        \imprimirpreambulo
    \end{minipage}%
   \end{center}

   \vspace*{1cm}
        
   \begin{center}
   Aprovado em 17 de julho de 2024.

   \vspace*{1cm}

   Banca examinadora:
   \end{center}

   \vspace*{-1cm}

   \assinatura{Prof. Ma. Caroline Santa Maria Rodrigues -- PUCRS}
   \assinatura{Prof. Dra. Clarissa Pinto Pizarro de Freitas -- PUCRS}
   \assinatura{\imprimirorientador -- PUCRS}
\end{folhadeaprovacao}
% ---


% ---
% Agradecimentos
% ---
\begin{agradecimentos}
Agradeço aos meus pais, Ademilson José Dias e Irinês Porfírio Dias, pelo apoio incondicional que me deram ao longo da vida,
incentivando mesmo as ideias questionáveis como aquela de fazer uma segunda graduação.

A minha irmã, Maria Eduarda Dias, cuja trajetória acadêmica é fonte de inspiração.

A minha namorada, Nathalia Fallavena Ceratti, que acompanhou a esta jornada de perto, oferecendo apoio e compreensão que
fizeram toda a diferença.

Aos meus amigos, Guilherme da Silva Petró, Norton Neves e Uilinton Luciano Marques da Silva, que foram compreensíveis nos
finais de semana em que não pude jogar RPG.

Ao professor Wagner de Lara Machado pela orientação e apoio no desenvolvimento deste trabalho.
\end{agradecimentos}
% ---


% ---
% RESUMOS
% ---
% Resumo em português
\setlength{\absparsep}{18pt} % ajusta o espaçamento dos parágrafos do resumo
\begin{resumo}
Desenvolvimentos na área de aprendizagem de máquina representam uma série de oportunidades para diversas áreas de pesquisa e de atuação profissional.
O uso de modelos de aprendizagem de máquina na predição de desfecho para tratamentos em saúde mental tem sido investigado e apresenta resultados promissores.
A aplicação desse tipo de modelo possibilitaria, por exemplo, maior assertividade na personalização de tratamentos psicoterápicos a nível individual.
Apesar dos benefícios potenciais, o tema ainda é pouco discutido nos espaços de pesquisa, formação e atuação em psicologia clínica. Este trabalho tem
por objetivo introduzir o tema de aprendizagem de máquina a psicólogos. Busca-se apresentar, de maneira acessível, os conceitos e terminologia básicos
em aprendizagem de máquina, proporcionando um panorama geral da área. Para complementar a exposição conceitual, é construído um exemplo a partir de dados
de um estudo sobre eficácia de intervenções em psicologia positiva para depressão via internet. Por fim, discutem-se algumas implicações do uso desse tipo
de tecnologia para a pesquisa e para a prática em psicoterapia.

 \textbf{Palavras-chave}: aprendizagem de máquina, psicoterapia, psicologia positiva, depressão.
\end{resumo}

% Resumo em inglês
\begin{resumo}[Abstract]
 \begin{otherlanguage*}{english}
   Advances in machine learning have opened up a multitude of opportunities for various research and professional fields. The use of machine learning models to
   predict treatment outcomes in mental health has been investigated showing promising results. The application of such models would allow greater assertiveness
   in tailoring psychotherapeutic treatments at the individual level. Despite its potential benefits, the topic remains unexplored in the fields of research, training
   and practice in clinical psychology. This work aims to introduce the topic of machine learning to psychologists. It presents the basic machine learning concepts and
   terminology in an accessible manner, providing a comprehensive overview of the field. To complement the conceptual discussion, an example is constructed using data
   from a study on the effectiveness of internet-based positive psychology interventions for depression. Finally, implications of using this type of technology in
   psychotherapy research and practice are discussed.
   \vspace{\onelineskip}
 
   \noindent 
   \textbf{Keywords}: machine learning, psychotherapy, positive psychology, depression.
 \end{otherlanguage*}
\end{resumo}


% ---
% LISTA DE ILUSTRAÇÕES
% ---
\pdfbookmark[0]{\listfigurename}{lof}
\listoffigures*
\cleardoublepage
% ---


% ---
% LISTA DE QUADROS
% ---
% \pdfbookmark[0]{\listofquadrosname}{loq}
% \listofquadros*
% \cleardoublepage
% ---


% ---
% LISTA DE TABELAS
% ---
% \pdfbookmark[0]{\listtablename}{lot}
% \listoftables*
% \cleardoublepage
% ---


% ---
% LISTA DE ABREVIATURAS E SIGLAS
% ---
% \begin{siglas}
%   \item[ABNT] Associação Brasileira de Normas Técnicas
%   \item[abnTeX] ABsurdas Normas para TeX
% \end{siglas}
% ---


% ---
% LISTA DE SÍMBOLOS
% ---
% \begin{simbolos}
%   \item[$ \Gamma $] Letra grega Gama
%   \item[$ \Lambda $] Lambda
%   \item[$ \zeta $] Letra grega minúscula zeta
%   \item[$ \in $] Pertence
% \end{simbolos}
% ---


% ---
% Sumário
% ---
\pdfbookmark[0]{\contentsname}{toc}
\tableofcontents*
\cleardoublepage
% ---


% ----------------------------------------------------------
% ELEMENTOS TEXTUAIS
% ----------------------------------------------------------
\textual

\include{./01-introducao/main.tex}

\include{./02-machine-learning/main.tex}

\include{./03-exemplo/main.tex}

\include{./04-conclusao/main.tex}

% ----------------------------------------------------------
% ELEMENTOS PÓS-TEXTUAIS
% ----------------------------------------------------------
\postextual
% ----------------------------------------------------------


% ----------------------------------------------------------
% REFERÊNCIAS
% ----------------------------------------------------------
\bibliography{references.bib}


% ----------------------------------------------------------
% GLOSSÁRIO
% ----------------------------------------------------------
%
% Consulte o manual da classe abntex2 para orientações sobre o glossário.
%
%\glossary

% ----------------------------------------------------------
% APÊNDICES
% ----------------------------------------------------------

% ---
% Inicia os apêndices
% ---
% \begin{apendicesenv}

% % Imprime uma página indicando o início dos apêndices
% \partapendices

% % ----------------------------------------------------------
% \chapter{Quisque libero justo}
% % ----------------------------------------------------------

% % ----------------------------------------------------------
% \chapter{Nullam elementum urna}
% % ----------------------------------------------------------

% \end{apendicesenv}
% ---


% ----------------------------------------------------------
% ANEXOS
% ----------------------------------------------------------

% ---
% Inicia os anexos
% ---
% \begin{anexosenv}

% % Imprime uma página indicando o início dos anexos
% \partanexos

% % ---
% \chapter{Morbi ultrices rutrum lorem}
% % ---

% % ---
% \chapter{Cras non urna sed feugiat cum}
% % ---


% % ---
% \chapter{Fusce facilisis lacinia dui}
% % ---

% \end{anexosenv}

%---------------------------------------------------------------------
% INDICE REMISSIVO
%---------------------------------------------------------------------
% \phantompart
% \printindex
%---------------------------------------------------------------------
\includepdf{./05-extras/contracapa.pdf}
\end{document}


\documentclass[
	% Opções da classe memoir
	12pt,				        % tamanho da fonte
	openany,			        % capítulos começam em pág ímpar (insere página vazia caso preciso)
	oneside,			        % para impressão em recto e verso. Oposto a oneside
	a4paper,			        % tamanho do papel. 
	% Opções da classe abntex2
	%chapter=TITLE,		        % títulos de capítulos convertidos em letras maiúsculas
	%section=TITLE,		        % títulos de seções convertidos em letras maiúsculas
	%subsection=TITLE,	        % títulos de subseções convertidos em letras maiúsculas
	%subsubsection=TITLE,       % títulos de subsubseções convertidos em letras maiúsculas
	% Opções do pacote babel
	english,			        % idioma adicional para hifenização
	brazil				        % o último idioma é o principal do documento
	]{abntex2}

% ---
% Pacotes básicos 
% ---
\usepackage{lmodern}			% Usa a fonte Latin Modern			
\usepackage[T1]{fontenc}		% Selecao de codigos de fonte.
\usepackage[utf8]{inputenc}		% Codificacao do documento (conversão automática dos acentos)
\usepackage{indentfirst}		% Indenta o primeiro parágrafo de cada seção.
\usepackage{color}				% Controle das cores
\usepackage{graphicx}			% Inclusão de gráficos
\usepackage{microtype} 			% para melhorias de justificação
% ---
		

% ---
% Pacotes de citações
% ---
% \usepackage[alf]{abntex2cite}						% Citações padrão ABNT
\usepackage[
  alf,
  abnt-nbr10520=1988,                     % Citações em sentence case
  abnt-emphasize=bf
  ]{abntex2cite}


% ---
% Pacotes adicionais
% ---
\usepackage{pdfpages}	% Inclusão de páginas em PDF


% ---
% Informações para CAPA e FOLHA DE ROSTO
% ---
\titulo{APRENDIZAGEM DE MÁQUINA PARA PREDIÇÃO DE DESFECHOS EM SAÚDE MENTAL: \textmd{UMA INTRODUÇÃO PARA PSICÓLOGOS}}
\autor{GUILHERME AUGUSTO DIAS}
\local{Porto Alegre}
\data{2024}
\orientador{Prof. Dr. Wagner de Lara Machado}
% \coorientador{}
% \instituicao{
%   Pontifícia Universidade Católica do Rio Grande do Sul
%   \par
%   Escola de Ciências da Saúde e da Vida
%   \par
%   Curso de Psicologia}
\tipotrabalho{TCC}
\preambulo{Trabalho de Conclusão de Curso apresentado como requisito para a obtenção do grau de Bacharel em Psicologia pela Escola de Ciências da Saúde e da Vida da Pontifícia Universidade Católica do Rio Grande do Sul.}
% ---

% ---
% Configurações de aparência do PDF final
% ---

% Alterando o aspecto da cor azul
\definecolor{blue}{RGB}{41,5,195}

% Informações do PDF
\makeatletter
\hypersetup{
		pdftitle={Aprendizagem de máquina para predição de desfechos em saúde mental: uma introdução para psicólogos}, 
		pdfauthor={Guilherme Augusto Dias},
    	pdfsubject={\imprimirpreambulo},
	    pdfcreator={LaTeX with abnTeX2},
		pdfkeywords={aprendizagem de máquina}{psicoterapia}{psicologia positiva}{depressão}, 
		colorlinks=true,		% true: colored links; false: boxed links
    	linkcolor=black,		% color of internal links
    	citecolor=black,		% color of links to bibliography
    	filecolor=black,		% color of file links
		urlcolor=black,
		bookmarksdepth=4
}
\makeatother
% --- 


% ---
% Posiciona figuras e tabelas no topo da página quando adicionadas sozinhas
% em um página em branco. Ver https://github.com/abntex/abntex2/issues/170
\makeatletter
\setlength{\@fptop}{5pt} % Set distance from top of page to first float
\makeatother
% ---

% ---
% Possibilita criação de Quadros e Lista de quadros.
% Ver https://github.com/abntex/abntex2/issues/176
%
\newcommand{\quadroname}{Quadro}
\newcommand{\listofquadrosname}{Lista de quadros}

\newfloat[chapter]{quadro}{loq}{\quadroname}
\newlistof{listofquadros}{loq}{\listofquadrosname}
\newlistentry{quadro}{loq}{0}

% Configurações para atender às regras da ABNT
\setfloatadjustment{quadro}{\centering}
\counterwithout{quadro}{chapter}
\renewcommand{\cftquadroname}{\quadroname\space} 
\renewcommand*{\cftquadroaftersnum}{\hfill--\hfill}

\setfloatlocations{quadro}{hbtp} % Ver https://github.com/abntex/abntex2/issues/176
% ---


% --- 
% Espaçamentos entre linhas e parágrafos 
% --- 

% O tamanho do parágrafo é dado por:
\setlength{\parindent}{1.3cm}

% Controle do espaçamento entre um parágrafo e outro:
\setlength{\parskip}{0.2cm}  % tente também \onelineskip

% ---
% Compila o índice
% ---
\makeindex
% ---


% ----
% Início do documento
% ----
\begin{document}

% Seleciona o idioma do documento (conforme pacotes do babel)
\selectlanguage{brazil}

% Retira espaço extra obsoleto entre as frases.
\frenchspacing 

% ----------------------------------------------------------
% ELEMENTOS PRÉ-TEXTUAIS
% ----------------------------------------------------------
\pretextual

% ---
% CAPA
% ---
\includepdf{./05-extras/capa.pdf}
% ---


% ---
% FOLHA DE ROSTO
% (o * indica que haverá a ficha bibliográfica)
% ---
\imprimirfolhaderosto
% ---


% ---
% Ficha bibliografica
% ---

% Isto é um exemplo de Ficha Catalográfica, ou ``Dados internacionais de
% catalogação-na-publicação''. Você pode utilizar este modelo como referência. 
% Porém, provavelmente a biblioteca da sua universidade lhe fornecerá um PDF
% com a ficha catalográfica definitiva após a defesa do trabalho. Quando estiver
% com o documento, salve-o como PDF no diretório do seu projeto e substitua todo
% o conteúdo de implementação deste arquivo pelo comando abaixo:
%
% \usepackage{pdfpages}		% necessário para comando \includepdf
% \begin{fichacatalografica}
%     \includepdf{fig_ficha_catalografica.pdf}
% \end{fichacatalografica}

% \begin{fichacatalografica}
% 	\sffamily
% 	\vspace*{\fill}                             % Posição vertical
% 	\begin{center}                              % Minipage Centralizado
% 	\fbox{\begin{minipage}[c][8cm]{13.5cm}		% Largura
% 	\small
% 	\imprimirautor
% 	%Sobrenome, Nome do autor
	
% 	\hspace{0.5cm} \imprimirtitulo  / \imprimirautor. --
% 	\imprimirlocal, \imprimirdata-
	
% 	\hspace{0.5cm} \thelastpage p. : il. (algumas color.) ; 30 cm.\\
	
% 	\hspace{0.5cm} \imprimirorientadorRotulo~\imprimirorientador\\
	
% 	\hspace{0.5cm}
% 	\parbox[t]{\textwidth}{\imprimirtipotrabalho~--~\imprimirinstituicao,
% 	\imprimirdata.}\\
	
% 	\hspace{0.5cm}
% 		1. Palavra-chave1.
% 		2. Palavra-chave2.
% 		2. Palavra-chave3.
% 		I. Orientador.
% 		II. Universidade xxx.
% 		III. Faculdade de xxx.
% 		IV. Título 			
% 	\end{minipage}}
% 	\end{center}
% \end{fichacatalografica}
% ---


% ---
% FOLHA DE APROVAÇÃO
% ---

% Isto é um exemplo de Folha de aprovação, elemento obrigatório da NBR
% 14724/2011 (seção 4.2.1.3). Você pode utilizar este modelo até a aprovação
% do trabalho. Após isso, substitua todo o conteúdo deste arquivo por uma
% imagem da página assinada pela banca com o comando abaixo:
%
% \begin{folhadeaprovacao}
% \includepdf{folhadeaprovacao_final.pdf}
% \end{folhadeaprovacao}
%
\setlength{\ABNTEXsignskip}{2cm}
\setlength{\ABNTEXsignwidth}{10cm}

\begin{folhadeaprovacao}
  \begin{center}
    {\ABNTEXchapterfont\large\imprimirautor}

    \vspace*{1cm}

    \begin{center}
      \ABNTEXchapterfont\bfseries\Large\imprimirtitulo
    \end{center}

    \vspace*{1cm}
    
    \hspace{.45\textwidth}
    \begin{minipage}{.5\textwidth}
        \imprimirpreambulo
    \end{minipage}%
   \end{center}

   \vspace*{1cm}
        
   \begin{center}
   Aprovado em 17 de julho de 2024.

   \vspace*{1cm}

   Banca examinadora:
   \end{center}

   \vspace*{-1cm}

   \assinatura{Prof. Ma. Caroline Santa Maria Rodrigues -- PUCRS}
   \assinatura{Prof. Dra. Clarissa Pinto Pizarro de Freitas -- PUCRS}
   \assinatura{\imprimirorientador -- PUCRS}
\end{folhadeaprovacao}
% ---


% ---
% Agradecimentos
% ---
\begin{agradecimentos}
Agradeço aos meus pais, Ademilson José Dias e Irinês Porfírio Dias, pelo apoio incondicional que me deram ao longo da vida,
incentivando mesmo as ideias questionáveis como aquela de fazer uma segunda graduação.

A minha irmã, Maria Eduarda Dias, cuja trajetória acadêmica é fonte de inspiração.

A minha namorada, Nathalia Fallavena Ceratti, que acompanhou a esta jornada de perto, oferecendo apoio e compreensão que
fizeram toda a diferença.

Aos meus amigos, Guilherme da Silva Petró, Norton Neves e Uilinton Luciano Marques da Silva, que foram compreensíveis nos
finais de semana em que não pude jogar RPG.

Ao professor Wagner de Lara Machado pela orientação e apoio no desenvolvimento deste trabalho.
\end{agradecimentos}
% ---


% ---
% RESUMOS
% ---
% Resumo em português
\setlength{\absparsep}{18pt} % ajusta o espaçamento dos parágrafos do resumo
\begin{resumo}
Desenvolvimentos na área de aprendizagem de máquina representam uma série de oportunidades para diversas áreas de pesquisa e de atuação profissional.
O uso de modelos de aprendizagem de máquina na predição de desfecho para tratamentos em saúde mental tem sido investigado e apresenta resultados promissores.
A aplicação desse tipo de modelo possibilitaria, por exemplo, maior assertividade na personalização de tratamentos psicoterápicos a nível individual.
Apesar dos benefícios potenciais, o tema ainda é pouco discutido nos espaços de pesquisa, formação e atuação em psicologia clínica. Este trabalho tem
por objetivo introduzir o tema de aprendizagem de máquina a psicólogos. Busca-se apresentar, de maneira acessível, os conceitos e terminologia básicos
em aprendizagem de máquina, proporcionando um panorama geral da área. Para complementar a exposição conceitual, é construído um exemplo a partir de dados
de um estudo sobre eficácia de intervenções em psicologia positiva para depressão via internet. Por fim, discutem-se algumas implicações do uso desse tipo
de tecnologia para a pesquisa e para a prática em psicoterapia.

 \textbf{Palavras-chave}: aprendizagem de máquina, psicoterapia, psicologia positiva, depressão.
\end{resumo}

% Resumo em inglês
\begin{resumo}[Abstract]
 \begin{otherlanguage*}{english}
   Advances in machine learning have opened up a multitude of opportunities for various research and professional fields. The use of machine learning models to
   predict treatment outcomes in mental health has been investigated showing promising results. The application of such models would allow greater assertiveness
   in tailoring psychotherapeutic treatments at the individual level. Despite its potential benefits, the topic remains unexplored in the fields of research, training
   and practice in clinical psychology. This work aims to introduce the topic of machine learning to psychologists. It presents the basic machine learning concepts and
   terminology in an accessible manner, providing a comprehensive overview of the field. To complement the conceptual discussion, an example is constructed using data
   from a study on the effectiveness of internet-based positive psychology interventions for depression. Finally, implications of using this type of technology in
   psychotherapy research and practice are discussed.
   \vspace{\onelineskip}
 
   \noindent 
   \textbf{Keywords}: machine learning, psychotherapy, positive psychology, depression.
 \end{otherlanguage*}
\end{resumo}


% ---
% LISTA DE ILUSTRAÇÕES
% ---
\pdfbookmark[0]{\listfigurename}{lof}
\listoffigures*
\cleardoublepage
% ---


% ---
% LISTA DE QUADROS
% ---
% \pdfbookmark[0]{\listofquadrosname}{loq}
% \listofquadros*
% \cleardoublepage
% ---


% ---
% LISTA DE TABELAS
% ---
% \pdfbookmark[0]{\listtablename}{lot}
% \listoftables*
% \cleardoublepage
% ---


% ---
% LISTA DE ABREVIATURAS E SIGLAS
% ---
% \begin{siglas}
%   \item[ABNT] Associação Brasileira de Normas Técnicas
%   \item[abnTeX] ABsurdas Normas para TeX
% \end{siglas}
% ---


% ---
% LISTA DE SÍMBOLOS
% ---
% \begin{simbolos}
%   \item[$ \Gamma $] Letra grega Gama
%   \item[$ \Lambda $] Lambda
%   \item[$ \zeta $] Letra grega minúscula zeta
%   \item[$ \in $] Pertence
% \end{simbolos}
% ---


% ---
% Sumário
% ---
\pdfbookmark[0]{\contentsname}{toc}
\tableofcontents*
\cleardoublepage
% ---


% ----------------------------------------------------------
% ELEMENTOS TEXTUAIS
% ----------------------------------------------------------
\textual

\include{./01-introducao/main.tex}

\include{./02-machine-learning/main.tex}

\include{./03-exemplo/main.tex}

\include{./04-conclusao/main.tex}

% ----------------------------------------------------------
% ELEMENTOS PÓS-TEXTUAIS
% ----------------------------------------------------------
\postextual
% ----------------------------------------------------------


% ----------------------------------------------------------
% REFERÊNCIAS
% ----------------------------------------------------------
\bibliography{references.bib}


% ----------------------------------------------------------
% GLOSSÁRIO
% ----------------------------------------------------------
%
% Consulte o manual da classe abntex2 para orientações sobre o glossário.
%
%\glossary

% ----------------------------------------------------------
% APÊNDICES
% ----------------------------------------------------------

% ---
% Inicia os apêndices
% ---
% \begin{apendicesenv}

% % Imprime uma página indicando o início dos apêndices
% \partapendices

% % ----------------------------------------------------------
% \chapter{Quisque libero justo}
% % ----------------------------------------------------------

% % ----------------------------------------------------------
% \chapter{Nullam elementum urna}
% % ----------------------------------------------------------

% \end{apendicesenv}
% ---


% ----------------------------------------------------------
% ANEXOS
% ----------------------------------------------------------

% ---
% Inicia os anexos
% ---
% \begin{anexosenv}

% % Imprime uma página indicando o início dos anexos
% \partanexos

% % ---
% \chapter{Morbi ultrices rutrum lorem}
% % ---

% % ---
% \chapter{Cras non urna sed feugiat cum}
% % ---


% % ---
% \chapter{Fusce facilisis lacinia dui}
% % ---

% \end{anexosenv}

%---------------------------------------------------------------------
% INDICE REMISSIVO
%---------------------------------------------------------------------
% \phantompart
% \printindex
%---------------------------------------------------------------------
\includepdf{./05-extras/contracapa.pdf}
\end{document}


\documentclass[
	% Opções da classe memoir
	12pt,				        % tamanho da fonte
	openany,			        % capítulos começam em pág ímpar (insere página vazia caso preciso)
	oneside,			        % para impressão em recto e verso. Oposto a oneside
	a4paper,			        % tamanho do papel. 
	% Opções da classe abntex2
	%chapter=TITLE,		        % títulos de capítulos convertidos em letras maiúsculas
	%section=TITLE,		        % títulos de seções convertidos em letras maiúsculas
	%subsection=TITLE,	        % títulos de subseções convertidos em letras maiúsculas
	%subsubsection=TITLE,       % títulos de subsubseções convertidos em letras maiúsculas
	% Opções do pacote babel
	english,			        % idioma adicional para hifenização
	brazil				        % o último idioma é o principal do documento
	]{abntex2}

% ---
% Pacotes básicos 
% ---
\usepackage{lmodern}			% Usa a fonte Latin Modern			
\usepackage[T1]{fontenc}		% Selecao de codigos de fonte.
\usepackage[utf8]{inputenc}		% Codificacao do documento (conversão automática dos acentos)
\usepackage{indentfirst}		% Indenta o primeiro parágrafo de cada seção.
\usepackage{color}				% Controle das cores
\usepackage{graphicx}			% Inclusão de gráficos
\usepackage{microtype} 			% para melhorias de justificação
% ---
		

% ---
% Pacotes de citações
% ---
% \usepackage[alf]{abntex2cite}						% Citações padrão ABNT
\usepackage[
  alf,
  abnt-nbr10520=1988,                     % Citações em sentence case
  abnt-emphasize=bf
  ]{abntex2cite}


% ---
% Pacotes adicionais
% ---
\usepackage{pdfpages}	% Inclusão de páginas em PDF


% ---
% Informações para CAPA e FOLHA DE ROSTO
% ---
\titulo{APRENDIZAGEM DE MÁQUINA PARA PREDIÇÃO DE DESFECHOS EM SAÚDE MENTAL: \textmd{UMA INTRODUÇÃO PARA PSICÓLOGOS}}
\autor{GUILHERME AUGUSTO DIAS}
\local{Porto Alegre}
\data{2024}
\orientador{Prof. Dr. Wagner de Lara Machado}
% \coorientador{}
% \instituicao{
%   Pontifícia Universidade Católica do Rio Grande do Sul
%   \par
%   Escola de Ciências da Saúde e da Vida
%   \par
%   Curso de Psicologia}
\tipotrabalho{TCC}
\preambulo{Trabalho de Conclusão de Curso apresentado como requisito para a obtenção do grau de Bacharel em Psicologia pela Escola de Ciências da Saúde e da Vida da Pontifícia Universidade Católica do Rio Grande do Sul.}
% ---

% ---
% Configurações de aparência do PDF final
% ---

% Alterando o aspecto da cor azul
\definecolor{blue}{RGB}{41,5,195}

% Informações do PDF
\makeatletter
\hypersetup{
		pdftitle={Aprendizagem de máquina para predição de desfechos em saúde mental: uma introdução para psicólogos}, 
		pdfauthor={Guilherme Augusto Dias},
    	pdfsubject={\imprimirpreambulo},
	    pdfcreator={LaTeX with abnTeX2},
		pdfkeywords={aprendizagem de máquina}{psicoterapia}{psicologia positiva}{depressão}, 
		colorlinks=true,		% true: colored links; false: boxed links
    	linkcolor=black,		% color of internal links
    	citecolor=black,		% color of links to bibliography
    	filecolor=black,		% color of file links
		urlcolor=black,
		bookmarksdepth=4
}
\makeatother
% --- 


% ---
% Posiciona figuras e tabelas no topo da página quando adicionadas sozinhas
% em um página em branco. Ver https://github.com/abntex/abntex2/issues/170
\makeatletter
\setlength{\@fptop}{5pt} % Set distance from top of page to first float
\makeatother
% ---

% ---
% Possibilita criação de Quadros e Lista de quadros.
% Ver https://github.com/abntex/abntex2/issues/176
%
\newcommand{\quadroname}{Quadro}
\newcommand{\listofquadrosname}{Lista de quadros}

\newfloat[chapter]{quadro}{loq}{\quadroname}
\newlistof{listofquadros}{loq}{\listofquadrosname}
\newlistentry{quadro}{loq}{0}

% Configurações para atender às regras da ABNT
\setfloatadjustment{quadro}{\centering}
\counterwithout{quadro}{chapter}
\renewcommand{\cftquadroname}{\quadroname\space} 
\renewcommand*{\cftquadroaftersnum}{\hfill--\hfill}

\setfloatlocations{quadro}{hbtp} % Ver https://github.com/abntex/abntex2/issues/176
% ---


% --- 
% Espaçamentos entre linhas e parágrafos 
% --- 

% O tamanho do parágrafo é dado por:
\setlength{\parindent}{1.3cm}

% Controle do espaçamento entre um parágrafo e outro:
\setlength{\parskip}{0.2cm}  % tente também \onelineskip

% ---
% Compila o índice
% ---
\makeindex
% ---


% ----
% Início do documento
% ----
\begin{document}

% Seleciona o idioma do documento (conforme pacotes do babel)
\selectlanguage{brazil}

% Retira espaço extra obsoleto entre as frases.
\frenchspacing 

% ----------------------------------------------------------
% ELEMENTOS PRÉ-TEXTUAIS
% ----------------------------------------------------------
\pretextual

% ---
% CAPA
% ---
\includepdf{./05-extras/capa.pdf}
% ---


% ---
% FOLHA DE ROSTO
% (o * indica que haverá a ficha bibliográfica)
% ---
\imprimirfolhaderosto
% ---


% ---
% Ficha bibliografica
% ---

% Isto é um exemplo de Ficha Catalográfica, ou ``Dados internacionais de
% catalogação-na-publicação''. Você pode utilizar este modelo como referência. 
% Porém, provavelmente a biblioteca da sua universidade lhe fornecerá um PDF
% com a ficha catalográfica definitiva após a defesa do trabalho. Quando estiver
% com o documento, salve-o como PDF no diretório do seu projeto e substitua todo
% o conteúdo de implementação deste arquivo pelo comando abaixo:
%
% \usepackage{pdfpages}		% necessário para comando \includepdf
% \begin{fichacatalografica}
%     \includepdf{fig_ficha_catalografica.pdf}
% \end{fichacatalografica}

% \begin{fichacatalografica}
% 	\sffamily
% 	\vspace*{\fill}                             % Posição vertical
% 	\begin{center}                              % Minipage Centralizado
% 	\fbox{\begin{minipage}[c][8cm]{13.5cm}		% Largura
% 	\small
% 	\imprimirautor
% 	%Sobrenome, Nome do autor
	
% 	\hspace{0.5cm} \imprimirtitulo  / \imprimirautor. --
% 	\imprimirlocal, \imprimirdata-
	
% 	\hspace{0.5cm} \thelastpage p. : il. (algumas color.) ; 30 cm.\\
	
% 	\hspace{0.5cm} \imprimirorientadorRotulo~\imprimirorientador\\
	
% 	\hspace{0.5cm}
% 	\parbox[t]{\textwidth}{\imprimirtipotrabalho~--~\imprimirinstituicao,
% 	\imprimirdata.}\\
	
% 	\hspace{0.5cm}
% 		1. Palavra-chave1.
% 		2. Palavra-chave2.
% 		2. Palavra-chave3.
% 		I. Orientador.
% 		II. Universidade xxx.
% 		III. Faculdade de xxx.
% 		IV. Título 			
% 	\end{minipage}}
% 	\end{center}
% \end{fichacatalografica}
% ---


% ---
% FOLHA DE APROVAÇÃO
% ---

% Isto é um exemplo de Folha de aprovação, elemento obrigatório da NBR
% 14724/2011 (seção 4.2.1.3). Você pode utilizar este modelo até a aprovação
% do trabalho. Após isso, substitua todo o conteúdo deste arquivo por uma
% imagem da página assinada pela banca com o comando abaixo:
%
% \begin{folhadeaprovacao}
% \includepdf{folhadeaprovacao_final.pdf}
% \end{folhadeaprovacao}
%
\setlength{\ABNTEXsignskip}{2cm}
\setlength{\ABNTEXsignwidth}{10cm}

\begin{folhadeaprovacao}
  \begin{center}
    {\ABNTEXchapterfont\large\imprimirautor}

    \vspace*{1cm}

    \begin{center}
      \ABNTEXchapterfont\bfseries\Large\imprimirtitulo
    \end{center}

    \vspace*{1cm}
    
    \hspace{.45\textwidth}
    \begin{minipage}{.5\textwidth}
        \imprimirpreambulo
    \end{minipage}%
   \end{center}

   \vspace*{1cm}
        
   \begin{center}
   Aprovado em 17 de julho de 2024.

   \vspace*{1cm}

   Banca examinadora:
   \end{center}

   \vspace*{-1cm}

   \assinatura{Prof. Ma. Caroline Santa Maria Rodrigues -- PUCRS}
   \assinatura{Prof. Dra. Clarissa Pinto Pizarro de Freitas -- PUCRS}
   \assinatura{\imprimirorientador -- PUCRS}
\end{folhadeaprovacao}
% ---


% ---
% Agradecimentos
% ---
\begin{agradecimentos}
Agradeço aos meus pais, Ademilson José Dias e Irinês Porfírio Dias, pelo apoio incondicional que me deram ao longo da vida,
incentivando mesmo as ideias questionáveis como aquela de fazer uma segunda graduação.

A minha irmã, Maria Eduarda Dias, cuja trajetória acadêmica é fonte de inspiração.

A minha namorada, Nathalia Fallavena Ceratti, que acompanhou a esta jornada de perto, oferecendo apoio e compreensão que
fizeram toda a diferença.

Aos meus amigos, Guilherme da Silva Petró, Norton Neves e Uilinton Luciano Marques da Silva, que foram compreensíveis nos
finais de semana em que não pude jogar RPG.

Ao professor Wagner de Lara Machado pela orientação e apoio no desenvolvimento deste trabalho.
\end{agradecimentos}
% ---


% ---
% RESUMOS
% ---
% Resumo em português
\setlength{\absparsep}{18pt} % ajusta o espaçamento dos parágrafos do resumo
\begin{resumo}
Desenvolvimentos na área de aprendizagem de máquina representam uma série de oportunidades para diversas áreas de pesquisa e de atuação profissional.
O uso de modelos de aprendizagem de máquina na predição de desfecho para tratamentos em saúde mental tem sido investigado e apresenta resultados promissores.
A aplicação desse tipo de modelo possibilitaria, por exemplo, maior assertividade na personalização de tratamentos psicoterápicos a nível individual.
Apesar dos benefícios potenciais, o tema ainda é pouco discutido nos espaços de pesquisa, formação e atuação em psicologia clínica. Este trabalho tem
por objetivo introduzir o tema de aprendizagem de máquina a psicólogos. Busca-se apresentar, de maneira acessível, os conceitos e terminologia básicos
em aprendizagem de máquina, proporcionando um panorama geral da área. Para complementar a exposição conceitual, é construído um exemplo a partir de dados
de um estudo sobre eficácia de intervenções em psicologia positiva para depressão via internet. Por fim, discutem-se algumas implicações do uso desse tipo
de tecnologia para a pesquisa e para a prática em psicoterapia.

 \textbf{Palavras-chave}: aprendizagem de máquina, psicoterapia, psicologia positiva, depressão.
\end{resumo}

% Resumo em inglês
\begin{resumo}[Abstract]
 \begin{otherlanguage*}{english}
   Advances in machine learning have opened up a multitude of opportunities for various research and professional fields. The use of machine learning models to
   predict treatment outcomes in mental health has been investigated showing promising results. The application of such models would allow greater assertiveness
   in tailoring psychotherapeutic treatments at the individual level. Despite its potential benefits, the topic remains unexplored in the fields of research, training
   and practice in clinical psychology. This work aims to introduce the topic of machine learning to psychologists. It presents the basic machine learning concepts and
   terminology in an accessible manner, providing a comprehensive overview of the field. To complement the conceptual discussion, an example is constructed using data
   from a study on the effectiveness of internet-based positive psychology interventions for depression. Finally, implications of using this type of technology in
   psychotherapy research and practice are discussed.
   \vspace{\onelineskip}
 
   \noindent 
   \textbf{Keywords}: machine learning, psychotherapy, positive psychology, depression.
 \end{otherlanguage*}
\end{resumo}


% ---
% LISTA DE ILUSTRAÇÕES
% ---
\pdfbookmark[0]{\listfigurename}{lof}
\listoffigures*
\cleardoublepage
% ---


% ---
% LISTA DE QUADROS
% ---
% \pdfbookmark[0]{\listofquadrosname}{loq}
% \listofquadros*
% \cleardoublepage
% ---


% ---
% LISTA DE TABELAS
% ---
% \pdfbookmark[0]{\listtablename}{lot}
% \listoftables*
% \cleardoublepage
% ---


% ---
% LISTA DE ABREVIATURAS E SIGLAS
% ---
% \begin{siglas}
%   \item[ABNT] Associação Brasileira de Normas Técnicas
%   \item[abnTeX] ABsurdas Normas para TeX
% \end{siglas}
% ---


% ---
% LISTA DE SÍMBOLOS
% ---
% \begin{simbolos}
%   \item[$ \Gamma $] Letra grega Gama
%   \item[$ \Lambda $] Lambda
%   \item[$ \zeta $] Letra grega minúscula zeta
%   \item[$ \in $] Pertence
% \end{simbolos}
% ---


% ---
% Sumário
% ---
\pdfbookmark[0]{\contentsname}{toc}
\tableofcontents*
\cleardoublepage
% ---


% ----------------------------------------------------------
% ELEMENTOS TEXTUAIS
% ----------------------------------------------------------
\textual

\include{./01-introducao/main.tex}

\include{./02-machine-learning/main.tex}

\include{./03-exemplo/main.tex}

\include{./04-conclusao/main.tex}

% ----------------------------------------------------------
% ELEMENTOS PÓS-TEXTUAIS
% ----------------------------------------------------------
\postextual
% ----------------------------------------------------------


% ----------------------------------------------------------
% REFERÊNCIAS
% ----------------------------------------------------------
\bibliography{references.bib}


% ----------------------------------------------------------
% GLOSSÁRIO
% ----------------------------------------------------------
%
% Consulte o manual da classe abntex2 para orientações sobre o glossário.
%
%\glossary

% ----------------------------------------------------------
% APÊNDICES
% ----------------------------------------------------------

% ---
% Inicia os apêndices
% ---
% \begin{apendicesenv}

% % Imprime uma página indicando o início dos apêndices
% \partapendices

% % ----------------------------------------------------------
% \chapter{Quisque libero justo}
% % ----------------------------------------------------------

% % ----------------------------------------------------------
% \chapter{Nullam elementum urna}
% % ----------------------------------------------------------

% \end{apendicesenv}
% ---


% ----------------------------------------------------------
% ANEXOS
% ----------------------------------------------------------

% ---
% Inicia os anexos
% ---
% \begin{anexosenv}

% % Imprime uma página indicando o início dos anexos
% \partanexos

% % ---
% \chapter{Morbi ultrices rutrum lorem}
% % ---

% % ---
% \chapter{Cras non urna sed feugiat cum}
% % ---


% % ---
% \chapter{Fusce facilisis lacinia dui}
% % ---

% \end{anexosenv}

%---------------------------------------------------------------------
% INDICE REMISSIVO
%---------------------------------------------------------------------
% \phantompart
% \printindex
%---------------------------------------------------------------------
\includepdf{./05-extras/contracapa.pdf}
\end{document}


\documentclass[
	% Opções da classe memoir
	12pt,				        % tamanho da fonte
	openany,			        % capítulos começam em pág ímpar (insere página vazia caso preciso)
	oneside,			        % para impressão em recto e verso. Oposto a oneside
	a4paper,			        % tamanho do papel. 
	% Opções da classe abntex2
	%chapter=TITLE,		        % títulos de capítulos convertidos em letras maiúsculas
	%section=TITLE,		        % títulos de seções convertidos em letras maiúsculas
	%subsection=TITLE,	        % títulos de subseções convertidos em letras maiúsculas
	%subsubsection=TITLE,       % títulos de subsubseções convertidos em letras maiúsculas
	% Opções do pacote babel
	english,			        % idioma adicional para hifenização
	brazil				        % o último idioma é o principal do documento
	]{abntex2}

% ---
% Pacotes básicos 
% ---
\usepackage{lmodern}			% Usa a fonte Latin Modern			
\usepackage[T1]{fontenc}		% Selecao de codigos de fonte.
\usepackage[utf8]{inputenc}		% Codificacao do documento (conversão automática dos acentos)
\usepackage{indentfirst}		% Indenta o primeiro parágrafo de cada seção.
\usepackage{color}				% Controle das cores
\usepackage{graphicx}			% Inclusão de gráficos
\usepackage{microtype} 			% para melhorias de justificação
% ---
		

% ---
% Pacotes de citações
% ---
% \usepackage[alf]{abntex2cite}						% Citações padrão ABNT
\usepackage[
  alf,
  abnt-nbr10520=1988,                     % Citações em sentence case
  abnt-emphasize=bf
  ]{abntex2cite}


% ---
% Pacotes adicionais
% ---
\usepackage{pdfpages}	% Inclusão de páginas em PDF


% ---
% Informações para CAPA e FOLHA DE ROSTO
% ---
\titulo{APRENDIZAGEM DE MÁQUINA PARA PREDIÇÃO DE DESFECHOS EM SAÚDE MENTAL: \textmd{UMA INTRODUÇÃO PARA PSICÓLOGOS}}
\autor{GUILHERME AUGUSTO DIAS}
\local{Porto Alegre}
\data{2024}
\orientador{Prof. Dr. Wagner de Lara Machado}
% \coorientador{}
% \instituicao{
%   Pontifícia Universidade Católica do Rio Grande do Sul
%   \par
%   Escola de Ciências da Saúde e da Vida
%   \par
%   Curso de Psicologia}
\tipotrabalho{TCC}
\preambulo{Trabalho de Conclusão de Curso apresentado como requisito para a obtenção do grau de Bacharel em Psicologia pela Escola de Ciências da Saúde e da Vida da Pontifícia Universidade Católica do Rio Grande do Sul.}
% ---

% ---
% Configurações de aparência do PDF final
% ---

% Alterando o aspecto da cor azul
\definecolor{blue}{RGB}{41,5,195}

% Informações do PDF
\makeatletter
\hypersetup{
		pdftitle={Aprendizagem de máquina para predição de desfechos em saúde mental: uma introdução para psicólogos}, 
		pdfauthor={Guilherme Augusto Dias},
    	pdfsubject={\imprimirpreambulo},
	    pdfcreator={LaTeX with abnTeX2},
		pdfkeywords={aprendizagem de máquina}{psicoterapia}{psicologia positiva}{depressão}, 
		colorlinks=true,		% true: colored links; false: boxed links
    	linkcolor=black,		% color of internal links
    	citecolor=black,		% color of links to bibliography
    	filecolor=black,		% color of file links
		urlcolor=black,
		bookmarksdepth=4
}
\makeatother
% --- 


% ---
% Posiciona figuras e tabelas no topo da página quando adicionadas sozinhas
% em um página em branco. Ver https://github.com/abntex/abntex2/issues/170
\makeatletter
\setlength{\@fptop}{5pt} % Set distance from top of page to first float
\makeatother
% ---

% ---
% Possibilita criação de Quadros e Lista de quadros.
% Ver https://github.com/abntex/abntex2/issues/176
%
\newcommand{\quadroname}{Quadro}
\newcommand{\listofquadrosname}{Lista de quadros}

\newfloat[chapter]{quadro}{loq}{\quadroname}
\newlistof{listofquadros}{loq}{\listofquadrosname}
\newlistentry{quadro}{loq}{0}

% Configurações para atender às regras da ABNT
\setfloatadjustment{quadro}{\centering}
\counterwithout{quadro}{chapter}
\renewcommand{\cftquadroname}{\quadroname\space} 
\renewcommand*{\cftquadroaftersnum}{\hfill--\hfill}

\setfloatlocations{quadro}{hbtp} % Ver https://github.com/abntex/abntex2/issues/176
% ---


% --- 
% Espaçamentos entre linhas e parágrafos 
% --- 

% O tamanho do parágrafo é dado por:
\setlength{\parindent}{1.3cm}

% Controle do espaçamento entre um parágrafo e outro:
\setlength{\parskip}{0.2cm}  % tente também \onelineskip

% ---
% Compila o índice
% ---
\makeindex
% ---


% ----
% Início do documento
% ----
\begin{document}

% Seleciona o idioma do documento (conforme pacotes do babel)
\selectlanguage{brazil}

% Retira espaço extra obsoleto entre as frases.
\frenchspacing 

% ----------------------------------------------------------
% ELEMENTOS PRÉ-TEXTUAIS
% ----------------------------------------------------------
\pretextual

% ---
% CAPA
% ---
\includepdf{./05-extras/capa.pdf}
% ---


% ---
% FOLHA DE ROSTO
% (o * indica que haverá a ficha bibliográfica)
% ---
\imprimirfolhaderosto
% ---


% ---
% Ficha bibliografica
% ---

% Isto é um exemplo de Ficha Catalográfica, ou ``Dados internacionais de
% catalogação-na-publicação''. Você pode utilizar este modelo como referência. 
% Porém, provavelmente a biblioteca da sua universidade lhe fornecerá um PDF
% com a ficha catalográfica definitiva após a defesa do trabalho. Quando estiver
% com o documento, salve-o como PDF no diretório do seu projeto e substitua todo
% o conteúdo de implementação deste arquivo pelo comando abaixo:
%
% \usepackage{pdfpages}		% necessário para comando \includepdf
% \begin{fichacatalografica}
%     \includepdf{fig_ficha_catalografica.pdf}
% \end{fichacatalografica}

% \begin{fichacatalografica}
% 	\sffamily
% 	\vspace*{\fill}                             % Posição vertical
% 	\begin{center}                              % Minipage Centralizado
% 	\fbox{\begin{minipage}[c][8cm]{13.5cm}		% Largura
% 	\small
% 	\imprimirautor
% 	%Sobrenome, Nome do autor
	
% 	\hspace{0.5cm} \imprimirtitulo  / \imprimirautor. --
% 	\imprimirlocal, \imprimirdata-
	
% 	\hspace{0.5cm} \thelastpage p. : il. (algumas color.) ; 30 cm.\\
	
% 	\hspace{0.5cm} \imprimirorientadorRotulo~\imprimirorientador\\
	
% 	\hspace{0.5cm}
% 	\parbox[t]{\textwidth}{\imprimirtipotrabalho~--~\imprimirinstituicao,
% 	\imprimirdata.}\\
	
% 	\hspace{0.5cm}
% 		1. Palavra-chave1.
% 		2. Palavra-chave2.
% 		2. Palavra-chave3.
% 		I. Orientador.
% 		II. Universidade xxx.
% 		III. Faculdade de xxx.
% 		IV. Título 			
% 	\end{minipage}}
% 	\end{center}
% \end{fichacatalografica}
% ---


% ---
% FOLHA DE APROVAÇÃO
% ---

% Isto é um exemplo de Folha de aprovação, elemento obrigatório da NBR
% 14724/2011 (seção 4.2.1.3). Você pode utilizar este modelo até a aprovação
% do trabalho. Após isso, substitua todo o conteúdo deste arquivo por uma
% imagem da página assinada pela banca com o comando abaixo:
%
% \begin{folhadeaprovacao}
% \includepdf{folhadeaprovacao_final.pdf}
% \end{folhadeaprovacao}
%
\setlength{\ABNTEXsignskip}{2cm}
\setlength{\ABNTEXsignwidth}{10cm}

\begin{folhadeaprovacao}
  \begin{center}
    {\ABNTEXchapterfont\large\imprimirautor}

    \vspace*{1cm}

    \begin{center}
      \ABNTEXchapterfont\bfseries\Large\imprimirtitulo
    \end{center}

    \vspace*{1cm}
    
    \hspace{.45\textwidth}
    \begin{minipage}{.5\textwidth}
        \imprimirpreambulo
    \end{minipage}%
   \end{center}

   \vspace*{1cm}
        
   \begin{center}
   Aprovado em 17 de julho de 2024.

   \vspace*{1cm}

   Banca examinadora:
   \end{center}

   \vspace*{-1cm}

   \assinatura{Prof. Ma. Caroline Santa Maria Rodrigues -- PUCRS}
   \assinatura{Prof. Dra. Clarissa Pinto Pizarro de Freitas -- PUCRS}
   \assinatura{\imprimirorientador -- PUCRS}
\end{folhadeaprovacao}
% ---


% ---
% Agradecimentos
% ---
\begin{agradecimentos}
Agradeço aos meus pais, Ademilson José Dias e Irinês Porfírio Dias, pelo apoio incondicional que me deram ao longo da vida,
incentivando mesmo as ideias questionáveis como aquela de fazer uma segunda graduação.

A minha irmã, Maria Eduarda Dias, cuja trajetória acadêmica é fonte de inspiração.

A minha namorada, Nathalia Fallavena Ceratti, que acompanhou a esta jornada de perto, oferecendo apoio e compreensão que
fizeram toda a diferença.

Aos meus amigos, Guilherme da Silva Petró, Norton Neves e Uilinton Luciano Marques da Silva, que foram compreensíveis nos
finais de semana em que não pude jogar RPG.

Ao professor Wagner de Lara Machado pela orientação e apoio no desenvolvimento deste trabalho.
\end{agradecimentos}
% ---


% ---
% RESUMOS
% ---
% Resumo em português
\setlength{\absparsep}{18pt} % ajusta o espaçamento dos parágrafos do resumo
\begin{resumo}
Desenvolvimentos na área de aprendizagem de máquina representam uma série de oportunidades para diversas áreas de pesquisa e de atuação profissional.
O uso de modelos de aprendizagem de máquina na predição de desfecho para tratamentos em saúde mental tem sido investigado e apresenta resultados promissores.
A aplicação desse tipo de modelo possibilitaria, por exemplo, maior assertividade na personalização de tratamentos psicoterápicos a nível individual.
Apesar dos benefícios potenciais, o tema ainda é pouco discutido nos espaços de pesquisa, formação e atuação em psicologia clínica. Este trabalho tem
por objetivo introduzir o tema de aprendizagem de máquina a psicólogos. Busca-se apresentar, de maneira acessível, os conceitos e terminologia básicos
em aprendizagem de máquina, proporcionando um panorama geral da área. Para complementar a exposição conceitual, é construído um exemplo a partir de dados
de um estudo sobre eficácia de intervenções em psicologia positiva para depressão via internet. Por fim, discutem-se algumas implicações do uso desse tipo
de tecnologia para a pesquisa e para a prática em psicoterapia.

 \textbf{Palavras-chave}: aprendizagem de máquina, psicoterapia, psicologia positiva, depressão.
\end{resumo}

% Resumo em inglês
\begin{resumo}[Abstract]
 \begin{otherlanguage*}{english}
   Advances in machine learning have opened up a multitude of opportunities for various research and professional fields. The use of machine learning models to
   predict treatment outcomes in mental health has been investigated showing promising results. The application of such models would allow greater assertiveness
   in tailoring psychotherapeutic treatments at the individual level. Despite its potential benefits, the topic remains unexplored in the fields of research, training
   and practice in clinical psychology. This work aims to introduce the topic of machine learning to psychologists. It presents the basic machine learning concepts and
   terminology in an accessible manner, providing a comprehensive overview of the field. To complement the conceptual discussion, an example is constructed using data
   from a study on the effectiveness of internet-based positive psychology interventions for depression. Finally, implications of using this type of technology in
   psychotherapy research and practice are discussed.
   \vspace{\onelineskip}
 
   \noindent 
   \textbf{Keywords}: machine learning, psychotherapy, positive psychology, depression.
 \end{otherlanguage*}
\end{resumo}


% ---
% LISTA DE ILUSTRAÇÕES
% ---
\pdfbookmark[0]{\listfigurename}{lof}
\listoffigures*
\cleardoublepage
% ---


% ---
% LISTA DE QUADROS
% ---
% \pdfbookmark[0]{\listofquadrosname}{loq}
% \listofquadros*
% \cleardoublepage
% ---


% ---
% LISTA DE TABELAS
% ---
% \pdfbookmark[0]{\listtablename}{lot}
% \listoftables*
% \cleardoublepage
% ---


% ---
% LISTA DE ABREVIATURAS E SIGLAS
% ---
% \begin{siglas}
%   \item[ABNT] Associação Brasileira de Normas Técnicas
%   \item[abnTeX] ABsurdas Normas para TeX
% \end{siglas}
% ---


% ---
% LISTA DE SÍMBOLOS
% ---
% \begin{simbolos}
%   \item[$ \Gamma $] Letra grega Gama
%   \item[$ \Lambda $] Lambda
%   \item[$ \zeta $] Letra grega minúscula zeta
%   \item[$ \in $] Pertence
% \end{simbolos}
% ---


% ---
% Sumário
% ---
\pdfbookmark[0]{\contentsname}{toc}
\tableofcontents*
\cleardoublepage
% ---


% ----------------------------------------------------------
% ELEMENTOS TEXTUAIS
% ----------------------------------------------------------
\textual

\include{./01-introducao/main.tex}

\include{./02-machine-learning/main.tex}

\include{./03-exemplo/main.tex}

\include{./04-conclusao/main.tex}

% ----------------------------------------------------------
% ELEMENTOS PÓS-TEXTUAIS
% ----------------------------------------------------------
\postextual
% ----------------------------------------------------------


% ----------------------------------------------------------
% REFERÊNCIAS
% ----------------------------------------------------------
\bibliography{references.bib}


% ----------------------------------------------------------
% GLOSSÁRIO
% ----------------------------------------------------------
%
% Consulte o manual da classe abntex2 para orientações sobre o glossário.
%
%\glossary

% ----------------------------------------------------------
% APÊNDICES
% ----------------------------------------------------------

% ---
% Inicia os apêndices
% ---
% \begin{apendicesenv}

% % Imprime uma página indicando o início dos apêndices
% \partapendices

% % ----------------------------------------------------------
% \chapter{Quisque libero justo}
% % ----------------------------------------------------------

% % ----------------------------------------------------------
% \chapter{Nullam elementum urna}
% % ----------------------------------------------------------

% \end{apendicesenv}
% ---


% ----------------------------------------------------------
% ANEXOS
% ----------------------------------------------------------

% ---
% Inicia os anexos
% ---
% \begin{anexosenv}

% % Imprime uma página indicando o início dos anexos
% \partanexos

% % ---
% \chapter{Morbi ultrices rutrum lorem}
% % ---

% % ---
% \chapter{Cras non urna sed feugiat cum}
% % ---


% % ---
% \chapter{Fusce facilisis lacinia dui}
% % ---

% \end{anexosenv}

%---------------------------------------------------------------------
% INDICE REMISSIVO
%---------------------------------------------------------------------
% \phantompart
% \printindex
%---------------------------------------------------------------------
\includepdf{./05-extras/contracapa.pdf}
\end{document}


% ----------------------------------------------------------
% ELEMENTOS PÓS-TEXTUAIS
% ----------------------------------------------------------
\postextual
% ----------------------------------------------------------


% ----------------------------------------------------------
% REFERÊNCIAS
% ----------------------------------------------------------
\bibliography{references.bib}


% ----------------------------------------------------------
% GLOSSÁRIO
% ----------------------------------------------------------
%
% Consulte o manual da classe abntex2 para orientações sobre o glossário.
%
%\glossary

% ----------------------------------------------------------
% APÊNDICES
% ----------------------------------------------------------

% ---
% Inicia os apêndices
% ---
% \begin{apendicesenv}

% % Imprime uma página indicando o início dos apêndices
% \partapendices

% % ----------------------------------------------------------
% \chapter{Quisque libero justo}
% % ----------------------------------------------------------

% % ----------------------------------------------------------
% \chapter{Nullam elementum urna}
% % ----------------------------------------------------------

% \end{apendicesenv}
% ---


% ----------------------------------------------------------
% ANEXOS
% ----------------------------------------------------------

% ---
% Inicia os anexos
% ---
% \begin{anexosenv}

% % Imprime uma página indicando o início dos anexos
% \partanexos

% % ---
% \chapter{Morbi ultrices rutrum lorem}
% % ---

% % ---
% \chapter{Cras non urna sed feugiat cum}
% % ---


% % ---
% \chapter{Fusce facilisis lacinia dui}
% % ---

% \end{anexosenv}

%---------------------------------------------------------------------
% INDICE REMISSIVO
%---------------------------------------------------------------------
% \phantompart
% \printindex
%---------------------------------------------------------------------
\includepdf{./05-extras/contracapa.pdf}
\end{document}


\documentclass[
	% Opções da classe memoir
	12pt,				        % tamanho da fonte
	openany,			        % capítulos começam em pág ímpar (insere página vazia caso preciso)
	oneside,			        % para impressão em recto e verso. Oposto a oneside
	a4paper,			        % tamanho do papel. 
	% Opções da classe abntex2
	%chapter=TITLE,		        % títulos de capítulos convertidos em letras maiúsculas
	%section=TITLE,		        % títulos de seções convertidos em letras maiúsculas
	%subsection=TITLE,	        % títulos de subseções convertidos em letras maiúsculas
	%subsubsection=TITLE,       % títulos de subsubseções convertidos em letras maiúsculas
	% Opções do pacote babel
	english,			        % idioma adicional para hifenização
	brazil				        % o último idioma é o principal do documento
	]{abntex2}

% ---
% Pacotes básicos 
% ---
\usepackage{lmodern}			% Usa a fonte Latin Modern			
\usepackage[T1]{fontenc}		% Selecao de codigos de fonte.
\usepackage[utf8]{inputenc}		% Codificacao do documento (conversão automática dos acentos)
\usepackage{indentfirst}		% Indenta o primeiro parágrafo de cada seção.
\usepackage{color}				% Controle das cores
\usepackage{graphicx}			% Inclusão de gráficos
\usepackage{microtype} 			% para melhorias de justificação
% ---
		

% ---
% Pacotes de citações
% ---
% \usepackage[alf]{abntex2cite}						% Citações padrão ABNT
\usepackage[
  alf,
  abnt-nbr10520=1988,                     % Citações em sentence case
  abnt-emphasize=bf
  ]{abntex2cite}


% ---
% Pacotes adicionais
% ---
\usepackage{pdfpages}	% Inclusão de páginas em PDF


% ---
% Informações para CAPA e FOLHA DE ROSTO
% ---
\titulo{APRENDIZAGEM DE MÁQUINA PARA PREDIÇÃO DE DESFECHOS EM SAÚDE MENTAL: \textmd{UMA INTRODUÇÃO PARA PSICÓLOGOS}}
\autor{GUILHERME AUGUSTO DIAS}
\local{Porto Alegre}
\data{2024}
\orientador{Prof. Dr. Wagner de Lara Machado}
% \coorientador{}
% \instituicao{
%   Pontifícia Universidade Católica do Rio Grande do Sul
%   \par
%   Escola de Ciências da Saúde e da Vida
%   \par
%   Curso de Psicologia}
\tipotrabalho{TCC}
\preambulo{Trabalho de Conclusão de Curso apresentado como requisito para a obtenção do grau de Bacharel em Psicologia pela Escola de Ciências da Saúde e da Vida da Pontifícia Universidade Católica do Rio Grande do Sul.}
% ---

% ---
% Configurações de aparência do PDF final
% ---

% Alterando o aspecto da cor azul
\definecolor{blue}{RGB}{41,5,195}

% Informações do PDF
\makeatletter
\hypersetup{
		pdftitle={Aprendizagem de máquina para predição de desfechos em saúde mental: uma introdução para psicólogos}, 
		pdfauthor={Guilherme Augusto Dias},
    	pdfsubject={\imprimirpreambulo},
	    pdfcreator={LaTeX with abnTeX2},
		pdfkeywords={aprendizagem de máquina}{psicoterapia}{psicologia positiva}{depressão}, 
		colorlinks=true,		% true: colored links; false: boxed links
    	linkcolor=black,		% color of internal links
    	citecolor=black,		% color of links to bibliography
    	filecolor=black,		% color of file links
		urlcolor=black,
		bookmarksdepth=4
}
\makeatother
% --- 


% ---
% Posiciona figuras e tabelas no topo da página quando adicionadas sozinhas
% em um página em branco. Ver https://github.com/abntex/abntex2/issues/170
\makeatletter
\setlength{\@fptop}{5pt} % Set distance from top of page to first float
\makeatother
% ---

% ---
% Possibilita criação de Quadros e Lista de quadros.
% Ver https://github.com/abntex/abntex2/issues/176
%
\newcommand{\quadroname}{Quadro}
\newcommand{\listofquadrosname}{Lista de quadros}

\newfloat[chapter]{quadro}{loq}{\quadroname}
\newlistof{listofquadros}{loq}{\listofquadrosname}
\newlistentry{quadro}{loq}{0}

% Configurações para atender às regras da ABNT
\setfloatadjustment{quadro}{\centering}
\counterwithout{quadro}{chapter}
\renewcommand{\cftquadroname}{\quadroname\space} 
\renewcommand*{\cftquadroaftersnum}{\hfill--\hfill}

\setfloatlocations{quadro}{hbtp} % Ver https://github.com/abntex/abntex2/issues/176
% ---


% --- 
% Espaçamentos entre linhas e parágrafos 
% --- 

% O tamanho do parágrafo é dado por:
\setlength{\parindent}{1.3cm}

% Controle do espaçamento entre um parágrafo e outro:
\setlength{\parskip}{0.2cm}  % tente também \onelineskip

% ---
% Compila o índice
% ---
\makeindex
% ---


% ----
% Início do documento
% ----
\begin{document}

% Seleciona o idioma do documento (conforme pacotes do babel)
\selectlanguage{brazil}

% Retira espaço extra obsoleto entre as frases.
\frenchspacing 

% ----------------------------------------------------------
% ELEMENTOS PRÉ-TEXTUAIS
% ----------------------------------------------------------
\pretextual

% ---
% CAPA
% ---
\includepdf{./05-extras/capa.pdf}
% ---


% ---
% FOLHA DE ROSTO
% (o * indica que haverá a ficha bibliográfica)
% ---
\imprimirfolhaderosto
% ---


% ---
% Ficha bibliografica
% ---

% Isto é um exemplo de Ficha Catalográfica, ou ``Dados internacionais de
% catalogação-na-publicação''. Você pode utilizar este modelo como referência. 
% Porém, provavelmente a biblioteca da sua universidade lhe fornecerá um PDF
% com a ficha catalográfica definitiva após a defesa do trabalho. Quando estiver
% com o documento, salve-o como PDF no diretório do seu projeto e substitua todo
% o conteúdo de implementação deste arquivo pelo comando abaixo:
%
% \usepackage{pdfpages}		% necessário para comando \includepdf
% \begin{fichacatalografica}
%     \includepdf{fig_ficha_catalografica.pdf}
% \end{fichacatalografica}

% \begin{fichacatalografica}
% 	\sffamily
% 	\vspace*{\fill}                             % Posição vertical
% 	\begin{center}                              % Minipage Centralizado
% 	\fbox{\begin{minipage}[c][8cm]{13.5cm}		% Largura
% 	\small
% 	\imprimirautor
% 	%Sobrenome, Nome do autor
	
% 	\hspace{0.5cm} \imprimirtitulo  / \imprimirautor. --
% 	\imprimirlocal, \imprimirdata-
	
% 	\hspace{0.5cm} \thelastpage p. : il. (algumas color.) ; 30 cm.\\
	
% 	\hspace{0.5cm} \imprimirorientadorRotulo~\imprimirorientador\\
	
% 	\hspace{0.5cm}
% 	\parbox[t]{\textwidth}{\imprimirtipotrabalho~--~\imprimirinstituicao,
% 	\imprimirdata.}\\
	
% 	\hspace{0.5cm}
% 		1. Palavra-chave1.
% 		2. Palavra-chave2.
% 		2. Palavra-chave3.
% 		I. Orientador.
% 		II. Universidade xxx.
% 		III. Faculdade de xxx.
% 		IV. Título 			
% 	\end{minipage}}
% 	\end{center}
% \end{fichacatalografica}
% ---


% ---
% FOLHA DE APROVAÇÃO
% ---

% Isto é um exemplo de Folha de aprovação, elemento obrigatório da NBR
% 14724/2011 (seção 4.2.1.3). Você pode utilizar este modelo até a aprovação
% do trabalho. Após isso, substitua todo o conteúdo deste arquivo por uma
% imagem da página assinada pela banca com o comando abaixo:
%
% \begin{folhadeaprovacao}
% \includepdf{folhadeaprovacao_final.pdf}
% \end{folhadeaprovacao}
%
\setlength{\ABNTEXsignskip}{2cm}
\setlength{\ABNTEXsignwidth}{10cm}

\begin{folhadeaprovacao}
  \begin{center}
    {\ABNTEXchapterfont\large\imprimirautor}

    \vspace*{1cm}

    \begin{center}
      \ABNTEXchapterfont\bfseries\Large\imprimirtitulo
    \end{center}

    \vspace*{1cm}
    
    \hspace{.45\textwidth}
    \begin{minipage}{.5\textwidth}
        \imprimirpreambulo
    \end{minipage}%
   \end{center}

   \vspace*{1cm}
        
   \begin{center}
   Aprovado em 17 de julho de 2024.

   \vspace*{1cm}

   Banca examinadora:
   \end{center}

   \vspace*{-1cm}

   \assinatura{Prof. Ma. Caroline Santa Maria Rodrigues -- PUCRS}
   \assinatura{Prof. Dra. Clarissa Pinto Pizarro de Freitas -- PUCRS}
   \assinatura{\imprimirorientador -- PUCRS}
\end{folhadeaprovacao}
% ---


% ---
% Agradecimentos
% ---
\begin{agradecimentos}
Agradeço aos meus pais, Ademilson José Dias e Irinês Porfírio Dias, pelo apoio incondicional que me deram ao longo da vida,
incentivando mesmo as ideias questionáveis como aquela de fazer uma segunda graduação.

A minha irmã, Maria Eduarda Dias, cuja trajetória acadêmica é fonte de inspiração.

A minha namorada, Nathalia Fallavena Ceratti, que acompanhou a esta jornada de perto, oferecendo apoio e compreensão que
fizeram toda a diferença.

Aos meus amigos, Guilherme da Silva Petró, Norton Neves e Uilinton Luciano Marques da Silva, que foram compreensíveis nos
finais de semana em que não pude jogar RPG.

Ao professor Wagner de Lara Machado pela orientação e apoio no desenvolvimento deste trabalho.
\end{agradecimentos}
% ---


% ---
% RESUMOS
% ---
% Resumo em português
\setlength{\absparsep}{18pt} % ajusta o espaçamento dos parágrafos do resumo
\begin{resumo}
Desenvolvimentos na área de aprendizagem de máquina representam uma série de oportunidades para diversas áreas de pesquisa e de atuação profissional.
O uso de modelos de aprendizagem de máquina na predição de desfecho para tratamentos em saúde mental tem sido investigado e apresenta resultados promissores.
A aplicação desse tipo de modelo possibilitaria, por exemplo, maior assertividade na personalização de tratamentos psicoterápicos a nível individual.
Apesar dos benefícios potenciais, o tema ainda é pouco discutido nos espaços de pesquisa, formação e atuação em psicologia clínica. Este trabalho tem
por objetivo introduzir o tema de aprendizagem de máquina a psicólogos. Busca-se apresentar, de maneira acessível, os conceitos e terminologia básicos
em aprendizagem de máquina, proporcionando um panorama geral da área. Para complementar a exposição conceitual, é construído um exemplo a partir de dados
de um estudo sobre eficácia de intervenções em psicologia positiva para depressão via internet. Por fim, discutem-se algumas implicações do uso desse tipo
de tecnologia para a pesquisa e para a prática em psicoterapia.

 \textbf{Palavras-chave}: aprendizagem de máquina, psicoterapia, psicologia positiva, depressão.
\end{resumo}

% Resumo em inglês
\begin{resumo}[Abstract]
 \begin{otherlanguage*}{english}
   Advances in machine learning have opened up a multitude of opportunities for various research and professional fields. The use of machine learning models to
   predict treatment outcomes in mental health has been investigated showing promising results. The application of such models would allow greater assertiveness
   in tailoring psychotherapeutic treatments at the individual level. Despite its potential benefits, the topic remains unexplored in the fields of research, training
   and practice in clinical psychology. This work aims to introduce the topic of machine learning to psychologists. It presents the basic machine learning concepts and
   terminology in an accessible manner, providing a comprehensive overview of the field. To complement the conceptual discussion, an example is constructed using data
   from a study on the effectiveness of internet-based positive psychology interventions for depression. Finally, implications of using this type of technology in
   psychotherapy research and practice are discussed.
   \vspace{\onelineskip}
 
   \noindent 
   \textbf{Keywords}: machine learning, psychotherapy, positive psychology, depression.
 \end{otherlanguage*}
\end{resumo}


% ---
% LISTA DE ILUSTRAÇÕES
% ---
\pdfbookmark[0]{\listfigurename}{lof}
\listoffigures*
\cleardoublepage
% ---


% ---
% LISTA DE QUADROS
% ---
% \pdfbookmark[0]{\listofquadrosname}{loq}
% \listofquadros*
% \cleardoublepage
% ---


% ---
% LISTA DE TABELAS
% ---
% \pdfbookmark[0]{\listtablename}{lot}
% \listoftables*
% \cleardoublepage
% ---


% ---
% LISTA DE ABREVIATURAS E SIGLAS
% ---
% \begin{siglas}
%   \item[ABNT] Associação Brasileira de Normas Técnicas
%   \item[abnTeX] ABsurdas Normas para TeX
% \end{siglas}
% ---


% ---
% LISTA DE SÍMBOLOS
% ---
% \begin{simbolos}
%   \item[$ \Gamma $] Letra grega Gama
%   \item[$ \Lambda $] Lambda
%   \item[$ \zeta $] Letra grega minúscula zeta
%   \item[$ \in $] Pertence
% \end{simbolos}
% ---


% ---
% Sumário
% ---
\pdfbookmark[0]{\contentsname}{toc}
\tableofcontents*
\cleardoublepage
% ---


% ----------------------------------------------------------
% ELEMENTOS TEXTUAIS
% ----------------------------------------------------------
\textual

\documentclass[
	% Opções da classe memoir
	12pt,				        % tamanho da fonte
	openany,			        % capítulos começam em pág ímpar (insere página vazia caso preciso)
	oneside,			        % para impressão em recto e verso. Oposto a oneside
	a4paper,			        % tamanho do papel. 
	% Opções da classe abntex2
	%chapter=TITLE,		        % títulos de capítulos convertidos em letras maiúsculas
	%section=TITLE,		        % títulos de seções convertidos em letras maiúsculas
	%subsection=TITLE,	        % títulos de subseções convertidos em letras maiúsculas
	%subsubsection=TITLE,       % títulos de subsubseções convertidos em letras maiúsculas
	% Opções do pacote babel
	english,			        % idioma adicional para hifenização
	brazil				        % o último idioma é o principal do documento
	]{abntex2}

% ---
% Pacotes básicos 
% ---
\usepackage{lmodern}			% Usa a fonte Latin Modern			
\usepackage[T1]{fontenc}		% Selecao de codigos de fonte.
\usepackage[utf8]{inputenc}		% Codificacao do documento (conversão automática dos acentos)
\usepackage{indentfirst}		% Indenta o primeiro parágrafo de cada seção.
\usepackage{color}				% Controle das cores
\usepackage{graphicx}			% Inclusão de gráficos
\usepackage{microtype} 			% para melhorias de justificação
% ---
		

% ---
% Pacotes de citações
% ---
% \usepackage[alf]{abntex2cite}						% Citações padrão ABNT
\usepackage[
  alf,
  abnt-nbr10520=1988,                     % Citações em sentence case
  abnt-emphasize=bf
  ]{abntex2cite}


% ---
% Pacotes adicionais
% ---
\usepackage{pdfpages}	% Inclusão de páginas em PDF


% ---
% Informações para CAPA e FOLHA DE ROSTO
% ---
\titulo{APRENDIZAGEM DE MÁQUINA PARA PREDIÇÃO DE DESFECHOS EM SAÚDE MENTAL: \textmd{UMA INTRODUÇÃO PARA PSICÓLOGOS}}
\autor{GUILHERME AUGUSTO DIAS}
\local{Porto Alegre}
\data{2024}
\orientador{Prof. Dr. Wagner de Lara Machado}
% \coorientador{}
% \instituicao{
%   Pontifícia Universidade Católica do Rio Grande do Sul
%   \par
%   Escola de Ciências da Saúde e da Vida
%   \par
%   Curso de Psicologia}
\tipotrabalho{TCC}
\preambulo{Trabalho de Conclusão de Curso apresentado como requisito para a obtenção do grau de Bacharel em Psicologia pela Escola de Ciências da Saúde e da Vida da Pontifícia Universidade Católica do Rio Grande do Sul.}
% ---

% ---
% Configurações de aparência do PDF final
% ---

% Alterando o aspecto da cor azul
\definecolor{blue}{RGB}{41,5,195}

% Informações do PDF
\makeatletter
\hypersetup{
		pdftitle={Aprendizagem de máquina para predição de desfechos em saúde mental: uma introdução para psicólogos}, 
		pdfauthor={Guilherme Augusto Dias},
    	pdfsubject={\imprimirpreambulo},
	    pdfcreator={LaTeX with abnTeX2},
		pdfkeywords={aprendizagem de máquina}{psicoterapia}{psicologia positiva}{depressão}, 
		colorlinks=true,		% true: colored links; false: boxed links
    	linkcolor=black,		% color of internal links
    	citecolor=black,		% color of links to bibliography
    	filecolor=black,		% color of file links
		urlcolor=black,
		bookmarksdepth=4
}
\makeatother
% --- 


% ---
% Posiciona figuras e tabelas no topo da página quando adicionadas sozinhas
% em um página em branco. Ver https://github.com/abntex/abntex2/issues/170
\makeatletter
\setlength{\@fptop}{5pt} % Set distance from top of page to first float
\makeatother
% ---

% ---
% Possibilita criação de Quadros e Lista de quadros.
% Ver https://github.com/abntex/abntex2/issues/176
%
\newcommand{\quadroname}{Quadro}
\newcommand{\listofquadrosname}{Lista de quadros}

\newfloat[chapter]{quadro}{loq}{\quadroname}
\newlistof{listofquadros}{loq}{\listofquadrosname}
\newlistentry{quadro}{loq}{0}

% Configurações para atender às regras da ABNT
\setfloatadjustment{quadro}{\centering}
\counterwithout{quadro}{chapter}
\renewcommand{\cftquadroname}{\quadroname\space} 
\renewcommand*{\cftquadroaftersnum}{\hfill--\hfill}

\setfloatlocations{quadro}{hbtp} % Ver https://github.com/abntex/abntex2/issues/176
% ---


% --- 
% Espaçamentos entre linhas e parágrafos 
% --- 

% O tamanho do parágrafo é dado por:
\setlength{\parindent}{1.3cm}

% Controle do espaçamento entre um parágrafo e outro:
\setlength{\parskip}{0.2cm}  % tente também \onelineskip

% ---
% Compila o índice
% ---
\makeindex
% ---


% ----
% Início do documento
% ----
\begin{document}

% Seleciona o idioma do documento (conforme pacotes do babel)
\selectlanguage{brazil}

% Retira espaço extra obsoleto entre as frases.
\frenchspacing 

% ----------------------------------------------------------
% ELEMENTOS PRÉ-TEXTUAIS
% ----------------------------------------------------------
\pretextual

% ---
% CAPA
% ---
\includepdf{./05-extras/capa.pdf}
% ---


% ---
% FOLHA DE ROSTO
% (o * indica que haverá a ficha bibliográfica)
% ---
\imprimirfolhaderosto
% ---


% ---
% Ficha bibliografica
% ---

% Isto é um exemplo de Ficha Catalográfica, ou ``Dados internacionais de
% catalogação-na-publicação''. Você pode utilizar este modelo como referência. 
% Porém, provavelmente a biblioteca da sua universidade lhe fornecerá um PDF
% com a ficha catalográfica definitiva após a defesa do trabalho. Quando estiver
% com o documento, salve-o como PDF no diretório do seu projeto e substitua todo
% o conteúdo de implementação deste arquivo pelo comando abaixo:
%
% \usepackage{pdfpages}		% necessário para comando \includepdf
% \begin{fichacatalografica}
%     \includepdf{fig_ficha_catalografica.pdf}
% \end{fichacatalografica}

% \begin{fichacatalografica}
% 	\sffamily
% 	\vspace*{\fill}                             % Posição vertical
% 	\begin{center}                              % Minipage Centralizado
% 	\fbox{\begin{minipage}[c][8cm]{13.5cm}		% Largura
% 	\small
% 	\imprimirautor
% 	%Sobrenome, Nome do autor
	
% 	\hspace{0.5cm} \imprimirtitulo  / \imprimirautor. --
% 	\imprimirlocal, \imprimirdata-
	
% 	\hspace{0.5cm} \thelastpage p. : il. (algumas color.) ; 30 cm.\\
	
% 	\hspace{0.5cm} \imprimirorientadorRotulo~\imprimirorientador\\
	
% 	\hspace{0.5cm}
% 	\parbox[t]{\textwidth}{\imprimirtipotrabalho~--~\imprimirinstituicao,
% 	\imprimirdata.}\\
	
% 	\hspace{0.5cm}
% 		1. Palavra-chave1.
% 		2. Palavra-chave2.
% 		2. Palavra-chave3.
% 		I. Orientador.
% 		II. Universidade xxx.
% 		III. Faculdade de xxx.
% 		IV. Título 			
% 	\end{minipage}}
% 	\end{center}
% \end{fichacatalografica}
% ---


% ---
% FOLHA DE APROVAÇÃO
% ---

% Isto é um exemplo de Folha de aprovação, elemento obrigatório da NBR
% 14724/2011 (seção 4.2.1.3). Você pode utilizar este modelo até a aprovação
% do trabalho. Após isso, substitua todo o conteúdo deste arquivo por uma
% imagem da página assinada pela banca com o comando abaixo:
%
% \begin{folhadeaprovacao}
% \includepdf{folhadeaprovacao_final.pdf}
% \end{folhadeaprovacao}
%
\setlength{\ABNTEXsignskip}{2cm}
\setlength{\ABNTEXsignwidth}{10cm}

\begin{folhadeaprovacao}
  \begin{center}
    {\ABNTEXchapterfont\large\imprimirautor}

    \vspace*{1cm}

    \begin{center}
      \ABNTEXchapterfont\bfseries\Large\imprimirtitulo
    \end{center}

    \vspace*{1cm}
    
    \hspace{.45\textwidth}
    \begin{minipage}{.5\textwidth}
        \imprimirpreambulo
    \end{minipage}%
   \end{center}

   \vspace*{1cm}
        
   \begin{center}
   Aprovado em 17 de julho de 2024.

   \vspace*{1cm}

   Banca examinadora:
   \end{center}

   \vspace*{-1cm}

   \assinatura{Prof. Ma. Caroline Santa Maria Rodrigues -- PUCRS}
   \assinatura{Prof. Dra. Clarissa Pinto Pizarro de Freitas -- PUCRS}
   \assinatura{\imprimirorientador -- PUCRS}
\end{folhadeaprovacao}
% ---


% ---
% Agradecimentos
% ---
\begin{agradecimentos}
Agradeço aos meus pais, Ademilson José Dias e Irinês Porfírio Dias, pelo apoio incondicional que me deram ao longo da vida,
incentivando mesmo as ideias questionáveis como aquela de fazer uma segunda graduação.

A minha irmã, Maria Eduarda Dias, cuja trajetória acadêmica é fonte de inspiração.

A minha namorada, Nathalia Fallavena Ceratti, que acompanhou a esta jornada de perto, oferecendo apoio e compreensão que
fizeram toda a diferença.

Aos meus amigos, Guilherme da Silva Petró, Norton Neves e Uilinton Luciano Marques da Silva, que foram compreensíveis nos
finais de semana em que não pude jogar RPG.

Ao professor Wagner de Lara Machado pela orientação e apoio no desenvolvimento deste trabalho.
\end{agradecimentos}
% ---


% ---
% RESUMOS
% ---
% Resumo em português
\setlength{\absparsep}{18pt} % ajusta o espaçamento dos parágrafos do resumo
\begin{resumo}
Desenvolvimentos na área de aprendizagem de máquina representam uma série de oportunidades para diversas áreas de pesquisa e de atuação profissional.
O uso de modelos de aprendizagem de máquina na predição de desfecho para tratamentos em saúde mental tem sido investigado e apresenta resultados promissores.
A aplicação desse tipo de modelo possibilitaria, por exemplo, maior assertividade na personalização de tratamentos psicoterápicos a nível individual.
Apesar dos benefícios potenciais, o tema ainda é pouco discutido nos espaços de pesquisa, formação e atuação em psicologia clínica. Este trabalho tem
por objetivo introduzir o tema de aprendizagem de máquina a psicólogos. Busca-se apresentar, de maneira acessível, os conceitos e terminologia básicos
em aprendizagem de máquina, proporcionando um panorama geral da área. Para complementar a exposição conceitual, é construído um exemplo a partir de dados
de um estudo sobre eficácia de intervenções em psicologia positiva para depressão via internet. Por fim, discutem-se algumas implicações do uso desse tipo
de tecnologia para a pesquisa e para a prática em psicoterapia.

 \textbf{Palavras-chave}: aprendizagem de máquina, psicoterapia, psicologia positiva, depressão.
\end{resumo}

% Resumo em inglês
\begin{resumo}[Abstract]
 \begin{otherlanguage*}{english}
   Advances in machine learning have opened up a multitude of opportunities for various research and professional fields. The use of machine learning models to
   predict treatment outcomes in mental health has been investigated showing promising results. The application of such models would allow greater assertiveness
   in tailoring psychotherapeutic treatments at the individual level. Despite its potential benefits, the topic remains unexplored in the fields of research, training
   and practice in clinical psychology. This work aims to introduce the topic of machine learning to psychologists. It presents the basic machine learning concepts and
   terminology in an accessible manner, providing a comprehensive overview of the field. To complement the conceptual discussion, an example is constructed using data
   from a study on the effectiveness of internet-based positive psychology interventions for depression. Finally, implications of using this type of technology in
   psychotherapy research and practice are discussed.
   \vspace{\onelineskip}
 
   \noindent 
   \textbf{Keywords}: machine learning, psychotherapy, positive psychology, depression.
 \end{otherlanguage*}
\end{resumo}


% ---
% LISTA DE ILUSTRAÇÕES
% ---
\pdfbookmark[0]{\listfigurename}{lof}
\listoffigures*
\cleardoublepage
% ---


% ---
% LISTA DE QUADROS
% ---
% \pdfbookmark[0]{\listofquadrosname}{loq}
% \listofquadros*
% \cleardoublepage
% ---


% ---
% LISTA DE TABELAS
% ---
% \pdfbookmark[0]{\listtablename}{lot}
% \listoftables*
% \cleardoublepage
% ---


% ---
% LISTA DE ABREVIATURAS E SIGLAS
% ---
% \begin{siglas}
%   \item[ABNT] Associação Brasileira de Normas Técnicas
%   \item[abnTeX] ABsurdas Normas para TeX
% \end{siglas}
% ---


% ---
% LISTA DE SÍMBOLOS
% ---
% \begin{simbolos}
%   \item[$ \Gamma $] Letra grega Gama
%   \item[$ \Lambda $] Lambda
%   \item[$ \zeta $] Letra grega minúscula zeta
%   \item[$ \in $] Pertence
% \end{simbolos}
% ---


% ---
% Sumário
% ---
\pdfbookmark[0]{\contentsname}{toc}
\tableofcontents*
\cleardoublepage
% ---


% ----------------------------------------------------------
% ELEMENTOS TEXTUAIS
% ----------------------------------------------------------
\textual

\include{./01-introducao/main.tex}

\include{./02-machine-learning/main.tex}

\include{./03-exemplo/main.tex}

\include{./04-conclusao/main.tex}

% ----------------------------------------------------------
% ELEMENTOS PÓS-TEXTUAIS
% ----------------------------------------------------------
\postextual
% ----------------------------------------------------------


% ----------------------------------------------------------
% REFERÊNCIAS
% ----------------------------------------------------------
\bibliography{references.bib}


% ----------------------------------------------------------
% GLOSSÁRIO
% ----------------------------------------------------------
%
% Consulte o manual da classe abntex2 para orientações sobre o glossário.
%
%\glossary

% ----------------------------------------------------------
% APÊNDICES
% ----------------------------------------------------------

% ---
% Inicia os apêndices
% ---
% \begin{apendicesenv}

% % Imprime uma página indicando o início dos apêndices
% \partapendices

% % ----------------------------------------------------------
% \chapter{Quisque libero justo}
% % ----------------------------------------------------------

% % ----------------------------------------------------------
% \chapter{Nullam elementum urna}
% % ----------------------------------------------------------

% \end{apendicesenv}
% ---


% ----------------------------------------------------------
% ANEXOS
% ----------------------------------------------------------

% ---
% Inicia os anexos
% ---
% \begin{anexosenv}

% % Imprime uma página indicando o início dos anexos
% \partanexos

% % ---
% \chapter{Morbi ultrices rutrum lorem}
% % ---

% % ---
% \chapter{Cras non urna sed feugiat cum}
% % ---


% % ---
% \chapter{Fusce facilisis lacinia dui}
% % ---

% \end{anexosenv}

%---------------------------------------------------------------------
% INDICE REMISSIVO
%---------------------------------------------------------------------
% \phantompart
% \printindex
%---------------------------------------------------------------------
\includepdf{./05-extras/contracapa.pdf}
\end{document}


\documentclass[
	% Opções da classe memoir
	12pt,				        % tamanho da fonte
	openany,			        % capítulos começam em pág ímpar (insere página vazia caso preciso)
	oneside,			        % para impressão em recto e verso. Oposto a oneside
	a4paper,			        % tamanho do papel. 
	% Opções da classe abntex2
	%chapter=TITLE,		        % títulos de capítulos convertidos em letras maiúsculas
	%section=TITLE,		        % títulos de seções convertidos em letras maiúsculas
	%subsection=TITLE,	        % títulos de subseções convertidos em letras maiúsculas
	%subsubsection=TITLE,       % títulos de subsubseções convertidos em letras maiúsculas
	% Opções do pacote babel
	english,			        % idioma adicional para hifenização
	brazil				        % o último idioma é o principal do documento
	]{abntex2}

% ---
% Pacotes básicos 
% ---
\usepackage{lmodern}			% Usa a fonte Latin Modern			
\usepackage[T1]{fontenc}		% Selecao de codigos de fonte.
\usepackage[utf8]{inputenc}		% Codificacao do documento (conversão automática dos acentos)
\usepackage{indentfirst}		% Indenta o primeiro parágrafo de cada seção.
\usepackage{color}				% Controle das cores
\usepackage{graphicx}			% Inclusão de gráficos
\usepackage{microtype} 			% para melhorias de justificação
% ---
		

% ---
% Pacotes de citações
% ---
% \usepackage[alf]{abntex2cite}						% Citações padrão ABNT
\usepackage[
  alf,
  abnt-nbr10520=1988,                     % Citações em sentence case
  abnt-emphasize=bf
  ]{abntex2cite}


% ---
% Pacotes adicionais
% ---
\usepackage{pdfpages}	% Inclusão de páginas em PDF


% ---
% Informações para CAPA e FOLHA DE ROSTO
% ---
\titulo{APRENDIZAGEM DE MÁQUINA PARA PREDIÇÃO DE DESFECHOS EM SAÚDE MENTAL: \textmd{UMA INTRODUÇÃO PARA PSICÓLOGOS}}
\autor{GUILHERME AUGUSTO DIAS}
\local{Porto Alegre}
\data{2024}
\orientador{Prof. Dr. Wagner de Lara Machado}
% \coorientador{}
% \instituicao{
%   Pontifícia Universidade Católica do Rio Grande do Sul
%   \par
%   Escola de Ciências da Saúde e da Vida
%   \par
%   Curso de Psicologia}
\tipotrabalho{TCC}
\preambulo{Trabalho de Conclusão de Curso apresentado como requisito para a obtenção do grau de Bacharel em Psicologia pela Escola de Ciências da Saúde e da Vida da Pontifícia Universidade Católica do Rio Grande do Sul.}
% ---

% ---
% Configurações de aparência do PDF final
% ---

% Alterando o aspecto da cor azul
\definecolor{blue}{RGB}{41,5,195}

% Informações do PDF
\makeatletter
\hypersetup{
		pdftitle={Aprendizagem de máquina para predição de desfechos em saúde mental: uma introdução para psicólogos}, 
		pdfauthor={Guilherme Augusto Dias},
    	pdfsubject={\imprimirpreambulo},
	    pdfcreator={LaTeX with abnTeX2},
		pdfkeywords={aprendizagem de máquina}{psicoterapia}{psicologia positiva}{depressão}, 
		colorlinks=true,		% true: colored links; false: boxed links
    	linkcolor=black,		% color of internal links
    	citecolor=black,		% color of links to bibliography
    	filecolor=black,		% color of file links
		urlcolor=black,
		bookmarksdepth=4
}
\makeatother
% --- 


% ---
% Posiciona figuras e tabelas no topo da página quando adicionadas sozinhas
% em um página em branco. Ver https://github.com/abntex/abntex2/issues/170
\makeatletter
\setlength{\@fptop}{5pt} % Set distance from top of page to first float
\makeatother
% ---

% ---
% Possibilita criação de Quadros e Lista de quadros.
% Ver https://github.com/abntex/abntex2/issues/176
%
\newcommand{\quadroname}{Quadro}
\newcommand{\listofquadrosname}{Lista de quadros}

\newfloat[chapter]{quadro}{loq}{\quadroname}
\newlistof{listofquadros}{loq}{\listofquadrosname}
\newlistentry{quadro}{loq}{0}

% Configurações para atender às regras da ABNT
\setfloatadjustment{quadro}{\centering}
\counterwithout{quadro}{chapter}
\renewcommand{\cftquadroname}{\quadroname\space} 
\renewcommand*{\cftquadroaftersnum}{\hfill--\hfill}

\setfloatlocations{quadro}{hbtp} % Ver https://github.com/abntex/abntex2/issues/176
% ---


% --- 
% Espaçamentos entre linhas e parágrafos 
% --- 

% O tamanho do parágrafo é dado por:
\setlength{\parindent}{1.3cm}

% Controle do espaçamento entre um parágrafo e outro:
\setlength{\parskip}{0.2cm}  % tente também \onelineskip

% ---
% Compila o índice
% ---
\makeindex
% ---


% ----
% Início do documento
% ----
\begin{document}

% Seleciona o idioma do documento (conforme pacotes do babel)
\selectlanguage{brazil}

% Retira espaço extra obsoleto entre as frases.
\frenchspacing 

% ----------------------------------------------------------
% ELEMENTOS PRÉ-TEXTUAIS
% ----------------------------------------------------------
\pretextual

% ---
% CAPA
% ---
\includepdf{./05-extras/capa.pdf}
% ---


% ---
% FOLHA DE ROSTO
% (o * indica que haverá a ficha bibliográfica)
% ---
\imprimirfolhaderosto
% ---


% ---
% Ficha bibliografica
% ---

% Isto é um exemplo de Ficha Catalográfica, ou ``Dados internacionais de
% catalogação-na-publicação''. Você pode utilizar este modelo como referência. 
% Porém, provavelmente a biblioteca da sua universidade lhe fornecerá um PDF
% com a ficha catalográfica definitiva após a defesa do trabalho. Quando estiver
% com o documento, salve-o como PDF no diretório do seu projeto e substitua todo
% o conteúdo de implementação deste arquivo pelo comando abaixo:
%
% \usepackage{pdfpages}		% necessário para comando \includepdf
% \begin{fichacatalografica}
%     \includepdf{fig_ficha_catalografica.pdf}
% \end{fichacatalografica}

% \begin{fichacatalografica}
% 	\sffamily
% 	\vspace*{\fill}                             % Posição vertical
% 	\begin{center}                              % Minipage Centralizado
% 	\fbox{\begin{minipage}[c][8cm]{13.5cm}		% Largura
% 	\small
% 	\imprimirautor
% 	%Sobrenome, Nome do autor
	
% 	\hspace{0.5cm} \imprimirtitulo  / \imprimirautor. --
% 	\imprimirlocal, \imprimirdata-
	
% 	\hspace{0.5cm} \thelastpage p. : il. (algumas color.) ; 30 cm.\\
	
% 	\hspace{0.5cm} \imprimirorientadorRotulo~\imprimirorientador\\
	
% 	\hspace{0.5cm}
% 	\parbox[t]{\textwidth}{\imprimirtipotrabalho~--~\imprimirinstituicao,
% 	\imprimirdata.}\\
	
% 	\hspace{0.5cm}
% 		1. Palavra-chave1.
% 		2. Palavra-chave2.
% 		2. Palavra-chave3.
% 		I. Orientador.
% 		II. Universidade xxx.
% 		III. Faculdade de xxx.
% 		IV. Título 			
% 	\end{minipage}}
% 	\end{center}
% \end{fichacatalografica}
% ---


% ---
% FOLHA DE APROVAÇÃO
% ---

% Isto é um exemplo de Folha de aprovação, elemento obrigatório da NBR
% 14724/2011 (seção 4.2.1.3). Você pode utilizar este modelo até a aprovação
% do trabalho. Após isso, substitua todo o conteúdo deste arquivo por uma
% imagem da página assinada pela banca com o comando abaixo:
%
% \begin{folhadeaprovacao}
% \includepdf{folhadeaprovacao_final.pdf}
% \end{folhadeaprovacao}
%
\setlength{\ABNTEXsignskip}{2cm}
\setlength{\ABNTEXsignwidth}{10cm}

\begin{folhadeaprovacao}
  \begin{center}
    {\ABNTEXchapterfont\large\imprimirautor}

    \vspace*{1cm}

    \begin{center}
      \ABNTEXchapterfont\bfseries\Large\imprimirtitulo
    \end{center}

    \vspace*{1cm}
    
    \hspace{.45\textwidth}
    \begin{minipage}{.5\textwidth}
        \imprimirpreambulo
    \end{minipage}%
   \end{center}

   \vspace*{1cm}
        
   \begin{center}
   Aprovado em 17 de julho de 2024.

   \vspace*{1cm}

   Banca examinadora:
   \end{center}

   \vspace*{-1cm}

   \assinatura{Prof. Ma. Caroline Santa Maria Rodrigues -- PUCRS}
   \assinatura{Prof. Dra. Clarissa Pinto Pizarro de Freitas -- PUCRS}
   \assinatura{\imprimirorientador -- PUCRS}
\end{folhadeaprovacao}
% ---


% ---
% Agradecimentos
% ---
\begin{agradecimentos}
Agradeço aos meus pais, Ademilson José Dias e Irinês Porfírio Dias, pelo apoio incondicional que me deram ao longo da vida,
incentivando mesmo as ideias questionáveis como aquela de fazer uma segunda graduação.

A minha irmã, Maria Eduarda Dias, cuja trajetória acadêmica é fonte de inspiração.

A minha namorada, Nathalia Fallavena Ceratti, que acompanhou a esta jornada de perto, oferecendo apoio e compreensão que
fizeram toda a diferença.

Aos meus amigos, Guilherme da Silva Petró, Norton Neves e Uilinton Luciano Marques da Silva, que foram compreensíveis nos
finais de semana em que não pude jogar RPG.

Ao professor Wagner de Lara Machado pela orientação e apoio no desenvolvimento deste trabalho.
\end{agradecimentos}
% ---


% ---
% RESUMOS
% ---
% Resumo em português
\setlength{\absparsep}{18pt} % ajusta o espaçamento dos parágrafos do resumo
\begin{resumo}
Desenvolvimentos na área de aprendizagem de máquina representam uma série de oportunidades para diversas áreas de pesquisa e de atuação profissional.
O uso de modelos de aprendizagem de máquina na predição de desfecho para tratamentos em saúde mental tem sido investigado e apresenta resultados promissores.
A aplicação desse tipo de modelo possibilitaria, por exemplo, maior assertividade na personalização de tratamentos psicoterápicos a nível individual.
Apesar dos benefícios potenciais, o tema ainda é pouco discutido nos espaços de pesquisa, formação e atuação em psicologia clínica. Este trabalho tem
por objetivo introduzir o tema de aprendizagem de máquina a psicólogos. Busca-se apresentar, de maneira acessível, os conceitos e terminologia básicos
em aprendizagem de máquina, proporcionando um panorama geral da área. Para complementar a exposição conceitual, é construído um exemplo a partir de dados
de um estudo sobre eficácia de intervenções em psicologia positiva para depressão via internet. Por fim, discutem-se algumas implicações do uso desse tipo
de tecnologia para a pesquisa e para a prática em psicoterapia.

 \textbf{Palavras-chave}: aprendizagem de máquina, psicoterapia, psicologia positiva, depressão.
\end{resumo}

% Resumo em inglês
\begin{resumo}[Abstract]
 \begin{otherlanguage*}{english}
   Advances in machine learning have opened up a multitude of opportunities for various research and professional fields. The use of machine learning models to
   predict treatment outcomes in mental health has been investigated showing promising results. The application of such models would allow greater assertiveness
   in tailoring psychotherapeutic treatments at the individual level. Despite its potential benefits, the topic remains unexplored in the fields of research, training
   and practice in clinical psychology. This work aims to introduce the topic of machine learning to psychologists. It presents the basic machine learning concepts and
   terminology in an accessible manner, providing a comprehensive overview of the field. To complement the conceptual discussion, an example is constructed using data
   from a study on the effectiveness of internet-based positive psychology interventions for depression. Finally, implications of using this type of technology in
   psychotherapy research and practice are discussed.
   \vspace{\onelineskip}
 
   \noindent 
   \textbf{Keywords}: machine learning, psychotherapy, positive psychology, depression.
 \end{otherlanguage*}
\end{resumo}


% ---
% LISTA DE ILUSTRAÇÕES
% ---
\pdfbookmark[0]{\listfigurename}{lof}
\listoffigures*
\cleardoublepage
% ---


% ---
% LISTA DE QUADROS
% ---
% \pdfbookmark[0]{\listofquadrosname}{loq}
% \listofquadros*
% \cleardoublepage
% ---


% ---
% LISTA DE TABELAS
% ---
% \pdfbookmark[0]{\listtablename}{lot}
% \listoftables*
% \cleardoublepage
% ---


% ---
% LISTA DE ABREVIATURAS E SIGLAS
% ---
% \begin{siglas}
%   \item[ABNT] Associação Brasileira de Normas Técnicas
%   \item[abnTeX] ABsurdas Normas para TeX
% \end{siglas}
% ---


% ---
% LISTA DE SÍMBOLOS
% ---
% \begin{simbolos}
%   \item[$ \Gamma $] Letra grega Gama
%   \item[$ \Lambda $] Lambda
%   \item[$ \zeta $] Letra grega minúscula zeta
%   \item[$ \in $] Pertence
% \end{simbolos}
% ---


% ---
% Sumário
% ---
\pdfbookmark[0]{\contentsname}{toc}
\tableofcontents*
\cleardoublepage
% ---


% ----------------------------------------------------------
% ELEMENTOS TEXTUAIS
% ----------------------------------------------------------
\textual

\include{./01-introducao/main.tex}

\include{./02-machine-learning/main.tex}

\include{./03-exemplo/main.tex}

\include{./04-conclusao/main.tex}

% ----------------------------------------------------------
% ELEMENTOS PÓS-TEXTUAIS
% ----------------------------------------------------------
\postextual
% ----------------------------------------------------------


% ----------------------------------------------------------
% REFERÊNCIAS
% ----------------------------------------------------------
\bibliography{references.bib}


% ----------------------------------------------------------
% GLOSSÁRIO
% ----------------------------------------------------------
%
% Consulte o manual da classe abntex2 para orientações sobre o glossário.
%
%\glossary

% ----------------------------------------------------------
% APÊNDICES
% ----------------------------------------------------------

% ---
% Inicia os apêndices
% ---
% \begin{apendicesenv}

% % Imprime uma página indicando o início dos apêndices
% \partapendices

% % ----------------------------------------------------------
% \chapter{Quisque libero justo}
% % ----------------------------------------------------------

% % ----------------------------------------------------------
% \chapter{Nullam elementum urna}
% % ----------------------------------------------------------

% \end{apendicesenv}
% ---


% ----------------------------------------------------------
% ANEXOS
% ----------------------------------------------------------

% ---
% Inicia os anexos
% ---
% \begin{anexosenv}

% % Imprime uma página indicando o início dos anexos
% \partanexos

% % ---
% \chapter{Morbi ultrices rutrum lorem}
% % ---

% % ---
% \chapter{Cras non urna sed feugiat cum}
% % ---


% % ---
% \chapter{Fusce facilisis lacinia dui}
% % ---

% \end{anexosenv}

%---------------------------------------------------------------------
% INDICE REMISSIVO
%---------------------------------------------------------------------
% \phantompart
% \printindex
%---------------------------------------------------------------------
\includepdf{./05-extras/contracapa.pdf}
\end{document}


\documentclass[
	% Opções da classe memoir
	12pt,				        % tamanho da fonte
	openany,			        % capítulos começam em pág ímpar (insere página vazia caso preciso)
	oneside,			        % para impressão em recto e verso. Oposto a oneside
	a4paper,			        % tamanho do papel. 
	% Opções da classe abntex2
	%chapter=TITLE,		        % títulos de capítulos convertidos em letras maiúsculas
	%section=TITLE,		        % títulos de seções convertidos em letras maiúsculas
	%subsection=TITLE,	        % títulos de subseções convertidos em letras maiúsculas
	%subsubsection=TITLE,       % títulos de subsubseções convertidos em letras maiúsculas
	% Opções do pacote babel
	english,			        % idioma adicional para hifenização
	brazil				        % o último idioma é o principal do documento
	]{abntex2}

% ---
% Pacotes básicos 
% ---
\usepackage{lmodern}			% Usa a fonte Latin Modern			
\usepackage[T1]{fontenc}		% Selecao de codigos de fonte.
\usepackage[utf8]{inputenc}		% Codificacao do documento (conversão automática dos acentos)
\usepackage{indentfirst}		% Indenta o primeiro parágrafo de cada seção.
\usepackage{color}				% Controle das cores
\usepackage{graphicx}			% Inclusão de gráficos
\usepackage{microtype} 			% para melhorias de justificação
% ---
		

% ---
% Pacotes de citações
% ---
% \usepackage[alf]{abntex2cite}						% Citações padrão ABNT
\usepackage[
  alf,
  abnt-nbr10520=1988,                     % Citações em sentence case
  abnt-emphasize=bf
  ]{abntex2cite}


% ---
% Pacotes adicionais
% ---
\usepackage{pdfpages}	% Inclusão de páginas em PDF


% ---
% Informações para CAPA e FOLHA DE ROSTO
% ---
\titulo{APRENDIZAGEM DE MÁQUINA PARA PREDIÇÃO DE DESFECHOS EM SAÚDE MENTAL: \textmd{UMA INTRODUÇÃO PARA PSICÓLOGOS}}
\autor{GUILHERME AUGUSTO DIAS}
\local{Porto Alegre}
\data{2024}
\orientador{Prof. Dr. Wagner de Lara Machado}
% \coorientador{}
% \instituicao{
%   Pontifícia Universidade Católica do Rio Grande do Sul
%   \par
%   Escola de Ciências da Saúde e da Vida
%   \par
%   Curso de Psicologia}
\tipotrabalho{TCC}
\preambulo{Trabalho de Conclusão de Curso apresentado como requisito para a obtenção do grau de Bacharel em Psicologia pela Escola de Ciências da Saúde e da Vida da Pontifícia Universidade Católica do Rio Grande do Sul.}
% ---

% ---
% Configurações de aparência do PDF final
% ---

% Alterando o aspecto da cor azul
\definecolor{blue}{RGB}{41,5,195}

% Informações do PDF
\makeatletter
\hypersetup{
		pdftitle={Aprendizagem de máquina para predição de desfechos em saúde mental: uma introdução para psicólogos}, 
		pdfauthor={Guilherme Augusto Dias},
    	pdfsubject={\imprimirpreambulo},
	    pdfcreator={LaTeX with abnTeX2},
		pdfkeywords={aprendizagem de máquina}{psicoterapia}{psicologia positiva}{depressão}, 
		colorlinks=true,		% true: colored links; false: boxed links
    	linkcolor=black,		% color of internal links
    	citecolor=black,		% color of links to bibliography
    	filecolor=black,		% color of file links
		urlcolor=black,
		bookmarksdepth=4
}
\makeatother
% --- 


% ---
% Posiciona figuras e tabelas no topo da página quando adicionadas sozinhas
% em um página em branco. Ver https://github.com/abntex/abntex2/issues/170
\makeatletter
\setlength{\@fptop}{5pt} % Set distance from top of page to first float
\makeatother
% ---

% ---
% Possibilita criação de Quadros e Lista de quadros.
% Ver https://github.com/abntex/abntex2/issues/176
%
\newcommand{\quadroname}{Quadro}
\newcommand{\listofquadrosname}{Lista de quadros}

\newfloat[chapter]{quadro}{loq}{\quadroname}
\newlistof{listofquadros}{loq}{\listofquadrosname}
\newlistentry{quadro}{loq}{0}

% Configurações para atender às regras da ABNT
\setfloatadjustment{quadro}{\centering}
\counterwithout{quadro}{chapter}
\renewcommand{\cftquadroname}{\quadroname\space} 
\renewcommand*{\cftquadroaftersnum}{\hfill--\hfill}

\setfloatlocations{quadro}{hbtp} % Ver https://github.com/abntex/abntex2/issues/176
% ---


% --- 
% Espaçamentos entre linhas e parágrafos 
% --- 

% O tamanho do parágrafo é dado por:
\setlength{\parindent}{1.3cm}

% Controle do espaçamento entre um parágrafo e outro:
\setlength{\parskip}{0.2cm}  % tente também \onelineskip

% ---
% Compila o índice
% ---
\makeindex
% ---


% ----
% Início do documento
% ----
\begin{document}

% Seleciona o idioma do documento (conforme pacotes do babel)
\selectlanguage{brazil}

% Retira espaço extra obsoleto entre as frases.
\frenchspacing 

% ----------------------------------------------------------
% ELEMENTOS PRÉ-TEXTUAIS
% ----------------------------------------------------------
\pretextual

% ---
% CAPA
% ---
\includepdf{./05-extras/capa.pdf}
% ---


% ---
% FOLHA DE ROSTO
% (o * indica que haverá a ficha bibliográfica)
% ---
\imprimirfolhaderosto
% ---


% ---
% Ficha bibliografica
% ---

% Isto é um exemplo de Ficha Catalográfica, ou ``Dados internacionais de
% catalogação-na-publicação''. Você pode utilizar este modelo como referência. 
% Porém, provavelmente a biblioteca da sua universidade lhe fornecerá um PDF
% com a ficha catalográfica definitiva após a defesa do trabalho. Quando estiver
% com o documento, salve-o como PDF no diretório do seu projeto e substitua todo
% o conteúdo de implementação deste arquivo pelo comando abaixo:
%
% \usepackage{pdfpages}		% necessário para comando \includepdf
% \begin{fichacatalografica}
%     \includepdf{fig_ficha_catalografica.pdf}
% \end{fichacatalografica}

% \begin{fichacatalografica}
% 	\sffamily
% 	\vspace*{\fill}                             % Posição vertical
% 	\begin{center}                              % Minipage Centralizado
% 	\fbox{\begin{minipage}[c][8cm]{13.5cm}		% Largura
% 	\small
% 	\imprimirautor
% 	%Sobrenome, Nome do autor
	
% 	\hspace{0.5cm} \imprimirtitulo  / \imprimirautor. --
% 	\imprimirlocal, \imprimirdata-
	
% 	\hspace{0.5cm} \thelastpage p. : il. (algumas color.) ; 30 cm.\\
	
% 	\hspace{0.5cm} \imprimirorientadorRotulo~\imprimirorientador\\
	
% 	\hspace{0.5cm}
% 	\parbox[t]{\textwidth}{\imprimirtipotrabalho~--~\imprimirinstituicao,
% 	\imprimirdata.}\\
	
% 	\hspace{0.5cm}
% 		1. Palavra-chave1.
% 		2. Palavra-chave2.
% 		2. Palavra-chave3.
% 		I. Orientador.
% 		II. Universidade xxx.
% 		III. Faculdade de xxx.
% 		IV. Título 			
% 	\end{minipage}}
% 	\end{center}
% \end{fichacatalografica}
% ---


% ---
% FOLHA DE APROVAÇÃO
% ---

% Isto é um exemplo de Folha de aprovação, elemento obrigatório da NBR
% 14724/2011 (seção 4.2.1.3). Você pode utilizar este modelo até a aprovação
% do trabalho. Após isso, substitua todo o conteúdo deste arquivo por uma
% imagem da página assinada pela banca com o comando abaixo:
%
% \begin{folhadeaprovacao}
% \includepdf{folhadeaprovacao_final.pdf}
% \end{folhadeaprovacao}
%
\setlength{\ABNTEXsignskip}{2cm}
\setlength{\ABNTEXsignwidth}{10cm}

\begin{folhadeaprovacao}
  \begin{center}
    {\ABNTEXchapterfont\large\imprimirautor}

    \vspace*{1cm}

    \begin{center}
      \ABNTEXchapterfont\bfseries\Large\imprimirtitulo
    \end{center}

    \vspace*{1cm}
    
    \hspace{.45\textwidth}
    \begin{minipage}{.5\textwidth}
        \imprimirpreambulo
    \end{minipage}%
   \end{center}

   \vspace*{1cm}
        
   \begin{center}
   Aprovado em 17 de julho de 2024.

   \vspace*{1cm}

   Banca examinadora:
   \end{center}

   \vspace*{-1cm}

   \assinatura{Prof. Ma. Caroline Santa Maria Rodrigues -- PUCRS}
   \assinatura{Prof. Dra. Clarissa Pinto Pizarro de Freitas -- PUCRS}
   \assinatura{\imprimirorientador -- PUCRS}
\end{folhadeaprovacao}
% ---


% ---
% Agradecimentos
% ---
\begin{agradecimentos}
Agradeço aos meus pais, Ademilson José Dias e Irinês Porfírio Dias, pelo apoio incondicional que me deram ao longo da vida,
incentivando mesmo as ideias questionáveis como aquela de fazer uma segunda graduação.

A minha irmã, Maria Eduarda Dias, cuja trajetória acadêmica é fonte de inspiração.

A minha namorada, Nathalia Fallavena Ceratti, que acompanhou a esta jornada de perto, oferecendo apoio e compreensão que
fizeram toda a diferença.

Aos meus amigos, Guilherme da Silva Petró, Norton Neves e Uilinton Luciano Marques da Silva, que foram compreensíveis nos
finais de semana em que não pude jogar RPG.

Ao professor Wagner de Lara Machado pela orientação e apoio no desenvolvimento deste trabalho.
\end{agradecimentos}
% ---


% ---
% RESUMOS
% ---
% Resumo em português
\setlength{\absparsep}{18pt} % ajusta o espaçamento dos parágrafos do resumo
\begin{resumo}
Desenvolvimentos na área de aprendizagem de máquina representam uma série de oportunidades para diversas áreas de pesquisa e de atuação profissional.
O uso de modelos de aprendizagem de máquina na predição de desfecho para tratamentos em saúde mental tem sido investigado e apresenta resultados promissores.
A aplicação desse tipo de modelo possibilitaria, por exemplo, maior assertividade na personalização de tratamentos psicoterápicos a nível individual.
Apesar dos benefícios potenciais, o tema ainda é pouco discutido nos espaços de pesquisa, formação e atuação em psicologia clínica. Este trabalho tem
por objetivo introduzir o tema de aprendizagem de máquina a psicólogos. Busca-se apresentar, de maneira acessível, os conceitos e terminologia básicos
em aprendizagem de máquina, proporcionando um panorama geral da área. Para complementar a exposição conceitual, é construído um exemplo a partir de dados
de um estudo sobre eficácia de intervenções em psicologia positiva para depressão via internet. Por fim, discutem-se algumas implicações do uso desse tipo
de tecnologia para a pesquisa e para a prática em psicoterapia.

 \textbf{Palavras-chave}: aprendizagem de máquina, psicoterapia, psicologia positiva, depressão.
\end{resumo}

% Resumo em inglês
\begin{resumo}[Abstract]
 \begin{otherlanguage*}{english}
   Advances in machine learning have opened up a multitude of opportunities for various research and professional fields. The use of machine learning models to
   predict treatment outcomes in mental health has been investigated showing promising results. The application of such models would allow greater assertiveness
   in tailoring psychotherapeutic treatments at the individual level. Despite its potential benefits, the topic remains unexplored in the fields of research, training
   and practice in clinical psychology. This work aims to introduce the topic of machine learning to psychologists. It presents the basic machine learning concepts and
   terminology in an accessible manner, providing a comprehensive overview of the field. To complement the conceptual discussion, an example is constructed using data
   from a study on the effectiveness of internet-based positive psychology interventions for depression. Finally, implications of using this type of technology in
   psychotherapy research and practice are discussed.
   \vspace{\onelineskip}
 
   \noindent 
   \textbf{Keywords}: machine learning, psychotherapy, positive psychology, depression.
 \end{otherlanguage*}
\end{resumo}


% ---
% LISTA DE ILUSTRAÇÕES
% ---
\pdfbookmark[0]{\listfigurename}{lof}
\listoffigures*
\cleardoublepage
% ---


% ---
% LISTA DE QUADROS
% ---
% \pdfbookmark[0]{\listofquadrosname}{loq}
% \listofquadros*
% \cleardoublepage
% ---


% ---
% LISTA DE TABELAS
% ---
% \pdfbookmark[0]{\listtablename}{lot}
% \listoftables*
% \cleardoublepage
% ---


% ---
% LISTA DE ABREVIATURAS E SIGLAS
% ---
% \begin{siglas}
%   \item[ABNT] Associação Brasileira de Normas Técnicas
%   \item[abnTeX] ABsurdas Normas para TeX
% \end{siglas}
% ---


% ---
% LISTA DE SÍMBOLOS
% ---
% \begin{simbolos}
%   \item[$ \Gamma $] Letra grega Gama
%   \item[$ \Lambda $] Lambda
%   \item[$ \zeta $] Letra grega minúscula zeta
%   \item[$ \in $] Pertence
% \end{simbolos}
% ---


% ---
% Sumário
% ---
\pdfbookmark[0]{\contentsname}{toc}
\tableofcontents*
\cleardoublepage
% ---


% ----------------------------------------------------------
% ELEMENTOS TEXTUAIS
% ----------------------------------------------------------
\textual

\include{./01-introducao/main.tex}

\include{./02-machine-learning/main.tex}

\include{./03-exemplo/main.tex}

\include{./04-conclusao/main.tex}

% ----------------------------------------------------------
% ELEMENTOS PÓS-TEXTUAIS
% ----------------------------------------------------------
\postextual
% ----------------------------------------------------------


% ----------------------------------------------------------
% REFERÊNCIAS
% ----------------------------------------------------------
\bibliography{references.bib}


% ----------------------------------------------------------
% GLOSSÁRIO
% ----------------------------------------------------------
%
% Consulte o manual da classe abntex2 para orientações sobre o glossário.
%
%\glossary

% ----------------------------------------------------------
% APÊNDICES
% ----------------------------------------------------------

% ---
% Inicia os apêndices
% ---
% \begin{apendicesenv}

% % Imprime uma página indicando o início dos apêndices
% \partapendices

% % ----------------------------------------------------------
% \chapter{Quisque libero justo}
% % ----------------------------------------------------------

% % ----------------------------------------------------------
% \chapter{Nullam elementum urna}
% % ----------------------------------------------------------

% \end{apendicesenv}
% ---


% ----------------------------------------------------------
% ANEXOS
% ----------------------------------------------------------

% ---
% Inicia os anexos
% ---
% \begin{anexosenv}

% % Imprime uma página indicando o início dos anexos
% \partanexos

% % ---
% \chapter{Morbi ultrices rutrum lorem}
% % ---

% % ---
% \chapter{Cras non urna sed feugiat cum}
% % ---


% % ---
% \chapter{Fusce facilisis lacinia dui}
% % ---

% \end{anexosenv}

%---------------------------------------------------------------------
% INDICE REMISSIVO
%---------------------------------------------------------------------
% \phantompart
% \printindex
%---------------------------------------------------------------------
\includepdf{./05-extras/contracapa.pdf}
\end{document}


\documentclass[
	% Opções da classe memoir
	12pt,				        % tamanho da fonte
	openany,			        % capítulos começam em pág ímpar (insere página vazia caso preciso)
	oneside,			        % para impressão em recto e verso. Oposto a oneside
	a4paper,			        % tamanho do papel. 
	% Opções da classe abntex2
	%chapter=TITLE,		        % títulos de capítulos convertidos em letras maiúsculas
	%section=TITLE,		        % títulos de seções convertidos em letras maiúsculas
	%subsection=TITLE,	        % títulos de subseções convertidos em letras maiúsculas
	%subsubsection=TITLE,       % títulos de subsubseções convertidos em letras maiúsculas
	% Opções do pacote babel
	english,			        % idioma adicional para hifenização
	brazil				        % o último idioma é o principal do documento
	]{abntex2}

% ---
% Pacotes básicos 
% ---
\usepackage{lmodern}			% Usa a fonte Latin Modern			
\usepackage[T1]{fontenc}		% Selecao de codigos de fonte.
\usepackage[utf8]{inputenc}		% Codificacao do documento (conversão automática dos acentos)
\usepackage{indentfirst}		% Indenta o primeiro parágrafo de cada seção.
\usepackage{color}				% Controle das cores
\usepackage{graphicx}			% Inclusão de gráficos
\usepackage{microtype} 			% para melhorias de justificação
% ---
		

% ---
% Pacotes de citações
% ---
% \usepackage[alf]{abntex2cite}						% Citações padrão ABNT
\usepackage[
  alf,
  abnt-nbr10520=1988,                     % Citações em sentence case
  abnt-emphasize=bf
  ]{abntex2cite}


% ---
% Pacotes adicionais
% ---
\usepackage{pdfpages}	% Inclusão de páginas em PDF


% ---
% Informações para CAPA e FOLHA DE ROSTO
% ---
\titulo{APRENDIZAGEM DE MÁQUINA PARA PREDIÇÃO DE DESFECHOS EM SAÚDE MENTAL: \textmd{UMA INTRODUÇÃO PARA PSICÓLOGOS}}
\autor{GUILHERME AUGUSTO DIAS}
\local{Porto Alegre}
\data{2024}
\orientador{Prof. Dr. Wagner de Lara Machado}
% \coorientador{}
% \instituicao{
%   Pontifícia Universidade Católica do Rio Grande do Sul
%   \par
%   Escola de Ciências da Saúde e da Vida
%   \par
%   Curso de Psicologia}
\tipotrabalho{TCC}
\preambulo{Trabalho de Conclusão de Curso apresentado como requisito para a obtenção do grau de Bacharel em Psicologia pela Escola de Ciências da Saúde e da Vida da Pontifícia Universidade Católica do Rio Grande do Sul.}
% ---

% ---
% Configurações de aparência do PDF final
% ---

% Alterando o aspecto da cor azul
\definecolor{blue}{RGB}{41,5,195}

% Informações do PDF
\makeatletter
\hypersetup{
		pdftitle={Aprendizagem de máquina para predição de desfechos em saúde mental: uma introdução para psicólogos}, 
		pdfauthor={Guilherme Augusto Dias},
    	pdfsubject={\imprimirpreambulo},
	    pdfcreator={LaTeX with abnTeX2},
		pdfkeywords={aprendizagem de máquina}{psicoterapia}{psicologia positiva}{depressão}, 
		colorlinks=true,		% true: colored links; false: boxed links
    	linkcolor=black,		% color of internal links
    	citecolor=black,		% color of links to bibliography
    	filecolor=black,		% color of file links
		urlcolor=black,
		bookmarksdepth=4
}
\makeatother
% --- 


% ---
% Posiciona figuras e tabelas no topo da página quando adicionadas sozinhas
% em um página em branco. Ver https://github.com/abntex/abntex2/issues/170
\makeatletter
\setlength{\@fptop}{5pt} % Set distance from top of page to first float
\makeatother
% ---

% ---
% Possibilita criação de Quadros e Lista de quadros.
% Ver https://github.com/abntex/abntex2/issues/176
%
\newcommand{\quadroname}{Quadro}
\newcommand{\listofquadrosname}{Lista de quadros}

\newfloat[chapter]{quadro}{loq}{\quadroname}
\newlistof{listofquadros}{loq}{\listofquadrosname}
\newlistentry{quadro}{loq}{0}

% Configurações para atender às regras da ABNT
\setfloatadjustment{quadro}{\centering}
\counterwithout{quadro}{chapter}
\renewcommand{\cftquadroname}{\quadroname\space} 
\renewcommand*{\cftquadroaftersnum}{\hfill--\hfill}

\setfloatlocations{quadro}{hbtp} % Ver https://github.com/abntex/abntex2/issues/176
% ---


% --- 
% Espaçamentos entre linhas e parágrafos 
% --- 

% O tamanho do parágrafo é dado por:
\setlength{\parindent}{1.3cm}

% Controle do espaçamento entre um parágrafo e outro:
\setlength{\parskip}{0.2cm}  % tente também \onelineskip

% ---
% Compila o índice
% ---
\makeindex
% ---


% ----
% Início do documento
% ----
\begin{document}

% Seleciona o idioma do documento (conforme pacotes do babel)
\selectlanguage{brazil}

% Retira espaço extra obsoleto entre as frases.
\frenchspacing 

% ----------------------------------------------------------
% ELEMENTOS PRÉ-TEXTUAIS
% ----------------------------------------------------------
\pretextual

% ---
% CAPA
% ---
\includepdf{./05-extras/capa.pdf}
% ---


% ---
% FOLHA DE ROSTO
% (o * indica que haverá a ficha bibliográfica)
% ---
\imprimirfolhaderosto
% ---


% ---
% Ficha bibliografica
% ---

% Isto é um exemplo de Ficha Catalográfica, ou ``Dados internacionais de
% catalogação-na-publicação''. Você pode utilizar este modelo como referência. 
% Porém, provavelmente a biblioteca da sua universidade lhe fornecerá um PDF
% com a ficha catalográfica definitiva após a defesa do trabalho. Quando estiver
% com o documento, salve-o como PDF no diretório do seu projeto e substitua todo
% o conteúdo de implementação deste arquivo pelo comando abaixo:
%
% \usepackage{pdfpages}		% necessário para comando \includepdf
% \begin{fichacatalografica}
%     \includepdf{fig_ficha_catalografica.pdf}
% \end{fichacatalografica}

% \begin{fichacatalografica}
% 	\sffamily
% 	\vspace*{\fill}                             % Posição vertical
% 	\begin{center}                              % Minipage Centralizado
% 	\fbox{\begin{minipage}[c][8cm]{13.5cm}		% Largura
% 	\small
% 	\imprimirautor
% 	%Sobrenome, Nome do autor
	
% 	\hspace{0.5cm} \imprimirtitulo  / \imprimirautor. --
% 	\imprimirlocal, \imprimirdata-
	
% 	\hspace{0.5cm} \thelastpage p. : il. (algumas color.) ; 30 cm.\\
	
% 	\hspace{0.5cm} \imprimirorientadorRotulo~\imprimirorientador\\
	
% 	\hspace{0.5cm}
% 	\parbox[t]{\textwidth}{\imprimirtipotrabalho~--~\imprimirinstituicao,
% 	\imprimirdata.}\\
	
% 	\hspace{0.5cm}
% 		1. Palavra-chave1.
% 		2. Palavra-chave2.
% 		2. Palavra-chave3.
% 		I. Orientador.
% 		II. Universidade xxx.
% 		III. Faculdade de xxx.
% 		IV. Título 			
% 	\end{minipage}}
% 	\end{center}
% \end{fichacatalografica}
% ---


% ---
% FOLHA DE APROVAÇÃO
% ---

% Isto é um exemplo de Folha de aprovação, elemento obrigatório da NBR
% 14724/2011 (seção 4.2.1.3). Você pode utilizar este modelo até a aprovação
% do trabalho. Após isso, substitua todo o conteúdo deste arquivo por uma
% imagem da página assinada pela banca com o comando abaixo:
%
% \begin{folhadeaprovacao}
% \includepdf{folhadeaprovacao_final.pdf}
% \end{folhadeaprovacao}
%
\setlength{\ABNTEXsignskip}{2cm}
\setlength{\ABNTEXsignwidth}{10cm}

\begin{folhadeaprovacao}
  \begin{center}
    {\ABNTEXchapterfont\large\imprimirautor}

    \vspace*{1cm}

    \begin{center}
      \ABNTEXchapterfont\bfseries\Large\imprimirtitulo
    \end{center}

    \vspace*{1cm}
    
    \hspace{.45\textwidth}
    \begin{minipage}{.5\textwidth}
        \imprimirpreambulo
    \end{minipage}%
   \end{center}

   \vspace*{1cm}
        
   \begin{center}
   Aprovado em 17 de julho de 2024.

   \vspace*{1cm}

   Banca examinadora:
   \end{center}

   \vspace*{-1cm}

   \assinatura{Prof. Ma. Caroline Santa Maria Rodrigues -- PUCRS}
   \assinatura{Prof. Dra. Clarissa Pinto Pizarro de Freitas -- PUCRS}
   \assinatura{\imprimirorientador -- PUCRS}
\end{folhadeaprovacao}
% ---


% ---
% Agradecimentos
% ---
\begin{agradecimentos}
Agradeço aos meus pais, Ademilson José Dias e Irinês Porfírio Dias, pelo apoio incondicional que me deram ao longo da vida,
incentivando mesmo as ideias questionáveis como aquela de fazer uma segunda graduação.

A minha irmã, Maria Eduarda Dias, cuja trajetória acadêmica é fonte de inspiração.

A minha namorada, Nathalia Fallavena Ceratti, que acompanhou a esta jornada de perto, oferecendo apoio e compreensão que
fizeram toda a diferença.

Aos meus amigos, Guilherme da Silva Petró, Norton Neves e Uilinton Luciano Marques da Silva, que foram compreensíveis nos
finais de semana em que não pude jogar RPG.

Ao professor Wagner de Lara Machado pela orientação e apoio no desenvolvimento deste trabalho.
\end{agradecimentos}
% ---


% ---
% RESUMOS
% ---
% Resumo em português
\setlength{\absparsep}{18pt} % ajusta o espaçamento dos parágrafos do resumo
\begin{resumo}
Desenvolvimentos na área de aprendizagem de máquina representam uma série de oportunidades para diversas áreas de pesquisa e de atuação profissional.
O uso de modelos de aprendizagem de máquina na predição de desfecho para tratamentos em saúde mental tem sido investigado e apresenta resultados promissores.
A aplicação desse tipo de modelo possibilitaria, por exemplo, maior assertividade na personalização de tratamentos psicoterápicos a nível individual.
Apesar dos benefícios potenciais, o tema ainda é pouco discutido nos espaços de pesquisa, formação e atuação em psicologia clínica. Este trabalho tem
por objetivo introduzir o tema de aprendizagem de máquina a psicólogos. Busca-se apresentar, de maneira acessível, os conceitos e terminologia básicos
em aprendizagem de máquina, proporcionando um panorama geral da área. Para complementar a exposição conceitual, é construído um exemplo a partir de dados
de um estudo sobre eficácia de intervenções em psicologia positiva para depressão via internet. Por fim, discutem-se algumas implicações do uso desse tipo
de tecnologia para a pesquisa e para a prática em psicoterapia.

 \textbf{Palavras-chave}: aprendizagem de máquina, psicoterapia, psicologia positiva, depressão.
\end{resumo}

% Resumo em inglês
\begin{resumo}[Abstract]
 \begin{otherlanguage*}{english}
   Advances in machine learning have opened up a multitude of opportunities for various research and professional fields. The use of machine learning models to
   predict treatment outcomes in mental health has been investigated showing promising results. The application of such models would allow greater assertiveness
   in tailoring psychotherapeutic treatments at the individual level. Despite its potential benefits, the topic remains unexplored in the fields of research, training
   and practice in clinical psychology. This work aims to introduce the topic of machine learning to psychologists. It presents the basic machine learning concepts and
   terminology in an accessible manner, providing a comprehensive overview of the field. To complement the conceptual discussion, an example is constructed using data
   from a study on the effectiveness of internet-based positive psychology interventions for depression. Finally, implications of using this type of technology in
   psychotherapy research and practice are discussed.
   \vspace{\onelineskip}
 
   \noindent 
   \textbf{Keywords}: machine learning, psychotherapy, positive psychology, depression.
 \end{otherlanguage*}
\end{resumo}


% ---
% LISTA DE ILUSTRAÇÕES
% ---
\pdfbookmark[0]{\listfigurename}{lof}
\listoffigures*
\cleardoublepage
% ---


% ---
% LISTA DE QUADROS
% ---
% \pdfbookmark[0]{\listofquadrosname}{loq}
% \listofquadros*
% \cleardoublepage
% ---


% ---
% LISTA DE TABELAS
% ---
% \pdfbookmark[0]{\listtablename}{lot}
% \listoftables*
% \cleardoublepage
% ---


% ---
% LISTA DE ABREVIATURAS E SIGLAS
% ---
% \begin{siglas}
%   \item[ABNT] Associação Brasileira de Normas Técnicas
%   \item[abnTeX] ABsurdas Normas para TeX
% \end{siglas}
% ---


% ---
% LISTA DE SÍMBOLOS
% ---
% \begin{simbolos}
%   \item[$ \Gamma $] Letra grega Gama
%   \item[$ \Lambda $] Lambda
%   \item[$ \zeta $] Letra grega minúscula zeta
%   \item[$ \in $] Pertence
% \end{simbolos}
% ---


% ---
% Sumário
% ---
\pdfbookmark[0]{\contentsname}{toc}
\tableofcontents*
\cleardoublepage
% ---


% ----------------------------------------------------------
% ELEMENTOS TEXTUAIS
% ----------------------------------------------------------
\textual

\include{./01-introducao/main.tex}

\include{./02-machine-learning/main.tex}

\include{./03-exemplo/main.tex}

\include{./04-conclusao/main.tex}

% ----------------------------------------------------------
% ELEMENTOS PÓS-TEXTUAIS
% ----------------------------------------------------------
\postextual
% ----------------------------------------------------------


% ----------------------------------------------------------
% REFERÊNCIAS
% ----------------------------------------------------------
\bibliography{references.bib}


% ----------------------------------------------------------
% GLOSSÁRIO
% ----------------------------------------------------------
%
% Consulte o manual da classe abntex2 para orientações sobre o glossário.
%
%\glossary

% ----------------------------------------------------------
% APÊNDICES
% ----------------------------------------------------------

% ---
% Inicia os apêndices
% ---
% \begin{apendicesenv}

% % Imprime uma página indicando o início dos apêndices
% \partapendices

% % ----------------------------------------------------------
% \chapter{Quisque libero justo}
% % ----------------------------------------------------------

% % ----------------------------------------------------------
% \chapter{Nullam elementum urna}
% % ----------------------------------------------------------

% \end{apendicesenv}
% ---


% ----------------------------------------------------------
% ANEXOS
% ----------------------------------------------------------

% ---
% Inicia os anexos
% ---
% \begin{anexosenv}

% % Imprime uma página indicando o início dos anexos
% \partanexos

% % ---
% \chapter{Morbi ultrices rutrum lorem}
% % ---

% % ---
% \chapter{Cras non urna sed feugiat cum}
% % ---


% % ---
% \chapter{Fusce facilisis lacinia dui}
% % ---

% \end{anexosenv}

%---------------------------------------------------------------------
% INDICE REMISSIVO
%---------------------------------------------------------------------
% \phantompart
% \printindex
%---------------------------------------------------------------------
\includepdf{./05-extras/contracapa.pdf}
\end{document}


% ----------------------------------------------------------
% ELEMENTOS PÓS-TEXTUAIS
% ----------------------------------------------------------
\postextual
% ----------------------------------------------------------


% ----------------------------------------------------------
% REFERÊNCIAS
% ----------------------------------------------------------
\bibliography{references.bib}


% ----------------------------------------------------------
% GLOSSÁRIO
% ----------------------------------------------------------
%
% Consulte o manual da classe abntex2 para orientações sobre o glossário.
%
%\glossary

% ----------------------------------------------------------
% APÊNDICES
% ----------------------------------------------------------

% ---
% Inicia os apêndices
% ---
% \begin{apendicesenv}

% % Imprime uma página indicando o início dos apêndices
% \partapendices

% % ----------------------------------------------------------
% \chapter{Quisque libero justo}
% % ----------------------------------------------------------

% % ----------------------------------------------------------
% \chapter{Nullam elementum urna}
% % ----------------------------------------------------------

% \end{apendicesenv}
% ---


% ----------------------------------------------------------
% ANEXOS
% ----------------------------------------------------------

% ---
% Inicia os anexos
% ---
% \begin{anexosenv}

% % Imprime uma página indicando o início dos anexos
% \partanexos

% % ---
% \chapter{Morbi ultrices rutrum lorem}
% % ---

% % ---
% \chapter{Cras non urna sed feugiat cum}
% % ---


% % ---
% \chapter{Fusce facilisis lacinia dui}
% % ---

% \end{anexosenv}

%---------------------------------------------------------------------
% INDICE REMISSIVO
%---------------------------------------------------------------------
% \phantompart
% \printindex
%---------------------------------------------------------------------
\includepdf{./05-extras/contracapa.pdf}
\end{document}


\documentclass[
	% Opções da classe memoir
	12pt,				        % tamanho da fonte
	openany,			        % capítulos começam em pág ímpar (insere página vazia caso preciso)
	oneside,			        % para impressão em recto e verso. Oposto a oneside
	a4paper,			        % tamanho do papel. 
	% Opções da classe abntex2
	%chapter=TITLE,		        % títulos de capítulos convertidos em letras maiúsculas
	%section=TITLE,		        % títulos de seções convertidos em letras maiúsculas
	%subsection=TITLE,	        % títulos de subseções convertidos em letras maiúsculas
	%subsubsection=TITLE,       % títulos de subsubseções convertidos em letras maiúsculas
	% Opções do pacote babel
	english,			        % idioma adicional para hifenização
	brazil				        % o último idioma é o principal do documento
	]{abntex2}

% ---
% Pacotes básicos 
% ---
\usepackage{lmodern}			% Usa a fonte Latin Modern			
\usepackage[T1]{fontenc}		% Selecao de codigos de fonte.
\usepackage[utf8]{inputenc}		% Codificacao do documento (conversão automática dos acentos)
\usepackage{indentfirst}		% Indenta o primeiro parágrafo de cada seção.
\usepackage{color}				% Controle das cores
\usepackage{graphicx}			% Inclusão de gráficos
\usepackage{microtype} 			% para melhorias de justificação
% ---
		

% ---
% Pacotes de citações
% ---
% \usepackage[alf]{abntex2cite}						% Citações padrão ABNT
\usepackage[
  alf,
  abnt-nbr10520=1988,                     % Citações em sentence case
  abnt-emphasize=bf
  ]{abntex2cite}


% ---
% Pacotes adicionais
% ---
\usepackage{pdfpages}	% Inclusão de páginas em PDF


% ---
% Informações para CAPA e FOLHA DE ROSTO
% ---
\titulo{APRENDIZAGEM DE MÁQUINA PARA PREDIÇÃO DE DESFECHOS EM SAÚDE MENTAL: \textmd{UMA INTRODUÇÃO PARA PSICÓLOGOS}}
\autor{GUILHERME AUGUSTO DIAS}
\local{Porto Alegre}
\data{2024}
\orientador{Prof. Dr. Wagner de Lara Machado}
% \coorientador{}
% \instituicao{
%   Pontifícia Universidade Católica do Rio Grande do Sul
%   \par
%   Escola de Ciências da Saúde e da Vida
%   \par
%   Curso de Psicologia}
\tipotrabalho{TCC}
\preambulo{Trabalho de Conclusão de Curso apresentado como requisito para a obtenção do grau de Bacharel em Psicologia pela Escola de Ciências da Saúde e da Vida da Pontifícia Universidade Católica do Rio Grande do Sul.}
% ---

% ---
% Configurações de aparência do PDF final
% ---

% Alterando o aspecto da cor azul
\definecolor{blue}{RGB}{41,5,195}

% Informações do PDF
\makeatletter
\hypersetup{
		pdftitle={Aprendizagem de máquina para predição de desfechos em saúde mental: uma introdução para psicólogos}, 
		pdfauthor={Guilherme Augusto Dias},
    	pdfsubject={\imprimirpreambulo},
	    pdfcreator={LaTeX with abnTeX2},
		pdfkeywords={aprendizagem de máquina}{psicoterapia}{psicologia positiva}{depressão}, 
		colorlinks=true,		% true: colored links; false: boxed links
    	linkcolor=black,		% color of internal links
    	citecolor=black,		% color of links to bibliography
    	filecolor=black,		% color of file links
		urlcolor=black,
		bookmarksdepth=4
}
\makeatother
% --- 


% ---
% Posiciona figuras e tabelas no topo da página quando adicionadas sozinhas
% em um página em branco. Ver https://github.com/abntex/abntex2/issues/170
\makeatletter
\setlength{\@fptop}{5pt} % Set distance from top of page to first float
\makeatother
% ---

% ---
% Possibilita criação de Quadros e Lista de quadros.
% Ver https://github.com/abntex/abntex2/issues/176
%
\newcommand{\quadroname}{Quadro}
\newcommand{\listofquadrosname}{Lista de quadros}

\newfloat[chapter]{quadro}{loq}{\quadroname}
\newlistof{listofquadros}{loq}{\listofquadrosname}
\newlistentry{quadro}{loq}{0}

% Configurações para atender às regras da ABNT
\setfloatadjustment{quadro}{\centering}
\counterwithout{quadro}{chapter}
\renewcommand{\cftquadroname}{\quadroname\space} 
\renewcommand*{\cftquadroaftersnum}{\hfill--\hfill}

\setfloatlocations{quadro}{hbtp} % Ver https://github.com/abntex/abntex2/issues/176
% ---


% --- 
% Espaçamentos entre linhas e parágrafos 
% --- 

% O tamanho do parágrafo é dado por:
\setlength{\parindent}{1.3cm}

% Controle do espaçamento entre um parágrafo e outro:
\setlength{\parskip}{0.2cm}  % tente também \onelineskip

% ---
% Compila o índice
% ---
\makeindex
% ---


% ----
% Início do documento
% ----
\begin{document}

% Seleciona o idioma do documento (conforme pacotes do babel)
\selectlanguage{brazil}

% Retira espaço extra obsoleto entre as frases.
\frenchspacing 

% ----------------------------------------------------------
% ELEMENTOS PRÉ-TEXTUAIS
% ----------------------------------------------------------
\pretextual

% ---
% CAPA
% ---
\includepdf{./05-extras/capa.pdf}
% ---


% ---
% FOLHA DE ROSTO
% (o * indica que haverá a ficha bibliográfica)
% ---
\imprimirfolhaderosto
% ---


% ---
% Ficha bibliografica
% ---

% Isto é um exemplo de Ficha Catalográfica, ou ``Dados internacionais de
% catalogação-na-publicação''. Você pode utilizar este modelo como referência. 
% Porém, provavelmente a biblioteca da sua universidade lhe fornecerá um PDF
% com a ficha catalográfica definitiva após a defesa do trabalho. Quando estiver
% com o documento, salve-o como PDF no diretório do seu projeto e substitua todo
% o conteúdo de implementação deste arquivo pelo comando abaixo:
%
% \usepackage{pdfpages}		% necessário para comando \includepdf
% \begin{fichacatalografica}
%     \includepdf{fig_ficha_catalografica.pdf}
% \end{fichacatalografica}

% \begin{fichacatalografica}
% 	\sffamily
% 	\vspace*{\fill}                             % Posição vertical
% 	\begin{center}                              % Minipage Centralizado
% 	\fbox{\begin{minipage}[c][8cm]{13.5cm}		% Largura
% 	\small
% 	\imprimirautor
% 	%Sobrenome, Nome do autor
	
% 	\hspace{0.5cm} \imprimirtitulo  / \imprimirautor. --
% 	\imprimirlocal, \imprimirdata-
	
% 	\hspace{0.5cm} \thelastpage p. : il. (algumas color.) ; 30 cm.\\
	
% 	\hspace{0.5cm} \imprimirorientadorRotulo~\imprimirorientador\\
	
% 	\hspace{0.5cm}
% 	\parbox[t]{\textwidth}{\imprimirtipotrabalho~--~\imprimirinstituicao,
% 	\imprimirdata.}\\
	
% 	\hspace{0.5cm}
% 		1. Palavra-chave1.
% 		2. Palavra-chave2.
% 		2. Palavra-chave3.
% 		I. Orientador.
% 		II. Universidade xxx.
% 		III. Faculdade de xxx.
% 		IV. Título 			
% 	\end{minipage}}
% 	\end{center}
% \end{fichacatalografica}
% ---


% ---
% FOLHA DE APROVAÇÃO
% ---

% Isto é um exemplo de Folha de aprovação, elemento obrigatório da NBR
% 14724/2011 (seção 4.2.1.3). Você pode utilizar este modelo até a aprovação
% do trabalho. Após isso, substitua todo o conteúdo deste arquivo por uma
% imagem da página assinada pela banca com o comando abaixo:
%
% \begin{folhadeaprovacao}
% \includepdf{folhadeaprovacao_final.pdf}
% \end{folhadeaprovacao}
%
\setlength{\ABNTEXsignskip}{2cm}
\setlength{\ABNTEXsignwidth}{10cm}

\begin{folhadeaprovacao}
  \begin{center}
    {\ABNTEXchapterfont\large\imprimirautor}

    \vspace*{1cm}

    \begin{center}
      \ABNTEXchapterfont\bfseries\Large\imprimirtitulo
    \end{center}

    \vspace*{1cm}
    
    \hspace{.45\textwidth}
    \begin{minipage}{.5\textwidth}
        \imprimirpreambulo
    \end{minipage}%
   \end{center}

   \vspace*{1cm}
        
   \begin{center}
   Aprovado em 17 de julho de 2024.

   \vspace*{1cm}

   Banca examinadora:
   \end{center}

   \vspace*{-1cm}

   \assinatura{Prof. Ma. Caroline Santa Maria Rodrigues -- PUCRS}
   \assinatura{Prof. Dra. Clarissa Pinto Pizarro de Freitas -- PUCRS}
   \assinatura{\imprimirorientador -- PUCRS}
\end{folhadeaprovacao}
% ---


% ---
% Agradecimentos
% ---
\begin{agradecimentos}
Agradeço aos meus pais, Ademilson José Dias e Irinês Porfírio Dias, pelo apoio incondicional que me deram ao longo da vida,
incentivando mesmo as ideias questionáveis como aquela de fazer uma segunda graduação.

A minha irmã, Maria Eduarda Dias, cuja trajetória acadêmica é fonte de inspiração.

A minha namorada, Nathalia Fallavena Ceratti, que acompanhou a esta jornada de perto, oferecendo apoio e compreensão que
fizeram toda a diferença.

Aos meus amigos, Guilherme da Silva Petró, Norton Neves e Uilinton Luciano Marques da Silva, que foram compreensíveis nos
finais de semana em que não pude jogar RPG.

Ao professor Wagner de Lara Machado pela orientação e apoio no desenvolvimento deste trabalho.
\end{agradecimentos}
% ---


% ---
% RESUMOS
% ---
% Resumo em português
\setlength{\absparsep}{18pt} % ajusta o espaçamento dos parágrafos do resumo
\begin{resumo}
Desenvolvimentos na área de aprendizagem de máquina representam uma série de oportunidades para diversas áreas de pesquisa e de atuação profissional.
O uso de modelos de aprendizagem de máquina na predição de desfecho para tratamentos em saúde mental tem sido investigado e apresenta resultados promissores.
A aplicação desse tipo de modelo possibilitaria, por exemplo, maior assertividade na personalização de tratamentos psicoterápicos a nível individual.
Apesar dos benefícios potenciais, o tema ainda é pouco discutido nos espaços de pesquisa, formação e atuação em psicologia clínica. Este trabalho tem
por objetivo introduzir o tema de aprendizagem de máquina a psicólogos. Busca-se apresentar, de maneira acessível, os conceitos e terminologia básicos
em aprendizagem de máquina, proporcionando um panorama geral da área. Para complementar a exposição conceitual, é construído um exemplo a partir de dados
de um estudo sobre eficácia de intervenções em psicologia positiva para depressão via internet. Por fim, discutem-se algumas implicações do uso desse tipo
de tecnologia para a pesquisa e para a prática em psicoterapia.

 \textbf{Palavras-chave}: aprendizagem de máquina, psicoterapia, psicologia positiva, depressão.
\end{resumo}

% Resumo em inglês
\begin{resumo}[Abstract]
 \begin{otherlanguage*}{english}
   Advances in machine learning have opened up a multitude of opportunities for various research and professional fields. The use of machine learning models to
   predict treatment outcomes in mental health has been investigated showing promising results. The application of such models would allow greater assertiveness
   in tailoring psychotherapeutic treatments at the individual level. Despite its potential benefits, the topic remains unexplored in the fields of research, training
   and practice in clinical psychology. This work aims to introduce the topic of machine learning to psychologists. It presents the basic machine learning concepts and
   terminology in an accessible manner, providing a comprehensive overview of the field. To complement the conceptual discussion, an example is constructed using data
   from a study on the effectiveness of internet-based positive psychology interventions for depression. Finally, implications of using this type of technology in
   psychotherapy research and practice are discussed.
   \vspace{\onelineskip}
 
   \noindent 
   \textbf{Keywords}: machine learning, psychotherapy, positive psychology, depression.
 \end{otherlanguage*}
\end{resumo}


% ---
% LISTA DE ILUSTRAÇÕES
% ---
\pdfbookmark[0]{\listfigurename}{lof}
\listoffigures*
\cleardoublepage
% ---


% ---
% LISTA DE QUADROS
% ---
% \pdfbookmark[0]{\listofquadrosname}{loq}
% \listofquadros*
% \cleardoublepage
% ---


% ---
% LISTA DE TABELAS
% ---
% \pdfbookmark[0]{\listtablename}{lot}
% \listoftables*
% \cleardoublepage
% ---


% ---
% LISTA DE ABREVIATURAS E SIGLAS
% ---
% \begin{siglas}
%   \item[ABNT] Associação Brasileira de Normas Técnicas
%   \item[abnTeX] ABsurdas Normas para TeX
% \end{siglas}
% ---


% ---
% LISTA DE SÍMBOLOS
% ---
% \begin{simbolos}
%   \item[$ \Gamma $] Letra grega Gama
%   \item[$ \Lambda $] Lambda
%   \item[$ \zeta $] Letra grega minúscula zeta
%   \item[$ \in $] Pertence
% \end{simbolos}
% ---


% ---
% Sumário
% ---
\pdfbookmark[0]{\contentsname}{toc}
\tableofcontents*
\cleardoublepage
% ---


% ----------------------------------------------------------
% ELEMENTOS TEXTUAIS
% ----------------------------------------------------------
\textual

\documentclass[
	% Opções da classe memoir
	12pt,				        % tamanho da fonte
	openany,			        % capítulos começam em pág ímpar (insere página vazia caso preciso)
	oneside,			        % para impressão em recto e verso. Oposto a oneside
	a4paper,			        % tamanho do papel. 
	% Opções da classe abntex2
	%chapter=TITLE,		        % títulos de capítulos convertidos em letras maiúsculas
	%section=TITLE,		        % títulos de seções convertidos em letras maiúsculas
	%subsection=TITLE,	        % títulos de subseções convertidos em letras maiúsculas
	%subsubsection=TITLE,       % títulos de subsubseções convertidos em letras maiúsculas
	% Opções do pacote babel
	english,			        % idioma adicional para hifenização
	brazil				        % o último idioma é o principal do documento
	]{abntex2}

% ---
% Pacotes básicos 
% ---
\usepackage{lmodern}			% Usa a fonte Latin Modern			
\usepackage[T1]{fontenc}		% Selecao de codigos de fonte.
\usepackage[utf8]{inputenc}		% Codificacao do documento (conversão automática dos acentos)
\usepackage{indentfirst}		% Indenta o primeiro parágrafo de cada seção.
\usepackage{color}				% Controle das cores
\usepackage{graphicx}			% Inclusão de gráficos
\usepackage{microtype} 			% para melhorias de justificação
% ---
		

% ---
% Pacotes de citações
% ---
% \usepackage[alf]{abntex2cite}						% Citações padrão ABNT
\usepackage[
  alf,
  abnt-nbr10520=1988,                     % Citações em sentence case
  abnt-emphasize=bf
  ]{abntex2cite}


% ---
% Pacotes adicionais
% ---
\usepackage{pdfpages}	% Inclusão de páginas em PDF


% ---
% Informações para CAPA e FOLHA DE ROSTO
% ---
\titulo{APRENDIZAGEM DE MÁQUINA PARA PREDIÇÃO DE DESFECHOS EM SAÚDE MENTAL: \textmd{UMA INTRODUÇÃO PARA PSICÓLOGOS}}
\autor{GUILHERME AUGUSTO DIAS}
\local{Porto Alegre}
\data{2024}
\orientador{Prof. Dr. Wagner de Lara Machado}
% \coorientador{}
% \instituicao{
%   Pontifícia Universidade Católica do Rio Grande do Sul
%   \par
%   Escola de Ciências da Saúde e da Vida
%   \par
%   Curso de Psicologia}
\tipotrabalho{TCC}
\preambulo{Trabalho de Conclusão de Curso apresentado como requisito para a obtenção do grau de Bacharel em Psicologia pela Escola de Ciências da Saúde e da Vida da Pontifícia Universidade Católica do Rio Grande do Sul.}
% ---

% ---
% Configurações de aparência do PDF final
% ---

% Alterando o aspecto da cor azul
\definecolor{blue}{RGB}{41,5,195}

% Informações do PDF
\makeatletter
\hypersetup{
		pdftitle={Aprendizagem de máquina para predição de desfechos em saúde mental: uma introdução para psicólogos}, 
		pdfauthor={Guilherme Augusto Dias},
    	pdfsubject={\imprimirpreambulo},
	    pdfcreator={LaTeX with abnTeX2},
		pdfkeywords={aprendizagem de máquina}{psicoterapia}{psicologia positiva}{depressão}, 
		colorlinks=true,		% true: colored links; false: boxed links
    	linkcolor=black,		% color of internal links
    	citecolor=black,		% color of links to bibliography
    	filecolor=black,		% color of file links
		urlcolor=black,
		bookmarksdepth=4
}
\makeatother
% --- 


% ---
% Posiciona figuras e tabelas no topo da página quando adicionadas sozinhas
% em um página em branco. Ver https://github.com/abntex/abntex2/issues/170
\makeatletter
\setlength{\@fptop}{5pt} % Set distance from top of page to first float
\makeatother
% ---

% ---
% Possibilita criação de Quadros e Lista de quadros.
% Ver https://github.com/abntex/abntex2/issues/176
%
\newcommand{\quadroname}{Quadro}
\newcommand{\listofquadrosname}{Lista de quadros}

\newfloat[chapter]{quadro}{loq}{\quadroname}
\newlistof{listofquadros}{loq}{\listofquadrosname}
\newlistentry{quadro}{loq}{0}

% Configurações para atender às regras da ABNT
\setfloatadjustment{quadro}{\centering}
\counterwithout{quadro}{chapter}
\renewcommand{\cftquadroname}{\quadroname\space} 
\renewcommand*{\cftquadroaftersnum}{\hfill--\hfill}

\setfloatlocations{quadro}{hbtp} % Ver https://github.com/abntex/abntex2/issues/176
% ---


% --- 
% Espaçamentos entre linhas e parágrafos 
% --- 

% O tamanho do parágrafo é dado por:
\setlength{\parindent}{1.3cm}

% Controle do espaçamento entre um parágrafo e outro:
\setlength{\parskip}{0.2cm}  % tente também \onelineskip

% ---
% Compila o índice
% ---
\makeindex
% ---


% ----
% Início do documento
% ----
\begin{document}

% Seleciona o idioma do documento (conforme pacotes do babel)
\selectlanguage{brazil}

% Retira espaço extra obsoleto entre as frases.
\frenchspacing 

% ----------------------------------------------------------
% ELEMENTOS PRÉ-TEXTUAIS
% ----------------------------------------------------------
\pretextual

% ---
% CAPA
% ---
\includepdf{./05-extras/capa.pdf}
% ---


% ---
% FOLHA DE ROSTO
% (o * indica que haverá a ficha bibliográfica)
% ---
\imprimirfolhaderosto
% ---


% ---
% Ficha bibliografica
% ---

% Isto é um exemplo de Ficha Catalográfica, ou ``Dados internacionais de
% catalogação-na-publicação''. Você pode utilizar este modelo como referência. 
% Porém, provavelmente a biblioteca da sua universidade lhe fornecerá um PDF
% com a ficha catalográfica definitiva após a defesa do trabalho. Quando estiver
% com o documento, salve-o como PDF no diretório do seu projeto e substitua todo
% o conteúdo de implementação deste arquivo pelo comando abaixo:
%
% \usepackage{pdfpages}		% necessário para comando \includepdf
% \begin{fichacatalografica}
%     \includepdf{fig_ficha_catalografica.pdf}
% \end{fichacatalografica}

% \begin{fichacatalografica}
% 	\sffamily
% 	\vspace*{\fill}                             % Posição vertical
% 	\begin{center}                              % Minipage Centralizado
% 	\fbox{\begin{minipage}[c][8cm]{13.5cm}		% Largura
% 	\small
% 	\imprimirautor
% 	%Sobrenome, Nome do autor
	
% 	\hspace{0.5cm} \imprimirtitulo  / \imprimirautor. --
% 	\imprimirlocal, \imprimirdata-
	
% 	\hspace{0.5cm} \thelastpage p. : il. (algumas color.) ; 30 cm.\\
	
% 	\hspace{0.5cm} \imprimirorientadorRotulo~\imprimirorientador\\
	
% 	\hspace{0.5cm}
% 	\parbox[t]{\textwidth}{\imprimirtipotrabalho~--~\imprimirinstituicao,
% 	\imprimirdata.}\\
	
% 	\hspace{0.5cm}
% 		1. Palavra-chave1.
% 		2. Palavra-chave2.
% 		2. Palavra-chave3.
% 		I. Orientador.
% 		II. Universidade xxx.
% 		III. Faculdade de xxx.
% 		IV. Título 			
% 	\end{minipage}}
% 	\end{center}
% \end{fichacatalografica}
% ---


% ---
% FOLHA DE APROVAÇÃO
% ---

% Isto é um exemplo de Folha de aprovação, elemento obrigatório da NBR
% 14724/2011 (seção 4.2.1.3). Você pode utilizar este modelo até a aprovação
% do trabalho. Após isso, substitua todo o conteúdo deste arquivo por uma
% imagem da página assinada pela banca com o comando abaixo:
%
% \begin{folhadeaprovacao}
% \includepdf{folhadeaprovacao_final.pdf}
% \end{folhadeaprovacao}
%
\setlength{\ABNTEXsignskip}{2cm}
\setlength{\ABNTEXsignwidth}{10cm}

\begin{folhadeaprovacao}
  \begin{center}
    {\ABNTEXchapterfont\large\imprimirautor}

    \vspace*{1cm}

    \begin{center}
      \ABNTEXchapterfont\bfseries\Large\imprimirtitulo
    \end{center}

    \vspace*{1cm}
    
    \hspace{.45\textwidth}
    \begin{minipage}{.5\textwidth}
        \imprimirpreambulo
    \end{minipage}%
   \end{center}

   \vspace*{1cm}
        
   \begin{center}
   Aprovado em 17 de julho de 2024.

   \vspace*{1cm}

   Banca examinadora:
   \end{center}

   \vspace*{-1cm}

   \assinatura{Prof. Ma. Caroline Santa Maria Rodrigues -- PUCRS}
   \assinatura{Prof. Dra. Clarissa Pinto Pizarro de Freitas -- PUCRS}
   \assinatura{\imprimirorientador -- PUCRS}
\end{folhadeaprovacao}
% ---


% ---
% Agradecimentos
% ---
\begin{agradecimentos}
Agradeço aos meus pais, Ademilson José Dias e Irinês Porfírio Dias, pelo apoio incondicional que me deram ao longo da vida,
incentivando mesmo as ideias questionáveis como aquela de fazer uma segunda graduação.

A minha irmã, Maria Eduarda Dias, cuja trajetória acadêmica é fonte de inspiração.

A minha namorada, Nathalia Fallavena Ceratti, que acompanhou a esta jornada de perto, oferecendo apoio e compreensão que
fizeram toda a diferença.

Aos meus amigos, Guilherme da Silva Petró, Norton Neves e Uilinton Luciano Marques da Silva, que foram compreensíveis nos
finais de semana em que não pude jogar RPG.

Ao professor Wagner de Lara Machado pela orientação e apoio no desenvolvimento deste trabalho.
\end{agradecimentos}
% ---


% ---
% RESUMOS
% ---
% Resumo em português
\setlength{\absparsep}{18pt} % ajusta o espaçamento dos parágrafos do resumo
\begin{resumo}
Desenvolvimentos na área de aprendizagem de máquina representam uma série de oportunidades para diversas áreas de pesquisa e de atuação profissional.
O uso de modelos de aprendizagem de máquina na predição de desfecho para tratamentos em saúde mental tem sido investigado e apresenta resultados promissores.
A aplicação desse tipo de modelo possibilitaria, por exemplo, maior assertividade na personalização de tratamentos psicoterápicos a nível individual.
Apesar dos benefícios potenciais, o tema ainda é pouco discutido nos espaços de pesquisa, formação e atuação em psicologia clínica. Este trabalho tem
por objetivo introduzir o tema de aprendizagem de máquina a psicólogos. Busca-se apresentar, de maneira acessível, os conceitos e terminologia básicos
em aprendizagem de máquina, proporcionando um panorama geral da área. Para complementar a exposição conceitual, é construído um exemplo a partir de dados
de um estudo sobre eficácia de intervenções em psicologia positiva para depressão via internet. Por fim, discutem-se algumas implicações do uso desse tipo
de tecnologia para a pesquisa e para a prática em psicoterapia.

 \textbf{Palavras-chave}: aprendizagem de máquina, psicoterapia, psicologia positiva, depressão.
\end{resumo}

% Resumo em inglês
\begin{resumo}[Abstract]
 \begin{otherlanguage*}{english}
   Advances in machine learning have opened up a multitude of opportunities for various research and professional fields. The use of machine learning models to
   predict treatment outcomes in mental health has been investigated showing promising results. The application of such models would allow greater assertiveness
   in tailoring psychotherapeutic treatments at the individual level. Despite its potential benefits, the topic remains unexplored in the fields of research, training
   and practice in clinical psychology. This work aims to introduce the topic of machine learning to psychologists. It presents the basic machine learning concepts and
   terminology in an accessible manner, providing a comprehensive overview of the field. To complement the conceptual discussion, an example is constructed using data
   from a study on the effectiveness of internet-based positive psychology interventions for depression. Finally, implications of using this type of technology in
   psychotherapy research and practice are discussed.
   \vspace{\onelineskip}
 
   \noindent 
   \textbf{Keywords}: machine learning, psychotherapy, positive psychology, depression.
 \end{otherlanguage*}
\end{resumo}


% ---
% LISTA DE ILUSTRAÇÕES
% ---
\pdfbookmark[0]{\listfigurename}{lof}
\listoffigures*
\cleardoublepage
% ---


% ---
% LISTA DE QUADROS
% ---
% \pdfbookmark[0]{\listofquadrosname}{loq}
% \listofquadros*
% \cleardoublepage
% ---


% ---
% LISTA DE TABELAS
% ---
% \pdfbookmark[0]{\listtablename}{lot}
% \listoftables*
% \cleardoublepage
% ---


% ---
% LISTA DE ABREVIATURAS E SIGLAS
% ---
% \begin{siglas}
%   \item[ABNT] Associação Brasileira de Normas Técnicas
%   \item[abnTeX] ABsurdas Normas para TeX
% \end{siglas}
% ---


% ---
% LISTA DE SÍMBOLOS
% ---
% \begin{simbolos}
%   \item[$ \Gamma $] Letra grega Gama
%   \item[$ \Lambda $] Lambda
%   \item[$ \zeta $] Letra grega minúscula zeta
%   \item[$ \in $] Pertence
% \end{simbolos}
% ---


% ---
% Sumário
% ---
\pdfbookmark[0]{\contentsname}{toc}
\tableofcontents*
\cleardoublepage
% ---


% ----------------------------------------------------------
% ELEMENTOS TEXTUAIS
% ----------------------------------------------------------
\textual

\include{./01-introducao/main.tex}

\include{./02-machine-learning/main.tex}

\include{./03-exemplo/main.tex}

\include{./04-conclusao/main.tex}

% ----------------------------------------------------------
% ELEMENTOS PÓS-TEXTUAIS
% ----------------------------------------------------------
\postextual
% ----------------------------------------------------------


% ----------------------------------------------------------
% REFERÊNCIAS
% ----------------------------------------------------------
\bibliography{references.bib}


% ----------------------------------------------------------
% GLOSSÁRIO
% ----------------------------------------------------------
%
% Consulte o manual da classe abntex2 para orientações sobre o glossário.
%
%\glossary

% ----------------------------------------------------------
% APÊNDICES
% ----------------------------------------------------------

% ---
% Inicia os apêndices
% ---
% \begin{apendicesenv}

% % Imprime uma página indicando o início dos apêndices
% \partapendices

% % ----------------------------------------------------------
% \chapter{Quisque libero justo}
% % ----------------------------------------------------------

% % ----------------------------------------------------------
% \chapter{Nullam elementum urna}
% % ----------------------------------------------------------

% \end{apendicesenv}
% ---


% ----------------------------------------------------------
% ANEXOS
% ----------------------------------------------------------

% ---
% Inicia os anexos
% ---
% \begin{anexosenv}

% % Imprime uma página indicando o início dos anexos
% \partanexos

% % ---
% \chapter{Morbi ultrices rutrum lorem}
% % ---

% % ---
% \chapter{Cras non urna sed feugiat cum}
% % ---


% % ---
% \chapter{Fusce facilisis lacinia dui}
% % ---

% \end{anexosenv}

%---------------------------------------------------------------------
% INDICE REMISSIVO
%---------------------------------------------------------------------
% \phantompart
% \printindex
%---------------------------------------------------------------------
\includepdf{./05-extras/contracapa.pdf}
\end{document}


\documentclass[
	% Opções da classe memoir
	12pt,				        % tamanho da fonte
	openany,			        % capítulos começam em pág ímpar (insere página vazia caso preciso)
	oneside,			        % para impressão em recto e verso. Oposto a oneside
	a4paper,			        % tamanho do papel. 
	% Opções da classe abntex2
	%chapter=TITLE,		        % títulos de capítulos convertidos em letras maiúsculas
	%section=TITLE,		        % títulos de seções convertidos em letras maiúsculas
	%subsection=TITLE,	        % títulos de subseções convertidos em letras maiúsculas
	%subsubsection=TITLE,       % títulos de subsubseções convertidos em letras maiúsculas
	% Opções do pacote babel
	english,			        % idioma adicional para hifenização
	brazil				        % o último idioma é o principal do documento
	]{abntex2}

% ---
% Pacotes básicos 
% ---
\usepackage{lmodern}			% Usa a fonte Latin Modern			
\usepackage[T1]{fontenc}		% Selecao de codigos de fonte.
\usepackage[utf8]{inputenc}		% Codificacao do documento (conversão automática dos acentos)
\usepackage{indentfirst}		% Indenta o primeiro parágrafo de cada seção.
\usepackage{color}				% Controle das cores
\usepackage{graphicx}			% Inclusão de gráficos
\usepackage{microtype} 			% para melhorias de justificação
% ---
		

% ---
% Pacotes de citações
% ---
% \usepackage[alf]{abntex2cite}						% Citações padrão ABNT
\usepackage[
  alf,
  abnt-nbr10520=1988,                     % Citações em sentence case
  abnt-emphasize=bf
  ]{abntex2cite}


% ---
% Pacotes adicionais
% ---
\usepackage{pdfpages}	% Inclusão de páginas em PDF


% ---
% Informações para CAPA e FOLHA DE ROSTO
% ---
\titulo{APRENDIZAGEM DE MÁQUINA PARA PREDIÇÃO DE DESFECHOS EM SAÚDE MENTAL: \textmd{UMA INTRODUÇÃO PARA PSICÓLOGOS}}
\autor{GUILHERME AUGUSTO DIAS}
\local{Porto Alegre}
\data{2024}
\orientador{Prof. Dr. Wagner de Lara Machado}
% \coorientador{}
% \instituicao{
%   Pontifícia Universidade Católica do Rio Grande do Sul
%   \par
%   Escola de Ciências da Saúde e da Vida
%   \par
%   Curso de Psicologia}
\tipotrabalho{TCC}
\preambulo{Trabalho de Conclusão de Curso apresentado como requisito para a obtenção do grau de Bacharel em Psicologia pela Escola de Ciências da Saúde e da Vida da Pontifícia Universidade Católica do Rio Grande do Sul.}
% ---

% ---
% Configurações de aparência do PDF final
% ---

% Alterando o aspecto da cor azul
\definecolor{blue}{RGB}{41,5,195}

% Informações do PDF
\makeatletter
\hypersetup{
		pdftitle={Aprendizagem de máquina para predição de desfechos em saúde mental: uma introdução para psicólogos}, 
		pdfauthor={Guilherme Augusto Dias},
    	pdfsubject={\imprimirpreambulo},
	    pdfcreator={LaTeX with abnTeX2},
		pdfkeywords={aprendizagem de máquina}{psicoterapia}{psicologia positiva}{depressão}, 
		colorlinks=true,		% true: colored links; false: boxed links
    	linkcolor=black,		% color of internal links
    	citecolor=black,		% color of links to bibliography
    	filecolor=black,		% color of file links
		urlcolor=black,
		bookmarksdepth=4
}
\makeatother
% --- 


% ---
% Posiciona figuras e tabelas no topo da página quando adicionadas sozinhas
% em um página em branco. Ver https://github.com/abntex/abntex2/issues/170
\makeatletter
\setlength{\@fptop}{5pt} % Set distance from top of page to first float
\makeatother
% ---

% ---
% Possibilita criação de Quadros e Lista de quadros.
% Ver https://github.com/abntex/abntex2/issues/176
%
\newcommand{\quadroname}{Quadro}
\newcommand{\listofquadrosname}{Lista de quadros}

\newfloat[chapter]{quadro}{loq}{\quadroname}
\newlistof{listofquadros}{loq}{\listofquadrosname}
\newlistentry{quadro}{loq}{0}

% Configurações para atender às regras da ABNT
\setfloatadjustment{quadro}{\centering}
\counterwithout{quadro}{chapter}
\renewcommand{\cftquadroname}{\quadroname\space} 
\renewcommand*{\cftquadroaftersnum}{\hfill--\hfill}

\setfloatlocations{quadro}{hbtp} % Ver https://github.com/abntex/abntex2/issues/176
% ---


% --- 
% Espaçamentos entre linhas e parágrafos 
% --- 

% O tamanho do parágrafo é dado por:
\setlength{\parindent}{1.3cm}

% Controle do espaçamento entre um parágrafo e outro:
\setlength{\parskip}{0.2cm}  % tente também \onelineskip

% ---
% Compila o índice
% ---
\makeindex
% ---


% ----
% Início do documento
% ----
\begin{document}

% Seleciona o idioma do documento (conforme pacotes do babel)
\selectlanguage{brazil}

% Retira espaço extra obsoleto entre as frases.
\frenchspacing 

% ----------------------------------------------------------
% ELEMENTOS PRÉ-TEXTUAIS
% ----------------------------------------------------------
\pretextual

% ---
% CAPA
% ---
\includepdf{./05-extras/capa.pdf}
% ---


% ---
% FOLHA DE ROSTO
% (o * indica que haverá a ficha bibliográfica)
% ---
\imprimirfolhaderosto
% ---


% ---
% Ficha bibliografica
% ---

% Isto é um exemplo de Ficha Catalográfica, ou ``Dados internacionais de
% catalogação-na-publicação''. Você pode utilizar este modelo como referência. 
% Porém, provavelmente a biblioteca da sua universidade lhe fornecerá um PDF
% com a ficha catalográfica definitiva após a defesa do trabalho. Quando estiver
% com o documento, salve-o como PDF no diretório do seu projeto e substitua todo
% o conteúdo de implementação deste arquivo pelo comando abaixo:
%
% \usepackage{pdfpages}		% necessário para comando \includepdf
% \begin{fichacatalografica}
%     \includepdf{fig_ficha_catalografica.pdf}
% \end{fichacatalografica}

% \begin{fichacatalografica}
% 	\sffamily
% 	\vspace*{\fill}                             % Posição vertical
% 	\begin{center}                              % Minipage Centralizado
% 	\fbox{\begin{minipage}[c][8cm]{13.5cm}		% Largura
% 	\small
% 	\imprimirautor
% 	%Sobrenome, Nome do autor
	
% 	\hspace{0.5cm} \imprimirtitulo  / \imprimirautor. --
% 	\imprimirlocal, \imprimirdata-
	
% 	\hspace{0.5cm} \thelastpage p. : il. (algumas color.) ; 30 cm.\\
	
% 	\hspace{0.5cm} \imprimirorientadorRotulo~\imprimirorientador\\
	
% 	\hspace{0.5cm}
% 	\parbox[t]{\textwidth}{\imprimirtipotrabalho~--~\imprimirinstituicao,
% 	\imprimirdata.}\\
	
% 	\hspace{0.5cm}
% 		1. Palavra-chave1.
% 		2. Palavra-chave2.
% 		2. Palavra-chave3.
% 		I. Orientador.
% 		II. Universidade xxx.
% 		III. Faculdade de xxx.
% 		IV. Título 			
% 	\end{minipage}}
% 	\end{center}
% \end{fichacatalografica}
% ---


% ---
% FOLHA DE APROVAÇÃO
% ---

% Isto é um exemplo de Folha de aprovação, elemento obrigatório da NBR
% 14724/2011 (seção 4.2.1.3). Você pode utilizar este modelo até a aprovação
% do trabalho. Após isso, substitua todo o conteúdo deste arquivo por uma
% imagem da página assinada pela banca com o comando abaixo:
%
% \begin{folhadeaprovacao}
% \includepdf{folhadeaprovacao_final.pdf}
% \end{folhadeaprovacao}
%
\setlength{\ABNTEXsignskip}{2cm}
\setlength{\ABNTEXsignwidth}{10cm}

\begin{folhadeaprovacao}
  \begin{center}
    {\ABNTEXchapterfont\large\imprimirautor}

    \vspace*{1cm}

    \begin{center}
      \ABNTEXchapterfont\bfseries\Large\imprimirtitulo
    \end{center}

    \vspace*{1cm}
    
    \hspace{.45\textwidth}
    \begin{minipage}{.5\textwidth}
        \imprimirpreambulo
    \end{minipage}%
   \end{center}

   \vspace*{1cm}
        
   \begin{center}
   Aprovado em 17 de julho de 2024.

   \vspace*{1cm}

   Banca examinadora:
   \end{center}

   \vspace*{-1cm}

   \assinatura{Prof. Ma. Caroline Santa Maria Rodrigues -- PUCRS}
   \assinatura{Prof. Dra. Clarissa Pinto Pizarro de Freitas -- PUCRS}
   \assinatura{\imprimirorientador -- PUCRS}
\end{folhadeaprovacao}
% ---


% ---
% Agradecimentos
% ---
\begin{agradecimentos}
Agradeço aos meus pais, Ademilson José Dias e Irinês Porfírio Dias, pelo apoio incondicional que me deram ao longo da vida,
incentivando mesmo as ideias questionáveis como aquela de fazer uma segunda graduação.

A minha irmã, Maria Eduarda Dias, cuja trajetória acadêmica é fonte de inspiração.

A minha namorada, Nathalia Fallavena Ceratti, que acompanhou a esta jornada de perto, oferecendo apoio e compreensão que
fizeram toda a diferença.

Aos meus amigos, Guilherme da Silva Petró, Norton Neves e Uilinton Luciano Marques da Silva, que foram compreensíveis nos
finais de semana em que não pude jogar RPG.

Ao professor Wagner de Lara Machado pela orientação e apoio no desenvolvimento deste trabalho.
\end{agradecimentos}
% ---


% ---
% RESUMOS
% ---
% Resumo em português
\setlength{\absparsep}{18pt} % ajusta o espaçamento dos parágrafos do resumo
\begin{resumo}
Desenvolvimentos na área de aprendizagem de máquina representam uma série de oportunidades para diversas áreas de pesquisa e de atuação profissional.
O uso de modelos de aprendizagem de máquina na predição de desfecho para tratamentos em saúde mental tem sido investigado e apresenta resultados promissores.
A aplicação desse tipo de modelo possibilitaria, por exemplo, maior assertividade na personalização de tratamentos psicoterápicos a nível individual.
Apesar dos benefícios potenciais, o tema ainda é pouco discutido nos espaços de pesquisa, formação e atuação em psicologia clínica. Este trabalho tem
por objetivo introduzir o tema de aprendizagem de máquina a psicólogos. Busca-se apresentar, de maneira acessível, os conceitos e terminologia básicos
em aprendizagem de máquina, proporcionando um panorama geral da área. Para complementar a exposição conceitual, é construído um exemplo a partir de dados
de um estudo sobre eficácia de intervenções em psicologia positiva para depressão via internet. Por fim, discutem-se algumas implicações do uso desse tipo
de tecnologia para a pesquisa e para a prática em psicoterapia.

 \textbf{Palavras-chave}: aprendizagem de máquina, psicoterapia, psicologia positiva, depressão.
\end{resumo}

% Resumo em inglês
\begin{resumo}[Abstract]
 \begin{otherlanguage*}{english}
   Advances in machine learning have opened up a multitude of opportunities for various research and professional fields. The use of machine learning models to
   predict treatment outcomes in mental health has been investigated showing promising results. The application of such models would allow greater assertiveness
   in tailoring psychotherapeutic treatments at the individual level. Despite its potential benefits, the topic remains unexplored in the fields of research, training
   and practice in clinical psychology. This work aims to introduce the topic of machine learning to psychologists. It presents the basic machine learning concepts and
   terminology in an accessible manner, providing a comprehensive overview of the field. To complement the conceptual discussion, an example is constructed using data
   from a study on the effectiveness of internet-based positive psychology interventions for depression. Finally, implications of using this type of technology in
   psychotherapy research and practice are discussed.
   \vspace{\onelineskip}
 
   \noindent 
   \textbf{Keywords}: machine learning, psychotherapy, positive psychology, depression.
 \end{otherlanguage*}
\end{resumo}


% ---
% LISTA DE ILUSTRAÇÕES
% ---
\pdfbookmark[0]{\listfigurename}{lof}
\listoffigures*
\cleardoublepage
% ---


% ---
% LISTA DE QUADROS
% ---
% \pdfbookmark[0]{\listofquadrosname}{loq}
% \listofquadros*
% \cleardoublepage
% ---


% ---
% LISTA DE TABELAS
% ---
% \pdfbookmark[0]{\listtablename}{lot}
% \listoftables*
% \cleardoublepage
% ---


% ---
% LISTA DE ABREVIATURAS E SIGLAS
% ---
% \begin{siglas}
%   \item[ABNT] Associação Brasileira de Normas Técnicas
%   \item[abnTeX] ABsurdas Normas para TeX
% \end{siglas}
% ---


% ---
% LISTA DE SÍMBOLOS
% ---
% \begin{simbolos}
%   \item[$ \Gamma $] Letra grega Gama
%   \item[$ \Lambda $] Lambda
%   \item[$ \zeta $] Letra grega minúscula zeta
%   \item[$ \in $] Pertence
% \end{simbolos}
% ---


% ---
% Sumário
% ---
\pdfbookmark[0]{\contentsname}{toc}
\tableofcontents*
\cleardoublepage
% ---


% ----------------------------------------------------------
% ELEMENTOS TEXTUAIS
% ----------------------------------------------------------
\textual

\include{./01-introducao/main.tex}

\include{./02-machine-learning/main.tex}

\include{./03-exemplo/main.tex}

\include{./04-conclusao/main.tex}

% ----------------------------------------------------------
% ELEMENTOS PÓS-TEXTUAIS
% ----------------------------------------------------------
\postextual
% ----------------------------------------------------------


% ----------------------------------------------------------
% REFERÊNCIAS
% ----------------------------------------------------------
\bibliography{references.bib}


% ----------------------------------------------------------
% GLOSSÁRIO
% ----------------------------------------------------------
%
% Consulte o manual da classe abntex2 para orientações sobre o glossário.
%
%\glossary

% ----------------------------------------------------------
% APÊNDICES
% ----------------------------------------------------------

% ---
% Inicia os apêndices
% ---
% \begin{apendicesenv}

% % Imprime uma página indicando o início dos apêndices
% \partapendices

% % ----------------------------------------------------------
% \chapter{Quisque libero justo}
% % ----------------------------------------------------------

% % ----------------------------------------------------------
% \chapter{Nullam elementum urna}
% % ----------------------------------------------------------

% \end{apendicesenv}
% ---


% ----------------------------------------------------------
% ANEXOS
% ----------------------------------------------------------

% ---
% Inicia os anexos
% ---
% \begin{anexosenv}

% % Imprime uma página indicando o início dos anexos
% \partanexos

% % ---
% \chapter{Morbi ultrices rutrum lorem}
% % ---

% % ---
% \chapter{Cras non urna sed feugiat cum}
% % ---


% % ---
% \chapter{Fusce facilisis lacinia dui}
% % ---

% \end{anexosenv}

%---------------------------------------------------------------------
% INDICE REMISSIVO
%---------------------------------------------------------------------
% \phantompart
% \printindex
%---------------------------------------------------------------------
\includepdf{./05-extras/contracapa.pdf}
\end{document}


\documentclass[
	% Opções da classe memoir
	12pt,				        % tamanho da fonte
	openany,			        % capítulos começam em pág ímpar (insere página vazia caso preciso)
	oneside,			        % para impressão em recto e verso. Oposto a oneside
	a4paper,			        % tamanho do papel. 
	% Opções da classe abntex2
	%chapter=TITLE,		        % títulos de capítulos convertidos em letras maiúsculas
	%section=TITLE,		        % títulos de seções convertidos em letras maiúsculas
	%subsection=TITLE,	        % títulos de subseções convertidos em letras maiúsculas
	%subsubsection=TITLE,       % títulos de subsubseções convertidos em letras maiúsculas
	% Opções do pacote babel
	english,			        % idioma adicional para hifenização
	brazil				        % o último idioma é o principal do documento
	]{abntex2}

% ---
% Pacotes básicos 
% ---
\usepackage{lmodern}			% Usa a fonte Latin Modern			
\usepackage[T1]{fontenc}		% Selecao de codigos de fonte.
\usepackage[utf8]{inputenc}		% Codificacao do documento (conversão automática dos acentos)
\usepackage{indentfirst}		% Indenta o primeiro parágrafo de cada seção.
\usepackage{color}				% Controle das cores
\usepackage{graphicx}			% Inclusão de gráficos
\usepackage{microtype} 			% para melhorias de justificação
% ---
		

% ---
% Pacotes de citações
% ---
% \usepackage[alf]{abntex2cite}						% Citações padrão ABNT
\usepackage[
  alf,
  abnt-nbr10520=1988,                     % Citações em sentence case
  abnt-emphasize=bf
  ]{abntex2cite}


% ---
% Pacotes adicionais
% ---
\usepackage{pdfpages}	% Inclusão de páginas em PDF


% ---
% Informações para CAPA e FOLHA DE ROSTO
% ---
\titulo{APRENDIZAGEM DE MÁQUINA PARA PREDIÇÃO DE DESFECHOS EM SAÚDE MENTAL: \textmd{UMA INTRODUÇÃO PARA PSICÓLOGOS}}
\autor{GUILHERME AUGUSTO DIAS}
\local{Porto Alegre}
\data{2024}
\orientador{Prof. Dr. Wagner de Lara Machado}
% \coorientador{}
% \instituicao{
%   Pontifícia Universidade Católica do Rio Grande do Sul
%   \par
%   Escola de Ciências da Saúde e da Vida
%   \par
%   Curso de Psicologia}
\tipotrabalho{TCC}
\preambulo{Trabalho de Conclusão de Curso apresentado como requisito para a obtenção do grau de Bacharel em Psicologia pela Escola de Ciências da Saúde e da Vida da Pontifícia Universidade Católica do Rio Grande do Sul.}
% ---

% ---
% Configurações de aparência do PDF final
% ---

% Alterando o aspecto da cor azul
\definecolor{blue}{RGB}{41,5,195}

% Informações do PDF
\makeatletter
\hypersetup{
		pdftitle={Aprendizagem de máquina para predição de desfechos em saúde mental: uma introdução para psicólogos}, 
		pdfauthor={Guilherme Augusto Dias},
    	pdfsubject={\imprimirpreambulo},
	    pdfcreator={LaTeX with abnTeX2},
		pdfkeywords={aprendizagem de máquina}{psicoterapia}{psicologia positiva}{depressão}, 
		colorlinks=true,		% true: colored links; false: boxed links
    	linkcolor=black,		% color of internal links
    	citecolor=black,		% color of links to bibliography
    	filecolor=black,		% color of file links
		urlcolor=black,
		bookmarksdepth=4
}
\makeatother
% --- 


% ---
% Posiciona figuras e tabelas no topo da página quando adicionadas sozinhas
% em um página em branco. Ver https://github.com/abntex/abntex2/issues/170
\makeatletter
\setlength{\@fptop}{5pt} % Set distance from top of page to first float
\makeatother
% ---

% ---
% Possibilita criação de Quadros e Lista de quadros.
% Ver https://github.com/abntex/abntex2/issues/176
%
\newcommand{\quadroname}{Quadro}
\newcommand{\listofquadrosname}{Lista de quadros}

\newfloat[chapter]{quadro}{loq}{\quadroname}
\newlistof{listofquadros}{loq}{\listofquadrosname}
\newlistentry{quadro}{loq}{0}

% Configurações para atender às regras da ABNT
\setfloatadjustment{quadro}{\centering}
\counterwithout{quadro}{chapter}
\renewcommand{\cftquadroname}{\quadroname\space} 
\renewcommand*{\cftquadroaftersnum}{\hfill--\hfill}

\setfloatlocations{quadro}{hbtp} % Ver https://github.com/abntex/abntex2/issues/176
% ---


% --- 
% Espaçamentos entre linhas e parágrafos 
% --- 

% O tamanho do parágrafo é dado por:
\setlength{\parindent}{1.3cm}

% Controle do espaçamento entre um parágrafo e outro:
\setlength{\parskip}{0.2cm}  % tente também \onelineskip

% ---
% Compila o índice
% ---
\makeindex
% ---


% ----
% Início do documento
% ----
\begin{document}

% Seleciona o idioma do documento (conforme pacotes do babel)
\selectlanguage{brazil}

% Retira espaço extra obsoleto entre as frases.
\frenchspacing 

% ----------------------------------------------------------
% ELEMENTOS PRÉ-TEXTUAIS
% ----------------------------------------------------------
\pretextual

% ---
% CAPA
% ---
\includepdf{./05-extras/capa.pdf}
% ---


% ---
% FOLHA DE ROSTO
% (o * indica que haverá a ficha bibliográfica)
% ---
\imprimirfolhaderosto
% ---


% ---
% Ficha bibliografica
% ---

% Isto é um exemplo de Ficha Catalográfica, ou ``Dados internacionais de
% catalogação-na-publicação''. Você pode utilizar este modelo como referência. 
% Porém, provavelmente a biblioteca da sua universidade lhe fornecerá um PDF
% com a ficha catalográfica definitiva após a defesa do trabalho. Quando estiver
% com o documento, salve-o como PDF no diretório do seu projeto e substitua todo
% o conteúdo de implementação deste arquivo pelo comando abaixo:
%
% \usepackage{pdfpages}		% necessário para comando \includepdf
% \begin{fichacatalografica}
%     \includepdf{fig_ficha_catalografica.pdf}
% \end{fichacatalografica}

% \begin{fichacatalografica}
% 	\sffamily
% 	\vspace*{\fill}                             % Posição vertical
% 	\begin{center}                              % Minipage Centralizado
% 	\fbox{\begin{minipage}[c][8cm]{13.5cm}		% Largura
% 	\small
% 	\imprimirautor
% 	%Sobrenome, Nome do autor
	
% 	\hspace{0.5cm} \imprimirtitulo  / \imprimirautor. --
% 	\imprimirlocal, \imprimirdata-
	
% 	\hspace{0.5cm} \thelastpage p. : il. (algumas color.) ; 30 cm.\\
	
% 	\hspace{0.5cm} \imprimirorientadorRotulo~\imprimirorientador\\
	
% 	\hspace{0.5cm}
% 	\parbox[t]{\textwidth}{\imprimirtipotrabalho~--~\imprimirinstituicao,
% 	\imprimirdata.}\\
	
% 	\hspace{0.5cm}
% 		1. Palavra-chave1.
% 		2. Palavra-chave2.
% 		2. Palavra-chave3.
% 		I. Orientador.
% 		II. Universidade xxx.
% 		III. Faculdade de xxx.
% 		IV. Título 			
% 	\end{minipage}}
% 	\end{center}
% \end{fichacatalografica}
% ---


% ---
% FOLHA DE APROVAÇÃO
% ---

% Isto é um exemplo de Folha de aprovação, elemento obrigatório da NBR
% 14724/2011 (seção 4.2.1.3). Você pode utilizar este modelo até a aprovação
% do trabalho. Após isso, substitua todo o conteúdo deste arquivo por uma
% imagem da página assinada pela banca com o comando abaixo:
%
% \begin{folhadeaprovacao}
% \includepdf{folhadeaprovacao_final.pdf}
% \end{folhadeaprovacao}
%
\setlength{\ABNTEXsignskip}{2cm}
\setlength{\ABNTEXsignwidth}{10cm}

\begin{folhadeaprovacao}
  \begin{center}
    {\ABNTEXchapterfont\large\imprimirautor}

    \vspace*{1cm}

    \begin{center}
      \ABNTEXchapterfont\bfseries\Large\imprimirtitulo
    \end{center}

    \vspace*{1cm}
    
    \hspace{.45\textwidth}
    \begin{minipage}{.5\textwidth}
        \imprimirpreambulo
    \end{minipage}%
   \end{center}

   \vspace*{1cm}
        
   \begin{center}
   Aprovado em 17 de julho de 2024.

   \vspace*{1cm}

   Banca examinadora:
   \end{center}

   \vspace*{-1cm}

   \assinatura{Prof. Ma. Caroline Santa Maria Rodrigues -- PUCRS}
   \assinatura{Prof. Dra. Clarissa Pinto Pizarro de Freitas -- PUCRS}
   \assinatura{\imprimirorientador -- PUCRS}
\end{folhadeaprovacao}
% ---


% ---
% Agradecimentos
% ---
\begin{agradecimentos}
Agradeço aos meus pais, Ademilson José Dias e Irinês Porfírio Dias, pelo apoio incondicional que me deram ao longo da vida,
incentivando mesmo as ideias questionáveis como aquela de fazer uma segunda graduação.

A minha irmã, Maria Eduarda Dias, cuja trajetória acadêmica é fonte de inspiração.

A minha namorada, Nathalia Fallavena Ceratti, que acompanhou a esta jornada de perto, oferecendo apoio e compreensão que
fizeram toda a diferença.

Aos meus amigos, Guilherme da Silva Petró, Norton Neves e Uilinton Luciano Marques da Silva, que foram compreensíveis nos
finais de semana em que não pude jogar RPG.

Ao professor Wagner de Lara Machado pela orientação e apoio no desenvolvimento deste trabalho.
\end{agradecimentos}
% ---


% ---
% RESUMOS
% ---
% Resumo em português
\setlength{\absparsep}{18pt} % ajusta o espaçamento dos parágrafos do resumo
\begin{resumo}
Desenvolvimentos na área de aprendizagem de máquina representam uma série de oportunidades para diversas áreas de pesquisa e de atuação profissional.
O uso de modelos de aprendizagem de máquina na predição de desfecho para tratamentos em saúde mental tem sido investigado e apresenta resultados promissores.
A aplicação desse tipo de modelo possibilitaria, por exemplo, maior assertividade na personalização de tratamentos psicoterápicos a nível individual.
Apesar dos benefícios potenciais, o tema ainda é pouco discutido nos espaços de pesquisa, formação e atuação em psicologia clínica. Este trabalho tem
por objetivo introduzir o tema de aprendizagem de máquina a psicólogos. Busca-se apresentar, de maneira acessível, os conceitos e terminologia básicos
em aprendizagem de máquina, proporcionando um panorama geral da área. Para complementar a exposição conceitual, é construído um exemplo a partir de dados
de um estudo sobre eficácia de intervenções em psicologia positiva para depressão via internet. Por fim, discutem-se algumas implicações do uso desse tipo
de tecnologia para a pesquisa e para a prática em psicoterapia.

 \textbf{Palavras-chave}: aprendizagem de máquina, psicoterapia, psicologia positiva, depressão.
\end{resumo}

% Resumo em inglês
\begin{resumo}[Abstract]
 \begin{otherlanguage*}{english}
   Advances in machine learning have opened up a multitude of opportunities for various research and professional fields. The use of machine learning models to
   predict treatment outcomes in mental health has been investigated showing promising results. The application of such models would allow greater assertiveness
   in tailoring psychotherapeutic treatments at the individual level. Despite its potential benefits, the topic remains unexplored in the fields of research, training
   and practice in clinical psychology. This work aims to introduce the topic of machine learning to psychologists. It presents the basic machine learning concepts and
   terminology in an accessible manner, providing a comprehensive overview of the field. To complement the conceptual discussion, an example is constructed using data
   from a study on the effectiveness of internet-based positive psychology interventions for depression. Finally, implications of using this type of technology in
   psychotherapy research and practice are discussed.
   \vspace{\onelineskip}
 
   \noindent 
   \textbf{Keywords}: machine learning, psychotherapy, positive psychology, depression.
 \end{otherlanguage*}
\end{resumo}


% ---
% LISTA DE ILUSTRAÇÕES
% ---
\pdfbookmark[0]{\listfigurename}{lof}
\listoffigures*
\cleardoublepage
% ---


% ---
% LISTA DE QUADROS
% ---
% \pdfbookmark[0]{\listofquadrosname}{loq}
% \listofquadros*
% \cleardoublepage
% ---


% ---
% LISTA DE TABELAS
% ---
% \pdfbookmark[0]{\listtablename}{lot}
% \listoftables*
% \cleardoublepage
% ---


% ---
% LISTA DE ABREVIATURAS E SIGLAS
% ---
% \begin{siglas}
%   \item[ABNT] Associação Brasileira de Normas Técnicas
%   \item[abnTeX] ABsurdas Normas para TeX
% \end{siglas}
% ---


% ---
% LISTA DE SÍMBOLOS
% ---
% \begin{simbolos}
%   \item[$ \Gamma $] Letra grega Gama
%   \item[$ \Lambda $] Lambda
%   \item[$ \zeta $] Letra grega minúscula zeta
%   \item[$ \in $] Pertence
% \end{simbolos}
% ---


% ---
% Sumário
% ---
\pdfbookmark[0]{\contentsname}{toc}
\tableofcontents*
\cleardoublepage
% ---


% ----------------------------------------------------------
% ELEMENTOS TEXTUAIS
% ----------------------------------------------------------
\textual

\include{./01-introducao/main.tex}

\include{./02-machine-learning/main.tex}

\include{./03-exemplo/main.tex}

\include{./04-conclusao/main.tex}

% ----------------------------------------------------------
% ELEMENTOS PÓS-TEXTUAIS
% ----------------------------------------------------------
\postextual
% ----------------------------------------------------------


% ----------------------------------------------------------
% REFERÊNCIAS
% ----------------------------------------------------------
\bibliography{references.bib}


% ----------------------------------------------------------
% GLOSSÁRIO
% ----------------------------------------------------------
%
% Consulte o manual da classe abntex2 para orientações sobre o glossário.
%
%\glossary

% ----------------------------------------------------------
% APÊNDICES
% ----------------------------------------------------------

% ---
% Inicia os apêndices
% ---
% \begin{apendicesenv}

% % Imprime uma página indicando o início dos apêndices
% \partapendices

% % ----------------------------------------------------------
% \chapter{Quisque libero justo}
% % ----------------------------------------------------------

% % ----------------------------------------------------------
% \chapter{Nullam elementum urna}
% % ----------------------------------------------------------

% \end{apendicesenv}
% ---


% ----------------------------------------------------------
% ANEXOS
% ----------------------------------------------------------

% ---
% Inicia os anexos
% ---
% \begin{anexosenv}

% % Imprime uma página indicando o início dos anexos
% \partanexos

% % ---
% \chapter{Morbi ultrices rutrum lorem}
% % ---

% % ---
% \chapter{Cras non urna sed feugiat cum}
% % ---


% % ---
% \chapter{Fusce facilisis lacinia dui}
% % ---

% \end{anexosenv}

%---------------------------------------------------------------------
% INDICE REMISSIVO
%---------------------------------------------------------------------
% \phantompart
% \printindex
%---------------------------------------------------------------------
\includepdf{./05-extras/contracapa.pdf}
\end{document}


\documentclass[
	% Opções da classe memoir
	12pt,				        % tamanho da fonte
	openany,			        % capítulos começam em pág ímpar (insere página vazia caso preciso)
	oneside,			        % para impressão em recto e verso. Oposto a oneside
	a4paper,			        % tamanho do papel. 
	% Opções da classe abntex2
	%chapter=TITLE,		        % títulos de capítulos convertidos em letras maiúsculas
	%section=TITLE,		        % títulos de seções convertidos em letras maiúsculas
	%subsection=TITLE,	        % títulos de subseções convertidos em letras maiúsculas
	%subsubsection=TITLE,       % títulos de subsubseções convertidos em letras maiúsculas
	% Opções do pacote babel
	english,			        % idioma adicional para hifenização
	brazil				        % o último idioma é o principal do documento
	]{abntex2}

% ---
% Pacotes básicos 
% ---
\usepackage{lmodern}			% Usa a fonte Latin Modern			
\usepackage[T1]{fontenc}		% Selecao de codigos de fonte.
\usepackage[utf8]{inputenc}		% Codificacao do documento (conversão automática dos acentos)
\usepackage{indentfirst}		% Indenta o primeiro parágrafo de cada seção.
\usepackage{color}				% Controle das cores
\usepackage{graphicx}			% Inclusão de gráficos
\usepackage{microtype} 			% para melhorias de justificação
% ---
		

% ---
% Pacotes de citações
% ---
% \usepackage[alf]{abntex2cite}						% Citações padrão ABNT
\usepackage[
  alf,
  abnt-nbr10520=1988,                     % Citações em sentence case
  abnt-emphasize=bf
  ]{abntex2cite}


% ---
% Pacotes adicionais
% ---
\usepackage{pdfpages}	% Inclusão de páginas em PDF


% ---
% Informações para CAPA e FOLHA DE ROSTO
% ---
\titulo{APRENDIZAGEM DE MÁQUINA PARA PREDIÇÃO DE DESFECHOS EM SAÚDE MENTAL: \textmd{UMA INTRODUÇÃO PARA PSICÓLOGOS}}
\autor{GUILHERME AUGUSTO DIAS}
\local{Porto Alegre}
\data{2024}
\orientador{Prof. Dr. Wagner de Lara Machado}
% \coorientador{}
% \instituicao{
%   Pontifícia Universidade Católica do Rio Grande do Sul
%   \par
%   Escola de Ciências da Saúde e da Vida
%   \par
%   Curso de Psicologia}
\tipotrabalho{TCC}
\preambulo{Trabalho de Conclusão de Curso apresentado como requisito para a obtenção do grau de Bacharel em Psicologia pela Escola de Ciências da Saúde e da Vida da Pontifícia Universidade Católica do Rio Grande do Sul.}
% ---

% ---
% Configurações de aparência do PDF final
% ---

% Alterando o aspecto da cor azul
\definecolor{blue}{RGB}{41,5,195}

% Informações do PDF
\makeatletter
\hypersetup{
		pdftitle={Aprendizagem de máquina para predição de desfechos em saúde mental: uma introdução para psicólogos}, 
		pdfauthor={Guilherme Augusto Dias},
    	pdfsubject={\imprimirpreambulo},
	    pdfcreator={LaTeX with abnTeX2},
		pdfkeywords={aprendizagem de máquina}{psicoterapia}{psicologia positiva}{depressão}, 
		colorlinks=true,		% true: colored links; false: boxed links
    	linkcolor=black,		% color of internal links
    	citecolor=black,		% color of links to bibliography
    	filecolor=black,		% color of file links
		urlcolor=black,
		bookmarksdepth=4
}
\makeatother
% --- 


% ---
% Posiciona figuras e tabelas no topo da página quando adicionadas sozinhas
% em um página em branco. Ver https://github.com/abntex/abntex2/issues/170
\makeatletter
\setlength{\@fptop}{5pt} % Set distance from top of page to first float
\makeatother
% ---

% ---
% Possibilita criação de Quadros e Lista de quadros.
% Ver https://github.com/abntex/abntex2/issues/176
%
\newcommand{\quadroname}{Quadro}
\newcommand{\listofquadrosname}{Lista de quadros}

\newfloat[chapter]{quadro}{loq}{\quadroname}
\newlistof{listofquadros}{loq}{\listofquadrosname}
\newlistentry{quadro}{loq}{0}

% Configurações para atender às regras da ABNT
\setfloatadjustment{quadro}{\centering}
\counterwithout{quadro}{chapter}
\renewcommand{\cftquadroname}{\quadroname\space} 
\renewcommand*{\cftquadroaftersnum}{\hfill--\hfill}

\setfloatlocations{quadro}{hbtp} % Ver https://github.com/abntex/abntex2/issues/176
% ---


% --- 
% Espaçamentos entre linhas e parágrafos 
% --- 

% O tamanho do parágrafo é dado por:
\setlength{\parindent}{1.3cm}

% Controle do espaçamento entre um parágrafo e outro:
\setlength{\parskip}{0.2cm}  % tente também \onelineskip

% ---
% Compila o índice
% ---
\makeindex
% ---


% ----
% Início do documento
% ----
\begin{document}

% Seleciona o idioma do documento (conforme pacotes do babel)
\selectlanguage{brazil}

% Retira espaço extra obsoleto entre as frases.
\frenchspacing 

% ----------------------------------------------------------
% ELEMENTOS PRÉ-TEXTUAIS
% ----------------------------------------------------------
\pretextual

% ---
% CAPA
% ---
\includepdf{./05-extras/capa.pdf}
% ---


% ---
% FOLHA DE ROSTO
% (o * indica que haverá a ficha bibliográfica)
% ---
\imprimirfolhaderosto
% ---


% ---
% Ficha bibliografica
% ---

% Isto é um exemplo de Ficha Catalográfica, ou ``Dados internacionais de
% catalogação-na-publicação''. Você pode utilizar este modelo como referência. 
% Porém, provavelmente a biblioteca da sua universidade lhe fornecerá um PDF
% com a ficha catalográfica definitiva após a defesa do trabalho. Quando estiver
% com o documento, salve-o como PDF no diretório do seu projeto e substitua todo
% o conteúdo de implementação deste arquivo pelo comando abaixo:
%
% \usepackage{pdfpages}		% necessário para comando \includepdf
% \begin{fichacatalografica}
%     \includepdf{fig_ficha_catalografica.pdf}
% \end{fichacatalografica}

% \begin{fichacatalografica}
% 	\sffamily
% 	\vspace*{\fill}                             % Posição vertical
% 	\begin{center}                              % Minipage Centralizado
% 	\fbox{\begin{minipage}[c][8cm]{13.5cm}		% Largura
% 	\small
% 	\imprimirautor
% 	%Sobrenome, Nome do autor
	
% 	\hspace{0.5cm} \imprimirtitulo  / \imprimirautor. --
% 	\imprimirlocal, \imprimirdata-
	
% 	\hspace{0.5cm} \thelastpage p. : il. (algumas color.) ; 30 cm.\\
	
% 	\hspace{0.5cm} \imprimirorientadorRotulo~\imprimirorientador\\
	
% 	\hspace{0.5cm}
% 	\parbox[t]{\textwidth}{\imprimirtipotrabalho~--~\imprimirinstituicao,
% 	\imprimirdata.}\\
	
% 	\hspace{0.5cm}
% 		1. Palavra-chave1.
% 		2. Palavra-chave2.
% 		2. Palavra-chave3.
% 		I. Orientador.
% 		II. Universidade xxx.
% 		III. Faculdade de xxx.
% 		IV. Título 			
% 	\end{minipage}}
% 	\end{center}
% \end{fichacatalografica}
% ---


% ---
% FOLHA DE APROVAÇÃO
% ---

% Isto é um exemplo de Folha de aprovação, elemento obrigatório da NBR
% 14724/2011 (seção 4.2.1.3). Você pode utilizar este modelo até a aprovação
% do trabalho. Após isso, substitua todo o conteúdo deste arquivo por uma
% imagem da página assinada pela banca com o comando abaixo:
%
% \begin{folhadeaprovacao}
% \includepdf{folhadeaprovacao_final.pdf}
% \end{folhadeaprovacao}
%
\setlength{\ABNTEXsignskip}{2cm}
\setlength{\ABNTEXsignwidth}{10cm}

\begin{folhadeaprovacao}
  \begin{center}
    {\ABNTEXchapterfont\large\imprimirautor}

    \vspace*{1cm}

    \begin{center}
      \ABNTEXchapterfont\bfseries\Large\imprimirtitulo
    \end{center}

    \vspace*{1cm}
    
    \hspace{.45\textwidth}
    \begin{minipage}{.5\textwidth}
        \imprimirpreambulo
    \end{minipage}%
   \end{center}

   \vspace*{1cm}
        
   \begin{center}
   Aprovado em 17 de julho de 2024.

   \vspace*{1cm}

   Banca examinadora:
   \end{center}

   \vspace*{-1cm}

   \assinatura{Prof. Ma. Caroline Santa Maria Rodrigues -- PUCRS}
   \assinatura{Prof. Dra. Clarissa Pinto Pizarro de Freitas -- PUCRS}
   \assinatura{\imprimirorientador -- PUCRS}
\end{folhadeaprovacao}
% ---


% ---
% Agradecimentos
% ---
\begin{agradecimentos}
Agradeço aos meus pais, Ademilson José Dias e Irinês Porfírio Dias, pelo apoio incondicional que me deram ao longo da vida,
incentivando mesmo as ideias questionáveis como aquela de fazer uma segunda graduação.

A minha irmã, Maria Eduarda Dias, cuja trajetória acadêmica é fonte de inspiração.

A minha namorada, Nathalia Fallavena Ceratti, que acompanhou a esta jornada de perto, oferecendo apoio e compreensão que
fizeram toda a diferença.

Aos meus amigos, Guilherme da Silva Petró, Norton Neves e Uilinton Luciano Marques da Silva, que foram compreensíveis nos
finais de semana em que não pude jogar RPG.

Ao professor Wagner de Lara Machado pela orientação e apoio no desenvolvimento deste trabalho.
\end{agradecimentos}
% ---


% ---
% RESUMOS
% ---
% Resumo em português
\setlength{\absparsep}{18pt} % ajusta o espaçamento dos parágrafos do resumo
\begin{resumo}
Desenvolvimentos na área de aprendizagem de máquina representam uma série de oportunidades para diversas áreas de pesquisa e de atuação profissional.
O uso de modelos de aprendizagem de máquina na predição de desfecho para tratamentos em saúde mental tem sido investigado e apresenta resultados promissores.
A aplicação desse tipo de modelo possibilitaria, por exemplo, maior assertividade na personalização de tratamentos psicoterápicos a nível individual.
Apesar dos benefícios potenciais, o tema ainda é pouco discutido nos espaços de pesquisa, formação e atuação em psicologia clínica. Este trabalho tem
por objetivo introduzir o tema de aprendizagem de máquina a psicólogos. Busca-se apresentar, de maneira acessível, os conceitos e terminologia básicos
em aprendizagem de máquina, proporcionando um panorama geral da área. Para complementar a exposição conceitual, é construído um exemplo a partir de dados
de um estudo sobre eficácia de intervenções em psicologia positiva para depressão via internet. Por fim, discutem-se algumas implicações do uso desse tipo
de tecnologia para a pesquisa e para a prática em psicoterapia.

 \textbf{Palavras-chave}: aprendizagem de máquina, psicoterapia, psicologia positiva, depressão.
\end{resumo}

% Resumo em inglês
\begin{resumo}[Abstract]
 \begin{otherlanguage*}{english}
   Advances in machine learning have opened up a multitude of opportunities for various research and professional fields. The use of machine learning models to
   predict treatment outcomes in mental health has been investigated showing promising results. The application of such models would allow greater assertiveness
   in tailoring psychotherapeutic treatments at the individual level. Despite its potential benefits, the topic remains unexplored in the fields of research, training
   and practice in clinical psychology. This work aims to introduce the topic of machine learning to psychologists. It presents the basic machine learning concepts and
   terminology in an accessible manner, providing a comprehensive overview of the field. To complement the conceptual discussion, an example is constructed using data
   from a study on the effectiveness of internet-based positive psychology interventions for depression. Finally, implications of using this type of technology in
   psychotherapy research and practice are discussed.
   \vspace{\onelineskip}
 
   \noindent 
   \textbf{Keywords}: machine learning, psychotherapy, positive psychology, depression.
 \end{otherlanguage*}
\end{resumo}


% ---
% LISTA DE ILUSTRAÇÕES
% ---
\pdfbookmark[0]{\listfigurename}{lof}
\listoffigures*
\cleardoublepage
% ---


% ---
% LISTA DE QUADROS
% ---
% \pdfbookmark[0]{\listofquadrosname}{loq}
% \listofquadros*
% \cleardoublepage
% ---


% ---
% LISTA DE TABELAS
% ---
% \pdfbookmark[0]{\listtablename}{lot}
% \listoftables*
% \cleardoublepage
% ---


% ---
% LISTA DE ABREVIATURAS E SIGLAS
% ---
% \begin{siglas}
%   \item[ABNT] Associação Brasileira de Normas Técnicas
%   \item[abnTeX] ABsurdas Normas para TeX
% \end{siglas}
% ---


% ---
% LISTA DE SÍMBOLOS
% ---
% \begin{simbolos}
%   \item[$ \Gamma $] Letra grega Gama
%   \item[$ \Lambda $] Lambda
%   \item[$ \zeta $] Letra grega minúscula zeta
%   \item[$ \in $] Pertence
% \end{simbolos}
% ---


% ---
% Sumário
% ---
\pdfbookmark[0]{\contentsname}{toc}
\tableofcontents*
\cleardoublepage
% ---


% ----------------------------------------------------------
% ELEMENTOS TEXTUAIS
% ----------------------------------------------------------
\textual

\include{./01-introducao/main.tex}

\include{./02-machine-learning/main.tex}

\include{./03-exemplo/main.tex}

\include{./04-conclusao/main.tex}

% ----------------------------------------------------------
% ELEMENTOS PÓS-TEXTUAIS
% ----------------------------------------------------------
\postextual
% ----------------------------------------------------------


% ----------------------------------------------------------
% REFERÊNCIAS
% ----------------------------------------------------------
\bibliography{references.bib}


% ----------------------------------------------------------
% GLOSSÁRIO
% ----------------------------------------------------------
%
% Consulte o manual da classe abntex2 para orientações sobre o glossário.
%
%\glossary

% ----------------------------------------------------------
% APÊNDICES
% ----------------------------------------------------------

% ---
% Inicia os apêndices
% ---
% \begin{apendicesenv}

% % Imprime uma página indicando o início dos apêndices
% \partapendices

% % ----------------------------------------------------------
% \chapter{Quisque libero justo}
% % ----------------------------------------------------------

% % ----------------------------------------------------------
% \chapter{Nullam elementum urna}
% % ----------------------------------------------------------

% \end{apendicesenv}
% ---


% ----------------------------------------------------------
% ANEXOS
% ----------------------------------------------------------

% ---
% Inicia os anexos
% ---
% \begin{anexosenv}

% % Imprime uma página indicando o início dos anexos
% \partanexos

% % ---
% \chapter{Morbi ultrices rutrum lorem}
% % ---

% % ---
% \chapter{Cras non urna sed feugiat cum}
% % ---


% % ---
% \chapter{Fusce facilisis lacinia dui}
% % ---

% \end{anexosenv}

%---------------------------------------------------------------------
% INDICE REMISSIVO
%---------------------------------------------------------------------
% \phantompart
% \printindex
%---------------------------------------------------------------------
\includepdf{./05-extras/contracapa.pdf}
\end{document}


% ----------------------------------------------------------
% ELEMENTOS PÓS-TEXTUAIS
% ----------------------------------------------------------
\postextual
% ----------------------------------------------------------


% ----------------------------------------------------------
% REFERÊNCIAS
% ----------------------------------------------------------
\bibliography{references.bib}


% ----------------------------------------------------------
% GLOSSÁRIO
% ----------------------------------------------------------
%
% Consulte o manual da classe abntex2 para orientações sobre o glossário.
%
%\glossary

% ----------------------------------------------------------
% APÊNDICES
% ----------------------------------------------------------

% ---
% Inicia os apêndices
% ---
% \begin{apendicesenv}

% % Imprime uma página indicando o início dos apêndices
% \partapendices

% % ----------------------------------------------------------
% \chapter{Quisque libero justo}
% % ----------------------------------------------------------

% % ----------------------------------------------------------
% \chapter{Nullam elementum urna}
% % ----------------------------------------------------------

% \end{apendicesenv}
% ---


% ----------------------------------------------------------
% ANEXOS
% ----------------------------------------------------------

% ---
% Inicia os anexos
% ---
% \begin{anexosenv}

% % Imprime uma página indicando o início dos anexos
% \partanexos

% % ---
% \chapter{Morbi ultrices rutrum lorem}
% % ---

% % ---
% \chapter{Cras non urna sed feugiat cum}
% % ---


% % ---
% \chapter{Fusce facilisis lacinia dui}
% % ---

% \end{anexosenv}

%---------------------------------------------------------------------
% INDICE REMISSIVO
%---------------------------------------------------------------------
% \phantompart
% \printindex
%---------------------------------------------------------------------
\includepdf{./05-extras/contracapa.pdf}
\end{document}


% ----------------------------------------------------------
% ELEMENTOS PÓS-TEXTUAIS
% ----------------------------------------------------------
\postextual
% ----------------------------------------------------------


% ----------------------------------------------------------
% REFERÊNCIAS
% ----------------------------------------------------------
\bibliography{references.bib}


% ----------------------------------------------------------
% GLOSSÁRIO
% ----------------------------------------------------------
%
% Consulte o manual da classe abntex2 para orientações sobre o glossário.
%
%\glossary

% ----------------------------------------------------------
% APÊNDICES
% ----------------------------------------------------------

% ---
% Inicia os apêndices
% ---
% \begin{apendicesenv}

% % Imprime uma página indicando o início dos apêndices
% \partapendices

% % ----------------------------------------------------------
% \chapter{Quisque libero justo}
% % ----------------------------------------------------------

% % ----------------------------------------------------------
% \chapter{Nullam elementum urna}
% % ----------------------------------------------------------

% \end{apendicesenv}
% ---


% ----------------------------------------------------------
% ANEXOS
% ----------------------------------------------------------

% ---
% Inicia os anexos
% ---
% \begin{anexosenv}

% % Imprime uma página indicando o início dos anexos
% \partanexos

% % ---
% \chapter{Morbi ultrices rutrum lorem}
% % ---

% % ---
% \chapter{Cras non urna sed feugiat cum}
% % ---


% % ---
% \chapter{Fusce facilisis lacinia dui}
% % ---

% \end{anexosenv}

%---------------------------------------------------------------------
% INDICE REMISSIVO
%---------------------------------------------------------------------
% \phantompart
% \printindex
%---------------------------------------------------------------------
\includepdf{./05-extras/contracapa.pdf}
\end{document}


% ----------------------------------------------------------
% ELEMENTOS PÓS-TEXTUAIS
% ----------------------------------------------------------
\postextual
% ----------------------------------------------------------


% ----------------------------------------------------------
% REFERÊNCIAS
% ----------------------------------------------------------
\bibliography{references.bib}


% ----------------------------------------------------------
% GLOSSÁRIO
% ----------------------------------------------------------
%
% Consulte o manual da classe abntex2 para orientações sobre o glossário.
%
%\glossary

% ----------------------------------------------------------
% APÊNDICES
% ----------------------------------------------------------

% ---
% Inicia os apêndices
% ---
% \begin{apendicesenv}

% % Imprime uma página indicando o início dos apêndices
% \partapendices

% % ----------------------------------------------------------
% \chapter{Quisque libero justo}
% % ----------------------------------------------------------

% % ----------------------------------------------------------
% \chapter{Nullam elementum urna}
% % ----------------------------------------------------------

% \end{apendicesenv}
% ---


% ----------------------------------------------------------
% ANEXOS
% ----------------------------------------------------------

% ---
% Inicia os anexos
% ---
% \begin{anexosenv}

% % Imprime uma página indicando o início dos anexos
% \partanexos

% % ---
% \chapter{Morbi ultrices rutrum lorem}
% % ---

% % ---
% \chapter{Cras non urna sed feugiat cum}
% % ---


% % ---
% \chapter{Fusce facilisis lacinia dui}
% % ---

% \end{anexosenv}

%---------------------------------------------------------------------
% INDICE REMISSIVO
%---------------------------------------------------------------------
% \phantompart
% \printindex
%---------------------------------------------------------------------
\includepdf{./05-extras/contracapa.pdf}
\end{document}
